\documentclass[10pt]{article}
\usepackage[margin=0.65in]{geometry}
\usepackage[T1]{fontenc}
\usepackage{hyperref}
\usepackage{xcolor}
\usepackage{listings}
\lstset{basicstyle=\ttfamily\scriptsize,breaklines=true,columns=fullflexible,keepspaces=true,showstringspaces=false}
\title{NeuraEdge CIFAR-10 Neuromorphic Accelerator: Full Pipeline and Line-by-Line Explainer}
\author{Generated from repository source and measured Vivado/XSim artifacts}
\date{\today}
\begin{document}
\maketitle
\section{What Was Built}
This document describes the implemented CIFAR-10 accelerator path. The trained Spikingformer backbone is used to export 384-dimensional fixed-point features and a trained classifier head. The FPGA accelerator stores one feature vector, applies a measured low-magnitude feature activity gate, computes ten fixed-point class scores with an on-chip weight memory and one DSP MAC datapath, returns the predicted label, and includes a margin-gated on-chip learning path that updates the true-label and predicted-label classifier rows in FPGA memory.
\section{Measured Verification Summary}
\begin{itemize}
\item Full inference XSim run: 10,000 CIFAR-10 samples, 95.51\% top-1 accuracy, exact score match against the fixed-point Python reference.
\item Feature activity gate: abs(feature) $\leq$ 2048 is treated as zero before inference and learning; the exported test set skips 57.81\% of feature-MAC opportunities.
\item Full on-chip-learning XSim run: 10,000 samples, 25 RTL weight/bias updates, zero score mismatches, zero prediction mismatches, zero update-flag mismatches, and zero update-count mismatches against the independent Python online-learning golden model.
\item Vivado implementation target: PYNQ Zynq-7020, xc7z020clg400-1, 100 MHz. Routed timing met with WNS 0.336 ns and WHS 0.100 ns.
\item Resource result: 1043 LUTs, 988 FFs, 6 RAMB36E1 blocks, and 1 DSP48E1.
\item Vectorless Vivado power: 0.123 W total, 0.018 W dynamic, 0.105 W static.
\end{itemize}
\section{Research Context And Comparison}
The model side follows the Spikingformer family of spike-driven transformer SNNs, whose reported CIFAR-10 accuracy is about 95.81--95.95\% in the cited paper/repository. The local fixed-point exported checkpoint reaches 95.51\% after mapping into this routed hardware head with the measured feature gate. Recent SNN-transformer targets to beat include Max-Former at 97.04\% on CIFAR-10 and alpha-XNOR spiking self-attention at 96.15\% for the reported multi-stage ViT setting. The implemented hardware contribution here is the verified FPGA path: local feature and weight memories, exact fixed-point RTL/golden agreement, activity-gated feature silence, and on-chip classifier learning.
\begin{itemize}
\item Spikingformer paper/repository: \url{https://arxiv.org/html/2304.11954v5}, \url{https://github.com/TheBrainLab/Spikingformer}.
\item Max-Former high-frequency spiking transformer reference: \url{https://arxiv.org/html/2505.18608v4}.
\item Alpha-XNOR spiking self-attention reference: \url{https://openaccess.thecvf.com/content/CVPR2025/papers/Xiao_Rethinking_Spiking_Self-Attention_Mechanism_Implementing_a-XNOR_Similarity_Calculation_in_Spiking_CVPR_2025_paper.pdf}.
\item TrueNorth reference: \url{https://research.ibm.com/publications/truenorth-design-and-tool-flow-of-a-65-mw-1-million-neuron-programmable-neurosynaptic-chip}.
\item Loihi on-chip learning reference: \url{https://redwood.berkeley.edu/wp-content/uploads/2021/08/Davies2018.pdf}.
\end{itemize}
\section{End-to-End Dataflow}
\begin{enumerate}
\item Python loads the trained Spikingformer checkpoint and CIFAR-10 data, runs the trained network, and exports fixed-point feature, label, weight, and bias memories.
\item XSim streams each 384-word feature vector into the RTL through the same handshake the PYNQ AXI-stream wrapper uses.
\item The RTL stores the feature vector locally, zeroes features whose magnitude is below the selected gate threshold, loops over the ten classes, reads each class row from BRAM, multiplies effective feature and weight in a DSP, accumulates a 64-bit score, and chooses argmax.
\item If training is enabled and the prediction is wrong but close enough by margin, the RTL latches the label, reads the true and predicted class rows, computes signed shifted deltas from the same effective feature path, saturates int16 weight updates, writes updated weights, and updates the two class biases.
\item Python validators independently replay the same fixed-point inference and online-learning rule from the exported memories, then compare every RTL score, prediction, update flag, and update count.
\item Vivado synthesizes and routes the PYNQ AXI-stream top, producing timing, utilization, power, and checkpoint artifacts.
\end{enumerate}
\section{Line-by-Line Source Explanation}
\clearpage
\subsection{cifar10\_neuromorphic\textbackslash{}rtl\textbackslash{}cifar10\_spikingformer\_head\_accelerator.sv}
Each entry below gives the exact source line followed by its hardware/software role in the implemented accelerator pipeline.
\paragraph{Line 1}
\noindent\colorbox{gray!8}{\parbox{\dimexpr\linewidth-2\fboxsep\relax}{\raggedright\ttfamily\scriptsize `timescale 1ns / 1ps}}
\par\smallskip
This simulator directive sets the time unit and precision used by testbenches and waveform timestamps. It keeps delays such as clock half-periods numerically meaningful in XSim.
\paragraph{Line 2}
\noindent\colorbox{gray!8}{\parbox{\dimexpr\linewidth-2\fboxsep\relax}{\raggedright\ttfamily\scriptsize `default\_nettype none}}
\par\smallskip
This disables implicit one-bit wires. It prevents spelling mistakes from silently creating unwanted hardware nets.
\paragraph{Line 3}
\noindent\colorbox{gray!8}{\parbox{\dimexpr\linewidth-2\fboxsep\relax}{\raggedright\ttfamily\scriptsize `include "rtl/cifar10\_spikingformer\_head\_params.svh"}}
\par\smallskip
This includes shared compile-time constants such as feature dimensions and fixed-point widths. The RTL and testbenches therefore agree on the same trained-model shape.
\paragraph{Line 4}
\noindent\colorbox{gray!8}{\parbox{\dimexpr\linewidth-2\fboxsep\relax}{\raggedright\ttfamily\scriptsize `ifndef CIFAR10\_SF\_HEAD\_INFER\_SKIP\_ABS\_THRESH}}
\par\smallskip
This line controls the feature-activity gate. When the absolute fixed-point feature is at or below this threshold, the RTL treats it as zero so the exported Spikingformer head ignores small noisy activations and records a measurable MAC-skip opportunity.
\paragraph{Line 5}
\noindent\colorbox{gray!8}{\parbox{\dimexpr\linewidth-2\fboxsep\relax}{\raggedright\ttfamily\scriptsize `define CIFAR10\_SF\_HEAD\_INFER\_SKIP\_ABS\_THRESH 2048}}
\par\smallskip
This line controls the feature-activity gate. When the absolute fixed-point feature is at or below this threshold, the RTL treats it as zero so the exported Spikingformer head ignores small noisy activations and records a measurable MAC-skip opportunity.
\paragraph{Line 6}
\noindent\colorbox{gray!8}{\parbox{\dimexpr\linewidth-2\fboxsep\relax}{\raggedright\ttfamily\scriptsize `endif}}
\par\smallskip
This SystemVerilog line contributes to the structural or sequential description of the accelerator. Its effect is resolved by XSim for verification and by Vivado during synthesis.
\paragraph{Line 7}
\noindent\colorbox{gray!8}{\parbox{\dimexpr\linewidth-2\fboxsep\relax}{\raggedright\ttfamily\scriptsize \emph{blank line}}}
\par\smallskip
This blank line is not hardware or software behavior. It separates nearby declarations or operations so the file remains readable and reviewable.
\paragraph{Line 8}
\noindent\colorbox{gray!8}{\parbox{\dimexpr\linewidth-2\fboxsep\relax}{\raggedright\ttfamily\scriptsize module cifar10\_spikingformer\_head\_accelerator \#(}}
\par\smallskip
This begins a hardware module, which is the FPGA equivalent of a reusable circuit block. Everything until the matching endmodule describes ports, storage, combinational equations, and clocked behavior for this block.
\paragraph{Line 9}
\noindent\colorbox{gray!8}{\parbox{\dimexpr\linewidth-2\fboxsep\relax}{\raggedright\ttfamily\scriptsize     parameter int FEATURE\_DIM = `CIFAR10\_SF\_HEAD\_FEATURE\_DIM,}}
\par\smallskip
This is a compile-time hardware setting. Vivado resolves it before synthesis, so it controls the generated circuit size, fixed-point width, or learning rule without runtime overhead.
\paragraph{Line 10}
\noindent\colorbox{gray!8}{\parbox{\dimexpr\linewidth-2\fboxsep\relax}{\raggedright\ttfamily\scriptsize     parameter int NUM\_CLASSES = `CIFAR10\_SF\_HEAD\_NUM\_CLASSES,}}
\par\smallskip
This is a compile-time hardware setting. Vivado resolves it before synthesis, so it controls the generated circuit size, fixed-point width, or learning rule without runtime overhead.
\paragraph{Line 11}
\noindent\colorbox{gray!8}{\parbox{\dimexpr\linewidth-2\fboxsep\relax}{\raggedright\ttfamily\scriptsize     parameter int FEATURE\_BITS = `CIFAR10\_SF\_HEAD\_FEATURE\_BITS,}}
\par\smallskip
This is a compile-time hardware setting. Vivado resolves it before synthesis, so it controls the generated circuit size, fixed-point width, or learning rule without runtime overhead.
\paragraph{Line 12}
\noindent\colorbox{gray!8}{\parbox{\dimexpr\linewidth-2\fboxsep\relax}{\raggedright\ttfamily\scriptsize     parameter int WEIGHT\_BITS = `CIFAR10\_SF\_HEAD\_WEIGHT\_BITS,}}
\par\smallskip
This is a compile-time hardware setting. Vivado resolves it before synthesis, so it controls the generated circuit size, fixed-point width, or learning rule without runtime overhead.
\paragraph{Line 13}
\noindent\colorbox{gray!8}{\parbox{\dimexpr\linewidth-2\fboxsep\relax}{\raggedright\ttfamily\scriptsize     parameter int ACC\_BITS = `CIFAR10\_SF\_HEAD\_ACC\_BITS,}}
\par\smallskip
This is a compile-time hardware setting. Vivado resolves it before synthesis, so it controls the generated circuit size, fixed-point width, or learning rule without runtime overhead.
\paragraph{Line 14}
\noindent\colorbox{gray!8}{\parbox{\dimexpr\linewidth-2\fboxsep\relax}{\raggedright\ttfamily\scriptsize     parameter int ENABLE\_ONCHIP\_LEARNING = 1,}}
\par\smallskip
This is a compile-time hardware setting. Vivado resolves it before synthesis, so it controls the generated circuit size, fixed-point width, or learning rule without runtime overhead.
\paragraph{Line 15}
\noindent\colorbox{gray!8}{\parbox{\dimexpr\linewidth-2\fboxsep\relax}{\raggedright\ttfamily\scriptsize     parameter int INFER\_SKIP\_ABS\_THRESH = `CIFAR10\_SF\_HEAD\_INFER\_SKIP\_ABS\_THRESH,}}
\par\smallskip
This line controls the feature-activity gate. When the absolute fixed-point feature is at or below this threshold, the RTL treats it as zero so the exported Spikingformer head ignores small noisy activations and records a measurable MAC-skip opportunity.
\paragraph{Line 16}
\noindent\colorbox{gray!8}{\parbox{\dimexpr\linewidth-2\fboxsep\relax}{\raggedright\ttfamily\scriptsize     parameter int LEARN\_SHIFT = 14,}}
\par\smallskip
This is a compile-time hardware setting. Vivado resolves it before synthesis, so it controls the generated circuit size, fixed-point width, or learning rule without runtime overhead.
\paragraph{Line 17}
\noindent\colorbox{gray!8}{\parbox{\dimexpr\linewidth-2\fboxsep\relax}{\raggedright\ttfamily\scriptsize     parameter int BIAS\_LEARN\_DELTA = 1,}}
\par\smallskip
This is a compile-time hardware setting. Vivado resolves it before synthesis, so it controls the generated circuit size, fixed-point width, or learning rule without runtime overhead.
\paragraph{Line 18}
\noindent\colorbox{gray!8}{\parbox{\dimexpr\linewidth-2\fboxsep\relax}{\raggedright\ttfamily\scriptsize     parameter logic signed [ACC\_BITS-1:0] LEARN\_MARGIN\_Q = 64'sd268435456}}
\par\smallskip
This is a compile-time hardware setting. Vivado resolves it before synthesis, so it controls the generated circuit size, fixed-point width, or learning rule without runtime overhead.
\paragraph{Line 19}
\noindent\colorbox{gray!8}{\parbox{\dimexpr\linewidth-2\fboxsep\relax}{\raggedright\ttfamily\scriptsize ) (}}
\par\smallskip
This SystemVerilog line contributes to the structural or sequential description of the accelerator. Its effect is resolved by XSim for verification and by Vivado during synthesis.
\paragraph{Line 20}
\noindent\colorbox{gray!8}{\parbox{\dimexpr\linewidth-2\fboxsep\relax}{\raggedright\ttfamily\scriptsize     input  wire clk,}}
\par\smallskip
This line declares a module port. A port is a physical or wrapper-visible signal used to move clocks, control pulses, features, labels, predictions, or debug information across module boundaries.
\paragraph{Line 21}
\noindent\colorbox{gray!8}{\parbox{\dimexpr\linewidth-2\fboxsep\relax}{\raggedright\ttfamily\scriptsize     input  wire rst\_n,}}
\par\smallskip
This line declares a module port. A port is a physical or wrapper-visible signal used to move clocks, control pulses, features, labels, predictions, or debug information across module boundaries.
\paragraph{Line 22}
\noindent\colorbox{gray!8}{\parbox{\dimexpr\linewidth-2\fboxsep\relax}{\raggedright\ttfamily\scriptsize \emph{blank line}}}
\par\smallskip
This blank line is not hardware or software behavior. It separates nearby declarations or operations so the file remains readable and reviewable.
\paragraph{Line 23}
\noindent\colorbox{gray!8}{\parbox{\dimexpr\linewidth-2\fboxsep\relax}{\raggedright\ttfamily\scriptsize     input  wire start,}}
\par\smallskip
This line declares a module port. A port is a physical or wrapper-visible signal used to move clocks, control pulses, features, labels, predictions, or debug information across module boundaries.
\paragraph{Line 24}
\noindent\colorbox{gray!8}{\parbox{\dimexpr\linewidth-2\fboxsep\relax}{\raggedright\ttfamily\scriptsize     input  wire feature\_valid,}}
\par\smallskip
This line declares a module port. A port is a physical or wrapper-visible signal used to move clocks, control pulses, features, labels, predictions, or debug information across module boundaries.
\paragraph{Line 25}
\noindent\colorbox{gray!8}{\parbox{\dimexpr\linewidth-2\fboxsep\relax}{\raggedright\ttfamily\scriptsize     output logic feature\_ready,}}
\par\smallskip
This line declares a module port. A port is a physical or wrapper-visible signal used to move clocks, control pulses, features, labels, predictions, or debug information across module boundaries.
\paragraph{Line 26}
\noindent\colorbox{gray!8}{\parbox{\dimexpr\linewidth-2\fboxsep\relax}{\raggedright\ttfamily\scriptsize     input  wire signed [FEATURE\_BITS-1:0] feature\_value,}}
\par\smallskip
This line declares a module port. A port is a physical or wrapper-visible signal used to move clocks, control pulses, features, labels, predictions, or debug information across module boundaries.
\paragraph{Line 27}
\noindent\colorbox{gray!8}{\parbox{\dimexpr\linewidth-2\fboxsep\relax}{\raggedright\ttfamily\scriptsize     input  wire frame\_done,}}
\par\smallskip
This line declares a module port. A port is a physical or wrapper-visible signal used to move clocks, control pulses, features, labels, predictions, or debug information across module boundaries.
\paragraph{Line 28}
\noindent\colorbox{gray!8}{\parbox{\dimexpr\linewidth-2\fboxsep\relax}{\raggedright\ttfamily\scriptsize     input  wire train\_enable,}}
\par\smallskip
This line declares a module port. A port is a physical or wrapper-visible signal used to move clocks, control pulses, features, labels, predictions, or debug information across module boundaries.
\paragraph{Line 29}
\noindent\colorbox{gray!8}{\parbox{\dimexpr\linewidth-2\fboxsep\relax}{\raggedright\ttfamily\scriptsize     input  wire [3:0] train\_label,}}
\par\smallskip
This line declares a module port. A port is a physical or wrapper-visible signal used to move clocks, control pulses, features, labels, predictions, or debug information across module boundaries.
\paragraph{Line 30}
\noindent\colorbox{gray!8}{\parbox{\dimexpr\linewidth-2\fboxsep\relax}{\raggedright\ttfamily\scriptsize \emph{blank line}}}
\par\smallskip
This blank line is not hardware or software behavior. It separates nearby declarations or operations so the file remains readable and reviewable.
\paragraph{Line 31}
\noindent\colorbox{gray!8}{\parbox{\dimexpr\linewidth-2\fboxsep\relax}{\raggedright\ttfamily\scriptsize     output logic busy,}}
\par\smallskip
This line declares a module port. A port is a physical or wrapper-visible signal used to move clocks, control pulses, features, labels, predictions, or debug information across module boundaries.
\paragraph{Line 32}
\noindent\colorbox{gray!8}{\parbox{\dimexpr\linewidth-2\fboxsep\relax}{\raggedright\ttfamily\scriptsize     output logic result\_valid,}}
\par\smallskip
This line declares a module port. A port is a physical or wrapper-visible signal used to move clocks, control pulses, features, labels, predictions, or debug information across module boundaries.
\paragraph{Line 33}
\noindent\colorbox{gray!8}{\parbox{\dimexpr\linewidth-2\fboxsep\relax}{\raggedright\ttfamily\scriptsize     output logic [3:0] predicted\_label,}}
\par\smallskip
This line declares a module port. A port is a physical or wrapper-visible signal used to move clocks, control pulses, features, labels, predictions, or debug information across module boundaries.
\paragraph{Line 34}
\noindent\colorbox{gray!8}{\parbox{\dimexpr\linewidth-2\fboxsep\relax}{\raggedright\ttfamily\scriptsize     output logic signed [ACC\_BITS-1:0] predicted\_score,}}
\par\smallskip
This line declares a module port. A port is a physical or wrapper-visible signal used to move clocks, control pulses, features, labels, predictions, or debug information across module boundaries.
\paragraph{Line 35}
\noindent\colorbox{gray!8}{\parbox{\dimexpr\linewidth-2\fboxsep\relax}{\raggedright\ttfamily\scriptsize     output logic train\_update\_valid,}}
\par\smallskip
This line declares a module port. A port is a physical or wrapper-visible signal used to move clocks, control pulses, features, labels, predictions, or debug information across module boundaries.
\paragraph{Line 36}
\noindent\colorbox{gray!8}{\parbox{\dimexpr\linewidth-2\fboxsep\relax}{\raggedright\ttfamily\scriptsize     output logic train\_update\_applied,}}
\par\smallskip
This line declares a module port. A port is a physical or wrapper-visible signal used to move clocks, control pulses, features, labels, predictions, or debug information across module boundaries.
\paragraph{Line 37}
\noindent\colorbox{gray!8}{\parbox{\dimexpr\linewidth-2\fboxsep\relax}{\raggedright\ttfamily\scriptsize \emph{blank line}}}
\par\smallskip
This blank line is not hardware or software behavior. It separates nearby declarations or operations so the file remains readable and reviewable.
\paragraph{Line 38}
\noindent\colorbox{gray!8}{\parbox{\dimexpr\linewidth-2\fboxsep\relax}{\raggedright\ttfamily\scriptsize     output logic [5:0] debug\_state,}}
\par\smallskip
This line declares a module port. A port is a physical or wrapper-visible signal used to move clocks, control pulses, features, labels, predictions, or debug information across module boundaries.
\paragraph{Line 39}
\noindent\colorbox{gray!8}{\parbox{\dimexpr\linewidth-2\fboxsep\relax}{\raggedright\ttfamily\scriptsize     output logic [15:0] debug\_feature,}}
\par\smallskip
This line declares a module port. A port is a physical or wrapper-visible signal used to move clocks, control pulses, features, labels, predictions, or debug information across module boundaries.
\paragraph{Line 40}
\noindent\colorbox{gray!8}{\parbox{\dimexpr\linewidth-2\fboxsep\relax}{\raggedright\ttfamily\scriptsize     output logic [7:0] debug\_class,}}
\par\smallskip
This line declares a module port. A port is a physical or wrapper-visible signal used to move clocks, control pulses, features, labels, predictions, or debug information across module boundaries.
\paragraph{Line 41}
\noindent\colorbox{gray!8}{\parbox{\dimexpr\linewidth-2\fboxsep\relax}{\raggedright\ttfamily\scriptsize     output logic [31:0] debug\_update\_count,}}
\par\smallskip
This line declares a module port. A port is a physical or wrapper-visible signal used to move clocks, control pulses, features, labels, predictions, or debug information across module boundaries.
\paragraph{Line 42}
\noindent\colorbox{gray!8}{\parbox{\dimexpr\linewidth-2\fboxsep\relax}{\raggedright\ttfamily\scriptsize     output logic [31:0] debug\_infer\_skip\_count,}}
\par\smallskip
This line exposes or updates a hardware activity counter. The counter makes the feature gate visible in waveforms and CSV traces, so the power-saving opportunity is observable from RTL rather than only from Python statistics.
\paragraph{Line 43}
\noindent\colorbox{gray!8}{\parbox{\dimexpr\linewidth-2\fboxsep\relax}{\raggedright\ttfamily\scriptsize     output logic [31:0] debug\_train\_skip\_count}}
\par\smallskip
This line exposes or updates a hardware activity counter. The counter makes the feature gate visible in waveforms and CSV traces, so the power-saving opportunity is observable from RTL rather than only from Python statistics.
\paragraph{Line 44}
\noindent\colorbox{gray!8}{\parbox{\dimexpr\linewidth-2\fboxsep\relax}{\raggedright\ttfamily\scriptsize );}}
\par\smallskip
This SystemVerilog line contributes to the structural or sequential description of the accelerator. Its effect is resolved by XSim for verification and by Vivado during synthesis.
\paragraph{Line 45}
\noindent\colorbox{gray!8}{\parbox{\dimexpr\linewidth-2\fboxsep\relax}{\raggedright\ttfamily\scriptsize     localparam int FEATURE\_W = \$clog2(FEATURE\_DIM + 1);}}
\par\smallskip
This is a local compile-time or procedural assignment. In context it either defines a helper value or connects a temporary expression used by nearby hardware logic.
\paragraph{Line 46}
\noindent\colorbox{gray!8}{\parbox{\dimexpr\linewidth-2\fboxsep\relax}{\raggedright\ttfamily\scriptsize     localparam int CLASS\_W = \$clog2(NUM\_CLASSES + 1);}}
\par\smallskip
This is a local compile-time or procedural assignment. In context it either defines a helper value or connects a temporary expression used by nearby hardware logic.
\paragraph{Line 47}
\noindent\colorbox{gray!8}{\parbox{\dimexpr\linewidth-2\fboxsep\relax}{\raggedright\ttfamily\scriptsize \emph{blank line}}}
\par\smallskip
This blank line is not hardware or software behavior. It separates nearby declarations or operations so the file remains readable and reviewable.
\paragraph{Line 48}
\noindent\colorbox{gray!8}{\parbox{\dimexpr\linewidth-2\fboxsep\relax}{\raggedright\ttfamily\scriptsize     typedef enum logic [5:0] \{}}
\par\smallskip
This declares the finite-state-machine state type. Using named states makes the cycle-by-cycle control flow explicit instead of relying on anonymous numeric constants.
\paragraph{Line 49}
\noindent\colorbox{gray!8}{\parbox{\dimexpr\linewidth-2\fboxsep\relax}{\raggedright\ttfamily\scriptsize         ST\_IDLE       = 6'd0,}}
\par\smallskip
This names one state in the control FSM. The numeric encoding is what appears in waveform traces, while the name explains the pipeline phase being executed.
\paragraph{Line 50}
\noindent\colorbox{gray!8}{\parbox{\dimexpr\linewidth-2\fboxsep\relax}{\raggedright\ttfamily\scriptsize         ST\_LOAD       = 6'd1,}}
\par\smallskip
This names one state in the control FSM. The numeric encoding is what appears in waveform traces, while the name explains the pipeline phase being executed.
\paragraph{Line 51}
\noindent\colorbox{gray!8}{\parbox{\dimexpr\linewidth-2\fboxsep\relax}{\raggedright\ttfamily\scriptsize         ST\_CLASS\_INIT = 6'd2,}}
\par\smallskip
This names one state in the control FSM. The numeric encoding is what appears in waveform traces, while the name explains the pipeline phase being executed.
\paragraph{Line 52}
\noindent\colorbox{gray!8}{\parbox{\dimexpr\linewidth-2\fboxsep\relax}{\raggedright\ttfamily\scriptsize         ST\_MAC\_READ   = 6'd3,}}
\par\smallskip
This names one state in the control FSM. The numeric encoding is what appears in waveform traces, while the name explains the pipeline phase being executed.
\paragraph{Line 53}
\noindent\colorbox{gray!8}{\parbox{\dimexpr\linewidth-2\fboxsep\relax}{\raggedright\ttfamily\scriptsize         ST\_MAC\_MUL    = 6'd4,}}
\par\smallskip
This names one state in the control FSM. The numeric encoding is what appears in waveform traces, while the name explains the pipeline phase being executed.
\paragraph{Line 54}
\noindent\colorbox{gray!8}{\parbox{\dimexpr\linewidth-2\fboxsep\relax}{\raggedright\ttfamily\scriptsize         ST\_MAC\_ACCUM  = 6'd5,}}
\par\smallskip
This names one state in the control FSM. The numeric encoding is what appears in waveform traces, while the name explains the pipeline phase being executed.
\paragraph{Line 55}
\noindent\colorbox{gray!8}{\parbox{\dimexpr\linewidth-2\fboxsep\relax}{\raggedright\ttfamily\scriptsize         ST\_CLASS\_DONE = 6'd6,}}
\par\smallskip
This names one state in the control FSM. The numeric encoding is what appears in waveform traces, while the name explains the pipeline phase being executed.
\paragraph{Line 56}
\noindent\colorbox{gray!8}{\parbox{\dimexpr\linewidth-2\fboxsep\relax}{\raggedright\ttfamily\scriptsize         ST\_RESULT     = 6'd7,}}
\par\smallskip
This names one state in the control FSM. The numeric encoding is what appears in waveform traces, while the name explains the pipeline phase being executed.
\paragraph{Line 57}
\noindent\colorbox{gray!8}{\parbox{\dimexpr\linewidth-2\fboxsep\relax}{\raggedright\ttfamily\scriptsize         ST\_TRAIN\_CHECK      = 6'd8,}}
\par\smallskip
This names one state in the control FSM. The numeric encoding is what appears in waveform traces, while the name explains the pipeline phase being executed.
\paragraph{Line 58}
\noindent\colorbox{gray!8}{\parbox{\dimexpr\linewidth-2\fboxsep\relax}{\raggedright\ttfamily\scriptsize         ST\_TRAIN\_READ\_TRUE  = 6'd9,}}
\par\smallskip
This names one state in the control FSM. The numeric encoding is what appears in waveform traces, while the name explains the pipeline phase being executed.
\paragraph{Line 59}
\noindent\colorbox{gray!8}{\parbox{\dimexpr\linewidth-2\fboxsep\relax}{\raggedright\ttfamily\scriptsize         ST\_TRAIN\_CALC\_TRUE  = 6'd10,}}
\par\smallskip
This names one state in the control FSM. The numeric encoding is what appears in waveform traces, while the name explains the pipeline phase being executed.
\paragraph{Line 60}
\noindent\colorbox{gray!8}{\parbox{\dimexpr\linewidth-2\fboxsep\relax}{\raggedright\ttfamily\scriptsize         ST\_TRAIN\_WRITE\_TRUE = 6'd11,}}
\par\smallskip
This names one state in the control FSM. The numeric encoding is what appears in waveform traces, while the name explains the pipeline phase being executed.
\paragraph{Line 61}
\noindent\colorbox{gray!8}{\parbox{\dimexpr\linewidth-2\fboxsep\relax}{\raggedright\ttfamily\scriptsize         ST\_TRAIN\_READ\_PRED  = 6'd12,}}
\par\smallskip
This names one state in the control FSM. The numeric encoding is what appears in waveform traces, while the name explains the pipeline phase being executed.
\paragraph{Line 62}
\noindent\colorbox{gray!8}{\parbox{\dimexpr\linewidth-2\fboxsep\relax}{\raggedright\ttfamily\scriptsize         ST\_TRAIN\_CALC\_PRED  = 6'd13,}}
\par\smallskip
This names one state in the control FSM. The numeric encoding is what appears in waveform traces, while the name explains the pipeline phase being executed.
\paragraph{Line 63}
\noindent\colorbox{gray!8}{\parbox{\dimexpr\linewidth-2\fboxsep\relax}{\raggedright\ttfamily\scriptsize         ST\_TRAIN\_WRITE\_PRED = 6'd14,}}
\par\smallskip
This names one state in the control FSM. The numeric encoding is what appears in waveform traces, while the name explains the pipeline phase being executed.
\paragraph{Line 64}
\noindent\colorbox{gray!8}{\parbox{\dimexpr\linewidth-2\fboxsep\relax}{\raggedright\ttfamily\scriptsize         ST\_TRAIN\_BIAS\_TRUE  = 6'd15,}}
\par\smallskip
This names one state in the control FSM. The numeric encoding is what appears in waveform traces, while the name explains the pipeline phase being executed.
\paragraph{Line 65}
\noindent\colorbox{gray!8}{\parbox{\dimexpr\linewidth-2\fboxsep\relax}{\raggedright\ttfamily\scriptsize         ST\_TRAIN\_BIAS\_PRED  = 6'd16}}
\par\smallskip
This names one state in the control FSM. The numeric encoding is what appears in waveform traces, while the name explains the pipeline phase being executed.
\paragraph{Line 66}
\noindent\colorbox{gray!8}{\parbox{\dimexpr\linewidth-2\fboxsep\relax}{\raggedright\ttfamily\scriptsize     \} state\_t;}}
\par\smallskip
This SystemVerilog line contributes to the structural or sequential description of the accelerator. Its effect is resolved by XSim for verification and by Vivado during synthesis.
\paragraph{Line 67}
\noindent\colorbox{gray!8}{\parbox{\dimexpr\linewidth-2\fboxsep\relax}{\raggedright\ttfamily\scriptsize \emph{blank line}}}
\par\smallskip
This blank line is not hardware or software behavior. It separates nearby declarations or operations so the file remains readable and reviewable.
\paragraph{Line 68}
\noindent\colorbox{gray!8}{\parbox{\dimexpr\linewidth-2\fboxsep\relax}{\raggedright\ttfamily\scriptsize     state\_t state;}}
\par\smallskip
This SystemVerilog line contributes to the structural or sequential description of the accelerator. Its effect is resolved by XSim for verification and by Vivado during synthesis.
\paragraph{Line 69}
\noindent\colorbox{gray!8}{\parbox{\dimexpr\linewidth-2\fboxsep\relax}{\raggedright\ttfamily\scriptsize \emph{blank line}}}
\par\smallskip
This blank line is not hardware or software behavior. It separates nearby declarations or operations so the file remains readable and reviewable.
\paragraph{Line 70}
\noindent\colorbox{gray!8}{\parbox{\dimexpr\linewidth-2\fboxsep\relax}{\raggedright\ttfamily\scriptsize     (* ram\_style = "distributed" *) logic signed [FEATURE\_BITS-1:0] feature\_mem [0:FEATURE\_DIM-1];}}
\par\smallskip
This line touches the feature memory that holds one 384-dimensional fixed-point Spikingformer feature vector for the current image. Keeping the feature vector local avoids repeatedly streaming the same values for every class score.
\paragraph{Line 71}
\noindent\colorbox{gray!8}{\parbox{\dimexpr\linewidth-2\fboxsep\relax}{\raggedright\ttfamily\scriptsize     (* ram\_style = "block" *) logic signed [WEIGHT\_BITS-1:0] weight\_mem [0:NUM\_CLASSES*FEATURE\_DIM-1];}}
\par\smallskip
This line touches the trained classifier weight store. In the final RTL the large 10 x 384 int16 matrix is inferred as true dual-port BRAM, so weights remain on chip and can also be updated by the learning state machine.
\paragraph{Line 72}
\noindent\colorbox{gray!8}{\parbox{\dimexpr\linewidth-2\fboxsep\relax}{\raggedright\ttfamily\scriptsize     logic signed [ACC\_BITS-1:0] bias\_mem [0:NUM\_CLASSES-1];}}
\par\smallskip
This line touches the class-bias store. Bias values are part of the affine classifier score and are updated on chip in separate one-write states to keep synthesis clean.
\paragraph{Line 73}
\noindent\colorbox{gray!8}{\parbox{\dimexpr\linewidth-2\fboxsep\relax}{\raggedright\ttfamily\scriptsize     logic signed [ACC\_BITS-1:0] class\_score\_mem [0:NUM\_CLASSES-1];}}
\par\smallskip
This stores or reads the ten class scores for the current sample. Those scores are needed for argmax prediction, CSV testbench reporting, and the margin test used before an online update.
\paragraph{Line 74}
\noindent\colorbox{gray!8}{\parbox{\dimexpr\linewidth-2\fboxsep\relax}{\raggedright\ttfamily\scriptsize \emph{blank line}}}
\par\smallskip
This blank line is not hardware or software behavior. It separates nearby declarations or operations so the file remains readable and reviewable.
\paragraph{Line 75}
\noindent\colorbox{gray!8}{\parbox{\dimexpr\linewidth-2\fboxsep\relax}{\raggedright\ttfamily\scriptsize     logic [FEATURE\_W-1:0] feature\_idx;}}
\par\smallskip
This SystemVerilog line contributes to the structural or sequential description of the accelerator. Its effect is resolved by XSim for verification and by Vivado during synthesis.
\paragraph{Line 76}
\noindent\colorbox{gray!8}{\parbox{\dimexpr\linewidth-2\fboxsep\relax}{\raggedright\ttfamily\scriptsize     logic [CLASS\_W-1:0] class\_idx;}}
\par\smallskip
This SystemVerilog line contributes to the structural or sequential description of the accelerator. Its effect is resolved by XSim for verification and by Vivado during synthesis.
\paragraph{Line 77}
\noindent\colorbox{gray!8}{\parbox{\dimexpr\linewidth-2\fboxsep\relax}{\raggedright\ttfamily\scriptsize     logic signed [ACC\_BITS-1:0] accum;}}
\par\smallskip
This SystemVerilog line contributes to the structural or sequential description of the accelerator. Its effect is resolved by XSim for verification and by Vivado during synthesis.
\paragraph{Line 78}
\noindent\colorbox{gray!8}{\parbox{\dimexpr\linewidth-2\fboxsep\relax}{\raggedright\ttfamily\scriptsize     logic signed [FEATURE\_BITS-1:0] feature\_eff\_r;}}
\par\smallskip
This SystemVerilog line contributes to the structural or sequential description of the accelerator. Its effect is resolved by XSim for verification and by Vivado during synthesis.
\paragraph{Line 79}
\noindent\colorbox{gray!8}{\parbox{\dimexpr\linewidth-2\fboxsep\relax}{\raggedright\ttfamily\scriptsize     logic signed [WEIGHT\_BITS-1:0] weight\_r;}}
\par\smallskip
This SystemVerilog line contributes to the structural or sequential description of the accelerator. Its effect is resolved by XSim for verification and by Vivado during synthesis.
\paragraph{Line 80}
\noindent\colorbox{gray!8}{\parbox{\dimexpr\linewidth-2\fboxsep\relax}{\raggedright\ttfamily\scriptsize     logic signed [ACC\_BITS-1:0] product\_r;}}
\par\smallskip
This SystemVerilog line contributes to the structural or sequential description of the accelerator. Its effect is resolved by XSim for verification and by Vivado during synthesis.
\paragraph{Line 81}
\noindent\colorbox{gray!8}{\parbox{\dimexpr\linewidth-2\fboxsep\relax}{\raggedright\ttfamily\scriptsize     logic signed [WEIGHT\_BITS-1:0] train\_true\_weight\_r;}}
\par\smallskip
This SystemVerilog line contributes to the structural or sequential description of the accelerator. Its effect is resolved by XSim for verification and by Vivado during synthesis.
\paragraph{Line 82}
\noindent\colorbox{gray!8}{\parbox{\dimexpr\linewidth-2\fboxsep\relax}{\raggedright\ttfamily\scriptsize     logic signed [WEIGHT\_BITS-1:0] train\_pred\_weight\_r;}}
\par\smallskip
This SystemVerilog line contributes to the structural or sequential description of the accelerator. Its effect is resolved by XSim for verification and by Vivado during synthesis.
\paragraph{Line 83}
\noindent\colorbox{gray!8}{\parbox{\dimexpr\linewidth-2\fboxsep\relax}{\raggedright\ttfamily\scriptsize     logic signed [WEIGHT\_BITS-1:0] train\_write\_weight\_r;}}
\par\smallskip
This SystemVerilog line contributes to the structural or sequential description of the accelerator. Its effect is resolved by XSim for verification and by Vivado during synthesis.
\paragraph{Line 84}
\noindent\colorbox{gray!8}{\parbox{\dimexpr\linewidth-2\fboxsep\relax}{\raggedright\ttfamily\scriptsize     logic [3:0] train\_label\_r;}}
\par\smallskip
This SystemVerilog line contributes to the structural or sequential description of the accelerator. Its effect is resolved by XSim for verification and by Vivado during synthesis.
\paragraph{Line 85}
\noindent\colorbox{gray!8}{\parbox{\dimexpr\linewidth-2\fboxsep\relax}{\raggedright\ttfamily\scriptsize     logic [3:0] train\_pred\_label\_r;}}
\par\smallskip
This SystemVerilog line contributes to the structural or sequential description of the accelerator. Its effect is resolved by XSim for verification and by Vivado during synthesis.
\paragraph{Line 86}
\noindent\colorbox{gray!8}{\parbox{\dimexpr\linewidth-2\fboxsep\relax}{\raggedright\ttfamily\scriptsize     logic [31:0] update\_count;}}
\par\smallskip
This SystemVerilog line contributes to the structural or sequential description of the accelerator. Its effect is resolved by XSim for verification and by Vivado during synthesis.
\paragraph{Line 87}
\noindent\colorbox{gray!8}{\parbox{\dimexpr\linewidth-2\fboxsep\relax}{\raggedright\ttfamily\scriptsize     logic [31:0] infer\_skip\_count;}}
\par\smallskip
This line exposes or updates a hardware activity counter. The counter makes the feature gate visible in waveforms and CSV traces, so the power-saving opportunity is observable from RTL rather than only from Python statistics.
\paragraph{Line 88}
\noindent\colorbox{gray!8}{\parbox{\dimexpr\linewidth-2\fboxsep\relax}{\raggedright\ttfamily\scriptsize     logic [31:0] train\_skip\_count;}}
\par\smallskip
This line exposes or updates a hardware activity counter. The counter makes the feature gate visible in waveforms and CSV traces, so the power-saving opportunity is observable from RTL rather than only from Python statistics.
\paragraph{Line 89}
\noindent\colorbox{gray!8}{\parbox{\dimexpr\linewidth-2\fboxsep\relax}{\raggedright\ttfamily\scriptsize \emph{blank line}}}
\par\smallskip
This blank line is not hardware or software behavior. It separates nearby declarations or operations so the file remains readable and reviewable.
\paragraph{Line 90}
\noindent\colorbox{gray!8}{\parbox{\dimexpr\linewidth-2\fboxsep\relax}{\raggedright\ttfamily\scriptsize     assign feature\_ready = (state == ST\_LOAD);}}
\par\smallskip
This is continuous combinational wiring. The output follows its source expression immediately in hardware rather than waiting for a clock edge.
\paragraph{Line 91}
\noindent\colorbox{gray!8}{\parbox{\dimexpr\linewidth-2\fboxsep\relax}{\raggedright\ttfamily\scriptsize     assign busy = (state != ST\_IDLE);}}
\par\smallskip
This is continuous combinational wiring. The output follows its source expression immediately in hardware rather than waiting for a clock edge.
\paragraph{Line 92}
\noindent\colorbox{gray!8}{\parbox{\dimexpr\linewidth-2\fboxsep\relax}{\raggedright\ttfamily\scriptsize     assign debug\_state = state;}}
\par\smallskip
This is continuous combinational wiring. The output follows its source expression immediately in hardware rather than waiting for a clock edge.
\paragraph{Line 93}
\noindent\colorbox{gray!8}{\parbox{\dimexpr\linewidth-2\fboxsep\relax}{\raggedright\ttfamily\scriptsize     assign debug\_feature = \{\{(16-FEATURE\_W)\{1'b0\}\}, feature\_idx\};}}
\par\smallskip
This is continuous combinational wiring. The output follows its source expression immediately in hardware rather than waiting for a clock edge.
\paragraph{Line 94}
\noindent\colorbox{gray!8}{\parbox{\dimexpr\linewidth-2\fboxsep\relax}{\raggedright\ttfamily\scriptsize     assign debug\_class = \{4'd0, 4'(class\_idx)\};}}
\par\smallskip
This is continuous combinational wiring. The output follows its source expression immediately in hardware rather than waiting for a clock edge.
\paragraph{Line 95}
\noindent\colorbox{gray!8}{\parbox{\dimexpr\linewidth-2\fboxsep\relax}{\raggedright\ttfamily\scriptsize     assign debug\_update\_count = update\_count;}}
\par\smallskip
This is continuous combinational wiring. The output follows its source expression immediately in hardware rather than waiting for a clock edge.
\paragraph{Line 96}
\noindent\colorbox{gray!8}{\parbox{\dimexpr\linewidth-2\fboxsep\relax}{\raggedright\ttfamily\scriptsize     assign debug\_infer\_skip\_count = infer\_skip\_count;}}
\par\smallskip
This line exposes or updates a hardware activity counter. The counter makes the feature gate visible in waveforms and CSV traces, so the power-saving opportunity is observable from RTL rather than only from Python statistics.
\paragraph{Line 97}
\noindent\colorbox{gray!8}{\parbox{\dimexpr\linewidth-2\fboxsep\relax}{\raggedright\ttfamily\scriptsize     assign debug\_train\_skip\_count = train\_skip\_count;}}
\par\smallskip
This line exposes or updates a hardware activity counter. The counter makes the feature gate visible in waveforms and CSV traces, so the power-saving opportunity is observable from RTL rather than only from Python statistics.
\paragraph{Line 98}
\noindent\colorbox{gray!8}{\parbox{\dimexpr\linewidth-2\fboxsep\relax}{\raggedright\ttfamily\scriptsize \emph{blank line}}}
\par\smallskip
This blank line is not hardware or software behavior. It separates nearby declarations or operations so the file remains readable and reviewable.
\paragraph{Line 99}
\noindent\colorbox{gray!8}{\parbox{\dimexpr\linewidth-2\fboxsep\relax}{\raggedright\ttfamily\scriptsize     initial begin}}
\par\smallskip
This simulation/synthesis initialization block loads trained memories from generated .mem files. On FPGA, Vivado maps these initial values into BRAM/LUTRAM initialization data.
\paragraph{Line 100}
\noindent\colorbox{gray!8}{\parbox{\dimexpr\linewidth-2\fboxsep\relax}{\raggedright\ttfamily\scriptsize         \$readmemh("rtl/cifar10\_spikingformer\_head\_weights.mem", weight\_mem);}}
\par\smallskip
This line touches the trained classifier weight store. In the final RTL the large 10 x 384 int16 matrix is inferred as true dual-port BRAM, so weights remain on chip and can also be updated by the learning state machine.
\paragraph{Line 101}
\noindent\colorbox{gray!8}{\parbox{\dimexpr\linewidth-2\fboxsep\relax}{\raggedright\ttfamily\scriptsize         \$readmemh("rtl/cifar10\_spikingformer\_head\_bias.mem", bias\_mem);}}
\par\smallskip
This line touches the class-bias store. Bias values are part of the affine classifier score and are updated on chip in separate one-write states to keep synthesis clean.
\paragraph{Line 102}
\noindent\colorbox{gray!8}{\parbox{\dimexpr\linewidth-2\fboxsep\relax}{\raggedright\ttfamily\scriptsize     end}}
\par\smallskip
This SystemVerilog line contributes to the structural or sequential description of the accelerator. Its effect is resolved by XSim for verification and by Vivado during synthesis.
\paragraph{Line 103}
\noindent\colorbox{gray!8}{\parbox{\dimexpr\linewidth-2\fboxsep\relax}{\raggedright\ttfamily\scriptsize \emph{blank line}}}
\par\smallskip
This blank line is not hardware or software behavior. It separates nearby declarations or operations so the file remains readable and reviewable.
\paragraph{Line 104}
\noindent\colorbox{gray!8}{\parbox{\dimexpr\linewidth-2\fboxsep\relax}{\raggedright\ttfamily\scriptsize     function automatic logic signed [WEIGHT\_BITS-1:0] sat\_weight(}}
\par\smallskip
This begins a pure helper function used inside the RTL. It describes fixed combinational math, not a software function running on a CPU.
\paragraph{Line 105}
\noindent\colorbox{gray!8}{\parbox{\dimexpr\linewidth-2\fboxsep\relax}{\raggedright\ttfamily\scriptsize         input logic signed [ACC\_BITS-1:0] value}}
\par\smallskip
This line declares a module port. A port is a physical or wrapper-visible signal used to move clocks, control pulses, features, labels, predictions, or debug information across module boundaries.
\paragraph{Line 106}
\noindent\colorbox{gray!8}{\parbox{\dimexpr\linewidth-2\fboxsep\relax}{\raggedright\ttfamily\scriptsize     );}}
\par\smallskip
This SystemVerilog line contributes to the structural or sequential description of the accelerator. Its effect is resolved by XSim for verification and by Vivado during synthesis.
\paragraph{Line 107}
\noindent\colorbox{gray!8}{\parbox{\dimexpr\linewidth-2\fboxsep\relax}{\raggedright\ttfamily\scriptsize         begin}}
\par\smallskip
This SystemVerilog line contributes to the structural or sequential description of the accelerator. Its effect is resolved by XSim for verification and by Vivado during synthesis.
\paragraph{Line 108}
\noindent\colorbox{gray!8}{\parbox{\dimexpr\linewidth-2\fboxsep\relax}{\raggedright\ttfamily\scriptsize             if (value > ((1 <<< (WEIGHT\_BITS - 1)) - 1)) begin}}
\par\smallskip
This SystemVerilog line contributes to the structural or sequential description of the accelerator. Its effect is resolved by XSim for verification and by Vivado during synthesis.
\paragraph{Line 109}
\noindent\colorbox{gray!8}{\parbox{\dimexpr\linewidth-2\fboxsep\relax}{\raggedright\ttfamily\scriptsize                 sat\_weight = \{1'b0, \{(WEIGHT\_BITS-1)\{1'b1\}\}\};}}
\par\smallskip
This line is part of signed saturation for 16-bit weights. Saturation prevents online learning from wrapping around when an update would exceed the representable fixed-point range.
\paragraph{Line 110}
\noindent\colorbox{gray!8}{\parbox{\dimexpr\linewidth-2\fboxsep\relax}{\raggedright\ttfamily\scriptsize             end else if (value < -(1 <<< (WEIGHT\_BITS - 1))) begin}}
\par\smallskip
This SystemVerilog line contributes to the structural or sequential description of the accelerator. Its effect is resolved by XSim for verification and by Vivado during synthesis.
\paragraph{Line 111}
\noindent\colorbox{gray!8}{\parbox{\dimexpr\linewidth-2\fboxsep\relax}{\raggedright\ttfamily\scriptsize                 sat\_weight = \{1'b1, \{(WEIGHT\_BITS-1)\{1'b0\}\}\};}}
\par\smallskip
This line is part of signed saturation for 16-bit weights. Saturation prevents online learning from wrapping around when an update would exceed the representable fixed-point range.
\paragraph{Line 112}
\noindent\colorbox{gray!8}{\parbox{\dimexpr\linewidth-2\fboxsep\relax}{\raggedright\ttfamily\scriptsize             end else begin}}
\par\smallskip
This SystemVerilog line contributes to the structural or sequential description of the accelerator. Its effect is resolved by XSim for verification and by Vivado during synthesis.
\paragraph{Line 113}
\noindent\colorbox{gray!8}{\parbox{\dimexpr\linewidth-2\fboxsep\relax}{\raggedright\ttfamily\scriptsize                 sat\_weight = value[WEIGHT\_BITS-1:0];}}
\par\smallskip
This line is part of signed saturation for 16-bit weights. Saturation prevents online learning from wrapping around when an update would exceed the representable fixed-point range.
\paragraph{Line 114}
\noindent\colorbox{gray!8}{\parbox{\dimexpr\linewidth-2\fboxsep\relax}{\raggedright\ttfamily\scriptsize             end}}
\par\smallskip
This SystemVerilog line contributes to the structural or sequential description of the accelerator. Its effect is resolved by XSim for verification and by Vivado during synthesis.
\paragraph{Line 115}
\noindent\colorbox{gray!8}{\parbox{\dimexpr\linewidth-2\fboxsep\relax}{\raggedright\ttfamily\scriptsize         end}}
\par\smallskip
This SystemVerilog line contributes to the structural or sequential description of the accelerator. Its effect is resolved by XSim for verification and by Vivado during synthesis.
\paragraph{Line 116}
\noindent\colorbox{gray!8}{\parbox{\dimexpr\linewidth-2\fboxsep\relax}{\raggedright\ttfamily\scriptsize     endfunction}}
\par\smallskip
This SystemVerilog line contributes to the structural or sequential description of the accelerator. Its effect is resolved by XSim for verification and by Vivado during synthesis.
\paragraph{Line 117}
\noindent\colorbox{gray!8}{\parbox{\dimexpr\linewidth-2\fboxsep\relax}{\raggedright\ttfamily\scriptsize \emph{blank line}}}
\par\smallskip
This blank line is not hardware or software behavior. It separates nearby declarations or operations so the file remains readable and reviewable.
\paragraph{Line 118}
\noindent\colorbox{gray!8}{\parbox{\dimexpr\linewidth-2\fboxsep\relax}{\raggedright\ttfamily\scriptsize     function automatic int weight\_addr(input int cls, input int feat);}}
\par\smallskip
This line declares a module port. A port is a physical or wrapper-visible signal used to move clocks, control pulses, features, labels, predictions, or debug information across module boundaries.
\paragraph{Line 119}
\noindent\colorbox{gray!8}{\parbox{\dimexpr\linewidth-2\fboxsep\relax}{\raggedright\ttfamily\scriptsize         begin}}
\par\smallskip
This SystemVerilog line contributes to the structural or sequential description of the accelerator. Its effect is resolved by XSim for verification and by Vivado during synthesis.
\paragraph{Line 120}
\noindent\colorbox{gray!8}{\parbox{\dimexpr\linewidth-2\fboxsep\relax}{\raggedright\ttfamily\scriptsize             weight\_addr = cls * FEATURE\_DIM + feat;}}
\par\smallskip
This is a local compile-time or procedural assignment. In context it either defines a helper value or connects a temporary expression used by nearby hardware logic.
\paragraph{Line 121}
\noindent\colorbox{gray!8}{\parbox{\dimexpr\linewidth-2\fboxsep\relax}{\raggedright\ttfamily\scriptsize         end}}
\par\smallskip
This SystemVerilog line contributes to the structural or sequential description of the accelerator. Its effect is resolved by XSim for verification and by Vivado during synthesis.
\paragraph{Line 122}
\noindent\colorbox{gray!8}{\parbox{\dimexpr\linewidth-2\fboxsep\relax}{\raggedright\ttfamily\scriptsize     endfunction}}
\par\smallskip
This SystemVerilog line contributes to the structural or sequential description of the accelerator. Its effect is resolved by XSim for verification and by Vivado during synthesis.
\paragraph{Line 123}
\noindent\colorbox{gray!8}{\parbox{\dimexpr\linewidth-2\fboxsep\relax}{\raggedright\ttfamily\scriptsize \emph{blank line}}}
\par\smallskip
This blank line is not hardware or software behavior. It separates nearby declarations or operations so the file remains readable and reviewable.
\paragraph{Line 124}
\noindent\colorbox{gray!8}{\parbox{\dimexpr\linewidth-2\fboxsep\relax}{\raggedright\ttfamily\scriptsize     function automatic logic signed [ACC\_BITS-1:0] learn\_delta(}}
\par\smallskip
This begins a pure helper function used inside the RTL. It describes fixed combinational math, not a software function running on a CPU.
\paragraph{Line 125}
\noindent\colorbox{gray!8}{\parbox{\dimexpr\linewidth-2\fboxsep\relax}{\raggedright\ttfamily\scriptsize         input logic signed [FEATURE\_BITS-1:0] value}}
\par\smallskip
This line declares a module port. A port is a physical or wrapper-visible signal used to move clocks, control pulses, features, labels, predictions, or debug information across module boundaries.
\paragraph{Line 126}
\noindent\colorbox{gray!8}{\parbox{\dimexpr\linewidth-2\fboxsep\relax}{\raggedright\ttfamily\scriptsize     );}}
\par\smallskip
This SystemVerilog line contributes to the structural or sequential description of the accelerator. Its effect is resolved by XSim for verification and by Vivado during synthesis.
\paragraph{Line 127}
\noindent\colorbox{gray!8}{\parbox{\dimexpr\linewidth-2\fboxsep\relax}{\raggedright\ttfamily\scriptsize         logic signed [ACC\_BITS-1:0] extended;}}
\par\smallskip
This SystemVerilog line contributes to the structural or sequential description of the accelerator. Its effect is resolved by XSim for verification and by Vivado during synthesis.
\paragraph{Line 128}
\noindent\colorbox{gray!8}{\parbox{\dimexpr\linewidth-2\fboxsep\relax}{\raggedright\ttfamily\scriptsize         begin}}
\par\smallskip
This SystemVerilog line contributes to the structural or sequential description of the accelerator. Its effect is resolved by XSim for verification and by Vivado during synthesis.
\paragraph{Line 129}
\noindent\colorbox{gray!8}{\parbox{\dimexpr\linewidth-2\fboxsep\relax}{\raggedright\ttfamily\scriptsize             extended = \{\{(ACC\_BITS-FEATURE\_BITS)\{value[FEATURE\_BITS-1]\}\}, value\};}}
\par\smallskip
This is a local compile-time or procedural assignment. In context it either defines a helper value or connects a temporary expression used by nearby hardware logic.
\paragraph{Line 130}
\noindent\colorbox{gray!8}{\parbox{\dimexpr\linewidth-2\fboxsep\relax}{\raggedright\ttfamily\scriptsize             learn\_delta = extended >>> LEARN\_SHIFT;}}
\par\smallskip
This line implements the hardware learning step size. The feature is sign-extended and arithmetically shifted so positive and negative features update weights symmetrically.
\paragraph{Line 131}
\noindent\colorbox{gray!8}{\parbox{\dimexpr\linewidth-2\fboxsep\relax}{\raggedright\ttfamily\scriptsize         end}}
\par\smallskip
This SystemVerilog line contributes to the structural or sequential description of the accelerator. Its effect is resolved by XSim for verification and by Vivado during synthesis.
\paragraph{Line 132}
\noindent\colorbox{gray!8}{\parbox{\dimexpr\linewidth-2\fboxsep\relax}{\raggedright\ttfamily\scriptsize     endfunction}}
\par\smallskip
This SystemVerilog line contributes to the structural or sequential description of the accelerator. Its effect is resolved by XSim for verification and by Vivado during synthesis.
\paragraph{Line 133}
\noindent\colorbox{gray!8}{\parbox{\dimexpr\linewidth-2\fboxsep\relax}{\raggedright\ttfamily\scriptsize \emph{blank line}}}
\par\smallskip
This blank line is not hardware or software behavior. It separates nearby declarations or operations so the file remains readable and reviewable.
\paragraph{Line 134}
\noindent\colorbox{gray!8}{\parbox{\dimexpr\linewidth-2\fboxsep\relax}{\raggedright\ttfamily\scriptsize     function automatic logic [FEATURE\_BITS:0] abs\_feature(}}
\par\smallskip
This begins a pure helper function used inside the RTL. It describes fixed combinational math, not a software function running on a CPU.
\paragraph{Line 135}
\noindent\colorbox{gray!8}{\parbox{\dimexpr\linewidth-2\fboxsep\relax}{\raggedright\ttfamily\scriptsize         input logic signed [FEATURE\_BITS-1:0] value}}
\par\smallskip
This line declares a module port. A port is a physical or wrapper-visible signal used to move clocks, control pulses, features, labels, predictions, or debug information across module boundaries.
\paragraph{Line 136}
\noindent\colorbox{gray!8}{\parbox{\dimexpr\linewidth-2\fboxsep\relax}{\raggedright\ttfamily\scriptsize     );}}
\par\smallskip
This SystemVerilog line contributes to the structural or sequential description of the accelerator. Its effect is resolved by XSim for verification and by Vivado during synthesis.
\paragraph{Line 137}
\noindent\colorbox{gray!8}{\parbox{\dimexpr\linewidth-2\fboxsep\relax}{\raggedright\ttfamily\scriptsize         logic signed [FEATURE\_BITS:0] extended;}}
\par\smallskip
This SystemVerilog line contributes to the structural or sequential description of the accelerator. Its effect is resolved by XSim for verification and by Vivado during synthesis.
\paragraph{Line 138}
\noindent\colorbox{gray!8}{\parbox{\dimexpr\linewidth-2\fboxsep\relax}{\raggedright\ttfamily\scriptsize         begin}}
\par\smallskip
This SystemVerilog line contributes to the structural or sequential description of the accelerator. Its effect is resolved by XSim for verification and by Vivado during synthesis.
\paragraph{Line 139}
\noindent\colorbox{gray!8}{\parbox{\dimexpr\linewidth-2\fboxsep\relax}{\raggedright\ttfamily\scriptsize             extended = \{value[FEATURE\_BITS-1], value\};}}
\par\smallskip
This is a local compile-time or procedural assignment. In context it either defines a helper value or connects a temporary expression used by nearby hardware logic.
\paragraph{Line 140}
\noindent\colorbox{gray!8}{\parbox{\dimexpr\linewidth-2\fboxsep\relax}{\raggedright\ttfamily\scriptsize             if (extended < 0) begin}}
\par\smallskip
This SystemVerilog line contributes to the structural or sequential description of the accelerator. Its effect is resolved by XSim for verification and by Vivado during synthesis.
\paragraph{Line 141}
\noindent\colorbox{gray!8}{\parbox{\dimexpr\linewidth-2\fboxsep\relax}{\raggedright\ttfamily\scriptsize                 abs\_feature = (0 - extended);}}
\par\smallskip
This is a local compile-time or procedural assignment. In context it either defines a helper value or connects a temporary expression used by nearby hardware logic.
\paragraph{Line 142}
\noindent\colorbox{gray!8}{\parbox{\dimexpr\linewidth-2\fboxsep\relax}{\raggedright\ttfamily\scriptsize             end else begin}}
\par\smallskip
This SystemVerilog line contributes to the structural or sequential description of the accelerator. Its effect is resolved by XSim for verification and by Vivado during synthesis.
\paragraph{Line 143}
\noindent\colorbox{gray!8}{\parbox{\dimexpr\linewidth-2\fboxsep\relax}{\raggedright\ttfamily\scriptsize                 abs\_feature = extended[FEATURE\_BITS:0];}}
\par\smallskip
This is a local compile-time or procedural assignment. In context it either defines a helper value or connects a temporary expression used by nearby hardware logic.
\paragraph{Line 144}
\noindent\colorbox{gray!8}{\parbox{\dimexpr\linewidth-2\fboxsep\relax}{\raggedright\ttfamily\scriptsize             end}}
\par\smallskip
This SystemVerilog line contributes to the structural or sequential description of the accelerator. Its effect is resolved by XSim for verification and by Vivado during synthesis.
\paragraph{Line 145}
\noindent\colorbox{gray!8}{\parbox{\dimexpr\linewidth-2\fboxsep\relax}{\raggedright\ttfamily\scriptsize         end}}
\par\smallskip
This SystemVerilog line contributes to the structural or sequential description of the accelerator. Its effect is resolved by XSim for verification and by Vivado during synthesis.
\paragraph{Line 146}
\noindent\colorbox{gray!8}{\parbox{\dimexpr\linewidth-2\fboxsep\relax}{\raggedright\ttfamily\scriptsize     endfunction}}
\par\smallskip
This SystemVerilog line contributes to the structural or sequential description of the accelerator. Its effect is resolved by XSim for verification and by Vivado during synthesis.
\paragraph{Line 147}
\noindent\colorbox{gray!8}{\parbox{\dimexpr\linewidth-2\fboxsep\relax}{\raggedright\ttfamily\scriptsize \emph{blank line}}}
\par\smallskip
This blank line is not hardware or software behavior. It separates nearby declarations or operations so the file remains readable and reviewable.
\paragraph{Line 148}
\noindent\colorbox{gray!8}{\parbox{\dimexpr\linewidth-2\fboxsep\relax}{\raggedright\ttfamily\scriptsize     function automatic logic skip\_feature(}}
\par\smallskip
This line is part of the combinational decision that detects low-magnitude features. It is the hardware equivalent of asking whether a small activation should be silent for this inference/update pass.
\paragraph{Line 149}
\noindent\colorbox{gray!8}{\parbox{\dimexpr\linewidth-2\fboxsep\relax}{\raggedright\ttfamily\scriptsize         input logic signed [FEATURE\_BITS-1:0] value}}
\par\smallskip
This line declares a module port. A port is a physical or wrapper-visible signal used to move clocks, control pulses, features, labels, predictions, or debug information across module boundaries.
\paragraph{Line 150}
\noindent\colorbox{gray!8}{\parbox{\dimexpr\linewidth-2\fboxsep\relax}{\raggedright\ttfamily\scriptsize     );}}
\par\smallskip
This SystemVerilog line contributes to the structural or sequential description of the accelerator. Its effect is resolved by XSim for verification and by Vivado during synthesis.
\paragraph{Line 151}
\noindent\colorbox{gray!8}{\parbox{\dimexpr\linewidth-2\fboxsep\relax}{\raggedright\ttfamily\scriptsize         begin}}
\par\smallskip
This SystemVerilog line contributes to the structural or sequential description of the accelerator. Its effect is resolved by XSim for verification and by Vivado during synthesis.
\paragraph{Line 152}
\noindent\colorbox{gray!8}{\parbox{\dimexpr\linewidth-2\fboxsep\relax}{\raggedright\ttfamily\scriptsize             if (INFER\_SKIP\_ABS\_THRESH <= 0) begin}}
\par\smallskip
This line controls the feature-activity gate. When the absolute fixed-point feature is at or below this threshold, the RTL treats it as zero so the exported Spikingformer head ignores small noisy activations and records a measurable MAC-skip opportunity.
\paragraph{Line 153}
\noindent\colorbox{gray!8}{\parbox{\dimexpr\linewidth-2\fboxsep\relax}{\raggedright\ttfamily\scriptsize                 skip\_feature = 1'b0;}}
\par\smallskip
This line is part of the combinational decision that detects low-magnitude features. It is the hardware equivalent of asking whether a small activation should be silent for this inference/update pass.
\paragraph{Line 154}
\noindent\colorbox{gray!8}{\parbox{\dimexpr\linewidth-2\fboxsep\relax}{\raggedright\ttfamily\scriptsize             end else begin}}
\par\smallskip
This SystemVerilog line contributes to the structural or sequential description of the accelerator. Its effect is resolved by XSim for verification and by Vivado during synthesis.
\paragraph{Line 155}
\noindent\colorbox{gray!8}{\parbox{\dimexpr\linewidth-2\fboxsep\relax}{\raggedright\ttfamily\scriptsize                 skip\_feature = abs\_feature(value) <= INFER\_SKIP\_ABS\_THRESH;}}
\par\smallskip
This line controls the feature-activity gate. When the absolute fixed-point feature is at or below this threshold, the RTL treats it as zero so the exported Spikingformer head ignores small noisy activations and records a measurable MAC-skip opportunity.
\paragraph{Line 156}
\noindent\colorbox{gray!8}{\parbox{\dimexpr\linewidth-2\fboxsep\relax}{\raggedright\ttfamily\scriptsize             end}}
\par\smallskip
This SystemVerilog line contributes to the structural or sequential description of the accelerator. Its effect is resolved by XSim for verification and by Vivado during synthesis.
\paragraph{Line 157}
\noindent\colorbox{gray!8}{\parbox{\dimexpr\linewidth-2\fboxsep\relax}{\raggedright\ttfamily\scriptsize         end}}
\par\smallskip
This SystemVerilog line contributes to the structural or sequential description of the accelerator. Its effect is resolved by XSim for verification and by Vivado during synthesis.
\paragraph{Line 158}
\noindent\colorbox{gray!8}{\parbox{\dimexpr\linewidth-2\fboxsep\relax}{\raggedright\ttfamily\scriptsize     endfunction}}
\par\smallskip
This SystemVerilog line contributes to the structural or sequential description of the accelerator. Its effect is resolved by XSim for verification and by Vivado during synthesis.
\paragraph{Line 159}
\noindent\colorbox{gray!8}{\parbox{\dimexpr\linewidth-2\fboxsep\relax}{\raggedright\ttfamily\scriptsize \emph{blank line}}}
\par\smallskip
This blank line is not hardware or software behavior. It separates nearby declarations or operations so the file remains readable and reviewable.
\paragraph{Line 160}
\noindent\colorbox{gray!8}{\parbox{\dimexpr\linewidth-2\fboxsep\relax}{\raggedright\ttfamily\scriptsize     function automatic logic signed [FEATURE\_BITS-1:0] effective\_feature(}}
\par\smallskip
This line is part of the denoised feature path. The returned value is either the original signed feature or exactly zero, and that registered value is what feeds both the DSP multiply and the online-learning delta.
\paragraph{Line 161}
\noindent\colorbox{gray!8}{\parbox{\dimexpr\linewidth-2\fboxsep\relax}{\raggedright\ttfamily\scriptsize         input logic signed [FEATURE\_BITS-1:0] value}}
\par\smallskip
This line declares a module port. A port is a physical or wrapper-visible signal used to move clocks, control pulses, features, labels, predictions, or debug information across module boundaries.
\paragraph{Line 162}
\noindent\colorbox{gray!8}{\parbox{\dimexpr\linewidth-2\fboxsep\relax}{\raggedright\ttfamily\scriptsize     );}}
\par\smallskip
This SystemVerilog line contributes to the structural or sequential description of the accelerator. Its effect is resolved by XSim for verification and by Vivado during synthesis.
\paragraph{Line 163}
\noindent\colorbox{gray!8}{\parbox{\dimexpr\linewidth-2\fboxsep\relax}{\raggedright\ttfamily\scriptsize         begin}}
\par\smallskip
This SystemVerilog line contributes to the structural or sequential description of the accelerator. Its effect is resolved by XSim for verification and by Vivado during synthesis.
\paragraph{Line 164}
\noindent\colorbox{gray!8}{\parbox{\dimexpr\linewidth-2\fboxsep\relax}{\raggedright\ttfamily\scriptsize             if (skip\_feature(value)) begin}}
\par\smallskip
This line is part of the combinational decision that detects low-magnitude features. It is the hardware equivalent of asking whether a small activation should be silent for this inference/update pass.
\paragraph{Line 165}
\noindent\colorbox{gray!8}{\parbox{\dimexpr\linewidth-2\fboxsep\relax}{\raggedright\ttfamily\scriptsize                 effective\_feature = '0;}}
\par\smallskip
This line is part of the denoised feature path. The returned value is either the original signed feature or exactly zero, and that registered value is what feeds both the DSP multiply and the online-learning delta.
\paragraph{Line 166}
\noindent\colorbox{gray!8}{\parbox{\dimexpr\linewidth-2\fboxsep\relax}{\raggedright\ttfamily\scriptsize             end else begin}}
\par\smallskip
This SystemVerilog line contributes to the structural or sequential description of the accelerator. Its effect is resolved by XSim for verification and by Vivado during synthesis.
\paragraph{Line 167}
\noindent\colorbox{gray!8}{\parbox{\dimexpr\linewidth-2\fboxsep\relax}{\raggedright\ttfamily\scriptsize                 effective\_feature = value;}}
\par\smallskip
This line is part of the denoised feature path. The returned value is either the original signed feature or exactly zero, and that registered value is what feeds both the DSP multiply and the online-learning delta.
\paragraph{Line 168}
\noindent\colorbox{gray!8}{\parbox{\dimexpr\linewidth-2\fboxsep\relax}{\raggedright\ttfamily\scriptsize             end}}
\par\smallskip
This SystemVerilog line contributes to the structural or sequential description of the accelerator. Its effect is resolved by XSim for verification and by Vivado during synthesis.
\paragraph{Line 169}
\noindent\colorbox{gray!8}{\parbox{\dimexpr\linewidth-2\fboxsep\relax}{\raggedright\ttfamily\scriptsize         end}}
\par\smallskip
This SystemVerilog line contributes to the structural or sequential description of the accelerator. Its effect is resolved by XSim for verification and by Vivado during synthesis.
\paragraph{Line 170}
\noindent\colorbox{gray!8}{\parbox{\dimexpr\linewidth-2\fboxsep\relax}{\raggedright\ttfamily\scriptsize     endfunction}}
\par\smallskip
This SystemVerilog line contributes to the structural or sequential description of the accelerator. Its effect is resolved by XSim for verification and by Vivado during synthesis.
\paragraph{Line 171}
\noindent\colorbox{gray!8}{\parbox{\dimexpr\linewidth-2\fboxsep\relax}{\raggedright\ttfamily\scriptsize \emph{blank line}}}
\par\smallskip
This blank line is not hardware or software behavior. It separates nearby declarations or operations so the file remains readable and reviewable.
\paragraph{Line 172}
\noindent\colorbox{gray!8}{\parbox{\dimexpr\linewidth-2\fboxsep\relax}{\raggedright\ttfamily\scriptsize     always\_ff @(posedge clk or negedge rst\_n) begin}}
\par\smallskip
This begins clocked sequential logic. Registers and RAM writes inside this block update only on a clock edge, which is what makes the design synthesizable and waveform-predictable.
\paragraph{Line 173}
\noindent\colorbox{gray!8}{\parbox{\dimexpr\linewidth-2\fboxsep\relax}{\raggedright\ttfamily\scriptsize         if (!rst\_n) begin}}
\par\smallskip
This is the active-low reset branch. It puts control registers and visible outputs into a known starting state before a frame is processed.
\paragraph{Line 174}
\noindent\colorbox{gray!8}{\parbox{\dimexpr\linewidth-2\fboxsep\relax}{\raggedright\ttfamily\scriptsize             state <= ST\_IDLE;}}
\par\smallskip
This schedules the next FSM state for the following clock edge. That transition is the hardware pipeline control path.
\paragraph{Line 175}
\noindent\colorbox{gray!8}{\parbox{\dimexpr\linewidth-2\fboxsep\relax}{\raggedright\ttfamily\scriptsize             feature\_idx <= '0;}}
\par\smallskip
This is a nonblocking register or memory assignment. The right side is sampled now and the left side updates together with other clocked assignments at the end of the clock edge.
\paragraph{Line 176}
\noindent\colorbox{gray!8}{\parbox{\dimexpr\linewidth-2\fboxsep\relax}{\raggedright\ttfamily\scriptsize             class\_idx <= '0;}}
\par\smallskip
This is a nonblocking register or memory assignment. The right side is sampled now and the left side updates together with other clocked assignments at the end of the clock edge.
\paragraph{Line 177}
\noindent\colorbox{gray!8}{\parbox{\dimexpr\linewidth-2\fboxsep\relax}{\raggedright\ttfamily\scriptsize             accum <= '0;}}
\par\smallskip
This is a nonblocking register or memory assignment. The right side is sampled now and the left side updates together with other clocked assignments at the end of the clock edge.
\paragraph{Line 178}
\noindent\colorbox{gray!8}{\parbox{\dimexpr\linewidth-2\fboxsep\relax}{\raggedright\ttfamily\scriptsize             product\_r <= '0;}}
\par\smallskip
This is a nonblocking register or memory assignment. The right side is sampled now and the left side updates together with other clocked assignments at the end of the clock edge.
\paragraph{Line 179}
\noindent\colorbox{gray!8}{\parbox{\dimexpr\linewidth-2\fboxsep\relax}{\raggedright\ttfamily\scriptsize             result\_valid <= 1'b0;}}
\par\smallskip
This is a nonblocking register or memory assignment. The right side is sampled now and the left side updates together with other clocked assignments at the end of the clock edge.
\paragraph{Line 180}
\noindent\colorbox{gray!8}{\parbox{\dimexpr\linewidth-2\fboxsep\relax}{\raggedright\ttfamily\scriptsize             train\_update\_valid <= 1'b0;}}
\par\smallskip
This is a nonblocking register or memory assignment. The right side is sampled now and the left side updates together with other clocked assignments at the end of the clock edge.
\paragraph{Line 181}
\noindent\colorbox{gray!8}{\parbox{\dimexpr\linewidth-2\fboxsep\relax}{\raggedright\ttfamily\scriptsize             train\_update\_applied <= 1'b0;}}
\par\smallskip
This is a nonblocking register or memory assignment. The right side is sampled now and the left side updates together with other clocked assignments at the end of the clock edge.
\paragraph{Line 182}
\noindent\colorbox{gray!8}{\parbox{\dimexpr\linewidth-2\fboxsep\relax}{\raggedright\ttfamily\scriptsize             predicted\_label <= 4'd0;}}
\par\smallskip
This is a nonblocking register or memory assignment. The right side is sampled now and the left side updates together with other clocked assignments at the end of the clock edge.
\paragraph{Line 183}
\noindent\colorbox{gray!8}{\parbox{\dimexpr\linewidth-2\fboxsep\relax}{\raggedright\ttfamily\scriptsize             predicted\_score <= '0;}}
\par\smallskip
This is a nonblocking register or memory assignment. The right side is sampled now and the left side updates together with other clocked assignments at the end of the clock edge.
\paragraph{Line 184}
\noindent\colorbox{gray!8}{\parbox{\dimexpr\linewidth-2\fboxsep\relax}{\raggedright\ttfamily\scriptsize             train\_write\_weight\_r <= '0;}}
\par\smallskip
This is a nonblocking register or memory assignment. The right side is sampled now and the left side updates together with other clocked assignments at the end of the clock edge.
\paragraph{Line 185}
\noindent\colorbox{gray!8}{\parbox{\dimexpr\linewidth-2\fboxsep\relax}{\raggedright\ttfamily\scriptsize             train\_label\_r <= 4'd0;}}
\par\smallskip
This is a nonblocking register or memory assignment. The right side is sampled now and the left side updates together with other clocked assignments at the end of the clock edge.
\paragraph{Line 186}
\noindent\colorbox{gray!8}{\parbox{\dimexpr\linewidth-2\fboxsep\relax}{\raggedright\ttfamily\scriptsize             train\_pred\_label\_r <= 4'd0;}}
\par\smallskip
This is a nonblocking register or memory assignment. The right side is sampled now and the left side updates together with other clocked assignments at the end of the clock edge.
\paragraph{Line 187}
\noindent\colorbox{gray!8}{\parbox{\dimexpr\linewidth-2\fboxsep\relax}{\raggedright\ttfamily\scriptsize             update\_count <= 32'd0;}}
\par\smallskip
This is a nonblocking register or memory assignment. The right side is sampled now and the left side updates together with other clocked assignments at the end of the clock edge.
\paragraph{Line 188}
\noindent\colorbox{gray!8}{\parbox{\dimexpr\linewidth-2\fboxsep\relax}{\raggedright\ttfamily\scriptsize             infer\_skip\_count <= 32'd0;}}
\par\smallskip
This line exposes or updates a hardware activity counter. The counter makes the feature gate visible in waveforms and CSV traces, so the power-saving opportunity is observable from RTL rather than only from Python statistics.
\paragraph{Line 189}
\noindent\colorbox{gray!8}{\parbox{\dimexpr\linewidth-2\fboxsep\relax}{\raggedright\ttfamily\scriptsize             train\_skip\_count <= 32'd0;}}
\par\smallskip
This line exposes or updates a hardware activity counter. The counter makes the feature gate visible in waveforms and CSV traces, so the power-saving opportunity is observable from RTL rather than only from Python statistics.
\paragraph{Line 190}
\noindent\colorbox{gray!8}{\parbox{\dimexpr\linewidth-2\fboxsep\relax}{\raggedright\ttfamily\scriptsize             for (int c = 0; c < NUM\_CLASSES; c++) begin}}
\par\smallskip
This is a local compile-time or procedural assignment. In context it either defines a helper value or connects a temporary expression used by nearby hardware logic.
\paragraph{Line 191}
\noindent\colorbox{gray!8}{\parbox{\dimexpr\linewidth-2\fboxsep\relax}{\raggedright\ttfamily\scriptsize                 class\_score\_mem[c] <= '0;}}
\par\smallskip
This stores or reads the ten class scores for the current sample. Those scores are needed for argmax prediction, CSV testbench reporting, and the margin test used before an online update.
\paragraph{Line 192}
\noindent\colorbox{gray!8}{\parbox{\dimexpr\linewidth-2\fboxsep\relax}{\raggedright\ttfamily\scriptsize             end}}
\par\smallskip
This SystemVerilog line contributes to the structural or sequential description of the accelerator. Its effect is resolved by XSim for verification and by Vivado during synthesis.
\paragraph{Line 193}
\noindent\colorbox{gray!8}{\parbox{\dimexpr\linewidth-2\fboxsep\relax}{\raggedright\ttfamily\scriptsize         end else begin}}
\par\smallskip
This SystemVerilog line contributes to the structural or sequential description of the accelerator. Its effect is resolved by XSim for verification and by Vivado during synthesis.
\paragraph{Line 194}
\noindent\colorbox{gray!8}{\parbox{\dimexpr\linewidth-2\fboxsep\relax}{\raggedright\ttfamily\scriptsize             result\_valid <= 1'b0;}}
\par\smallskip
This is a nonblocking register or memory assignment. The right side is sampled now and the left side updates together with other clocked assignments at the end of the clock edge.
\paragraph{Line 195}
\noindent\colorbox{gray!8}{\parbox{\dimexpr\linewidth-2\fboxsep\relax}{\raggedright\ttfamily\scriptsize             train\_update\_valid <= 1'b0;}}
\par\smallskip
This is a nonblocking register or memory assignment. The right side is sampled now and the left side updates together with other clocked assignments at the end of the clock edge.
\paragraph{Line 196}
\noindent\colorbox{gray!8}{\parbox{\dimexpr\linewidth-2\fboxsep\relax}{\raggedright\ttfamily\scriptsize             train\_update\_applied <= 1'b0;}}
\par\smallskip
This is a nonblocking register or memory assignment. The right side is sampled now and the left side updates together with other clocked assignments at the end of the clock edge.
\paragraph{Line 197}
\noindent\colorbox{gray!8}{\parbox{\dimexpr\linewidth-2\fboxsep\relax}{\raggedright\ttfamily\scriptsize             unique case (state)}}
\par\smallskip
This starts a state-machine branch table. The unique qualifier tells tools that one legal state should match at a time, improving diagnostics and optimization.
\paragraph{Line 198}
\noindent\colorbox{gray!8}{\parbox{\dimexpr\linewidth-2\fboxsep\relax}{\raggedright\ttfamily\scriptsize                 ST\_IDLE: begin}}
\par\smallskip
This names one state in the control FSM. The numeric encoding is what appears in waveform traces, while the name explains the pipeline phase being executed.
\paragraph{Line 199}
\noindent\colorbox{gray!8}{\parbox{\dimexpr\linewidth-2\fboxsep\relax}{\raggedright\ttfamily\scriptsize                     if (start) begin}}
\par\smallskip
This SystemVerilog line contributes to the structural or sequential description of the accelerator. Its effect is resolved by XSim for verification and by Vivado during synthesis.
\paragraph{Line 200}
\noindent\colorbox{gray!8}{\parbox{\dimexpr\linewidth-2\fboxsep\relax}{\raggedright\ttfamily\scriptsize                         feature\_idx <= '0;}}
\par\smallskip
This is a nonblocking register or memory assignment. The right side is sampled now and the left side updates together with other clocked assignments at the end of the clock edge.
\paragraph{Line 201}
\noindent\colorbox{gray!8}{\parbox{\dimexpr\linewidth-2\fboxsep\relax}{\raggedright\ttfamily\scriptsize                         state <= ST\_LOAD;}}
\par\smallskip
This schedules the next FSM state for the following clock edge. That transition is the hardware pipeline control path.
\paragraph{Line 202}
\noindent\colorbox{gray!8}{\parbox{\dimexpr\linewidth-2\fboxsep\relax}{\raggedright\ttfamily\scriptsize                     end}}
\par\smallskip
This SystemVerilog line contributes to the structural or sequential description of the accelerator. Its effect is resolved by XSim for verification and by Vivado during synthesis.
\paragraph{Line 203}
\noindent\colorbox{gray!8}{\parbox{\dimexpr\linewidth-2\fboxsep\relax}{\raggedright\ttfamily\scriptsize                 end}}
\par\smallskip
This SystemVerilog line contributes to the structural or sequential description of the accelerator. Its effect is resolved by XSim for verification and by Vivado during synthesis.
\paragraph{Line 204}
\noindent\colorbox{gray!8}{\parbox{\dimexpr\linewidth-2\fboxsep\relax}{\raggedright\ttfamily\scriptsize \emph{blank line}}}
\par\smallskip
This blank line is not hardware or software behavior. It separates nearby declarations or operations so the file remains readable and reviewable.
\paragraph{Line 205}
\noindent\colorbox{gray!8}{\parbox{\dimexpr\linewidth-2\fboxsep\relax}{\raggedright\ttfamily\scriptsize                 ST\_LOAD: begin}}
\par\smallskip
This names one state in the control FSM. The numeric encoding is what appears in waveform traces, while the name explains the pipeline phase being executed.
\paragraph{Line 206}
\noindent\colorbox{gray!8}{\parbox{\dimexpr\linewidth-2\fboxsep\relax}{\raggedright\ttfamily\scriptsize                     if (feature\_valid \&\& feature\_ready) begin}}
\par\smallskip
This SystemVerilog line contributes to the structural or sequential description of the accelerator. Its effect is resolved by XSim for verification and by Vivado during synthesis.
\paragraph{Line 207}
\noindent\colorbox{gray!8}{\parbox{\dimexpr\linewidth-2\fboxsep\relax}{\raggedright\ttfamily\scriptsize                         if (feature\_idx == FEATURE\_DIM - 1) begin}}
\par\smallskip
This is a local compile-time or procedural assignment. In context it either defines a helper value or connects a temporary expression used by nearby hardware logic.
\paragraph{Line 208}
\noindent\colorbox{gray!8}{\parbox{\dimexpr\linewidth-2\fboxsep\relax}{\raggedright\ttfamily\scriptsize                             feature\_idx <= '0;}}
\par\smallskip
This is a nonblocking register or memory assignment. The right side is sampled now and the left side updates together with other clocked assignments at the end of the clock edge.
\paragraph{Line 209}
\noindent\colorbox{gray!8}{\parbox{\dimexpr\linewidth-2\fboxsep\relax}{\raggedright\ttfamily\scriptsize                             class\_idx <= '0;}}
\par\smallskip
This is a nonblocking register or memory assignment. The right side is sampled now and the left side updates together with other clocked assignments at the end of the clock edge.
\paragraph{Line 210}
\noindent\colorbox{gray!8}{\parbox{\dimexpr\linewidth-2\fboxsep\relax}{\raggedright\ttfamily\scriptsize                             state <= ST\_CLASS\_INIT;}}
\par\smallskip
This schedules the next FSM state for the following clock edge. That transition is the hardware pipeline control path.
\paragraph{Line 211}
\noindent\colorbox{gray!8}{\parbox{\dimexpr\linewidth-2\fboxsep\relax}{\raggedright\ttfamily\scriptsize                         end else begin}}
\par\smallskip
This SystemVerilog line contributes to the structural or sequential description of the accelerator. Its effect is resolved by XSim for verification and by Vivado during synthesis.
\paragraph{Line 212}
\noindent\colorbox{gray!8}{\parbox{\dimexpr\linewidth-2\fboxsep\relax}{\raggedright\ttfamily\scriptsize                             feature\_idx <= feature\_idx + 1'b1;}}
\par\smallskip
This is a nonblocking register or memory assignment. The right side is sampled now and the left side updates together with other clocked assignments at the end of the clock edge.
\paragraph{Line 213}
\noindent\colorbox{gray!8}{\parbox{\dimexpr\linewidth-2\fboxsep\relax}{\raggedright\ttfamily\scriptsize                         end}}
\par\smallskip
This SystemVerilog line contributes to the structural or sequential description of the accelerator. Its effect is resolved by XSim for verification and by Vivado during synthesis.
\paragraph{Line 214}
\noindent\colorbox{gray!8}{\parbox{\dimexpr\linewidth-2\fboxsep\relax}{\raggedright\ttfamily\scriptsize                     end else if (frame\_done \&\& feature\_idx == FEATURE\_DIM) begin}}
\par\smallskip
This is a local compile-time or procedural assignment. In context it either defines a helper value or connects a temporary expression used by nearby hardware logic.
\paragraph{Line 215}
\noindent\colorbox{gray!8}{\parbox{\dimexpr\linewidth-2\fboxsep\relax}{\raggedright\ttfamily\scriptsize                         feature\_idx <= '0;}}
\par\smallskip
This is a nonblocking register or memory assignment. The right side is sampled now and the left side updates together with other clocked assignments at the end of the clock edge.
\paragraph{Line 216}
\noindent\colorbox{gray!8}{\parbox{\dimexpr\linewidth-2\fboxsep\relax}{\raggedright\ttfamily\scriptsize                         class\_idx <= '0;}}
\par\smallskip
This is a nonblocking register or memory assignment. The right side is sampled now and the left side updates together with other clocked assignments at the end of the clock edge.
\paragraph{Line 217}
\noindent\colorbox{gray!8}{\parbox{\dimexpr\linewidth-2\fboxsep\relax}{\raggedright\ttfamily\scriptsize                         state <= ST\_CLASS\_INIT;}}
\par\smallskip
This schedules the next FSM state for the following clock edge. That transition is the hardware pipeline control path.
\paragraph{Line 218}
\noindent\colorbox{gray!8}{\parbox{\dimexpr\linewidth-2\fboxsep\relax}{\raggedright\ttfamily\scriptsize                     end}}
\par\smallskip
This SystemVerilog line contributes to the structural or sequential description of the accelerator. Its effect is resolved by XSim for verification and by Vivado during synthesis.
\paragraph{Line 219}
\noindent\colorbox{gray!8}{\parbox{\dimexpr\linewidth-2\fboxsep\relax}{\raggedright\ttfamily\scriptsize                 end}}
\par\smallskip
This SystemVerilog line contributes to the structural or sequential description of the accelerator. Its effect is resolved by XSim for verification and by Vivado during synthesis.
\paragraph{Line 220}
\noindent\colorbox{gray!8}{\parbox{\dimexpr\linewidth-2\fboxsep\relax}{\raggedright\ttfamily\scriptsize \emph{blank line}}}
\par\smallskip
This blank line is not hardware or software behavior. It separates nearby declarations or operations so the file remains readable and reviewable.
\paragraph{Line 221}
\noindent\colorbox{gray!8}{\parbox{\dimexpr\linewidth-2\fboxsep\relax}{\raggedright\ttfamily\scriptsize                 ST\_CLASS\_INIT: begin}}
\par\smallskip
This names one state in the control FSM. The numeric encoding is what appears in waveform traces, while the name explains the pipeline phase being executed.
\paragraph{Line 222}
\noindent\colorbox{gray!8}{\parbox{\dimexpr\linewidth-2\fboxsep\relax}{\raggedright\ttfamily\scriptsize                     accum <= bias\_mem[class\_idx];}}
\par\smallskip
This line touches the class-bias store. Bias values are part of the affine classifier score and are updated on chip in separate one-write states to keep synthesis clean.
\paragraph{Line 223}
\noindent\colorbox{gray!8}{\parbox{\dimexpr\linewidth-2\fboxsep\relax}{\raggedright\ttfamily\scriptsize                     feature\_idx <= '0;}}
\par\smallskip
This is a nonblocking register or memory assignment. The right side is sampled now and the left side updates together with other clocked assignments at the end of the clock edge.
\paragraph{Line 224}
\noindent\colorbox{gray!8}{\parbox{\dimexpr\linewidth-2\fboxsep\relax}{\raggedright\ttfamily\scriptsize                     state <= ST\_MAC\_READ;}}
\par\smallskip
This schedules the next FSM state for the following clock edge. That transition is the hardware pipeline control path.
\paragraph{Line 225}
\noindent\colorbox{gray!8}{\parbox{\dimexpr\linewidth-2\fboxsep\relax}{\raggedright\ttfamily\scriptsize                 end}}
\par\smallskip
This SystemVerilog line contributes to the structural or sequential description of the accelerator. Its effect is resolved by XSim for verification and by Vivado during synthesis.
\paragraph{Line 226}
\noindent\colorbox{gray!8}{\parbox{\dimexpr\linewidth-2\fboxsep\relax}{\raggedright\ttfamily\scriptsize \emph{blank line}}}
\par\smallskip
This blank line is not hardware or software behavior. It separates nearby declarations or operations so the file remains readable and reviewable.
\paragraph{Line 227}
\noindent\colorbox{gray!8}{\parbox{\dimexpr\linewidth-2\fboxsep\relax}{\raggedright\ttfamily\scriptsize                 ST\_MAC\_READ: begin}}
\par\smallskip
This names one state in the control FSM. The numeric encoding is what appears in waveform traces, while the name explains the pipeline phase being executed.
\paragraph{Line 228}
\noindent\colorbox{gray!8}{\parbox{\dimexpr\linewidth-2\fboxsep\relax}{\raggedright\ttfamily\scriptsize                     if (skip\_feature(feature\_mem[feature\_idx])) begin}}
\par\smallskip
This line is part of the combinational decision that detects low-magnitude features. It is the hardware equivalent of asking whether a small activation should be silent for this inference/update pass.
\paragraph{Line 229}
\noindent\colorbox{gray!8}{\parbox{\dimexpr\linewidth-2\fboxsep\relax}{\raggedright\ttfamily\scriptsize                         infer\_skip\_count <= infer\_skip\_count + 1'b1;}}
\par\smallskip
This line exposes or updates a hardware activity counter. The counter makes the feature gate visible in waveforms and CSV traces, so the power-saving opportunity is observable from RTL rather than only from Python statistics.
\paragraph{Line 230}
\noindent\colorbox{gray!8}{\parbox{\dimexpr\linewidth-2\fboxsep\relax}{\raggedright\ttfamily\scriptsize                         state <= ST\_CLASS\_DONE;}}
\par\smallskip
This schedules the next FSM state for the following clock edge. That transition is the hardware pipeline control path.
\paragraph{Line 231}
\noindent\colorbox{gray!8}{\parbox{\dimexpr\linewidth-2\fboxsep\relax}{\raggedright\ttfamily\scriptsize                     end else begin}}
\par\smallskip
This SystemVerilog line contributes to the structural or sequential description of the accelerator. Its effect is resolved by XSim for verification and by Vivado during synthesis.
\paragraph{Line 232}
\noindent\colorbox{gray!8}{\parbox{\dimexpr\linewidth-2\fboxsep\relax}{\raggedright\ttfamily\scriptsize                         state <= ST\_MAC\_MUL;}}
\par\smallskip
This schedules the next FSM state for the following clock edge. That transition is the hardware pipeline control path.
\paragraph{Line 233}
\noindent\colorbox{gray!8}{\parbox{\dimexpr\linewidth-2\fboxsep\relax}{\raggedright\ttfamily\scriptsize                     end}}
\par\smallskip
This SystemVerilog line contributes to the structural or sequential description of the accelerator. Its effect is resolved by XSim for verification and by Vivado during synthesis.
\paragraph{Line 234}
\noindent\colorbox{gray!8}{\parbox{\dimexpr\linewidth-2\fboxsep\relax}{\raggedright\ttfamily\scriptsize                 end}}
\par\smallskip
This SystemVerilog line contributes to the structural or sequential description of the accelerator. Its effect is resolved by XSim for verification and by Vivado during synthesis.
\paragraph{Line 235}
\noindent\colorbox{gray!8}{\parbox{\dimexpr\linewidth-2\fboxsep\relax}{\raggedright\ttfamily\scriptsize \emph{blank line}}}
\par\smallskip
This blank line is not hardware or software behavior. It separates nearby declarations or operations so the file remains readable and reviewable.
\paragraph{Line 236}
\noindent\colorbox{gray!8}{\parbox{\dimexpr\linewidth-2\fboxsep\relax}{\raggedright\ttfamily\scriptsize                 ST\_MAC\_MUL: begin}}
\par\smallskip
This names one state in the control FSM. The numeric encoding is what appears in waveform traces, while the name explains the pipeline phase being executed.
\paragraph{Line 237}
\noindent\colorbox{gray!8}{\parbox{\dimexpr\linewidth-2\fboxsep\relax}{\raggedright\ttfamily\scriptsize                     product\_r <= \$signed(feature\_eff\_r) * \$signed(weight\_r);}}
\par\smallskip
This is a nonblocking register or memory assignment. The right side is sampled now and the left side updates together with other clocked assignments at the end of the clock edge.
\paragraph{Line 238}
\noindent\colorbox{gray!8}{\parbox{\dimexpr\linewidth-2\fboxsep\relax}{\raggedright\ttfamily\scriptsize                     state <= ST\_MAC\_ACCUM;}}
\par\smallskip
This schedules the next FSM state for the following clock edge. That transition is the hardware pipeline control path.
\paragraph{Line 239}
\noindent\colorbox{gray!8}{\parbox{\dimexpr\linewidth-2\fboxsep\relax}{\raggedright\ttfamily\scriptsize                 end}}
\par\smallskip
This SystemVerilog line contributes to the structural or sequential description of the accelerator. Its effect is resolved by XSim for verification and by Vivado during synthesis.
\paragraph{Line 240}
\noindent\colorbox{gray!8}{\parbox{\dimexpr\linewidth-2\fboxsep\relax}{\raggedright\ttfamily\scriptsize \emph{blank line}}}
\par\smallskip
This blank line is not hardware or software behavior. It separates nearby declarations or operations so the file remains readable and reviewable.
\paragraph{Line 241}
\noindent\colorbox{gray!8}{\parbox{\dimexpr\linewidth-2\fboxsep\relax}{\raggedright\ttfamily\scriptsize                 ST\_MAC\_ACCUM: begin}}
\par\smallskip
This names one state in the control FSM. The numeric encoding is what appears in waveform traces, while the name explains the pipeline phase being executed.
\paragraph{Line 242}
\noindent\colorbox{gray!8}{\parbox{\dimexpr\linewidth-2\fboxsep\relax}{\raggedright\ttfamily\scriptsize                     accum <= accum + product\_r;}}
\par\smallskip
This is a nonblocking register or memory assignment. The right side is sampled now and the left side updates together with other clocked assignments at the end of the clock edge.
\paragraph{Line 243}
\noindent\colorbox{gray!8}{\parbox{\dimexpr\linewidth-2\fboxsep\relax}{\raggedright\ttfamily\scriptsize                     state <= ST\_CLASS\_DONE;}}
\par\smallskip
This schedules the next FSM state for the following clock edge. That transition is the hardware pipeline control path.
\paragraph{Line 244}
\noindent\colorbox{gray!8}{\parbox{\dimexpr\linewidth-2\fboxsep\relax}{\raggedright\ttfamily\scriptsize                 end}}
\par\smallskip
This SystemVerilog line contributes to the structural or sequential description of the accelerator. Its effect is resolved by XSim for verification and by Vivado during synthesis.
\paragraph{Line 245}
\noindent\colorbox{gray!8}{\parbox{\dimexpr\linewidth-2\fboxsep\relax}{\raggedright\ttfamily\scriptsize \emph{blank line}}}
\par\smallskip
This blank line is not hardware or software behavior. It separates nearby declarations or operations so the file remains readable and reviewable.
\paragraph{Line 246}
\noindent\colorbox{gray!8}{\parbox{\dimexpr\linewidth-2\fboxsep\relax}{\raggedright\ttfamily\scriptsize                 ST\_CLASS\_DONE: begin}}
\par\smallskip
This names one state in the control FSM. The numeric encoding is what appears in waveform traces, while the name explains the pipeline phase being executed.
\paragraph{Line 247}
\noindent\colorbox{gray!8}{\parbox{\dimexpr\linewidth-2\fboxsep\relax}{\raggedright\ttfamily\scriptsize                     if (feature\_idx == FEATURE\_DIM - 1) begin}}
\par\smallskip
This is a local compile-time or procedural assignment. In context it either defines a helper value or connects a temporary expression used by nearby hardware logic.
\paragraph{Line 248}
\noindent\colorbox{gray!8}{\parbox{\dimexpr\linewidth-2\fboxsep\relax}{\raggedright\ttfamily\scriptsize                         class\_score\_mem[class\_idx] <= accum;}}
\par\smallskip
This stores or reads the ten class scores for the current sample. Those scores are needed for argmax prediction, CSV testbench reporting, and the margin test used before an online update.
\paragraph{Line 249}
\noindent\colorbox{gray!8}{\parbox{\dimexpr\linewidth-2\fboxsep\relax}{\raggedright\ttfamily\scriptsize                         if (class\_idx == 0) begin}}
\par\smallskip
This is a local compile-time or procedural assignment. In context it either defines a helper value or connects a temporary expression used by nearby hardware logic.
\paragraph{Line 250}
\noindent\colorbox{gray!8}{\parbox{\dimexpr\linewidth-2\fboxsep\relax}{\raggedright\ttfamily\scriptsize                             predicted\_label <= 4'd0;}}
\par\smallskip
This is a nonblocking register or memory assignment. The right side is sampled now and the left side updates together with other clocked assignments at the end of the clock edge.
\paragraph{Line 251}
\noindent\colorbox{gray!8}{\parbox{\dimexpr\linewidth-2\fboxsep\relax}{\raggedright\ttfamily\scriptsize                             predicted\_score <= accum;}}
\par\smallskip
This is a nonblocking register or memory assignment. The right side is sampled now and the left side updates together with other clocked assignments at the end of the clock edge.
\paragraph{Line 252}
\noindent\colorbox{gray!8}{\parbox{\dimexpr\linewidth-2\fboxsep\relax}{\raggedright\ttfamily\scriptsize                         end else if (accum > predicted\_score) begin}}
\par\smallskip
This SystemVerilog line contributes to the structural or sequential description of the accelerator. Its effect is resolved by XSim for verification and by Vivado during synthesis.
\paragraph{Line 253}
\noindent\colorbox{gray!8}{\parbox{\dimexpr\linewidth-2\fboxsep\relax}{\raggedright\ttfamily\scriptsize                             predicted\_label <= 4'(class\_idx);}}
\par\smallskip
This is a nonblocking register or memory assignment. The right side is sampled now and the left side updates together with other clocked assignments at the end of the clock edge.
\paragraph{Line 254}
\noindent\colorbox{gray!8}{\parbox{\dimexpr\linewidth-2\fboxsep\relax}{\raggedright\ttfamily\scriptsize                             predicted\_score <= accum;}}
\par\smallskip
This is a nonblocking register or memory assignment. The right side is sampled now and the left side updates together with other clocked assignments at the end of the clock edge.
\paragraph{Line 255}
\noindent\colorbox{gray!8}{\parbox{\dimexpr\linewidth-2\fboxsep\relax}{\raggedright\ttfamily\scriptsize                         end}}
\par\smallskip
This SystemVerilog line contributes to the structural or sequential description of the accelerator. Its effect is resolved by XSim for verification and by Vivado during synthesis.
\paragraph{Line 256}
\noindent\colorbox{gray!8}{\parbox{\dimexpr\linewidth-2\fboxsep\relax}{\raggedright\ttfamily\scriptsize                         if (class\_idx == NUM\_CLASSES - 1) begin}}
\par\smallskip
This is a local compile-time or procedural assignment. In context it either defines a helper value or connects a temporary expression used by nearby hardware logic.
\paragraph{Line 257}
\noindent\colorbox{gray!8}{\parbox{\dimexpr\linewidth-2\fboxsep\relax}{\raggedright\ttfamily\scriptsize                             if (ENABLE\_ONCHIP\_LEARNING != 0) begin}}
\par\smallskip
This is a local compile-time or procedural assignment. In context it either defines a helper value or connects a temporary expression used by nearby hardware logic.
\paragraph{Line 258}
\noindent\colorbox{gray!8}{\parbox{\dimexpr\linewidth-2\fboxsep\relax}{\raggedright\ttfamily\scriptsize                                 state <= ST\_TRAIN\_CHECK;}}
\par\smallskip
This schedules the next FSM state for the following clock edge. That transition is the hardware pipeline control path.
\paragraph{Line 259}
\noindent\colorbox{gray!8}{\parbox{\dimexpr\linewidth-2\fboxsep\relax}{\raggedright\ttfamily\scriptsize                             end else begin}}
\par\smallskip
This SystemVerilog line contributes to the structural or sequential description of the accelerator. Its effect is resolved by XSim for verification and by Vivado during synthesis.
\paragraph{Line 260}
\noindent\colorbox{gray!8}{\parbox{\dimexpr\linewidth-2\fboxsep\relax}{\raggedright\ttfamily\scriptsize                                 state <= ST\_RESULT;}}
\par\smallskip
This schedules the next FSM state for the following clock edge. That transition is the hardware pipeline control path.
\paragraph{Line 261}
\noindent\colorbox{gray!8}{\parbox{\dimexpr\linewidth-2\fboxsep\relax}{\raggedright\ttfamily\scriptsize                             end}}
\par\smallskip
This SystemVerilog line contributes to the structural or sequential description of the accelerator. Its effect is resolved by XSim for verification and by Vivado during synthesis.
\paragraph{Line 262}
\noindent\colorbox{gray!8}{\parbox{\dimexpr\linewidth-2\fboxsep\relax}{\raggedright\ttfamily\scriptsize                         end else begin}}
\par\smallskip
This SystemVerilog line contributes to the structural or sequential description of the accelerator. Its effect is resolved by XSim for verification and by Vivado during synthesis.
\paragraph{Line 263}
\noindent\colorbox{gray!8}{\parbox{\dimexpr\linewidth-2\fboxsep\relax}{\raggedright\ttfamily\scriptsize                             class\_idx <= class\_idx + 1'b1;}}
\par\smallskip
This is a nonblocking register or memory assignment. The right side is sampled now and the left side updates together with other clocked assignments at the end of the clock edge.
\paragraph{Line 264}
\noindent\colorbox{gray!8}{\parbox{\dimexpr\linewidth-2\fboxsep\relax}{\raggedright\ttfamily\scriptsize                             state <= ST\_CLASS\_INIT;}}
\par\smallskip
This schedules the next FSM state for the following clock edge. That transition is the hardware pipeline control path.
\paragraph{Line 265}
\noindent\colorbox{gray!8}{\parbox{\dimexpr\linewidth-2\fboxsep\relax}{\raggedright\ttfamily\scriptsize                         end}}
\par\smallskip
This SystemVerilog line contributes to the structural or sequential description of the accelerator. Its effect is resolved by XSim for verification and by Vivado during synthesis.
\paragraph{Line 266}
\noindent\colorbox{gray!8}{\parbox{\dimexpr\linewidth-2\fboxsep\relax}{\raggedright\ttfamily\scriptsize                     end else begin}}
\par\smallskip
This SystemVerilog line contributes to the structural or sequential description of the accelerator. Its effect is resolved by XSim for verification and by Vivado during synthesis.
\paragraph{Line 267}
\noindent\colorbox{gray!8}{\parbox{\dimexpr\linewidth-2\fboxsep\relax}{\raggedright\ttfamily\scriptsize                         feature\_idx <= feature\_idx + 1'b1;}}
\par\smallskip
This is a nonblocking register or memory assignment. The right side is sampled now and the left side updates together with other clocked assignments at the end of the clock edge.
\paragraph{Line 268}
\noindent\colorbox{gray!8}{\parbox{\dimexpr\linewidth-2\fboxsep\relax}{\raggedright\ttfamily\scriptsize                         state <= ST\_MAC\_READ;}}
\par\smallskip
This schedules the next FSM state for the following clock edge. That transition is the hardware pipeline control path.
\paragraph{Line 269}
\noindent\colorbox{gray!8}{\parbox{\dimexpr\linewidth-2\fboxsep\relax}{\raggedright\ttfamily\scriptsize                     end}}
\par\smallskip
This SystemVerilog line contributes to the structural or sequential description of the accelerator. Its effect is resolved by XSim for verification and by Vivado during synthesis.
\paragraph{Line 270}
\noindent\colorbox{gray!8}{\parbox{\dimexpr\linewidth-2\fboxsep\relax}{\raggedright\ttfamily\scriptsize                 end}}
\par\smallskip
This SystemVerilog line contributes to the structural or sequential description of the accelerator. Its effect is resolved by XSim for verification and by Vivado during synthesis.
\paragraph{Line 271}
\noindent\colorbox{gray!8}{\parbox{\dimexpr\linewidth-2\fboxsep\relax}{\raggedright\ttfamily\scriptsize \emph{blank line}}}
\par\smallskip
This blank line is not hardware or software behavior. It separates nearby declarations or operations so the file remains readable and reviewable.
\paragraph{Line 272}
\noindent\colorbox{gray!8}{\parbox{\dimexpr\linewidth-2\fboxsep\relax}{\raggedright\ttfamily\scriptsize                 ST\_TRAIN\_CHECK: begin}}
\par\smallskip
This names one state in the control FSM. The numeric encoding is what appears in waveform traces, while the name explains the pipeline phase being executed.
\paragraph{Line 273}
\noindent\colorbox{gray!8}{\parbox{\dimexpr\linewidth-2\fboxsep\relax}{\raggedright\ttfamily\scriptsize                     train\_update\_valid <= 1'b1;}}
\par\smallskip
This is a nonblocking register or memory assignment. The right side is sampled now and the left side updates together with other clocked assignments at the end of the clock edge.
\paragraph{Line 274}
\noindent\colorbox{gray!8}{\parbox{\dimexpr\linewidth-2\fboxsep\relax}{\raggedright\ttfamily\scriptsize                     if (train\_enable}}
\par\smallskip
This SystemVerilog line contributes to the structural or sequential description of the accelerator. Its effect is resolved by XSim for verification and by Vivado during synthesis.
\paragraph{Line 275}
\noindent\colorbox{gray!8}{\parbox{\dimexpr\linewidth-2\fboxsep\relax}{\raggedright\ttfamily\scriptsize                             \&\& train\_label < NUM\_CLASSES}}
\par\smallskip
This SystemVerilog line contributes to the structural or sequential description of the accelerator. Its effect is resolved by XSim for verification and by Vivado during synthesis.
\paragraph{Line 276}
\noindent\colorbox{gray!8}{\parbox{\dimexpr\linewidth-2\fboxsep\relax}{\raggedright\ttfamily\scriptsize                             \&\& train\_label != predicted\_label}}
\par\smallskip
This is a local compile-time or procedural assignment. In context it either defines a helper value or connects a temporary expression used by nearby hardware logic.
\paragraph{Line 277}
\noindent\colorbox{gray!8}{\parbox{\dimexpr\linewidth-2\fboxsep\relax}{\raggedright\ttfamily\scriptsize                             \&\& (predicted\_score - class\_score\_mem[train\_label]) <= LEARN\_MARGIN\_Q) begin}}
\par\smallskip
This stores or reads the ten class scores for the current sample. Those scores are needed for argmax prediction, CSV testbench reporting, and the margin test used before an online update.
\paragraph{Line 278}
\noindent\colorbox{gray!8}{\parbox{\dimexpr\linewidth-2\fboxsep\relax}{\raggedright\ttfamily\scriptsize                         feature\_idx <= '0;}}
\par\smallskip
This is a nonblocking register or memory assignment. The right side is sampled now and the left side updates together with other clocked assignments at the end of the clock edge.
\paragraph{Line 279}
\noindent\colorbox{gray!8}{\parbox{\dimexpr\linewidth-2\fboxsep\relax}{\raggedright\ttfamily\scriptsize                         train\_label\_r <= train\_label;}}
\par\smallskip
This is a nonblocking register or memory assignment. The right side is sampled now and the left side updates together with other clocked assignments at the end of the clock edge.
\paragraph{Line 280}
\noindent\colorbox{gray!8}{\parbox{\dimexpr\linewidth-2\fboxsep\relax}{\raggedright\ttfamily\scriptsize                         train\_pred\_label\_r <= predicted\_label;}}
\par\smallskip
This is a nonblocking register or memory assignment. The right side is sampled now and the left side updates together with other clocked assignments at the end of the clock edge.
\paragraph{Line 281}
\noindent\colorbox{gray!8}{\parbox{\dimexpr\linewidth-2\fboxsep\relax}{\raggedright\ttfamily\scriptsize                         train\_update\_applied <= 1'b1;}}
\par\smallskip
This is a nonblocking register or memory assignment. The right side is sampled now and the left side updates together with other clocked assignments at the end of the clock edge.
\paragraph{Line 282}
\noindent\colorbox{gray!8}{\parbox{\dimexpr\linewidth-2\fboxsep\relax}{\raggedright\ttfamily\scriptsize                         update\_count <= update\_count + 1'b1;}}
\par\smallskip
This is a nonblocking register or memory assignment. The right side is sampled now and the left side updates together with other clocked assignments at the end of the clock edge.
\paragraph{Line 283}
\noindent\colorbox{gray!8}{\parbox{\dimexpr\linewidth-2\fboxsep\relax}{\raggedright\ttfamily\scriptsize                         state <= ST\_TRAIN\_READ\_TRUE;}}
\par\smallskip
This schedules the next FSM state for the following clock edge. That transition is the hardware pipeline control path.
\paragraph{Line 284}
\noindent\colorbox{gray!8}{\parbox{\dimexpr\linewidth-2\fboxsep\relax}{\raggedright\ttfamily\scriptsize                     end else begin}}
\par\smallskip
This SystemVerilog line contributes to the structural or sequential description of the accelerator. Its effect is resolved by XSim for verification and by Vivado during synthesis.
\paragraph{Line 285}
\noindent\colorbox{gray!8}{\parbox{\dimexpr\linewidth-2\fboxsep\relax}{\raggedright\ttfamily\scriptsize                         state <= ST\_RESULT;}}
\par\smallskip
This schedules the next FSM state for the following clock edge. That transition is the hardware pipeline control path.
\paragraph{Line 286}
\noindent\colorbox{gray!8}{\parbox{\dimexpr\linewidth-2\fboxsep\relax}{\raggedright\ttfamily\scriptsize                     end}}
\par\smallskip
This SystemVerilog line contributes to the structural or sequential description of the accelerator. Its effect is resolved by XSim for verification and by Vivado during synthesis.
\paragraph{Line 287}
\noindent\colorbox{gray!8}{\parbox{\dimexpr\linewidth-2\fboxsep\relax}{\raggedright\ttfamily\scriptsize                 end}}
\par\smallskip
This SystemVerilog line contributes to the structural or sequential description of the accelerator. Its effect is resolved by XSim for verification and by Vivado during synthesis.
\paragraph{Line 288}
\noindent\colorbox{gray!8}{\parbox{\dimexpr\linewidth-2\fboxsep\relax}{\raggedright\ttfamily\scriptsize \emph{blank line}}}
\par\smallskip
This blank line is not hardware or software behavior. It separates nearby declarations or operations so the file remains readable and reviewable.
\paragraph{Line 289}
\noindent\colorbox{gray!8}{\parbox{\dimexpr\linewidth-2\fboxsep\relax}{\raggedright\ttfamily\scriptsize                 ST\_TRAIN\_READ\_TRUE: begin}}
\par\smallskip
This names one state in the control FSM. The numeric encoding is what appears in waveform traces, while the name explains the pipeline phase being executed.
\paragraph{Line 290}
\noindent\colorbox{gray!8}{\parbox{\dimexpr\linewidth-2\fboxsep\relax}{\raggedright\ttfamily\scriptsize                     if (skip\_feature(feature\_mem[feature\_idx])) begin}}
\par\smallskip
This line is part of the combinational decision that detects low-magnitude features. It is the hardware equivalent of asking whether a small activation should be silent for this inference/update pass.
\paragraph{Line 291}
\noindent\colorbox{gray!8}{\parbox{\dimexpr\linewidth-2\fboxsep\relax}{\raggedright\ttfamily\scriptsize                         train\_skip\_count <= train\_skip\_count + 1'b1;}}
\par\smallskip
This line exposes or updates a hardware activity counter. The counter makes the feature gate visible in waveforms and CSV traces, so the power-saving opportunity is observable from RTL rather than only from Python statistics.
\paragraph{Line 292}
\noindent\colorbox{gray!8}{\parbox{\dimexpr\linewidth-2\fboxsep\relax}{\raggedright\ttfamily\scriptsize                         if (feature\_idx == FEATURE\_DIM - 1) begin}}
\par\smallskip
This is a local compile-time or procedural assignment. In context it either defines a helper value or connects a temporary expression used by nearby hardware logic.
\paragraph{Line 293}
\noindent\colorbox{gray!8}{\parbox{\dimexpr\linewidth-2\fboxsep\relax}{\raggedright\ttfamily\scriptsize                             state <= ST\_TRAIN\_BIAS\_TRUE;}}
\par\smallskip
This schedules the next FSM state for the following clock edge. That transition is the hardware pipeline control path.
\paragraph{Line 294}
\noindent\colorbox{gray!8}{\parbox{\dimexpr\linewidth-2\fboxsep\relax}{\raggedright\ttfamily\scriptsize                         end else begin}}
\par\smallskip
This SystemVerilog line contributes to the structural or sequential description of the accelerator. Its effect is resolved by XSim for verification and by Vivado during synthesis.
\paragraph{Line 295}
\noindent\colorbox{gray!8}{\parbox{\dimexpr\linewidth-2\fboxsep\relax}{\raggedright\ttfamily\scriptsize                             feature\_idx <= feature\_idx + 1'b1;}}
\par\smallskip
This is a nonblocking register or memory assignment. The right side is sampled now and the left side updates together with other clocked assignments at the end of the clock edge.
\paragraph{Line 296}
\noindent\colorbox{gray!8}{\parbox{\dimexpr\linewidth-2\fboxsep\relax}{\raggedright\ttfamily\scriptsize                             state <= ST\_TRAIN\_READ\_TRUE;}}
\par\smallskip
This schedules the next FSM state for the following clock edge. That transition is the hardware pipeline control path.
\paragraph{Line 297}
\noindent\colorbox{gray!8}{\parbox{\dimexpr\linewidth-2\fboxsep\relax}{\raggedright\ttfamily\scriptsize                         end}}
\par\smallskip
This SystemVerilog line contributes to the structural or sequential description of the accelerator. Its effect is resolved by XSim for verification and by Vivado during synthesis.
\paragraph{Line 298}
\noindent\colorbox{gray!8}{\parbox{\dimexpr\linewidth-2\fboxsep\relax}{\raggedright\ttfamily\scriptsize                     end else begin}}
\par\smallskip
This SystemVerilog line contributes to the structural or sequential description of the accelerator. Its effect is resolved by XSim for verification and by Vivado during synthesis.
\paragraph{Line 299}
\noindent\colorbox{gray!8}{\parbox{\dimexpr\linewidth-2\fboxsep\relax}{\raggedright\ttfamily\scriptsize                         state <= ST\_TRAIN\_CALC\_TRUE;}}
\par\smallskip
This schedules the next FSM state for the following clock edge. That transition is the hardware pipeline control path.
\paragraph{Line 300}
\noindent\colorbox{gray!8}{\parbox{\dimexpr\linewidth-2\fboxsep\relax}{\raggedright\ttfamily\scriptsize                     end}}
\par\smallskip
This SystemVerilog line contributes to the structural or sequential description of the accelerator. Its effect is resolved by XSim for verification and by Vivado during synthesis.
\paragraph{Line 301}
\noindent\colorbox{gray!8}{\parbox{\dimexpr\linewidth-2\fboxsep\relax}{\raggedright\ttfamily\scriptsize                 end}}
\par\smallskip
This SystemVerilog line contributes to the structural or sequential description of the accelerator. Its effect is resolved by XSim for verification and by Vivado during synthesis.
\paragraph{Line 302}
\noindent\colorbox{gray!8}{\parbox{\dimexpr\linewidth-2\fboxsep\relax}{\raggedright\ttfamily\scriptsize \emph{blank line}}}
\par\smallskip
This blank line is not hardware or software behavior. It separates nearby declarations or operations so the file remains readable and reviewable.
\paragraph{Line 303}
\noindent\colorbox{gray!8}{\parbox{\dimexpr\linewidth-2\fboxsep\relax}{\raggedright\ttfamily\scriptsize                 ST\_TRAIN\_CALC\_TRUE: begin}}
\par\smallskip
This names one state in the control FSM. The numeric encoding is what appears in waveform traces, while the name explains the pipeline phase being executed.
\paragraph{Line 304}
\noindent\colorbox{gray!8}{\parbox{\dimexpr\linewidth-2\fboxsep\relax}{\raggedright\ttfamily\scriptsize                     train\_write\_weight\_r <= sat\_weight(}}
\par\smallskip
This line is part of signed saturation for 16-bit weights. Saturation prevents online learning from wrapping around when an update would exceed the representable fixed-point range.
\paragraph{Line 305}
\noindent\colorbox{gray!8}{\parbox{\dimexpr\linewidth-2\fboxsep\relax}{\raggedright\ttfamily\scriptsize                         \{\{(ACC\_BITS-WEIGHT\_BITS)\{train\_true\_weight\_r[WEIGHT\_BITS-1]\}\}, train\_true\_weight\_r\}}}
\par\smallskip
This SystemVerilog line contributes to the structural or sequential description of the accelerator. Its effect is resolved by XSim for verification and by Vivado during synthesis.
\paragraph{Line 306}
\noindent\colorbox{gray!8}{\parbox{\dimexpr\linewidth-2\fboxsep\relax}{\raggedright\ttfamily\scriptsize                         + learn\_delta(feature\_eff\_r)}}
\par\smallskip
This line implements the hardware learning step size. The feature is sign-extended and arithmetically shifted so positive and negative features update weights symmetrically.
\paragraph{Line 307}
\noindent\colorbox{gray!8}{\parbox{\dimexpr\linewidth-2\fboxsep\relax}{\raggedright\ttfamily\scriptsize                     );}}
\par\smallskip
This SystemVerilog line contributes to the structural or sequential description of the accelerator. Its effect is resolved by XSim for verification and by Vivado during synthesis.
\paragraph{Line 308}
\noindent\colorbox{gray!8}{\parbox{\dimexpr\linewidth-2\fboxsep\relax}{\raggedright\ttfamily\scriptsize                     state <= ST\_TRAIN\_WRITE\_TRUE;}}
\par\smallskip
This schedules the next FSM state for the following clock edge. That transition is the hardware pipeline control path.
\paragraph{Line 309}
\noindent\colorbox{gray!8}{\parbox{\dimexpr\linewidth-2\fboxsep\relax}{\raggedright\ttfamily\scriptsize                 end}}
\par\smallskip
This SystemVerilog line contributes to the structural or sequential description of the accelerator. Its effect is resolved by XSim for verification and by Vivado during synthesis.
\paragraph{Line 310}
\noindent\colorbox{gray!8}{\parbox{\dimexpr\linewidth-2\fboxsep\relax}{\raggedright\ttfamily\scriptsize \emph{blank line}}}
\par\smallskip
This blank line is not hardware or software behavior. It separates nearby declarations or operations so the file remains readable and reviewable.
\paragraph{Line 311}
\noindent\colorbox{gray!8}{\parbox{\dimexpr\linewidth-2\fboxsep\relax}{\raggedright\ttfamily\scriptsize                 ST\_TRAIN\_WRITE\_TRUE: begin}}
\par\smallskip
This names one state in the control FSM. The numeric encoding is what appears in waveform traces, while the name explains the pipeline phase being executed.
\paragraph{Line 312}
\noindent\colorbox{gray!8}{\parbox{\dimexpr\linewidth-2\fboxsep\relax}{\raggedright\ttfamily\scriptsize                     state <= ST\_TRAIN\_READ\_PRED;}}
\par\smallskip
This schedules the next FSM state for the following clock edge. That transition is the hardware pipeline control path.
\paragraph{Line 313}
\noindent\colorbox{gray!8}{\parbox{\dimexpr\linewidth-2\fboxsep\relax}{\raggedright\ttfamily\scriptsize                 end}}
\par\smallskip
This SystemVerilog line contributes to the structural or sequential description of the accelerator. Its effect is resolved by XSim for verification and by Vivado during synthesis.
\paragraph{Line 314}
\noindent\colorbox{gray!8}{\parbox{\dimexpr\linewidth-2\fboxsep\relax}{\raggedright\ttfamily\scriptsize \emph{blank line}}}
\par\smallskip
This blank line is not hardware or software behavior. It separates nearby declarations or operations so the file remains readable and reviewable.
\paragraph{Line 315}
\noindent\colorbox{gray!8}{\parbox{\dimexpr\linewidth-2\fboxsep\relax}{\raggedright\ttfamily\scriptsize                 ST\_TRAIN\_READ\_PRED: begin}}
\par\smallskip
This names one state in the control FSM. The numeric encoding is what appears in waveform traces, while the name explains the pipeline phase being executed.
\paragraph{Line 316}
\noindent\colorbox{gray!8}{\parbox{\dimexpr\linewidth-2\fboxsep\relax}{\raggedright\ttfamily\scriptsize                     state <= ST\_TRAIN\_CALC\_PRED;}}
\par\smallskip
This schedules the next FSM state for the following clock edge. That transition is the hardware pipeline control path.
\paragraph{Line 317}
\noindent\colorbox{gray!8}{\parbox{\dimexpr\linewidth-2\fboxsep\relax}{\raggedright\ttfamily\scriptsize                 end}}
\par\smallskip
This SystemVerilog line contributes to the structural or sequential description of the accelerator. Its effect is resolved by XSim for verification and by Vivado during synthesis.
\paragraph{Line 318}
\noindent\colorbox{gray!8}{\parbox{\dimexpr\linewidth-2\fboxsep\relax}{\raggedright\ttfamily\scriptsize \emph{blank line}}}
\par\smallskip
This blank line is not hardware or software behavior. It separates nearby declarations or operations so the file remains readable and reviewable.
\paragraph{Line 319}
\noindent\colorbox{gray!8}{\parbox{\dimexpr\linewidth-2\fboxsep\relax}{\raggedright\ttfamily\scriptsize                 ST\_TRAIN\_CALC\_PRED: begin}}
\par\smallskip
This names one state in the control FSM. The numeric encoding is what appears in waveform traces, while the name explains the pipeline phase being executed.
\paragraph{Line 320}
\noindent\colorbox{gray!8}{\parbox{\dimexpr\linewidth-2\fboxsep\relax}{\raggedright\ttfamily\scriptsize                     train\_write\_weight\_r <= sat\_weight(}}
\par\smallskip
This line is part of signed saturation for 16-bit weights. Saturation prevents online learning from wrapping around when an update would exceed the representable fixed-point range.
\paragraph{Line 321}
\noindent\colorbox{gray!8}{\parbox{\dimexpr\linewidth-2\fboxsep\relax}{\raggedright\ttfamily\scriptsize                         \{\{(ACC\_BITS-WEIGHT\_BITS)\{train\_pred\_weight\_r[WEIGHT\_BITS-1]\}\}, train\_pred\_weight\_r\}}}
\par\smallskip
This SystemVerilog line contributes to the structural or sequential description of the accelerator. Its effect is resolved by XSim for verification and by Vivado during synthesis.
\paragraph{Line 322}
\noindent\colorbox{gray!8}{\parbox{\dimexpr\linewidth-2\fboxsep\relax}{\raggedright\ttfamily\scriptsize                         - learn\_delta(feature\_eff\_r)}}
\par\smallskip
This line implements the hardware learning step size. The feature is sign-extended and arithmetically shifted so positive and negative features update weights symmetrically.
\paragraph{Line 323}
\noindent\colorbox{gray!8}{\parbox{\dimexpr\linewidth-2\fboxsep\relax}{\raggedright\ttfamily\scriptsize                     );}}
\par\smallskip
This SystemVerilog line contributes to the structural or sequential description of the accelerator. Its effect is resolved by XSim for verification and by Vivado during synthesis.
\paragraph{Line 324}
\noindent\colorbox{gray!8}{\parbox{\dimexpr\linewidth-2\fboxsep\relax}{\raggedright\ttfamily\scriptsize                     state <= ST\_TRAIN\_WRITE\_PRED;}}
\par\smallskip
This schedules the next FSM state for the following clock edge. That transition is the hardware pipeline control path.
\paragraph{Line 325}
\noindent\colorbox{gray!8}{\parbox{\dimexpr\linewidth-2\fboxsep\relax}{\raggedright\ttfamily\scriptsize                 end}}
\par\smallskip
This SystemVerilog line contributes to the structural or sequential description of the accelerator. Its effect is resolved by XSim for verification and by Vivado during synthesis.
\paragraph{Line 326}
\noindent\colorbox{gray!8}{\parbox{\dimexpr\linewidth-2\fboxsep\relax}{\raggedright\ttfamily\scriptsize \emph{blank line}}}
\par\smallskip
This blank line is not hardware or software behavior. It separates nearby declarations or operations so the file remains readable and reviewable.
\paragraph{Line 327}
\noindent\colorbox{gray!8}{\parbox{\dimexpr\linewidth-2\fboxsep\relax}{\raggedright\ttfamily\scriptsize                 ST\_TRAIN\_WRITE\_PRED: begin}}
\par\smallskip
This names one state in the control FSM. The numeric encoding is what appears in waveform traces, while the name explains the pipeline phase being executed.
\paragraph{Line 328}
\noindent\colorbox{gray!8}{\parbox{\dimexpr\linewidth-2\fboxsep\relax}{\raggedright\ttfamily\scriptsize                     if (feature\_idx == FEATURE\_DIM - 1) begin}}
\par\smallskip
This is a local compile-time or procedural assignment. In context it either defines a helper value or connects a temporary expression used by nearby hardware logic.
\paragraph{Line 329}
\noindent\colorbox{gray!8}{\parbox{\dimexpr\linewidth-2\fboxsep\relax}{\raggedright\ttfamily\scriptsize                         state <= ST\_TRAIN\_BIAS\_TRUE;}}
\par\smallskip
This schedules the next FSM state for the following clock edge. That transition is the hardware pipeline control path.
\paragraph{Line 330}
\noindent\colorbox{gray!8}{\parbox{\dimexpr\linewidth-2\fboxsep\relax}{\raggedright\ttfamily\scriptsize                     end else begin}}
\par\smallskip
This SystemVerilog line contributes to the structural or sequential description of the accelerator. Its effect is resolved by XSim for verification and by Vivado during synthesis.
\paragraph{Line 331}
\noindent\colorbox{gray!8}{\parbox{\dimexpr\linewidth-2\fboxsep\relax}{\raggedright\ttfamily\scriptsize                         feature\_idx <= feature\_idx + 1'b1;}}
\par\smallskip
This is a nonblocking register or memory assignment. The right side is sampled now and the left side updates together with other clocked assignments at the end of the clock edge.
\paragraph{Line 332}
\noindent\colorbox{gray!8}{\parbox{\dimexpr\linewidth-2\fboxsep\relax}{\raggedright\ttfamily\scriptsize                         state <= ST\_TRAIN\_READ\_TRUE;}}
\par\smallskip
This schedules the next FSM state for the following clock edge. That transition is the hardware pipeline control path.
\paragraph{Line 333}
\noindent\colorbox{gray!8}{\parbox{\dimexpr\linewidth-2\fboxsep\relax}{\raggedright\ttfamily\scriptsize                     end}}
\par\smallskip
This SystemVerilog line contributes to the structural or sequential description of the accelerator. Its effect is resolved by XSim for verification and by Vivado during synthesis.
\paragraph{Line 334}
\noindent\colorbox{gray!8}{\parbox{\dimexpr\linewidth-2\fboxsep\relax}{\raggedright\ttfamily\scriptsize                 end}}
\par\smallskip
This SystemVerilog line contributes to the structural or sequential description of the accelerator. Its effect is resolved by XSim for verification and by Vivado during synthesis.
\paragraph{Line 335}
\noindent\colorbox{gray!8}{\parbox{\dimexpr\linewidth-2\fboxsep\relax}{\raggedright\ttfamily\scriptsize \emph{blank line}}}
\par\smallskip
This blank line is not hardware or software behavior. It separates nearby declarations or operations so the file remains readable and reviewable.
\paragraph{Line 336}
\noindent\colorbox{gray!8}{\parbox{\dimexpr\linewidth-2\fboxsep\relax}{\raggedright\ttfamily\scriptsize                 ST\_TRAIN\_BIAS\_TRUE: begin}}
\par\smallskip
This names one state in the control FSM. The numeric encoding is what appears in waveform traces, while the name explains the pipeline phase being executed.
\paragraph{Line 337}
\noindent\colorbox{gray!8}{\parbox{\dimexpr\linewidth-2\fboxsep\relax}{\raggedright\ttfamily\scriptsize                     state <= ST\_TRAIN\_BIAS\_PRED;}}
\par\smallskip
This schedules the next FSM state for the following clock edge. That transition is the hardware pipeline control path.
\paragraph{Line 338}
\noindent\colorbox{gray!8}{\parbox{\dimexpr\linewidth-2\fboxsep\relax}{\raggedright\ttfamily\scriptsize                 end}}
\par\smallskip
This SystemVerilog line contributes to the structural or sequential description of the accelerator. Its effect is resolved by XSim for verification and by Vivado during synthesis.
\paragraph{Line 339}
\noindent\colorbox{gray!8}{\parbox{\dimexpr\linewidth-2\fboxsep\relax}{\raggedright\ttfamily\scriptsize \emph{blank line}}}
\par\smallskip
This blank line is not hardware or software behavior. It separates nearby declarations or operations so the file remains readable and reviewable.
\paragraph{Line 340}
\noindent\colorbox{gray!8}{\parbox{\dimexpr\linewidth-2\fboxsep\relax}{\raggedright\ttfamily\scriptsize                 ST\_TRAIN\_BIAS\_PRED: begin}}
\par\smallskip
This names one state in the control FSM. The numeric encoding is what appears in waveform traces, while the name explains the pipeline phase being executed.
\paragraph{Line 341}
\noindent\colorbox{gray!8}{\parbox{\dimexpr\linewidth-2\fboxsep\relax}{\raggedright\ttfamily\scriptsize                     state <= ST\_RESULT;}}
\par\smallskip
This schedules the next FSM state for the following clock edge. That transition is the hardware pipeline control path.
\paragraph{Line 342}
\noindent\colorbox{gray!8}{\parbox{\dimexpr\linewidth-2\fboxsep\relax}{\raggedright\ttfamily\scriptsize                 end}}
\par\smallskip
This SystemVerilog line contributes to the structural or sequential description of the accelerator. Its effect is resolved by XSim for verification and by Vivado during synthesis.
\paragraph{Line 343}
\noindent\colorbox{gray!8}{\parbox{\dimexpr\linewidth-2\fboxsep\relax}{\raggedright\ttfamily\scriptsize \emph{blank line}}}
\par\smallskip
This blank line is not hardware or software behavior. It separates nearby declarations or operations so the file remains readable and reviewable.
\paragraph{Line 344}
\noindent\colorbox{gray!8}{\parbox{\dimexpr\linewidth-2\fboxsep\relax}{\raggedright\ttfamily\scriptsize                 ST\_RESULT: begin}}
\par\smallskip
This names one state in the control FSM. The numeric encoding is what appears in waveform traces, while the name explains the pipeline phase being executed.
\paragraph{Line 345}
\noindent\colorbox{gray!8}{\parbox{\dimexpr\linewidth-2\fboxsep\relax}{\raggedright\ttfamily\scriptsize                     result\_valid <= 1'b1;}}
\par\smallskip
This is a nonblocking register or memory assignment. The right side is sampled now and the left side updates together with other clocked assignments at the end of the clock edge.
\paragraph{Line 346}
\noindent\colorbox{gray!8}{\parbox{\dimexpr\linewidth-2\fboxsep\relax}{\raggedright\ttfamily\scriptsize                     state <= ST\_IDLE;}}
\par\smallskip
This schedules the next FSM state for the following clock edge. That transition is the hardware pipeline control path.
\paragraph{Line 347}
\noindent\colorbox{gray!8}{\parbox{\dimexpr\linewidth-2\fboxsep\relax}{\raggedright\ttfamily\scriptsize                 end}}
\par\smallskip
This SystemVerilog line contributes to the structural or sequential description of the accelerator. Its effect is resolved by XSim for verification and by Vivado during synthesis.
\paragraph{Line 348}
\noindent\colorbox{gray!8}{\parbox{\dimexpr\linewidth-2\fboxsep\relax}{\raggedright\ttfamily\scriptsize \emph{blank line}}}
\par\smallskip
This blank line is not hardware or software behavior. It separates nearby declarations or operations so the file remains readable and reviewable.
\paragraph{Line 349}
\noindent\colorbox{gray!8}{\parbox{\dimexpr\linewidth-2\fboxsep\relax}{\raggedright\ttfamily\scriptsize                 default: begin}}
\par\smallskip
This SystemVerilog line contributes to the structural or sequential description of the accelerator. Its effect is resolved by XSim for verification and by Vivado during synthesis.
\paragraph{Line 350}
\noindent\colorbox{gray!8}{\parbox{\dimexpr\linewidth-2\fboxsep\relax}{\raggedright\ttfamily\scriptsize                     state <= ST\_IDLE;}}
\par\smallskip
This schedules the next FSM state for the following clock edge. That transition is the hardware pipeline control path.
\paragraph{Line 351}
\noindent\colorbox{gray!8}{\parbox{\dimexpr\linewidth-2\fboxsep\relax}{\raggedright\ttfamily\scriptsize                 end}}
\par\smallskip
This SystemVerilog line contributes to the structural or sequential description of the accelerator. Its effect is resolved by XSim for verification and by Vivado during synthesis.
\paragraph{Line 352}
\noindent\colorbox{gray!8}{\parbox{\dimexpr\linewidth-2\fboxsep\relax}{\raggedright\ttfamily\scriptsize             endcase}}
\par\smallskip
This SystemVerilog line contributes to the structural or sequential description of the accelerator. Its effect is resolved by XSim for verification and by Vivado during synthesis.
\paragraph{Line 353}
\noindent\colorbox{gray!8}{\parbox{\dimexpr\linewidth-2\fboxsep\relax}{\raggedright\ttfamily\scriptsize         end}}
\par\smallskip
This SystemVerilog line contributes to the structural or sequential description of the accelerator. Its effect is resolved by XSim for verification and by Vivado during synthesis.
\paragraph{Line 354}
\noindent\colorbox{gray!8}{\parbox{\dimexpr\linewidth-2\fboxsep\relax}{\raggedright\ttfamily\scriptsize     end}}
\par\smallskip
This SystemVerilog line contributes to the structural or sequential description of the accelerator. Its effect is resolved by XSim for verification and by Vivado during synthesis.
\paragraph{Line 355}
\noindent\colorbox{gray!8}{\parbox{\dimexpr\linewidth-2\fboxsep\relax}{\raggedright\ttfamily\scriptsize \emph{blank line}}}
\par\smallskip
This blank line is not hardware or software behavior. It separates nearby declarations or operations so the file remains readable and reviewable.
\paragraph{Line 356}
\noindent\colorbox{gray!8}{\parbox{\dimexpr\linewidth-2\fboxsep\relax}{\raggedright\ttfamily\scriptsize     always\_ff @(posedge clk) begin}}
\par\smallskip
This begins clocked sequential logic. Registers and RAM writes inside this block update only on a clock edge, which is what makes the design synthesizable and waveform-predictable.
\paragraph{Line 357}
\noindent\colorbox{gray!8}{\parbox{\dimexpr\linewidth-2\fboxsep\relax}{\raggedright\ttfamily\scriptsize         unique case (state)}}
\par\smallskip
This starts a state-machine branch table. The unique qualifier tells tools that one legal state should match at a time, improving diagnostics and optimization.
\paragraph{Line 358}
\noindent\colorbox{gray!8}{\parbox{\dimexpr\linewidth-2\fboxsep\relax}{\raggedright\ttfamily\scriptsize             ST\_LOAD: begin}}
\par\smallskip
This names one state in the control FSM. The numeric encoding is what appears in waveform traces, while the name explains the pipeline phase being executed.
\paragraph{Line 359}
\noindent\colorbox{gray!8}{\parbox{\dimexpr\linewidth-2\fboxsep\relax}{\raggedright\ttfamily\scriptsize                 if (feature\_valid \&\& feature\_ready) begin}}
\par\smallskip
This SystemVerilog line contributes to the structural or sequential description of the accelerator. Its effect is resolved by XSim for verification and by Vivado during synthesis.
\paragraph{Line 360}
\noindent\colorbox{gray!8}{\parbox{\dimexpr\linewidth-2\fboxsep\relax}{\raggedright\ttfamily\scriptsize                     feature\_mem[feature\_idx] <= feature\_value;}}
\par\smallskip
This line touches the feature memory that holds one 384-dimensional fixed-point Spikingformer feature vector for the current image. Keeping the feature vector local avoids repeatedly streaming the same values for every class score.
\paragraph{Line 361}
\noindent\colorbox{gray!8}{\parbox{\dimexpr\linewidth-2\fboxsep\relax}{\raggedright\ttfamily\scriptsize                 end}}
\par\smallskip
This SystemVerilog line contributes to the structural or sequential description of the accelerator. Its effect is resolved by XSim for verification and by Vivado during synthesis.
\paragraph{Line 362}
\noindent\colorbox{gray!8}{\parbox{\dimexpr\linewidth-2\fboxsep\relax}{\raggedright\ttfamily\scriptsize             end}}
\par\smallskip
This SystemVerilog line contributes to the structural or sequential description of the accelerator. Its effect is resolved by XSim for verification and by Vivado during synthesis.
\paragraph{Line 363}
\noindent\colorbox{gray!8}{\parbox{\dimexpr\linewidth-2\fboxsep\relax}{\raggedright\ttfamily\scriptsize \emph{blank line}}}
\par\smallskip
This blank line is not hardware or software behavior. It separates nearby declarations or operations so the file remains readable and reviewable.
\paragraph{Line 364}
\noindent\colorbox{gray!8}{\parbox{\dimexpr\linewidth-2\fboxsep\relax}{\raggedright\ttfamily\scriptsize             ST\_MAC\_READ: begin}}
\par\smallskip
This names one state in the control FSM. The numeric encoding is what appears in waveform traces, while the name explains the pipeline phase being executed.
\paragraph{Line 365}
\noindent\colorbox{gray!8}{\parbox{\dimexpr\linewidth-2\fboxsep\relax}{\raggedright\ttfamily\scriptsize                 feature\_eff\_r <= effective\_feature(feature\_mem[feature\_idx]);}}
\par\smallskip
This line is part of the denoised feature path. The returned value is either the original signed feature or exactly zero, and that registered value is what feeds both the DSP multiply and the online-learning delta.
\paragraph{Line 366}
\noindent\colorbox{gray!8}{\parbox{\dimexpr\linewidth-2\fboxsep\relax}{\raggedright\ttfamily\scriptsize                 if (!skip\_feature(feature\_mem[feature\_idx])) begin}}
\par\smallskip
This line is part of the combinational decision that detects low-magnitude features. It is the hardware equivalent of asking whether a small activation should be silent for this inference/update pass.
\paragraph{Line 367}
\noindent\colorbox{gray!8}{\parbox{\dimexpr\linewidth-2\fboxsep\relax}{\raggedright\ttfamily\scriptsize                     weight\_r <= weight\_mem[weight\_addr(int'(class\_idx), int'(feature\_idx))];}}
\par\smallskip
This line touches the trained classifier weight store. In the final RTL the large 10 x 384 int16 matrix is inferred as true dual-port BRAM, so weights remain on chip and can also be updated by the learning state machine.
\paragraph{Line 368}
\noindent\colorbox{gray!8}{\parbox{\dimexpr\linewidth-2\fboxsep\relax}{\raggedright\ttfamily\scriptsize                 end}}
\par\smallskip
This SystemVerilog line contributes to the structural or sequential description of the accelerator. Its effect is resolved by XSim for verification and by Vivado during synthesis.
\paragraph{Line 369}
\noindent\colorbox{gray!8}{\parbox{\dimexpr\linewidth-2\fboxsep\relax}{\raggedright\ttfamily\scriptsize             end}}
\par\smallskip
This SystemVerilog line contributes to the structural or sequential description of the accelerator. Its effect is resolved by XSim for verification and by Vivado during synthesis.
\paragraph{Line 370}
\noindent\colorbox{gray!8}{\parbox{\dimexpr\linewidth-2\fboxsep\relax}{\raggedright\ttfamily\scriptsize \emph{blank line}}}
\par\smallskip
This blank line is not hardware or software behavior. It separates nearby declarations or operations so the file remains readable and reviewable.
\paragraph{Line 371}
\noindent\colorbox{gray!8}{\parbox{\dimexpr\linewidth-2\fboxsep\relax}{\raggedright\ttfamily\scriptsize             ST\_TRAIN\_READ\_TRUE: begin}}
\par\smallskip
This names one state in the control FSM. The numeric encoding is what appears in waveform traces, while the name explains the pipeline phase being executed.
\paragraph{Line 372}
\noindent\colorbox{gray!8}{\parbox{\dimexpr\linewidth-2\fboxsep\relax}{\raggedright\ttfamily\scriptsize                 feature\_eff\_r <= effective\_feature(feature\_mem[feature\_idx]);}}
\par\smallskip
This line is part of the denoised feature path. The returned value is either the original signed feature or exactly zero, and that registered value is what feeds both the DSP multiply and the online-learning delta.
\paragraph{Line 373}
\noindent\colorbox{gray!8}{\parbox{\dimexpr\linewidth-2\fboxsep\relax}{\raggedright\ttfamily\scriptsize                 if (!skip\_feature(feature\_mem[feature\_idx])) begin}}
\par\smallskip
This line is part of the combinational decision that detects low-magnitude features. It is the hardware equivalent of asking whether a small activation should be silent for this inference/update pass.
\paragraph{Line 374}
\noindent\colorbox{gray!8}{\parbox{\dimexpr\linewidth-2\fboxsep\relax}{\raggedright\ttfamily\scriptsize                     train\_true\_weight\_r <= weight\_mem[weight\_addr(int'(train\_label\_r), int'(feature\_idx))];}}
\par\smallskip
This line touches the trained classifier weight store. In the final RTL the large 10 x 384 int16 matrix is inferred as true dual-port BRAM, so weights remain on chip and can also be updated by the learning state machine.
\paragraph{Line 375}
\noindent\colorbox{gray!8}{\parbox{\dimexpr\linewidth-2\fboxsep\relax}{\raggedright\ttfamily\scriptsize                 end}}
\par\smallskip
This SystemVerilog line contributes to the structural or sequential description of the accelerator. Its effect is resolved by XSim for verification and by Vivado during synthesis.
\paragraph{Line 376}
\noindent\colorbox{gray!8}{\parbox{\dimexpr\linewidth-2\fboxsep\relax}{\raggedright\ttfamily\scriptsize             end}}
\par\smallskip
This SystemVerilog line contributes to the structural or sequential description of the accelerator. Its effect is resolved by XSim for verification and by Vivado during synthesis.
\paragraph{Line 377}
\noindent\colorbox{gray!8}{\parbox{\dimexpr\linewidth-2\fboxsep\relax}{\raggedright\ttfamily\scriptsize \emph{blank line}}}
\par\smallskip
This blank line is not hardware or software behavior. It separates nearby declarations or operations so the file remains readable and reviewable.
\paragraph{Line 378}
\noindent\colorbox{gray!8}{\parbox{\dimexpr\linewidth-2\fboxsep\relax}{\raggedright\ttfamily\scriptsize             ST\_TRAIN\_WRITE\_TRUE: begin}}
\par\smallskip
This names one state in the control FSM. The numeric encoding is what appears in waveform traces, while the name explains the pipeline phase being executed.
\paragraph{Line 379}
\noindent\colorbox{gray!8}{\parbox{\dimexpr\linewidth-2\fboxsep\relax}{\raggedright\ttfamily\scriptsize                 weight\_mem[weight\_addr(int'(train\_label\_r), int'(feature\_idx))] <= train\_write\_weight\_r;}}
\par\smallskip
This line touches the trained classifier weight store. In the final RTL the large 10 x 384 int16 matrix is inferred as true dual-port BRAM, so weights remain on chip and can also be updated by the learning state machine.
\paragraph{Line 380}
\noindent\colorbox{gray!8}{\parbox{\dimexpr\linewidth-2\fboxsep\relax}{\raggedright\ttfamily\scriptsize             end}}
\par\smallskip
This SystemVerilog line contributes to the structural or sequential description of the accelerator. Its effect is resolved by XSim for verification and by Vivado during synthesis.
\paragraph{Line 381}
\noindent\colorbox{gray!8}{\parbox{\dimexpr\linewidth-2\fboxsep\relax}{\raggedright\ttfamily\scriptsize \emph{blank line}}}
\par\smallskip
This blank line is not hardware or software behavior. It separates nearby declarations or operations so the file remains readable and reviewable.
\paragraph{Line 382}
\noindent\colorbox{gray!8}{\parbox{\dimexpr\linewidth-2\fboxsep\relax}{\raggedright\ttfamily\scriptsize             ST\_TRAIN\_READ\_PRED: begin}}
\par\smallskip
This names one state in the control FSM. The numeric encoding is what appears in waveform traces, while the name explains the pipeline phase being executed.
\paragraph{Line 383}
\noindent\colorbox{gray!8}{\parbox{\dimexpr\linewidth-2\fboxsep\relax}{\raggedright\ttfamily\scriptsize                 train\_pred\_weight\_r <= weight\_mem[weight\_addr(int'(train\_pred\_label\_r), int'(feature\_idx))];}}
\par\smallskip
This line touches the trained classifier weight store. In the final RTL the large 10 x 384 int16 matrix is inferred as true dual-port BRAM, so weights remain on chip and can also be updated by the learning state machine.
\paragraph{Line 384}
\noindent\colorbox{gray!8}{\parbox{\dimexpr\linewidth-2\fboxsep\relax}{\raggedright\ttfamily\scriptsize             end}}
\par\smallskip
This SystemVerilog line contributes to the structural or sequential description of the accelerator. Its effect is resolved by XSim for verification and by Vivado during synthesis.
\paragraph{Line 385}
\noindent\colorbox{gray!8}{\parbox{\dimexpr\linewidth-2\fboxsep\relax}{\raggedright\ttfamily\scriptsize \emph{blank line}}}
\par\smallskip
This blank line is not hardware or software behavior. It separates nearby declarations or operations so the file remains readable and reviewable.
\paragraph{Line 386}
\noindent\colorbox{gray!8}{\parbox{\dimexpr\linewidth-2\fboxsep\relax}{\raggedright\ttfamily\scriptsize             ST\_TRAIN\_WRITE\_PRED: begin}}
\par\smallskip
This names one state in the control FSM. The numeric encoding is what appears in waveform traces, while the name explains the pipeline phase being executed.
\paragraph{Line 387}
\noindent\colorbox{gray!8}{\parbox{\dimexpr\linewidth-2\fboxsep\relax}{\raggedright\ttfamily\scriptsize                 weight\_mem[weight\_addr(int'(train\_pred\_label\_r), int'(feature\_idx))] <= train\_write\_weight\_r;}}
\par\smallskip
This line touches the trained classifier weight store. In the final RTL the large 10 x 384 int16 matrix is inferred as true dual-port BRAM, so weights remain on chip and can also be updated by the learning state machine.
\paragraph{Line 388}
\noindent\colorbox{gray!8}{\parbox{\dimexpr\linewidth-2\fboxsep\relax}{\raggedright\ttfamily\scriptsize             end}}
\par\smallskip
This SystemVerilog line contributes to the structural or sequential description of the accelerator. Its effect is resolved by XSim for verification and by Vivado during synthesis.
\paragraph{Line 389}
\noindent\colorbox{gray!8}{\parbox{\dimexpr\linewidth-2\fboxsep\relax}{\raggedright\ttfamily\scriptsize \emph{blank line}}}
\par\smallskip
This blank line is not hardware or software behavior. It separates nearby declarations or operations so the file remains readable and reviewable.
\paragraph{Line 390}
\noindent\colorbox{gray!8}{\parbox{\dimexpr\linewidth-2\fboxsep\relax}{\raggedright\ttfamily\scriptsize             ST\_TRAIN\_BIAS\_TRUE: begin}}
\par\smallskip
This names one state in the control FSM. The numeric encoding is what appears in waveform traces, while the name explains the pipeline phase being executed.
\paragraph{Line 391}
\noindent\colorbox{gray!8}{\parbox{\dimexpr\linewidth-2\fboxsep\relax}{\raggedright\ttfamily\scriptsize                 bias\_mem[train\_label\_r] <= bias\_mem[train\_label\_r] + BIAS\_LEARN\_DELTA;}}
\par\smallskip
This line touches the class-bias store. Bias values are part of the affine classifier score and are updated on chip in separate one-write states to keep synthesis clean.
\paragraph{Line 392}
\noindent\colorbox{gray!8}{\parbox{\dimexpr\linewidth-2\fboxsep\relax}{\raggedright\ttfamily\scriptsize             end}}
\par\smallskip
This SystemVerilog line contributes to the structural or sequential description of the accelerator. Its effect is resolved by XSim for verification and by Vivado during synthesis.
\paragraph{Line 393}
\noindent\colorbox{gray!8}{\parbox{\dimexpr\linewidth-2\fboxsep\relax}{\raggedright\ttfamily\scriptsize \emph{blank line}}}
\par\smallskip
This blank line is not hardware or software behavior. It separates nearby declarations or operations so the file remains readable and reviewable.
\paragraph{Line 394}
\noindent\colorbox{gray!8}{\parbox{\dimexpr\linewidth-2\fboxsep\relax}{\raggedright\ttfamily\scriptsize             ST\_TRAIN\_BIAS\_PRED: begin}}
\par\smallskip
This names one state in the control FSM. The numeric encoding is what appears in waveform traces, while the name explains the pipeline phase being executed.
\paragraph{Line 395}
\noindent\colorbox{gray!8}{\parbox{\dimexpr\linewidth-2\fboxsep\relax}{\raggedright\ttfamily\scriptsize                 bias\_mem[train\_pred\_label\_r] <= bias\_mem[train\_pred\_label\_r] - BIAS\_LEARN\_DELTA;}}
\par\smallskip
This line touches the class-bias store. Bias values are part of the affine classifier score and are updated on chip in separate one-write states to keep synthesis clean.
\paragraph{Line 396}
\noindent\colorbox{gray!8}{\parbox{\dimexpr\linewidth-2\fboxsep\relax}{\raggedright\ttfamily\scriptsize             end}}
\par\smallskip
This SystemVerilog line contributes to the structural or sequential description of the accelerator. Its effect is resolved by XSim for verification and by Vivado during synthesis.
\paragraph{Line 397}
\noindent\colorbox{gray!8}{\parbox{\dimexpr\linewidth-2\fboxsep\relax}{\raggedright\ttfamily\scriptsize \emph{blank line}}}
\par\smallskip
This blank line is not hardware or software behavior. It separates nearby declarations or operations so the file remains readable and reviewable.
\paragraph{Line 398}
\noindent\colorbox{gray!8}{\parbox{\dimexpr\linewidth-2\fboxsep\relax}{\raggedright\ttfamily\scriptsize             default: begin}}
\par\smallskip
This SystemVerilog line contributes to the structural or sequential description of the accelerator. Its effect is resolved by XSim for verification and by Vivado during synthesis.
\paragraph{Line 399}
\noindent\colorbox{gray!8}{\parbox{\dimexpr\linewidth-2\fboxsep\relax}{\raggedright\ttfamily\scriptsize             end}}
\par\smallskip
This SystemVerilog line contributes to the structural or sequential description of the accelerator. Its effect is resolved by XSim for verification and by Vivado during synthesis.
\paragraph{Line 400}
\noindent\colorbox{gray!8}{\parbox{\dimexpr\linewidth-2\fboxsep\relax}{\raggedright\ttfamily\scriptsize         endcase}}
\par\smallskip
This SystemVerilog line contributes to the structural or sequential description of the accelerator. Its effect is resolved by XSim for verification and by Vivado during synthesis.
\paragraph{Line 401}
\noindent\colorbox{gray!8}{\parbox{\dimexpr\linewidth-2\fboxsep\relax}{\raggedright\ttfamily\scriptsize     end}}
\par\smallskip
This SystemVerilog line contributes to the structural or sequential description of the accelerator. Its effect is resolved by XSim for verification and by Vivado during synthesis.
\paragraph{Line 402}
\noindent\colorbox{gray!8}{\parbox{\dimexpr\linewidth-2\fboxsep\relax}{\raggedright\ttfamily\scriptsize endmodule}}
\par\smallskip
This closes the module definition. Synthesis stops adding logic to this circuit block at this point.
\paragraph{Line 403}
\noindent\colorbox{gray!8}{\parbox{\dimexpr\linewidth-2\fboxsep\relax}{\raggedright\ttfamily\scriptsize \emph{blank line}}}
\par\smallskip
This blank line is not hardware or software behavior. It separates nearby declarations or operations so the file remains readable and reviewable.
\paragraph{Line 404}
\noindent\colorbox{gray!8}{\parbox{\dimexpr\linewidth-2\fboxsep\relax}{\raggedright\ttfamily\scriptsize `default\_nettype wire}}
\par\smallskip
This SystemVerilog line contributes to the structural or sequential description of the accelerator. Its effect is resolved by XSim for verification and by Vivado during synthesis.
\clearpage
\subsection{cifar10\_neuromorphic\textbackslash{}rtl\textbackslash{}cifar10\_spikingformer\_head\_pynq\_axi\_stream\_top.sv}
Each entry below gives the exact source line followed by its hardware/software role in the implemented accelerator pipeline.
\paragraph{Line 1}
\noindent\colorbox{gray!8}{\parbox{\dimexpr\linewidth-2\fboxsep\relax}{\raggedright\ttfamily\scriptsize `timescale 1ns / 1ps}}
\par\smallskip
This simulator directive sets the time unit and precision used by testbenches and waveform timestamps. It keeps delays such as clock half-periods numerically meaningful in XSim.
\paragraph{Line 2}
\noindent\colorbox{gray!8}{\parbox{\dimexpr\linewidth-2\fboxsep\relax}{\raggedright\ttfamily\scriptsize `default\_nettype none}}
\par\smallskip
This disables implicit one-bit wires. It prevents spelling mistakes from silently creating unwanted hardware nets.
\paragraph{Line 3}
\noindent\colorbox{gray!8}{\parbox{\dimexpr\linewidth-2\fboxsep\relax}{\raggedright\ttfamily\scriptsize \emph{blank line}}}
\par\smallskip
This blank line is not hardware or software behavior. It separates nearby declarations or operations so the file remains readable and reviewable.
\paragraph{Line 4}
\noindent\colorbox{gray!8}{\parbox{\dimexpr\linewidth-2\fboxsep\relax}{\raggedright\ttfamily\scriptsize module cifar10\_spikingformer\_head\_pynq\_axi\_stream\_top (}}
\par\smallskip
This begins a hardware module, which is the FPGA equivalent of a reusable circuit block. Everything until the matching endmodule describes ports, storage, combinational equations, and clocked behavior for this block.
\paragraph{Line 5}
\noindent\colorbox{gray!8}{\parbox{\dimexpr\linewidth-2\fboxsep\relax}{\raggedright\ttfamily\scriptsize     input  wire ap\_clk,}}
\par\smallskip
This line declares a module port. A port is a physical or wrapper-visible signal used to move clocks, control pulses, features, labels, predictions, or debug information across module boundaries.
\paragraph{Line 6}
\noindent\colorbox{gray!8}{\parbox{\dimexpr\linewidth-2\fboxsep\relax}{\raggedright\ttfamily\scriptsize     input  wire ap\_rst\_n,}}
\par\smallskip
This line declares a module port. A port is a physical or wrapper-visible signal used to move clocks, control pulses, features, labels, predictions, or debug information across module boundaries.
\paragraph{Line 7}
\noindent\colorbox{gray!8}{\parbox{\dimexpr\linewidth-2\fboxsep\relax}{\raggedright\ttfamily\scriptsize \emph{blank line}}}
\par\smallskip
This blank line is not hardware or software behavior. It separates nearby declarations or operations so the file remains readable and reviewable.
\paragraph{Line 8}
\noindent\colorbox{gray!8}{\parbox{\dimexpr\linewidth-2\fboxsep\relax}{\raggedright\ttfamily\scriptsize     input  wire s\_axis\_feature\_tvalid,}}
\par\smallskip
This line declares a module port. A port is a physical or wrapper-visible signal used to move clocks, control pulses, features, labels, predictions, or debug information across module boundaries.
\paragraph{Line 9}
\noindent\colorbox{gray!8}{\parbox{\dimexpr\linewidth-2\fboxsep\relax}{\raggedright\ttfamily\scriptsize     output logic s\_axis\_feature\_tready,}}
\par\smallskip
This line declares a module port. A port is a physical or wrapper-visible signal used to move clocks, control pulses, features, labels, predictions, or debug information across module boundaries.
\paragraph{Line 10}
\noindent\colorbox{gray!8}{\parbox{\dimexpr\linewidth-2\fboxsep\relax}{\raggedright\ttfamily\scriptsize     input  wire signed [15:0] s\_axis\_feature\_tdata,}}
\par\smallskip
This line declares a module port. A port is a physical or wrapper-visible signal used to move clocks, control pulses, features, labels, predictions, or debug information across module boundaries.
\paragraph{Line 11}
\noindent\colorbox{gray!8}{\parbox{\dimexpr\linewidth-2\fboxsep\relax}{\raggedright\ttfamily\scriptsize     input  wire s\_axis\_feature\_tlast,}}
\par\smallskip
This line declares a module port. A port is a physical or wrapper-visible signal used to move clocks, control pulses, features, labels, predictions, or debug information across module boundaries.
\paragraph{Line 12}
\noindent\colorbox{gray!8}{\parbox{\dimexpr\linewidth-2\fboxsep\relax}{\raggedright\ttfamily\scriptsize \emph{blank line}}}
\par\smallskip
This blank line is not hardware or software behavior. It separates nearby declarations or operations so the file remains readable and reviewable.
\paragraph{Line 13}
\noindent\colorbox{gray!8}{\parbox{\dimexpr\linewidth-2\fboxsep\relax}{\raggedright\ttfamily\scriptsize     input  wire train\_enable,}}
\par\smallskip
This line declares a module port. A port is a physical or wrapper-visible signal used to move clocks, control pulses, features, labels, predictions, or debug information across module boundaries.
\paragraph{Line 14}
\noindent\colorbox{gray!8}{\parbox{\dimexpr\linewidth-2\fboxsep\relax}{\raggedright\ttfamily\scriptsize     input  wire [3:0] train\_label,}}
\par\smallskip
This line declares a module port. A port is a physical or wrapper-visible signal used to move clocks, control pulses, features, labels, predictions, or debug information across module boundaries.
\paragraph{Line 15}
\noindent\colorbox{gray!8}{\parbox{\dimexpr\linewidth-2\fboxsep\relax}{\raggedright\ttfamily\scriptsize \emph{blank line}}}
\par\smallskip
This blank line is not hardware or software behavior. It separates nearby declarations or operations so the file remains readable and reviewable.
\paragraph{Line 16}
\noindent\colorbox{gray!8}{\parbox{\dimexpr\linewidth-2\fboxsep\relax}{\raggedright\ttfamily\scriptsize     output logic m\_axis\_result\_tvalid,}}
\par\smallskip
This line declares a module port. A port is a physical or wrapper-visible signal used to move clocks, control pulses, features, labels, predictions, or debug information across module boundaries.
\paragraph{Line 17}
\noindent\colorbox{gray!8}{\parbox{\dimexpr\linewidth-2\fboxsep\relax}{\raggedright\ttfamily\scriptsize     input  wire m\_axis\_result\_tready,}}
\par\smallskip
This line declares a module port. A port is a physical or wrapper-visible signal used to move clocks, control pulses, features, labels, predictions, or debug information across module boundaries.
\paragraph{Line 18}
\noindent\colorbox{gray!8}{\parbox{\dimexpr\linewidth-2\fboxsep\relax}{\raggedright\ttfamily\scriptsize     output logic [31:0] m\_axis\_result\_tdata,}}
\par\smallskip
This line declares a module port. A port is a physical or wrapper-visible signal used to move clocks, control pulses, features, labels, predictions, or debug information across module boundaries.
\paragraph{Line 19}
\noindent\colorbox{gray!8}{\parbox{\dimexpr\linewidth-2\fboxsep\relax}{\raggedright\ttfamily\scriptsize \emph{blank line}}}
\par\smallskip
This blank line is not hardware or software behavior. It separates nearby declarations or operations so the file remains readable and reviewable.
\paragraph{Line 20}
\noindent\colorbox{gray!8}{\parbox{\dimexpr\linewidth-2\fboxsep\relax}{\raggedright\ttfamily\scriptsize     output logic busy,}}
\par\smallskip
This line declares a module port. A port is a physical or wrapper-visible signal used to move clocks, control pulses, features, labels, predictions, or debug information across module boundaries.
\paragraph{Line 21}
\noindent\colorbox{gray!8}{\parbox{\dimexpr\linewidth-2\fboxsep\relax}{\raggedright\ttfamily\scriptsize     output logic [31:0] debug\_infer\_skip\_count,}}
\par\smallskip
This line exposes or updates a hardware activity counter. The counter makes the feature gate visible in waveforms and CSV traces, so the power-saving opportunity is observable from RTL rather than only from Python statistics.
\paragraph{Line 22}
\noindent\colorbox{gray!8}{\parbox{\dimexpr\linewidth-2\fboxsep\relax}{\raggedright\ttfamily\scriptsize     output logic [31:0] debug\_train\_skip\_count}}
\par\smallskip
This line exposes or updates a hardware activity counter. The counter makes the feature gate visible in waveforms and CSV traces, so the power-saving opportunity is observable from RTL rather than only from Python statistics.
\paragraph{Line 23}
\noindent\colorbox{gray!8}{\parbox{\dimexpr\linewidth-2\fboxsep\relax}{\raggedright\ttfamily\scriptsize );}}
\par\smallskip
This SystemVerilog line contributes to the structural or sequential description of the accelerator. Its effect is resolved by XSim for verification and by Vivado during synthesis.
\paragraph{Line 24}
\noindent\colorbox{gray!8}{\parbox{\dimexpr\linewidth-2\fboxsep\relax}{\raggedright\ttfamily\scriptsize     typedef enum logic [2:0] \{}}
\par\smallskip
This declares the finite-state-machine state type. Using named states makes the cycle-by-cycle control flow explicit instead of relying on anonymous numeric constants.
\paragraph{Line 25}
\noindent\colorbox{gray!8}{\parbox{\dimexpr\linewidth-2\fboxsep\relax}{\raggedright\ttfamily\scriptsize         WR\_IDLE        = 3'd0,}}
\par\smallskip
This names one state in the control FSM. The numeric encoding is what appears in waveform traces, while the name explains the pipeline phase being executed.
\paragraph{Line 26}
\noindent\colorbox{gray!8}{\parbox{\dimexpr\linewidth-2\fboxsep\relax}{\raggedright\ttfamily\scriptsize         WR\_WAIT\_READY  = 3'd1,}}
\par\smallskip
This names one state in the control FSM. The numeric encoding is what appears in waveform traces, while the name explains the pipeline phase being executed.
\paragraph{Line 27}
\noindent\colorbox{gray!8}{\parbox{\dimexpr\linewidth-2\fboxsep\relax}{\raggedright\ttfamily\scriptsize         WR\_SEND\_BUFFER = 3'd2,}}
\par\smallskip
This names one state in the control FSM. The numeric encoding is what appears in waveform traces, while the name explains the pipeline phase being executed.
\paragraph{Line 28}
\noindent\colorbox{gray!8}{\parbox{\dimexpr\linewidth-2\fboxsep\relax}{\raggedright\ttfamily\scriptsize         WR\_STREAM      = 3'd3,}}
\par\smallskip
This names one state in the control FSM. The numeric encoding is what appears in waveform traces, while the name explains the pipeline phase being executed.
\paragraph{Line 29}
\noindent\colorbox{gray!8}{\parbox{\dimexpr\linewidth-2\fboxsep\relax}{\raggedright\ttfamily\scriptsize         WR\_FRAME\_DONE  = 3'd4,}}
\par\smallskip
This names one state in the control FSM. The numeric encoding is what appears in waveform traces, while the name explains the pipeline phase being executed.
\paragraph{Line 30}
\noindent\colorbox{gray!8}{\parbox{\dimexpr\linewidth-2\fboxsep\relax}{\raggedright\ttfamily\scriptsize         WR\_WAIT\_RESULT = 3'd5,}}
\par\smallskip
This names one state in the control FSM. The numeric encoding is what appears in waveform traces, while the name explains the pipeline phase being executed.
\paragraph{Line 31}
\noindent\colorbox{gray!8}{\parbox{\dimexpr\linewidth-2\fboxsep\relax}{\raggedright\ttfamily\scriptsize         WR\_HOLD\_RESULT = 3'd6}}
\par\smallskip
This names one state in the control FSM. The numeric encoding is what appears in waveform traces, while the name explains the pipeline phase being executed.
\paragraph{Line 32}
\noindent\colorbox{gray!8}{\parbox{\dimexpr\linewidth-2\fboxsep\relax}{\raggedright\ttfamily\scriptsize     \} wrap\_state\_t;}}
\par\smallskip
This SystemVerilog line contributes to the structural or sequential description of the accelerator. Its effect is resolved by XSim for verification and by Vivado during synthesis.
\paragraph{Line 33}
\noindent\colorbox{gray!8}{\parbox{\dimexpr\linewidth-2\fboxsep\relax}{\raggedright\ttfamily\scriptsize \emph{blank line}}}
\par\smallskip
This blank line is not hardware or software behavior. It separates nearby declarations or operations so the file remains readable and reviewable.
\paragraph{Line 34}
\noindent\colorbox{gray!8}{\parbox{\dimexpr\linewidth-2\fboxsep\relax}{\raggedright\ttfamily\scriptsize     wrap\_state\_t wrap\_state;}}
\par\smallskip
This SystemVerilog line contributes to the structural or sequential description of the accelerator. Its effect is resolved by XSim for verification and by Vivado during synthesis.
\paragraph{Line 35}
\noindent\colorbox{gray!8}{\parbox{\dimexpr\linewidth-2\fboxsep\relax}{\raggedright\ttfamily\scriptsize \emph{blank line}}}
\par\smallskip
This blank line is not hardware or software behavior. It separates nearby declarations or operations so the file remains readable and reviewable.
\paragraph{Line 36}
\noindent\colorbox{gray!8}{\parbox{\dimexpr\linewidth-2\fboxsep\relax}{\raggedright\ttfamily\scriptsize     logic core\_start;}}
\par\smallskip
This SystemVerilog line contributes to the structural or sequential description of the accelerator. Its effect is resolved by XSim for verification and by Vivado during synthesis.
\paragraph{Line 37}
\noindent\colorbox{gray!8}{\parbox{\dimexpr\linewidth-2\fboxsep\relax}{\raggedright\ttfamily\scriptsize     logic core\_feature\_valid;}}
\par\smallskip
This SystemVerilog line contributes to the structural or sequential description of the accelerator. Its effect is resolved by XSim for verification and by Vivado during synthesis.
\paragraph{Line 38}
\noindent\colorbox{gray!8}{\parbox{\dimexpr\linewidth-2\fboxsep\relax}{\raggedright\ttfamily\scriptsize     logic core\_feature\_ready;}}
\par\smallskip
This SystemVerilog line contributes to the structural or sequential description of the accelerator. Its effect is resolved by XSim for verification and by Vivado during synthesis.
\paragraph{Line 39}
\noindent\colorbox{gray!8}{\parbox{\dimexpr\linewidth-2\fboxsep\relax}{\raggedright\ttfamily\scriptsize     logic signed [15:0] core\_feature\_value;}}
\par\smallskip
This SystemVerilog line contributes to the structural or sequential description of the accelerator. Its effect is resolved by XSim for verification and by Vivado during synthesis.
\paragraph{Line 40}
\noindent\colorbox{gray!8}{\parbox{\dimexpr\linewidth-2\fboxsep\relax}{\raggedright\ttfamily\scriptsize     logic core\_frame\_done;}}
\par\smallskip
This SystemVerilog line contributes to the structural or sequential description of the accelerator. Its effect is resolved by XSim for verification and by Vivado during synthesis.
\paragraph{Line 41}
\noindent\colorbox{gray!8}{\parbox{\dimexpr\linewidth-2\fboxsep\relax}{\raggedright\ttfamily\scriptsize     logic core\_busy;}}
\par\smallskip
This SystemVerilog line contributes to the structural or sequential description of the accelerator. Its effect is resolved by XSim for verification and by Vivado during synthesis.
\paragraph{Line 42}
\noindent\colorbox{gray!8}{\parbox{\dimexpr\linewidth-2\fboxsep\relax}{\raggedright\ttfamily\scriptsize     logic core\_result\_valid;}}
\par\smallskip
This SystemVerilog line contributes to the structural or sequential description of the accelerator. Its effect is resolved by XSim for verification and by Vivado during synthesis.
\paragraph{Line 43}
\noindent\colorbox{gray!8}{\parbox{\dimexpr\linewidth-2\fboxsep\relax}{\raggedright\ttfamily\scriptsize     logic [3:0] core\_predicted\_label;}}
\par\smallskip
This SystemVerilog line contributes to the structural or sequential description of the accelerator. Its effect is resolved by XSim for verification and by Vivado during synthesis.
\paragraph{Line 44}
\noindent\colorbox{gray!8}{\parbox{\dimexpr\linewidth-2\fboxsep\relax}{\raggedright\ttfamily\scriptsize     logic signed [63:0] core\_predicted\_score;}}
\par\smallskip
This SystemVerilog line contributes to the structural or sequential description of the accelerator. Its effect is resolved by XSim for verification and by Vivado during synthesis.
\paragraph{Line 45}
\noindent\colorbox{gray!8}{\parbox{\dimexpr\linewidth-2\fboxsep\relax}{\raggedright\ttfamily\scriptsize     logic core\_train\_update\_valid;}}
\par\smallskip
This SystemVerilog line contributes to the structural or sequential description of the accelerator. Its effect is resolved by XSim for verification and by Vivado during synthesis.
\paragraph{Line 46}
\noindent\colorbox{gray!8}{\parbox{\dimexpr\linewidth-2\fboxsep\relax}{\raggedright\ttfamily\scriptsize     logic core\_train\_update\_applied;}}
\par\smallskip
This SystemVerilog line contributes to the structural or sequential description of the accelerator. Its effect is resolved by XSim for verification and by Vivado during synthesis.
\paragraph{Line 47}
\noindent\colorbox{gray!8}{\parbox{\dimexpr\linewidth-2\fboxsep\relax}{\raggedright\ttfamily\scriptsize     logic [5:0] unused\_debug\_state;}}
\par\smallskip
This SystemVerilog line contributes to the structural or sequential description of the accelerator. Its effect is resolved by XSim for verification and by Vivado during synthesis.
\paragraph{Line 48}
\noindent\colorbox{gray!8}{\parbox{\dimexpr\linewidth-2\fboxsep\relax}{\raggedright\ttfamily\scriptsize     logic [15:0] unused\_debug\_feature;}}
\par\smallskip
This SystemVerilog line contributes to the structural or sequential description of the accelerator. Its effect is resolved by XSim for verification and by Vivado during synthesis.
\paragraph{Line 49}
\noindent\colorbox{gray!8}{\parbox{\dimexpr\linewidth-2\fboxsep\relax}{\raggedright\ttfamily\scriptsize     logic [7:0] unused\_debug\_class;}}
\par\smallskip
This SystemVerilog line contributes to the structural or sequential description of the accelerator. Its effect is resolved by XSim for verification and by Vivado during synthesis.
\paragraph{Line 50}
\noindent\colorbox{gray!8}{\parbox{\dimexpr\linewidth-2\fboxsep\relax}{\raggedright\ttfamily\scriptsize     logic [31:0] core\_update\_count;}}
\par\smallskip
This SystemVerilog line contributes to the structural or sequential description of the accelerator. Its effect is resolved by XSim for verification and by Vivado during synthesis.
\paragraph{Line 51}
\noindent\colorbox{gray!8}{\parbox{\dimexpr\linewidth-2\fboxsep\relax}{\raggedright\ttfamily\scriptsize     logic [31:0] core\_infer\_skip\_count;}}
\par\smallskip
This line exposes or updates a hardware activity counter. The counter makes the feature gate visible in waveforms and CSV traces, so the power-saving opportunity is observable from RTL rather than only from Python statistics.
\paragraph{Line 52}
\noindent\colorbox{gray!8}{\parbox{\dimexpr\linewidth-2\fboxsep\relax}{\raggedright\ttfamily\scriptsize     logic [31:0] core\_train\_skip\_count;}}
\par\smallskip
This line exposes or updates a hardware activity counter. The counter makes the feature gate visible in waveforms and CSV traces, so the power-saving opportunity is observable from RTL rather than only from Python statistics.
\paragraph{Line 53}
\noindent\colorbox{gray!8}{\parbox{\dimexpr\linewidth-2\fboxsep\relax}{\raggedright\ttfamily\scriptsize     logic signed [15:0] buffered\_feature;}}
\par\smallskip
This SystemVerilog line contributes to the structural or sequential description of the accelerator. Its effect is resolved by XSim for verification and by Vivado during synthesis.
\paragraph{Line 54}
\noindent\colorbox{gray!8}{\parbox{\dimexpr\linewidth-2\fboxsep\relax}{\raggedright\ttfamily\scriptsize     logic buffered\_last;}}
\par\smallskip
This SystemVerilog line contributes to the structural or sequential description of the accelerator. Its effect is resolved by XSim for verification and by Vivado during synthesis.
\paragraph{Line 55}
\noindent\colorbox{gray!8}{\parbox{\dimexpr\linewidth-2\fboxsep\relax}{\raggedright\ttfamily\scriptsize     logic update\_applied\_latched;}}
\par\smallskip
This SystemVerilog line contributes to the structural or sequential description of the accelerator. Its effect is resolved by XSim for verification and by Vivado during synthesis.
\paragraph{Line 56}
\noindent\colorbox{gray!8}{\parbox{\dimexpr\linewidth-2\fboxsep\relax}{\raggedright\ttfamily\scriptsize \emph{blank line}}}
\par\smallskip
This blank line is not hardware or software behavior. It separates nearby declarations or operations so the file remains readable and reviewable.
\paragraph{Line 57}
\noindent\colorbox{gray!8}{\parbox{\dimexpr\linewidth-2\fboxsep\relax}{\raggedright\ttfamily\scriptsize     assign busy = core\_busy || (wrap\_state != WR\_IDLE);}}
\par\smallskip
This is continuous combinational wiring. The output follows its source expression immediately in hardware rather than waiting for a clock edge.
\paragraph{Line 58}
\noindent\colorbox{gray!8}{\parbox{\dimexpr\linewidth-2\fboxsep\relax}{\raggedright\ttfamily\scriptsize     assign s\_axis\_feature\_tready = (wrap\_state == WR\_IDLE) || (wrap\_state == WR\_STREAM \&\& core\_feature\_ready);}}
\par\smallskip
This is continuous combinational wiring. The output follows its source expression immediately in hardware rather than waiting for a clock edge.
\paragraph{Line 59}
\noindent\colorbox{gray!8}{\parbox{\dimexpr\linewidth-2\fboxsep\relax}{\raggedright\ttfamily\scriptsize \emph{blank line}}}
\par\smallskip
This blank line is not hardware or software behavior. It separates nearby declarations or operations so the file remains readable and reviewable.
\paragraph{Line 60}
\noindent\colorbox{gray!8}{\parbox{\dimexpr\linewidth-2\fboxsep\relax}{\raggedright\ttfamily\scriptsize     cifar10\_spikingformer\_head\_accelerator u\_head (}}
\par\smallskip
This SystemVerilog line contributes to the structural or sequential description of the accelerator. Its effect is resolved by XSim for verification and by Vivado during synthesis.
\paragraph{Line 61}
\noindent\colorbox{gray!8}{\parbox{\dimexpr\linewidth-2\fboxsep\relax}{\raggedright\ttfamily\scriptsize         .clk(ap\_clk),}}
\par\smallskip
This SystemVerilog line contributes to the structural or sequential description of the accelerator. Its effect is resolved by XSim for verification and by Vivado during synthesis.
\paragraph{Line 62}
\noindent\colorbox{gray!8}{\parbox{\dimexpr\linewidth-2\fboxsep\relax}{\raggedright\ttfamily\scriptsize         .rst\_n(ap\_rst\_n),}}
\par\smallskip
This SystemVerilog line contributes to the structural or sequential description of the accelerator. Its effect is resolved by XSim for verification and by Vivado during synthesis.
\paragraph{Line 63}
\noindent\colorbox{gray!8}{\parbox{\dimexpr\linewidth-2\fboxsep\relax}{\raggedright\ttfamily\scriptsize         .start(core\_start),}}
\par\smallskip
This SystemVerilog line contributes to the structural or sequential description of the accelerator. Its effect is resolved by XSim for verification and by Vivado during synthesis.
\paragraph{Line 64}
\noindent\colorbox{gray!8}{\parbox{\dimexpr\linewidth-2\fboxsep\relax}{\raggedright\ttfamily\scriptsize         .feature\_valid(core\_feature\_valid),}}
\par\smallskip
This SystemVerilog line contributes to the structural or sequential description of the accelerator. Its effect is resolved by XSim for verification and by Vivado during synthesis.
\paragraph{Line 65}
\noindent\colorbox{gray!8}{\parbox{\dimexpr\linewidth-2\fboxsep\relax}{\raggedright\ttfamily\scriptsize         .feature\_ready(core\_feature\_ready),}}
\par\smallskip
This SystemVerilog line contributes to the structural or sequential description of the accelerator. Its effect is resolved by XSim for verification and by Vivado during synthesis.
\paragraph{Line 66}
\noindent\colorbox{gray!8}{\parbox{\dimexpr\linewidth-2\fboxsep\relax}{\raggedright\ttfamily\scriptsize         .feature\_value(core\_feature\_value),}}
\par\smallskip
This SystemVerilog line contributes to the structural or sequential description of the accelerator. Its effect is resolved by XSim for verification and by Vivado during synthesis.
\paragraph{Line 67}
\noindent\colorbox{gray!8}{\parbox{\dimexpr\linewidth-2\fboxsep\relax}{\raggedright\ttfamily\scriptsize         .frame\_done(core\_frame\_done),}}
\par\smallskip
This SystemVerilog line contributes to the structural or sequential description of the accelerator. Its effect is resolved by XSim for verification and by Vivado during synthesis.
\paragraph{Line 68}
\noindent\colorbox{gray!8}{\parbox{\dimexpr\linewidth-2\fboxsep\relax}{\raggedright\ttfamily\scriptsize         .train\_enable(train\_enable),}}
\par\smallskip
This SystemVerilog line contributes to the structural or sequential description of the accelerator. Its effect is resolved by XSim for verification and by Vivado during synthesis.
\paragraph{Line 69}
\noindent\colorbox{gray!8}{\parbox{\dimexpr\linewidth-2\fboxsep\relax}{\raggedright\ttfamily\scriptsize         .train\_label(train\_label),}}
\par\smallskip
This SystemVerilog line contributes to the structural or sequential description of the accelerator. Its effect is resolved by XSim for verification and by Vivado during synthesis.
\paragraph{Line 70}
\noindent\colorbox{gray!8}{\parbox{\dimexpr\linewidth-2\fboxsep\relax}{\raggedright\ttfamily\scriptsize         .busy(core\_busy),}}
\par\smallskip
This SystemVerilog line contributes to the structural or sequential description of the accelerator. Its effect is resolved by XSim for verification and by Vivado during synthesis.
\paragraph{Line 71}
\noindent\colorbox{gray!8}{\parbox{\dimexpr\linewidth-2\fboxsep\relax}{\raggedright\ttfamily\scriptsize         .result\_valid(core\_result\_valid),}}
\par\smallskip
This SystemVerilog line contributes to the structural or sequential description of the accelerator. Its effect is resolved by XSim for verification and by Vivado during synthesis.
\paragraph{Line 72}
\noindent\colorbox{gray!8}{\parbox{\dimexpr\linewidth-2\fboxsep\relax}{\raggedright\ttfamily\scriptsize         .predicted\_label(core\_predicted\_label),}}
\par\smallskip
This SystemVerilog line contributes to the structural or sequential description of the accelerator. Its effect is resolved by XSim for verification and by Vivado during synthesis.
\paragraph{Line 73}
\noindent\colorbox{gray!8}{\parbox{\dimexpr\linewidth-2\fboxsep\relax}{\raggedright\ttfamily\scriptsize         .predicted\_score(core\_predicted\_score),}}
\par\smallskip
This SystemVerilog line contributes to the structural or sequential description of the accelerator. Its effect is resolved by XSim for verification and by Vivado during synthesis.
\paragraph{Line 74}
\noindent\colorbox{gray!8}{\parbox{\dimexpr\linewidth-2\fboxsep\relax}{\raggedright\ttfamily\scriptsize         .train\_update\_valid(core\_train\_update\_valid),}}
\par\smallskip
This SystemVerilog line contributes to the structural or sequential description of the accelerator. Its effect is resolved by XSim for verification and by Vivado during synthesis.
\paragraph{Line 75}
\noindent\colorbox{gray!8}{\parbox{\dimexpr\linewidth-2\fboxsep\relax}{\raggedright\ttfamily\scriptsize         .train\_update\_applied(core\_train\_update\_applied),}}
\par\smallskip
This SystemVerilog line contributes to the structural or sequential description of the accelerator. Its effect is resolved by XSim for verification and by Vivado during synthesis.
\paragraph{Line 76}
\noindent\colorbox{gray!8}{\parbox{\dimexpr\linewidth-2\fboxsep\relax}{\raggedright\ttfamily\scriptsize         .debug\_state(unused\_debug\_state),}}
\par\smallskip
This SystemVerilog line contributes to the structural or sequential description of the accelerator. Its effect is resolved by XSim for verification and by Vivado during synthesis.
\paragraph{Line 77}
\noindent\colorbox{gray!8}{\parbox{\dimexpr\linewidth-2\fboxsep\relax}{\raggedright\ttfamily\scriptsize         .debug\_feature(unused\_debug\_feature),}}
\par\smallskip
This SystemVerilog line contributes to the structural or sequential description of the accelerator. Its effect is resolved by XSim for verification and by Vivado during synthesis.
\paragraph{Line 78}
\noindent\colorbox{gray!8}{\parbox{\dimexpr\linewidth-2\fboxsep\relax}{\raggedright\ttfamily\scriptsize         .debug\_class(unused\_debug\_class),}}
\par\smallskip
This SystemVerilog line contributes to the structural or sequential description of the accelerator. Its effect is resolved by XSim for verification and by Vivado during synthesis.
\paragraph{Line 79}
\noindent\colorbox{gray!8}{\parbox{\dimexpr\linewidth-2\fboxsep\relax}{\raggedright\ttfamily\scriptsize         .debug\_update\_count(core\_update\_count),}}
\par\smallskip
This SystemVerilog line contributes to the structural or sequential description of the accelerator. Its effect is resolved by XSim for verification and by Vivado during synthesis.
\paragraph{Line 80}
\noindent\colorbox{gray!8}{\parbox{\dimexpr\linewidth-2\fboxsep\relax}{\raggedright\ttfamily\scriptsize         .debug\_infer\_skip\_count(core\_infer\_skip\_count),}}
\par\smallskip
This line exposes or updates a hardware activity counter. The counter makes the feature gate visible in waveforms and CSV traces, so the power-saving opportunity is observable from RTL rather than only from Python statistics.
\paragraph{Line 81}
\noindent\colorbox{gray!8}{\parbox{\dimexpr\linewidth-2\fboxsep\relax}{\raggedright\ttfamily\scriptsize         .debug\_train\_skip\_count(core\_train\_skip\_count)}}
\par\smallskip
This line exposes or updates a hardware activity counter. The counter makes the feature gate visible in waveforms and CSV traces, so the power-saving opportunity is observable from RTL rather than only from Python statistics.
\paragraph{Line 82}
\noindent\colorbox{gray!8}{\parbox{\dimexpr\linewidth-2\fboxsep\relax}{\raggedright\ttfamily\scriptsize     );}}
\par\smallskip
This SystemVerilog line contributes to the structural or sequential description of the accelerator. Its effect is resolved by XSim for verification and by Vivado during synthesis.
\paragraph{Line 83}
\noindent\colorbox{gray!8}{\parbox{\dimexpr\linewidth-2\fboxsep\relax}{\raggedright\ttfamily\scriptsize \emph{blank line}}}
\par\smallskip
This blank line is not hardware or software behavior. It separates nearby declarations or operations so the file remains readable and reviewable.
\paragraph{Line 84}
\noindent\colorbox{gray!8}{\parbox{\dimexpr\linewidth-2\fboxsep\relax}{\raggedright\ttfamily\scriptsize     assign debug\_infer\_skip\_count = core\_infer\_skip\_count;}}
\par\smallskip
This line exposes or updates a hardware activity counter. The counter makes the feature gate visible in waveforms and CSV traces, so the power-saving opportunity is observable from RTL rather than only from Python statistics.
\paragraph{Line 85}
\noindent\colorbox{gray!8}{\parbox{\dimexpr\linewidth-2\fboxsep\relax}{\raggedright\ttfamily\scriptsize     assign debug\_train\_skip\_count = core\_train\_skip\_count;}}
\par\smallskip
This line exposes or updates a hardware activity counter. The counter makes the feature gate visible in waveforms and CSV traces, so the power-saving opportunity is observable from RTL rather than only from Python statistics.
\paragraph{Line 86}
\noindent\colorbox{gray!8}{\parbox{\dimexpr\linewidth-2\fboxsep\relax}{\raggedright\ttfamily\scriptsize \emph{blank line}}}
\par\smallskip
This blank line is not hardware or software behavior. It separates nearby declarations or operations so the file remains readable and reviewable.
\paragraph{Line 87}
\noindent\colorbox{gray!8}{\parbox{\dimexpr\linewidth-2\fboxsep\relax}{\raggedright\ttfamily\scriptsize     always\_ff @(posedge ap\_clk or negedge ap\_rst\_n) begin}}
\par\smallskip
This begins clocked sequential logic. Registers and RAM writes inside this block update only on a clock edge, which is what makes the design synthesizable and waveform-predictable.
\paragraph{Line 88}
\noindent\colorbox{gray!8}{\parbox{\dimexpr\linewidth-2\fboxsep\relax}{\raggedright\ttfamily\scriptsize         if (!ap\_rst\_n) begin}}
\par\smallskip
This SystemVerilog line contributes to the structural or sequential description of the accelerator. Its effect is resolved by XSim for verification and by Vivado during synthesis.
\paragraph{Line 89}
\noindent\colorbox{gray!8}{\parbox{\dimexpr\linewidth-2\fboxsep\relax}{\raggedright\ttfamily\scriptsize             wrap\_state <= WR\_IDLE;}}
\par\smallskip
This schedules the next FSM state for the following clock edge. That transition is the hardware pipeline control path.
\paragraph{Line 90}
\noindent\colorbox{gray!8}{\parbox{\dimexpr\linewidth-2\fboxsep\relax}{\raggedright\ttfamily\scriptsize             core\_start <= 1'b0;}}
\par\smallskip
This is a nonblocking register or memory assignment. The right side is sampled now and the left side updates together with other clocked assignments at the end of the clock edge.
\paragraph{Line 91}
\noindent\colorbox{gray!8}{\parbox{\dimexpr\linewidth-2\fboxsep\relax}{\raggedright\ttfamily\scriptsize             core\_feature\_valid <= 1'b0;}}
\par\smallskip
This is a nonblocking register or memory assignment. The right side is sampled now and the left side updates together with other clocked assignments at the end of the clock edge.
\paragraph{Line 92}
\noindent\colorbox{gray!8}{\parbox{\dimexpr\linewidth-2\fboxsep\relax}{\raggedright\ttfamily\scriptsize             core\_feature\_value <= '0;}}
\par\smallskip
This is a nonblocking register or memory assignment. The right side is sampled now and the left side updates together with other clocked assignments at the end of the clock edge.
\paragraph{Line 93}
\noindent\colorbox{gray!8}{\parbox{\dimexpr\linewidth-2\fboxsep\relax}{\raggedright\ttfamily\scriptsize             core\_frame\_done <= 1'b0;}}
\par\smallskip
This is a nonblocking register or memory assignment. The right side is sampled now and the left side updates together with other clocked assignments at the end of the clock edge.
\paragraph{Line 94}
\noindent\colorbox{gray!8}{\parbox{\dimexpr\linewidth-2\fboxsep\relax}{\raggedright\ttfamily\scriptsize             buffered\_feature <= '0;}}
\par\smallskip
This is a nonblocking register or memory assignment. The right side is sampled now and the left side updates together with other clocked assignments at the end of the clock edge.
\paragraph{Line 95}
\noindent\colorbox{gray!8}{\parbox{\dimexpr\linewidth-2\fboxsep\relax}{\raggedright\ttfamily\scriptsize             buffered\_last <= 1'b0;}}
\par\smallskip
This is a nonblocking register or memory assignment. The right side is sampled now and the left side updates together with other clocked assignments at the end of the clock edge.
\paragraph{Line 96}
\noindent\colorbox{gray!8}{\parbox{\dimexpr\linewidth-2\fboxsep\relax}{\raggedright\ttfamily\scriptsize             update\_applied\_latched <= 1'b0;}}
\par\smallskip
This is a nonblocking register or memory assignment. The right side is sampled now and the left side updates together with other clocked assignments at the end of the clock edge.
\paragraph{Line 97}
\noindent\colorbox{gray!8}{\parbox{\dimexpr\linewidth-2\fboxsep\relax}{\raggedright\ttfamily\scriptsize             m\_axis\_result\_tvalid <= 1'b0;}}
\par\smallskip
This is a nonblocking register or memory assignment. The right side is sampled now and the left side updates together with other clocked assignments at the end of the clock edge.
\paragraph{Line 98}
\noindent\colorbox{gray!8}{\parbox{\dimexpr\linewidth-2\fboxsep\relax}{\raggedright\ttfamily\scriptsize             m\_axis\_result\_tdata <= 32'd0;}}
\par\smallskip
This is a nonblocking register or memory assignment. The right side is sampled now and the left side updates together with other clocked assignments at the end of the clock edge.
\paragraph{Line 99}
\noindent\colorbox{gray!8}{\parbox{\dimexpr\linewidth-2\fboxsep\relax}{\raggedright\ttfamily\scriptsize         end else begin}}
\par\smallskip
This SystemVerilog line contributes to the structural or sequential description of the accelerator. Its effect is resolved by XSim for verification and by Vivado during synthesis.
\paragraph{Line 100}
\noindent\colorbox{gray!8}{\parbox{\dimexpr\linewidth-2\fboxsep\relax}{\raggedright\ttfamily\scriptsize             core\_start <= 1'b0;}}
\par\smallskip
This is a nonblocking register or memory assignment. The right side is sampled now and the left side updates together with other clocked assignments at the end of the clock edge.
\paragraph{Line 101}
\noindent\colorbox{gray!8}{\parbox{\dimexpr\linewidth-2\fboxsep\relax}{\raggedright\ttfamily\scriptsize             core\_feature\_valid <= 1'b0;}}
\par\smallskip
This is a nonblocking register or memory assignment. The right side is sampled now and the left side updates together with other clocked assignments at the end of the clock edge.
\paragraph{Line 102}
\noindent\colorbox{gray!8}{\parbox{\dimexpr\linewidth-2\fboxsep\relax}{\raggedright\ttfamily\scriptsize             core\_frame\_done <= 1'b0;}}
\par\smallskip
This is a nonblocking register or memory assignment. The right side is sampled now and the left side updates together with other clocked assignments at the end of the clock edge.
\paragraph{Line 103}
\noindent\colorbox{gray!8}{\parbox{\dimexpr\linewidth-2\fboxsep\relax}{\raggedright\ttfamily\scriptsize \emph{blank line}}}
\par\smallskip
This blank line is not hardware or software behavior. It separates nearby declarations or operations so the file remains readable and reviewable.
\paragraph{Line 104}
\noindent\colorbox{gray!8}{\parbox{\dimexpr\linewidth-2\fboxsep\relax}{\raggedright\ttfamily\scriptsize             if (core\_train\_update\_applied) begin}}
\par\smallskip
This SystemVerilog line contributes to the structural or sequential description of the accelerator. Its effect is resolved by XSim for verification and by Vivado during synthesis.
\paragraph{Line 105}
\noindent\colorbox{gray!8}{\parbox{\dimexpr\linewidth-2\fboxsep\relax}{\raggedright\ttfamily\scriptsize                 update\_applied\_latched <= 1'b1;}}
\par\smallskip
This is a nonblocking register or memory assignment. The right side is sampled now and the left side updates together with other clocked assignments at the end of the clock edge.
\paragraph{Line 106}
\noindent\colorbox{gray!8}{\parbox{\dimexpr\linewidth-2\fboxsep\relax}{\raggedright\ttfamily\scriptsize             end}}
\par\smallskip
This SystemVerilog line contributes to the structural or sequential description of the accelerator. Its effect is resolved by XSim for verification and by Vivado during synthesis.
\paragraph{Line 107}
\noindent\colorbox{gray!8}{\parbox{\dimexpr\linewidth-2\fboxsep\relax}{\raggedright\ttfamily\scriptsize \emph{blank line}}}
\par\smallskip
This blank line is not hardware or software behavior. It separates nearby declarations or operations so the file remains readable and reviewable.
\paragraph{Line 108}
\noindent\colorbox{gray!8}{\parbox{\dimexpr\linewidth-2\fboxsep\relax}{\raggedright\ttfamily\scriptsize             unique case (wrap\_state)}}
\par\smallskip
This starts a state-machine branch table. The unique qualifier tells tools that one legal state should match at a time, improving diagnostics and optimization.
\paragraph{Line 109}
\noindent\colorbox{gray!8}{\parbox{\dimexpr\linewidth-2\fboxsep\relax}{\raggedright\ttfamily\scriptsize                 WR\_IDLE: begin}}
\par\smallskip
This names one state in the control FSM. The numeric encoding is what appears in waveform traces, while the name explains the pipeline phase being executed.
\paragraph{Line 110}
\noindent\colorbox{gray!8}{\parbox{\dimexpr\linewidth-2\fboxsep\relax}{\raggedright\ttfamily\scriptsize                     update\_applied\_latched <= 1'b0;}}
\par\smallskip
This is a nonblocking register or memory assignment. The right side is sampled now and the left side updates together with other clocked assignments at the end of the clock edge.
\paragraph{Line 111}
\noindent\colorbox{gray!8}{\parbox{\dimexpr\linewidth-2\fboxsep\relax}{\raggedright\ttfamily\scriptsize                     if (s\_axis\_feature\_tvalid \&\& s\_axis\_feature\_tready) begin}}
\par\smallskip
This SystemVerilog line contributes to the structural or sequential description of the accelerator. Its effect is resolved by XSim for verification and by Vivado during synthesis.
\paragraph{Line 112}
\noindent\colorbox{gray!8}{\parbox{\dimexpr\linewidth-2\fboxsep\relax}{\raggedright\ttfamily\scriptsize                         buffered\_feature <= s\_axis\_feature\_tdata;}}
\par\smallskip
This is a nonblocking register or memory assignment. The right side is sampled now and the left side updates together with other clocked assignments at the end of the clock edge.
\paragraph{Line 113}
\noindent\colorbox{gray!8}{\parbox{\dimexpr\linewidth-2\fboxsep\relax}{\raggedright\ttfamily\scriptsize                         buffered\_last <= s\_axis\_feature\_tlast;}}
\par\smallskip
This is a nonblocking register or memory assignment. The right side is sampled now and the left side updates together with other clocked assignments at the end of the clock edge.
\paragraph{Line 114}
\noindent\colorbox{gray!8}{\parbox{\dimexpr\linewidth-2\fboxsep\relax}{\raggedright\ttfamily\scriptsize                         core\_start <= 1'b1;}}
\par\smallskip
This is a nonblocking register or memory assignment. The right side is sampled now and the left side updates together with other clocked assignments at the end of the clock edge.
\paragraph{Line 115}
\noindent\colorbox{gray!8}{\parbox{\dimexpr\linewidth-2\fboxsep\relax}{\raggedright\ttfamily\scriptsize                         wrap\_state <= WR\_WAIT\_READY;}}
\par\smallskip
This schedules the next FSM state for the following clock edge. That transition is the hardware pipeline control path.
\paragraph{Line 116}
\noindent\colorbox{gray!8}{\parbox{\dimexpr\linewidth-2\fboxsep\relax}{\raggedright\ttfamily\scriptsize                     end}}
\par\smallskip
This SystemVerilog line contributes to the structural or sequential description of the accelerator. Its effect is resolved by XSim for verification and by Vivado during synthesis.
\paragraph{Line 117}
\noindent\colorbox{gray!8}{\parbox{\dimexpr\linewidth-2\fboxsep\relax}{\raggedright\ttfamily\scriptsize                 end}}
\par\smallskip
This SystemVerilog line contributes to the structural or sequential description of the accelerator. Its effect is resolved by XSim for verification and by Vivado during synthesis.
\paragraph{Line 118}
\noindent\colorbox{gray!8}{\parbox{\dimexpr\linewidth-2\fboxsep\relax}{\raggedright\ttfamily\scriptsize \emph{blank line}}}
\par\smallskip
This blank line is not hardware or software behavior. It separates nearby declarations or operations so the file remains readable and reviewable.
\paragraph{Line 119}
\noindent\colorbox{gray!8}{\parbox{\dimexpr\linewidth-2\fboxsep\relax}{\raggedright\ttfamily\scriptsize                 WR\_WAIT\_READY: begin}}
\par\smallskip
This names one state in the control FSM. The numeric encoding is what appears in waveform traces, while the name explains the pipeline phase being executed.
\paragraph{Line 120}
\noindent\colorbox{gray!8}{\parbox{\dimexpr\linewidth-2\fboxsep\relax}{\raggedright\ttfamily\scriptsize                     if (core\_feature\_ready) begin}}
\par\smallskip
This SystemVerilog line contributes to the structural or sequential description of the accelerator. Its effect is resolved by XSim for verification and by Vivado during synthesis.
\paragraph{Line 121}
\noindent\colorbox{gray!8}{\parbox{\dimexpr\linewidth-2\fboxsep\relax}{\raggedright\ttfamily\scriptsize                         core\_feature\_value <= buffered\_feature;}}
\par\smallskip
This is a nonblocking register or memory assignment. The right side is sampled now and the left side updates together with other clocked assignments at the end of the clock edge.
\paragraph{Line 122}
\noindent\colorbox{gray!8}{\parbox{\dimexpr\linewidth-2\fboxsep\relax}{\raggedright\ttfamily\scriptsize                         core\_feature\_valid <= 1'b1;}}
\par\smallskip
This is a nonblocking register or memory assignment. The right side is sampled now and the left side updates together with other clocked assignments at the end of the clock edge.
\paragraph{Line 123}
\noindent\colorbox{gray!8}{\parbox{\dimexpr\linewidth-2\fboxsep\relax}{\raggedright\ttfamily\scriptsize                         wrap\_state <= WR\_SEND\_BUFFER;}}
\par\smallskip
This schedules the next FSM state for the following clock edge. That transition is the hardware pipeline control path.
\paragraph{Line 124}
\noindent\colorbox{gray!8}{\parbox{\dimexpr\linewidth-2\fboxsep\relax}{\raggedright\ttfamily\scriptsize                     end}}
\par\smallskip
This SystemVerilog line contributes to the structural or sequential description of the accelerator. Its effect is resolved by XSim for verification and by Vivado during synthesis.
\paragraph{Line 125}
\noindent\colorbox{gray!8}{\parbox{\dimexpr\linewidth-2\fboxsep\relax}{\raggedright\ttfamily\scriptsize                 end}}
\par\smallskip
This SystemVerilog line contributes to the structural or sequential description of the accelerator. Its effect is resolved by XSim for verification and by Vivado during synthesis.
\paragraph{Line 126}
\noindent\colorbox{gray!8}{\parbox{\dimexpr\linewidth-2\fboxsep\relax}{\raggedright\ttfamily\scriptsize \emph{blank line}}}
\par\smallskip
This blank line is not hardware or software behavior. It separates nearby declarations or operations so the file remains readable and reviewable.
\paragraph{Line 127}
\noindent\colorbox{gray!8}{\parbox{\dimexpr\linewidth-2\fboxsep\relax}{\raggedright\ttfamily\scriptsize                 WR\_SEND\_BUFFER: begin}}
\par\smallskip
This names one state in the control FSM. The numeric encoding is what appears in waveform traces, while the name explains the pipeline phase being executed.
\paragraph{Line 128}
\noindent\colorbox{gray!8}{\parbox{\dimexpr\linewidth-2\fboxsep\relax}{\raggedright\ttfamily\scriptsize                     if (buffered\_last) begin}}
\par\smallskip
This SystemVerilog line contributes to the structural or sequential description of the accelerator. Its effect is resolved by XSim for verification and by Vivado during synthesis.
\paragraph{Line 129}
\noindent\colorbox{gray!8}{\parbox{\dimexpr\linewidth-2\fboxsep\relax}{\raggedright\ttfamily\scriptsize                         wrap\_state <= WR\_FRAME\_DONE;}}
\par\smallskip
This schedules the next FSM state for the following clock edge. That transition is the hardware pipeline control path.
\paragraph{Line 130}
\noindent\colorbox{gray!8}{\parbox{\dimexpr\linewidth-2\fboxsep\relax}{\raggedright\ttfamily\scriptsize                     end else begin}}
\par\smallskip
This SystemVerilog line contributes to the structural or sequential description of the accelerator. Its effect is resolved by XSim for verification and by Vivado during synthesis.
\paragraph{Line 131}
\noindent\colorbox{gray!8}{\parbox{\dimexpr\linewidth-2\fboxsep\relax}{\raggedright\ttfamily\scriptsize                         wrap\_state <= WR\_STREAM;}}
\par\smallskip
This schedules the next FSM state for the following clock edge. That transition is the hardware pipeline control path.
\paragraph{Line 132}
\noindent\colorbox{gray!8}{\parbox{\dimexpr\linewidth-2\fboxsep\relax}{\raggedright\ttfamily\scriptsize                     end}}
\par\smallskip
This SystemVerilog line contributes to the structural or sequential description of the accelerator. Its effect is resolved by XSim for verification and by Vivado during synthesis.
\paragraph{Line 133}
\noindent\colorbox{gray!8}{\parbox{\dimexpr\linewidth-2\fboxsep\relax}{\raggedright\ttfamily\scriptsize                 end}}
\par\smallskip
This SystemVerilog line contributes to the structural or sequential description of the accelerator. Its effect is resolved by XSim for verification and by Vivado during synthesis.
\paragraph{Line 134}
\noindent\colorbox{gray!8}{\parbox{\dimexpr\linewidth-2\fboxsep\relax}{\raggedright\ttfamily\scriptsize \emph{blank line}}}
\par\smallskip
This blank line is not hardware or software behavior. It separates nearby declarations or operations so the file remains readable and reviewable.
\paragraph{Line 135}
\noindent\colorbox{gray!8}{\parbox{\dimexpr\linewidth-2\fboxsep\relax}{\raggedright\ttfamily\scriptsize                 WR\_STREAM: begin}}
\par\smallskip
This names one state in the control FSM. The numeric encoding is what appears in waveform traces, while the name explains the pipeline phase being executed.
\paragraph{Line 136}
\noindent\colorbox{gray!8}{\parbox{\dimexpr\linewidth-2\fboxsep\relax}{\raggedright\ttfamily\scriptsize                     if (s\_axis\_feature\_tvalid \&\& s\_axis\_feature\_tready) begin}}
\par\smallskip
This SystemVerilog line contributes to the structural or sequential description of the accelerator. Its effect is resolved by XSim for verification and by Vivado during synthesis.
\paragraph{Line 137}
\noindent\colorbox{gray!8}{\parbox{\dimexpr\linewidth-2\fboxsep\relax}{\raggedright\ttfamily\scriptsize                         core\_feature\_value <= s\_axis\_feature\_tdata;}}
\par\smallskip
This is a nonblocking register or memory assignment. The right side is sampled now and the left side updates together with other clocked assignments at the end of the clock edge.
\paragraph{Line 138}
\noindent\colorbox{gray!8}{\parbox{\dimexpr\linewidth-2\fboxsep\relax}{\raggedright\ttfamily\scriptsize                         core\_feature\_valid <= 1'b1;}}
\par\smallskip
This is a nonblocking register or memory assignment. The right side is sampled now and the left side updates together with other clocked assignments at the end of the clock edge.
\paragraph{Line 139}
\noindent\colorbox{gray!8}{\parbox{\dimexpr\linewidth-2\fboxsep\relax}{\raggedright\ttfamily\scriptsize                         if (s\_axis\_feature\_tlast) begin}}
\par\smallskip
This SystemVerilog line contributes to the structural or sequential description of the accelerator. Its effect is resolved by XSim for verification and by Vivado during synthesis.
\paragraph{Line 140}
\noindent\colorbox{gray!8}{\parbox{\dimexpr\linewidth-2\fboxsep\relax}{\raggedright\ttfamily\scriptsize                             wrap\_state <= WR\_FRAME\_DONE;}}
\par\smallskip
This schedules the next FSM state for the following clock edge. That transition is the hardware pipeline control path.
\paragraph{Line 141}
\noindent\colorbox{gray!8}{\parbox{\dimexpr\linewidth-2\fboxsep\relax}{\raggedright\ttfamily\scriptsize                         end}}
\par\smallskip
This SystemVerilog line contributes to the structural or sequential description of the accelerator. Its effect is resolved by XSim for verification and by Vivado during synthesis.
\paragraph{Line 142}
\noindent\colorbox{gray!8}{\parbox{\dimexpr\linewidth-2\fboxsep\relax}{\raggedright\ttfamily\scriptsize                     end}}
\par\smallskip
This SystemVerilog line contributes to the structural or sequential description of the accelerator. Its effect is resolved by XSim for verification and by Vivado during synthesis.
\paragraph{Line 143}
\noindent\colorbox{gray!8}{\parbox{\dimexpr\linewidth-2\fboxsep\relax}{\raggedright\ttfamily\scriptsize                 end}}
\par\smallskip
This SystemVerilog line contributes to the structural or sequential description of the accelerator. Its effect is resolved by XSim for verification and by Vivado during synthesis.
\paragraph{Line 144}
\noindent\colorbox{gray!8}{\parbox{\dimexpr\linewidth-2\fboxsep\relax}{\raggedright\ttfamily\scriptsize \emph{blank line}}}
\par\smallskip
This blank line is not hardware or software behavior. It separates nearby declarations or operations so the file remains readable and reviewable.
\paragraph{Line 145}
\noindent\colorbox{gray!8}{\parbox{\dimexpr\linewidth-2\fboxsep\relax}{\raggedright\ttfamily\scriptsize                 WR\_FRAME\_DONE: begin}}
\par\smallskip
This names one state in the control FSM. The numeric encoding is what appears in waveform traces, while the name explains the pipeline phase being executed.
\paragraph{Line 146}
\noindent\colorbox{gray!8}{\parbox{\dimexpr\linewidth-2\fboxsep\relax}{\raggedright\ttfamily\scriptsize                     core\_frame\_done <= 1'b1;}}
\par\smallskip
This is a nonblocking register or memory assignment. The right side is sampled now and the left side updates together with other clocked assignments at the end of the clock edge.
\paragraph{Line 147}
\noindent\colorbox{gray!8}{\parbox{\dimexpr\linewidth-2\fboxsep\relax}{\raggedright\ttfamily\scriptsize                     wrap\_state <= WR\_WAIT\_RESULT;}}
\par\smallskip
This schedules the next FSM state for the following clock edge. That transition is the hardware pipeline control path.
\paragraph{Line 148}
\noindent\colorbox{gray!8}{\parbox{\dimexpr\linewidth-2\fboxsep\relax}{\raggedright\ttfamily\scriptsize                 end}}
\par\smallskip
This SystemVerilog line contributes to the structural or sequential description of the accelerator. Its effect is resolved by XSim for verification and by Vivado during synthesis.
\paragraph{Line 149}
\noindent\colorbox{gray!8}{\parbox{\dimexpr\linewidth-2\fboxsep\relax}{\raggedright\ttfamily\scriptsize \emph{blank line}}}
\par\smallskip
This blank line is not hardware or software behavior. It separates nearby declarations or operations so the file remains readable and reviewable.
\paragraph{Line 150}
\noindent\colorbox{gray!8}{\parbox{\dimexpr\linewidth-2\fboxsep\relax}{\raggedright\ttfamily\scriptsize                 WR\_WAIT\_RESULT: begin}}
\par\smallskip
This names one state in the control FSM. The numeric encoding is what appears in waveform traces, while the name explains the pipeline phase being executed.
\paragraph{Line 151}
\noindent\colorbox{gray!8}{\parbox{\dimexpr\linewidth-2\fboxsep\relax}{\raggedright\ttfamily\scriptsize                     if (core\_result\_valid) begin}}
\par\smallskip
This SystemVerilog line contributes to the structural or sequential description of the accelerator. Its effect is resolved by XSim for verification and by Vivado during synthesis.
\paragraph{Line 152}
\noindent\colorbox{gray!8}{\parbox{\dimexpr\linewidth-2\fboxsep\relax}{\raggedright\ttfamily\scriptsize                         m\_axis\_result\_tdata <= \{}}
\par\smallskip
This is a nonblocking register or memory assignment. The right side is sampled now and the left side updates together with other clocked assignments at the end of the clock edge.
\paragraph{Line 153}
\noindent\colorbox{gray!8}{\parbox{\dimexpr\linewidth-2\fboxsep\relax}{\raggedright\ttfamily\scriptsize                             core\_predicted\_score[15:0],}}
\par\smallskip
This SystemVerilog line contributes to the structural or sequential description of the accelerator. Its effect is resolved by XSim for verification and by Vivado during synthesis.
\paragraph{Line 154}
\noindent\colorbox{gray!8}{\parbox{\dimexpr\linewidth-2\fboxsep\relax}{\raggedright\ttfamily\scriptsize                             core\_update\_count[10:0],}}
\par\smallskip
This SystemVerilog line contributes to the structural or sequential description of the accelerator. Its effect is resolved by XSim for verification and by Vivado during synthesis.
\paragraph{Line 155}
\noindent\colorbox{gray!8}{\parbox{\dimexpr\linewidth-2\fboxsep\relax}{\raggedright\ttfamily\scriptsize                             update\_applied\_latched,}}
\par\smallskip
This SystemVerilog line contributes to the structural or sequential description of the accelerator. Its effect is resolved by XSim for verification and by Vivado during synthesis.
\paragraph{Line 156}
\noindent\colorbox{gray!8}{\parbox{\dimexpr\linewidth-2\fboxsep\relax}{\raggedright\ttfamily\scriptsize                             core\_predicted\_label}}
\par\smallskip
This SystemVerilog line contributes to the structural or sequential description of the accelerator. Its effect is resolved by XSim for verification and by Vivado during synthesis.
\paragraph{Line 157}
\noindent\colorbox{gray!8}{\parbox{\dimexpr\linewidth-2\fboxsep\relax}{\raggedright\ttfamily\scriptsize                         \};}}
\par\smallskip
This SystemVerilog line contributes to the structural or sequential description of the accelerator. Its effect is resolved by XSim for verification and by Vivado during synthesis.
\paragraph{Line 158}
\noindent\colorbox{gray!8}{\parbox{\dimexpr\linewidth-2\fboxsep\relax}{\raggedright\ttfamily\scriptsize                         m\_axis\_result\_tvalid <= 1'b1;}}
\par\smallskip
This is a nonblocking register or memory assignment. The right side is sampled now and the left side updates together with other clocked assignments at the end of the clock edge.
\paragraph{Line 159}
\noindent\colorbox{gray!8}{\parbox{\dimexpr\linewidth-2\fboxsep\relax}{\raggedright\ttfamily\scriptsize                         wrap\_state <= WR\_HOLD\_RESULT;}}
\par\smallskip
This schedules the next FSM state for the following clock edge. That transition is the hardware pipeline control path.
\paragraph{Line 160}
\noindent\colorbox{gray!8}{\parbox{\dimexpr\linewidth-2\fboxsep\relax}{\raggedright\ttfamily\scriptsize                     end}}
\par\smallskip
This SystemVerilog line contributes to the structural or sequential description of the accelerator. Its effect is resolved by XSim for verification and by Vivado during synthesis.
\paragraph{Line 161}
\noindent\colorbox{gray!8}{\parbox{\dimexpr\linewidth-2\fboxsep\relax}{\raggedright\ttfamily\scriptsize                 end}}
\par\smallskip
This SystemVerilog line contributes to the structural or sequential description of the accelerator. Its effect is resolved by XSim for verification and by Vivado during synthesis.
\paragraph{Line 162}
\noindent\colorbox{gray!8}{\parbox{\dimexpr\linewidth-2\fboxsep\relax}{\raggedright\ttfamily\scriptsize \emph{blank line}}}
\par\smallskip
This blank line is not hardware or software behavior. It separates nearby declarations or operations so the file remains readable and reviewable.
\paragraph{Line 163}
\noindent\colorbox{gray!8}{\parbox{\dimexpr\linewidth-2\fboxsep\relax}{\raggedright\ttfamily\scriptsize                 WR\_HOLD\_RESULT: begin}}
\par\smallskip
This names one state in the control FSM. The numeric encoding is what appears in waveform traces, while the name explains the pipeline phase being executed.
\paragraph{Line 164}
\noindent\colorbox{gray!8}{\parbox{\dimexpr\linewidth-2\fboxsep\relax}{\raggedright\ttfamily\scriptsize                     if (m\_axis\_result\_tvalid \&\& m\_axis\_result\_tready) begin}}
\par\smallskip
This SystemVerilog line contributes to the structural or sequential description of the accelerator. Its effect is resolved by XSim for verification and by Vivado during synthesis.
\paragraph{Line 165}
\noindent\colorbox{gray!8}{\parbox{\dimexpr\linewidth-2\fboxsep\relax}{\raggedright\ttfamily\scriptsize                         m\_axis\_result\_tvalid <= 1'b0;}}
\par\smallskip
This is a nonblocking register or memory assignment. The right side is sampled now and the left side updates together with other clocked assignments at the end of the clock edge.
\paragraph{Line 166}
\noindent\colorbox{gray!8}{\parbox{\dimexpr\linewidth-2\fboxsep\relax}{\raggedright\ttfamily\scriptsize                         wrap\_state <= WR\_IDLE;}}
\par\smallskip
This schedules the next FSM state for the following clock edge. That transition is the hardware pipeline control path.
\paragraph{Line 167}
\noindent\colorbox{gray!8}{\parbox{\dimexpr\linewidth-2\fboxsep\relax}{\raggedright\ttfamily\scriptsize                     end}}
\par\smallskip
This SystemVerilog line contributes to the structural or sequential description of the accelerator. Its effect is resolved by XSim for verification and by Vivado during synthesis.
\paragraph{Line 168}
\noindent\colorbox{gray!8}{\parbox{\dimexpr\linewidth-2\fboxsep\relax}{\raggedright\ttfamily\scriptsize                 end}}
\par\smallskip
This SystemVerilog line contributes to the structural or sequential description of the accelerator. Its effect is resolved by XSim for verification and by Vivado during synthesis.
\paragraph{Line 169}
\noindent\colorbox{gray!8}{\parbox{\dimexpr\linewidth-2\fboxsep\relax}{\raggedright\ttfamily\scriptsize \emph{blank line}}}
\par\smallskip
This blank line is not hardware or software behavior. It separates nearby declarations or operations so the file remains readable and reviewable.
\paragraph{Line 170}
\noindent\colorbox{gray!8}{\parbox{\dimexpr\linewidth-2\fboxsep\relax}{\raggedright\ttfamily\scriptsize                 default: begin}}
\par\smallskip
This SystemVerilog line contributes to the structural or sequential description of the accelerator. Its effect is resolved by XSim for verification and by Vivado during synthesis.
\paragraph{Line 171}
\noindent\colorbox{gray!8}{\parbox{\dimexpr\linewidth-2\fboxsep\relax}{\raggedright\ttfamily\scriptsize                     wrap\_state <= WR\_IDLE;}}
\par\smallskip
This schedules the next FSM state for the following clock edge. That transition is the hardware pipeline control path.
\paragraph{Line 172}
\noindent\colorbox{gray!8}{\parbox{\dimexpr\linewidth-2\fboxsep\relax}{\raggedright\ttfamily\scriptsize                 end}}
\par\smallskip
This SystemVerilog line contributes to the structural or sequential description of the accelerator. Its effect is resolved by XSim for verification and by Vivado during synthesis.
\paragraph{Line 173}
\noindent\colorbox{gray!8}{\parbox{\dimexpr\linewidth-2\fboxsep\relax}{\raggedright\ttfamily\scriptsize             endcase}}
\par\smallskip
This SystemVerilog line contributes to the structural or sequential description of the accelerator. Its effect is resolved by XSim for verification and by Vivado during synthesis.
\paragraph{Line 174}
\noindent\colorbox{gray!8}{\parbox{\dimexpr\linewidth-2\fboxsep\relax}{\raggedright\ttfamily\scriptsize         end}}
\par\smallskip
This SystemVerilog line contributes to the structural or sequential description of the accelerator. Its effect is resolved by XSim for verification and by Vivado during synthesis.
\paragraph{Line 175}
\noindent\colorbox{gray!8}{\parbox{\dimexpr\linewidth-2\fboxsep\relax}{\raggedright\ttfamily\scriptsize     end}}
\par\smallskip
This SystemVerilog line contributes to the structural or sequential description of the accelerator. Its effect is resolved by XSim for verification and by Vivado during synthesis.
\paragraph{Line 176}
\noindent\colorbox{gray!8}{\parbox{\dimexpr\linewidth-2\fboxsep\relax}{\raggedright\ttfamily\scriptsize endmodule}}
\par\smallskip
This closes the module definition. Synthesis stops adding logic to this circuit block at this point.
\paragraph{Line 177}
\noindent\colorbox{gray!8}{\parbox{\dimexpr\linewidth-2\fboxsep\relax}{\raggedright\ttfamily\scriptsize \emph{blank line}}}
\par\smallskip
This blank line is not hardware or software behavior. It separates nearby declarations or operations so the file remains readable and reviewable.
\paragraph{Line 178}
\noindent\colorbox{gray!8}{\parbox{\dimexpr\linewidth-2\fboxsep\relax}{\raggedright\ttfamily\scriptsize `default\_nettype wire}}
\par\smallskip
This SystemVerilog line contributes to the structural or sequential description of the accelerator. Its effect is resolved by XSim for verification and by Vivado during synthesis.
\clearpage
\subsection{cifar10\_neuromorphic\textbackslash{}rtl\textbackslash{}tb\_spikingformer\_head\_learning\_unit.sv}
Each entry below gives the exact source line followed by its hardware/software role in the implemented accelerator pipeline.
\paragraph{Line 1}
\noindent\colorbox{gray!8}{\parbox{\dimexpr\linewidth-2\fboxsep\relax}{\raggedright\ttfamily\scriptsize `timescale 1ns / 1ps}}
\par\smallskip
This simulator directive sets the time unit and precision used by testbenches and waveform timestamps. It keeps delays such as clock half-periods numerically meaningful in XSim.
\paragraph{Line 2}
\noindent\colorbox{gray!8}{\parbox{\dimexpr\linewidth-2\fboxsep\relax}{\raggedright\ttfamily\scriptsize \emph{blank line}}}
\par\smallskip
This blank line is not hardware or software behavior. It separates nearby declarations or operations so the file remains readable and reviewable.
\paragraph{Line 3}
\noindent\colorbox{gray!8}{\parbox{\dimexpr\linewidth-2\fboxsep\relax}{\raggedright\ttfamily\scriptsize module tb\_spikingformer\_head\_learning\_unit;}}
\par\smallskip
This begins a hardware module, which is the FPGA equivalent of a reusable circuit block. Everything until the matching endmodule describes ports, storage, combinational equations, and clocked behavior for this block.
\paragraph{Line 4}
\noindent\colorbox{gray!8}{\parbox{\dimexpr\linewidth-2\fboxsep\relax}{\raggedright\ttfamily\scriptsize     localparam int FEATURE\_DIM = 4;}}
\par\smallskip
This is a local compile-time or procedural assignment. In context it either defines a helper value or connects a temporary expression used by nearby hardware logic.
\paragraph{Line 5}
\noindent\colorbox{gray!8}{\parbox{\dimexpr\linewidth-2\fboxsep\relax}{\raggedright\ttfamily\scriptsize     localparam int NUM\_CLASSES = 3;}}
\par\smallskip
This is a local compile-time or procedural assignment. In context it either defines a helper value or connects a temporary expression used by nearby hardware logic.
\paragraph{Line 6}
\noindent\colorbox{gray!8}{\parbox{\dimexpr\linewidth-2\fboxsep\relax}{\raggedright\ttfamily\scriptsize \emph{blank line}}}
\par\smallskip
This blank line is not hardware or software behavior. It separates nearby declarations or operations so the file remains readable and reviewable.
\paragraph{Line 7}
\noindent\colorbox{gray!8}{\parbox{\dimexpr\linewidth-2\fboxsep\relax}{\raggedright\ttfamily\scriptsize     logic clk = 1'b0;}}
\par\smallskip
This is a local compile-time or procedural assignment. In context it either defines a helper value or connects a temporary expression used by nearby hardware logic.
\paragraph{Line 8}
\noindent\colorbox{gray!8}{\parbox{\dimexpr\linewidth-2\fboxsep\relax}{\raggedright\ttfamily\scriptsize     logic rst\_n = 1'b0;}}
\par\smallskip
This is a local compile-time or procedural assignment. In context it either defines a helper value or connects a temporary expression used by nearby hardware logic.
\paragraph{Line 9}
\noindent\colorbox{gray!8}{\parbox{\dimexpr\linewidth-2\fboxsep\relax}{\raggedright\ttfamily\scriptsize     logic start;}}
\par\smallskip
This SystemVerilog line contributes to the structural or sequential description of the accelerator. Its effect is resolved by XSim for verification and by Vivado during synthesis.
\paragraph{Line 10}
\noindent\colorbox{gray!8}{\parbox{\dimexpr\linewidth-2\fboxsep\relax}{\raggedright\ttfamily\scriptsize     logic feature\_valid;}}
\par\smallskip
This SystemVerilog line contributes to the structural or sequential description of the accelerator. Its effect is resolved by XSim for verification and by Vivado during synthesis.
\paragraph{Line 11}
\noindent\colorbox{gray!8}{\parbox{\dimexpr\linewidth-2\fboxsep\relax}{\raggedright\ttfamily\scriptsize     logic feature\_ready;}}
\par\smallskip
This SystemVerilog line contributes to the structural or sequential description of the accelerator. Its effect is resolved by XSim for verification and by Vivado during synthesis.
\paragraph{Line 12}
\noindent\colorbox{gray!8}{\parbox{\dimexpr\linewidth-2\fboxsep\relax}{\raggedright\ttfamily\scriptsize     logic signed [15:0] feature\_value;}}
\par\smallskip
This SystemVerilog line contributes to the structural or sequential description of the accelerator. Its effect is resolved by XSim for verification and by Vivado during synthesis.
\paragraph{Line 13}
\noindent\colorbox{gray!8}{\parbox{\dimexpr\linewidth-2\fboxsep\relax}{\raggedright\ttfamily\scriptsize     logic frame\_done;}}
\par\smallskip
This SystemVerilog line contributes to the structural or sequential description of the accelerator. Its effect is resolved by XSim for verification and by Vivado during synthesis.
\paragraph{Line 14}
\noindent\colorbox{gray!8}{\parbox{\dimexpr\linewidth-2\fboxsep\relax}{\raggedright\ttfamily\scriptsize     logic train\_enable;}}
\par\smallskip
This SystemVerilog line contributes to the structural or sequential description of the accelerator. Its effect is resolved by XSim for verification and by Vivado during synthesis.
\paragraph{Line 15}
\noindent\colorbox{gray!8}{\parbox{\dimexpr\linewidth-2\fboxsep\relax}{\raggedright\ttfamily\scriptsize     logic [3:0] train\_label;}}
\par\smallskip
This SystemVerilog line contributes to the structural or sequential description of the accelerator. Its effect is resolved by XSim for verification and by Vivado during synthesis.
\paragraph{Line 16}
\noindent\colorbox{gray!8}{\parbox{\dimexpr\linewidth-2\fboxsep\relax}{\raggedright\ttfamily\scriptsize     logic busy;}}
\par\smallskip
This SystemVerilog line contributes to the structural or sequential description of the accelerator. Its effect is resolved by XSim for verification and by Vivado during synthesis.
\paragraph{Line 17}
\noindent\colorbox{gray!8}{\parbox{\dimexpr\linewidth-2\fboxsep\relax}{\raggedright\ttfamily\scriptsize     logic result\_valid;}}
\par\smallskip
This SystemVerilog line contributes to the structural or sequential description of the accelerator. Its effect is resolved by XSim for verification and by Vivado during synthesis.
\paragraph{Line 18}
\noindent\colorbox{gray!8}{\parbox{\dimexpr\linewidth-2\fboxsep\relax}{\raggedright\ttfamily\scriptsize     logic [3:0] predicted\_label;}}
\par\smallskip
This SystemVerilog line contributes to the structural or sequential description of the accelerator. Its effect is resolved by XSim for verification and by Vivado during synthesis.
\paragraph{Line 19}
\noindent\colorbox{gray!8}{\parbox{\dimexpr\linewidth-2\fboxsep\relax}{\raggedright\ttfamily\scriptsize     logic signed [63:0] predicted\_score;}}
\par\smallskip
This SystemVerilog line contributes to the structural or sequential description of the accelerator. Its effect is resolved by XSim for verification and by Vivado during synthesis.
\paragraph{Line 20}
\noindent\colorbox{gray!8}{\parbox{\dimexpr\linewidth-2\fboxsep\relax}{\raggedright\ttfamily\scriptsize     logic train\_update\_valid;}}
\par\smallskip
This SystemVerilog line contributes to the structural or sequential description of the accelerator. Its effect is resolved by XSim for verification and by Vivado during synthesis.
\paragraph{Line 21}
\noindent\colorbox{gray!8}{\parbox{\dimexpr\linewidth-2\fboxsep\relax}{\raggedright\ttfamily\scriptsize     logic train\_update\_applied;}}
\par\smallskip
This SystemVerilog line contributes to the structural or sequential description of the accelerator. Its effect is resolved by XSim for verification and by Vivado during synthesis.
\paragraph{Line 22}
\noindent\colorbox{gray!8}{\parbox{\dimexpr\linewidth-2\fboxsep\relax}{\raggedright\ttfamily\scriptsize     logic [5:0] debug\_state;}}
\par\smallskip
This SystemVerilog line contributes to the structural or sequential description of the accelerator. Its effect is resolved by XSim for verification and by Vivado during synthesis.
\paragraph{Line 23}
\noindent\colorbox{gray!8}{\parbox{\dimexpr\linewidth-2\fboxsep\relax}{\raggedright\ttfamily\scriptsize     logic [15:0] debug\_feature;}}
\par\smallskip
This SystemVerilog line contributes to the structural or sequential description of the accelerator. Its effect is resolved by XSim for verification and by Vivado during synthesis.
\paragraph{Line 24}
\noindent\colorbox{gray!8}{\parbox{\dimexpr\linewidth-2\fboxsep\relax}{\raggedright\ttfamily\scriptsize     logic [7:0] debug\_class;}}
\par\smallskip
This SystemVerilog line contributes to the structural or sequential description of the accelerator. Its effect is resolved by XSim for verification and by Vivado during synthesis.
\paragraph{Line 25}
\noindent\colorbox{gray!8}{\parbox{\dimexpr\linewidth-2\fboxsep\relax}{\raggedright\ttfamily\scriptsize     logic [31:0] debug\_update\_count;}}
\par\smallskip
This SystemVerilog line contributes to the structural or sequential description of the accelerator. Its effect is resolved by XSim for verification and by Vivado during synthesis.
\paragraph{Line 26}
\noindent\colorbox{gray!8}{\parbox{\dimexpr\linewidth-2\fboxsep\relax}{\raggedright\ttfamily\scriptsize     logic [31:0] debug\_infer\_skip\_count;}}
\par\smallskip
This line exposes or updates a hardware activity counter. The counter makes the feature gate visible in waveforms and CSV traces, so the power-saving opportunity is observable from RTL rather than only from Python statistics.
\paragraph{Line 27}
\noindent\colorbox{gray!8}{\parbox{\dimexpr\linewidth-2\fboxsep\relax}{\raggedright\ttfamily\scriptsize     logic [31:0] debug\_train\_skip\_count;}}
\par\smallskip
This line exposes or updates a hardware activity counter. The counter makes the feature gate visible in waveforms and CSV traces, so the power-saving opportunity is observable from RTL rather than only from Python statistics.
\paragraph{Line 28}
\noindent\colorbox{gray!8}{\parbox{\dimexpr\linewidth-2\fboxsep\relax}{\raggedright\ttfamily\scriptsize \emph{blank line}}}
\par\smallskip
This blank line is not hardware or software behavior. It separates nearby declarations or operations so the file remains readable and reviewable.
\paragraph{Line 29}
\noindent\colorbox{gray!8}{\parbox{\dimexpr\linewidth-2\fboxsep\relax}{\raggedright\ttfamily\scriptsize     logic signed [15:0] x [0:FEATURE\_DIM-1];}}
\par\smallskip
This SystemVerilog line contributes to the structural or sequential description of the accelerator. Its effect is resolved by XSim for verification and by Vivado during synthesis.
\paragraph{Line 30}
\noindent\colorbox{gray!8}{\parbox{\dimexpr\linewidth-2\fboxsep\relax}{\raggedright\ttfamily\scriptsize     int errors = 0;}}
\par\smallskip
This is a local compile-time or procedural assignment. In context it either defines a helper value or connects a temporary expression used by nearby hardware logic.
\paragraph{Line 31}
\noindent\colorbox{gray!8}{\parbox{\dimexpr\linewidth-2\fboxsep\relax}{\raggedright\ttfamily\scriptsize \emph{blank line}}}
\par\smallskip
This blank line is not hardware or software behavior. It separates nearby declarations or operations so the file remains readable and reviewable.
\paragraph{Line 32}
\noindent\colorbox{gray!8}{\parbox{\dimexpr\linewidth-2\fboxsep\relax}{\raggedright\ttfamily\scriptsize     cifar10\_spikingformer\_head\_accelerator \#(}}
\par\smallskip
This SystemVerilog line contributes to the structural or sequential description of the accelerator. Its effect is resolved by XSim for verification and by Vivado during synthesis.
\paragraph{Line 33}
\noindent\colorbox{gray!8}{\parbox{\dimexpr\linewidth-2\fboxsep\relax}{\raggedright\ttfamily\scriptsize         .FEATURE\_DIM(FEATURE\_DIM),}}
\par\smallskip
This SystemVerilog line contributes to the structural or sequential description of the accelerator. Its effect is resolved by XSim for verification and by Vivado during synthesis.
\paragraph{Line 34}
\noindent\colorbox{gray!8}{\parbox{\dimexpr\linewidth-2\fboxsep\relax}{\raggedright\ttfamily\scriptsize         .NUM\_CLASSES(NUM\_CLASSES),}}
\par\smallskip
This SystemVerilog line contributes to the structural or sequential description of the accelerator. Its effect is resolved by XSim for verification and by Vivado during synthesis.
\paragraph{Line 35}
\noindent\colorbox{gray!8}{\parbox{\dimexpr\linewidth-2\fboxsep\relax}{\raggedright\ttfamily\scriptsize         .FEATURE\_BITS(16),}}
\par\smallskip
This SystemVerilog line contributes to the structural or sequential description of the accelerator. Its effect is resolved by XSim for verification and by Vivado during synthesis.
\paragraph{Line 36}
\noindent\colorbox{gray!8}{\parbox{\dimexpr\linewidth-2\fboxsep\relax}{\raggedright\ttfamily\scriptsize         .WEIGHT\_BITS(16),}}
\par\smallskip
This SystemVerilog line contributes to the structural or sequential description of the accelerator. Its effect is resolved by XSim for verification and by Vivado during synthesis.
\paragraph{Line 37}
\noindent\colorbox{gray!8}{\parbox{\dimexpr\linewidth-2\fboxsep\relax}{\raggedright\ttfamily\scriptsize         .ACC\_BITS(64),}}
\par\smallskip
This SystemVerilog line contributes to the structural or sequential description of the accelerator. Its effect is resolved by XSim for verification and by Vivado during synthesis.
\paragraph{Line 38}
\noindent\colorbox{gray!8}{\parbox{\dimexpr\linewidth-2\fboxsep\relax}{\raggedright\ttfamily\scriptsize         .ENABLE\_ONCHIP\_LEARNING(1),}}
\par\smallskip
This SystemVerilog line contributes to the structural or sequential description of the accelerator. Its effect is resolved by XSim for verification and by Vivado during synthesis.
\paragraph{Line 39}
\noindent\colorbox{gray!8}{\parbox{\dimexpr\linewidth-2\fboxsep\relax}{\raggedright\ttfamily\scriptsize         .INFER\_SKIP\_ABS\_THRESH(0),}}
\par\smallskip
This line controls the feature-activity gate. When the absolute fixed-point feature is at or below this threshold, the RTL treats it as zero so the exported Spikingformer head ignores small noisy activations and records a measurable MAC-skip opportunity.
\paragraph{Line 40}
\noindent\colorbox{gray!8}{\parbox{\dimexpr\linewidth-2\fboxsep\relax}{\raggedright\ttfamily\scriptsize         .LEARN\_SHIFT(1),}}
\par\smallskip
This SystemVerilog line contributes to the structural or sequential description of the accelerator. Its effect is resolved by XSim for verification and by Vivado during synthesis.
\paragraph{Line 41}
\noindent\colorbox{gray!8}{\parbox{\dimexpr\linewidth-2\fboxsep\relax}{\raggedright\ttfamily\scriptsize         .BIAS\_LEARN\_DELTA(3),}}
\par\smallskip
This SystemVerilog line contributes to the structural or sequential description of the accelerator. Its effect is resolved by XSim for verification and by Vivado during synthesis.
\paragraph{Line 42}
\noindent\colorbox{gray!8}{\parbox{\dimexpr\linewidth-2\fboxsep\relax}{\raggedright\ttfamily\scriptsize         .LEARN\_MARGIN\_Q(64'sd1000000)}}
\par\smallskip
This SystemVerilog line contributes to the structural or sequential description of the accelerator. Its effect is resolved by XSim for verification and by Vivado during synthesis.
\paragraph{Line 43}
\noindent\colorbox{gray!8}{\parbox{\dimexpr\linewidth-2\fboxsep\relax}{\raggedright\ttfamily\scriptsize     ) dut (}}
\par\smallskip
This SystemVerilog line contributes to the structural or sequential description of the accelerator. Its effect is resolved by XSim for verification and by Vivado during synthesis.
\paragraph{Line 44}
\noindent\colorbox{gray!8}{\parbox{\dimexpr\linewidth-2\fboxsep\relax}{\raggedright\ttfamily\scriptsize         .clk(clk),}}
\par\smallskip
This SystemVerilog line contributes to the structural or sequential description of the accelerator. Its effect is resolved by XSim for verification and by Vivado during synthesis.
\paragraph{Line 45}
\noindent\colorbox{gray!8}{\parbox{\dimexpr\linewidth-2\fboxsep\relax}{\raggedright\ttfamily\scriptsize         .rst\_n(rst\_n),}}
\par\smallskip
This SystemVerilog line contributes to the structural or sequential description of the accelerator. Its effect is resolved by XSim for verification and by Vivado during synthesis.
\paragraph{Line 46}
\noindent\colorbox{gray!8}{\parbox{\dimexpr\linewidth-2\fboxsep\relax}{\raggedright\ttfamily\scriptsize         .start(start),}}
\par\smallskip
This SystemVerilog line contributes to the structural or sequential description of the accelerator. Its effect is resolved by XSim for verification and by Vivado during synthesis.
\paragraph{Line 47}
\noindent\colorbox{gray!8}{\parbox{\dimexpr\linewidth-2\fboxsep\relax}{\raggedright\ttfamily\scriptsize         .feature\_valid(feature\_valid),}}
\par\smallskip
This SystemVerilog line contributes to the structural or sequential description of the accelerator. Its effect is resolved by XSim for verification and by Vivado during synthesis.
\paragraph{Line 48}
\noindent\colorbox{gray!8}{\parbox{\dimexpr\linewidth-2\fboxsep\relax}{\raggedright\ttfamily\scriptsize         .feature\_ready(feature\_ready),}}
\par\smallskip
This SystemVerilog line contributes to the structural or sequential description of the accelerator. Its effect is resolved by XSim for verification and by Vivado during synthesis.
\paragraph{Line 49}
\noindent\colorbox{gray!8}{\parbox{\dimexpr\linewidth-2\fboxsep\relax}{\raggedright\ttfamily\scriptsize         .feature\_value(feature\_value),}}
\par\smallskip
This SystemVerilog line contributes to the structural or sequential description of the accelerator. Its effect is resolved by XSim for verification and by Vivado during synthesis.
\paragraph{Line 50}
\noindent\colorbox{gray!8}{\parbox{\dimexpr\linewidth-2\fboxsep\relax}{\raggedright\ttfamily\scriptsize         .frame\_done(frame\_done),}}
\par\smallskip
This SystemVerilog line contributes to the structural or sequential description of the accelerator. Its effect is resolved by XSim for verification and by Vivado during synthesis.
\paragraph{Line 51}
\noindent\colorbox{gray!8}{\parbox{\dimexpr\linewidth-2\fboxsep\relax}{\raggedright\ttfamily\scriptsize         .train\_enable(train\_enable),}}
\par\smallskip
This SystemVerilog line contributes to the structural or sequential description of the accelerator. Its effect is resolved by XSim for verification and by Vivado during synthesis.
\paragraph{Line 52}
\noindent\colorbox{gray!8}{\parbox{\dimexpr\linewidth-2\fboxsep\relax}{\raggedright\ttfamily\scriptsize         .train\_label(train\_label),}}
\par\smallskip
This SystemVerilog line contributes to the structural or sequential description of the accelerator. Its effect is resolved by XSim for verification and by Vivado during synthesis.
\paragraph{Line 53}
\noindent\colorbox{gray!8}{\parbox{\dimexpr\linewidth-2\fboxsep\relax}{\raggedright\ttfamily\scriptsize         .busy(busy),}}
\par\smallskip
This SystemVerilog line contributes to the structural or sequential description of the accelerator. Its effect is resolved by XSim for verification and by Vivado during synthesis.
\paragraph{Line 54}
\noindent\colorbox{gray!8}{\parbox{\dimexpr\linewidth-2\fboxsep\relax}{\raggedright\ttfamily\scriptsize         .result\_valid(result\_valid),}}
\par\smallskip
This SystemVerilog line contributes to the structural or sequential description of the accelerator. Its effect is resolved by XSim for verification and by Vivado during synthesis.
\paragraph{Line 55}
\noindent\colorbox{gray!8}{\parbox{\dimexpr\linewidth-2\fboxsep\relax}{\raggedright\ttfamily\scriptsize         .predicted\_label(predicted\_label),}}
\par\smallskip
This SystemVerilog line contributes to the structural or sequential description of the accelerator. Its effect is resolved by XSim for verification and by Vivado during synthesis.
\paragraph{Line 56}
\noindent\colorbox{gray!8}{\parbox{\dimexpr\linewidth-2\fboxsep\relax}{\raggedright\ttfamily\scriptsize         .predicted\_score(predicted\_score),}}
\par\smallskip
This SystemVerilog line contributes to the structural or sequential description of the accelerator. Its effect is resolved by XSim for verification and by Vivado during synthesis.
\paragraph{Line 57}
\noindent\colorbox{gray!8}{\parbox{\dimexpr\linewidth-2\fboxsep\relax}{\raggedright\ttfamily\scriptsize         .train\_update\_valid(train\_update\_valid),}}
\par\smallskip
This SystemVerilog line contributes to the structural or sequential description of the accelerator. Its effect is resolved by XSim for verification and by Vivado during synthesis.
\paragraph{Line 58}
\noindent\colorbox{gray!8}{\parbox{\dimexpr\linewidth-2\fboxsep\relax}{\raggedright\ttfamily\scriptsize         .train\_update\_applied(train\_update\_applied),}}
\par\smallskip
This SystemVerilog line contributes to the structural or sequential description of the accelerator. Its effect is resolved by XSim for verification and by Vivado during synthesis.
\paragraph{Line 59}
\noindent\colorbox{gray!8}{\parbox{\dimexpr\linewidth-2\fboxsep\relax}{\raggedright\ttfamily\scriptsize         .debug\_state(debug\_state),}}
\par\smallskip
This SystemVerilog line contributes to the structural or sequential description of the accelerator. Its effect is resolved by XSim for verification and by Vivado during synthesis.
\paragraph{Line 60}
\noindent\colorbox{gray!8}{\parbox{\dimexpr\linewidth-2\fboxsep\relax}{\raggedright\ttfamily\scriptsize         .debug\_feature(debug\_feature),}}
\par\smallskip
This SystemVerilog line contributes to the structural or sequential description of the accelerator. Its effect is resolved by XSim for verification and by Vivado during synthesis.
\paragraph{Line 61}
\noindent\colorbox{gray!8}{\parbox{\dimexpr\linewidth-2\fboxsep\relax}{\raggedright\ttfamily\scriptsize         .debug\_class(debug\_class),}}
\par\smallskip
This SystemVerilog line contributes to the structural or sequential description of the accelerator. Its effect is resolved by XSim for verification and by Vivado during synthesis.
\paragraph{Line 62}
\noindent\colorbox{gray!8}{\parbox{\dimexpr\linewidth-2\fboxsep\relax}{\raggedright\ttfamily\scriptsize         .debug\_update\_count(debug\_update\_count),}}
\par\smallskip
This SystemVerilog line contributes to the structural or sequential description of the accelerator. Its effect is resolved by XSim for verification and by Vivado during synthesis.
\paragraph{Line 63}
\noindent\colorbox{gray!8}{\parbox{\dimexpr\linewidth-2\fboxsep\relax}{\raggedright\ttfamily\scriptsize         .debug\_infer\_skip\_count(debug\_infer\_skip\_count),}}
\par\smallskip
This line exposes or updates a hardware activity counter. The counter makes the feature gate visible in waveforms and CSV traces, so the power-saving opportunity is observable from RTL rather than only from Python statistics.
\paragraph{Line 64}
\noindent\colorbox{gray!8}{\parbox{\dimexpr\linewidth-2\fboxsep\relax}{\raggedright\ttfamily\scriptsize         .debug\_train\_skip\_count(debug\_train\_skip\_count)}}
\par\smallskip
This line exposes or updates a hardware activity counter. The counter makes the feature gate visible in waveforms and CSV traces, so the power-saving opportunity is observable from RTL rather than only from Python statistics.
\paragraph{Line 65}
\noindent\colorbox{gray!8}{\parbox{\dimexpr\linewidth-2\fboxsep\relax}{\raggedright\ttfamily\scriptsize     );}}
\par\smallskip
This SystemVerilog line contributes to the structural or sequential description of the accelerator. Its effect is resolved by XSim for verification and by Vivado during synthesis.
\paragraph{Line 66}
\noindent\colorbox{gray!8}{\parbox{\dimexpr\linewidth-2\fboxsep\relax}{\raggedright\ttfamily\scriptsize \emph{blank line}}}
\par\smallskip
This blank line is not hardware or software behavior. It separates nearby declarations or operations so the file remains readable and reviewable.
\paragraph{Line 67}
\noindent\colorbox{gray!8}{\parbox{\dimexpr\linewidth-2\fboxsep\relax}{\raggedright\ttfamily\scriptsize     always \#5 clk = \textasciitilde{}clk;}}
\par\smallskip
This is a local compile-time or procedural assignment. In context it either defines a helper value or connects a temporary expression used by nearby hardware logic.
\paragraph{Line 68}
\noindent\colorbox{gray!8}{\parbox{\dimexpr\linewidth-2\fboxsep\relax}{\raggedright\ttfamily\scriptsize \emph{blank line}}}
\par\smallskip
This blank line is not hardware or software behavior. It separates nearby declarations or operations so the file remains readable and reviewable.
\paragraph{Line 69}
\noindent\colorbox{gray!8}{\parbox{\dimexpr\linewidth-2\fboxsep\relax}{\raggedright\ttfamily\scriptsize     task automatic expect\_i16(input string name, input int got, input int expected);}}
\par\smallskip
This line declares a module port. A port is a physical or wrapper-visible signal used to move clocks, control pulses, features, labels, predictions, or debug information across module boundaries.
\paragraph{Line 70}
\noindent\colorbox{gray!8}{\parbox{\dimexpr\linewidth-2\fboxsep\relax}{\raggedright\ttfamily\scriptsize         begin}}
\par\smallskip
This SystemVerilog line contributes to the structural or sequential description of the accelerator. Its effect is resolved by XSim for verification and by Vivado during synthesis.
\paragraph{Line 71}
\noindent\colorbox{gray!8}{\parbox{\dimexpr\linewidth-2\fboxsep\relax}{\raggedright\ttfamily\scriptsize             if (got !== expected) begin}}
\par\smallskip
This is a local compile-time or procedural assignment. In context it either defines a helper value or connects a temporary expression used by nearby hardware logic.
\paragraph{Line 72}
\noindent\colorbox{gray!8}{\parbox{\dimexpr\linewidth-2\fboxsep\relax}{\raggedright\ttfamily\scriptsize                 \$display("ERROR: \%s got=\%0d expected=\%0d", name, got, expected);}}
\par\smallskip
This is a local compile-time or procedural assignment. In context it either defines a helper value or connects a temporary expression used by nearby hardware logic.
\paragraph{Line 73}
\noindent\colorbox{gray!8}{\parbox{\dimexpr\linewidth-2\fboxsep\relax}{\raggedright\ttfamily\scriptsize                 errors++;}}
\par\smallskip
This SystemVerilog line contributes to the structural or sequential description of the accelerator. Its effect is resolved by XSim for verification and by Vivado during synthesis.
\paragraph{Line 74}
\noindent\colorbox{gray!8}{\parbox{\dimexpr\linewidth-2\fboxsep\relax}{\raggedright\ttfamily\scriptsize             end}}
\par\smallskip
This SystemVerilog line contributes to the structural or sequential description of the accelerator. Its effect is resolved by XSim for verification and by Vivado during synthesis.
\paragraph{Line 75}
\noindent\colorbox{gray!8}{\parbox{\dimexpr\linewidth-2\fboxsep\relax}{\raggedright\ttfamily\scriptsize         end}}
\par\smallskip
This SystemVerilog line contributes to the structural or sequential description of the accelerator. Its effect is resolved by XSim for verification and by Vivado during synthesis.
\paragraph{Line 76}
\noindent\colorbox{gray!8}{\parbox{\dimexpr\linewidth-2\fboxsep\relax}{\raggedright\ttfamily\scriptsize     endtask}}
\par\smallskip
This SystemVerilog line contributes to the structural or sequential description of the accelerator. Its effect is resolved by XSim for verification and by Vivado during synthesis.
\paragraph{Line 77}
\noindent\colorbox{gray!8}{\parbox{\dimexpr\linewidth-2\fboxsep\relax}{\raggedright\ttfamily\scriptsize \emph{blank line}}}
\par\smallskip
This blank line is not hardware or software behavior. It separates nearby declarations or operations so the file remains readable and reviewable.
\paragraph{Line 78}
\noindent\colorbox{gray!8}{\parbox{\dimexpr\linewidth-2\fboxsep\relax}{\raggedright\ttfamily\scriptsize     initial begin}}
\par\smallskip
This simulation/synthesis initialization block loads trained memories from generated .mem files. On FPGA, Vivado maps these initial values into BRAM/LUTRAM initialization data.
\paragraph{Line 79}
\noindent\colorbox{gray!8}{\parbox{\dimexpr\linewidth-2\fboxsep\relax}{\raggedright\ttfamily\scriptsize         x[0] = 16'sd4;}}
\par\smallskip
This is a local compile-time or procedural assignment. In context it either defines a helper value or connects a temporary expression used by nearby hardware logic.
\paragraph{Line 80}
\noindent\colorbox{gray!8}{\parbox{\dimexpr\linewidth-2\fboxsep\relax}{\raggedright\ttfamily\scriptsize         x[1] = -16'sd2;}}
\par\smallskip
This is a local compile-time or procedural assignment. In context it either defines a helper value or connects a temporary expression used by nearby hardware logic.
\paragraph{Line 81}
\noindent\colorbox{gray!8}{\parbox{\dimexpr\linewidth-2\fboxsep\relax}{\raggedright\ttfamily\scriptsize         x[2] = 16'sd6;}}
\par\smallskip
This is a local compile-time or procedural assignment. In context it either defines a helper value or connects a temporary expression used by nearby hardware logic.
\paragraph{Line 82}
\noindent\colorbox{gray!8}{\parbox{\dimexpr\linewidth-2\fboxsep\relax}{\raggedright\ttfamily\scriptsize         x[3] = -16'sd8;}}
\par\smallskip
This is a local compile-time or procedural assignment. In context it either defines a helper value or connects a temporary expression used by nearby hardware logic.
\paragraph{Line 83}
\noindent\colorbox{gray!8}{\parbox{\dimexpr\linewidth-2\fboxsep\relax}{\raggedright\ttfamily\scriptsize \emph{blank line}}}
\par\smallskip
This blank line is not hardware or software behavior. It separates nearby declarations or operations so the file remains readable and reviewable.
\paragraph{Line 84}
\noindent\colorbox{gray!8}{\parbox{\dimexpr\linewidth-2\fboxsep\relax}{\raggedright\ttfamily\scriptsize         start = 1'b0;}}
\par\smallskip
This is a local compile-time or procedural assignment. In context it either defines a helper value or connects a temporary expression used by nearby hardware logic.
\paragraph{Line 85}
\noindent\colorbox{gray!8}{\parbox{\dimexpr\linewidth-2\fboxsep\relax}{\raggedright\ttfamily\scriptsize         feature\_valid = 1'b0;}}
\par\smallskip
This is a local compile-time or procedural assignment. In context it either defines a helper value or connects a temporary expression used by nearby hardware logic.
\paragraph{Line 86}
\noindent\colorbox{gray!8}{\parbox{\dimexpr\linewidth-2\fboxsep\relax}{\raggedright\ttfamily\scriptsize         feature\_value = 16'sd0;}}
\par\smallskip
This is a local compile-time or procedural assignment. In context it either defines a helper value or connects a temporary expression used by nearby hardware logic.
\paragraph{Line 87}
\noindent\colorbox{gray!8}{\parbox{\dimexpr\linewidth-2\fboxsep\relax}{\raggedright\ttfamily\scriptsize         frame\_done = 1'b0;}}
\par\smallskip
This is a local compile-time or procedural assignment. In context it either defines a helper value or connects a temporary expression used by nearby hardware logic.
\paragraph{Line 88}
\noindent\colorbox{gray!8}{\parbox{\dimexpr\linewidth-2\fboxsep\relax}{\raggedright\ttfamily\scriptsize         train\_enable = 1'b0;}}
\par\smallskip
This is a local compile-time or procedural assignment. In context it either defines a helper value or connects a temporary expression used by nearby hardware logic.
\paragraph{Line 89}
\noindent\colorbox{gray!8}{\parbox{\dimexpr\linewidth-2\fboxsep\relax}{\raggedright\ttfamily\scriptsize         train\_label = 4'd1;}}
\par\smallskip
This is a local compile-time or procedural assignment. In context it either defines a helper value or connects a temporary expression used by nearby hardware logic.
\paragraph{Line 90}
\noindent\colorbox{gray!8}{\parbox{\dimexpr\linewidth-2\fboxsep\relax}{\raggedright\ttfamily\scriptsize \emph{blank line}}}
\par\smallskip
This blank line is not hardware or software behavior. It separates nearby declarations or operations so the file remains readable and reviewable.
\paragraph{Line 91}
\noindent\colorbox{gray!8}{\parbox{\dimexpr\linewidth-2\fboxsep\relax}{\raggedright\ttfamily\scriptsize         repeat (8) @(negedge clk);}}
\par\smallskip
This SystemVerilog line contributes to the structural or sequential description of the accelerator. Its effect is resolved by XSim for verification and by Vivado during synthesis.
\paragraph{Line 92}
\noindent\colorbox{gray!8}{\parbox{\dimexpr\linewidth-2\fboxsep\relax}{\raggedright\ttfamily\scriptsize         for (int i = 0; i < NUM\_CLASSES * FEATURE\_DIM; i++) begin}}
\par\smallskip
This is a local compile-time or procedural assignment. In context it either defines a helper value or connects a temporary expression used by nearby hardware logic.
\paragraph{Line 93}
\noindent\colorbox{gray!8}{\parbox{\dimexpr\linewidth-2\fboxsep\relax}{\raggedright\ttfamily\scriptsize             dut.weight\_mem[i] = 16'sd0;}}
\par\smallskip
This line touches the trained classifier weight store. In the final RTL the large 10 x 384 int16 matrix is inferred as true dual-port BRAM, so weights remain on chip and can also be updated by the learning state machine.
\paragraph{Line 94}
\noindent\colorbox{gray!8}{\parbox{\dimexpr\linewidth-2\fboxsep\relax}{\raggedright\ttfamily\scriptsize         end}}
\par\smallskip
This SystemVerilog line contributes to the structural or sequential description of the accelerator. Its effect is resolved by XSim for verification and by Vivado during synthesis.
\paragraph{Line 95}
\noindent\colorbox{gray!8}{\parbox{\dimexpr\linewidth-2\fboxsep\relax}{\raggedright\ttfamily\scriptsize         dut.bias\_mem[0] = 64'sd0;}}
\par\smallskip
This line touches the class-bias store. Bias values are part of the affine classifier score and are updated on chip in separate one-write states to keep synthesis clean.
\paragraph{Line 96}
\noindent\colorbox{gray!8}{\parbox{\dimexpr\linewidth-2\fboxsep\relax}{\raggedright\ttfamily\scriptsize         dut.bias\_mem[1] = 64'sd0;}}
\par\smallskip
This line touches the class-bias store. Bias values are part of the affine classifier score and are updated on chip in separate one-write states to keep synthesis clean.
\paragraph{Line 97}
\noindent\colorbox{gray!8}{\parbox{\dimexpr\linewidth-2\fboxsep\relax}{\raggedright\ttfamily\scriptsize         dut.bias\_mem[2] = 64'sd20;}}
\par\smallskip
This line touches the class-bias store. Bias values are part of the affine classifier score and are updated on chip in separate one-write states to keep synthesis clean.
\paragraph{Line 98}
\noindent\colorbox{gray!8}{\parbox{\dimexpr\linewidth-2\fboxsep\relax}{\raggedright\ttfamily\scriptsize         rst\_n = 1'b1;}}
\par\smallskip
This is a local compile-time or procedural assignment. In context it either defines a helper value or connects a temporary expression used by nearby hardware logic.
\paragraph{Line 99}
\noindent\colorbox{gray!8}{\parbox{\dimexpr\linewidth-2\fboxsep\relax}{\raggedright\ttfamily\scriptsize         repeat (4) @(negedge clk);}}
\par\smallskip
This SystemVerilog line contributes to the structural or sequential description of the accelerator. Its effect is resolved by XSim for verification and by Vivado during synthesis.
\paragraph{Line 100}
\noindent\colorbox{gray!8}{\parbox{\dimexpr\linewidth-2\fboxsep\relax}{\raggedright\ttfamily\scriptsize \emph{blank line}}}
\par\smallskip
This blank line is not hardware or software behavior. It separates nearby declarations or operations so the file remains readable and reviewable.
\paragraph{Line 101}
\noindent\colorbox{gray!8}{\parbox{\dimexpr\linewidth-2\fboxsep\relax}{\raggedright\ttfamily\scriptsize         @(negedge clk);}}
\par\smallskip
This SystemVerilog line contributes to the structural or sequential description of the accelerator. Its effect is resolved by XSim for verification and by Vivado during synthesis.
\paragraph{Line 102}
\noindent\colorbox{gray!8}{\parbox{\dimexpr\linewidth-2\fboxsep\relax}{\raggedright\ttfamily\scriptsize         train\_enable = 1'b1;}}
\par\smallskip
This is a local compile-time or procedural assignment. In context it either defines a helper value or connects a temporary expression used by nearby hardware logic.
\paragraph{Line 103}
\noindent\colorbox{gray!8}{\parbox{\dimexpr\linewidth-2\fboxsep\relax}{\raggedright\ttfamily\scriptsize         train\_label = 4'd1;}}
\par\smallskip
This is a local compile-time or procedural assignment. In context it either defines a helper value or connects a temporary expression used by nearby hardware logic.
\paragraph{Line 104}
\noindent\colorbox{gray!8}{\parbox{\dimexpr\linewidth-2\fboxsep\relax}{\raggedright\ttfamily\scriptsize         start = 1'b1;}}
\par\smallskip
This is a local compile-time or procedural assignment. In context it either defines a helper value or connects a temporary expression used by nearby hardware logic.
\paragraph{Line 105}
\noindent\colorbox{gray!8}{\parbox{\dimexpr\linewidth-2\fboxsep\relax}{\raggedright\ttfamily\scriptsize         @(negedge clk);}}
\par\smallskip
This SystemVerilog line contributes to the structural or sequential description of the accelerator. Its effect is resolved by XSim for verification and by Vivado during synthesis.
\paragraph{Line 106}
\noindent\colorbox{gray!8}{\parbox{\dimexpr\linewidth-2\fboxsep\relax}{\raggedright\ttfamily\scriptsize         start = 1'b0;}}
\par\smallskip
This is a local compile-time or procedural assignment. In context it either defines a helper value or connects a temporary expression used by nearby hardware logic.
\paragraph{Line 107}
\noindent\colorbox{gray!8}{\parbox{\dimexpr\linewidth-2\fboxsep\relax}{\raggedright\ttfamily\scriptsize \emph{blank line}}}
\par\smallskip
This blank line is not hardware or software behavior. It separates nearby declarations or operations so the file remains readable and reviewable.
\paragraph{Line 108}
\noindent\colorbox{gray!8}{\parbox{\dimexpr\linewidth-2\fboxsep\relax}{\raggedright\ttfamily\scriptsize         while (!feature\_ready) @(posedge clk);}}
\par\smallskip
This SystemVerilog line contributes to the structural or sequential description of the accelerator. Its effect is resolved by XSim for verification and by Vivado during synthesis.
\paragraph{Line 109}
\noindent\colorbox{gray!8}{\parbox{\dimexpr\linewidth-2\fboxsep\relax}{\raggedright\ttfamily\scriptsize         for (int f = 0; f < FEATURE\_DIM; f++) begin}}
\par\smallskip
This is a local compile-time or procedural assignment. In context it either defines a helper value or connects a temporary expression used by nearby hardware logic.
\paragraph{Line 110}
\noindent\colorbox{gray!8}{\parbox{\dimexpr\linewidth-2\fboxsep\relax}{\raggedright\ttfamily\scriptsize             @(negedge clk);}}
\par\smallskip
This SystemVerilog line contributes to the structural or sequential description of the accelerator. Its effect is resolved by XSim for verification and by Vivado during synthesis.
\paragraph{Line 111}
\noindent\colorbox{gray!8}{\parbox{\dimexpr\linewidth-2\fboxsep\relax}{\raggedright\ttfamily\scriptsize             feature\_value = x[f];}}
\par\smallskip
This is a local compile-time or procedural assignment. In context it either defines a helper value or connects a temporary expression used by nearby hardware logic.
\paragraph{Line 112}
\noindent\colorbox{gray!8}{\parbox{\dimexpr\linewidth-2\fboxsep\relax}{\raggedright\ttfamily\scriptsize             feature\_valid = 1'b1;}}
\par\smallskip
This is a local compile-time or procedural assignment. In context it either defines a helper value or connects a temporary expression used by nearby hardware logic.
\paragraph{Line 113}
\noindent\colorbox{gray!8}{\parbox{\dimexpr\linewidth-2\fboxsep\relax}{\raggedright\ttfamily\scriptsize             @(posedge clk);}}
\par\smallskip
This SystemVerilog line contributes to the structural or sequential description of the accelerator. Its effect is resolved by XSim for verification and by Vivado during synthesis.
\paragraph{Line 114}
\noindent\colorbox{gray!8}{\parbox{\dimexpr\linewidth-2\fboxsep\relax}{\raggedright\ttfamily\scriptsize         end}}
\par\smallskip
This SystemVerilog line contributes to the structural or sequential description of the accelerator. Its effect is resolved by XSim for verification and by Vivado during synthesis.
\paragraph{Line 115}
\noindent\colorbox{gray!8}{\parbox{\dimexpr\linewidth-2\fboxsep\relax}{\raggedright\ttfamily\scriptsize         @(negedge clk);}}
\par\smallskip
This SystemVerilog line contributes to the structural or sequential description of the accelerator. Its effect is resolved by XSim for verification and by Vivado during synthesis.
\paragraph{Line 116}
\noindent\colorbox{gray!8}{\parbox{\dimexpr\linewidth-2\fboxsep\relax}{\raggedright\ttfamily\scriptsize         feature\_valid = 1'b0;}}
\par\smallskip
This is a local compile-time or procedural assignment. In context it either defines a helper value or connects a temporary expression used by nearby hardware logic.
\paragraph{Line 117}
\noindent\colorbox{gray!8}{\parbox{\dimexpr\linewidth-2\fboxsep\relax}{\raggedright\ttfamily\scriptsize         frame\_done = 1'b1;}}
\par\smallskip
This is a local compile-time or procedural assignment. In context it either defines a helper value or connects a temporary expression used by nearby hardware logic.
\paragraph{Line 118}
\noindent\colorbox{gray!8}{\parbox{\dimexpr\linewidth-2\fboxsep\relax}{\raggedright\ttfamily\scriptsize         @(negedge clk);}}
\par\smallskip
This SystemVerilog line contributes to the structural or sequential description of the accelerator. Its effect is resolved by XSim for verification and by Vivado during synthesis.
\paragraph{Line 119}
\noindent\colorbox{gray!8}{\parbox{\dimexpr\linewidth-2\fboxsep\relax}{\raggedright\ttfamily\scriptsize         frame\_done = 1'b0;}}
\par\smallskip
This is a local compile-time or procedural assignment. In context it either defines a helper value or connects a temporary expression used by nearby hardware logic.
\paragraph{Line 120}
\noindent\colorbox{gray!8}{\parbox{\dimexpr\linewidth-2\fboxsep\relax}{\raggedright\ttfamily\scriptsize \emph{blank line}}}
\par\smallskip
This blank line is not hardware or software behavior. It separates nearby declarations or operations so the file remains readable and reviewable.
\paragraph{Line 121}
\noindent\colorbox{gray!8}{\parbox{\dimexpr\linewidth-2\fboxsep\relax}{\raggedright\ttfamily\scriptsize         while (!result\_valid) @(posedge clk);}}
\par\smallskip
This SystemVerilog line contributes to the structural or sequential description of the accelerator. Its effect is resolved by XSim for verification and by Vivado during synthesis.
\paragraph{Line 122}
\noindent\colorbox{gray!8}{\parbox{\dimexpr\linewidth-2\fboxsep\relax}{\raggedright\ttfamily\scriptsize         if (predicted\_label !== 4'd2) begin}}
\par\smallskip
This is a local compile-time or procedural assignment. In context it either defines a helper value or connects a temporary expression used by nearby hardware logic.
\paragraph{Line 123}
\noindent\colorbox{gray!8}{\parbox{\dimexpr\linewidth-2\fboxsep\relax}{\raggedright\ttfamily\scriptsize             \$display("ERROR: predicted\_label got=\%0d expected=2", predicted\_label);}}
\par\smallskip
This is a local compile-time or procedural assignment. In context it either defines a helper value or connects a temporary expression used by nearby hardware logic.
\paragraph{Line 124}
\noindent\colorbox{gray!8}{\parbox{\dimexpr\linewidth-2\fboxsep\relax}{\raggedright\ttfamily\scriptsize             errors++;}}
\par\smallskip
This SystemVerilog line contributes to the structural or sequential description of the accelerator. Its effect is resolved by XSim for verification and by Vivado during synthesis.
\paragraph{Line 125}
\noindent\colorbox{gray!8}{\parbox{\dimexpr\linewidth-2\fboxsep\relax}{\raggedright\ttfamily\scriptsize         end}}
\par\smallskip
This SystemVerilog line contributes to the structural or sequential description of the accelerator. Its effect is resolved by XSim for verification and by Vivado during synthesis.
\paragraph{Line 126}
\noindent\colorbox{gray!8}{\parbox{\dimexpr\linewidth-2\fboxsep\relax}{\raggedright\ttfamily\scriptsize         if (!train\_update\_applied \&\& debug\_update\_count != 32'd1) begin}}
\par\smallskip
This is a local compile-time or procedural assignment. In context it either defines a helper value or connects a temporary expression used by nearby hardware logic.
\paragraph{Line 127}
\noindent\colorbox{gray!8}{\parbox{\dimexpr\linewidth-2\fboxsep\relax}{\raggedright\ttfamily\scriptsize             \$display("ERROR: update did not apply");}}
\par\smallskip
This SystemVerilog line contributes to the structural or sequential description of the accelerator. Its effect is resolved by XSim for verification and by Vivado during synthesis.
\paragraph{Line 128}
\noindent\colorbox{gray!8}{\parbox{\dimexpr\linewidth-2\fboxsep\relax}{\raggedright\ttfamily\scriptsize             errors++;}}
\par\smallskip
This SystemVerilog line contributes to the structural or sequential description of the accelerator. Its effect is resolved by XSim for verification and by Vivado during synthesis.
\paragraph{Line 129}
\noindent\colorbox{gray!8}{\parbox{\dimexpr\linewidth-2\fboxsep\relax}{\raggedright\ttfamily\scriptsize         end}}
\par\smallskip
This SystemVerilog line contributes to the structural or sequential description of the accelerator. Its effect is resolved by XSim for verification and by Vivado during synthesis.
\paragraph{Line 130}
\noindent\colorbox{gray!8}{\parbox{\dimexpr\linewidth-2\fboxsep\relax}{\raggedright\ttfamily\scriptsize \emph{blank line}}}
\par\smallskip
This blank line is not hardware or software behavior. It separates nearby declarations or operations so the file remains readable and reviewable.
\paragraph{Line 131}
\noindent\colorbox{gray!8}{\parbox{\dimexpr\linewidth-2\fboxsep\relax}{\raggedright\ttfamily\scriptsize         while (busy) @(posedge clk);}}
\par\smallskip
This SystemVerilog line contributes to the structural or sequential description of the accelerator. Its effect is resolved by XSim for verification and by Vivado during synthesis.
\paragraph{Line 132}
\noindent\colorbox{gray!8}{\parbox{\dimexpr\linewidth-2\fboxsep\relax}{\raggedright\ttfamily\scriptsize         expect\_i16("class1 weight0", dut.weight\_mem[1 * FEATURE\_DIM + 0], 2);}}
\par\smallskip
This line touches the trained classifier weight store. In the final RTL the large 10 x 384 int16 matrix is inferred as true dual-port BRAM, so weights remain on chip and can also be updated by the learning state machine.
\paragraph{Line 133}
\noindent\colorbox{gray!8}{\parbox{\dimexpr\linewidth-2\fboxsep\relax}{\raggedright\ttfamily\scriptsize         expect\_i16("class1 weight1", dut.weight\_mem[1 * FEATURE\_DIM + 1], -1);}}
\par\smallskip
This line touches the trained classifier weight store. In the final RTL the large 10 x 384 int16 matrix is inferred as true dual-port BRAM, so weights remain on chip and can also be updated by the learning state machine.
\paragraph{Line 134}
\noindent\colorbox{gray!8}{\parbox{\dimexpr\linewidth-2\fboxsep\relax}{\raggedright\ttfamily\scriptsize         expect\_i16("class1 weight2", dut.weight\_mem[1 * FEATURE\_DIM + 2], 3);}}
\par\smallskip
This line touches the trained classifier weight store. In the final RTL the large 10 x 384 int16 matrix is inferred as true dual-port BRAM, so weights remain on chip and can also be updated by the learning state machine.
\paragraph{Line 135}
\noindent\colorbox{gray!8}{\parbox{\dimexpr\linewidth-2\fboxsep\relax}{\raggedright\ttfamily\scriptsize         expect\_i16("class1 weight3", dut.weight\_mem[1 * FEATURE\_DIM + 3], -4);}}
\par\smallskip
This line touches the trained classifier weight store. In the final RTL the large 10 x 384 int16 matrix is inferred as true dual-port BRAM, so weights remain on chip and can also be updated by the learning state machine.
\paragraph{Line 136}
\noindent\colorbox{gray!8}{\parbox{\dimexpr\linewidth-2\fboxsep\relax}{\raggedright\ttfamily\scriptsize         expect\_i16("class2 weight0", dut.weight\_mem[2 * FEATURE\_DIM + 0], -2);}}
\par\smallskip
This line touches the trained classifier weight store. In the final RTL the large 10 x 384 int16 matrix is inferred as true dual-port BRAM, so weights remain on chip and can also be updated by the learning state machine.
\paragraph{Line 137}
\noindent\colorbox{gray!8}{\parbox{\dimexpr\linewidth-2\fboxsep\relax}{\raggedright\ttfamily\scriptsize         expect\_i16("class2 weight1", dut.weight\_mem[2 * FEATURE\_DIM + 1], 1);}}
\par\smallskip
This line touches the trained classifier weight store. In the final RTL the large 10 x 384 int16 matrix is inferred as true dual-port BRAM, so weights remain on chip and can also be updated by the learning state machine.
\paragraph{Line 138}
\noindent\colorbox{gray!8}{\parbox{\dimexpr\linewidth-2\fboxsep\relax}{\raggedright\ttfamily\scriptsize         expect\_i16("class2 weight2", dut.weight\_mem[2 * FEATURE\_DIM + 2], -3);}}
\par\smallskip
This line touches the trained classifier weight store. In the final RTL the large 10 x 384 int16 matrix is inferred as true dual-port BRAM, so weights remain on chip and can also be updated by the learning state machine.
\paragraph{Line 139}
\noindent\colorbox{gray!8}{\parbox{\dimexpr\linewidth-2\fboxsep\relax}{\raggedright\ttfamily\scriptsize         expect\_i16("class2 weight3", dut.weight\_mem[2 * FEATURE\_DIM + 3], 4);}}
\par\smallskip
This line touches the trained classifier weight store. In the final RTL the large 10 x 384 int16 matrix is inferred as true dual-port BRAM, so weights remain on chip and can also be updated by the learning state machine.
\paragraph{Line 140}
\noindent\colorbox{gray!8}{\parbox{\dimexpr\linewidth-2\fboxsep\relax}{\raggedright\ttfamily\scriptsize         if (\$signed(dut.bias\_mem[1]) !== 64'sd3 || \$signed(dut.bias\_mem[2]) !== 64'sd17) begin}}
\par\smallskip
This line touches the class-bias store. Bias values are part of the affine classifier score and are updated on chip in separate one-write states to keep synthesis clean.
\paragraph{Line 141}
\noindent\colorbox{gray!8}{\parbox{\dimexpr\linewidth-2\fboxsep\relax}{\raggedright\ttfamily\scriptsize             \$display("ERROR: bias update got class1=\%0d class2=\%0d", \$signed(dut.bias\_mem[1]), \$signed(dut.bias\_mem[2]));}}
\par\smallskip
This line touches the class-bias store. Bias values are part of the affine classifier score and are updated on chip in separate one-write states to keep synthesis clean.
\paragraph{Line 142}
\noindent\colorbox{gray!8}{\parbox{\dimexpr\linewidth-2\fboxsep\relax}{\raggedright\ttfamily\scriptsize             errors++;}}
\par\smallskip
This SystemVerilog line contributes to the structural or sequential description of the accelerator. Its effect is resolved by XSim for verification and by Vivado during synthesis.
\paragraph{Line 143}
\noindent\colorbox{gray!8}{\parbox{\dimexpr\linewidth-2\fboxsep\relax}{\raggedright\ttfamily\scriptsize         end}}
\par\smallskip
This SystemVerilog line contributes to the structural or sequential description of the accelerator. Its effect is resolved by XSim for verification and by Vivado during synthesis.
\paragraph{Line 144}
\noindent\colorbox{gray!8}{\parbox{\dimexpr\linewidth-2\fboxsep\relax}{\raggedright\ttfamily\scriptsize \emph{blank line}}}
\par\smallskip
This blank line is not hardware or software behavior. It separates nearby declarations or operations so the file remains readable and reviewable.
\paragraph{Line 145}
\noindent\colorbox{gray!8}{\parbox{\dimexpr\linewidth-2\fboxsep\relax}{\raggedright\ttfamily\scriptsize         \$display("tb\_spikingformer\_head\_learning\_unit errors=\%0d", errors);}}
\par\smallskip
This is a local compile-time or procedural assignment. In context it either defines a helper value or connects a temporary expression used by nearby hardware logic.
\paragraph{Line 146}
\noindent\colorbox{gray!8}{\parbox{\dimexpr\linewidth-2\fboxsep\relax}{\raggedright\ttfamily\scriptsize         if (errors != 0) \$fatal;}}
\par\smallskip
This is a local compile-time or procedural assignment. In context it either defines a helper value or connects a temporary expression used by nearby hardware logic.
\paragraph{Line 147}
\noindent\colorbox{gray!8}{\parbox{\dimexpr\linewidth-2\fboxsep\relax}{\raggedright\ttfamily\scriptsize         \$finish;}}
\par\smallskip
This SystemVerilog line contributes to the structural or sequential description of the accelerator. Its effect is resolved by XSim for verification and by Vivado during synthesis.
\paragraph{Line 148}
\noindent\colorbox{gray!8}{\parbox{\dimexpr\linewidth-2\fboxsep\relax}{\raggedright\ttfamily\scriptsize     end}}
\par\smallskip
This SystemVerilog line contributes to the structural or sequential description of the accelerator. Its effect is resolved by XSim for verification and by Vivado during synthesis.
\paragraph{Line 149}
\noindent\colorbox{gray!8}{\parbox{\dimexpr\linewidth-2\fboxsep\relax}{\raggedright\ttfamily\scriptsize endmodule}}
\par\smallskip
This closes the module definition. Synthesis stops adding logic to this circuit block at this point.
\clearpage
\subsection{cifar10\_neuromorphic\textbackslash{}rtl\textbackslash{}tb\_cifar10\_spikingformer\_head\_accelerator.sv}
Each entry below gives the exact source line followed by its hardware/software role in the implemented accelerator pipeline.
\paragraph{Line 1}
\noindent\colorbox{gray!8}{\parbox{\dimexpr\linewidth-2\fboxsep\relax}{\raggedright\ttfamily\scriptsize `timescale 1ns / 1ps}}
\par\smallskip
This simulator directive sets the time unit and precision used by testbenches and waveform timestamps. It keeps delays such as clock half-periods numerically meaningful in XSim.
\paragraph{Line 2}
\noindent\colorbox{gray!8}{\parbox{\dimexpr\linewidth-2\fboxsep\relax}{\raggedright\ttfamily\scriptsize `include "rtl/cifar10\_spikingformer\_head\_params.svh"}}
\par\smallskip
This includes shared compile-time constants such as feature dimensions and fixed-point widths. The RTL and testbenches therefore agree on the same trained-model shape.
\paragraph{Line 3}
\noindent\colorbox{gray!8}{\parbox{\dimexpr\linewidth-2\fboxsep\relax}{\raggedright\ttfamily\scriptsize \emph{blank line}}}
\par\smallskip
This blank line is not hardware or software behavior. It separates nearby declarations or operations so the file remains readable and reviewable.
\paragraph{Line 4}
\noindent\colorbox{gray!8}{\parbox{\dimexpr\linewidth-2\fboxsep\relax}{\raggedright\ttfamily\scriptsize module tb\_cifar10\_spikingformer\_head\_accelerator;}}
\par\smallskip
This begins a hardware module, which is the FPGA equivalent of a reusable circuit block. Everything until the matching endmodule describes ports, storage, combinational equations, and clocked behavior for this block.
\paragraph{Line 5}
\noindent\colorbox{gray!8}{\parbox{\dimexpr\linewidth-2\fboxsep\relax}{\raggedright\ttfamily\scriptsize     parameter int NUM\_SAMPLES = `CIFAR10\_SF\_HEAD\_NUM\_SAMPLES;}}
\par\smallskip
This is a compile-time hardware setting. Vivado resolves it before synthesis, so it controls the generated circuit size, fixed-point width, or learning rule without runtime overhead.
\paragraph{Line 6}
\noindent\colorbox{gray!8}{\parbox{\dimexpr\linewidth-2\fboxsep\relax}{\raggedright\ttfamily\scriptsize     parameter int TRACE\_SAMPLES = `CIFAR10\_SF\_HEAD\_TRACE\_SAMPLES;}}
\par\smallskip
This is a compile-time hardware setting. Vivado resolves it before synthesis, so it controls the generated circuit size, fixed-point width, or learning rule without runtime overhead.
\paragraph{Line 7}
\noindent\colorbox{gray!8}{\parbox{\dimexpr\linewidth-2\fboxsep\relax}{\raggedright\ttfamily\scriptsize     parameter int FEATURE\_DIM = `CIFAR10\_SF\_HEAD\_FEATURE\_DIM;}}
\par\smallskip
This is a compile-time hardware setting. Vivado resolves it before synthesis, so it controls the generated circuit size, fixed-point width, or learning rule without runtime overhead.
\paragraph{Line 8}
\noindent\colorbox{gray!8}{\parbox{\dimexpr\linewidth-2\fboxsep\relax}{\raggedright\ttfamily\scriptsize     parameter int MAX\_SAMPLES = 10000;}}
\par\smallskip
This is a compile-time hardware setting. Vivado resolves it before synthesis, so it controls the generated circuit size, fixed-point width, or learning rule without runtime overhead.
\paragraph{Line 9}
\noindent\colorbox{gray!8}{\parbox{\dimexpr\linewidth-2\fboxsep\relax}{\raggedright\ttfamily\scriptsize     parameter int MAX\_FEATURE\_WORDS = MAX\_SAMPLES * FEATURE\_DIM;}}
\par\smallskip
This is a compile-time hardware setting. Vivado resolves it before synthesis, so it controls the generated circuit size, fixed-point width, or learning rule without runtime overhead.
\paragraph{Line 10}
\noindent\colorbox{gray!8}{\parbox{\dimexpr\linewidth-2\fboxsep\relax}{\raggedright\ttfamily\scriptsize \emph{blank line}}}
\par\smallskip
This blank line is not hardware or software behavior. It separates nearby declarations or operations so the file remains readable and reviewable.
\paragraph{Line 11}
\noindent\colorbox{gray!8}{\parbox{\dimexpr\linewidth-2\fboxsep\relax}{\raggedright\ttfamily\scriptsize     logic clk = 1'b0;}}
\par\smallskip
This is a local compile-time or procedural assignment. In context it either defines a helper value or connects a temporary expression used by nearby hardware logic.
\paragraph{Line 12}
\noindent\colorbox{gray!8}{\parbox{\dimexpr\linewidth-2\fboxsep\relax}{\raggedright\ttfamily\scriptsize     logic rst\_n = 1'b0;}}
\par\smallskip
This is a local compile-time or procedural assignment. In context it either defines a helper value or connects a temporary expression used by nearby hardware logic.
\paragraph{Line 13}
\noindent\colorbox{gray!8}{\parbox{\dimexpr\linewidth-2\fboxsep\relax}{\raggedright\ttfamily\scriptsize     logic start;}}
\par\smallskip
This SystemVerilog line contributes to the structural or sequential description of the accelerator. Its effect is resolved by XSim for verification and by Vivado during synthesis.
\paragraph{Line 14}
\noindent\colorbox{gray!8}{\parbox{\dimexpr\linewidth-2\fboxsep\relax}{\raggedright\ttfamily\scriptsize     logic feature\_valid;}}
\par\smallskip
This SystemVerilog line contributes to the structural or sequential description of the accelerator. Its effect is resolved by XSim for verification and by Vivado during synthesis.
\paragraph{Line 15}
\noindent\colorbox{gray!8}{\parbox{\dimexpr\linewidth-2\fboxsep\relax}{\raggedright\ttfamily\scriptsize     logic feature\_ready;}}
\par\smallskip
This SystemVerilog line contributes to the structural or sequential description of the accelerator. Its effect is resolved by XSim for verification and by Vivado during synthesis.
\paragraph{Line 16}
\noindent\colorbox{gray!8}{\parbox{\dimexpr\linewidth-2\fboxsep\relax}{\raggedright\ttfamily\scriptsize     logic signed [`CIFAR10\_SF\_HEAD\_FEATURE\_BITS-1:0] feature\_value;}}
\par\smallskip
This SystemVerilog line contributes to the structural or sequential description of the accelerator. Its effect is resolved by XSim for verification and by Vivado during synthesis.
\paragraph{Line 17}
\noindent\colorbox{gray!8}{\parbox{\dimexpr\linewidth-2\fboxsep\relax}{\raggedright\ttfamily\scriptsize     logic frame\_done;}}
\par\smallskip
This SystemVerilog line contributes to the structural or sequential description of the accelerator. Its effect is resolved by XSim for verification and by Vivado during synthesis.
\paragraph{Line 18}
\noindent\colorbox{gray!8}{\parbox{\dimexpr\linewidth-2\fboxsep\relax}{\raggedright\ttfamily\scriptsize     logic train\_enable;}}
\par\smallskip
This SystemVerilog line contributes to the structural or sequential description of the accelerator. Its effect is resolved by XSim for verification and by Vivado during synthesis.
\paragraph{Line 19}
\noindent\colorbox{gray!8}{\parbox{\dimexpr\linewidth-2\fboxsep\relax}{\raggedright\ttfamily\scriptsize     logic [3:0] train\_label;}}
\par\smallskip
This SystemVerilog line contributes to the structural or sequential description of the accelerator. Its effect is resolved by XSim for verification and by Vivado during synthesis.
\paragraph{Line 20}
\noindent\colorbox{gray!8}{\parbox{\dimexpr\linewidth-2\fboxsep\relax}{\raggedright\ttfamily\scriptsize     logic busy;}}
\par\smallskip
This SystemVerilog line contributes to the structural or sequential description of the accelerator. Its effect is resolved by XSim for verification and by Vivado during synthesis.
\paragraph{Line 21}
\noindent\colorbox{gray!8}{\parbox{\dimexpr\linewidth-2\fboxsep\relax}{\raggedright\ttfamily\scriptsize     logic result\_valid;}}
\par\smallskip
This SystemVerilog line contributes to the structural or sequential description of the accelerator. Its effect is resolved by XSim for verification and by Vivado during synthesis.
\paragraph{Line 22}
\noindent\colorbox{gray!8}{\parbox{\dimexpr\linewidth-2\fboxsep\relax}{\raggedright\ttfamily\scriptsize     logic [3:0] predicted\_label;}}
\par\smallskip
This SystemVerilog line contributes to the structural or sequential description of the accelerator. Its effect is resolved by XSim for verification and by Vivado during synthesis.
\paragraph{Line 23}
\noindent\colorbox{gray!8}{\parbox{\dimexpr\linewidth-2\fboxsep\relax}{\raggedright\ttfamily\scriptsize     logic signed [`CIFAR10\_SF\_HEAD\_ACC\_BITS-1:0] predicted\_score;}}
\par\smallskip
This SystemVerilog line contributes to the structural or sequential description of the accelerator. Its effect is resolved by XSim for verification and by Vivado during synthesis.
\paragraph{Line 24}
\noindent\colorbox{gray!8}{\parbox{\dimexpr\linewidth-2\fboxsep\relax}{\raggedright\ttfamily\scriptsize     logic train\_update\_valid;}}
\par\smallskip
This SystemVerilog line contributes to the structural or sequential description of the accelerator. Its effect is resolved by XSim for verification and by Vivado during synthesis.
\paragraph{Line 25}
\noindent\colorbox{gray!8}{\parbox{\dimexpr\linewidth-2\fboxsep\relax}{\raggedright\ttfamily\scriptsize     logic train\_update\_applied;}}
\par\smallskip
This SystemVerilog line contributes to the structural or sequential description of the accelerator. Its effect is resolved by XSim for verification and by Vivado during synthesis.
\paragraph{Line 26}
\noindent\colorbox{gray!8}{\parbox{\dimexpr\linewidth-2\fboxsep\relax}{\raggedright\ttfamily\scriptsize     logic [5:0] debug\_state;}}
\par\smallskip
This SystemVerilog line contributes to the structural or sequential description of the accelerator. Its effect is resolved by XSim for verification and by Vivado during synthesis.
\paragraph{Line 27}
\noindent\colorbox{gray!8}{\parbox{\dimexpr\linewidth-2\fboxsep\relax}{\raggedright\ttfamily\scriptsize     logic [15:0] debug\_feature;}}
\par\smallskip
This SystemVerilog line contributes to the structural or sequential description of the accelerator. Its effect is resolved by XSim for verification and by Vivado during synthesis.
\paragraph{Line 28}
\noindent\colorbox{gray!8}{\parbox{\dimexpr\linewidth-2\fboxsep\relax}{\raggedright\ttfamily\scriptsize     logic [7:0] debug\_class;}}
\par\smallskip
This SystemVerilog line contributes to the structural or sequential description of the accelerator. Its effect is resolved by XSim for verification and by Vivado during synthesis.
\paragraph{Line 29}
\noindent\colorbox{gray!8}{\parbox{\dimexpr\linewidth-2\fboxsep\relax}{\raggedright\ttfamily\scriptsize     logic [31:0] debug\_update\_count;}}
\par\smallskip
This SystemVerilog line contributes to the structural or sequential description of the accelerator. Its effect is resolved by XSim for verification and by Vivado during synthesis.
\paragraph{Line 30}
\noindent\colorbox{gray!8}{\parbox{\dimexpr\linewidth-2\fboxsep\relax}{\raggedright\ttfamily\scriptsize     logic [31:0] debug\_infer\_skip\_count;}}
\par\smallskip
This line exposes or updates a hardware activity counter. The counter makes the feature gate visible in waveforms and CSV traces, so the power-saving opportunity is observable from RTL rather than only from Python statistics.
\paragraph{Line 31}
\noindent\colorbox{gray!8}{\parbox{\dimexpr\linewidth-2\fboxsep\relax}{\raggedright\ttfamily\scriptsize     logic [31:0] debug\_train\_skip\_count;}}
\par\smallskip
This line exposes or updates a hardware activity counter. The counter makes the feature gate visible in waveforms and CSV traces, so the power-saving opportunity is observable from RTL rather than only from Python statistics.
\paragraph{Line 32}
\noindent\colorbox{gray!8}{\parbox{\dimexpr\linewidth-2\fboxsep\relax}{\raggedright\ttfamily\scriptsize \emph{blank line}}}
\par\smallskip
This blank line is not hardware or software behavior. It separates nearby declarations or operations so the file remains readable and reviewable.
\paragraph{Line 33}
\noindent\colorbox{gray!8}{\parbox{\dimexpr\linewidth-2\fboxsep\relax}{\raggedright\ttfamily\scriptsize     logic signed [`CIFAR10\_SF\_HEAD\_FEATURE\_BITS-1:0] feature\_mem [0:MAX\_FEATURE\_WORDS-1];}}
\par\smallskip
This line touches the feature memory that holds one 384-dimensional fixed-point Spikingformer feature vector for the current image. Keeping the feature vector local avoids repeatedly streaming the same values for every class score.
\paragraph{Line 34}
\noindent\colorbox{gray!8}{\parbox{\dimexpr\linewidth-2\fboxsep\relax}{\raggedright\ttfamily\scriptsize     logic [3:0] label\_mem [0:MAX\_SAMPLES-1];}}
\par\smallskip
This SystemVerilog line contributes to the structural or sequential description of the accelerator. Its effect is resolved by XSim for verification and by Vivado during synthesis.
\paragraph{Line 35}
\noindent\colorbox{gray!8}{\parbox{\dimexpr\linewidth-2\fboxsep\relax}{\raggedright\ttfamily\scriptsize \emph{blank line}}}
\par\smallskip
This blank line is not hardware or software behavior. It separates nearby declarations or operations so the file remains readable and reviewable.
\paragraph{Line 36}
\noindent\colorbox{gray!8}{\parbox{\dimexpr\linewidth-2\fboxsep\relax}{\raggedright\ttfamily\scriptsize     int feature\_ptr = 0;}}
\par\smallskip
This is a local compile-time or procedural assignment. In context it either defines a helper value or connects a temporary expression used by nearby hardware logic.
\paragraph{Line 37}
\noindent\colorbox{gray!8}{\parbox{\dimexpr\linewidth-2\fboxsep\relax}{\raggedright\ttfamily\scriptsize     int correct = 0;}}
\par\smallskip
This is a local compile-time or procedural assignment. In context it either defines a helper value or connects a temporary expression used by nearby hardware logic.
\paragraph{Line 38}
\noindent\colorbox{gray!8}{\parbox{\dimexpr\linewidth-2\fboxsep\relax}{\raggedright\ttfamily\scriptsize     int errors = 0;}}
\par\smallskip
This is a local compile-time or procedural assignment. In context it either defines a helper value or connects a temporary expression used by nearby hardware logic.
\paragraph{Line 39}
\noindent\colorbox{gray!8}{\parbox{\dimexpr\linewidth-2\fboxsep\relax}{\raggedright\ttfamily\scriptsize     int cycle = 0;}}
\par\smallskip
This is a local compile-time or procedural assignment. In context it either defines a helper value or connects a temporary expression used by nearby hardware logic.
\paragraph{Line 40}
\noindent\colorbox{gray!8}{\parbox{\dimexpr\linewidth-2\fboxsep\relax}{\raggedright\ttfamily\scriptsize     int trace\_fh;}}
\par\smallskip
This SystemVerilog line contributes to the structural or sequential description of the accelerator. Its effect is resolved by XSim for verification and by Vivado during synthesis.
\paragraph{Line 41}
\noindent\colorbox{gray!8}{\parbox{\dimexpr\linewidth-2\fboxsep\relax}{\raggedright\ttfamily\scriptsize     int result\_fh;}}
\par\smallskip
This SystemVerilog line contributes to the structural or sequential description of the accelerator. Its effect is resolved by XSim for verification and by Vivado during synthesis.
\paragraph{Line 42}
\noindent\colorbox{gray!8}{\parbox{\dimexpr\linewidth-2\fboxsep\relax}{\raggedright\ttfamily\scriptsize     bit trace\_enable = 1'b0;}}
\par\smallskip
This is a local compile-time or procedural assignment. In context it either defines a helper value or connects a temporary expression used by nearby hardware logic.
\paragraph{Line 43}
\noindent\colorbox{gray!8}{\parbox{\dimexpr\linewidth-2\fboxsep\relax}{\raggedright\ttfamily\scriptsize     bit vcd\_enabled = 1'b1;}}
\par\smallskip
This is a local compile-time or procedural assignment. In context it either defines a helper value or connects a temporary expression used by nearby hardware logic.
\paragraph{Line 44}
\noindent\colorbox{gray!8}{\parbox{\dimexpr\linewidth-2\fboxsep\relax}{\raggedright\ttfamily\scriptsize \emph{blank line}}}
\par\smallskip
This blank line is not hardware or software behavior. It separates nearby declarations or operations so the file remains readable and reviewable.
\paragraph{Line 45}
\noindent\colorbox{gray!8}{\parbox{\dimexpr\linewidth-2\fboxsep\relax}{\raggedright\ttfamily\scriptsize     cifar10\_spikingformer\_head\_accelerator dut (}}
\par\smallskip
This SystemVerilog line contributes to the structural or sequential description of the accelerator. Its effect is resolved by XSim for verification and by Vivado during synthesis.
\paragraph{Line 46}
\noindent\colorbox{gray!8}{\parbox{\dimexpr\linewidth-2\fboxsep\relax}{\raggedright\ttfamily\scriptsize         .clk(clk),}}
\par\smallskip
This SystemVerilog line contributes to the structural or sequential description of the accelerator. Its effect is resolved by XSim for verification and by Vivado during synthesis.
\paragraph{Line 47}
\noindent\colorbox{gray!8}{\parbox{\dimexpr\linewidth-2\fboxsep\relax}{\raggedright\ttfamily\scriptsize         .rst\_n(rst\_n),}}
\par\smallskip
This SystemVerilog line contributes to the structural or sequential description of the accelerator. Its effect is resolved by XSim for verification and by Vivado during synthesis.
\paragraph{Line 48}
\noindent\colorbox{gray!8}{\parbox{\dimexpr\linewidth-2\fboxsep\relax}{\raggedright\ttfamily\scriptsize         .start(start),}}
\par\smallskip
This SystemVerilog line contributes to the structural or sequential description of the accelerator. Its effect is resolved by XSim for verification and by Vivado during synthesis.
\paragraph{Line 49}
\noindent\colorbox{gray!8}{\parbox{\dimexpr\linewidth-2\fboxsep\relax}{\raggedright\ttfamily\scriptsize         .feature\_valid(feature\_valid),}}
\par\smallskip
This SystemVerilog line contributes to the structural or sequential description of the accelerator. Its effect is resolved by XSim for verification and by Vivado during synthesis.
\paragraph{Line 50}
\noindent\colorbox{gray!8}{\parbox{\dimexpr\linewidth-2\fboxsep\relax}{\raggedright\ttfamily\scriptsize         .feature\_ready(feature\_ready),}}
\par\smallskip
This SystemVerilog line contributes to the structural or sequential description of the accelerator. Its effect is resolved by XSim for verification and by Vivado during synthesis.
\paragraph{Line 51}
\noindent\colorbox{gray!8}{\parbox{\dimexpr\linewidth-2\fboxsep\relax}{\raggedright\ttfamily\scriptsize         .feature\_value(feature\_value),}}
\par\smallskip
This SystemVerilog line contributes to the structural or sequential description of the accelerator. Its effect is resolved by XSim for verification and by Vivado during synthesis.
\paragraph{Line 52}
\noindent\colorbox{gray!8}{\parbox{\dimexpr\linewidth-2\fboxsep\relax}{\raggedright\ttfamily\scriptsize         .frame\_done(frame\_done),}}
\par\smallskip
This SystemVerilog line contributes to the structural or sequential description of the accelerator. Its effect is resolved by XSim for verification and by Vivado during synthesis.
\paragraph{Line 53}
\noindent\colorbox{gray!8}{\parbox{\dimexpr\linewidth-2\fboxsep\relax}{\raggedright\ttfamily\scriptsize         .train\_enable(train\_enable),}}
\par\smallskip
This SystemVerilog line contributes to the structural or sequential description of the accelerator. Its effect is resolved by XSim for verification and by Vivado during synthesis.
\paragraph{Line 54}
\noindent\colorbox{gray!8}{\parbox{\dimexpr\linewidth-2\fboxsep\relax}{\raggedright\ttfamily\scriptsize         .train\_label(train\_label),}}
\par\smallskip
This SystemVerilog line contributes to the structural or sequential description of the accelerator. Its effect is resolved by XSim for verification and by Vivado during synthesis.
\paragraph{Line 55}
\noindent\colorbox{gray!8}{\parbox{\dimexpr\linewidth-2\fboxsep\relax}{\raggedright\ttfamily\scriptsize         .busy(busy),}}
\par\smallskip
This SystemVerilog line contributes to the structural or sequential description of the accelerator. Its effect is resolved by XSim for verification and by Vivado during synthesis.
\paragraph{Line 56}
\noindent\colorbox{gray!8}{\parbox{\dimexpr\linewidth-2\fboxsep\relax}{\raggedright\ttfamily\scriptsize         .result\_valid(result\_valid),}}
\par\smallskip
This SystemVerilog line contributes to the structural or sequential description of the accelerator. Its effect is resolved by XSim for verification and by Vivado during synthesis.
\paragraph{Line 57}
\noindent\colorbox{gray!8}{\parbox{\dimexpr\linewidth-2\fboxsep\relax}{\raggedright\ttfamily\scriptsize         .predicted\_label(predicted\_label),}}
\par\smallskip
This SystemVerilog line contributes to the structural or sequential description of the accelerator. Its effect is resolved by XSim for verification and by Vivado during synthesis.
\paragraph{Line 58}
\noindent\colorbox{gray!8}{\parbox{\dimexpr\linewidth-2\fboxsep\relax}{\raggedright\ttfamily\scriptsize         .predicted\_score(predicted\_score),}}
\par\smallskip
This SystemVerilog line contributes to the structural or sequential description of the accelerator. Its effect is resolved by XSim for verification and by Vivado during synthesis.
\paragraph{Line 59}
\noindent\colorbox{gray!8}{\parbox{\dimexpr\linewidth-2\fboxsep\relax}{\raggedright\ttfamily\scriptsize         .train\_update\_valid(train\_update\_valid),}}
\par\smallskip
This SystemVerilog line contributes to the structural or sequential description of the accelerator. Its effect is resolved by XSim for verification and by Vivado during synthesis.
\paragraph{Line 60}
\noindent\colorbox{gray!8}{\parbox{\dimexpr\linewidth-2\fboxsep\relax}{\raggedright\ttfamily\scriptsize         .train\_update\_applied(train\_update\_applied),}}
\par\smallskip
This SystemVerilog line contributes to the structural or sequential description of the accelerator. Its effect is resolved by XSim for verification and by Vivado during synthesis.
\paragraph{Line 61}
\noindent\colorbox{gray!8}{\parbox{\dimexpr\linewidth-2\fboxsep\relax}{\raggedright\ttfamily\scriptsize         .debug\_state(debug\_state),}}
\par\smallskip
This SystemVerilog line contributes to the structural or sequential description of the accelerator. Its effect is resolved by XSim for verification and by Vivado during synthesis.
\paragraph{Line 62}
\noindent\colorbox{gray!8}{\parbox{\dimexpr\linewidth-2\fboxsep\relax}{\raggedright\ttfamily\scriptsize         .debug\_feature(debug\_feature),}}
\par\smallskip
This SystemVerilog line contributes to the structural or sequential description of the accelerator. Its effect is resolved by XSim for verification and by Vivado during synthesis.
\paragraph{Line 63}
\noindent\colorbox{gray!8}{\parbox{\dimexpr\linewidth-2\fboxsep\relax}{\raggedright\ttfamily\scriptsize         .debug\_class(debug\_class),}}
\par\smallskip
This SystemVerilog line contributes to the structural or sequential description of the accelerator. Its effect is resolved by XSim for verification and by Vivado during synthesis.
\paragraph{Line 64}
\noindent\colorbox{gray!8}{\parbox{\dimexpr\linewidth-2\fboxsep\relax}{\raggedright\ttfamily\scriptsize         .debug\_update\_count(debug\_update\_count),}}
\par\smallskip
This SystemVerilog line contributes to the structural or sequential description of the accelerator. Its effect is resolved by XSim for verification and by Vivado during synthesis.
\paragraph{Line 65}
\noindent\colorbox{gray!8}{\parbox{\dimexpr\linewidth-2\fboxsep\relax}{\raggedright\ttfamily\scriptsize         .debug\_infer\_skip\_count(debug\_infer\_skip\_count),}}
\par\smallskip
This line exposes or updates a hardware activity counter. The counter makes the feature gate visible in waveforms and CSV traces, so the power-saving opportunity is observable from RTL rather than only from Python statistics.
\paragraph{Line 66}
\noindent\colorbox{gray!8}{\parbox{\dimexpr\linewidth-2\fboxsep\relax}{\raggedright\ttfamily\scriptsize         .debug\_train\_skip\_count(debug\_train\_skip\_count)}}
\par\smallskip
This line exposes or updates a hardware activity counter. The counter makes the feature gate visible in waveforms and CSV traces, so the power-saving opportunity is observable from RTL rather than only from Python statistics.
\paragraph{Line 67}
\noindent\colorbox{gray!8}{\parbox{\dimexpr\linewidth-2\fboxsep\relax}{\raggedright\ttfamily\scriptsize     );}}
\par\smallskip
This SystemVerilog line contributes to the structural or sequential description of the accelerator. Its effect is resolved by XSim for verification and by Vivado during synthesis.
\paragraph{Line 68}
\noindent\colorbox{gray!8}{\parbox{\dimexpr\linewidth-2\fboxsep\relax}{\raggedright\ttfamily\scriptsize \emph{blank line}}}
\par\smallskip
This blank line is not hardware or software behavior. It separates nearby declarations or operations so the file remains readable and reviewable.
\paragraph{Line 69}
\noindent\colorbox{gray!8}{\parbox{\dimexpr\linewidth-2\fboxsep\relax}{\raggedright\ttfamily\scriptsize     always \#5 clk = \textasciitilde{}clk;}}
\par\smallskip
This is a local compile-time or procedural assignment. In context it either defines a helper value or connects a temporary expression used by nearby hardware logic.
\paragraph{Line 70}
\noindent\colorbox{gray!8}{\parbox{\dimexpr\linewidth-2\fboxsep\relax}{\raggedright\ttfamily\scriptsize \emph{blank line}}}
\par\smallskip
This blank line is not hardware or software behavior. It separates nearby declarations or operations so the file remains readable and reviewable.
\paragraph{Line 71}
\noindent\colorbox{gray!8}{\parbox{\dimexpr\linewidth-2\fboxsep\relax}{\raggedright\ttfamily\scriptsize     task automatic require\_file(input string path);}}
\par\smallskip
This line declares a module port. A port is a physical or wrapper-visible signal used to move clocks, control pulses, features, labels, predictions, or debug information across module boundaries.
\paragraph{Line 72}
\noindent\colorbox{gray!8}{\parbox{\dimexpr\linewidth-2\fboxsep\relax}{\raggedright\ttfamily\scriptsize         int fh;}}
\par\smallskip
This SystemVerilog line contributes to the structural or sequential description of the accelerator. Its effect is resolved by XSim for verification and by Vivado during synthesis.
\paragraph{Line 73}
\noindent\colorbox{gray!8}{\parbox{\dimexpr\linewidth-2\fboxsep\relax}{\raggedright\ttfamily\scriptsize         begin}}
\par\smallskip
This SystemVerilog line contributes to the structural or sequential description of the accelerator. Its effect is resolved by XSim for verification and by Vivado during synthesis.
\paragraph{Line 74}
\noindent\colorbox{gray!8}{\parbox{\dimexpr\linewidth-2\fboxsep\relax}{\raggedright\ttfamily\scriptsize             fh = \$fopen(path, "r");}}
\par\smallskip
This is a local compile-time or procedural assignment. In context it either defines a helper value or connects a temporary expression used by nearby hardware logic.
\paragraph{Line 75}
\noindent\colorbox{gray!8}{\parbox{\dimexpr\linewidth-2\fboxsep\relax}{\raggedright\ttfamily\scriptsize             if (fh == 0) begin}}
\par\smallskip
This is a local compile-time or procedural assignment. In context it either defines a helper value or connects a temporary expression used by nearby hardware logic.
\paragraph{Line 76}
\noindent\colorbox{gray!8}{\parbox{\dimexpr\linewidth-2\fboxsep\relax}{\raggedright\ttfamily\scriptsize                 \$display("ERROR: missing required file \%s", path);}}
\par\smallskip
This SystemVerilog line contributes to the structural or sequential description of the accelerator. Its effect is resolved by XSim for verification and by Vivado during synthesis.
\paragraph{Line 77}
\noindent\colorbox{gray!8}{\parbox{\dimexpr\linewidth-2\fboxsep\relax}{\raggedright\ttfamily\scriptsize                 \$fatal;}}
\par\smallskip
This SystemVerilog line contributes to the structural or sequential description of the accelerator. Its effect is resolved by XSim for verification and by Vivado during synthesis.
\paragraph{Line 78}
\noindent\colorbox{gray!8}{\parbox{\dimexpr\linewidth-2\fboxsep\relax}{\raggedright\ttfamily\scriptsize             end}}
\par\smallskip
This SystemVerilog line contributes to the structural or sequential description of the accelerator. Its effect is resolved by XSim for verification and by Vivado during synthesis.
\paragraph{Line 79}
\noindent\colorbox{gray!8}{\parbox{\dimexpr\linewidth-2\fboxsep\relax}{\raggedright\ttfamily\scriptsize             \$fclose(fh);}}
\par\smallskip
This SystemVerilog line contributes to the structural or sequential description of the accelerator. Its effect is resolved by XSim for verification and by Vivado during synthesis.
\paragraph{Line 80}
\noindent\colorbox{gray!8}{\parbox{\dimexpr\linewidth-2\fboxsep\relax}{\raggedright\ttfamily\scriptsize         end}}
\par\smallskip
This SystemVerilog line contributes to the structural or sequential description of the accelerator. Its effect is resolved by XSim for verification and by Vivado during synthesis.
\paragraph{Line 81}
\noindent\colorbox{gray!8}{\parbox{\dimexpr\linewidth-2\fboxsep\relax}{\raggedright\ttfamily\scriptsize     endtask}}
\par\smallskip
This SystemVerilog line contributes to the structural or sequential description of the accelerator. Its effect is resolved by XSim for verification and by Vivado during synthesis.
\paragraph{Line 82}
\noindent\colorbox{gray!8}{\parbox{\dimexpr\linewidth-2\fboxsep\relax}{\raggedright\ttfamily\scriptsize \emph{blank line}}}
\par\smallskip
This blank line is not hardware or software behavior. It separates nearby declarations or operations so the file remains readable and reviewable.
\paragraph{Line 83}
\noindent\colorbox{gray!8}{\parbox{\dimexpr\linewidth-2\fboxsep\relax}{\raggedright\ttfamily\scriptsize     task automatic send\_sample(input int sample\_idx);}}
\par\smallskip
This line declares a module port. A port is a physical or wrapper-visible signal used to move clocks, control pulses, features, labels, predictions, or debug information across module boundaries.
\paragraph{Line 84}
\noindent\colorbox{gray!8}{\parbox{\dimexpr\linewidth-2\fboxsep\relax}{\raggedright\ttfamily\scriptsize         int waited;}}
\par\smallskip
This SystemVerilog line contributes to the structural or sequential description of the accelerator. Its effect is resolved by XSim for verification and by Vivado during synthesis.
\paragraph{Line 85}
\noindent\colorbox{gray!8}{\parbox{\dimexpr\linewidth-2\fboxsep\relax}{\raggedright\ttfamily\scriptsize         begin}}
\par\smallskip
This SystemVerilog line contributes to the structural or sequential description of the accelerator. Its effect is resolved by XSim for verification and by Vivado during synthesis.
\paragraph{Line 86}
\noindent\colorbox{gray!8}{\parbox{\dimexpr\linewidth-2\fboxsep\relax}{\raggedright\ttfamily\scriptsize             trace\_enable = (sample\_idx < TRACE\_SAMPLES);}}
\par\smallskip
This is a local compile-time or procedural assignment. In context it either defines a helper value or connects a temporary expression used by nearby hardware logic.
\paragraph{Line 87}
\noindent\colorbox{gray!8}{\parbox{\dimexpr\linewidth-2\fboxsep\relax}{\raggedright\ttfamily\scriptsize             if (sample\_idx == TRACE\_SAMPLES \&\& vcd\_enabled) begin}}
\par\smallskip
This is a local compile-time or procedural assignment. In context it either defines a helper value or connects a temporary expression used by nearby hardware logic.
\paragraph{Line 88}
\noindent\colorbox{gray!8}{\parbox{\dimexpr\linewidth-2\fboxsep\relax}{\raggedright\ttfamily\scriptsize                 \$dumpoff;}}
\par\smallskip
This SystemVerilog line contributes to the structural or sequential description of the accelerator. Its effect is resolved by XSim for verification and by Vivado during synthesis.
\paragraph{Line 89}
\noindent\colorbox{gray!8}{\parbox{\dimexpr\linewidth-2\fboxsep\relax}{\raggedright\ttfamily\scriptsize                 vcd\_enabled = 1'b0;}}
\par\smallskip
This is a local compile-time or procedural assignment. In context it either defines a helper value or connects a temporary expression used by nearby hardware logic.
\paragraph{Line 90}
\noindent\colorbox{gray!8}{\parbox{\dimexpr\linewidth-2\fboxsep\relax}{\raggedright\ttfamily\scriptsize             end}}
\par\smallskip
This SystemVerilog line contributes to the structural or sequential description of the accelerator. Its effect is resolved by XSim for verification and by Vivado during synthesis.
\paragraph{Line 91}
\noindent\colorbox{gray!8}{\parbox{\dimexpr\linewidth-2\fboxsep\relax}{\raggedright\ttfamily\scriptsize \emph{blank line}}}
\par\smallskip
This blank line is not hardware or software behavior. It separates nearby declarations or operations so the file remains readable and reviewable.
\paragraph{Line 92}
\noindent\colorbox{gray!8}{\parbox{\dimexpr\linewidth-2\fboxsep\relax}{\raggedright\ttfamily\scriptsize             @(negedge clk);}}
\par\smallskip
This SystemVerilog line contributes to the structural or sequential description of the accelerator. Its effect is resolved by XSim for verification and by Vivado during synthesis.
\paragraph{Line 93}
\noindent\colorbox{gray!8}{\parbox{\dimexpr\linewidth-2\fboxsep\relax}{\raggedright\ttfamily\scriptsize             start = 1'b1;}}
\par\smallskip
This is a local compile-time or procedural assignment. In context it either defines a helper value or connects a temporary expression used by nearby hardware logic.
\paragraph{Line 94}
\noindent\colorbox{gray!8}{\parbox{\dimexpr\linewidth-2\fboxsep\relax}{\raggedright\ttfamily\scriptsize             @(negedge clk);}}
\par\smallskip
This SystemVerilog line contributes to the structural or sequential description of the accelerator. Its effect is resolved by XSim for verification and by Vivado during synthesis.
\paragraph{Line 95}
\noindent\colorbox{gray!8}{\parbox{\dimexpr\linewidth-2\fboxsep\relax}{\raggedright\ttfamily\scriptsize             start = 1'b0;}}
\par\smallskip
This is a local compile-time or procedural assignment. In context it either defines a helper value or connects a temporary expression used by nearby hardware logic.
\paragraph{Line 96}
\noindent\colorbox{gray!8}{\parbox{\dimexpr\linewidth-2\fboxsep\relax}{\raggedright\ttfamily\scriptsize \emph{blank line}}}
\par\smallskip
This blank line is not hardware or software behavior. It separates nearby declarations or operations so the file remains readable and reviewable.
\paragraph{Line 97}
\noindent\colorbox{gray!8}{\parbox{\dimexpr\linewidth-2\fboxsep\relax}{\raggedright\ttfamily\scriptsize             waited = 0;}}
\par\smallskip
This is a local compile-time or procedural assignment. In context it either defines a helper value or connects a temporary expression used by nearby hardware logic.
\paragraph{Line 98}
\noindent\colorbox{gray!8}{\parbox{\dimexpr\linewidth-2\fboxsep\relax}{\raggedright\ttfamily\scriptsize             while (!feature\_ready \&\& waited < 10000) begin}}
\par\smallskip
This SystemVerilog line contributes to the structural or sequential description of the accelerator. Its effect is resolved by XSim for verification and by Vivado during synthesis.
\paragraph{Line 99}
\noindent\colorbox{gray!8}{\parbox{\dimexpr\linewidth-2\fboxsep\relax}{\raggedright\ttfamily\scriptsize                 @(posedge clk);}}
\par\smallskip
This SystemVerilog line contributes to the structural or sequential description of the accelerator. Its effect is resolved by XSim for verification and by Vivado during synthesis.
\paragraph{Line 100}
\noindent\colorbox{gray!8}{\parbox{\dimexpr\linewidth-2\fboxsep\relax}{\raggedright\ttfamily\scriptsize                 waited++;}}
\par\smallskip
This SystemVerilog line contributes to the structural or sequential description of the accelerator. Its effect is resolved by XSim for verification and by Vivado during synthesis.
\paragraph{Line 101}
\noindent\colorbox{gray!8}{\parbox{\dimexpr\linewidth-2\fboxsep\relax}{\raggedright\ttfamily\scriptsize             end}}
\par\smallskip
This SystemVerilog line contributes to the structural or sequential description of the accelerator. Its effect is resolved by XSim for verification and by Vivado during synthesis.
\paragraph{Line 102}
\noindent\colorbox{gray!8}{\parbox{\dimexpr\linewidth-2\fboxsep\relax}{\raggedright\ttfamily\scriptsize             if (!feature\_ready) begin}}
\par\smallskip
This SystemVerilog line contributes to the structural or sequential description of the accelerator. Its effect is resolved by XSim for verification and by Vivado during synthesis.
\paragraph{Line 103}
\noindent\colorbox{gray!8}{\parbox{\dimexpr\linewidth-2\fboxsep\relax}{\raggedright\ttfamily\scriptsize                 \$display("ERROR: sample \%0d never reached feature\_ready", sample\_idx);}}
\par\smallskip
This SystemVerilog line contributes to the structural or sequential description of the accelerator. Its effect is resolved by XSim for verification and by Vivado during synthesis.
\paragraph{Line 104}
\noindent\colorbox{gray!8}{\parbox{\dimexpr\linewidth-2\fboxsep\relax}{\raggedright\ttfamily\scriptsize                 errors++;}}
\par\smallskip
This SystemVerilog line contributes to the structural or sequential description of the accelerator. Its effect is resolved by XSim for verification and by Vivado during synthesis.
\paragraph{Line 105}
\noindent\colorbox{gray!8}{\parbox{\dimexpr\linewidth-2\fboxsep\relax}{\raggedright\ttfamily\scriptsize             end}}
\par\smallskip
This SystemVerilog line contributes to the structural or sequential description of the accelerator. Its effect is resolved by XSim for verification and by Vivado during synthesis.
\paragraph{Line 106}
\noindent\colorbox{gray!8}{\parbox{\dimexpr\linewidth-2\fboxsep\relax}{\raggedright\ttfamily\scriptsize \emph{blank line}}}
\par\smallskip
This blank line is not hardware or software behavior. It separates nearby declarations or operations so the file remains readable and reviewable.
\paragraph{Line 107}
\noindent\colorbox{gray!8}{\parbox{\dimexpr\linewidth-2\fboxsep\relax}{\raggedright\ttfamily\scriptsize             for (int f = 0; f < FEATURE\_DIM; f++) begin}}
\par\smallskip
This is a local compile-time or procedural assignment. In context it either defines a helper value or connects a temporary expression used by nearby hardware logic.
\paragraph{Line 108}
\noindent\colorbox{gray!8}{\parbox{\dimexpr\linewidth-2\fboxsep\relax}{\raggedright\ttfamily\scriptsize                 @(negedge clk);}}
\par\smallskip
This SystemVerilog line contributes to the structural or sequential description of the accelerator. Its effect is resolved by XSim for verification and by Vivado during synthesis.
\paragraph{Line 109}
\noindent\colorbox{gray!8}{\parbox{\dimexpr\linewidth-2\fboxsep\relax}{\raggedright\ttfamily\scriptsize                 feature\_value = feature\_mem[feature\_ptr];}}
\par\smallskip
This line touches the feature memory that holds one 384-dimensional fixed-point Spikingformer feature vector for the current image. Keeping the feature vector local avoids repeatedly streaming the same values for every class score.
\paragraph{Line 110}
\noindent\colorbox{gray!8}{\parbox{\dimexpr\linewidth-2\fboxsep\relax}{\raggedright\ttfamily\scriptsize                 feature\_ptr++;}}
\par\smallskip
This SystemVerilog line contributes to the structural or sequential description of the accelerator. Its effect is resolved by XSim for verification and by Vivado during synthesis.
\paragraph{Line 111}
\noindent\colorbox{gray!8}{\parbox{\dimexpr\linewidth-2\fboxsep\relax}{\raggedright\ttfamily\scriptsize                 feature\_valid = 1'b1;}}
\par\smallskip
This is a local compile-time or procedural assignment. In context it either defines a helper value or connects a temporary expression used by nearby hardware logic.
\paragraph{Line 112}
\noindent\colorbox{gray!8}{\parbox{\dimexpr\linewidth-2\fboxsep\relax}{\raggedright\ttfamily\scriptsize                 @(posedge clk);}}
\par\smallskip
This SystemVerilog line contributes to the structural or sequential description of the accelerator. Its effect is resolved by XSim for verification and by Vivado during synthesis.
\paragraph{Line 113}
\noindent\colorbox{gray!8}{\parbox{\dimexpr\linewidth-2\fboxsep\relax}{\raggedright\ttfamily\scriptsize                 \#1;}}
\par\smallskip
This is a human-facing comment. It documents intent for the reader and is ignored by the compiler or interpreter.
\paragraph{Line 114}
\noindent\colorbox{gray!8}{\parbox{\dimexpr\linewidth-2\fboxsep\relax}{\raggedright\ttfamily\scriptsize                 if (f < FEATURE\_DIM - 1 \&\& !feature\_ready) begin}}
\par\smallskip
This SystemVerilog line contributes to the structural or sequential description of the accelerator. Its effect is resolved by XSim for verification and by Vivado during synthesis.
\paragraph{Line 115}
\noindent\colorbox{gray!8}{\parbox{\dimexpr\linewidth-2\fboxsep\relax}{\raggedright\ttfamily\scriptsize                     \$display("ERROR: feature\_ready dropped during sample \%0d feature \%0d", sample\_idx, f);}}
\par\smallskip
This SystemVerilog line contributes to the structural or sequential description of the accelerator. Its effect is resolved by XSim for verification and by Vivado during synthesis.
\paragraph{Line 116}
\noindent\colorbox{gray!8}{\parbox{\dimexpr\linewidth-2\fboxsep\relax}{\raggedright\ttfamily\scriptsize                     errors++;}}
\par\smallskip
This SystemVerilog line contributes to the structural or sequential description of the accelerator. Its effect is resolved by XSim for verification and by Vivado during synthesis.
\paragraph{Line 117}
\noindent\colorbox{gray!8}{\parbox{\dimexpr\linewidth-2\fboxsep\relax}{\raggedright\ttfamily\scriptsize                 end}}
\par\smallskip
This SystemVerilog line contributes to the structural or sequential description of the accelerator. Its effect is resolved by XSim for verification and by Vivado during synthesis.
\paragraph{Line 118}
\noindent\colorbox{gray!8}{\parbox{\dimexpr\linewidth-2\fboxsep\relax}{\raggedright\ttfamily\scriptsize             end}}
\par\smallskip
This SystemVerilog line contributes to the structural or sequential description of the accelerator. Its effect is resolved by XSim for verification and by Vivado during synthesis.
\paragraph{Line 119}
\noindent\colorbox{gray!8}{\parbox{\dimexpr\linewidth-2\fboxsep\relax}{\raggedright\ttfamily\scriptsize \emph{blank line}}}
\par\smallskip
This blank line is not hardware or software behavior. It separates nearby declarations or operations so the file remains readable and reviewable.
\paragraph{Line 120}
\noindent\colorbox{gray!8}{\parbox{\dimexpr\linewidth-2\fboxsep\relax}{\raggedright\ttfamily\scriptsize             @(negedge clk);}}
\par\smallskip
This SystemVerilog line contributes to the structural or sequential description of the accelerator. Its effect is resolved by XSim for verification and by Vivado during synthesis.
\paragraph{Line 121}
\noindent\colorbox{gray!8}{\parbox{\dimexpr\linewidth-2\fboxsep\relax}{\raggedright\ttfamily\scriptsize             feature\_valid = 1'b0;}}
\par\smallskip
This is a local compile-time or procedural assignment. In context it either defines a helper value or connects a temporary expression used by nearby hardware logic.
\paragraph{Line 122}
\noindent\colorbox{gray!8}{\parbox{\dimexpr\linewidth-2\fboxsep\relax}{\raggedright\ttfamily\scriptsize             frame\_done = 1'b1;}}
\par\smallskip
This is a local compile-time or procedural assignment. In context it either defines a helper value or connects a temporary expression used by nearby hardware logic.
\paragraph{Line 123}
\noindent\colorbox{gray!8}{\parbox{\dimexpr\linewidth-2\fboxsep\relax}{\raggedright\ttfamily\scriptsize             @(negedge clk);}}
\par\smallskip
This SystemVerilog line contributes to the structural or sequential description of the accelerator. Its effect is resolved by XSim for verification and by Vivado during synthesis.
\paragraph{Line 124}
\noindent\colorbox{gray!8}{\parbox{\dimexpr\linewidth-2\fboxsep\relax}{\raggedright\ttfamily\scriptsize             frame\_done = 1'b0;}}
\par\smallskip
This is a local compile-time or procedural assignment. In context it either defines a helper value or connects a temporary expression used by nearby hardware logic.
\paragraph{Line 125}
\noindent\colorbox{gray!8}{\parbox{\dimexpr\linewidth-2\fboxsep\relax}{\raggedright\ttfamily\scriptsize \emph{blank line}}}
\par\smallskip
This blank line is not hardware or software behavior. It separates nearby declarations or operations so the file remains readable and reviewable.
\paragraph{Line 126}
\noindent\colorbox{gray!8}{\parbox{\dimexpr\linewidth-2\fboxsep\relax}{\raggedright\ttfamily\scriptsize             waited = 0;}}
\par\smallskip
This is a local compile-time or procedural assignment. In context it either defines a helper value or connects a temporary expression used by nearby hardware logic.
\paragraph{Line 127}
\noindent\colorbox{gray!8}{\parbox{\dimexpr\linewidth-2\fboxsep\relax}{\raggedright\ttfamily\scriptsize             while (!result\_valid \&\& waited < 100000) begin}}
\par\smallskip
This SystemVerilog line contributes to the structural or sequential description of the accelerator. Its effect is resolved by XSim for verification and by Vivado during synthesis.
\paragraph{Line 128}
\noindent\colorbox{gray!8}{\parbox{\dimexpr\linewidth-2\fboxsep\relax}{\raggedright\ttfamily\scriptsize                 @(posedge clk);}}
\par\smallskip
This SystemVerilog line contributes to the structural or sequential description of the accelerator. Its effect is resolved by XSim for verification and by Vivado during synthesis.
\paragraph{Line 129}
\noindent\colorbox{gray!8}{\parbox{\dimexpr\linewidth-2\fboxsep\relax}{\raggedright\ttfamily\scriptsize                 waited++;}}
\par\smallskip
This SystemVerilog line contributes to the structural or sequential description of the accelerator. Its effect is resolved by XSim for verification and by Vivado during synthesis.
\paragraph{Line 130}
\noindent\colorbox{gray!8}{\parbox{\dimexpr\linewidth-2\fboxsep\relax}{\raggedright\ttfamily\scriptsize             end}}
\par\smallskip
This SystemVerilog line contributes to the structural or sequential description of the accelerator. Its effect is resolved by XSim for verification and by Vivado during synthesis.
\paragraph{Line 131}
\noindent\colorbox{gray!8}{\parbox{\dimexpr\linewidth-2\fboxsep\relax}{\raggedright\ttfamily\scriptsize             if (!result\_valid) begin}}
\par\smallskip
This SystemVerilog line contributes to the structural or sequential description of the accelerator. Its effect is resolved by XSim for verification and by Vivado during synthesis.
\paragraph{Line 132}
\noindent\colorbox{gray!8}{\parbox{\dimexpr\linewidth-2\fboxsep\relax}{\raggedright\ttfamily\scriptsize                 \$display("ERROR: sample \%0d no result\_valid", sample\_idx);}}
\par\smallskip
This SystemVerilog line contributes to the structural or sequential description of the accelerator. Its effect is resolved by XSim for verification and by Vivado during synthesis.
\paragraph{Line 133}
\noindent\colorbox{gray!8}{\parbox{\dimexpr\linewidth-2\fboxsep\relax}{\raggedright\ttfamily\scriptsize                 errors++;}}
\par\smallskip
This SystemVerilog line contributes to the structural or sequential description of the accelerator. Its effect is resolved by XSim for verification and by Vivado during synthesis.
\paragraph{Line 134}
\noindent\colorbox{gray!8}{\parbox{\dimexpr\linewidth-2\fboxsep\relax}{\raggedright\ttfamily\scriptsize             end else begin}}
\par\smallskip
This SystemVerilog line contributes to the structural or sequential description of the accelerator. Its effect is resolved by XSim for verification and by Vivado during synthesis.
\paragraph{Line 135}
\noindent\colorbox{gray!8}{\parbox{\dimexpr\linewidth-2\fboxsep\relax}{\raggedright\ttfamily\scriptsize                 if (predicted\_label == label\_mem[sample\_idx]) begin}}
\par\smallskip
This is a local compile-time or procedural assignment. In context it either defines a helper value or connects a temporary expression used by nearby hardware logic.
\paragraph{Line 136}
\noindent\colorbox{gray!8}{\parbox{\dimexpr\linewidth-2\fboxsep\relax}{\raggedright\ttfamily\scriptsize                     correct++;}}
\par\smallskip
This SystemVerilog line contributes to the structural or sequential description of the accelerator. Its effect is resolved by XSim for verification and by Vivado during synthesis.
\paragraph{Line 137}
\noindent\colorbox{gray!8}{\parbox{\dimexpr\linewidth-2\fboxsep\relax}{\raggedright\ttfamily\scriptsize                 end}}
\par\smallskip
This SystemVerilog line contributes to the structural or sequential description of the accelerator. Its effect is resolved by XSim for verification and by Vivado during synthesis.
\paragraph{Line 138}
\noindent\colorbox{gray!8}{\parbox{\dimexpr\linewidth-2\fboxsep\relax}{\raggedright\ttfamily\scriptsize                 \$fwrite(result\_fh, "\%0d,\%0d,\%0d,\%0d,\%0d,\%0d,\%0d",}}
\par\smallskip
This SystemVerilog line contributes to the structural or sequential description of the accelerator. Its effect is resolved by XSim for verification and by Vivado during synthesis.
\paragraph{Line 139}
\noindent\colorbox{gray!8}{\parbox{\dimexpr\linewidth-2\fboxsep\relax}{\raggedright\ttfamily\scriptsize                         sample\_idx, label\_mem[sample\_idx], predicted\_label, predicted\_score,}}
\par\smallskip
This SystemVerilog line contributes to the structural or sequential description of the accelerator. Its effect is resolved by XSim for verification and by Vivado during synthesis.
\paragraph{Line 140}
\noindent\colorbox{gray!8}{\parbox{\dimexpr\linewidth-2\fboxsep\relax}{\raggedright\ttfamily\scriptsize                         predicted\_label == label\_mem[sample\_idx], debug\_infer\_skip\_count, debug\_train\_skip\_count);}}
\par\smallskip
This line exposes or updates a hardware activity counter. The counter makes the feature gate visible in waveforms and CSV traces, so the power-saving opportunity is observable from RTL rather than only from Python statistics.
\paragraph{Line 141}
\noindent\colorbox{gray!8}{\parbox{\dimexpr\linewidth-2\fboxsep\relax}{\raggedright\ttfamily\scriptsize                 for (int c = 0; c < 10; c++) begin}}
\par\smallskip
This is a local compile-time or procedural assignment. In context it either defines a helper value or connects a temporary expression used by nearby hardware logic.
\paragraph{Line 142}
\noindent\colorbox{gray!8}{\parbox{\dimexpr\linewidth-2\fboxsep\relax}{\raggedright\ttfamily\scriptsize                     \$fwrite(result\_fh, ",\%0d", \$signed(dut.class\_score\_mem[c]));}}
\par\smallskip
This stores or reads the ten class scores for the current sample. Those scores are needed for argmax prediction, CSV testbench reporting, and the margin test used before an online update.
\paragraph{Line 143}
\noindent\colorbox{gray!8}{\parbox{\dimexpr\linewidth-2\fboxsep\relax}{\raggedright\ttfamily\scriptsize                 end}}
\par\smallskip
This SystemVerilog line contributes to the structural or sequential description of the accelerator. Its effect is resolved by XSim for verification and by Vivado during synthesis.
\paragraph{Line 144}
\noindent\colorbox{gray!8}{\parbox{\dimexpr\linewidth-2\fboxsep\relax}{\raggedright\ttfamily\scriptsize                 \$fwrite(result\_fh, "\textbackslash{}n");}}
\par\smallskip
This SystemVerilog line contributes to the structural or sequential description of the accelerator. Its effect is resolved by XSim for verification and by Vivado during synthesis.
\paragraph{Line 145}
\noindent\colorbox{gray!8}{\parbox{\dimexpr\linewidth-2\fboxsep\relax}{\raggedright\ttfamily\scriptsize                 if (sample\_idx < TRACE\_SAMPLES || (sample\_idx \% 1000) == 0 || sample\_idx == NUM\_SAMPLES - 1) begin}}
\par\smallskip
This is a local compile-time or procedural assignment. In context it either defines a helper value or connects a temporary expression used by nearby hardware logic.
\paragraph{Line 146}
\noindent\colorbox{gray!8}{\parbox{\dimexpr\linewidth-2\fboxsep\relax}{\raggedright\ttfamily\scriptsize                     \$display("sample=\%0d label=\%0d pred=\%0d correct=\%0d",}}
\par\smallskip
This is a local compile-time or procedural assignment. In context it either defines a helper value or connects a temporary expression used by nearby hardware logic.
\paragraph{Line 147}
\noindent\colorbox{gray!8}{\parbox{\dimexpr\linewidth-2\fboxsep\relax}{\raggedright\ttfamily\scriptsize                              sample\_idx, label\_mem[sample\_idx], predicted\_label,}}
\par\smallskip
This SystemVerilog line contributes to the structural or sequential description of the accelerator. Its effect is resolved by XSim for verification and by Vivado during synthesis.
\paragraph{Line 148}
\noindent\colorbox{gray!8}{\parbox{\dimexpr\linewidth-2\fboxsep\relax}{\raggedright\ttfamily\scriptsize                              predicted\_label == label\_mem[sample\_idx]);}}
\par\smallskip
This is a local compile-time or procedural assignment. In context it either defines a helper value or connects a temporary expression used by nearby hardware logic.
\paragraph{Line 149}
\noindent\colorbox{gray!8}{\parbox{\dimexpr\linewidth-2\fboxsep\relax}{\raggedright\ttfamily\scriptsize                 end}}
\par\smallskip
This SystemVerilog line contributes to the structural or sequential description of the accelerator. Its effect is resolved by XSim for verification and by Vivado during synthesis.
\paragraph{Line 150}
\noindent\colorbox{gray!8}{\parbox{\dimexpr\linewidth-2\fboxsep\relax}{\raggedright\ttfamily\scriptsize             end}}
\par\smallskip
This SystemVerilog line contributes to the structural or sequential description of the accelerator. Its effect is resolved by XSim for verification and by Vivado during synthesis.
\paragraph{Line 151}
\noindent\colorbox{gray!8}{\parbox{\dimexpr\linewidth-2\fboxsep\relax}{\raggedright\ttfamily\scriptsize \emph{blank line}}}
\par\smallskip
This blank line is not hardware or software behavior. It separates nearby declarations or operations so the file remains readable and reviewable.
\paragraph{Line 152}
\noindent\colorbox{gray!8}{\parbox{\dimexpr\linewidth-2\fboxsep\relax}{\raggedright\ttfamily\scriptsize             waited = 0;}}
\par\smallskip
This is a local compile-time or procedural assignment. In context it either defines a helper value or connects a temporary expression used by nearby hardware logic.
\paragraph{Line 153}
\noindent\colorbox{gray!8}{\parbox{\dimexpr\linewidth-2\fboxsep\relax}{\raggedright\ttfamily\scriptsize             while (busy \&\& waited < 1000) begin}}
\par\smallskip
This SystemVerilog line contributes to the structural or sequential description of the accelerator. Its effect is resolved by XSim for verification and by Vivado during synthesis.
\paragraph{Line 154}
\noindent\colorbox{gray!8}{\parbox{\dimexpr\linewidth-2\fboxsep\relax}{\raggedright\ttfamily\scriptsize                 @(posedge clk);}}
\par\smallskip
This SystemVerilog line contributes to the structural or sequential description of the accelerator. Its effect is resolved by XSim for verification and by Vivado during synthesis.
\paragraph{Line 155}
\noindent\colorbox{gray!8}{\parbox{\dimexpr\linewidth-2\fboxsep\relax}{\raggedright\ttfamily\scriptsize                 waited++;}}
\par\smallskip
This SystemVerilog line contributes to the structural or sequential description of the accelerator. Its effect is resolved by XSim for verification and by Vivado during synthesis.
\paragraph{Line 156}
\noindent\colorbox{gray!8}{\parbox{\dimexpr\linewidth-2\fboxsep\relax}{\raggedright\ttfamily\scriptsize             end}}
\par\smallskip
This SystemVerilog line contributes to the structural or sequential description of the accelerator. Its effect is resolved by XSim for verification and by Vivado during synthesis.
\paragraph{Line 157}
\noindent\colorbox{gray!8}{\parbox{\dimexpr\linewidth-2\fboxsep\relax}{\raggedright\ttfamily\scriptsize             trace\_enable = 1'b0;}}
\par\smallskip
This is a local compile-time or procedural assignment. In context it either defines a helper value or connects a temporary expression used by nearby hardware logic.
\paragraph{Line 158}
\noindent\colorbox{gray!8}{\parbox{\dimexpr\linewidth-2\fboxsep\relax}{\raggedright\ttfamily\scriptsize         end}}
\par\smallskip
This SystemVerilog line contributes to the structural or sequential description of the accelerator. Its effect is resolved by XSim for verification and by Vivado during synthesis.
\paragraph{Line 159}
\noindent\colorbox{gray!8}{\parbox{\dimexpr\linewidth-2\fboxsep\relax}{\raggedright\ttfamily\scriptsize     endtask}}
\par\smallskip
This SystemVerilog line contributes to the structural or sequential description of the accelerator. Its effect is resolved by XSim for verification and by Vivado during synthesis.
\paragraph{Line 160}
\noindent\colorbox{gray!8}{\parbox{\dimexpr\linewidth-2\fboxsep\relax}{\raggedright\ttfamily\scriptsize \emph{blank line}}}
\par\smallskip
This blank line is not hardware or software behavior. It separates nearby declarations or operations so the file remains readable and reviewable.
\paragraph{Line 161}
\noindent\colorbox{gray!8}{\parbox{\dimexpr\linewidth-2\fboxsep\relax}{\raggedright\ttfamily\scriptsize     always @(posedge clk) begin}}
\par\smallskip
This SystemVerilog line contributes to the structural or sequential description of the accelerator. Its effect is resolved by XSim for verification and by Vivado during synthesis.
\paragraph{Line 162}
\noindent\colorbox{gray!8}{\parbox{\dimexpr\linewidth-2\fboxsep\relax}{\raggedright\ttfamily\scriptsize         if (trace\_fh != 0 \&\& trace\_enable) begin}}
\par\smallskip
This is a local compile-time or procedural assignment. In context it either defines a helper value or connects a temporary expression used by nearby hardware logic.
\paragraph{Line 163}
\noindent\colorbox{gray!8}{\parbox{\dimexpr\linewidth-2\fboxsep\relax}{\raggedright\ttfamily\scriptsize             \$fwrite(trace\_fh, "\%0t,\%0d,\%0d,\%0d,\%0d,\%0d,\%0d,\%0d,\%0d,\%0d,\%0d,\%0d,\%0d\textbackslash{}n",}}
\par\smallskip
This SystemVerilog line contributes to the structural or sequential description of the accelerator. Its effect is resolved by XSim for verification and by Vivado during synthesis.
\paragraph{Line 164}
\noindent\colorbox{gray!8}{\parbox{\dimexpr\linewidth-2\fboxsep\relax}{\raggedright\ttfamily\scriptsize                     \$time, cycle, rst\_n, debug\_state, start, feature\_valid, feature\_ready,}}
\par\smallskip
This SystemVerilog line contributes to the structural or sequential description of the accelerator. Its effect is resolved by XSim for verification and by Vivado during synthesis.
\paragraph{Line 165}
\noindent\colorbox{gray!8}{\parbox{\dimexpr\linewidth-2\fboxsep\relax}{\raggedright\ttfamily\scriptsize                     debug\_feature, debug\_class, result\_valid, predicted\_label,}}
\par\smallskip
This SystemVerilog line contributes to the structural or sequential description of the accelerator. Its effect is resolved by XSim for verification and by Vivado during synthesis.
\paragraph{Line 166}
\noindent\colorbox{gray!8}{\parbox{\dimexpr\linewidth-2\fboxsep\relax}{\raggedright\ttfamily\scriptsize                     debug\_infer\_skip\_count, debug\_train\_skip\_count);}}
\par\smallskip
This line exposes or updates a hardware activity counter. The counter makes the feature gate visible in waveforms and CSV traces, so the power-saving opportunity is observable from RTL rather than only from Python statistics.
\paragraph{Line 167}
\noindent\colorbox{gray!8}{\parbox{\dimexpr\linewidth-2\fboxsep\relax}{\raggedright\ttfamily\scriptsize         end}}
\par\smallskip
This SystemVerilog line contributes to the structural or sequential description of the accelerator. Its effect is resolved by XSim for verification and by Vivado during synthesis.
\paragraph{Line 168}
\noindent\colorbox{gray!8}{\parbox{\dimexpr\linewidth-2\fboxsep\relax}{\raggedright\ttfamily\scriptsize         cycle++;}}
\par\smallskip
This SystemVerilog line contributes to the structural or sequential description of the accelerator. Its effect is resolved by XSim for verification and by Vivado during synthesis.
\paragraph{Line 169}
\noindent\colorbox{gray!8}{\parbox{\dimexpr\linewidth-2\fboxsep\relax}{\raggedright\ttfamily\scriptsize     end}}
\par\smallskip
This SystemVerilog line contributes to the structural or sequential description of the accelerator. Its effect is resolved by XSim for verification and by Vivado during synthesis.
\paragraph{Line 170}
\noindent\colorbox{gray!8}{\parbox{\dimexpr\linewidth-2\fboxsep\relax}{\raggedright\ttfamily\scriptsize \emph{blank line}}}
\par\smallskip
This blank line is not hardware or software behavior. It separates nearby declarations or operations so the file remains readable and reviewable.
\paragraph{Line 171}
\noindent\colorbox{gray!8}{\parbox{\dimexpr\linewidth-2\fboxsep\relax}{\raggedright\ttfamily\scriptsize     initial begin}}
\par\smallskip
This simulation/synthesis initialization block loads trained memories from generated .mem files. On FPGA, Vivado maps these initial values into BRAM/LUTRAM initialization data.
\paragraph{Line 172}
\noindent\colorbox{gray!8}{\parbox{\dimexpr\linewidth-2\fboxsep\relax}{\raggedright\ttfamily\scriptsize         if (NUM\_SAMPLES > MAX\_SAMPLES) begin}}
\par\smallskip
This SystemVerilog line contributes to the structural or sequential description of the accelerator. Its effect is resolved by XSim for verification and by Vivado during synthesis.
\paragraph{Line 173}
\noindent\colorbox{gray!8}{\parbox{\dimexpr\linewidth-2\fboxsep\relax}{\raggedright\ttfamily\scriptsize             \$display("ERROR: NUM\_SAMPLES=\%0d exceeds MAX\_SAMPLES=\%0d", NUM\_SAMPLES, MAX\_SAMPLES);}}
\par\smallskip
This is a local compile-time or procedural assignment. In context it either defines a helper value or connects a temporary expression used by nearby hardware logic.
\paragraph{Line 174}
\noindent\colorbox{gray!8}{\parbox{\dimexpr\linewidth-2\fboxsep\relax}{\raggedright\ttfamily\scriptsize             \$fatal;}}
\par\smallskip
This SystemVerilog line contributes to the structural or sequential description of the accelerator. Its effect is resolved by XSim for verification and by Vivado during synthesis.
\paragraph{Line 175}
\noindent\colorbox{gray!8}{\parbox{\dimexpr\linewidth-2\fboxsep\relax}{\raggedright\ttfamily\scriptsize         end}}
\par\smallskip
This SystemVerilog line contributes to the structural or sequential description of the accelerator. Its effect is resolved by XSim for verification and by Vivado during synthesis.
\paragraph{Line 176}
\noindent\colorbox{gray!8}{\parbox{\dimexpr\linewidth-2\fboxsep\relax}{\raggedright\ttfamily\scriptsize         require\_file("rtl/cifar10\_spikingformer\_head\_features.mem");}}
\par\smallskip
This SystemVerilog line contributes to the structural or sequential description of the accelerator. Its effect is resolved by XSim for verification and by Vivado during synthesis.
\paragraph{Line 177}
\noindent\colorbox{gray!8}{\parbox{\dimexpr\linewidth-2\fboxsep\relax}{\raggedright\ttfamily\scriptsize         require\_file("rtl/cifar10\_spikingformer\_head\_weights.mem");}}
\par\smallskip
This SystemVerilog line contributes to the structural or sequential description of the accelerator. Its effect is resolved by XSim for verification and by Vivado during synthesis.
\paragraph{Line 178}
\noindent\colorbox{gray!8}{\parbox{\dimexpr\linewidth-2\fboxsep\relax}{\raggedright\ttfamily\scriptsize         require\_file("rtl/cifar10\_spikingformer\_head\_bias.mem");}}
\par\smallskip
This SystemVerilog line contributes to the structural or sequential description of the accelerator. Its effect is resolved by XSim for verification and by Vivado during synthesis.
\paragraph{Line 179}
\noindent\colorbox{gray!8}{\parbox{\dimexpr\linewidth-2\fboxsep\relax}{\raggedright\ttfamily\scriptsize         require\_file("rtl/cifar10\_spikingformer\_head\_labels.mem");}}
\par\smallskip
This SystemVerilog line contributes to the structural or sequential description of the accelerator. Its effect is resolved by XSim for verification and by Vivado during synthesis.
\paragraph{Line 180}
\noindent\colorbox{gray!8}{\parbox{\dimexpr\linewidth-2\fboxsep\relax}{\raggedright\ttfamily\scriptsize         \$readmemh("rtl/cifar10\_spikingformer\_head\_features.mem", feature\_mem);}}
\par\smallskip
This line touches the feature memory that holds one 384-dimensional fixed-point Spikingformer feature vector for the current image. Keeping the feature vector local avoids repeatedly streaming the same values for every class score.
\paragraph{Line 181}
\noindent\colorbox{gray!8}{\parbox{\dimexpr\linewidth-2\fboxsep\relax}{\raggedright\ttfamily\scriptsize         \$readmemh("rtl/cifar10\_spikingformer\_head\_labels.mem", label\_mem);}}
\par\smallskip
This reads a hexadecimal memory file into an RTL array. It is how exported trained weights and biases become deterministic FPGA memory contents.
\paragraph{Line 182}
\noindent\colorbox{gray!8}{\parbox{\dimexpr\linewidth-2\fboxsep\relax}{\raggedright\ttfamily\scriptsize \emph{blank line}}}
\par\smallskip
This blank line is not hardware or software behavior. It separates nearby declarations or operations so the file remains readable and reviewable.
\paragraph{Line 183}
\noindent\colorbox{gray!8}{\parbox{\dimexpr\linewidth-2\fboxsep\relax}{\raggedright\ttfamily\scriptsize         \$dumpfile("cifar10\_spikingformer\_head\_accelerator.vcd");}}
\par\smallskip
This SystemVerilog line contributes to the structural or sequential description of the accelerator. Its effect is resolved by XSim for verification and by Vivado during synthesis.
\paragraph{Line 184}
\noindent\colorbox{gray!8}{\parbox{\dimexpr\linewidth-2\fboxsep\relax}{\raggedright\ttfamily\scriptsize         \$dumpvars(0, tb\_cifar10\_spikingformer\_head\_accelerator);}}
\par\smallskip
This SystemVerilog line contributes to the structural or sequential description of the accelerator. Its effect is resolved by XSim for verification and by Vivado during synthesis.
\paragraph{Line 185}
\noindent\colorbox{gray!8}{\parbox{\dimexpr\linewidth-2\fboxsep\relax}{\raggedright\ttfamily\scriptsize         trace\_fh = \$fopen("cifar10\_spikingformer\_head\_trace.csv", "w");}}
\par\smallskip
This is a local compile-time or procedural assignment. In context it either defines a helper value or connects a temporary expression used by nearby hardware logic.
\paragraph{Line 186}
\noindent\colorbox{gray!8}{\parbox{\dimexpr\linewidth-2\fboxsep\relax}{\raggedright\ttfamily\scriptsize         result\_fh = \$fopen("cifar10\_spikingformer\_head\_results.csv", "w");}}
\par\smallskip
This is a local compile-time or procedural assignment. In context it either defines a helper value or connects a temporary expression used by nearby hardware logic.
\paragraph{Line 187}
\noindent\colorbox{gray!8}{\parbox{\dimexpr\linewidth-2\fboxsep\relax}{\raggedright\ttfamily\scriptsize         \$fwrite(trace\_fh, "time\_ps,cycle,rst\_n,state,start,feature\_valid,feature\_ready,feature\_idx,class\_idx,result\_valid,predicted\_label,infer\_skip\_count,train\_skip\_count\textbackslash{}n");}}
\par\smallskip
This line exposes or updates a hardware activity counter. The counter makes the feature gate visible in waveforms and CSV traces, so the power-saving opportunity is observable from RTL rather than only from Python statistics.
\paragraph{Line 188}
\noindent\colorbox{gray!8}{\parbox{\dimexpr\linewidth-2\fboxsep\relax}{\raggedright\ttfamily\scriptsize         \$fwrite(result\_fh, "sample,label,predicted,score,correct,infer\_skip\_count,train\_skip\_count");}}
\par\smallskip
This line exposes or updates a hardware activity counter. The counter makes the feature gate visible in waveforms and CSV traces, so the power-saving opportunity is observable from RTL rather than only from Python statistics.
\paragraph{Line 189}
\noindent\colorbox{gray!8}{\parbox{\dimexpr\linewidth-2\fboxsep\relax}{\raggedright\ttfamily\scriptsize         for (int c = 0; c < 10; c++) begin}}
\par\smallskip
This is a local compile-time or procedural assignment. In context it either defines a helper value or connects a temporary expression used by nearby hardware logic.
\paragraph{Line 190}
\noindent\colorbox{gray!8}{\parbox{\dimexpr\linewidth-2\fboxsep\relax}{\raggedright\ttfamily\scriptsize             \$fwrite(result\_fh, ",score\%0d", c);}}
\par\smallskip
This SystemVerilog line contributes to the structural or sequential description of the accelerator. Its effect is resolved by XSim for verification and by Vivado during synthesis.
\paragraph{Line 191}
\noindent\colorbox{gray!8}{\parbox{\dimexpr\linewidth-2\fboxsep\relax}{\raggedright\ttfamily\scriptsize         end}}
\par\smallskip
This SystemVerilog line contributes to the structural or sequential description of the accelerator. Its effect is resolved by XSim for verification and by Vivado during synthesis.
\paragraph{Line 192}
\noindent\colorbox{gray!8}{\parbox{\dimexpr\linewidth-2\fboxsep\relax}{\raggedright\ttfamily\scriptsize         \$fwrite(result\_fh, "\textbackslash{}n");}}
\par\smallskip
This SystemVerilog line contributes to the structural or sequential description of the accelerator. Its effect is resolved by XSim for verification and by Vivado during synthesis.
\paragraph{Line 193}
\noindent\colorbox{gray!8}{\parbox{\dimexpr\linewidth-2\fboxsep\relax}{\raggedright\ttfamily\scriptsize \emph{blank line}}}
\par\smallskip
This blank line is not hardware or software behavior. It separates nearby declarations or operations so the file remains readable and reviewable.
\paragraph{Line 194}
\noindent\colorbox{gray!8}{\parbox{\dimexpr\linewidth-2\fboxsep\relax}{\raggedright\ttfamily\scriptsize         start = 1'b0;}}
\par\smallskip
This is a local compile-time or procedural assignment. In context it either defines a helper value or connects a temporary expression used by nearby hardware logic.
\paragraph{Line 195}
\noindent\colorbox{gray!8}{\parbox{\dimexpr\linewidth-2\fboxsep\relax}{\raggedright\ttfamily\scriptsize         feature\_valid = 1'b0;}}
\par\smallskip
This is a local compile-time or procedural assignment. In context it either defines a helper value or connects a temporary expression used by nearby hardware logic.
\paragraph{Line 196}
\noindent\colorbox{gray!8}{\parbox{\dimexpr\linewidth-2\fboxsep\relax}{\raggedright\ttfamily\scriptsize         feature\_value = '0;}}
\par\smallskip
This is a local compile-time or procedural assignment. In context it either defines a helper value or connects a temporary expression used by nearby hardware logic.
\paragraph{Line 197}
\noindent\colorbox{gray!8}{\parbox{\dimexpr\linewidth-2\fboxsep\relax}{\raggedright\ttfamily\scriptsize         frame\_done = 1'b0;}}
\par\smallskip
This is a local compile-time or procedural assignment. In context it either defines a helper value or connects a temporary expression used by nearby hardware logic.
\paragraph{Line 198}
\noindent\colorbox{gray!8}{\parbox{\dimexpr\linewidth-2\fboxsep\relax}{\raggedright\ttfamily\scriptsize         train\_enable = 1'b0;}}
\par\smallskip
This is a local compile-time or procedural assignment. In context it either defines a helper value or connects a temporary expression used by nearby hardware logic.
\paragraph{Line 199}
\noindent\colorbox{gray!8}{\parbox{\dimexpr\linewidth-2\fboxsep\relax}{\raggedright\ttfamily\scriptsize         train\_label = 4'd0;}}
\par\smallskip
This is a local compile-time or procedural assignment. In context it either defines a helper value or connects a temporary expression used by nearby hardware logic.
\paragraph{Line 200}
\noindent\colorbox{gray!8}{\parbox{\dimexpr\linewidth-2\fboxsep\relax}{\raggedright\ttfamily\scriptsize \emph{blank line}}}
\par\smallskip
This blank line is not hardware or software behavior. It separates nearby declarations or operations so the file remains readable and reviewable.
\paragraph{Line 201}
\noindent\colorbox{gray!8}{\parbox{\dimexpr\linewidth-2\fboxsep\relax}{\raggedright\ttfamily\scriptsize         repeat (8) @(negedge clk);}}
\par\smallskip
This SystemVerilog line contributes to the structural or sequential description of the accelerator. Its effect is resolved by XSim for verification and by Vivado during synthesis.
\paragraph{Line 202}
\noindent\colorbox{gray!8}{\parbox{\dimexpr\linewidth-2\fboxsep\relax}{\raggedright\ttfamily\scriptsize         rst\_n = 1'b1;}}
\par\smallskip
This is a local compile-time or procedural assignment. In context it either defines a helper value or connects a temporary expression used by nearby hardware logic.
\paragraph{Line 203}
\noindent\colorbox{gray!8}{\parbox{\dimexpr\linewidth-2\fboxsep\relax}{\raggedright\ttfamily\scriptsize         repeat (4) @(negedge clk);}}
\par\smallskip
This SystemVerilog line contributes to the structural or sequential description of the accelerator. Its effect is resolved by XSim for verification and by Vivado during synthesis.
\paragraph{Line 204}
\noindent\colorbox{gray!8}{\parbox{\dimexpr\linewidth-2\fboxsep\relax}{\raggedright\ttfamily\scriptsize \emph{blank line}}}
\par\smallskip
This blank line is not hardware or software behavior. It separates nearby declarations or operations so the file remains readable and reviewable.
\paragraph{Line 205}
\noindent\colorbox{gray!8}{\parbox{\dimexpr\linewidth-2\fboxsep\relax}{\raggedright\ttfamily\scriptsize         for (int s = 0; s < NUM\_SAMPLES; s++) begin}}
\par\smallskip
This is a local compile-time or procedural assignment. In context it either defines a helper value or connects a temporary expression used by nearby hardware logic.
\paragraph{Line 206}
\noindent\colorbox{gray!8}{\parbox{\dimexpr\linewidth-2\fboxsep\relax}{\raggedright\ttfamily\scriptsize             send\_sample(s);}}
\par\smallskip
This SystemVerilog line contributes to the structural or sequential description of the accelerator. Its effect is resolved by XSim for verification and by Vivado during synthesis.
\paragraph{Line 207}
\noindent\colorbox{gray!8}{\parbox{\dimexpr\linewidth-2\fboxsep\relax}{\raggedright\ttfamily\scriptsize         end}}
\par\smallskip
This SystemVerilog line contributes to the structural or sequential description of the accelerator. Its effect is resolved by XSim for verification and by Vivado during synthesis.
\paragraph{Line 208}
\noindent\colorbox{gray!8}{\parbox{\dimexpr\linewidth-2\fboxsep\relax}{\raggedright\ttfamily\scriptsize \emph{blank line}}}
\par\smallskip
This blank line is not hardware or software behavior. It separates nearby declarations or operations so the file remains readable and reviewable.
\paragraph{Line 209}
\noindent\colorbox{gray!8}{\parbox{\dimexpr\linewidth-2\fboxsep\relax}{\raggedright\ttfamily\scriptsize         repeat (10) @(posedge clk);}}
\par\smallskip
This SystemVerilog line contributes to the structural or sequential description of the accelerator. Its effect is resolved by XSim for verification and by Vivado during synthesis.
\paragraph{Line 210}
\noindent\colorbox{gray!8}{\parbox{\dimexpr\linewidth-2\fboxsep\relax}{\raggedright\ttfamily\scriptsize         \$fclose(trace\_fh);}}
\par\smallskip
This SystemVerilog line contributes to the structural or sequential description of the accelerator. Its effect is resolved by XSim for verification and by Vivado during synthesis.
\paragraph{Line 211}
\noindent\colorbox{gray!8}{\parbox{\dimexpr\linewidth-2\fboxsep\relax}{\raggedright\ttfamily\scriptsize         \$fclose(result\_fh);}}
\par\smallskip
This SystemVerilog line contributes to the structural or sequential description of the accelerator. Its effect is resolved by XSim for verification and by Vivado during synthesis.
\paragraph{Line 212}
\noindent\colorbox{gray!8}{\parbox{\dimexpr\linewidth-2\fboxsep\relax}{\raggedright\ttfamily\scriptsize         \$display("tb\_cifar10\_spikingformer\_head\_accelerator samples=\%0d correct=\%0d errors=\%0d", NUM\_SAMPLES, correct, errors);}}
\par\smallskip
This is a local compile-time or procedural assignment. In context it either defines a helper value or connects a temporary expression used by nearby hardware logic.
\paragraph{Line 213}
\noindent\colorbox{gray!8}{\parbox{\dimexpr\linewidth-2\fboxsep\relax}{\raggedright\ttfamily\scriptsize         if (errors != 0) begin}}
\par\smallskip
This is a local compile-time or procedural assignment. In context it either defines a helper value or connects a temporary expression used by nearby hardware logic.
\paragraph{Line 214}
\noindent\colorbox{gray!8}{\parbox{\dimexpr\linewidth-2\fboxsep\relax}{\raggedright\ttfamily\scriptsize             \$fatal;}}
\par\smallskip
This SystemVerilog line contributes to the structural or sequential description of the accelerator. Its effect is resolved by XSim for verification and by Vivado during synthesis.
\paragraph{Line 215}
\noindent\colorbox{gray!8}{\parbox{\dimexpr\linewidth-2\fboxsep\relax}{\raggedright\ttfamily\scriptsize         end}}
\par\smallskip
This SystemVerilog line contributes to the structural or sequential description of the accelerator. Its effect is resolved by XSim for verification and by Vivado during synthesis.
\paragraph{Line 216}
\noindent\colorbox{gray!8}{\parbox{\dimexpr\linewidth-2\fboxsep\relax}{\raggedright\ttfamily\scriptsize         \$finish;}}
\par\smallskip
This SystemVerilog line contributes to the structural or sequential description of the accelerator. Its effect is resolved by XSim for verification and by Vivado during synthesis.
\paragraph{Line 217}
\noindent\colorbox{gray!8}{\parbox{\dimexpr\linewidth-2\fboxsep\relax}{\raggedright\ttfamily\scriptsize     end}}
\par\smallskip
This SystemVerilog line contributes to the structural or sequential description of the accelerator. Its effect is resolved by XSim for verification and by Vivado during synthesis.
\paragraph{Line 218}
\noindent\colorbox{gray!8}{\parbox{\dimexpr\linewidth-2\fboxsep\relax}{\raggedright\ttfamily\scriptsize \emph{blank line}}}
\par\smallskip
This blank line is not hardware or software behavior. It separates nearby declarations or operations so the file remains readable and reviewable.
\paragraph{Line 219}
\noindent\colorbox{gray!8}{\parbox{\dimexpr\linewidth-2\fboxsep\relax}{\raggedright\ttfamily\scriptsize     initial begin}}
\par\smallskip
This simulation/synthesis initialization block loads trained memories from generated .mem files. On FPGA, Vivado maps these initial values into BRAM/LUTRAM initialization data.
\paragraph{Line 220}
\noindent\colorbox{gray!8}{\parbox{\dimexpr\linewidth-2\fboxsep\relax}{\raggedright\ttfamily\scriptsize         \#2\_000\_000\_000;}}
\par\smallskip
This is a human-facing comment. It documents intent for the reader and is ignored by the compiler or interpreter.
\paragraph{Line 221}
\noindent\colorbox{gray!8}{\parbox{\dimexpr\linewidth-2\fboxsep\relax}{\raggedright\ttfamily\scriptsize         \$display("ERROR: tb\_cifar10\_spikingformer\_head\_accelerator watchdog timeout");}}
\par\smallskip
This SystemVerilog line contributes to the structural or sequential description of the accelerator. Its effect is resolved by XSim for verification and by Vivado during synthesis.
\paragraph{Line 222}
\noindent\colorbox{gray!8}{\parbox{\dimexpr\linewidth-2\fboxsep\relax}{\raggedright\ttfamily\scriptsize         \$fatal;}}
\par\smallskip
This SystemVerilog line contributes to the structural or sequential description of the accelerator. Its effect is resolved by XSim for verification and by Vivado during synthesis.
\paragraph{Line 223}
\noindent\colorbox{gray!8}{\parbox{\dimexpr\linewidth-2\fboxsep\relax}{\raggedright\ttfamily\scriptsize     end}}
\par\smallskip
This SystemVerilog line contributes to the structural or sequential description of the accelerator. Its effect is resolved by XSim for verification and by Vivado during synthesis.
\paragraph{Line 224}
\noindent\colorbox{gray!8}{\parbox{\dimexpr\linewidth-2\fboxsep\relax}{\raggedright\ttfamily\scriptsize endmodule}}
\par\smallskip
This closes the module definition. Synthesis stops adding logic to this circuit block at this point.
\clearpage
\subsection{cifar10\_neuromorphic\textbackslash{}rtl\textbackslash{}tb\_cifar10\_spikingformer\_head\_onchip\_learning.sv}
Each entry below gives the exact source line followed by its hardware/software role in the implemented accelerator pipeline.
\paragraph{Line 1}
\noindent\colorbox{gray!8}{\parbox{\dimexpr\linewidth-2\fboxsep\relax}{\raggedright\ttfamily\scriptsize `timescale 1ns / 1ps}}
\par\smallskip
This simulator directive sets the time unit and precision used by testbenches and waveform timestamps. It keeps delays such as clock half-periods numerically meaningful in XSim.
\paragraph{Line 2}
\noindent\colorbox{gray!8}{\parbox{\dimexpr\linewidth-2\fboxsep\relax}{\raggedright\ttfamily\scriptsize `include "rtl/cifar10\_spikingformer\_head\_params.svh"}}
\par\smallskip
This includes shared compile-time constants such as feature dimensions and fixed-point widths. The RTL and testbenches therefore agree on the same trained-model shape.
\paragraph{Line 3}
\noindent\colorbox{gray!8}{\parbox{\dimexpr\linewidth-2\fboxsep\relax}{\raggedright\ttfamily\scriptsize \emph{blank line}}}
\par\smallskip
This blank line is not hardware or software behavior. It separates nearby declarations or operations so the file remains readable and reviewable.
\paragraph{Line 4}
\noindent\colorbox{gray!8}{\parbox{\dimexpr\linewidth-2\fboxsep\relax}{\raggedright\ttfamily\scriptsize module tb\_cifar10\_spikingformer\_head\_onchip\_learning;}}
\par\smallskip
This begins a hardware module, which is the FPGA equivalent of a reusable circuit block. Everything until the matching endmodule describes ports, storage, combinational equations, and clocked behavior for this block.
\paragraph{Line 5}
\noindent\colorbox{gray!8}{\parbox{\dimexpr\linewidth-2\fboxsep\relax}{\raggedright\ttfamily\scriptsize     parameter int NUM\_SAMPLES = `CIFAR10\_SF\_HEAD\_NUM\_SAMPLES;}}
\par\smallskip
This is a compile-time hardware setting. Vivado resolves it before synthesis, so it controls the generated circuit size, fixed-point width, or learning rule without runtime overhead.
\paragraph{Line 6}
\noindent\colorbox{gray!8}{\parbox{\dimexpr\linewidth-2\fboxsep\relax}{\raggedright\ttfamily\scriptsize     parameter int TRACE\_SAMPLES = `CIFAR10\_SF\_HEAD\_TRACE\_SAMPLES;}}
\par\smallskip
This is a compile-time hardware setting. Vivado resolves it before synthesis, so it controls the generated circuit size, fixed-point width, or learning rule without runtime overhead.
\paragraph{Line 7}
\noindent\colorbox{gray!8}{\parbox{\dimexpr\linewidth-2\fboxsep\relax}{\raggedright\ttfamily\scriptsize     parameter int FEATURE\_DIM = `CIFAR10\_SF\_HEAD\_FEATURE\_DIM;}}
\par\smallskip
This is a compile-time hardware setting. Vivado resolves it before synthesis, so it controls the generated circuit size, fixed-point width, or learning rule without runtime overhead.
\paragraph{Line 8}
\noindent\colorbox{gray!8}{\parbox{\dimexpr\linewidth-2\fboxsep\relax}{\raggedright\ttfamily\scriptsize     parameter int MAX\_SAMPLES = 10000;}}
\par\smallskip
This is a compile-time hardware setting. Vivado resolves it before synthesis, so it controls the generated circuit size, fixed-point width, or learning rule without runtime overhead.
\paragraph{Line 9}
\noindent\colorbox{gray!8}{\parbox{\dimexpr\linewidth-2\fboxsep\relax}{\raggedright\ttfamily\scriptsize     parameter int MAX\_FEATURE\_WORDS = MAX\_SAMPLES * FEATURE\_DIM;}}
\par\smallskip
This is a compile-time hardware setting. Vivado resolves it before synthesis, so it controls the generated circuit size, fixed-point width, or learning rule without runtime overhead.
\paragraph{Line 10}
\noindent\colorbox{gray!8}{\parbox{\dimexpr\linewidth-2\fboxsep\relax}{\raggedright\ttfamily\scriptsize \emph{blank line}}}
\par\smallskip
This blank line is not hardware or software behavior. It separates nearby declarations or operations so the file remains readable and reviewable.
\paragraph{Line 11}
\noindent\colorbox{gray!8}{\parbox{\dimexpr\linewidth-2\fboxsep\relax}{\raggedright\ttfamily\scriptsize     logic clk = 1'b0;}}
\par\smallskip
This is a local compile-time or procedural assignment. In context it either defines a helper value or connects a temporary expression used by nearby hardware logic.
\paragraph{Line 12}
\noindent\colorbox{gray!8}{\parbox{\dimexpr\linewidth-2\fboxsep\relax}{\raggedright\ttfamily\scriptsize     logic rst\_n = 1'b0;}}
\par\smallskip
This is a local compile-time or procedural assignment. In context it either defines a helper value or connects a temporary expression used by nearby hardware logic.
\paragraph{Line 13}
\noindent\colorbox{gray!8}{\parbox{\dimexpr\linewidth-2\fboxsep\relax}{\raggedright\ttfamily\scriptsize     logic start;}}
\par\smallskip
This SystemVerilog line contributes to the structural or sequential description of the accelerator. Its effect is resolved by XSim for verification and by Vivado during synthesis.
\paragraph{Line 14}
\noindent\colorbox{gray!8}{\parbox{\dimexpr\linewidth-2\fboxsep\relax}{\raggedright\ttfamily\scriptsize     logic feature\_valid;}}
\par\smallskip
This SystemVerilog line contributes to the structural or sequential description of the accelerator. Its effect is resolved by XSim for verification and by Vivado during synthesis.
\paragraph{Line 15}
\noindent\colorbox{gray!8}{\parbox{\dimexpr\linewidth-2\fboxsep\relax}{\raggedright\ttfamily\scriptsize     logic feature\_ready;}}
\par\smallskip
This SystemVerilog line contributes to the structural or sequential description of the accelerator. Its effect is resolved by XSim for verification and by Vivado during synthesis.
\paragraph{Line 16}
\noindent\colorbox{gray!8}{\parbox{\dimexpr\linewidth-2\fboxsep\relax}{\raggedright\ttfamily\scriptsize     logic signed [`CIFAR10\_SF\_HEAD\_FEATURE\_BITS-1:0] feature\_value;}}
\par\smallskip
This SystemVerilog line contributes to the structural or sequential description of the accelerator. Its effect is resolved by XSim for verification and by Vivado during synthesis.
\paragraph{Line 17}
\noindent\colorbox{gray!8}{\parbox{\dimexpr\linewidth-2\fboxsep\relax}{\raggedright\ttfamily\scriptsize     logic frame\_done;}}
\par\smallskip
This SystemVerilog line contributes to the structural or sequential description of the accelerator. Its effect is resolved by XSim for verification and by Vivado during synthesis.
\paragraph{Line 18}
\noindent\colorbox{gray!8}{\parbox{\dimexpr\linewidth-2\fboxsep\relax}{\raggedright\ttfamily\scriptsize     logic train\_enable;}}
\par\smallskip
This SystemVerilog line contributes to the structural or sequential description of the accelerator. Its effect is resolved by XSim for verification and by Vivado during synthesis.
\paragraph{Line 19}
\noindent\colorbox{gray!8}{\parbox{\dimexpr\linewidth-2\fboxsep\relax}{\raggedright\ttfamily\scriptsize     logic [3:0] train\_label;}}
\par\smallskip
This SystemVerilog line contributes to the structural or sequential description of the accelerator. Its effect is resolved by XSim for verification and by Vivado during synthesis.
\paragraph{Line 20}
\noindent\colorbox{gray!8}{\parbox{\dimexpr\linewidth-2\fboxsep\relax}{\raggedright\ttfamily\scriptsize     logic busy;}}
\par\smallskip
This SystemVerilog line contributes to the structural or sequential description of the accelerator. Its effect is resolved by XSim for verification and by Vivado during synthesis.
\paragraph{Line 21}
\noindent\colorbox{gray!8}{\parbox{\dimexpr\linewidth-2\fboxsep\relax}{\raggedright\ttfamily\scriptsize     logic result\_valid;}}
\par\smallskip
This SystemVerilog line contributes to the structural or sequential description of the accelerator. Its effect is resolved by XSim for verification and by Vivado during synthesis.
\paragraph{Line 22}
\noindent\colorbox{gray!8}{\parbox{\dimexpr\linewidth-2\fboxsep\relax}{\raggedright\ttfamily\scriptsize     logic [3:0] predicted\_label;}}
\par\smallskip
This SystemVerilog line contributes to the structural or sequential description of the accelerator. Its effect is resolved by XSim for verification and by Vivado during synthesis.
\paragraph{Line 23}
\noindent\colorbox{gray!8}{\parbox{\dimexpr\linewidth-2\fboxsep\relax}{\raggedright\ttfamily\scriptsize     logic signed [`CIFAR10\_SF\_HEAD\_ACC\_BITS-1:0] predicted\_score;}}
\par\smallskip
This SystemVerilog line contributes to the structural or sequential description of the accelerator. Its effect is resolved by XSim for verification and by Vivado during synthesis.
\paragraph{Line 24}
\noindent\colorbox{gray!8}{\parbox{\dimexpr\linewidth-2\fboxsep\relax}{\raggedright\ttfamily\scriptsize     logic train\_update\_valid;}}
\par\smallskip
This SystemVerilog line contributes to the structural or sequential description of the accelerator. Its effect is resolved by XSim for verification and by Vivado during synthesis.
\paragraph{Line 25}
\noindent\colorbox{gray!8}{\parbox{\dimexpr\linewidth-2\fboxsep\relax}{\raggedright\ttfamily\scriptsize     logic train\_update\_applied;}}
\par\smallskip
This SystemVerilog line contributes to the structural or sequential description of the accelerator. Its effect is resolved by XSim for verification and by Vivado during synthesis.
\paragraph{Line 26}
\noindent\colorbox{gray!8}{\parbox{\dimexpr\linewidth-2\fboxsep\relax}{\raggedright\ttfamily\scriptsize     logic [5:0] debug\_state;}}
\par\smallskip
This SystemVerilog line contributes to the structural or sequential description of the accelerator. Its effect is resolved by XSim for verification and by Vivado during synthesis.
\paragraph{Line 27}
\noindent\colorbox{gray!8}{\parbox{\dimexpr\linewidth-2\fboxsep\relax}{\raggedright\ttfamily\scriptsize     logic [15:0] debug\_feature;}}
\par\smallskip
This SystemVerilog line contributes to the structural or sequential description of the accelerator. Its effect is resolved by XSim for verification and by Vivado during synthesis.
\paragraph{Line 28}
\noindent\colorbox{gray!8}{\parbox{\dimexpr\linewidth-2\fboxsep\relax}{\raggedright\ttfamily\scriptsize     logic [7:0] debug\_class;}}
\par\smallskip
This SystemVerilog line contributes to the structural or sequential description of the accelerator. Its effect is resolved by XSim for verification and by Vivado during synthesis.
\paragraph{Line 29}
\noindent\colorbox{gray!8}{\parbox{\dimexpr\linewidth-2\fboxsep\relax}{\raggedright\ttfamily\scriptsize     logic [31:0] debug\_update\_count;}}
\par\smallskip
This SystemVerilog line contributes to the structural or sequential description of the accelerator. Its effect is resolved by XSim for verification and by Vivado during synthesis.
\paragraph{Line 30}
\noindent\colorbox{gray!8}{\parbox{\dimexpr\linewidth-2\fboxsep\relax}{\raggedright\ttfamily\scriptsize     logic [31:0] debug\_infer\_skip\_count;}}
\par\smallskip
This line exposes or updates a hardware activity counter. The counter makes the feature gate visible in waveforms and CSV traces, so the power-saving opportunity is observable from RTL rather than only from Python statistics.
\paragraph{Line 31}
\noindent\colorbox{gray!8}{\parbox{\dimexpr\linewidth-2\fboxsep\relax}{\raggedright\ttfamily\scriptsize     logic [31:0] debug\_train\_skip\_count;}}
\par\smallskip
This line exposes or updates a hardware activity counter. The counter makes the feature gate visible in waveforms and CSV traces, so the power-saving opportunity is observable from RTL rather than only from Python statistics.
\paragraph{Line 32}
\noindent\colorbox{gray!8}{\parbox{\dimexpr\linewidth-2\fboxsep\relax}{\raggedright\ttfamily\scriptsize \emph{blank line}}}
\par\smallskip
This blank line is not hardware or software behavior. It separates nearby declarations or operations so the file remains readable and reviewable.
\paragraph{Line 33}
\noindent\colorbox{gray!8}{\parbox{\dimexpr\linewidth-2\fboxsep\relax}{\raggedright\ttfamily\scriptsize     logic signed [`CIFAR10\_SF\_HEAD\_FEATURE\_BITS-1:0] feature\_mem [0:MAX\_FEATURE\_WORDS-1];}}
\par\smallskip
This line touches the feature memory that holds one 384-dimensional fixed-point Spikingformer feature vector for the current image. Keeping the feature vector local avoids repeatedly streaming the same values for every class score.
\paragraph{Line 34}
\noindent\colorbox{gray!8}{\parbox{\dimexpr\linewidth-2\fboxsep\relax}{\raggedright\ttfamily\scriptsize     logic [3:0] label\_mem [0:MAX\_SAMPLES-1];}}
\par\smallskip
This SystemVerilog line contributes to the structural or sequential description of the accelerator. Its effect is resolved by XSim for verification and by Vivado during synthesis.
\paragraph{Line 35}
\noindent\colorbox{gray!8}{\parbox{\dimexpr\linewidth-2\fboxsep\relax}{\raggedright\ttfamily\scriptsize \emph{blank line}}}
\par\smallskip
This blank line is not hardware or software behavior. It separates nearby declarations or operations so the file remains readable and reviewable.
\paragraph{Line 36}
\noindent\colorbox{gray!8}{\parbox{\dimexpr\linewidth-2\fboxsep\relax}{\raggedright\ttfamily\scriptsize     int feature\_ptr = 0;}}
\par\smallskip
This is a local compile-time or procedural assignment. In context it either defines a helper value or connects a temporary expression used by nearby hardware logic.
\paragraph{Line 37}
\noindent\colorbox{gray!8}{\parbox{\dimexpr\linewidth-2\fboxsep\relax}{\raggedright\ttfamily\scriptsize     int correct = 0;}}
\par\smallskip
This is a local compile-time or procedural assignment. In context it either defines a helper value or connects a temporary expression used by nearby hardware logic.
\paragraph{Line 38}
\noindent\colorbox{gray!8}{\parbox{\dimexpr\linewidth-2\fboxsep\relax}{\raggedright\ttfamily\scriptsize     int errors = 0;}}
\par\smallskip
This is a local compile-time or procedural assignment. In context it either defines a helper value or connects a temporary expression used by nearby hardware logic.
\paragraph{Line 39}
\noindent\colorbox{gray!8}{\parbox{\dimexpr\linewidth-2\fboxsep\relax}{\raggedright\ttfamily\scriptsize     int updates = 0;}}
\par\smallskip
This is a local compile-time or procedural assignment. In context it either defines a helper value or connects a temporary expression used by nearby hardware logic.
\paragraph{Line 40}
\noindent\colorbox{gray!8}{\parbox{\dimexpr\linewidth-2\fboxsep\relax}{\raggedright\ttfamily\scriptsize     int cycle = 0;}}
\par\smallskip
This is a local compile-time or procedural assignment. In context it either defines a helper value or connects a temporary expression used by nearby hardware logic.
\paragraph{Line 41}
\noindent\colorbox{gray!8}{\parbox{\dimexpr\linewidth-2\fboxsep\relax}{\raggedright\ttfamily\scriptsize     int trace\_fh;}}
\par\smallskip
This SystemVerilog line contributes to the structural or sequential description of the accelerator. Its effect is resolved by XSim for verification and by Vivado during synthesis.
\paragraph{Line 42}
\noindent\colorbox{gray!8}{\parbox{\dimexpr\linewidth-2\fboxsep\relax}{\raggedright\ttfamily\scriptsize     int result\_fh;}}
\par\smallskip
This SystemVerilog line contributes to the structural or sequential description of the accelerator. Its effect is resolved by XSim for verification and by Vivado during synthesis.
\paragraph{Line 43}
\noindent\colorbox{gray!8}{\parbox{\dimexpr\linewidth-2\fboxsep\relax}{\raggedright\ttfamily\scriptsize     int weight\_fh;}}
\par\smallskip
This SystemVerilog line contributes to the structural or sequential description of the accelerator. Its effect is resolved by XSim for verification and by Vivado during synthesis.
\paragraph{Line 44}
\noindent\colorbox{gray!8}{\parbox{\dimexpr\linewidth-2\fboxsep\relax}{\raggedright\ttfamily\scriptsize     int bias\_fh;}}
\par\smallskip
This SystemVerilog line contributes to the structural or sequential description of the accelerator. Its effect is resolved by XSim for verification and by Vivado during synthesis.
\paragraph{Line 45}
\noindent\colorbox{gray!8}{\parbox{\dimexpr\linewidth-2\fboxsep\relax}{\raggedright\ttfamily\scriptsize     bit trace\_enable = 1'b0;}}
\par\smallskip
This is a local compile-time or procedural assignment. In context it either defines a helper value or connects a temporary expression used by nearby hardware logic.
\paragraph{Line 46}
\noindent\colorbox{gray!8}{\parbox{\dimexpr\linewidth-2\fboxsep\relax}{\raggedright\ttfamily\scriptsize     bit vcd\_enabled = 1'b1;}}
\par\smallskip
This is a local compile-time or procedural assignment. In context it either defines a helper value or connects a temporary expression used by nearby hardware logic.
\paragraph{Line 47}
\noindent\colorbox{gray!8}{\parbox{\dimexpr\linewidth-2\fboxsep\relax}{\raggedright\ttfamily\scriptsize \emph{blank line}}}
\par\smallskip
This blank line is not hardware or software behavior. It separates nearby declarations or operations so the file remains readable and reviewable.
\paragraph{Line 48}
\noindent\colorbox{gray!8}{\parbox{\dimexpr\linewidth-2\fboxsep\relax}{\raggedright\ttfamily\scriptsize     cifar10\_spikingformer\_head\_accelerator \#(}}
\par\smallskip
This SystemVerilog line contributes to the structural or sequential description of the accelerator. Its effect is resolved by XSim for verification and by Vivado during synthesis.
\paragraph{Line 49}
\noindent\colorbox{gray!8}{\parbox{\dimexpr\linewidth-2\fboxsep\relax}{\raggedright\ttfamily\scriptsize         .ENABLE\_ONCHIP\_LEARNING(1),}}
\par\smallskip
This SystemVerilog line contributes to the structural or sequential description of the accelerator. Its effect is resolved by XSim for verification and by Vivado during synthesis.
\paragraph{Line 50}
\noindent\colorbox{gray!8}{\parbox{\dimexpr\linewidth-2\fboxsep\relax}{\raggedright\ttfamily\scriptsize         .LEARN\_SHIFT(14),}}
\par\smallskip
This SystemVerilog line contributes to the structural or sequential description of the accelerator. Its effect is resolved by XSim for verification and by Vivado during synthesis.
\paragraph{Line 51}
\noindent\colorbox{gray!8}{\parbox{\dimexpr\linewidth-2\fboxsep\relax}{\raggedright\ttfamily\scriptsize         .BIAS\_LEARN\_DELTA(1),}}
\par\smallskip
This SystemVerilog line contributes to the structural or sequential description of the accelerator. Its effect is resolved by XSim for verification and by Vivado during synthesis.
\paragraph{Line 52}
\noindent\colorbox{gray!8}{\parbox{\dimexpr\linewidth-2\fboxsep\relax}{\raggedright\ttfamily\scriptsize         .LEARN\_MARGIN\_Q(64'sd268435456)}}
\par\smallskip
This SystemVerilog line contributes to the structural or sequential description of the accelerator. Its effect is resolved by XSim for verification and by Vivado during synthesis.
\paragraph{Line 53}
\noindent\colorbox{gray!8}{\parbox{\dimexpr\linewidth-2\fboxsep\relax}{\raggedright\ttfamily\scriptsize     ) dut (}}
\par\smallskip
This SystemVerilog line contributes to the structural or sequential description of the accelerator. Its effect is resolved by XSim for verification and by Vivado during synthesis.
\paragraph{Line 54}
\noindent\colorbox{gray!8}{\parbox{\dimexpr\linewidth-2\fboxsep\relax}{\raggedright\ttfamily\scriptsize         .clk(clk),}}
\par\smallskip
This SystemVerilog line contributes to the structural or sequential description of the accelerator. Its effect is resolved by XSim for verification and by Vivado during synthesis.
\paragraph{Line 55}
\noindent\colorbox{gray!8}{\parbox{\dimexpr\linewidth-2\fboxsep\relax}{\raggedright\ttfamily\scriptsize         .rst\_n(rst\_n),}}
\par\smallskip
This SystemVerilog line contributes to the structural or sequential description of the accelerator. Its effect is resolved by XSim for verification and by Vivado during synthesis.
\paragraph{Line 56}
\noindent\colorbox{gray!8}{\parbox{\dimexpr\linewidth-2\fboxsep\relax}{\raggedright\ttfamily\scriptsize         .start(start),}}
\par\smallskip
This SystemVerilog line contributes to the structural or sequential description of the accelerator. Its effect is resolved by XSim for verification and by Vivado during synthesis.
\paragraph{Line 57}
\noindent\colorbox{gray!8}{\parbox{\dimexpr\linewidth-2\fboxsep\relax}{\raggedright\ttfamily\scriptsize         .feature\_valid(feature\_valid),}}
\par\smallskip
This SystemVerilog line contributes to the structural or sequential description of the accelerator. Its effect is resolved by XSim for verification and by Vivado during synthesis.
\paragraph{Line 58}
\noindent\colorbox{gray!8}{\parbox{\dimexpr\linewidth-2\fboxsep\relax}{\raggedright\ttfamily\scriptsize         .feature\_ready(feature\_ready),}}
\par\smallskip
This SystemVerilog line contributes to the structural or sequential description of the accelerator. Its effect is resolved by XSim for verification and by Vivado during synthesis.
\paragraph{Line 59}
\noindent\colorbox{gray!8}{\parbox{\dimexpr\linewidth-2\fboxsep\relax}{\raggedright\ttfamily\scriptsize         .feature\_value(feature\_value),}}
\par\smallskip
This SystemVerilog line contributes to the structural or sequential description of the accelerator. Its effect is resolved by XSim for verification and by Vivado during synthesis.
\paragraph{Line 60}
\noindent\colorbox{gray!8}{\parbox{\dimexpr\linewidth-2\fboxsep\relax}{\raggedright\ttfamily\scriptsize         .frame\_done(frame\_done),}}
\par\smallskip
This SystemVerilog line contributes to the structural or sequential description of the accelerator. Its effect is resolved by XSim for verification and by Vivado during synthesis.
\paragraph{Line 61}
\noindent\colorbox{gray!8}{\parbox{\dimexpr\linewidth-2\fboxsep\relax}{\raggedright\ttfamily\scriptsize         .train\_enable(train\_enable),}}
\par\smallskip
This SystemVerilog line contributes to the structural or sequential description of the accelerator. Its effect is resolved by XSim for verification and by Vivado during synthesis.
\paragraph{Line 62}
\noindent\colorbox{gray!8}{\parbox{\dimexpr\linewidth-2\fboxsep\relax}{\raggedright\ttfamily\scriptsize         .train\_label(train\_label),}}
\par\smallskip
This SystemVerilog line contributes to the structural or sequential description of the accelerator. Its effect is resolved by XSim for verification and by Vivado during synthesis.
\paragraph{Line 63}
\noindent\colorbox{gray!8}{\parbox{\dimexpr\linewidth-2\fboxsep\relax}{\raggedright\ttfamily\scriptsize         .busy(busy),}}
\par\smallskip
This SystemVerilog line contributes to the structural or sequential description of the accelerator. Its effect is resolved by XSim for verification and by Vivado during synthesis.
\paragraph{Line 64}
\noindent\colorbox{gray!8}{\parbox{\dimexpr\linewidth-2\fboxsep\relax}{\raggedright\ttfamily\scriptsize         .result\_valid(result\_valid),}}
\par\smallskip
This SystemVerilog line contributes to the structural or sequential description of the accelerator. Its effect is resolved by XSim for verification and by Vivado during synthesis.
\paragraph{Line 65}
\noindent\colorbox{gray!8}{\parbox{\dimexpr\linewidth-2\fboxsep\relax}{\raggedright\ttfamily\scriptsize         .predicted\_label(predicted\_label),}}
\par\smallskip
This SystemVerilog line contributes to the structural or sequential description of the accelerator. Its effect is resolved by XSim for verification and by Vivado during synthesis.
\paragraph{Line 66}
\noindent\colorbox{gray!8}{\parbox{\dimexpr\linewidth-2\fboxsep\relax}{\raggedright\ttfamily\scriptsize         .predicted\_score(predicted\_score),}}
\par\smallskip
This SystemVerilog line contributes to the structural or sequential description of the accelerator. Its effect is resolved by XSim for verification and by Vivado during synthesis.
\paragraph{Line 67}
\noindent\colorbox{gray!8}{\parbox{\dimexpr\linewidth-2\fboxsep\relax}{\raggedright\ttfamily\scriptsize         .train\_update\_valid(train\_update\_valid),}}
\par\smallskip
This SystemVerilog line contributes to the structural or sequential description of the accelerator. Its effect is resolved by XSim for verification and by Vivado during synthesis.
\paragraph{Line 68}
\noindent\colorbox{gray!8}{\parbox{\dimexpr\linewidth-2\fboxsep\relax}{\raggedright\ttfamily\scriptsize         .train\_update\_applied(train\_update\_applied),}}
\par\smallskip
This SystemVerilog line contributes to the structural or sequential description of the accelerator. Its effect is resolved by XSim for verification and by Vivado during synthesis.
\paragraph{Line 69}
\noindent\colorbox{gray!8}{\parbox{\dimexpr\linewidth-2\fboxsep\relax}{\raggedright\ttfamily\scriptsize         .debug\_state(debug\_state),}}
\par\smallskip
This SystemVerilog line contributes to the structural or sequential description of the accelerator. Its effect is resolved by XSim for verification and by Vivado during synthesis.
\paragraph{Line 70}
\noindent\colorbox{gray!8}{\parbox{\dimexpr\linewidth-2\fboxsep\relax}{\raggedright\ttfamily\scriptsize         .debug\_feature(debug\_feature),}}
\par\smallskip
This SystemVerilog line contributes to the structural or sequential description of the accelerator. Its effect is resolved by XSim for verification and by Vivado during synthesis.
\paragraph{Line 71}
\noindent\colorbox{gray!8}{\parbox{\dimexpr\linewidth-2\fboxsep\relax}{\raggedright\ttfamily\scriptsize         .debug\_class(debug\_class),}}
\par\smallskip
This SystemVerilog line contributes to the structural or sequential description of the accelerator. Its effect is resolved by XSim for verification and by Vivado during synthesis.
\paragraph{Line 72}
\noindent\colorbox{gray!8}{\parbox{\dimexpr\linewidth-2\fboxsep\relax}{\raggedright\ttfamily\scriptsize         .debug\_update\_count(debug\_update\_count),}}
\par\smallskip
This SystemVerilog line contributes to the structural or sequential description of the accelerator. Its effect is resolved by XSim for verification and by Vivado during synthesis.
\paragraph{Line 73}
\noindent\colorbox{gray!8}{\parbox{\dimexpr\linewidth-2\fboxsep\relax}{\raggedright\ttfamily\scriptsize         .debug\_infer\_skip\_count(debug\_infer\_skip\_count),}}
\par\smallskip
This line exposes or updates a hardware activity counter. The counter makes the feature gate visible in waveforms and CSV traces, so the power-saving opportunity is observable from RTL rather than only from Python statistics.
\paragraph{Line 74}
\noindent\colorbox{gray!8}{\parbox{\dimexpr\linewidth-2\fboxsep\relax}{\raggedright\ttfamily\scriptsize         .debug\_train\_skip\_count(debug\_train\_skip\_count)}}
\par\smallskip
This line exposes or updates a hardware activity counter. The counter makes the feature gate visible in waveforms and CSV traces, so the power-saving opportunity is observable from RTL rather than only from Python statistics.
\paragraph{Line 75}
\noindent\colorbox{gray!8}{\parbox{\dimexpr\linewidth-2\fboxsep\relax}{\raggedright\ttfamily\scriptsize     );}}
\par\smallskip
This SystemVerilog line contributes to the structural or sequential description of the accelerator. Its effect is resolved by XSim for verification and by Vivado during synthesis.
\paragraph{Line 76}
\noindent\colorbox{gray!8}{\parbox{\dimexpr\linewidth-2\fboxsep\relax}{\raggedright\ttfamily\scriptsize \emph{blank line}}}
\par\smallskip
This blank line is not hardware or software behavior. It separates nearby declarations or operations so the file remains readable and reviewable.
\paragraph{Line 77}
\noindent\colorbox{gray!8}{\parbox{\dimexpr\linewidth-2\fboxsep\relax}{\raggedright\ttfamily\scriptsize     always \#5 clk = \textasciitilde{}clk;}}
\par\smallskip
This is a local compile-time or procedural assignment. In context it either defines a helper value or connects a temporary expression used by nearby hardware logic.
\paragraph{Line 78}
\noindent\colorbox{gray!8}{\parbox{\dimexpr\linewidth-2\fboxsep\relax}{\raggedright\ttfamily\scriptsize \emph{blank line}}}
\par\smallskip
This blank line is not hardware or software behavior. It separates nearby declarations or operations so the file remains readable and reviewable.
\paragraph{Line 79}
\noindent\colorbox{gray!8}{\parbox{\dimexpr\linewidth-2\fboxsep\relax}{\raggedright\ttfamily\scriptsize     task automatic require\_file(input string path);}}
\par\smallskip
This line declares a module port. A port is a physical or wrapper-visible signal used to move clocks, control pulses, features, labels, predictions, or debug information across module boundaries.
\paragraph{Line 80}
\noindent\colorbox{gray!8}{\parbox{\dimexpr\linewidth-2\fboxsep\relax}{\raggedright\ttfamily\scriptsize         int fh;}}
\par\smallskip
This SystemVerilog line contributes to the structural or sequential description of the accelerator. Its effect is resolved by XSim for verification and by Vivado during synthesis.
\paragraph{Line 81}
\noindent\colorbox{gray!8}{\parbox{\dimexpr\linewidth-2\fboxsep\relax}{\raggedright\ttfamily\scriptsize         begin}}
\par\smallskip
This SystemVerilog line contributes to the structural or sequential description of the accelerator. Its effect is resolved by XSim for verification and by Vivado during synthesis.
\paragraph{Line 82}
\noindent\colorbox{gray!8}{\parbox{\dimexpr\linewidth-2\fboxsep\relax}{\raggedright\ttfamily\scriptsize             fh = \$fopen(path, "r");}}
\par\smallskip
This is a local compile-time or procedural assignment. In context it either defines a helper value or connects a temporary expression used by nearby hardware logic.
\paragraph{Line 83}
\noindent\colorbox{gray!8}{\parbox{\dimexpr\linewidth-2\fboxsep\relax}{\raggedright\ttfamily\scriptsize             if (fh == 0) begin}}
\par\smallskip
This is a local compile-time or procedural assignment. In context it either defines a helper value or connects a temporary expression used by nearby hardware logic.
\paragraph{Line 84}
\noindent\colorbox{gray!8}{\parbox{\dimexpr\linewidth-2\fboxsep\relax}{\raggedright\ttfamily\scriptsize                 \$display("ERROR: missing required file \%s", path);}}
\par\smallskip
This SystemVerilog line contributes to the structural or sequential description of the accelerator. Its effect is resolved by XSim for verification and by Vivado during synthesis.
\paragraph{Line 85}
\noindent\colorbox{gray!8}{\parbox{\dimexpr\linewidth-2\fboxsep\relax}{\raggedright\ttfamily\scriptsize                 \$fatal;}}
\par\smallskip
This SystemVerilog line contributes to the structural or sequential description of the accelerator. Its effect is resolved by XSim for verification and by Vivado during synthesis.
\paragraph{Line 86}
\noindent\colorbox{gray!8}{\parbox{\dimexpr\linewidth-2\fboxsep\relax}{\raggedright\ttfamily\scriptsize             end}}
\par\smallskip
This SystemVerilog line contributes to the structural or sequential description of the accelerator. Its effect is resolved by XSim for verification and by Vivado during synthesis.
\paragraph{Line 87}
\noindent\colorbox{gray!8}{\parbox{\dimexpr\linewidth-2\fboxsep\relax}{\raggedright\ttfamily\scriptsize             \$fclose(fh);}}
\par\smallskip
This SystemVerilog line contributes to the structural or sequential description of the accelerator. Its effect is resolved by XSim for verification and by Vivado during synthesis.
\paragraph{Line 88}
\noindent\colorbox{gray!8}{\parbox{\dimexpr\linewidth-2\fboxsep\relax}{\raggedright\ttfamily\scriptsize         end}}
\par\smallskip
This SystemVerilog line contributes to the structural or sequential description of the accelerator. Its effect is resolved by XSim for verification and by Vivado during synthesis.
\paragraph{Line 89}
\noindent\colorbox{gray!8}{\parbox{\dimexpr\linewidth-2\fboxsep\relax}{\raggedright\ttfamily\scriptsize     endtask}}
\par\smallskip
This SystemVerilog line contributes to the structural or sequential description of the accelerator. Its effect is resolved by XSim for verification and by Vivado during synthesis.
\paragraph{Line 90}
\noindent\colorbox{gray!8}{\parbox{\dimexpr\linewidth-2\fboxsep\relax}{\raggedright\ttfamily\scriptsize \emph{blank line}}}
\par\smallskip
This blank line is not hardware or software behavior. It separates nearby declarations or operations so the file remains readable and reviewable.
\paragraph{Line 91}
\noindent\colorbox{gray!8}{\parbox{\dimexpr\linewidth-2\fboxsep\relax}{\raggedright\ttfamily\scriptsize     task automatic send\_sample(input int sample\_idx);}}
\par\smallskip
This line declares a module port. A port is a physical or wrapper-visible signal used to move clocks, control pulses, features, labels, predictions, or debug information across module boundaries.
\paragraph{Line 92}
\noindent\colorbox{gray!8}{\parbox{\dimexpr\linewidth-2\fboxsep\relax}{\raggedright\ttfamily\scriptsize         int waited;}}
\par\smallskip
This SystemVerilog line contributes to the structural or sequential description of the accelerator. Its effect is resolved by XSim for verification and by Vivado during synthesis.
\paragraph{Line 93}
\noindent\colorbox{gray!8}{\parbox{\dimexpr\linewidth-2\fboxsep\relax}{\raggedright\ttfamily\scriptsize         bit update\_seen;}}
\par\smallskip
This SystemVerilog line contributes to the structural or sequential description of the accelerator. Its effect is resolved by XSim for verification and by Vivado during synthesis.
\paragraph{Line 94}
\noindent\colorbox{gray!8}{\parbox{\dimexpr\linewidth-2\fboxsep\relax}{\raggedright\ttfamily\scriptsize         begin}}
\par\smallskip
This SystemVerilog line contributes to the structural or sequential description of the accelerator. Its effect is resolved by XSim for verification and by Vivado during synthesis.
\paragraph{Line 95}
\noindent\colorbox{gray!8}{\parbox{\dimexpr\linewidth-2\fboxsep\relax}{\raggedright\ttfamily\scriptsize             trace\_enable = (sample\_idx < TRACE\_SAMPLES);}}
\par\smallskip
This is a local compile-time or procedural assignment. In context it either defines a helper value or connects a temporary expression used by nearby hardware logic.
\paragraph{Line 96}
\noindent\colorbox{gray!8}{\parbox{\dimexpr\linewidth-2\fboxsep\relax}{\raggedright\ttfamily\scriptsize             if (sample\_idx == TRACE\_SAMPLES \&\& vcd\_enabled) begin}}
\par\smallskip
This is a local compile-time or procedural assignment. In context it either defines a helper value or connects a temporary expression used by nearby hardware logic.
\paragraph{Line 97}
\noindent\colorbox{gray!8}{\parbox{\dimexpr\linewidth-2\fboxsep\relax}{\raggedright\ttfamily\scriptsize                 \$dumpoff;}}
\par\smallskip
This SystemVerilog line contributes to the structural or sequential description of the accelerator. Its effect is resolved by XSim for verification and by Vivado during synthesis.
\paragraph{Line 98}
\noindent\colorbox{gray!8}{\parbox{\dimexpr\linewidth-2\fboxsep\relax}{\raggedright\ttfamily\scriptsize                 vcd\_enabled = 1'b0;}}
\par\smallskip
This is a local compile-time or procedural assignment. In context it either defines a helper value or connects a temporary expression used by nearby hardware logic.
\paragraph{Line 99}
\noindent\colorbox{gray!8}{\parbox{\dimexpr\linewidth-2\fboxsep\relax}{\raggedright\ttfamily\scriptsize             end}}
\par\smallskip
This SystemVerilog line contributes to the structural or sequential description of the accelerator. Its effect is resolved by XSim for verification and by Vivado during synthesis.
\paragraph{Line 100}
\noindent\colorbox{gray!8}{\parbox{\dimexpr\linewidth-2\fboxsep\relax}{\raggedright\ttfamily\scriptsize \emph{blank line}}}
\par\smallskip
This blank line is not hardware or software behavior. It separates nearby declarations or operations so the file remains readable and reviewable.
\paragraph{Line 101}
\noindent\colorbox{gray!8}{\parbox{\dimexpr\linewidth-2\fboxsep\relax}{\raggedright\ttfamily\scriptsize             @(negedge clk);}}
\par\smallskip
This SystemVerilog line contributes to the structural or sequential description of the accelerator. Its effect is resolved by XSim for verification and by Vivado during synthesis.
\paragraph{Line 102}
\noindent\colorbox{gray!8}{\parbox{\dimexpr\linewidth-2\fboxsep\relax}{\raggedright\ttfamily\scriptsize             train\_label = label\_mem[sample\_idx];}}
\par\smallskip
This is a local compile-time or procedural assignment. In context it either defines a helper value or connects a temporary expression used by nearby hardware logic.
\paragraph{Line 103}
\noindent\colorbox{gray!8}{\parbox{\dimexpr\linewidth-2\fboxsep\relax}{\raggedright\ttfamily\scriptsize             train\_enable = 1'b1;}}
\par\smallskip
This is a local compile-time or procedural assignment. In context it either defines a helper value or connects a temporary expression used by nearby hardware logic.
\paragraph{Line 104}
\noindent\colorbox{gray!8}{\parbox{\dimexpr\linewidth-2\fboxsep\relax}{\raggedright\ttfamily\scriptsize             start = 1'b1;}}
\par\smallskip
This is a local compile-time or procedural assignment. In context it either defines a helper value or connects a temporary expression used by nearby hardware logic.
\paragraph{Line 105}
\noindent\colorbox{gray!8}{\parbox{\dimexpr\linewidth-2\fboxsep\relax}{\raggedright\ttfamily\scriptsize             @(negedge clk);}}
\par\smallskip
This SystemVerilog line contributes to the structural or sequential description of the accelerator. Its effect is resolved by XSim for verification and by Vivado during synthesis.
\paragraph{Line 106}
\noindent\colorbox{gray!8}{\parbox{\dimexpr\linewidth-2\fboxsep\relax}{\raggedright\ttfamily\scriptsize             start = 1'b0;}}
\par\smallskip
This is a local compile-time or procedural assignment. In context it either defines a helper value or connects a temporary expression used by nearby hardware logic.
\paragraph{Line 107}
\noindent\colorbox{gray!8}{\parbox{\dimexpr\linewidth-2\fboxsep\relax}{\raggedright\ttfamily\scriptsize \emph{blank line}}}
\par\smallskip
This blank line is not hardware or software behavior. It separates nearby declarations or operations so the file remains readable and reviewable.
\paragraph{Line 108}
\noindent\colorbox{gray!8}{\parbox{\dimexpr\linewidth-2\fboxsep\relax}{\raggedright\ttfamily\scriptsize             waited = 0;}}
\par\smallskip
This is a local compile-time or procedural assignment. In context it either defines a helper value or connects a temporary expression used by nearby hardware logic.
\paragraph{Line 109}
\noindent\colorbox{gray!8}{\parbox{\dimexpr\linewidth-2\fboxsep\relax}{\raggedright\ttfamily\scriptsize             while (!feature\_ready \&\& waited < 10000) begin}}
\par\smallskip
This SystemVerilog line contributes to the structural or sequential description of the accelerator. Its effect is resolved by XSim for verification and by Vivado during synthesis.
\paragraph{Line 110}
\noindent\colorbox{gray!8}{\parbox{\dimexpr\linewidth-2\fboxsep\relax}{\raggedright\ttfamily\scriptsize                 @(posedge clk);}}
\par\smallskip
This SystemVerilog line contributes to the structural or sequential description of the accelerator. Its effect is resolved by XSim for verification and by Vivado during synthesis.
\paragraph{Line 111}
\noindent\colorbox{gray!8}{\parbox{\dimexpr\linewidth-2\fboxsep\relax}{\raggedright\ttfamily\scriptsize                 waited++;}}
\par\smallskip
This SystemVerilog line contributes to the structural or sequential description of the accelerator. Its effect is resolved by XSim for verification and by Vivado during synthesis.
\paragraph{Line 112}
\noindent\colorbox{gray!8}{\parbox{\dimexpr\linewidth-2\fboxsep\relax}{\raggedright\ttfamily\scriptsize             end}}
\par\smallskip
This SystemVerilog line contributes to the structural or sequential description of the accelerator. Its effect is resolved by XSim for verification and by Vivado during synthesis.
\paragraph{Line 113}
\noindent\colorbox{gray!8}{\parbox{\dimexpr\linewidth-2\fboxsep\relax}{\raggedright\ttfamily\scriptsize             if (!feature\_ready) begin}}
\par\smallskip
This SystemVerilog line contributes to the structural or sequential description of the accelerator. Its effect is resolved by XSim for verification and by Vivado during synthesis.
\paragraph{Line 114}
\noindent\colorbox{gray!8}{\parbox{\dimexpr\linewidth-2\fboxsep\relax}{\raggedright\ttfamily\scriptsize                 \$display("ERROR: sample \%0d never reached feature\_ready", sample\_idx);}}
\par\smallskip
This SystemVerilog line contributes to the structural or sequential description of the accelerator. Its effect is resolved by XSim for verification and by Vivado during synthesis.
\paragraph{Line 115}
\noindent\colorbox{gray!8}{\parbox{\dimexpr\linewidth-2\fboxsep\relax}{\raggedright\ttfamily\scriptsize                 errors++;}}
\par\smallskip
This SystemVerilog line contributes to the structural or sequential description of the accelerator. Its effect is resolved by XSim for verification and by Vivado during synthesis.
\paragraph{Line 116}
\noindent\colorbox{gray!8}{\parbox{\dimexpr\linewidth-2\fboxsep\relax}{\raggedright\ttfamily\scriptsize             end}}
\par\smallskip
This SystemVerilog line contributes to the structural or sequential description of the accelerator. Its effect is resolved by XSim for verification and by Vivado during synthesis.
\paragraph{Line 117}
\noindent\colorbox{gray!8}{\parbox{\dimexpr\linewidth-2\fboxsep\relax}{\raggedright\ttfamily\scriptsize \emph{blank line}}}
\par\smallskip
This blank line is not hardware or software behavior. It separates nearby declarations or operations so the file remains readable and reviewable.
\paragraph{Line 118}
\noindent\colorbox{gray!8}{\parbox{\dimexpr\linewidth-2\fboxsep\relax}{\raggedright\ttfamily\scriptsize             for (int f = 0; f < FEATURE\_DIM; f++) begin}}
\par\smallskip
This is a local compile-time or procedural assignment. In context it either defines a helper value or connects a temporary expression used by nearby hardware logic.
\paragraph{Line 119}
\noindent\colorbox{gray!8}{\parbox{\dimexpr\linewidth-2\fboxsep\relax}{\raggedright\ttfamily\scriptsize                 @(negedge clk);}}
\par\smallskip
This SystemVerilog line contributes to the structural or sequential description of the accelerator. Its effect is resolved by XSim for verification and by Vivado during synthesis.
\paragraph{Line 120}
\noindent\colorbox{gray!8}{\parbox{\dimexpr\linewidth-2\fboxsep\relax}{\raggedright\ttfamily\scriptsize                 feature\_value = feature\_mem[feature\_ptr];}}
\par\smallskip
This line touches the feature memory that holds one 384-dimensional fixed-point Spikingformer feature vector for the current image. Keeping the feature vector local avoids repeatedly streaming the same values for every class score.
\paragraph{Line 121}
\noindent\colorbox{gray!8}{\parbox{\dimexpr\linewidth-2\fboxsep\relax}{\raggedright\ttfamily\scriptsize                 feature\_ptr++;}}
\par\smallskip
This SystemVerilog line contributes to the structural or sequential description of the accelerator. Its effect is resolved by XSim for verification and by Vivado during synthesis.
\paragraph{Line 122}
\noindent\colorbox{gray!8}{\parbox{\dimexpr\linewidth-2\fboxsep\relax}{\raggedright\ttfamily\scriptsize                 feature\_valid = 1'b1;}}
\par\smallskip
This is a local compile-time or procedural assignment. In context it either defines a helper value or connects a temporary expression used by nearby hardware logic.
\paragraph{Line 123}
\noindent\colorbox{gray!8}{\parbox{\dimexpr\linewidth-2\fboxsep\relax}{\raggedright\ttfamily\scriptsize                 @(posedge clk);}}
\par\smallskip
This SystemVerilog line contributes to the structural or sequential description of the accelerator. Its effect is resolved by XSim for verification and by Vivado during synthesis.
\paragraph{Line 124}
\noindent\colorbox{gray!8}{\parbox{\dimexpr\linewidth-2\fboxsep\relax}{\raggedright\ttfamily\scriptsize                 \#1;}}
\par\smallskip
This is a human-facing comment. It documents intent for the reader and is ignored by the compiler or interpreter.
\paragraph{Line 125}
\noindent\colorbox{gray!8}{\parbox{\dimexpr\linewidth-2\fboxsep\relax}{\raggedright\ttfamily\scriptsize                 if (f < FEATURE\_DIM - 1 \&\& !feature\_ready) begin}}
\par\smallskip
This SystemVerilog line contributes to the structural or sequential description of the accelerator. Its effect is resolved by XSim for verification and by Vivado during synthesis.
\paragraph{Line 126}
\noindent\colorbox{gray!8}{\parbox{\dimexpr\linewidth-2\fboxsep\relax}{\raggedright\ttfamily\scriptsize                     \$display("ERROR: feature\_ready dropped during sample \%0d feature \%0d", sample\_idx, f);}}
\par\smallskip
This SystemVerilog line contributes to the structural or sequential description of the accelerator. Its effect is resolved by XSim for verification and by Vivado during synthesis.
\paragraph{Line 127}
\noindent\colorbox{gray!8}{\parbox{\dimexpr\linewidth-2\fboxsep\relax}{\raggedright\ttfamily\scriptsize                     errors++;}}
\par\smallskip
This SystemVerilog line contributes to the structural or sequential description of the accelerator. Its effect is resolved by XSim for verification and by Vivado during synthesis.
\paragraph{Line 128}
\noindent\colorbox{gray!8}{\parbox{\dimexpr\linewidth-2\fboxsep\relax}{\raggedright\ttfamily\scriptsize                 end}}
\par\smallskip
This SystemVerilog line contributes to the structural or sequential description of the accelerator. Its effect is resolved by XSim for verification and by Vivado during synthesis.
\paragraph{Line 129}
\noindent\colorbox{gray!8}{\parbox{\dimexpr\linewidth-2\fboxsep\relax}{\raggedright\ttfamily\scriptsize             end}}
\par\smallskip
This SystemVerilog line contributes to the structural or sequential description of the accelerator. Its effect is resolved by XSim for verification and by Vivado during synthesis.
\paragraph{Line 130}
\noindent\colorbox{gray!8}{\parbox{\dimexpr\linewidth-2\fboxsep\relax}{\raggedright\ttfamily\scriptsize \emph{blank line}}}
\par\smallskip
This blank line is not hardware or software behavior. It separates nearby declarations or operations so the file remains readable and reviewable.
\paragraph{Line 131}
\noindent\colorbox{gray!8}{\parbox{\dimexpr\linewidth-2\fboxsep\relax}{\raggedright\ttfamily\scriptsize             @(negedge clk);}}
\par\smallskip
This SystemVerilog line contributes to the structural or sequential description of the accelerator. Its effect is resolved by XSim for verification and by Vivado during synthesis.
\paragraph{Line 132}
\noindent\colorbox{gray!8}{\parbox{\dimexpr\linewidth-2\fboxsep\relax}{\raggedright\ttfamily\scriptsize             feature\_valid = 1'b0;}}
\par\smallskip
This is a local compile-time or procedural assignment. In context it either defines a helper value or connects a temporary expression used by nearby hardware logic.
\paragraph{Line 133}
\noindent\colorbox{gray!8}{\parbox{\dimexpr\linewidth-2\fboxsep\relax}{\raggedright\ttfamily\scriptsize             frame\_done = 1'b1;}}
\par\smallskip
This is a local compile-time or procedural assignment. In context it either defines a helper value or connects a temporary expression used by nearby hardware logic.
\paragraph{Line 134}
\noindent\colorbox{gray!8}{\parbox{\dimexpr\linewidth-2\fboxsep\relax}{\raggedright\ttfamily\scriptsize             @(negedge clk);}}
\par\smallskip
This SystemVerilog line contributes to the structural or sequential description of the accelerator. Its effect is resolved by XSim for verification and by Vivado during synthesis.
\paragraph{Line 135}
\noindent\colorbox{gray!8}{\parbox{\dimexpr\linewidth-2\fboxsep\relax}{\raggedright\ttfamily\scriptsize             frame\_done = 1'b0;}}
\par\smallskip
This is a local compile-time or procedural assignment. In context it either defines a helper value or connects a temporary expression used by nearby hardware logic.
\paragraph{Line 136}
\noindent\colorbox{gray!8}{\parbox{\dimexpr\linewidth-2\fboxsep\relax}{\raggedright\ttfamily\scriptsize \emph{blank line}}}
\par\smallskip
This blank line is not hardware or software behavior. It separates nearby declarations or operations so the file remains readable and reviewable.
\paragraph{Line 137}
\noindent\colorbox{gray!8}{\parbox{\dimexpr\linewidth-2\fboxsep\relax}{\raggedright\ttfamily\scriptsize             waited = 0;}}
\par\smallskip
This is a local compile-time or procedural assignment. In context it either defines a helper value or connects a temporary expression used by nearby hardware logic.
\paragraph{Line 138}
\noindent\colorbox{gray!8}{\parbox{\dimexpr\linewidth-2\fboxsep\relax}{\raggedright\ttfamily\scriptsize             update\_seen = 1'b0;}}
\par\smallskip
This is a local compile-time or procedural assignment. In context it either defines a helper value or connects a temporary expression used by nearby hardware logic.
\paragraph{Line 139}
\noindent\colorbox{gray!8}{\parbox{\dimexpr\linewidth-2\fboxsep\relax}{\raggedright\ttfamily\scriptsize             while (!result\_valid \&\& waited < 200000) begin}}
\par\smallskip
This SystemVerilog line contributes to the structural or sequential description of the accelerator. Its effect is resolved by XSim for verification and by Vivado during synthesis.
\paragraph{Line 140}
\noindent\colorbox{gray!8}{\parbox{\dimexpr\linewidth-2\fboxsep\relax}{\raggedright\ttfamily\scriptsize                 @(posedge clk);}}
\par\smallskip
This SystemVerilog line contributes to the structural or sequential description of the accelerator. Its effect is resolved by XSim for verification and by Vivado during synthesis.
\paragraph{Line 141}
\noindent\colorbox{gray!8}{\parbox{\dimexpr\linewidth-2\fboxsep\relax}{\raggedright\ttfamily\scriptsize                 if (train\_update\_applied) begin}}
\par\smallskip
This SystemVerilog line contributes to the structural or sequential description of the accelerator. Its effect is resolved by XSim for verification and by Vivado during synthesis.
\paragraph{Line 142}
\noindent\colorbox{gray!8}{\parbox{\dimexpr\linewidth-2\fboxsep\relax}{\raggedright\ttfamily\scriptsize                     update\_seen = 1'b1;}}
\par\smallskip
This is a local compile-time or procedural assignment. In context it either defines a helper value or connects a temporary expression used by nearby hardware logic.
\paragraph{Line 143}
\noindent\colorbox{gray!8}{\parbox{\dimexpr\linewidth-2\fboxsep\relax}{\raggedright\ttfamily\scriptsize                 end}}
\par\smallskip
This SystemVerilog line contributes to the structural or sequential description of the accelerator. Its effect is resolved by XSim for verification and by Vivado during synthesis.
\paragraph{Line 144}
\noindent\colorbox{gray!8}{\parbox{\dimexpr\linewidth-2\fboxsep\relax}{\raggedright\ttfamily\scriptsize                 waited++;}}
\par\smallskip
This SystemVerilog line contributes to the structural or sequential description of the accelerator. Its effect is resolved by XSim for verification and by Vivado during synthesis.
\paragraph{Line 145}
\noindent\colorbox{gray!8}{\parbox{\dimexpr\linewidth-2\fboxsep\relax}{\raggedright\ttfamily\scriptsize             end}}
\par\smallskip
This SystemVerilog line contributes to the structural or sequential description of the accelerator. Its effect is resolved by XSim for verification and by Vivado during synthesis.
\paragraph{Line 146}
\noindent\colorbox{gray!8}{\parbox{\dimexpr\linewidth-2\fboxsep\relax}{\raggedright\ttfamily\scriptsize             if (!result\_valid) begin}}
\par\smallskip
This SystemVerilog line contributes to the structural or sequential description of the accelerator. Its effect is resolved by XSim for verification and by Vivado during synthesis.
\paragraph{Line 147}
\noindent\colorbox{gray!8}{\parbox{\dimexpr\linewidth-2\fboxsep\relax}{\raggedright\ttfamily\scriptsize                 \$display("ERROR: sample \%0d no result\_valid", sample\_idx);}}
\par\smallskip
This SystemVerilog line contributes to the structural or sequential description of the accelerator. Its effect is resolved by XSim for verification and by Vivado during synthesis.
\paragraph{Line 148}
\noindent\colorbox{gray!8}{\parbox{\dimexpr\linewidth-2\fboxsep\relax}{\raggedright\ttfamily\scriptsize                 errors++;}}
\par\smallskip
This SystemVerilog line contributes to the structural or sequential description of the accelerator. Its effect is resolved by XSim for verification and by Vivado during synthesis.
\paragraph{Line 149}
\noindent\colorbox{gray!8}{\parbox{\dimexpr\linewidth-2\fboxsep\relax}{\raggedright\ttfamily\scriptsize             end else begin}}
\par\smallskip
This SystemVerilog line contributes to the structural or sequential description of the accelerator. Its effect is resolved by XSim for verification and by Vivado during synthesis.
\paragraph{Line 150}
\noindent\colorbox{gray!8}{\parbox{\dimexpr\linewidth-2\fboxsep\relax}{\raggedright\ttfamily\scriptsize                 if (predicted\_label == label\_mem[sample\_idx]) begin}}
\par\smallskip
This is a local compile-time or procedural assignment. In context it either defines a helper value or connects a temporary expression used by nearby hardware logic.
\paragraph{Line 151}
\noindent\colorbox{gray!8}{\parbox{\dimexpr\linewidth-2\fboxsep\relax}{\raggedright\ttfamily\scriptsize                     correct++;}}
\par\smallskip
This SystemVerilog line contributes to the structural or sequential description of the accelerator. Its effect is resolved by XSim for verification and by Vivado during synthesis.
\paragraph{Line 152}
\noindent\colorbox{gray!8}{\parbox{\dimexpr\linewidth-2\fboxsep\relax}{\raggedright\ttfamily\scriptsize                 end}}
\par\smallskip
This SystemVerilog line contributes to the structural or sequential description of the accelerator. Its effect is resolved by XSim for verification and by Vivado during synthesis.
\paragraph{Line 153}
\noindent\colorbox{gray!8}{\parbox{\dimexpr\linewidth-2\fboxsep\relax}{\raggedright\ttfamily\scriptsize                 if (update\_seen) begin}}
\par\smallskip
This SystemVerilog line contributes to the structural or sequential description of the accelerator. Its effect is resolved by XSim for verification and by Vivado during synthesis.
\paragraph{Line 154}
\noindent\colorbox{gray!8}{\parbox{\dimexpr\linewidth-2\fboxsep\relax}{\raggedright\ttfamily\scriptsize                     updates++;}}
\par\smallskip
This SystemVerilog line contributes to the structural or sequential description of the accelerator. Its effect is resolved by XSim for verification and by Vivado during synthesis.
\paragraph{Line 155}
\noindent\colorbox{gray!8}{\parbox{\dimexpr\linewidth-2\fboxsep\relax}{\raggedright\ttfamily\scriptsize                 end}}
\par\smallskip
This SystemVerilog line contributes to the structural or sequential description of the accelerator. Its effect is resolved by XSim for verification and by Vivado during synthesis.
\paragraph{Line 156}
\noindent\colorbox{gray!8}{\parbox{\dimexpr\linewidth-2\fboxsep\relax}{\raggedright\ttfamily\scriptsize                 \$fwrite(result\_fh, "\%0d,\%0d,\%0d,\%0d,\%0d,\%0d,\%0d,\%0d,\%0d",}}
\par\smallskip
This SystemVerilog line contributes to the structural or sequential description of the accelerator. Its effect is resolved by XSim for verification and by Vivado during synthesis.
\paragraph{Line 157}
\noindent\colorbox{gray!8}{\parbox{\dimexpr\linewidth-2\fboxsep\relax}{\raggedright\ttfamily\scriptsize                         sample\_idx, label\_mem[sample\_idx], predicted\_label, predicted\_score,}}
\par\smallskip
This SystemVerilog line contributes to the structural or sequential description of the accelerator. Its effect is resolved by XSim for verification and by Vivado during synthesis.
\paragraph{Line 158}
\noindent\colorbox{gray!8}{\parbox{\dimexpr\linewidth-2\fboxsep\relax}{\raggedright\ttfamily\scriptsize                         predicted\_label == label\_mem[sample\_idx], update\_seen, debug\_update\_count,}}
\par\smallskip
This is a local compile-time or procedural assignment. In context it either defines a helper value or connects a temporary expression used by nearby hardware logic.
\paragraph{Line 159}
\noindent\colorbox{gray!8}{\parbox{\dimexpr\linewidth-2\fboxsep\relax}{\raggedright\ttfamily\scriptsize                         debug\_infer\_skip\_count, debug\_train\_skip\_count);}}
\par\smallskip
This line exposes or updates a hardware activity counter. The counter makes the feature gate visible in waveforms and CSV traces, so the power-saving opportunity is observable from RTL rather than only from Python statistics.
\paragraph{Line 160}
\noindent\colorbox{gray!8}{\parbox{\dimexpr\linewidth-2\fboxsep\relax}{\raggedright\ttfamily\scriptsize                 for (int c = 0; c < 10; c++) begin}}
\par\smallskip
This is a local compile-time or procedural assignment. In context it either defines a helper value or connects a temporary expression used by nearby hardware logic.
\paragraph{Line 161}
\noindent\colorbox{gray!8}{\parbox{\dimexpr\linewidth-2\fboxsep\relax}{\raggedright\ttfamily\scriptsize                     \$fwrite(result\_fh, ",\%0d", \$signed(dut.class\_score\_mem[c]));}}
\par\smallskip
This stores or reads the ten class scores for the current sample. Those scores are needed for argmax prediction, CSV testbench reporting, and the margin test used before an online update.
\paragraph{Line 162}
\noindent\colorbox{gray!8}{\parbox{\dimexpr\linewidth-2\fboxsep\relax}{\raggedright\ttfamily\scriptsize                 end}}
\par\smallskip
This SystemVerilog line contributes to the structural or sequential description of the accelerator. Its effect is resolved by XSim for verification and by Vivado during synthesis.
\paragraph{Line 163}
\noindent\colorbox{gray!8}{\parbox{\dimexpr\linewidth-2\fboxsep\relax}{\raggedright\ttfamily\scriptsize                 \$fwrite(result\_fh, "\textbackslash{}n");}}
\par\smallskip
This SystemVerilog line contributes to the structural or sequential description of the accelerator. Its effect is resolved by XSim for verification and by Vivado during synthesis.
\paragraph{Line 164}
\noindent\colorbox{gray!8}{\parbox{\dimexpr\linewidth-2\fboxsep\relax}{\raggedright\ttfamily\scriptsize                 if (sample\_idx < TRACE\_SAMPLES || (sample\_idx \% 1000) == 0 || sample\_idx == NUM\_SAMPLES - 1) begin}}
\par\smallskip
This is a local compile-time or procedural assignment. In context it either defines a helper value or connects a temporary expression used by nearby hardware logic.
\paragraph{Line 165}
\noindent\colorbox{gray!8}{\parbox{\dimexpr\linewidth-2\fboxsep\relax}{\raggedright\ttfamily\scriptsize                     \$display("learn\_sample=\%0d label=\%0d pred=\%0d correct=\%0d update=\%0d",}}
\par\smallskip
This is a local compile-time or procedural assignment. In context it either defines a helper value or connects a temporary expression used by nearby hardware logic.
\paragraph{Line 166}
\noindent\colorbox{gray!8}{\parbox{\dimexpr\linewidth-2\fboxsep\relax}{\raggedright\ttfamily\scriptsize                              sample\_idx, label\_mem[sample\_idx], predicted\_label,}}
\par\smallskip
This SystemVerilog line contributes to the structural or sequential description of the accelerator. Its effect is resolved by XSim for verification and by Vivado during synthesis.
\paragraph{Line 167}
\noindent\colorbox{gray!8}{\parbox{\dimexpr\linewidth-2\fboxsep\relax}{\raggedright\ttfamily\scriptsize                              predicted\_label == label\_mem[sample\_idx], update\_seen);}}
\par\smallskip
This is a local compile-time or procedural assignment. In context it either defines a helper value or connects a temporary expression used by nearby hardware logic.
\paragraph{Line 168}
\noindent\colorbox{gray!8}{\parbox{\dimexpr\linewidth-2\fboxsep\relax}{\raggedright\ttfamily\scriptsize                 end}}
\par\smallskip
This SystemVerilog line contributes to the structural or sequential description of the accelerator. Its effect is resolved by XSim for verification and by Vivado during synthesis.
\paragraph{Line 169}
\noindent\colorbox{gray!8}{\parbox{\dimexpr\linewidth-2\fboxsep\relax}{\raggedright\ttfamily\scriptsize             end}}
\par\smallskip
This SystemVerilog line contributes to the structural or sequential description of the accelerator. Its effect is resolved by XSim for verification and by Vivado during synthesis.
\paragraph{Line 170}
\noindent\colorbox{gray!8}{\parbox{\dimexpr\linewidth-2\fboxsep\relax}{\raggedright\ttfamily\scriptsize \emph{blank line}}}
\par\smallskip
This blank line is not hardware or software behavior. It separates nearby declarations or operations so the file remains readable and reviewable.
\paragraph{Line 171}
\noindent\colorbox{gray!8}{\parbox{\dimexpr\linewidth-2\fboxsep\relax}{\raggedright\ttfamily\scriptsize             waited = 0;}}
\par\smallskip
This is a local compile-time or procedural assignment. In context it either defines a helper value or connects a temporary expression used by nearby hardware logic.
\paragraph{Line 172}
\noindent\colorbox{gray!8}{\parbox{\dimexpr\linewidth-2\fboxsep\relax}{\raggedright\ttfamily\scriptsize             while (busy \&\& waited < 1000) begin}}
\par\smallskip
This SystemVerilog line contributes to the structural or sequential description of the accelerator. Its effect is resolved by XSim for verification and by Vivado during synthesis.
\paragraph{Line 173}
\noindent\colorbox{gray!8}{\parbox{\dimexpr\linewidth-2\fboxsep\relax}{\raggedright\ttfamily\scriptsize                 @(posedge clk);}}
\par\smallskip
This SystemVerilog line contributes to the structural or sequential description of the accelerator. Its effect is resolved by XSim for verification and by Vivado during synthesis.
\paragraph{Line 174}
\noindent\colorbox{gray!8}{\parbox{\dimexpr\linewidth-2\fboxsep\relax}{\raggedright\ttfamily\scriptsize                 waited++;}}
\par\smallskip
This SystemVerilog line contributes to the structural or sequential description of the accelerator. Its effect is resolved by XSim for verification and by Vivado during synthesis.
\paragraph{Line 175}
\noindent\colorbox{gray!8}{\parbox{\dimexpr\linewidth-2\fboxsep\relax}{\raggedright\ttfamily\scriptsize             end}}
\par\smallskip
This SystemVerilog line contributes to the structural or sequential description of the accelerator. Its effect is resolved by XSim for verification and by Vivado during synthesis.
\paragraph{Line 176}
\noindent\colorbox{gray!8}{\parbox{\dimexpr\linewidth-2\fboxsep\relax}{\raggedright\ttfamily\scriptsize             train\_enable = 1'b0;}}
\par\smallskip
This is a local compile-time or procedural assignment. In context it either defines a helper value or connects a temporary expression used by nearby hardware logic.
\paragraph{Line 177}
\noindent\colorbox{gray!8}{\parbox{\dimexpr\linewidth-2\fboxsep\relax}{\raggedright\ttfamily\scriptsize             trace\_enable = 1'b0;}}
\par\smallskip
This is a local compile-time or procedural assignment. In context it either defines a helper value or connects a temporary expression used by nearby hardware logic.
\paragraph{Line 178}
\noindent\colorbox{gray!8}{\parbox{\dimexpr\linewidth-2\fboxsep\relax}{\raggedright\ttfamily\scriptsize         end}}
\par\smallskip
This SystemVerilog line contributes to the structural or sequential description of the accelerator. Its effect is resolved by XSim for verification and by Vivado during synthesis.
\paragraph{Line 179}
\noindent\colorbox{gray!8}{\parbox{\dimexpr\linewidth-2\fboxsep\relax}{\raggedright\ttfamily\scriptsize     endtask}}
\par\smallskip
This SystemVerilog line contributes to the structural or sequential description of the accelerator. Its effect is resolved by XSim for verification and by Vivado during synthesis.
\paragraph{Line 180}
\noindent\colorbox{gray!8}{\parbox{\dimexpr\linewidth-2\fboxsep\relax}{\raggedright\ttfamily\scriptsize \emph{blank line}}}
\par\smallskip
This blank line is not hardware or software behavior. It separates nearby declarations or operations so the file remains readable and reviewable.
\paragraph{Line 181}
\noindent\colorbox{gray!8}{\parbox{\dimexpr\linewidth-2\fboxsep\relax}{\raggedright\ttfamily\scriptsize     always @(posedge clk) begin}}
\par\smallskip
This SystemVerilog line contributes to the structural or sequential description of the accelerator. Its effect is resolved by XSim for verification and by Vivado during synthesis.
\paragraph{Line 182}
\noindent\colorbox{gray!8}{\parbox{\dimexpr\linewidth-2\fboxsep\relax}{\raggedright\ttfamily\scriptsize         if (trace\_fh != 0 \&\& trace\_enable) begin}}
\par\smallskip
This is a local compile-time or procedural assignment. In context it either defines a helper value or connects a temporary expression used by nearby hardware logic.
\paragraph{Line 183}
\noindent\colorbox{gray!8}{\parbox{\dimexpr\linewidth-2\fboxsep\relax}{\raggedright\ttfamily\scriptsize             \$fwrite(trace\_fh, "\%0t,\%0d,\%0d,\%0d,\%0d,\%0d,\%0d,\%0d,\%0d,\%0d,\%0d,\%0d,\%0d,\%0d,\%0d\textbackslash{}n",}}
\par\smallskip
This SystemVerilog line contributes to the structural or sequential description of the accelerator. Its effect is resolved by XSim for verification and by Vivado during synthesis.
\paragraph{Line 184}
\noindent\colorbox{gray!8}{\parbox{\dimexpr\linewidth-2\fboxsep\relax}{\raggedright\ttfamily\scriptsize                     \$time, cycle, rst\_n, debug\_state, start, feature\_valid, feature\_ready,}}
\par\smallskip
This SystemVerilog line contributes to the structural or sequential description of the accelerator. Its effect is resolved by XSim for verification and by Vivado during synthesis.
\paragraph{Line 185}
\noindent\colorbox{gray!8}{\parbox{\dimexpr\linewidth-2\fboxsep\relax}{\raggedright\ttfamily\scriptsize                     debug\_feature, debug\_class, result\_valid, predicted\_label,}}
\par\smallskip
This SystemVerilog line contributes to the structural or sequential description of the accelerator. Its effect is resolved by XSim for verification and by Vivado during synthesis.
\paragraph{Line 186}
\noindent\colorbox{gray!8}{\parbox{\dimexpr\linewidth-2\fboxsep\relax}{\raggedright\ttfamily\scriptsize                     train\_update\_applied, debug\_update\_count, debug\_infer\_skip\_count, debug\_train\_skip\_count);}}
\par\smallskip
This line exposes or updates a hardware activity counter. The counter makes the feature gate visible in waveforms and CSV traces, so the power-saving opportunity is observable from RTL rather than only from Python statistics.
\paragraph{Line 187}
\noindent\colorbox{gray!8}{\parbox{\dimexpr\linewidth-2\fboxsep\relax}{\raggedright\ttfamily\scriptsize         end}}
\par\smallskip
This SystemVerilog line contributes to the structural or sequential description of the accelerator. Its effect is resolved by XSim for verification and by Vivado during synthesis.
\paragraph{Line 188}
\noindent\colorbox{gray!8}{\parbox{\dimexpr\linewidth-2\fboxsep\relax}{\raggedright\ttfamily\scriptsize         cycle++;}}
\par\smallskip
This SystemVerilog line contributes to the structural or sequential description of the accelerator. Its effect is resolved by XSim for verification and by Vivado during synthesis.
\paragraph{Line 189}
\noindent\colorbox{gray!8}{\parbox{\dimexpr\linewidth-2\fboxsep\relax}{\raggedright\ttfamily\scriptsize     end}}
\par\smallskip
This SystemVerilog line contributes to the structural or sequential description of the accelerator. Its effect is resolved by XSim for verification and by Vivado during synthesis.
\paragraph{Line 190}
\noindent\colorbox{gray!8}{\parbox{\dimexpr\linewidth-2\fboxsep\relax}{\raggedright\ttfamily\scriptsize \emph{blank line}}}
\par\smallskip
This blank line is not hardware or software behavior. It separates nearby declarations or operations so the file remains readable and reviewable.
\paragraph{Line 191}
\noindent\colorbox{gray!8}{\parbox{\dimexpr\linewidth-2\fboxsep\relax}{\raggedright\ttfamily\scriptsize     initial begin}}
\par\smallskip
This simulation/synthesis initialization block loads trained memories from generated .mem files. On FPGA, Vivado maps these initial values into BRAM/LUTRAM initialization data.
\paragraph{Line 192}
\noindent\colorbox{gray!8}{\parbox{\dimexpr\linewidth-2\fboxsep\relax}{\raggedright\ttfamily\scriptsize         if (NUM\_SAMPLES > MAX\_SAMPLES) begin}}
\par\smallskip
This SystemVerilog line contributes to the structural or sequential description of the accelerator. Its effect is resolved by XSim for verification and by Vivado during synthesis.
\paragraph{Line 193}
\noindent\colorbox{gray!8}{\parbox{\dimexpr\linewidth-2\fboxsep\relax}{\raggedright\ttfamily\scriptsize             \$display("ERROR: NUM\_SAMPLES=\%0d exceeds MAX\_SAMPLES=\%0d", NUM\_SAMPLES, MAX\_SAMPLES);}}
\par\smallskip
This is a local compile-time or procedural assignment. In context it either defines a helper value or connects a temporary expression used by nearby hardware logic.
\paragraph{Line 194}
\noindent\colorbox{gray!8}{\parbox{\dimexpr\linewidth-2\fboxsep\relax}{\raggedright\ttfamily\scriptsize             \$fatal;}}
\par\smallskip
This SystemVerilog line contributes to the structural or sequential description of the accelerator. Its effect is resolved by XSim for verification and by Vivado during synthesis.
\paragraph{Line 195}
\noindent\colorbox{gray!8}{\parbox{\dimexpr\linewidth-2\fboxsep\relax}{\raggedright\ttfamily\scriptsize         end}}
\par\smallskip
This SystemVerilog line contributes to the structural or sequential description of the accelerator. Its effect is resolved by XSim for verification and by Vivado during synthesis.
\paragraph{Line 196}
\noindent\colorbox{gray!8}{\parbox{\dimexpr\linewidth-2\fboxsep\relax}{\raggedright\ttfamily\scriptsize         require\_file("rtl/cifar10\_spikingformer\_head\_features.mem");}}
\par\smallskip
This SystemVerilog line contributes to the structural or sequential description of the accelerator. Its effect is resolved by XSim for verification and by Vivado during synthesis.
\paragraph{Line 197}
\noindent\colorbox{gray!8}{\parbox{\dimexpr\linewidth-2\fboxsep\relax}{\raggedright\ttfamily\scriptsize         require\_file("rtl/cifar10\_spikingformer\_head\_weights.mem");}}
\par\smallskip
This SystemVerilog line contributes to the structural or sequential description of the accelerator. Its effect is resolved by XSim for verification and by Vivado during synthesis.
\paragraph{Line 198}
\noindent\colorbox{gray!8}{\parbox{\dimexpr\linewidth-2\fboxsep\relax}{\raggedright\ttfamily\scriptsize         require\_file("rtl/cifar10\_spikingformer\_head\_bias.mem");}}
\par\smallskip
This SystemVerilog line contributes to the structural or sequential description of the accelerator. Its effect is resolved by XSim for verification and by Vivado during synthesis.
\paragraph{Line 199}
\noindent\colorbox{gray!8}{\parbox{\dimexpr\linewidth-2\fboxsep\relax}{\raggedright\ttfamily\scriptsize         require\_file("rtl/cifar10\_spikingformer\_head\_labels.mem");}}
\par\smallskip
This SystemVerilog line contributes to the structural or sequential description of the accelerator. Its effect is resolved by XSim for verification and by Vivado during synthesis.
\paragraph{Line 200}
\noindent\colorbox{gray!8}{\parbox{\dimexpr\linewidth-2\fboxsep\relax}{\raggedright\ttfamily\scriptsize         \$readmemh("rtl/cifar10\_spikingformer\_head\_features.mem", feature\_mem);}}
\par\smallskip
This line touches the feature memory that holds one 384-dimensional fixed-point Spikingformer feature vector for the current image. Keeping the feature vector local avoids repeatedly streaming the same values for every class score.
\paragraph{Line 201}
\noindent\colorbox{gray!8}{\parbox{\dimexpr\linewidth-2\fboxsep\relax}{\raggedright\ttfamily\scriptsize         \$readmemh("rtl/cifar10\_spikingformer\_head\_labels.mem", label\_mem);}}
\par\smallskip
This reads a hexadecimal memory file into an RTL array. It is how exported trained weights and biases become deterministic FPGA memory contents.
\paragraph{Line 202}
\noindent\colorbox{gray!8}{\parbox{\dimexpr\linewidth-2\fboxsep\relax}{\raggedright\ttfamily\scriptsize \emph{blank line}}}
\par\smallskip
This blank line is not hardware or software behavior. It separates nearby declarations or operations so the file remains readable and reviewable.
\paragraph{Line 203}
\noindent\colorbox{gray!8}{\parbox{\dimexpr\linewidth-2\fboxsep\relax}{\raggedright\ttfamily\scriptsize         \$dumpfile("cifar10\_spikingformer\_head\_onchip\_learning.vcd");}}
\par\smallskip
This SystemVerilog line contributes to the structural or sequential description of the accelerator. Its effect is resolved by XSim for verification and by Vivado during synthesis.
\paragraph{Line 204}
\noindent\colorbox{gray!8}{\parbox{\dimexpr\linewidth-2\fboxsep\relax}{\raggedright\ttfamily\scriptsize         \$dumpvars(0, tb\_cifar10\_spikingformer\_head\_onchip\_learning);}}
\par\smallskip
This SystemVerilog line contributes to the structural or sequential description of the accelerator. Its effect is resolved by XSim for verification and by Vivado during synthesis.
\paragraph{Line 205}
\noindent\colorbox{gray!8}{\parbox{\dimexpr\linewidth-2\fboxsep\relax}{\raggedright\ttfamily\scriptsize         trace\_fh = \$fopen("cifar10\_spikingformer\_head\_onchip\_learning\_trace.csv", "w");}}
\par\smallskip
This is a local compile-time or procedural assignment. In context it either defines a helper value or connects a temporary expression used by nearby hardware logic.
\paragraph{Line 206}
\noindent\colorbox{gray!8}{\parbox{\dimexpr\linewidth-2\fboxsep\relax}{\raggedright\ttfamily\scriptsize         result\_fh = \$fopen("cifar10\_spikingformer\_head\_onchip\_learning\_results.csv", "w");}}
\par\smallskip
This is a local compile-time or procedural assignment. In context it either defines a helper value or connects a temporary expression used by nearby hardware logic.
\paragraph{Line 207}
\noindent\colorbox{gray!8}{\parbox{\dimexpr\linewidth-2\fboxsep\relax}{\raggedright\ttfamily\scriptsize         \$fwrite(trace\_fh, "time\_ps,cycle,rst\_n,state,start,feature\_valid,feature\_ready,feature\_idx,class\_idx,result\_valid,predicted\_label,train\_update\_applied,update\_count,infer\_skip\_count,train\_skip\_count\textbackslash{}n");}}
\par\smallskip
This line exposes or updates a hardware activity counter. The counter makes the feature gate visible in waveforms and CSV traces, so the power-saving opportunity is observable from RTL rather than only from Python statistics.
\paragraph{Line 208}
\noindent\colorbox{gray!8}{\parbox{\dimexpr\linewidth-2\fboxsep\relax}{\raggedright\ttfamily\scriptsize         \$fwrite(result\_fh, "sample,label,predicted,score,correct,update\_applied,update\_count,infer\_skip\_count,train\_skip\_count");}}
\par\smallskip
This line exposes or updates a hardware activity counter. The counter makes the feature gate visible in waveforms and CSV traces, so the power-saving opportunity is observable from RTL rather than only from Python statistics.
\paragraph{Line 209}
\noindent\colorbox{gray!8}{\parbox{\dimexpr\linewidth-2\fboxsep\relax}{\raggedright\ttfamily\scriptsize         for (int c = 0; c < 10; c++) begin}}
\par\smallskip
This is a local compile-time or procedural assignment. In context it either defines a helper value or connects a temporary expression used by nearby hardware logic.
\paragraph{Line 210}
\noindent\colorbox{gray!8}{\parbox{\dimexpr\linewidth-2\fboxsep\relax}{\raggedright\ttfamily\scriptsize             \$fwrite(result\_fh, ",score\%0d", c);}}
\par\smallskip
This SystemVerilog line contributes to the structural or sequential description of the accelerator. Its effect is resolved by XSim for verification and by Vivado during synthesis.
\paragraph{Line 211}
\noindent\colorbox{gray!8}{\parbox{\dimexpr\linewidth-2\fboxsep\relax}{\raggedright\ttfamily\scriptsize         end}}
\par\smallskip
This SystemVerilog line contributes to the structural or sequential description of the accelerator. Its effect is resolved by XSim for verification and by Vivado during synthesis.
\paragraph{Line 212}
\noindent\colorbox{gray!8}{\parbox{\dimexpr\linewidth-2\fboxsep\relax}{\raggedright\ttfamily\scriptsize         \$fwrite(result\_fh, "\textbackslash{}n");}}
\par\smallskip
This SystemVerilog line contributes to the structural or sequential description of the accelerator. Its effect is resolved by XSim for verification and by Vivado during synthesis.
\paragraph{Line 213}
\noindent\colorbox{gray!8}{\parbox{\dimexpr\linewidth-2\fboxsep\relax}{\raggedright\ttfamily\scriptsize \emph{blank line}}}
\par\smallskip
This blank line is not hardware or software behavior. It separates nearby declarations or operations so the file remains readable and reviewable.
\paragraph{Line 214}
\noindent\colorbox{gray!8}{\parbox{\dimexpr\linewidth-2\fboxsep\relax}{\raggedright\ttfamily\scriptsize         start = 1'b0;}}
\par\smallskip
This is a local compile-time or procedural assignment. In context it either defines a helper value or connects a temporary expression used by nearby hardware logic.
\paragraph{Line 215}
\noindent\colorbox{gray!8}{\parbox{\dimexpr\linewidth-2\fboxsep\relax}{\raggedright\ttfamily\scriptsize         feature\_valid = 1'b0;}}
\par\smallskip
This is a local compile-time or procedural assignment. In context it either defines a helper value or connects a temporary expression used by nearby hardware logic.
\paragraph{Line 216}
\noindent\colorbox{gray!8}{\parbox{\dimexpr\linewidth-2\fboxsep\relax}{\raggedright\ttfamily\scriptsize         feature\_value = '0;}}
\par\smallskip
This is a local compile-time or procedural assignment. In context it either defines a helper value or connects a temporary expression used by nearby hardware logic.
\paragraph{Line 217}
\noindent\colorbox{gray!8}{\parbox{\dimexpr\linewidth-2\fboxsep\relax}{\raggedright\ttfamily\scriptsize         frame\_done = 1'b0;}}
\par\smallskip
This is a local compile-time or procedural assignment. In context it either defines a helper value or connects a temporary expression used by nearby hardware logic.
\paragraph{Line 218}
\noindent\colorbox{gray!8}{\parbox{\dimexpr\linewidth-2\fboxsep\relax}{\raggedright\ttfamily\scriptsize         train\_enable = 1'b0;}}
\par\smallskip
This is a local compile-time or procedural assignment. In context it either defines a helper value or connects a temporary expression used by nearby hardware logic.
\paragraph{Line 219}
\noindent\colorbox{gray!8}{\parbox{\dimexpr\linewidth-2\fboxsep\relax}{\raggedright\ttfamily\scriptsize         train\_label = 4'd0;}}
\par\smallskip
This is a local compile-time or procedural assignment. In context it either defines a helper value or connects a temporary expression used by nearby hardware logic.
\paragraph{Line 220}
\noindent\colorbox{gray!8}{\parbox{\dimexpr\linewidth-2\fboxsep\relax}{\raggedright\ttfamily\scriptsize \emph{blank line}}}
\par\smallskip
This blank line is not hardware or software behavior. It separates nearby declarations or operations so the file remains readable and reviewable.
\paragraph{Line 221}
\noindent\colorbox{gray!8}{\parbox{\dimexpr\linewidth-2\fboxsep\relax}{\raggedright\ttfamily\scriptsize         repeat (8) @(negedge clk);}}
\par\smallskip
This SystemVerilog line contributes to the structural or sequential description of the accelerator. Its effect is resolved by XSim for verification and by Vivado during synthesis.
\paragraph{Line 222}
\noindent\colorbox{gray!8}{\parbox{\dimexpr\linewidth-2\fboxsep\relax}{\raggedright\ttfamily\scriptsize         rst\_n = 1'b1;}}
\par\smallskip
This is a local compile-time or procedural assignment. In context it either defines a helper value or connects a temporary expression used by nearby hardware logic.
\paragraph{Line 223}
\noindent\colorbox{gray!8}{\parbox{\dimexpr\linewidth-2\fboxsep\relax}{\raggedright\ttfamily\scriptsize         repeat (4) @(negedge clk);}}
\par\smallskip
This SystemVerilog line contributes to the structural or sequential description of the accelerator. Its effect is resolved by XSim for verification and by Vivado during synthesis.
\paragraph{Line 224}
\noindent\colorbox{gray!8}{\parbox{\dimexpr\linewidth-2\fboxsep\relax}{\raggedright\ttfamily\scriptsize \emph{blank line}}}
\par\smallskip
This blank line is not hardware or software behavior. It separates nearby declarations or operations so the file remains readable and reviewable.
\paragraph{Line 225}
\noindent\colorbox{gray!8}{\parbox{\dimexpr\linewidth-2\fboxsep\relax}{\raggedright\ttfamily\scriptsize         for (int s = 0; s < NUM\_SAMPLES; s++) begin}}
\par\smallskip
This is a local compile-time or procedural assignment. In context it either defines a helper value or connects a temporary expression used by nearby hardware logic.
\paragraph{Line 226}
\noindent\colorbox{gray!8}{\parbox{\dimexpr\linewidth-2\fboxsep\relax}{\raggedright\ttfamily\scriptsize             send\_sample(s);}}
\par\smallskip
This SystemVerilog line contributes to the structural or sequential description of the accelerator. Its effect is resolved by XSim for verification and by Vivado during synthesis.
\paragraph{Line 227}
\noindent\colorbox{gray!8}{\parbox{\dimexpr\linewidth-2\fboxsep\relax}{\raggedright\ttfamily\scriptsize         end}}
\par\smallskip
This SystemVerilog line contributes to the structural or sequential description of the accelerator. Its effect is resolved by XSim for verification and by Vivado during synthesis.
\paragraph{Line 228}
\noindent\colorbox{gray!8}{\parbox{\dimexpr\linewidth-2\fboxsep\relax}{\raggedright\ttfamily\scriptsize \emph{blank line}}}
\par\smallskip
This blank line is not hardware or software behavior. It separates nearby declarations or operations so the file remains readable and reviewable.
\paragraph{Line 229}
\noindent\colorbox{gray!8}{\parbox{\dimexpr\linewidth-2\fboxsep\relax}{\raggedright\ttfamily\scriptsize         repeat (10) @(posedge clk);}}
\par\smallskip
This SystemVerilog line contributes to the structural or sequential description of the accelerator. Its effect is resolved by XSim for verification and by Vivado during synthesis.
\paragraph{Line 230}
\noindent\colorbox{gray!8}{\parbox{\dimexpr\linewidth-2\fboxsep\relax}{\raggedright\ttfamily\scriptsize         weight\_fh = \$fopen("cifar10\_spikingformer\_head\_onchip\_learning\_weights\_after.mem", "w");}}
\par\smallskip
This is a local compile-time or procedural assignment. In context it either defines a helper value or connects a temporary expression used by nearby hardware logic.
\paragraph{Line 231}
\noindent\colorbox{gray!8}{\parbox{\dimexpr\linewidth-2\fboxsep\relax}{\raggedright\ttfamily\scriptsize         bias\_fh = \$fopen("cifar10\_spikingformer\_head\_onchip\_learning\_bias\_after.mem", "w");}}
\par\smallskip
This is a local compile-time or procedural assignment. In context it either defines a helper value or connects a temporary expression used by nearby hardware logic.
\paragraph{Line 232}
\noindent\colorbox{gray!8}{\parbox{\dimexpr\linewidth-2\fboxsep\relax}{\raggedright\ttfamily\scriptsize         for (int c = 0; c < 10; c++) begin}}
\par\smallskip
This is a local compile-time or procedural assignment. In context it either defines a helper value or connects a temporary expression used by nearby hardware logic.
\paragraph{Line 233}
\noindent\colorbox{gray!8}{\parbox{\dimexpr\linewidth-2\fboxsep\relax}{\raggedright\ttfamily\scriptsize             for (int f = 0; f < FEATURE\_DIM; f++) begin}}
\par\smallskip
This is a local compile-time or procedural assignment. In context it either defines a helper value or connects a temporary expression used by nearby hardware logic.
\paragraph{Line 234}
\noindent\colorbox{gray!8}{\parbox{\dimexpr\linewidth-2\fboxsep\relax}{\raggedright\ttfamily\scriptsize                 \$fwrite(weight\_fh, "\%0d\textbackslash{}n", \$signed(dut.weight\_mem[c * FEATURE\_DIM + f]));}}
\par\smallskip
This line touches the trained classifier weight store. In the final RTL the large 10 x 384 int16 matrix is inferred as true dual-port BRAM, so weights remain on chip and can also be updated by the learning state machine.
\paragraph{Line 235}
\noindent\colorbox{gray!8}{\parbox{\dimexpr\linewidth-2\fboxsep\relax}{\raggedright\ttfamily\scriptsize             end}}
\par\smallskip
This SystemVerilog line contributes to the structural or sequential description of the accelerator. Its effect is resolved by XSim for verification and by Vivado during synthesis.
\paragraph{Line 236}
\noindent\colorbox{gray!8}{\parbox{\dimexpr\linewidth-2\fboxsep\relax}{\raggedright\ttfamily\scriptsize             \$fwrite(bias\_fh, "\%0d\textbackslash{}n", \$signed(dut.bias\_mem[c]));}}
\par\smallskip
This line touches the class-bias store. Bias values are part of the affine classifier score and are updated on chip in separate one-write states to keep synthesis clean.
\paragraph{Line 237}
\noindent\colorbox{gray!8}{\parbox{\dimexpr\linewidth-2\fboxsep\relax}{\raggedright\ttfamily\scriptsize         end}}
\par\smallskip
This SystemVerilog line contributes to the structural or sequential description of the accelerator. Its effect is resolved by XSim for verification and by Vivado during synthesis.
\paragraph{Line 238}
\noindent\colorbox{gray!8}{\parbox{\dimexpr\linewidth-2\fboxsep\relax}{\raggedright\ttfamily\scriptsize         \$fclose(weight\_fh);}}
\par\smallskip
This SystemVerilog line contributes to the structural or sequential description of the accelerator. Its effect is resolved by XSim for verification and by Vivado during synthesis.
\paragraph{Line 239}
\noindent\colorbox{gray!8}{\parbox{\dimexpr\linewidth-2\fboxsep\relax}{\raggedright\ttfamily\scriptsize         \$fclose(bias\_fh);}}
\par\smallskip
This SystemVerilog line contributes to the structural or sequential description of the accelerator. Its effect is resolved by XSim for verification and by Vivado during synthesis.
\paragraph{Line 240}
\noindent\colorbox{gray!8}{\parbox{\dimexpr\linewidth-2\fboxsep\relax}{\raggedright\ttfamily\scriptsize         \$fclose(trace\_fh);}}
\par\smallskip
This SystemVerilog line contributes to the structural or sequential description of the accelerator. Its effect is resolved by XSim for verification and by Vivado during synthesis.
\paragraph{Line 241}
\noindent\colorbox{gray!8}{\parbox{\dimexpr\linewidth-2\fboxsep\relax}{\raggedright\ttfamily\scriptsize         \$fclose(result\_fh);}}
\par\smallskip
This SystemVerilog line contributes to the structural or sequential description of the accelerator. Its effect is resolved by XSim for verification and by Vivado during synthesis.
\paragraph{Line 242}
\noindent\colorbox{gray!8}{\parbox{\dimexpr\linewidth-2\fboxsep\relax}{\raggedright\ttfamily\scriptsize         \$display("tb\_cifar10\_spikingformer\_head\_onchip\_learning samples=\%0d correct=\%0d updates=\%0d errors=\%0d", NUM\_SAMPLES, correct, updates, errors);}}
\par\smallskip
This is a local compile-time or procedural assignment. In context it either defines a helper value or connects a temporary expression used by nearby hardware logic.
\paragraph{Line 243}
\noindent\colorbox{gray!8}{\parbox{\dimexpr\linewidth-2\fboxsep\relax}{\raggedright\ttfamily\scriptsize         if (errors != 0) begin}}
\par\smallskip
This is a local compile-time or procedural assignment. In context it either defines a helper value or connects a temporary expression used by nearby hardware logic.
\paragraph{Line 244}
\noindent\colorbox{gray!8}{\parbox{\dimexpr\linewidth-2\fboxsep\relax}{\raggedright\ttfamily\scriptsize             \$fatal;}}
\par\smallskip
This SystemVerilog line contributes to the structural or sequential description of the accelerator. Its effect is resolved by XSim for verification and by Vivado during synthesis.
\paragraph{Line 245}
\noindent\colorbox{gray!8}{\parbox{\dimexpr\linewidth-2\fboxsep\relax}{\raggedright\ttfamily\scriptsize         end}}
\par\smallskip
This SystemVerilog line contributes to the structural or sequential description of the accelerator. Its effect is resolved by XSim for verification and by Vivado during synthesis.
\paragraph{Line 246}
\noindent\colorbox{gray!8}{\parbox{\dimexpr\linewidth-2\fboxsep\relax}{\raggedright\ttfamily\scriptsize         \$finish;}}
\par\smallskip
This SystemVerilog line contributes to the structural or sequential description of the accelerator. Its effect is resolved by XSim for verification and by Vivado during synthesis.
\paragraph{Line 247}
\noindent\colorbox{gray!8}{\parbox{\dimexpr\linewidth-2\fboxsep\relax}{\raggedright\ttfamily\scriptsize     end}}
\par\smallskip
This SystemVerilog line contributes to the structural or sequential description of the accelerator. Its effect is resolved by XSim for verification and by Vivado during synthesis.
\paragraph{Line 248}
\noindent\colorbox{gray!8}{\parbox{\dimexpr\linewidth-2\fboxsep\relax}{\raggedright\ttfamily\scriptsize \emph{blank line}}}
\par\smallskip
This blank line is not hardware or software behavior. It separates nearby declarations or operations so the file remains readable and reviewable.
\paragraph{Line 249}
\noindent\colorbox{gray!8}{\parbox{\dimexpr\linewidth-2\fboxsep\relax}{\raggedright\ttfamily\scriptsize     initial begin}}
\par\smallskip
This simulation/synthesis initialization block loads trained memories from generated .mem files. On FPGA, Vivado maps these initial values into BRAM/LUTRAM initialization data.
\paragraph{Line 250}
\noindent\colorbox{gray!8}{\parbox{\dimexpr\linewidth-2\fboxsep\relax}{\raggedright\ttfamily\scriptsize         \#2\_000\_000\_000;}}
\par\smallskip
This is a human-facing comment. It documents intent for the reader and is ignored by the compiler or interpreter.
\paragraph{Line 251}
\noindent\colorbox{gray!8}{\parbox{\dimexpr\linewidth-2\fboxsep\relax}{\raggedright\ttfamily\scriptsize         \$display("ERROR: tb\_cifar10\_spikingformer\_head\_onchip\_learning watchdog timeout");}}
\par\smallskip
This SystemVerilog line contributes to the structural or sequential description of the accelerator. Its effect is resolved by XSim for verification and by Vivado during synthesis.
\paragraph{Line 252}
\noindent\colorbox{gray!8}{\parbox{\dimexpr\linewidth-2\fboxsep\relax}{\raggedright\ttfamily\scriptsize         \$fatal;}}
\par\smallskip
This SystemVerilog line contributes to the structural or sequential description of the accelerator. Its effect is resolved by XSim for verification and by Vivado during synthesis.
\paragraph{Line 253}
\noindent\colorbox{gray!8}{\parbox{\dimexpr\linewidth-2\fboxsep\relax}{\raggedright\ttfamily\scriptsize     end}}
\par\smallskip
This SystemVerilog line contributes to the structural or sequential description of the accelerator. Its effect is resolved by XSim for verification and by Vivado during synthesis.
\paragraph{Line 254}
\noindent\colorbox{gray!8}{\parbox{\dimexpr\linewidth-2\fboxsep\relax}{\raggedright\ttfamily\scriptsize endmodule}}
\par\smallskip
This closes the module definition. Synthesis stops adding logic to this circuit block at this point.
\clearpage
\subsection{cifar10\_neuromorphic\textbackslash{}export\_spikingformer\_head\_accelerator.py}
Each entry below gives the exact source line followed by its hardware/software role in the implemented accelerator pipeline.
\paragraph{Line 1}
\noindent\colorbox{gray!8}{\parbox{\dimexpr\linewidth-2\fboxsep\relax}{\raggedright\ttfamily\scriptsize from \_\_future\_\_ import annotations}}
\par\smallskip
This imports a Python library or symbol. The scripts use standard parsing, numeric arrays, plotting, and file APIs to make the hardware artifacts reproducible.
\paragraph{Line 2}
\noindent\colorbox{gray!8}{\parbox{\dimexpr\linewidth-2\fboxsep\relax}{\raggedright\ttfamily\scriptsize \emph{blank line}}}
\par\smallskip
This blank line is not hardware or software behavior. It separates nearby declarations or operations so the file remains readable and reviewable.
\paragraph{Line 3}
\noindent\colorbox{gray!8}{\parbox{\dimexpr\linewidth-2\fboxsep\relax}{\raggedright\ttfamily\scriptsize import argparse}}
\par\smallskip
This imports a Python library or symbol. The scripts use standard parsing, numeric arrays, plotting, and file APIs to make the hardware artifacts reproducible.
\paragraph{Line 4}
\noindent\colorbox{gray!8}{\parbox{\dimexpr\linewidth-2\fboxsep\relax}{\raggedright\ttfamily\scriptsize import csv}}
\par\smallskip
This imports a Python library or symbol. The scripts use standard parsing, numeric arrays, plotting, and file APIs to make the hardware artifacts reproducible.
\paragraph{Line 5}
\noindent\colorbox{gray!8}{\parbox{\dimexpr\linewidth-2\fboxsep\relax}{\raggedright\ttfamily\scriptsize import json}}
\par\smallskip
This imports a Python library or symbol. The scripts use standard parsing, numeric arrays, plotting, and file APIs to make the hardware artifacts reproducible.
\paragraph{Line 6}
\noindent\colorbox{gray!8}{\parbox{\dimexpr\linewidth-2\fboxsep\relax}{\raggedright\ttfamily\scriptsize import os}}
\par\smallskip
This imports a Python library or symbol. The scripts use standard parsing, numeric arrays, plotting, and file APIs to make the hardware artifacts reproducible.
\paragraph{Line 7}
\noindent\colorbox{gray!8}{\parbox{\dimexpr\linewidth-2\fboxsep\relax}{\raggedright\ttfamily\scriptsize import sys}}
\par\smallskip
This imports a Python library or symbol. The scripts use standard parsing, numeric arrays, plotting, and file APIs to make the hardware artifacts reproducible.
\paragraph{Line 8}
\noindent\colorbox{gray!8}{\parbox{\dimexpr\linewidth-2\fboxsep\relax}{\raggedright\ttfamily\scriptsize from contextlib import nullcontext}}
\par\smallskip
This imports a Python library or symbol. The scripts use standard parsing, numeric arrays, plotting, and file APIs to make the hardware artifacts reproducible.
\paragraph{Line 9}
\noindent\colorbox{gray!8}{\parbox{\dimexpr\linewidth-2\fboxsep\relax}{\raggedright\ttfamily\scriptsize from functools import partial}}
\par\smallskip
This imports a Python library or symbol. The scripts use standard parsing, numeric arrays, plotting, and file APIs to make the hardware artifacts reproducible.
\paragraph{Line 10}
\noindent\colorbox{gray!8}{\parbox{\dimexpr\linewidth-2\fboxsep\relax}{\raggedright\ttfamily\scriptsize from pathlib import Path}}
\par\smallskip
This imports a Python library or symbol. The scripts use standard parsing, numeric arrays, plotting, and file APIs to make the hardware artifacts reproducible.
\paragraph{Line 11}
\noindent\colorbox{gray!8}{\parbox{\dimexpr\linewidth-2\fboxsep\relax}{\raggedright\ttfamily\scriptsize \emph{blank line}}}
\par\smallskip
This blank line is not hardware or software behavior. It separates nearby declarations or operations so the file remains readable and reviewable.
\paragraph{Line 12}
\noindent\colorbox{gray!8}{\parbox{\dimexpr\linewidth-2\fboxsep\relax}{\raggedright\ttfamily\scriptsize import matplotlib.pyplot as plt}}
\par\smallskip
This imports a Python library or symbol. The scripts use standard parsing, numeric arrays, plotting, and file APIs to make the hardware artifacts reproducible.
\paragraph{Line 13}
\noindent\colorbox{gray!8}{\parbox{\dimexpr\linewidth-2\fboxsep\relax}{\raggedright\ttfamily\scriptsize import numpy as np}}
\par\smallskip
This imports a Python library or symbol. The scripts use standard parsing, numeric arrays, plotting, and file APIs to make the hardware artifacts reproducible.
\paragraph{Line 14}
\noindent\colorbox{gray!8}{\parbox{\dimexpr\linewidth-2\fboxsep\relax}{\raggedright\ttfamily\scriptsize import torch}}
\par\smallskip
This imports a Python library or symbol. The scripts use standard parsing, numeric arrays, plotting, and file APIs to make the hardware artifacts reproducible.
\paragraph{Line 15}
\noindent\colorbox{gray!8}{\parbox{\dimexpr\linewidth-2\fboxsep\relax}{\raggedright\ttfamily\scriptsize import torch.nn as nn}}
\par\smallskip
This imports a Python library or symbol. The scripts use standard parsing, numeric arrays, plotting, and file APIs to make the hardware artifacts reproducible.
\paragraph{Line 16}
\noindent\colorbox{gray!8}{\parbox{\dimexpr\linewidth-2\fboxsep\relax}{\raggedright\ttfamily\scriptsize from matplotlib.backends.backend\_pdf import PdfPages}}
\par\smallskip
This imports a Python library or symbol. The scripts use standard parsing, numeric arrays, plotting, and file APIs to make the hardware artifacts reproducible.
\paragraph{Line 17}
\noindent\colorbox{gray!8}{\parbox{\dimexpr\linewidth-2\fboxsep\relax}{\raggedright\ttfamily\scriptsize from timm.data import create\_transform}}
\par\smallskip
This imports a Python library or symbol. The scripts use standard parsing, numeric arrays, plotting, and file APIs to make the hardware artifacts reproducible.
\paragraph{Line 18}
\noindent\colorbox{gray!8}{\parbox{\dimexpr\linewidth-2\fboxsep\relax}{\raggedright\ttfamily\scriptsize from timm.models import create\_model}}
\par\smallskip
This imports a Python library or symbol. The scripts use standard parsing, numeric arrays, plotting, and file APIs to make the hardware artifacts reproducible.
\paragraph{Line 19}
\noindent\colorbox{gray!8}{\parbox{\dimexpr\linewidth-2\fboxsep\relax}{\raggedright\ttfamily\scriptsize from torch.utils.data import DataLoader, Subset}}
\par\smallskip
This imports a Python library or symbol. The scripts use standard parsing, numeric arrays, plotting, and file APIs to make the hardware artifacts reproducible.
\paragraph{Line 20}
\noindent\colorbox{gray!8}{\parbox{\dimexpr\linewidth-2\fboxsep\relax}{\raggedright\ttfamily\scriptsize from torchvision import datasets}}
\par\smallskip
This imports a Python library or symbol. The scripts use standard parsing, numeric arrays, plotting, and file APIs to make the hardware artifacts reproducible.
\paragraph{Line 21}
\noindent\colorbox{gray!8}{\parbox{\dimexpr\linewidth-2\fboxsep\relax}{\raggedright\ttfamily\scriptsize \emph{blank line}}}
\par\smallskip
This blank line is not hardware or software behavior. It separates nearby declarations or operations so the file remains readable and reviewable.
\paragraph{Line 22}
\noindent\colorbox{gray!8}{\parbox{\dimexpr\linewidth-2\fboxsep\relax}{\raggedright\ttfamily\scriptsize \emph{blank line}}}
\par\smallskip
This blank line is not hardware or software behavior. It separates nearby declarations or operations so the file remains readable and reviewable.
\paragraph{Line 23}
\noindent\colorbox{gray!8}{\parbox{\dimexpr\linewidth-2\fboxsep\relax}{\raggedright\ttfamily\scriptsize CIFAR10\_CLASSES = [}}
\par\smallskip
This Python line participates in the export, validation, or report-generation pipeline. It is executed on the host PC to prepare data for RTL or to check RTL outputs.
\paragraph{Line 24}
\noindent\colorbox{gray!8}{\parbox{\dimexpr\linewidth-2\fboxsep\relax}{\raggedright\ttfamily\scriptsize     "airplane",}}
\par\smallskip
This Python line participates in the export, validation, or report-generation pipeline. It is executed on the host PC to prepare data for RTL or to check RTL outputs.
\paragraph{Line 25}
\noindent\colorbox{gray!8}{\parbox{\dimexpr\linewidth-2\fboxsep\relax}{\raggedright\ttfamily\scriptsize     "automobile",}}
\par\smallskip
This Python line participates in the export, validation, or report-generation pipeline. It is executed on the host PC to prepare data for RTL or to check RTL outputs.
\paragraph{Line 26}
\noindent\colorbox{gray!8}{\parbox{\dimexpr\linewidth-2\fboxsep\relax}{\raggedright\ttfamily\scriptsize     "bird",}}
\par\smallskip
This Python line participates in the export, validation, or report-generation pipeline. It is executed on the host PC to prepare data for RTL or to check RTL outputs.
\paragraph{Line 27}
\noindent\colorbox{gray!8}{\parbox{\dimexpr\linewidth-2\fboxsep\relax}{\raggedright\ttfamily\scriptsize     "cat",}}
\par\smallskip
This Python line participates in the export, validation, or report-generation pipeline. It is executed on the host PC to prepare data for RTL or to check RTL outputs.
\paragraph{Line 28}
\noindent\colorbox{gray!8}{\parbox{\dimexpr\linewidth-2\fboxsep\relax}{\raggedright\ttfamily\scriptsize     "deer",}}
\par\smallskip
This Python line participates in the export, validation, or report-generation pipeline. It is executed on the host PC to prepare data for RTL or to check RTL outputs.
\paragraph{Line 29}
\noindent\colorbox{gray!8}{\parbox{\dimexpr\linewidth-2\fboxsep\relax}{\raggedright\ttfamily\scriptsize     "dog",}}
\par\smallskip
This Python line participates in the export, validation, or report-generation pipeline. It is executed on the host PC to prepare data for RTL or to check RTL outputs.
\paragraph{Line 30}
\noindent\colorbox{gray!8}{\parbox{\dimexpr\linewidth-2\fboxsep\relax}{\raggedright\ttfamily\scriptsize     "frog",}}
\par\smallskip
This Python line participates in the export, validation, or report-generation pipeline. It is executed on the host PC to prepare data for RTL or to check RTL outputs.
\paragraph{Line 31}
\noindent\colorbox{gray!8}{\parbox{\dimexpr\linewidth-2\fboxsep\relax}{\raggedright\ttfamily\scriptsize     "horse",}}
\par\smallskip
This Python line participates in the export, validation, or report-generation pipeline. It is executed on the host PC to prepare data for RTL or to check RTL outputs.
\paragraph{Line 32}
\noindent\colorbox{gray!8}{\parbox{\dimexpr\linewidth-2\fboxsep\relax}{\raggedright\ttfamily\scriptsize     "ship",}}
\par\smallskip
This Python line participates in the export, validation, or report-generation pipeline. It is executed on the host PC to prepare data for RTL or to check RTL outputs.
\paragraph{Line 33}
\noindent\colorbox{gray!8}{\parbox{\dimexpr\linewidth-2\fboxsep\relax}{\raggedright\ttfamily\scriptsize     "truck",}}
\par\smallskip
This Python line participates in the export, validation, or report-generation pipeline. It is executed on the host PC to prepare data for RTL or to check RTL outputs.
\paragraph{Line 34}
\noindent\colorbox{gray!8}{\parbox{\dimexpr\linewidth-2\fboxsep\relax}{\raggedright\ttfamily\scriptsize ]}}
\par\smallskip
This Python line participates in the export, validation, or report-generation pipeline. It is executed on the host PC to prepare data for RTL or to check RTL outputs.
\paragraph{Line 35}
\noindent\colorbox{gray!8}{\parbox{\dimexpr\linewidth-2\fboxsep\relax}{\raggedright\ttfamily\scriptsize CIFAR10\_MEAN = (0.4914, 0.4822, 0.4465)}}
\par\smallskip
This Python line participates in the export, validation, or report-generation pipeline. It is executed on the host PC to prepare data for RTL or to check RTL outputs.
\paragraph{Line 36}
\noindent\colorbox{gray!8}{\parbox{\dimexpr\linewidth-2\fboxsep\relax}{\raggedright\ttfamily\scriptsize CIFAR10\_STD = (0.2470, 0.2435, 0.2616)}}
\par\smallskip
This Python line participates in the export, validation, or report-generation pipeline. It is executed on the host PC to prepare data for RTL or to check RTL outputs.
\paragraph{Line 37}
\noindent\colorbox{gray!8}{\parbox{\dimexpr\linewidth-2\fboxsep\relax}{\raggedright\ttfamily\scriptsize NUM\_CLASSES = 10}}
\par\smallskip
This Python line participates in the export, validation, or report-generation pipeline. It is executed on the host PC to prepare data for RTL or to check RTL outputs.
\paragraph{Line 38}
\noindent\colorbox{gray!8}{\parbox{\dimexpr\linewidth-2\fboxsep\relax}{\raggedright\ttfamily\scriptsize FEATURE\_DIM = 384}}
\par\smallskip
This Python line participates in the export, validation, or report-generation pipeline. It is executed on the host PC to prepare data for RTL or to check RTL outputs.
\paragraph{Line 39}
\noindent\colorbox{gray!8}{\parbox{\dimexpr\linewidth-2\fboxsep\relax}{\raggedright\ttfamily\scriptsize \emph{blank line}}}
\par\smallskip
This blank line is not hardware or software behavior. It separates nearby declarations or operations so the file remains readable and reviewable.
\paragraph{Line 40}
\noindent\colorbox{gray!8}{\parbox{\dimexpr\linewidth-2\fboxsep\relax}{\raggedright\ttfamily\scriptsize \emph{blank line}}}
\par\smallskip
This blank line is not hardware or software behavior. It separates nearby declarations or operations so the file remains readable and reviewable.
\paragraph{Line 41}
\noindent\colorbox{gray!8}{\parbox{\dimexpr\linewidth-2\fboxsep\relax}{\raggedright\ttfamily\scriptsize def repo\_root() -> Path:}}
\par\smallskip
This defines a Python function. The function packages one repeatable step such as loading a memory file, simulating the golden learning rule, validating CSVs, or generating plots.
\paragraph{Line 42}
\noindent\colorbox{gray!8}{\parbox{\dimexpr\linewidth-2\fboxsep\relax}{\raggedright\ttfamily\scriptsize     return Path(\_\_file\_\_).resolve().parents[1]}}
\par\smallskip
This uses path objects so generated artifacts are addressed consistently across Windows shells and Python code.
\paragraph{Line 43}
\noindent\colorbox{gray!8}{\parbox{\dimexpr\linewidth-2\fboxsep\relax}{\raggedright\ttfamily\scriptsize \emph{blank line}}}
\par\smallskip
This blank line is not hardware or software behavior. It separates nearby declarations or operations so the file remains readable and reviewable.
\paragraph{Line 44}
\noindent\colorbox{gray!8}{\parbox{\dimexpr\linewidth-2\fboxsep\relax}{\raggedright\ttfamily\scriptsize \emph{blank line}}}
\par\smallskip
This blank line is not hardware or software behavior. It separates nearby declarations or operations so the file remains readable and reviewable.
\paragraph{Line 45}
\noindent\colorbox{gray!8}{\parbox{\dimexpr\linewidth-2\fboxsep\relax}{\raggedright\ttfamily\scriptsize def add\_runtime\_dlls() -> None:}}
\par\smallskip
This defines a Python function. The function packages one repeatable step such as loading a memory file, simulating the golden learning rule, validating CSVs, or generating plots.
\paragraph{Line 46}
\noindent\colorbox{gray!8}{\parbox{\dimexpr\linewidth-2\fboxsep\relax}{\raggedright\ttfamily\scriptsize     env\_root = repo\_root() / "external\_research" / "spikingformer\_py310"}}
\par\smallskip
This Python line participates in the export, validation, or report-generation pipeline. It is executed on the host PC to prepare data for RTL or to check RTL outputs.
\paragraph{Line 47}
\noindent\colorbox{gray!8}{\parbox{\dimexpr\linewidth-2\fboxsep\relax}{\raggedright\ttfamily\scriptsize     for path in (}}
\par\smallskip
This Python line participates in the export, validation, or report-generation pipeline. It is executed on the host PC to prepare data for RTL or to check RTL outputs.
\paragraph{Line 48}
\noindent\colorbox{gray!8}{\parbox{\dimexpr\linewidth-2\fboxsep\relax}{\raggedright\ttfamily\scriptsize         env\_root / "Lib" / "site-packages" / "torch" / "lib",}}
\par\smallskip
This Python line participates in the export, validation, or report-generation pipeline. It is executed on the host PC to prepare data for RTL or to check RTL outputs.
\paragraph{Line 49}
\noindent\colorbox{gray!8}{\parbox{\dimexpr\linewidth-2\fboxsep\relax}{\raggedright\ttfamily\scriptsize         env\_root / "Lib" / "site-packages" / "nvidia" / "cuda\_runtime" / "bin",}}
\par\smallskip
This Python line participates in the export, validation, or report-generation pipeline. It is executed on the host PC to prepare data for RTL or to check RTL outputs.
\paragraph{Line 50}
\noindent\colorbox{gray!8}{\parbox{\dimexpr\linewidth-2\fboxsep\relax}{\raggedright\ttfamily\scriptsize         env\_root / "Lib" / "site-packages" / "nvidia" / "cuda\_nvrtc" / "bin",}}
\par\smallskip
This Python line participates in the export, validation, or report-generation pipeline. It is executed on the host PC to prepare data for RTL or to check RTL outputs.
\paragraph{Line 51}
\noindent\colorbox{gray!8}{\parbox{\dimexpr\linewidth-2\fboxsep\relax}{\raggedright\ttfamily\scriptsize     ):}}
\par\smallskip
This Python line participates in the export, validation, or report-generation pipeline. It is executed on the host PC to prepare data for RTL or to check RTL outputs.
\paragraph{Line 52}
\noindent\colorbox{gray!8}{\parbox{\dimexpr\linewidth-2\fboxsep\relax}{\raggedright\ttfamily\scriptsize         if path.exists():}}
\par\smallskip
This Python line participates in the export, validation, or report-generation pipeline. It is executed on the host PC to prepare data for RTL or to check RTL outputs.
\paragraph{Line 53}
\noindent\colorbox{gray!8}{\parbox{\dimexpr\linewidth-2\fboxsep\relax}{\raggedright\ttfamily\scriptsize             try:}}
\par\smallskip
This Python line participates in the export, validation, or report-generation pipeline. It is executed on the host PC to prepare data for RTL or to check RTL outputs.
\paragraph{Line 54}
\noindent\colorbox{gray!8}{\parbox{\dimexpr\linewidth-2\fboxsep\relax}{\raggedright\ttfamily\scriptsize                 os.add\_dll\_directory(str(path))}}
\par\smallskip
This Python line participates in the export, validation, or report-generation pipeline. It is executed on the host PC to prepare data for RTL or to check RTL outputs.
\paragraph{Line 55}
\noindent\colorbox{gray!8}{\parbox{\dimexpr\linewidth-2\fboxsep\relax}{\raggedright\ttfamily\scriptsize             except (AttributeError, OSError):}}
\par\smallskip
This Python line participates in the export, validation, or report-generation pipeline. It is executed on the host PC to prepare data for RTL or to check RTL outputs.
\paragraph{Line 56}
\noindent\colorbox{gray!8}{\parbox{\dimexpr\linewidth-2\fboxsep\relax}{\raggedright\ttfamily\scriptsize                 pass}}
\par\smallskip
This Python line participates in the export, validation, or report-generation pipeline. It is executed on the host PC to prepare data for RTL or to check RTL outputs.
\paragraph{Line 57}
\noindent\colorbox{gray!8}{\parbox{\dimexpr\linewidth-2\fboxsep\relax}{\raggedright\ttfamily\scriptsize             os.environ["PATH"] = str(path) + os.pathsep + os.environ.get("PATH", "")}}
\par\smallskip
This Python line participates in the export, validation, or report-generation pipeline. It is executed on the host PC to prepare data for RTL or to check RTL outputs.
\paragraph{Line 58}
\noindent\colorbox{gray!8}{\parbox{\dimexpr\linewidth-2\fboxsep\relax}{\raggedright\ttfamily\scriptsize \emph{blank line}}}
\par\smallskip
This blank line is not hardware or software behavior. It separates nearby declarations or operations so the file remains readable and reviewable.
\paragraph{Line 59}
\noindent\colorbox{gray!8}{\parbox{\dimexpr\linewidth-2\fboxsep\relax}{\raggedright\ttfamily\scriptsize \emph{blank line}}}
\par\smallskip
This blank line is not hardware or software behavior. It separates nearby declarations or operations so the file remains readable and reviewable.
\paragraph{Line 60}
\noindent\colorbox{gray!8}{\parbox{\dimexpr\linewidth-2\fboxsep\relax}{\raggedright\ttfamily\scriptsize def import\_upstream\_model() -> None:}}
\par\smallskip
This defines a Python function. The function packages one repeatable step such as loading a memory file, simulating the golden learning rule, validating CSVs, or generating plots.
\paragraph{Line 61}
\noindent\colorbox{gray!8}{\parbox{\dimexpr\linewidth-2\fboxsep\relax}{\raggedright\ttfamily\scriptsize     upstream = repo\_root() / "external\_research" / "Spikingformer" / "cifar10"}}
\par\smallskip
This Python line participates in the export, validation, or report-generation pipeline. It is executed on the host PC to prepare data for RTL or to check RTL outputs.
\paragraph{Line 62}
\noindent\colorbox{gray!8}{\parbox{\dimexpr\linewidth-2\fboxsep\relax}{\raggedright\ttfamily\scriptsize     if not upstream.exists():}}
\par\smallskip
This Python line participates in the export, validation, or report-generation pipeline. It is executed on the host PC to prepare data for RTL or to check RTL outputs.
\paragraph{Line 63}
\noindent\colorbox{gray!8}{\parbox{\dimexpr\linewidth-2\fboxsep\relax}{\raggedright\ttfamily\scriptsize         raise FileNotFoundError(f"Missing upstream Spikingformer source: \{upstream\}")}}
\par\smallskip
This Python line participates in the export, validation, or report-generation pipeline. It is executed on the host PC to prepare data for RTL or to check RTL outputs.
\paragraph{Line 64}
\noindent\colorbox{gray!8}{\parbox{\dimexpr\linewidth-2\fboxsep\relax}{\raggedright\ttfamily\scriptsize     sys.path.insert(0, str(upstream))}}
\par\smallskip
This Python line participates in the export, validation, or report-generation pipeline. It is executed on the host PC to prepare data for RTL or to check RTL outputs.
\paragraph{Line 65}
\noindent\colorbox{gray!8}{\parbox{\dimexpr\linewidth-2\fboxsep\relax}{\raggedright\ttfamily\scriptsize     import model  \# noqa: F401}}
\par\smallskip
This imports a Python library or symbol. The scripts use standard parsing, numeric arrays, plotting, and file APIs to make the hardware artifacts reproducible.
\paragraph{Line 66}
\noindent\colorbox{gray!8}{\parbox{\dimexpr\linewidth-2\fboxsep\relax}{\raggedright\ttfamily\scriptsize \emph{blank line}}}
\par\smallskip
This blank line is not hardware or software behavior. It separates nearby declarations or operations so the file remains readable and reviewable.
\paragraph{Line 67}
\noindent\colorbox{gray!8}{\parbox{\dimexpr\linewidth-2\fboxsep\relax}{\raggedright\ttfamily\scriptsize \emph{blank line}}}
\par\smallskip
This blank line is not hardware or software behavior. It separates nearby declarations or operations so the file remains readable and reviewable.
\paragraph{Line 68}
\noindent\colorbox{gray!8}{\parbox{\dimexpr\linewidth-2\fboxsep\relax}{\raggedright\ttfamily\scriptsize def signed\_hex(value: int, bits: int) -> str:}}
\par\smallskip
This defines a Python function. The function packages one repeatable step such as loading a memory file, simulating the golden learning rule, validating CSVs, or generating plots.
\paragraph{Line 69}
\noindent\colorbox{gray!8}{\parbox{\dimexpr\linewidth-2\fboxsep\relax}{\raggedright\ttfamily\scriptsize     return f"\{value \& ((1 << bits) - 1):0\{bits // 4\}x\}"}}
\par\smallskip
This Python line participates in the export, validation, or report-generation pipeline. It is executed on the host PC to prepare data for RTL or to check RTL outputs.
\paragraph{Line 70}
\noindent\colorbox{gray!8}{\parbox{\dimexpr\linewidth-2\fboxsep\relax}{\raggedright\ttfamily\scriptsize \emph{blank line}}}
\par\smallskip
This blank line is not hardware or software behavior. It separates nearby declarations or operations so the file remains readable and reviewable.
\paragraph{Line 71}
\noindent\colorbox{gray!8}{\parbox{\dimexpr\linewidth-2\fboxsep\relax}{\raggedright\ttfamily\scriptsize \emph{blank line}}}
\par\smallskip
This blank line is not hardware or software behavior. It separates nearby declarations or operations so the file remains readable and reviewable.
\paragraph{Line 72}
\noindent\colorbox{gray!8}{\parbox{\dimexpr\linewidth-2\fboxsep\relax}{\raggedright\ttfamily\scriptsize def choose\_scale(values: torch.Tensor, bits: int) -> float:}}
\par\smallskip
This defines a Python function. The function packages one repeatable step such as loading a memory file, simulating the golden learning rule, validating CSVs, or generating plots.
\paragraph{Line 73}
\noindent\colorbox{gray!8}{\parbox{\dimexpr\linewidth-2\fboxsep\relax}{\raggedright\ttfamily\scriptsize     qmax = float((1 << (bits - 1)) - 1)}}
\par\smallskip
This Python line participates in the export, validation, or report-generation pipeline. It is executed on the host PC to prepare data for RTL or to check RTL outputs.
\paragraph{Line 74}
\noindent\colorbox{gray!8}{\parbox{\dimexpr\linewidth-2\fboxsep\relax}{\raggedright\ttfamily\scriptsize     max\_abs = float(values.detach().abs().max().cpu())}}
\par\smallskip
This Python line participates in the export, validation, or report-generation pipeline. It is executed on the host PC to prepare data for RTL or to check RTL outputs.
\paragraph{Line 75}
\noindent\colorbox{gray!8}{\parbox{\dimexpr\linewidth-2\fboxsep\relax}{\raggedright\ttfamily\scriptsize     return max(max\_abs / qmax, 1.0e-12)}}
\par\smallskip
This Python line participates in the export, validation, or report-generation pipeline. It is executed on the host PC to prepare data for RTL or to check RTL outputs.
\paragraph{Line 76}
\noindent\colorbox{gray!8}{\parbox{\dimexpr\linewidth-2\fboxsep\relax}{\raggedright\ttfamily\scriptsize \emph{blank line}}}
\par\smallskip
This blank line is not hardware or software behavior. It separates nearby declarations or operations so the file remains readable and reviewable.
\paragraph{Line 77}
\noindent\colorbox{gray!8}{\parbox{\dimexpr\linewidth-2\fboxsep\relax}{\raggedright\ttfamily\scriptsize \emph{blank line}}}
\par\smallskip
This blank line is not hardware or software behavior. It separates nearby declarations or operations so the file remains readable and reviewable.
\paragraph{Line 78}
\noindent\colorbox{gray!8}{\parbox{\dimexpr\linewidth-2\fboxsep\relax}{\raggedright\ttfamily\scriptsize def quantize\_signed(values: torch.Tensor, scale: float, bits: int) -> torch.Tensor:}}
\par\smallskip
This defines a Python function. The function packages one repeatable step such as loading a memory file, simulating the golden learning rule, validating CSVs, or generating plots.
\paragraph{Line 79}
\noindent\colorbox{gray!8}{\parbox{\dimexpr\linewidth-2\fboxsep\relax}{\raggedright\ttfamily\scriptsize     qmin = -(1 << (bits - 1))}}
\par\smallskip
This Python line participates in the export, validation, or report-generation pipeline. It is executed on the host PC to prepare data for RTL or to check RTL outputs.
\paragraph{Line 80}
\noindent\colorbox{gray!8}{\parbox{\dimexpr\linewidth-2\fboxsep\relax}{\raggedright\ttfamily\scriptsize     qmax = (1 << (bits - 1)) - 1}}
\par\smallskip
This Python line participates in the export, validation, or report-generation pipeline. It is executed on the host PC to prepare data for RTL or to check RTL outputs.
\paragraph{Line 81}
\noindent\colorbox{gray!8}{\parbox{\dimexpr\linewidth-2\fboxsep\relax}{\raggedright\ttfamily\scriptsize     return torch.round(values / scale).clamp(qmin, qmax).to(torch.int64)}}
\par\smallskip
This Python line participates in the export, validation, or report-generation pipeline. It is executed on the host PC to prepare data for RTL or to check RTL outputs.
\paragraph{Line 82}
\noindent\colorbox{gray!8}{\parbox{\dimexpr\linewidth-2\fboxsep\relax}{\raggedright\ttfamily\scriptsize \emph{blank line}}}
\par\smallskip
This blank line is not hardware or software behavior. It separates nearby declarations or operations so the file remains readable and reviewable.
\paragraph{Line 83}
\noindent\colorbox{gray!8}{\parbox{\dimexpr\linewidth-2\fboxsep\relax}{\raggedright\ttfamily\scriptsize \emph{blank line}}}
\par\smallskip
This blank line is not hardware or software behavior. It separates nearby declarations or operations so the file remains readable and reviewable.
\paragraph{Line 84}
\noindent\colorbox{gray!8}{\parbox{\dimexpr\linewidth-2\fboxsep\relax}{\raggedright\ttfamily\scriptsize def load\_checkpoint(path: Path) -> dict[str, torch.Tensor]:}}
\par\smallskip
This defines a Python function. The function packages one repeatable step such as loading a memory file, simulating the golden learning rule, validating CSVs, or generating plots.
\paragraph{Line 85}
\noindent\colorbox{gray!8}{\parbox{\dimexpr\linewidth-2\fboxsep\relax}{\raggedright\ttfamily\scriptsize     checkpoint = torch.load(path, map\_location="cpu")}}
\par\smallskip
This Python line participates in the export, validation, or report-generation pipeline. It is executed on the host PC to prepare data for RTL or to check RTL outputs.
\paragraph{Line 86}
\noindent\colorbox{gray!8}{\parbox{\dimexpr\linewidth-2\fboxsep\relax}{\raggedright\ttfamily\scriptsize     state = checkpoint.get("state\_dict", checkpoint)}}
\par\smallskip
This Python line participates in the export, validation, or report-generation pipeline. It is executed on the host PC to prepare data for RTL or to check RTL outputs.
\paragraph{Line 87}
\noindent\colorbox{gray!8}{\parbox{\dimexpr\linewidth-2\fboxsep\relax}{\raggedright\ttfamily\scriptsize     clean: dict[str, torch.Tensor] = \{\}}}
\par\smallskip
This Python line participates in the export, validation, or report-generation pipeline. It is executed on the host PC to prepare data for RTL or to check RTL outputs.
\paragraph{Line 88}
\noindent\colorbox{gray!8}{\parbox{\dimexpr\linewidth-2\fboxsep\relax}{\raggedright\ttfamily\scriptsize     for key, value in state.items():}}
\par\smallskip
This Python line participates in the export, validation, or report-generation pipeline. It is executed on the host PC to prepare data for RTL or to check RTL outputs.
\paragraph{Line 89}
\noindent\colorbox{gray!8}{\parbox{\dimexpr\linewidth-2\fboxsep\relax}{\raggedright\ttfamily\scriptsize         clean[key.removeprefix("module.")] = value}}
\par\smallskip
This Python line participates in the export, validation, or report-generation pipeline. It is executed on the host PC to prepare data for RTL or to check RTL outputs.
\paragraph{Line 90}
\noindent\colorbox{gray!8}{\parbox{\dimexpr\linewidth-2\fboxsep\relax}{\raggedright\ttfamily\scriptsize     return clean}}
\par\smallskip
This Python line participates in the export, validation, or report-generation pipeline. It is executed on the host PC to prepare data for RTL or to check RTL outputs.
\paragraph{Line 91}
\noindent\colorbox{gray!8}{\parbox{\dimexpr\linewidth-2\fboxsep\relax}{\raggedright\ttfamily\scriptsize \emph{blank line}}}
\par\smallskip
This blank line is not hardware or software behavior. It separates nearby declarations or operations so the file remains readable and reviewable.
\paragraph{Line 92}
\noindent\colorbox{gray!8}{\parbox{\dimexpr\linewidth-2\fboxsep\relax}{\raggedright\ttfamily\scriptsize \emph{blank line}}}
\par\smallskip
This blank line is not hardware or software behavior. It separates nearby declarations or operations so the file remains readable and reviewable.
\paragraph{Line 93}
\noindent\colorbox{gray!8}{\parbox{\dimexpr\linewidth-2\fboxsep\relax}{\raggedright\ttfamily\scriptsize def build\_model(device: torch.device) -> nn.Module:}}
\par\smallskip
This defines a Python function. The function packages one repeatable step such as loading a memory file, simulating the golden learning rule, validating CSVs, or generating plots.
\paragraph{Line 94}
\noindent\colorbox{gray!8}{\parbox{\dimexpr\linewidth-2\fboxsep\relax}{\raggedright\ttfamily\scriptsize     model = create\_model(}}
\par\smallskip
This Python line participates in the export, validation, or report-generation pipeline. It is executed on the host PC to prepare data for RTL or to check RTL outputs.
\paragraph{Line 95}
\noindent\colorbox{gray!8}{\parbox{\dimexpr\linewidth-2\fboxsep\relax}{\raggedright\ttfamily\scriptsize         "Spikingformer",}}
\par\smallskip
This Python line participates in the export, validation, or report-generation pipeline. It is executed on the host PC to prepare data for RTL or to check RTL outputs.
\paragraph{Line 96}
\noindent\colorbox{gray!8}{\parbox{\dimexpr\linewidth-2\fboxsep\relax}{\raggedright\ttfamily\scriptsize         pretrained=False,}}
\par\smallskip
This Python line participates in the export, validation, or report-generation pipeline. It is executed on the host PC to prepare data for RTL or to check RTL outputs.
\paragraph{Line 97}
\noindent\colorbox{gray!8}{\parbox{\dimexpr\linewidth-2\fboxsep\relax}{\raggedright\ttfamily\scriptsize         drop\_rate=0.0,}}
\par\smallskip
This Python line participates in the export, validation, or report-generation pipeline. It is executed on the host PC to prepare data for RTL or to check RTL outputs.
\paragraph{Line 98}
\noindent\colorbox{gray!8}{\parbox{\dimexpr\linewidth-2\fboxsep\relax}{\raggedright\ttfamily\scriptsize         drop\_path\_rate=0.0,}}
\par\smallskip
This Python line participates in the export, validation, or report-generation pipeline. It is executed on the host PC to prepare data for RTL or to check RTL outputs.
\paragraph{Line 99}
\noindent\colorbox{gray!8}{\parbox{\dimexpr\linewidth-2\fboxsep\relax}{\raggedright\ttfamily\scriptsize         drop\_block\_rate=None,}}
\par\smallskip
This Python line participates in the export, validation, or report-generation pipeline. It is executed on the host PC to prepare data for RTL or to check RTL outputs.
\paragraph{Line 100}
\noindent\colorbox{gray!8}{\parbox{\dimexpr\linewidth-2\fboxsep\relax}{\raggedright\ttfamily\scriptsize         img\_size\_h=32,}}
\par\smallskip
This Python line participates in the export, validation, or report-generation pipeline. It is executed on the host PC to prepare data for RTL or to check RTL outputs.
\paragraph{Line 101}
\noindent\colorbox{gray!8}{\parbox{\dimexpr\linewidth-2\fboxsep\relax}{\raggedright\ttfamily\scriptsize         img\_size\_w=32,}}
\par\smallskip
This Python line participates in the export, validation, or report-generation pipeline. It is executed on the host PC to prepare data for RTL or to check RTL outputs.
\paragraph{Line 102}
\noindent\colorbox{gray!8}{\parbox{\dimexpr\linewidth-2\fboxsep\relax}{\raggedright\ttfamily\scriptsize         patch\_size=4,}}
\par\smallskip
This Python line participates in the export, validation, or report-generation pipeline. It is executed on the host PC to prepare data for RTL or to check RTL outputs.
\paragraph{Line 103}
\noindent\colorbox{gray!8}{\parbox{\dimexpr\linewidth-2\fboxsep\relax}{\raggedright\ttfamily\scriptsize         embed\_dims=FEATURE\_DIM,}}
\par\smallskip
This Python line participates in the export, validation, or report-generation pipeline. It is executed on the host PC to prepare data for RTL or to check RTL outputs.
\paragraph{Line 104}
\noindent\colorbox{gray!8}{\parbox{\dimexpr\linewidth-2\fboxsep\relax}{\raggedright\ttfamily\scriptsize         num\_heads=12,}}
\par\smallskip
This Python line participates in the export, validation, or report-generation pipeline. It is executed on the host PC to prepare data for RTL or to check RTL outputs.
\paragraph{Line 105}
\noindent\colorbox{gray!8}{\parbox{\dimexpr\linewidth-2\fboxsep\relax}{\raggedright\ttfamily\scriptsize         mlp\_ratios=4,}}
\par\smallskip
This Python line participates in the export, validation, or report-generation pipeline. It is executed on the host PC to prepare data for RTL or to check RTL outputs.
\paragraph{Line 106}
\noindent\colorbox{gray!8}{\parbox{\dimexpr\linewidth-2\fboxsep\relax}{\raggedright\ttfamily\scriptsize         in\_channels=3,}}
\par\smallskip
This Python line participates in the export, validation, or report-generation pipeline. It is executed on the host PC to prepare data for RTL or to check RTL outputs.
\paragraph{Line 107}
\noindent\colorbox{gray!8}{\parbox{\dimexpr\linewidth-2\fboxsep\relax}{\raggedright\ttfamily\scriptsize         num\_classes=NUM\_CLASSES,}}
\par\smallskip
This Python line participates in the export, validation, or report-generation pipeline. It is executed on the host PC to prepare data for RTL or to check RTL outputs.
\paragraph{Line 108}
\noindent\colorbox{gray!8}{\parbox{\dimexpr\linewidth-2\fboxsep\relax}{\raggedright\ttfamily\scriptsize         qkv\_bias=False,}}
\par\smallskip
This line handles the classifier parameters. Those arrays mirror the RTL memories, so mismatches reveal an export, quantization, or hardware-update bug.
\paragraph{Line 109}
\noindent\colorbox{gray!8}{\parbox{\dimexpr\linewidth-2\fboxsep\relax}{\raggedright\ttfamily\scriptsize         norm\_layer=partial(nn.LayerNorm, eps=1.0e-6),}}
\par\smallskip
This Python line participates in the export, validation, or report-generation pipeline. It is executed on the host PC to prepare data for RTL or to check RTL outputs.
\paragraph{Line 110}
\noindent\colorbox{gray!8}{\parbox{\dimexpr\linewidth-2\fboxsep\relax}{\raggedright\ttfamily\scriptsize         depths=4,}}
\par\smallskip
This Python line participates in the export, validation, or report-generation pipeline. It is executed on the host PC to prepare data for RTL or to check RTL outputs.
\paragraph{Line 111}
\noindent\colorbox{gray!8}{\parbox{\dimexpr\linewidth-2\fboxsep\relax}{\raggedright\ttfamily\scriptsize         sr\_ratios=1,}}
\par\smallskip
This Python line participates in the export, validation, or report-generation pipeline. It is executed on the host PC to prepare data for RTL or to check RTL outputs.
\paragraph{Line 112}
\noindent\colorbox{gray!8}{\parbox{\dimexpr\linewidth-2\fboxsep\relax}{\raggedright\ttfamily\scriptsize         T=4,}}
\par\smallskip
This Python line participates in the export, validation, or report-generation pipeline. It is executed on the host PC to prepare data for RTL or to check RTL outputs.
\paragraph{Line 113}
\noindent\colorbox{gray!8}{\parbox{\dimexpr\linewidth-2\fboxsep\relax}{\raggedright\ttfamily\scriptsize     )}}
\par\smallskip
This Python line participates in the export, validation, or report-generation pipeline. It is executed on the host PC to prepare data for RTL or to check RTL outputs.
\paragraph{Line 114}
\noindent\colorbox{gray!8}{\parbox{\dimexpr\linewidth-2\fboxsep\relax}{\raggedright\ttfamily\scriptsize     return model.to(device)}}
\par\smallskip
This Python line participates in the export, validation, or report-generation pipeline. It is executed on the host PC to prepare data for RTL or to check RTL outputs.
\paragraph{Line 115}
\noindent\colorbox{gray!8}{\parbox{\dimexpr\linewidth-2\fboxsep\relax}{\raggedright\ttfamily\scriptsize \emph{blank line}}}
\par\smallskip
This blank line is not hardware or software behavior. It separates nearby declarations or operations so the file remains readable and reviewable.
\paragraph{Line 116}
\noindent\colorbox{gray!8}{\parbox{\dimexpr\linewidth-2\fboxsep\relax}{\raggedright\ttfamily\scriptsize \emph{blank line}}}
\par\smallskip
This blank line is not hardware or software behavior. It separates nearby declarations or operations so the file remains readable and reviewable.
\paragraph{Line 117}
\noindent\colorbox{gray!8}{\parbox{\dimexpr\linewidth-2\fboxsep\relax}{\raggedright\ttfamily\scriptsize def make\_loader(args: argparse.Namespace) -> DataLoader:}}
\par\smallskip
This defines a Python function. The function packages one repeatable step such as loading a memory file, simulating the golden learning rule, validating CSVs, or generating plots.
\paragraph{Line 118}
\noindent\colorbox{gray!8}{\parbox{\dimexpr\linewidth-2\fboxsep\relax}{\raggedright\ttfamily\scriptsize     transform = create\_transform(}}
\par\smallskip
This Python line participates in the export, validation, or report-generation pipeline. It is executed on the host PC to prepare data for RTL or to check RTL outputs.
\paragraph{Line 119}
\noindent\colorbox{gray!8}{\parbox{\dimexpr\linewidth-2\fboxsep\relax}{\raggedright\ttfamily\scriptsize         input\_size=(3, 32, 32),}}
\par\smallskip
This Python line participates in the export, validation, or report-generation pipeline. It is executed on the host PC to prepare data for RTL or to check RTL outputs.
\paragraph{Line 120}
\noindent\colorbox{gray!8}{\parbox{\dimexpr\linewidth-2\fboxsep\relax}{\raggedright\ttfamily\scriptsize         is\_training=False,}}
\par\smallskip
This Python line participates in the export, validation, or report-generation pipeline. It is executed on the host PC to prepare data for RTL or to check RTL outputs.
\paragraph{Line 121}
\noindent\colorbox{gray!8}{\parbox{\dimexpr\linewidth-2\fboxsep\relax}{\raggedright\ttfamily\scriptsize         interpolation="bicubic",}}
\par\smallskip
This Python line participates in the export, validation, or report-generation pipeline. It is executed on the host PC to prepare data for RTL or to check RTL outputs.
\paragraph{Line 122}
\noindent\colorbox{gray!8}{\parbox{\dimexpr\linewidth-2\fboxsep\relax}{\raggedright\ttfamily\scriptsize         mean=CIFAR10\_MEAN,}}
\par\smallskip
This Python line participates in the export, validation, or report-generation pipeline. It is executed on the host PC to prepare data for RTL or to check RTL outputs.
\paragraph{Line 123}
\noindent\colorbox{gray!8}{\parbox{\dimexpr\linewidth-2\fboxsep\relax}{\raggedright\ttfamily\scriptsize         std=CIFAR10\_STD,}}
\par\smallskip
This Python line participates in the export, validation, or report-generation pipeline. It is executed on the host PC to prepare data for RTL or to check RTL outputs.
\paragraph{Line 124}
\noindent\colorbox{gray!8}{\parbox{\dimexpr\linewidth-2\fboxsep\relax}{\raggedright\ttfamily\scriptsize         crop\_pct=1.0,}}
\par\smallskip
This Python line participates in the export, validation, or report-generation pipeline. It is executed on the host PC to prepare data for RTL or to check RTL outputs.
\paragraph{Line 125}
\noindent\colorbox{gray!8}{\parbox{\dimexpr\linewidth-2\fboxsep\relax}{\raggedright\ttfamily\scriptsize     )}}
\par\smallskip
This Python line participates in the export, validation, or report-generation pipeline. It is executed on the host PC to prepare data for RTL or to check RTL outputs.
\paragraph{Line 126}
\noindent\colorbox{gray!8}{\parbox{\dimexpr\linewidth-2\fboxsep\relax}{\raggedright\ttfamily\scriptsize     dataset = datasets.CIFAR10(root=str(args.data\_dir), train=False, download=True, transform=transform)}}
\par\smallskip
This Python line participates in the export, validation, or report-generation pipeline. It is executed on the host PC to prepare data for RTL or to check RTL outputs.
\paragraph{Line 127}
\noindent\colorbox{gray!8}{\parbox{\dimexpr\linewidth-2\fboxsep\relax}{\raggedright\ttfamily\scriptsize     if args.samples > 0:}}
\par\smallskip
This Python line participates in the export, validation, or report-generation pipeline. It is executed on the host PC to prepare data for RTL or to check RTL outputs.
\paragraph{Line 128}
\noindent\colorbox{gray!8}{\parbox{\dimexpr\linewidth-2\fboxsep\relax}{\raggedright\ttfamily\scriptsize         dataset = Subset(dataset, range(min(args.samples, len(dataset))))}}
\par\smallskip
This Python line participates in the export, validation, or report-generation pipeline. It is executed on the host PC to prepare data for RTL or to check RTL outputs.
\paragraph{Line 129}
\noindent\colorbox{gray!8}{\parbox{\dimexpr\linewidth-2\fboxsep\relax}{\raggedright\ttfamily\scriptsize     return DataLoader(}}
\par\smallskip
This Python line participates in the export, validation, or report-generation pipeline. It is executed on the host PC to prepare data for RTL or to check RTL outputs.
\paragraph{Line 130}
\noindent\colorbox{gray!8}{\parbox{\dimexpr\linewidth-2\fboxsep\relax}{\raggedright\ttfamily\scriptsize         dataset,}}
\par\smallskip
This Python line participates in the export, validation, or report-generation pipeline. It is executed on the host PC to prepare data for RTL or to check RTL outputs.
\paragraph{Line 131}
\noindent\colorbox{gray!8}{\parbox{\dimexpr\linewidth-2\fboxsep\relax}{\raggedright\ttfamily\scriptsize         batch\_size=args.batch\_size,}}
\par\smallskip
This Python line participates in the export, validation, or report-generation pipeline. It is executed on the host PC to prepare data for RTL or to check RTL outputs.
\paragraph{Line 132}
\noindent\colorbox{gray!8}{\parbox{\dimexpr\linewidth-2\fboxsep\relax}{\raggedright\ttfamily\scriptsize         shuffle=False,}}
\par\smallskip
This Python line participates in the export, validation, or report-generation pipeline. It is executed on the host PC to prepare data for RTL or to check RTL outputs.
\paragraph{Line 133}
\noindent\colorbox{gray!8}{\parbox{\dimexpr\linewidth-2\fboxsep\relax}{\raggedright\ttfamily\scriptsize         num\_workers=args.workers,}}
\par\smallskip
This Python line participates in the export, validation, or report-generation pipeline. It is executed on the host PC to prepare data for RTL or to check RTL outputs.
\paragraph{Line 134}
\noindent\colorbox{gray!8}{\parbox{\dimexpr\linewidth-2\fboxsep\relax}{\raggedright\ttfamily\scriptsize         pin\_memory=torch.cuda.is\_available(),}}
\par\smallskip
This Python line participates in the export, validation, or report-generation pipeline. It is executed on the host PC to prepare data for RTL or to check RTL outputs.
\paragraph{Line 135}
\noindent\colorbox{gray!8}{\parbox{\dimexpr\linewidth-2\fboxsep\relax}{\raggedright\ttfamily\scriptsize         persistent\_workers=args.workers > 0,}}
\par\smallskip
This Python line participates in the export, validation, or report-generation pipeline. It is executed on the host PC to prepare data for RTL or to check RTL outputs.
\paragraph{Line 136}
\noindent\colorbox{gray!8}{\parbox{\dimexpr\linewidth-2\fboxsep\relax}{\raggedright\ttfamily\scriptsize     )}}
\par\smallskip
This Python line participates in the export, validation, or report-generation pipeline. It is executed on the host PC to prepare data for RTL or to check RTL outputs.
\paragraph{Line 137}
\noindent\colorbox{gray!8}{\parbox{\dimexpr\linewidth-2\fboxsep\relax}{\raggedright\ttfamily\scriptsize \emph{blank line}}}
\par\smallskip
This blank line is not hardware or software behavior. It separates nearby declarations or operations so the file remains readable and reviewable.
\paragraph{Line 138}
\noindent\colorbox{gray!8}{\parbox{\dimexpr\linewidth-2\fboxsep\relax}{\raggedright\ttfamily\scriptsize \emph{blank line}}}
\par\smallskip
This blank line is not hardware or software behavior. It separates nearby declarations or operations so the file remains readable and reviewable.
\paragraph{Line 139}
\noindent\colorbox{gray!8}{\parbox{\dimexpr\linewidth-2\fboxsep\relax}{\raggedright\ttfamily\scriptsize def autocast\_context(device: torch.device, enabled: bool):}}
\par\smallskip
This defines a Python function. The function packages one repeatable step such as loading a memory file, simulating the golden learning rule, validating CSVs, or generating plots.
\paragraph{Line 140}
\noindent\colorbox{gray!8}{\parbox{\dimexpr\linewidth-2\fboxsep\relax}{\raggedright\ttfamily\scriptsize     if device.type == "cuda":}}
\par\smallskip
This Python line participates in the export, validation, or report-generation pipeline. It is executed on the host PC to prepare data for RTL or to check RTL outputs.
\paragraph{Line 141}
\noindent\colorbox{gray!8}{\parbox{\dimexpr\linewidth-2\fboxsep\relax}{\raggedright\ttfamily\scriptsize         return torch.amp.autocast(device\_type="cuda", enabled=enabled)}}
\par\smallskip
This Python line participates in the export, validation, or report-generation pipeline. It is executed on the host PC to prepare data for RTL or to check RTL outputs.
\paragraph{Line 142}
\noindent\colorbox{gray!8}{\parbox{\dimexpr\linewidth-2\fboxsep\relax}{\raggedright\ttfamily\scriptsize     return nullcontext()}}
\par\smallskip
This Python line participates in the export, validation, or report-generation pipeline. It is executed on the host PC to prepare data for RTL or to check RTL outputs.
\paragraph{Line 143}
\noindent\colorbox{gray!8}{\parbox{\dimexpr\linewidth-2\fboxsep\relax}{\raggedright\ttfamily\scriptsize \emph{blank line}}}
\par\smallskip
This blank line is not hardware or software behavior. It separates nearby declarations or operations so the file remains readable and reviewable.
\paragraph{Line 144}
\noindent\colorbox{gray!8}{\parbox{\dimexpr\linewidth-2\fboxsep\relax}{\raggedright\ttfamily\scriptsize \emph{blank line}}}
\par\smallskip
This blank line is not hardware or software behavior. It separates nearby declarations or operations so the file remains readable and reviewable.
\paragraph{Line 145}
\noindent\colorbox{gray!8}{\parbox{\dimexpr\linewidth-2\fboxsep\relax}{\raggedright\ttfamily\scriptsize def collect\_features(args: argparse.Namespace) -> tuple[torch.Tensor, torch.Tensor, torch.Tensor]:}}
\par\smallskip
This defines a Python function. The function packages one repeatable step such as loading a memory file, simulating the golden learning rule, validating CSVs, or generating plots.
\paragraph{Line 146}
\noindent\colorbox{gray!8}{\parbox{\dimexpr\linewidth-2\fboxsep\relax}{\raggedright\ttfamily\scriptsize     add\_runtime\_dlls()}}
\par\smallskip
This Python line participates in the export, validation, or report-generation pipeline. It is executed on the host PC to prepare data for RTL or to check RTL outputs.
\paragraph{Line 147}
\noindent\colorbox{gray!8}{\parbox{\dimexpr\linewidth-2\fboxsep\relax}{\raggedright\ttfamily\scriptsize     import\_upstream\_model()}}
\par\smallskip
This Python line participates in the export, validation, or report-generation pipeline. It is executed on the host PC to prepare data for RTL or to check RTL outputs.
\paragraph{Line 148}
\noindent\colorbox{gray!8}{\parbox{\dimexpr\linewidth-2\fboxsep\relax}{\raggedright\ttfamily\scriptsize     from spikingjelly.clock\_driven import functional}}
\par\smallskip
This imports a Python library or symbol. The scripts use standard parsing, numeric arrays, plotting, and file APIs to make the hardware artifacts reproducible.
\paragraph{Line 149}
\noindent\colorbox{gray!8}{\parbox{\dimexpr\linewidth-2\fboxsep\relax}{\raggedright\ttfamily\scriptsize \emph{blank line}}}
\par\smallskip
This blank line is not hardware or software behavior. It separates nearby declarations or operations so the file remains readable and reviewable.
\paragraph{Line 150}
\noindent\colorbox{gray!8}{\parbox{\dimexpr\linewidth-2\fboxsep\relax}{\raggedright\ttfamily\scriptsize     device\_name = "cuda" if args.device == "auto" and torch.cuda.is\_available() else ("cpu" if args.device == "auto" else args.device)}}
\par\smallskip
This Python line participates in the export, validation, or report-generation pipeline. It is executed on the host PC to prepare data for RTL or to check RTL outputs.
\paragraph{Line 151}
\noindent\colorbox{gray!8}{\parbox{\dimexpr\linewidth-2\fboxsep\relax}{\raggedright\ttfamily\scriptsize     device = torch.device(device\_name)}}
\par\smallskip
This Python line participates in the export, validation, or report-generation pipeline. It is executed on the host PC to prepare data for RTL or to check RTL outputs.
\paragraph{Line 152}
\noindent\colorbox{gray!8}{\parbox{\dimexpr\linewidth-2\fboxsep\relax}{\raggedright\ttfamily\scriptsize     if device.type == "cuda":}}
\par\smallskip
This Python line participates in the export, validation, or report-generation pipeline. It is executed on the host PC to prepare data for RTL or to check RTL outputs.
\paragraph{Line 153}
\noindent\colorbox{gray!8}{\parbox{\dimexpr\linewidth-2\fboxsep\relax}{\raggedright\ttfamily\scriptsize         torch.backends.cudnn.benchmark = True}}
\par\smallskip
This Python line participates in the export, validation, or report-generation pipeline. It is executed on the host PC to prepare data for RTL or to check RTL outputs.
\paragraph{Line 154}
\noindent\colorbox{gray!8}{\parbox{\dimexpr\linewidth-2\fboxsep\relax}{\raggedright\ttfamily\scriptsize     model = build\_model(device)}}
\par\smallskip
This Python line participates in the export, validation, or report-generation pipeline. It is executed on the host PC to prepare data for RTL or to check RTL outputs.
\paragraph{Line 155}
\noindent\colorbox{gray!8}{\parbox{\dimexpr\linewidth-2\fboxsep\relax}{\raggedright\ttfamily\scriptsize     model.load\_state\_dict(load\_checkpoint(args.checkpoint), strict=True)}}
\par\smallskip
This Python line participates in the export, validation, or report-generation pipeline. It is executed on the host PC to prepare data for RTL or to check RTL outputs.
\paragraph{Line 156}
\noindent\colorbox{gray!8}{\parbox{\dimexpr\linewidth-2\fboxsep\relax}{\raggedright\ttfamily\scriptsize     model.eval()}}
\par\smallskip
This Python line participates in the export, validation, or report-generation pipeline. It is executed on the host PC to prepare data for RTL or to check RTL outputs.
\paragraph{Line 157}
\noindent\colorbox{gray!8}{\parbox{\dimexpr\linewidth-2\fboxsep\relax}{\raggedright\ttfamily\scriptsize     loader = make\_loader(args)}}
\par\smallskip
This Python line participates in the export, validation, or report-generation pipeline. It is executed on the host PC to prepare data for RTL or to check RTL outputs.
\paragraph{Line 158}
\noindent\colorbox{gray!8}{\parbox{\dimexpr\linewidth-2\fboxsep\relax}{\raggedright\ttfamily\scriptsize \emph{blank line}}}
\par\smallskip
This blank line is not hardware or software behavior. It separates nearby declarations or operations so the file remains readable and reviewable.
\paragraph{Line 159}
\noindent\colorbox{gray!8}{\parbox{\dimexpr\linewidth-2\fboxsep\relax}{\raggedright\ttfamily\scriptsize     all\_features: list[torch.Tensor] = []}}
\par\smallskip
This Python line participates in the export, validation, or report-generation pipeline. It is executed on the host PC to prepare data for RTL or to check RTL outputs.
\paragraph{Line 160}
\noindent\colorbox{gray!8}{\parbox{\dimexpr\linewidth-2\fboxsep\relax}{\raggedright\ttfamily\scriptsize     all\_logits: list[torch.Tensor] = []}}
\par\smallskip
This Python line participates in the export, validation, or report-generation pipeline. It is executed on the host PC to prepare data for RTL or to check RTL outputs.
\paragraph{Line 161}
\noindent\colorbox{gray!8}{\parbox{\dimexpr\linewidth-2\fboxsep\relax}{\raggedright\ttfamily\scriptsize     all\_labels: list[torch.Tensor] = []}}
\par\smallskip
This Python line participates in the export, validation, or report-generation pipeline. It is executed on the host PC to prepare data for RTL or to check RTL outputs.
\paragraph{Line 162}
\noindent\colorbox{gray!8}{\parbox{\dimexpr\linewidth-2\fboxsep\relax}{\raggedright\ttfamily\scriptsize     with torch.no\_grad():}}
\par\smallskip
This Python line participates in the export, validation, or report-generation pipeline. It is executed on the host PC to prepare data for RTL or to check RTL outputs.
\paragraph{Line 163}
\noindent\colorbox{gray!8}{\parbox{\dimexpr\linewidth-2\fboxsep\relax}{\raggedright\ttfamily\scriptsize         for images, labels in loader:}}
\par\smallskip
This Python line participates in the export, validation, or report-generation pipeline. It is executed on the host PC to prepare data for RTL or to check RTL outputs.
\paragraph{Line 164}
\noindent\colorbox{gray!8}{\parbox{\dimexpr\linewidth-2\fboxsep\relax}{\raggedright\ttfamily\scriptsize             images = images.to(device, non\_blocking=True)}}
\par\smallskip
This Python line participates in the export, validation, or report-generation pipeline. It is executed on the host PC to prepare data for RTL or to check RTL outputs.
\paragraph{Line 165}
\noindent\colorbox{gray!8}{\parbox{\dimexpr\linewidth-2\fboxsep\relax}{\raggedright\ttfamily\scriptsize             with autocast\_context(device, args.amp):}}
\par\smallskip
This Python line participates in the export, validation, or report-generation pipeline. It is executed on the host PC to prepare data for RTL or to check RTL outputs.
\paragraph{Line 166}
\noindent\colorbox{gray!8}{\parbox{\dimexpr\linewidth-2\fboxsep\relax}{\raggedright\ttfamily\scriptsize                 time\_images = images.unsqueeze(0).repeat(model.T, 1, 1, 1, 1)}}
\par\smallskip
This Python line participates in the export, validation, or report-generation pipeline. It is executed on the host PC to prepare data for RTL or to check RTL outputs.
\paragraph{Line 167}
\noindent\colorbox{gray!8}{\parbox{\dimexpr\linewidth-2\fboxsep\relax}{\raggedright\ttfamily\scriptsize                 time\_features = model.forward\_features(time\_images)}}
\par\smallskip
This Python line participates in the export, validation, or report-generation pipeline. It is executed on the host PC to prepare data for RTL or to check RTL outputs.
\paragraph{Line 168}
\noindent\colorbox{gray!8}{\parbox{\dimexpr\linewidth-2\fboxsep\relax}{\raggedright\ttfamily\scriptsize                 features = time\_features.mean(0)}}
\par\smallskip
This Python line participates in the export, validation, or report-generation pipeline. It is executed on the host PC to prepare data for RTL or to check RTL outputs.
\paragraph{Line 169}
\noindent\colorbox{gray!8}{\parbox{\dimexpr\linewidth-2\fboxsep\relax}{\raggedright\ttfamily\scriptsize                 logits = model.head(features)}}
\par\smallskip
This Python line participates in the export, validation, or report-generation pipeline. It is executed on the host PC to prepare data for RTL or to check RTL outputs.
\paragraph{Line 170}
\noindent\colorbox{gray!8}{\parbox{\dimexpr\linewidth-2\fboxsep\relax}{\raggedright\ttfamily\scriptsize             functional.reset\_net(model)}}
\par\smallskip
This Python line participates in the export, validation, or report-generation pipeline. It is executed on the host PC to prepare data for RTL or to check RTL outputs.
\paragraph{Line 171}
\noindent\colorbox{gray!8}{\parbox{\dimexpr\linewidth-2\fboxsep\relax}{\raggedright\ttfamily\scriptsize             all\_features.append(features.detach().float().cpu())}}
\par\smallskip
This Python line participates in the export, validation, or report-generation pipeline. It is executed on the host PC to prepare data for RTL or to check RTL outputs.
\paragraph{Line 172}
\noindent\colorbox{gray!8}{\parbox{\dimexpr\linewidth-2\fboxsep\relax}{\raggedright\ttfamily\scriptsize             all\_logits.append(logits.detach().float().cpu())}}
\par\smallskip
This Python line participates in the export, validation, or report-generation pipeline. It is executed on the host PC to prepare data for RTL or to check RTL outputs.
\paragraph{Line 173}
\noindent\colorbox{gray!8}{\parbox{\dimexpr\linewidth-2\fboxsep\relax}{\raggedright\ttfamily\scriptsize             all\_labels.append(labels.detach().cpu())}}
\par\smallskip
This Python line participates in the export, validation, or report-generation pipeline. It is executed on the host PC to prepare data for RTL or to check RTL outputs.
\paragraph{Line 174}
\noindent\colorbox{gray!8}{\parbox{\dimexpr\linewidth-2\fboxsep\relax}{\raggedright\ttfamily\scriptsize     return torch.cat(all\_features), torch.cat(all\_logits), torch.cat(all\_labels)}}
\par\smallskip
This Python line participates in the export, validation, or report-generation pipeline. It is executed on the host PC to prepare data for RTL or to check RTL outputs.
\paragraph{Line 175}
\noindent\colorbox{gray!8}{\parbox{\dimexpr\linewidth-2\fboxsep\relax}{\raggedright\ttfamily\scriptsize \emph{blank line}}}
\par\smallskip
This blank line is not hardware or software behavior. It separates nearby declarations or operations so the file remains readable and reviewable.
\paragraph{Line 176}
\noindent\colorbox{gray!8}{\parbox{\dimexpr\linewidth-2\fboxsep\relax}{\raggedright\ttfamily\scriptsize \emph{blank line}}}
\par\smallskip
This blank line is not hardware or software behavior. It separates nearby declarations or operations so the file remains readable and reviewable.
\paragraph{Line 177}
\noindent\colorbox{gray!8}{\parbox{\dimexpr\linewidth-2\fboxsep\relax}{\raggedright\ttfamily\scriptsize def fold\_batch\_norm(weight: torch.Tensor, bias: torch.Tensor | None, bn\_prefix: str, state: dict[str, torch.Tensor]) -> tuple[torch.Tensor, torch.Tensor]:}}
\par\smallskip
This defines a Python function. The function packages one repeatable step such as loading a memory file, simulating the golden learning rule, validating CSVs, or generating plots.
\paragraph{Line 178}
\noindent\colorbox{gray!8}{\parbox{\dimexpr\linewidth-2\fboxsep\relax}{\raggedright\ttfamily\scriptsize     out\_channels = weight.shape[0]}}
\par\smallskip
This Python line participates in the export, validation, or report-generation pipeline. It is executed on the host PC to prepare data for RTL or to check RTL outputs.
\paragraph{Line 179}
\noindent\colorbox{gray!8}{\parbox{\dimexpr\linewidth-2\fboxsep\relax}{\raggedright\ttfamily\scriptsize     if bias is None:}}
\par\smallskip
This line handles the classifier parameters. Those arrays mirror the RTL memories, so mismatches reveal an export, quantization, or hardware-update bug.
\paragraph{Line 180}
\noindent\colorbox{gray!8}{\parbox{\dimexpr\linewidth-2\fboxsep\relax}{\raggedright\ttfamily\scriptsize         bias = torch.zeros(out\_channels, dtype=weight.dtype)}}
\par\smallskip
This line handles the classifier parameters. Those arrays mirror the RTL memories, so mismatches reveal an export, quantization, or hardware-update bug.
\paragraph{Line 181}
\noindent\colorbox{gray!8}{\parbox{\dimexpr\linewidth-2\fboxsep\relax}{\raggedright\ttfamily\scriptsize     gamma = state[f"\{bn\_prefix\}.weight"]}}
\par\smallskip
This Python line participates in the export, validation, or report-generation pipeline. It is executed on the host PC to prepare data for RTL or to check RTL outputs.
\paragraph{Line 182}
\noindent\colorbox{gray!8}{\parbox{\dimexpr\linewidth-2\fboxsep\relax}{\raggedright\ttfamily\scriptsize     beta = state[f"\{bn\_prefix\}.bias"]}}
\par\smallskip
This line handles the classifier parameters. Those arrays mirror the RTL memories, so mismatches reveal an export, quantization, or hardware-update bug.
\paragraph{Line 183}
\noindent\colorbox{gray!8}{\parbox{\dimexpr\linewidth-2\fboxsep\relax}{\raggedright\ttfamily\scriptsize     mean = state[f"\{bn\_prefix\}.running\_mean"]}}
\par\smallskip
This Python line participates in the export, validation, or report-generation pipeline. It is executed on the host PC to prepare data for RTL or to check RTL outputs.
\paragraph{Line 184}
\noindent\colorbox{gray!8}{\parbox{\dimexpr\linewidth-2\fboxsep\relax}{\raggedright\ttfamily\scriptsize     var = state[f"\{bn\_prefix\}.running\_var"]}}
\par\smallskip
This Python line participates in the export, validation, or report-generation pipeline. It is executed on the host PC to prepare data for RTL or to check RTL outputs.
\paragraph{Line 185}
\noindent\colorbox{gray!8}{\parbox{\dimexpr\linewidth-2\fboxsep\relax}{\raggedright\ttfamily\scriptsize     scale = gamma / torch.sqrt(var + 1.0e-5)}}
\par\smallskip
This Python line participates in the export, validation, or report-generation pipeline. It is executed on the host PC to prepare data for RTL or to check RTL outputs.
\paragraph{Line 186}
\noindent\colorbox{gray!8}{\parbox{\dimexpr\linewidth-2\fboxsep\relax}{\raggedright\ttfamily\scriptsize     folded\_weight = weight * scale.reshape(-1, *([1] * (weight.ndim - 1)))}}
\par\smallskip
This Python line participates in the export, validation, or report-generation pipeline. It is executed on the host PC to prepare data for RTL or to check RTL outputs.
\paragraph{Line 187}
\noindent\colorbox{gray!8}{\parbox{\dimexpr\linewidth-2\fboxsep\relax}{\raggedright\ttfamily\scriptsize     folded\_bias = (bias - mean) * scale + beta}}
\par\smallskip
This line handles the classifier parameters. Those arrays mirror the RTL memories, so mismatches reveal an export, quantization, or hardware-update bug.
\paragraph{Line 188}
\noindent\colorbox{gray!8}{\parbox{\dimexpr\linewidth-2\fboxsep\relax}{\raggedright\ttfamily\scriptsize     return folded\_weight, folded\_bias}}
\par\smallskip
This line handles the classifier parameters. Those arrays mirror the RTL memories, so mismatches reveal an export, quantization, or hardware-update bug.
\paragraph{Line 189}
\noindent\colorbox{gray!8}{\parbox{\dimexpr\linewidth-2\fboxsep\relax}{\raggedright\ttfamily\scriptsize \emph{blank line}}}
\par\smallskip
This blank line is not hardware or software behavior. It separates nearby declarations or operations so the file remains readable and reviewable.
\paragraph{Line 190}
\noindent\colorbox{gray!8}{\parbox{\dimexpr\linewidth-2\fboxsep\relax}{\raggedright\ttfamily\scriptsize \emph{blank line}}}
\par\smallskip
This blank line is not hardware or software behavior. It separates nearby declarations or operations so the file remains readable and reviewable.
\paragraph{Line 191}
\noindent\colorbox{gray!8}{\parbox{\dimexpr\linewidth-2\fboxsep\relax}{\raggedright\ttfamily\scriptsize def tensor\_stats(tensor: torch.Tensor) -> dict[str, float | list[int]]:}}
\par\smallskip
This defines a Python function. The function packages one repeatable step such as loading a memory file, simulating the golden learning rule, validating CSVs, or generating plots.
\paragraph{Line 192}
\noindent\colorbox{gray!8}{\parbox{\dimexpr\linewidth-2\fboxsep\relax}{\raggedright\ttfamily\scriptsize     t = tensor.detach().float()}}
\par\smallskip
This Python line participates in the export, validation, or report-generation pipeline. It is executed on the host PC to prepare data for RTL or to check RTL outputs.
\paragraph{Line 193}
\noindent\colorbox{gray!8}{\parbox{\dimexpr\linewidth-2\fboxsep\relax}{\raggedright\ttfamily\scriptsize     return \{}}
\par\smallskip
This Python line participates in the export, validation, or report-generation pipeline. It is executed on the host PC to prepare data for RTL or to check RTL outputs.
\paragraph{Line 194}
\noindent\colorbox{gray!8}{\parbox{\dimexpr\linewidth-2\fboxsep\relax}{\raggedright\ttfamily\scriptsize         "shape": list(t.shape),}}
\par\smallskip
This Python line participates in the export, validation, or report-generation pipeline. It is executed on the host PC to prepare data for RTL or to check RTL outputs.
\paragraph{Line 195}
\noindent\colorbox{gray!8}{\parbox{\dimexpr\linewidth-2\fboxsep\relax}{\raggedright\ttfamily\scriptsize         "min": float(t.min()),}}
\par\smallskip
This Python line participates in the export, validation, or report-generation pipeline. It is executed on the host PC to prepare data for RTL or to check RTL outputs.
\paragraph{Line 196}
\noindent\colorbox{gray!8}{\parbox{\dimexpr\linewidth-2\fboxsep\relax}{\raggedright\ttfamily\scriptsize         "max": float(t.max()),}}
\par\smallskip
This Python line participates in the export, validation, or report-generation pipeline. It is executed on the host PC to prepare data for RTL or to check RTL outputs.
\paragraph{Line 197}
\noindent\colorbox{gray!8}{\parbox{\dimexpr\linewidth-2\fboxsep\relax}{\raggedright\ttfamily\scriptsize         "mean": float(t.mean()),}}
\par\smallskip
This Python line participates in the export, validation, or report-generation pipeline. It is executed on the host PC to prepare data for RTL or to check RTL outputs.
\paragraph{Line 198}
\noindent\colorbox{gray!8}{\parbox{\dimexpr\linewidth-2\fboxsep\relax}{\raggedright\ttfamily\scriptsize         "std": float(t.std(unbiased=False)),}}
\par\smallskip
This line handles the classifier parameters. Those arrays mirror the RTL memories, so mismatches reveal an export, quantization, or hardware-update bug.
\paragraph{Line 199}
\noindent\colorbox{gray!8}{\parbox{\dimexpr\linewidth-2\fboxsep\relax}{\raggedright\ttfamily\scriptsize     \}}}
\par\smallskip
This Python line participates in the export, validation, or report-generation pipeline. It is executed on the host PC to prepare data for RTL or to check RTL outputs.
\paragraph{Line 200}
\noindent\colorbox{gray!8}{\parbox{\dimexpr\linewidth-2\fboxsep\relax}{\raggedright\ttfamily\scriptsize \emph{blank line}}}
\par\smallskip
This blank line is not hardware or software behavior. It separates nearby declarations or operations so the file remains readable and reviewable.
\paragraph{Line 201}
\noindent\colorbox{gray!8}{\parbox{\dimexpr\linewidth-2\fboxsep\relax}{\raggedright\ttfamily\scriptsize \emph{blank line}}}
\par\smallskip
This blank line is not hardware or software behavior. It separates nearby declarations or operations so the file remains readable and reviewable.
\paragraph{Line 202}
\noindent\colorbox{gray!8}{\parbox{\dimexpr\linewidth-2\fboxsep\relax}{\raggedright\ttfamily\scriptsize def full\_checkpoint\_mapping(args: argparse.Namespace) -> dict[str, object]:}}
\par\smallskip
This defines a Python function. The function packages one repeatable step such as loading a memory file, simulating the golden learning rule, validating CSVs, or generating plots.
\paragraph{Line 203}
\noindent\colorbox{gray!8}{\parbox{\dimexpr\linewidth-2\fboxsep\relax}{\raggedright\ttfamily\scriptsize     state = load\_checkpoint(args.checkpoint)}}
\par\smallskip
This Python line participates in the export, validation, or report-generation pipeline. It is executed on the host PC to prepare data for RTL or to check RTL outputs.
\paragraph{Line 204}
\noindent\colorbox{gray!8}{\parbox{\dimexpr\linewidth-2\fboxsep\relax}{\raggedright\ttfamily\scriptsize     layers: list[dict[str, object]] = []}}
\par\smallskip
This Python line participates in the export, validation, or report-generation pipeline. It is executed on the host PC to prepare data for RTL or to check RTL outputs.
\paragraph{Line 205}
\noindent\colorbox{gray!8}{\parbox{\dimexpr\linewidth-2\fboxsep\relax}{\raggedright\ttfamily\scriptsize \emph{blank line}}}
\par\smallskip
This blank line is not hardware or software behavior. It separates nearby declarations or operations so the file remains readable and reviewable.
\paragraph{Line 206}
\noindent\colorbox{gray!8}{\parbox{\dimexpr\linewidth-2\fboxsep\relax}{\raggedright\ttfamily\scriptsize     conv\_bn\_pairs = [}}
\par\smallskip
This Python line participates in the export, validation, or report-generation pipeline. It is executed on the host PC to prepare data for RTL or to check RTL outputs.
\paragraph{Line 207}
\noindent\colorbox{gray!8}{\parbox{\dimexpr\linewidth-2\fboxsep\relax}{\raggedright\ttfamily\scriptsize         ("patch\_embed.block0\_conv", "patch\_embed.block0\_bn"),}}
\par\smallskip
This Python line participates in the export, validation, or report-generation pipeline. It is executed on the host PC to prepare data for RTL or to check RTL outputs.
\paragraph{Line 208}
\noindent\colorbox{gray!8}{\parbox{\dimexpr\linewidth-2\fboxsep\relax}{\raggedright\ttfamily\scriptsize         ("patch\_embed.block1\_conv", "patch\_embed.block1\_bn"),}}
\par\smallskip
This Python line participates in the export, validation, or report-generation pipeline. It is executed on the host PC to prepare data for RTL or to check RTL outputs.
\paragraph{Line 209}
\noindent\colorbox{gray!8}{\parbox{\dimexpr\linewidth-2\fboxsep\relax}{\raggedright\ttfamily\scriptsize         ("patch\_embed.block2\_conv", "patch\_embed.block2\_bn"),}}
\par\smallskip
This Python line participates in the export, validation, or report-generation pipeline. It is executed on the host PC to prepare data for RTL or to check RTL outputs.
\paragraph{Line 210}
\noindent\colorbox{gray!8}{\parbox{\dimexpr\linewidth-2\fboxsep\relax}{\raggedright\ttfamily\scriptsize         ("patch\_embed.block3\_conv", "patch\_embed.block3\_bn"),}}
\par\smallskip
This Python line participates in the export, validation, or report-generation pipeline. It is executed on the host PC to prepare data for RTL or to check RTL outputs.
\paragraph{Line 211}
\noindent\colorbox{gray!8}{\parbox{\dimexpr\linewidth-2\fboxsep\relax}{\raggedright\ttfamily\scriptsize         ("patch\_embed.block4\_conv", "patch\_embed.block4\_bn"),}}
\par\smallskip
This Python line participates in the export, validation, or report-generation pipeline. It is executed on the host PC to prepare data for RTL or to check RTL outputs.
\paragraph{Line 212}
\noindent\colorbox{gray!8}{\parbox{\dimexpr\linewidth-2\fboxsep\relax}{\raggedright\ttfamily\scriptsize     ]}}
\par\smallskip
This Python line participates in the export, validation, or report-generation pipeline. It is executed on the host PC to prepare data for RTL or to check RTL outputs.
\paragraph{Line 213}
\noindent\colorbox{gray!8}{\parbox{\dimexpr\linewidth-2\fboxsep\relax}{\raggedright\ttfamily\scriptsize     for block in range(4):}}
\par\smallskip
This Python line participates in the export, validation, or report-generation pipeline. It is executed on the host PC to prepare data for RTL or to check RTL outputs.
\paragraph{Line 214}
\noindent\colorbox{gray!8}{\parbox{\dimexpr\linewidth-2\fboxsep\relax}{\raggedright\ttfamily\scriptsize         base = f"block.\{block\}"}}
\par\smallskip
This Python line participates in the export, validation, or report-generation pipeline. It is executed on the host PC to prepare data for RTL or to check RTL outputs.
\paragraph{Line 215}
\noindent\colorbox{gray!8}{\parbox{\dimexpr\linewidth-2\fboxsep\relax}{\raggedright\ttfamily\scriptsize         conv\_bn\_pairs.extend(}}
\par\smallskip
This Python line participates in the export, validation, or report-generation pipeline. It is executed on the host PC to prepare data for RTL or to check RTL outputs.
\paragraph{Line 216}
\noindent\colorbox{gray!8}{\parbox{\dimexpr\linewidth-2\fboxsep\relax}{\raggedright\ttfamily\scriptsize             [}}
\par\smallskip
This Python line participates in the export, validation, or report-generation pipeline. It is executed on the host PC to prepare data for RTL or to check RTL outputs.
\paragraph{Line 217}
\noindent\colorbox{gray!8}{\parbox{\dimexpr\linewidth-2\fboxsep\relax}{\raggedright\ttfamily\scriptsize                 (f"\{base\}.attn.q\_conv", f"\{base\}.attn.q\_bn"),}}
\par\smallskip
This Python line participates in the export, validation, or report-generation pipeline. It is executed on the host PC to prepare data for RTL or to check RTL outputs.
\paragraph{Line 218}
\noindent\colorbox{gray!8}{\parbox{\dimexpr\linewidth-2\fboxsep\relax}{\raggedright\ttfamily\scriptsize                 (f"\{base\}.attn.k\_conv", f"\{base\}.attn.k\_bn"),}}
\par\smallskip
This Python line participates in the export, validation, or report-generation pipeline. It is executed on the host PC to prepare data for RTL or to check RTL outputs.
\paragraph{Line 219}
\noindent\colorbox{gray!8}{\parbox{\dimexpr\linewidth-2\fboxsep\relax}{\raggedright\ttfamily\scriptsize                 (f"\{base\}.attn.v\_conv", f"\{base\}.attn.v\_bn"),}}
\par\smallskip
This Python line participates in the export, validation, or report-generation pipeline. It is executed on the host PC to prepare data for RTL or to check RTL outputs.
\paragraph{Line 220}
\noindent\colorbox{gray!8}{\parbox{\dimexpr\linewidth-2\fboxsep\relax}{\raggedright\ttfamily\scriptsize                 (f"\{base\}.attn.proj\_conv", f"\{base\}.attn.proj\_bn"),}}
\par\smallskip
This Python line participates in the export, validation, or report-generation pipeline. It is executed on the host PC to prepare data for RTL or to check RTL outputs.
\paragraph{Line 221}
\noindent\colorbox{gray!8}{\parbox{\dimexpr\linewidth-2\fboxsep\relax}{\raggedright\ttfamily\scriptsize                 (f"\{base\}.mlp.mlp1\_conv", f"\{base\}.mlp.mlp1\_bn"),}}
\par\smallskip
This Python line participates in the export, validation, or report-generation pipeline. It is executed on the host PC to prepare data for RTL or to check RTL outputs.
\paragraph{Line 222}
\noindent\colorbox{gray!8}{\parbox{\dimexpr\linewidth-2\fboxsep\relax}{\raggedright\ttfamily\scriptsize                 (f"\{base\}.mlp.mlp2\_conv", f"\{base\}.mlp.mlp2\_bn"),}}
\par\smallskip
This Python line participates in the export, validation, or report-generation pipeline. It is executed on the host PC to prepare data for RTL or to check RTL outputs.
\paragraph{Line 223}
\noindent\colorbox{gray!8}{\parbox{\dimexpr\linewidth-2\fboxsep\relax}{\raggedright\ttfamily\scriptsize             ]}}
\par\smallskip
This Python line participates in the export, validation, or report-generation pipeline. It is executed on the host PC to prepare data for RTL or to check RTL outputs.
\paragraph{Line 224}
\noindent\colorbox{gray!8}{\parbox{\dimexpr\linewidth-2\fboxsep\relax}{\raggedright\ttfamily\scriptsize         )}}
\par\smallskip
This Python line participates in the export, validation, or report-generation pipeline. It is executed on the host PC to prepare data for RTL or to check RTL outputs.
\paragraph{Line 225}
\noindent\colorbox{gray!8}{\parbox{\dimexpr\linewidth-2\fboxsep\relax}{\raggedright\ttfamily\scriptsize \emph{blank line}}}
\par\smallskip
This blank line is not hardware or software behavior. It separates nearby declarations or operations so the file remains readable and reviewable.
\paragraph{Line 226}
\noindent\colorbox{gray!8}{\parbox{\dimexpr\linewidth-2\fboxsep\relax}{\raggedright\ttfamily\scriptsize     total\_weights = 0}}
\par\smallskip
This line handles the classifier parameters. Those arrays mirror the RTL memories, so mismatches reveal an export, quantization, or hardware-update bug.
\paragraph{Line 227}
\noindent\colorbox{gray!8}{\parbox{\dimexpr\linewidth-2\fboxsep\relax}{\raggedright\ttfamily\scriptsize     total\_biases = 0}}
\par\smallskip
This line handles the classifier parameters. Those arrays mirror the RTL memories, so mismatches reveal an export, quantization, or hardware-update bug.
\paragraph{Line 228}
\noindent\colorbox{gray!8}{\parbox{\dimexpr\linewidth-2\fboxsep\relax}{\raggedright\ttfamily\scriptsize     for conv\_prefix, bn\_prefix in conv\_bn\_pairs:}}
\par\smallskip
This Python line participates in the export, validation, or report-generation pipeline. It is executed on the host PC to prepare data for RTL or to check RTL outputs.
\paragraph{Line 229}
\noindent\colorbox{gray!8}{\parbox{\dimexpr\linewidth-2\fboxsep\relax}{\raggedright\ttfamily\scriptsize         weight = state[f"\{conv\_prefix\}.weight"].float()}}
\par\smallskip
This Python line participates in the export, validation, or report-generation pipeline. It is executed on the host PC to prepare data for RTL or to check RTL outputs.
\paragraph{Line 230}
\noindent\colorbox{gray!8}{\parbox{\dimexpr\linewidth-2\fboxsep\relax}{\raggedright\ttfamily\scriptsize         bias = state.get(f"\{conv\_prefix\}.bias")}}
\par\smallskip
This line handles the classifier parameters. Those arrays mirror the RTL memories, so mismatches reveal an export, quantization, or hardware-update bug.
\paragraph{Line 231}
\noindent\colorbox{gray!8}{\parbox{\dimexpr\linewidth-2\fboxsep\relax}{\raggedright\ttfamily\scriptsize         if bias is not None:}}
\par\smallskip
This line handles the classifier parameters. Those arrays mirror the RTL memories, so mismatches reveal an export, quantization, or hardware-update bug.
\paragraph{Line 232}
\noindent\colorbox{gray!8}{\parbox{\dimexpr\linewidth-2\fboxsep\relax}{\raggedright\ttfamily\scriptsize             bias = bias.float()}}
\par\smallskip
This line handles the classifier parameters. Those arrays mirror the RTL memories, so mismatches reveal an export, quantization, or hardware-update bug.
\paragraph{Line 233}
\noindent\colorbox{gray!8}{\parbox{\dimexpr\linewidth-2\fboxsep\relax}{\raggedright\ttfamily\scriptsize         folded\_weight, folded\_bias = fold\_batch\_norm(weight, bias, bn\_prefix, state)}}
\par\smallskip
This line handles the classifier parameters. Those arrays mirror the RTL memories, so mismatches reveal an export, quantization, or hardware-update bug.
\paragraph{Line 234}
\noindent\colorbox{gray!8}{\parbox{\dimexpr\linewidth-2\fboxsep\relax}{\raggedright\ttfamily\scriptsize         total\_weights += folded\_weight.numel()}}
\par\smallskip
This line handles the classifier parameters. Those arrays mirror the RTL memories, so mismatches reveal an export, quantization, or hardware-update bug.
\paragraph{Line 235}
\noindent\colorbox{gray!8}{\parbox{\dimexpr\linewidth-2\fboxsep\relax}{\raggedright\ttfamily\scriptsize         total\_biases += folded\_bias.numel()}}
\par\smallskip
This line handles the classifier parameters. Those arrays mirror the RTL memories, so mismatches reveal an export, quantization, or hardware-update bug.
\paragraph{Line 236}
\noindent\colorbox{gray!8}{\parbox{\dimexpr\linewidth-2\fboxsep\relax}{\raggedright\ttfamily\scriptsize         layers.append(}}
\par\smallskip
This Python line participates in the export, validation, or report-generation pipeline. It is executed on the host PC to prepare data for RTL or to check RTL outputs.
\paragraph{Line 237}
\noindent\colorbox{gray!8}{\parbox{\dimexpr\linewidth-2\fboxsep\relax}{\raggedright\ttfamily\scriptsize             \{}}
\par\smallskip
This Python line participates in the export, validation, or report-generation pipeline. It is executed on the host PC to prepare data for RTL or to check RTL outputs.
\paragraph{Line 238}
\noindent\colorbox{gray!8}{\parbox{\dimexpr\linewidth-2\fboxsep\relax}{\raggedright\ttfamily\scriptsize                 "name": conv\_prefix,}}
\par\smallskip
This Python line participates in the export, validation, or report-generation pipeline. It is executed on the host PC to prepare data for RTL or to check RTL outputs.
\paragraph{Line 239}
\noindent\colorbox{gray!8}{\parbox{\dimexpr\linewidth-2\fboxsep\relax}{\raggedright\ttfamily\scriptsize                 "kind": "conv\_or\_pointwise\_with\_batchnorm\_folded",}}
\par\smallskip
This Python line participates in the export, validation, or report-generation pipeline. It is executed on the host PC to prepare data for RTL or to check RTL outputs.
\paragraph{Line 240}
\noindent\colorbox{gray!8}{\parbox{\dimexpr\linewidth-2\fboxsep\relax}{\raggedright\ttfamily\scriptsize                 "weight": tensor\_stats(folded\_weight),}}
\par\smallskip
This Python line participates in the export, validation, or report-generation pipeline. It is executed on the host PC to prepare data for RTL or to check RTL outputs.
\paragraph{Line 241}
\noindent\colorbox{gray!8}{\parbox{\dimexpr\linewidth-2\fboxsep\relax}{\raggedright\ttfamily\scriptsize                 "bias": tensor\_stats(folded\_bias),}}
\par\smallskip
This line handles the classifier parameters. Those arrays mirror the RTL memories, so mismatches reveal an export, quantization, or hardware-update bug.
\paragraph{Line 242}
\noindent\colorbox{gray!8}{\parbox{\dimexpr\linewidth-2\fboxsep\relax}{\raggedright\ttfamily\scriptsize                 "weight\_values": int(folded\_weight.numel()),}}
\par\smallskip
This Python line participates in the export, validation, or report-generation pipeline. It is executed on the host PC to prepare data for RTL or to check RTL outputs.
\paragraph{Line 243}
\noindent\colorbox{gray!8}{\parbox{\dimexpr\linewidth-2\fboxsep\relax}{\raggedright\ttfamily\scriptsize                 "bias\_values": int(folded\_bias.numel()),}}
\par\smallskip
This line handles the classifier parameters. Those arrays mirror the RTL memories, so mismatches reveal an export, quantization, or hardware-update bug.
\paragraph{Line 244}
\noindent\colorbox{gray!8}{\parbox{\dimexpr\linewidth-2\fboxsep\relax}{\raggedright\ttfamily\scriptsize             \}}}
\par\smallskip
This Python line participates in the export, validation, or report-generation pipeline. It is executed on the host PC to prepare data for RTL or to check RTL outputs.
\paragraph{Line 245}
\noindent\colorbox{gray!8}{\parbox{\dimexpr\linewidth-2\fboxsep\relax}{\raggedright\ttfamily\scriptsize         )}}
\par\smallskip
This Python line participates in the export, validation, or report-generation pipeline. It is executed on the host PC to prepare data for RTL or to check RTL outputs.
\paragraph{Line 246}
\noindent\colorbox{gray!8}{\parbox{\dimexpr\linewidth-2\fboxsep\relax}{\raggedright\ttfamily\scriptsize \emph{blank line}}}
\par\smallskip
This blank line is not hardware or software behavior. It separates nearby declarations or operations so the file remains readable and reviewable.
\paragraph{Line 247}
\noindent\colorbox{gray!8}{\parbox{\dimexpr\linewidth-2\fboxsep\relax}{\raggedright\ttfamily\scriptsize     head\_weight = state["head.weight"].float()}}
\par\smallskip
This Python line participates in the export, validation, or report-generation pipeline. It is executed on the host PC to prepare data for RTL or to check RTL outputs.
\paragraph{Line 248}
\noindent\colorbox{gray!8}{\parbox{\dimexpr\linewidth-2\fboxsep\relax}{\raggedright\ttfamily\scriptsize     head\_bias = state["head.bias"].float()}}
\par\smallskip
This line handles the classifier parameters. Those arrays mirror the RTL memories, so mismatches reveal an export, quantization, or hardware-update bug.
\paragraph{Line 249}
\noindent\colorbox{gray!8}{\parbox{\dimexpr\linewidth-2\fboxsep\relax}{\raggedright\ttfamily\scriptsize     total\_weights += head\_weight.numel()}}
\par\smallskip
This line handles the classifier parameters. Those arrays mirror the RTL memories, so mismatches reveal an export, quantization, or hardware-update bug.
\paragraph{Line 250}
\noindent\colorbox{gray!8}{\parbox{\dimexpr\linewidth-2\fboxsep\relax}{\raggedright\ttfamily\scriptsize     total\_biases += head\_bias.numel()}}
\par\smallskip
This line handles the classifier parameters. Those arrays mirror the RTL memories, so mismatches reveal an export, quantization, or hardware-update bug.
\paragraph{Line 251}
\noindent\colorbox{gray!8}{\parbox{\dimexpr\linewidth-2\fboxsep\relax}{\raggedright\ttfamily\scriptsize     layers.append(}}
\par\smallskip
This Python line participates in the export, validation, or report-generation pipeline. It is executed on the host PC to prepare data for RTL or to check RTL outputs.
\paragraph{Line 252}
\noindent\colorbox{gray!8}{\parbox{\dimexpr\linewidth-2\fboxsep\relax}{\raggedright\ttfamily\scriptsize         \{}}
\par\smallskip
This Python line participates in the export, validation, or report-generation pipeline. It is executed on the host PC to prepare data for RTL or to check RTL outputs.
\paragraph{Line 253}
\noindent\colorbox{gray!8}{\parbox{\dimexpr\linewidth-2\fboxsep\relax}{\raggedright\ttfamily\scriptsize             "name": "head",}}
\par\smallskip
This Python line participates in the export, validation, or report-generation pipeline. It is executed on the host PC to prepare data for RTL or to check RTL outputs.
\paragraph{Line 254}
\noindent\colorbox{gray!8}{\parbox{\dimexpr\linewidth-2\fboxsep\relax}{\raggedright\ttfamily\scriptsize             "kind": "linear\_classifier",}}
\par\smallskip
This Python line participates in the export, validation, or report-generation pipeline. It is executed on the host PC to prepare data for RTL or to check RTL outputs.
\paragraph{Line 255}
\noindent\colorbox{gray!8}{\parbox{\dimexpr\linewidth-2\fboxsep\relax}{\raggedright\ttfamily\scriptsize             "weight": tensor\_stats(head\_weight),}}
\par\smallskip
This Python line participates in the export, validation, or report-generation pipeline. It is executed on the host PC to prepare data for RTL or to check RTL outputs.
\paragraph{Line 256}
\noindent\colorbox{gray!8}{\parbox{\dimexpr\linewidth-2\fboxsep\relax}{\raggedright\ttfamily\scriptsize             "bias": tensor\_stats(head\_bias),}}
\par\smallskip
This line handles the classifier parameters. Those arrays mirror the RTL memories, so mismatches reveal an export, quantization, or hardware-update bug.
\paragraph{Line 257}
\noindent\colorbox{gray!8}{\parbox{\dimexpr\linewidth-2\fboxsep\relax}{\raggedright\ttfamily\scriptsize             "weight\_values": int(head\_weight.numel()),}}
\par\smallskip
This Python line participates in the export, validation, or report-generation pipeline. It is executed on the host PC to prepare data for RTL or to check RTL outputs.
\paragraph{Line 258}
\noindent\colorbox{gray!8}{\parbox{\dimexpr\linewidth-2\fboxsep\relax}{\raggedright\ttfamily\scriptsize             "bias\_values": int(head\_bias.numel()),}}
\par\smallskip
This line handles the classifier parameters. Those arrays mirror the RTL memories, so mismatches reveal an export, quantization, or hardware-update bug.
\paragraph{Line 259}
\noindent\colorbox{gray!8}{\parbox{\dimexpr\linewidth-2\fboxsep\relax}{\raggedright\ttfamily\scriptsize         \}}}
\par\smallskip
This Python line participates in the export, validation, or report-generation pipeline. It is executed on the host PC to prepare data for RTL or to check RTL outputs.
\paragraph{Line 260}
\noindent\colorbox{gray!8}{\parbox{\dimexpr\linewidth-2\fboxsep\relax}{\raggedright\ttfamily\scriptsize     )}}
\par\smallskip
This Python line participates in the export, validation, or report-generation pipeline. It is executed on the host PC to prepare data for RTL or to check RTL outputs.
\paragraph{Line 261}
\noindent\colorbox{gray!8}{\parbox{\dimexpr\linewidth-2\fboxsep\relax}{\raggedright\ttfamily\scriptsize \emph{blank line}}}
\par\smallskip
This blank line is not hardware or software behavior. It separates nearby declarations or operations so the file remains readable and reviewable.
\paragraph{Line 262}
\noindent\colorbox{gray!8}{\parbox{\dimexpr\linewidth-2\fboxsep\relax}{\raggedright\ttfamily\scriptsize     return \{}}
\par\smallskip
This Python line participates in the export, validation, or report-generation pipeline. It is executed on the host PC to prepare data for RTL or to check RTL outputs.
\paragraph{Line 263}
\noindent\colorbox{gray!8}{\parbox{\dimexpr\linewidth-2\fboxsep\relax}{\raggedright\ttfamily\scriptsize         "checkpoint": str(args.checkpoint),}}
\par\smallskip
This Python line participates in the export, validation, or report-generation pipeline. It is executed on the host PC to prepare data for RTL or to check RTL outputs.
\paragraph{Line 264}
\noindent\colorbox{gray!8}{\parbox{\dimexpr\linewidth-2\fboxsep\relax}{\raggedright\ttfamily\scriptsize         "checkpoint\_metric\_top1": None,}}
\par\smallskip
This Python line participates in the export, validation, or report-generation pipeline. It is executed on the host PC to prepare data for RTL or to check RTL outputs.
\paragraph{Line 265}
\noindent\colorbox{gray!8}{\parbox{\dimexpr\linewidth-2\fboxsep\relax}{\raggedright\ttfamily\scriptsize         "checkpoint\_metric\_note": "Current measured accuracy is written after export as exported\_float\_top1\_pct and exported\_fixed\_top1\_pct.",}}
\par\smallskip
This Python line participates in the export, validation, or report-generation pipeline. It is executed on the host PC to prepare data for RTL or to check RTL outputs.
\paragraph{Line 266}
\noindent\colorbox{gray!8}{\parbox{\dimexpr\linewidth-2\fboxsep\relax}{\raggedright\ttfamily\scriptsize         "source\_model": "external\_research/Spikingformer/cifar10/model.py",}}
\par\smallskip
This Python line participates in the export, validation, or report-generation pipeline. It is executed on the host PC to prepare data for RTL or to check RTL outputs.
\paragraph{Line 267}
\noindent\colorbox{gray!8}{\parbox{\dimexpr\linewidth-2\fboxsep\relax}{\raggedright\ttfamily\scriptsize         "topology": \{}}
\par\smallskip
This Python line participates in the export, validation, or report-generation pipeline. It is executed on the host PC to prepare data for RTL or to check RTL outputs.
\paragraph{Line 268}
\noindent\colorbox{gray!8}{\parbox{\dimexpr\linewidth-2\fboxsep\relax}{\raggedright\ttfamily\scriptsize             "time\_steps": 4,}}
\par\smallskip
This Python line participates in the export, validation, or report-generation pipeline. It is executed on the host PC to prepare data for RTL or to check RTL outputs.
\paragraph{Line 269}
\noindent\colorbox{gray!8}{\parbox{\dimexpr\linewidth-2\fboxsep\relax}{\raggedright\ttfamily\scriptsize             "image": [3, 32, 32],}}
\par\smallskip
This Python line participates in the export, validation, or report-generation pipeline. It is executed on the host PC to prepare data for RTL or to check RTL outputs.
\paragraph{Line 270}
\noindent\colorbox{gray!8}{\parbox{\dimexpr\linewidth-2\fboxsep\relax}{\raggedright\ttfamily\scriptsize             "tokenizer": "5 convolution stages with LIF gates after stages 0-3 and 2 max-pool stages",}}
\par\smallskip
This Python line participates in the export, validation, or report-generation pipeline. It is executed on the host PC to prepare data for RTL or to check RTL outputs.
\paragraph{Line 271}
\noindent\colorbox{gray!8}{\parbox{\dimexpr\linewidth-2\fboxsep\relax}{\raggedright\ttfamily\scriptsize             "transformer\_blocks": 4,}}
\par\smallskip
This Python line participates in the export, validation, or report-generation pipeline. It is executed on the host PC to prepare data for RTL or to check RTL outputs.
\paragraph{Line 272}
\noindent\colorbox{gray!8}{\parbox{\dimexpr\linewidth-2\fboxsep\relax}{\raggedright\ttfamily\scriptsize             "embed\_dim": FEATURE\_DIM,}}
\par\smallskip
This Python line participates in the export, validation, or report-generation pipeline. It is executed on the host PC to prepare data for RTL or to check RTL outputs.
\paragraph{Line 273}
\noindent\colorbox{gray!8}{\parbox{\dimexpr\linewidth-2\fboxsep\relax}{\raggedright\ttfamily\scriptsize             "attention\_heads": 12,}}
\par\smallskip
This Python line participates in the export, validation, or report-generation pipeline. It is executed on the host PC to prepare data for RTL or to check RTL outputs.
\paragraph{Line 274}
\noindent\colorbox{gray!8}{\parbox{\dimexpr\linewidth-2\fboxsep\relax}{\raggedright\ttfamily\scriptsize             "mlp\_ratio": 4,}}
\par\smallskip
This Python line participates in the export, validation, or report-generation pipeline. It is executed on the host PC to prepare data for RTL or to check RTL outputs.
\paragraph{Line 275}
\noindent\colorbox{gray!8}{\parbox{\dimexpr\linewidth-2\fboxsep\relax}{\raggedright\ttfamily\scriptsize             "classifier": [FEATURE\_DIM, NUM\_CLASSES],}}
\par\smallskip
This Python line participates in the export, validation, or report-generation pipeline. It is executed on the host PC to prepare data for RTL or to check RTL outputs.
\paragraph{Line 276}
\noindent\colorbox{gray!8}{\parbox{\dimexpr\linewidth-2\fboxsep\relax}{\raggedright\ttfamily\scriptsize         \},}}
\par\smallskip
This Python line participates in the export, validation, or report-generation pipeline. It is executed on the host PC to prepare data for RTL or to check RTL outputs.
\paragraph{Line 277}
\noindent\colorbox{gray!8}{\parbox{\dimexpr\linewidth-2\fboxsep\relax}{\raggedright\ttfamily\scriptsize         "mapped\_layers": layers,}}
\par\smallskip
This Python line participates in the export, validation, or report-generation pipeline. It is executed on the host PC to prepare data for RTL or to check RTL outputs.
\paragraph{Line 278}
\noindent\colorbox{gray!8}{\parbox{\dimexpr\linewidth-2\fboxsep\relax}{\raggedright\ttfamily\scriptsize         "parameter\_counts": \{}}
\par\smallskip
This Python line participates in the export, validation, or report-generation pipeline. It is executed on the host PC to prepare data for RTL or to check RTL outputs.
\paragraph{Line 279}
\noindent\colorbox{gray!8}{\parbox{\dimexpr\linewidth-2\fboxsep\relax}{\raggedright\ttfamily\scriptsize             "weights": int(total\_weights),}}
\par\smallskip
This line handles the classifier parameters. Those arrays mirror the RTL memories, so mismatches reveal an export, quantization, or hardware-update bug.
\paragraph{Line 280}
\noindent\colorbox{gray!8}{\parbox{\dimexpr\linewidth-2\fboxsep\relax}{\raggedright\ttfamily\scriptsize             "biases": int(total\_biases),}}
\par\smallskip
This line handles the classifier parameters. Those arrays mirror the RTL memories, so mismatches reveal an export, quantization, or hardware-update bug.
\paragraph{Line 281}
\noindent\colorbox{gray!8}{\parbox{\dimexpr\linewidth-2\fboxsep\relax}{\raggedright\ttfamily\scriptsize             "total": int(total\_weights + total\_biases),}}
\par\smallskip
This line handles the classifier parameters. Those arrays mirror the RTL memories, so mismatches reveal an export, quantization, or hardware-update bug.
\paragraph{Line 282}
\noindent\colorbox{gray!8}{\parbox{\dimexpr\linewidth-2\fboxsep\relax}{\raggedright\ttfamily\scriptsize             "int8\_weight\_bytes": int(total\_weights),}}
\par\smallskip
This line handles the classifier parameters. Those arrays mirror the RTL memories, so mismatches reveal an export, quantization, or hardware-update bug.
\paragraph{Line 283}
\noindent\colorbox{gray!8}{\parbox{\dimexpr\linewidth-2\fboxsep\relax}{\raggedright\ttfamily\scriptsize             "int16\_weight\_bytes": int(total\_weights * 2),}}
\par\smallskip
This line handles the classifier parameters. Those arrays mirror the RTL memories, so mismatches reveal an export, quantization, or hardware-update bug.
\paragraph{Line 284}
\noindent\colorbox{gray!8}{\parbox{\dimexpr\linewidth-2\fboxsep\relax}{\raggedright\ttfamily\scriptsize             "float32\_checkpoint\_bytes": int((total\_weights + total\_biases) * 4),}}
\par\smallskip
This line handles the classifier parameters. Those arrays mirror the RTL memories, so mismatches reveal an export, quantization, or hardware-update bug.
\paragraph{Line 285}
\noindent\colorbox{gray!8}{\parbox{\dimexpr\linewidth-2\fboxsep\relax}{\raggedright\ttfamily\scriptsize         \},}}
\par\smallskip
This Python line participates in the export, validation, or report-generation pipeline. It is executed on the host PC to prepare data for RTL or to check RTL outputs.
\paragraph{Line 286}
\noindent\colorbox{gray!8}{\parbox{\dimexpr\linewidth-2\fboxsep\relax}{\raggedright\ttfamily\scriptsize         "implementation\_note": (}}
\par\smallskip
This Python line participates in the export, validation, or report-generation pipeline. It is executed on the host PC to prepare data for RTL or to check RTL outputs.
\paragraph{Line 287}
\noindent\colorbox{gray!8}{\parbox{\dimexpr\linewidth-2\fboxsep\relax}{\raggedright\ttfamily\scriptsize             "The shipped head RTL consumes fixed-point Spikingformer features and the trained classifier weights. "}}
\par\smallskip
This line handles the classifier parameters. Those arrays mirror the RTL memories, so mismatches reveal an export, quantization, or hardware-update bug.
\paragraph{Line 288}
\noindent\colorbox{gray!8}{\parbox{\dimexpr\linewidth-2\fboxsep\relax}{\raggedright\ttfamily\scriptsize             "The manifest above enumerates every trained backbone tensor that must be implemented by a DDR-backed "}}
\par\smallskip
This Python line participates in the export, validation, or report-generation pipeline. It is executed on the host PC to prepare data for RTL or to check RTL outputs.
\paragraph{Line 289}
\noindent\colorbox{gray!8}{\parbox{\dimexpr\linewidth-2\fboxsep\relax}{\raggedright\ttfamily\scriptsize             "streaming convolution/attention engine for a full bitstream-resident Spikingformer."}}
\par\smallskip
This Python line participates in the export, validation, or report-generation pipeline. It is executed on the host PC to prepare data for RTL or to check RTL outputs.
\paragraph{Line 290}
\noindent\colorbox{gray!8}{\parbox{\dimexpr\linewidth-2\fboxsep\relax}{\raggedright\ttfamily\scriptsize         ),}}
\par\smallskip
This Python line participates in the export, validation, or report-generation pipeline. It is executed on the host PC to prepare data for RTL or to check RTL outputs.
\paragraph{Line 291}
\noindent\colorbox{gray!8}{\parbox{\dimexpr\linewidth-2\fboxsep\relax}{\raggedright\ttfamily\scriptsize     \}}}
\par\smallskip
This Python line participates in the export, validation, or report-generation pipeline. It is executed on the host PC to prepare data for RTL or to check RTL outputs.
\paragraph{Line 292}
\noindent\colorbox{gray!8}{\parbox{\dimexpr\linewidth-2\fboxsep\relax}{\raggedright\ttfamily\scriptsize \emph{blank line}}}
\par\smallskip
This blank line is not hardware or software behavior. It separates nearby declarations or operations so the file remains readable and reviewable.
\paragraph{Line 293}
\noindent\colorbox{gray!8}{\parbox{\dimexpr\linewidth-2\fboxsep\relax}{\raggedright\ttfamily\scriptsize \emph{blank line}}}
\par\smallskip
This blank line is not hardware or software behavior. It separates nearby declarations or operations so the file remains readable and reviewable.
\paragraph{Line 294}
\noindent\colorbox{gray!8}{\parbox{\dimexpr\linewidth-2\fboxsep\relax}{\raggedright\ttfamily\scriptsize def write\_head\_memories(}}
\par\smallskip
This defines a Python function. The function packages one repeatable step such as loading a memory file, simulating the golden learning rule, validating CSVs, or generating plots.
\paragraph{Line 295}
\noindent\colorbox{gray!8}{\parbox{\dimexpr\linewidth-2\fboxsep\relax}{\raggedright\ttfamily\scriptsize     args: argparse.Namespace,}}
\par\smallskip
This contributes to the command-line interface. It lets the rerun flow choose paths, sample counts, or output locations without editing code.
\paragraph{Line 296}
\noindent\colorbox{gray!8}{\parbox{\dimexpr\linewidth-2\fboxsep\relax}{\raggedright\ttfamily\scriptsize     features: torch.Tensor,}}
\par\smallskip
This Python line participates in the export, validation, or report-generation pipeline. It is executed on the host PC to prepare data for RTL or to check RTL outputs.
\paragraph{Line 297}
\noindent\colorbox{gray!8}{\parbox{\dimexpr\linewidth-2\fboxsep\relax}{\raggedright\ttfamily\scriptsize     logits: torch.Tensor,}}
\par\smallskip
This Python line participates in the export, validation, or report-generation pipeline. It is executed on the host PC to prepare data for RTL or to check RTL outputs.
\paragraph{Line 298}
\noindent\colorbox{gray!8}{\parbox{\dimexpr\linewidth-2\fboxsep\relax}{\raggedright\ttfamily\scriptsize     labels: torch.Tensor,}}
\par\smallskip
This Python line participates in the export, validation, or report-generation pipeline. It is executed on the host PC to prepare data for RTL or to check RTL outputs.
\paragraph{Line 299}
\noindent\colorbox{gray!8}{\parbox{\dimexpr\linewidth-2\fboxsep\relax}{\raggedright\ttfamily\scriptsize ) -> dict[str, object]:}}
\par\smallskip
This Python line participates in the export, validation, or report-generation pipeline. It is executed on the host PC to prepare data for RTL or to check RTL outputs.
\paragraph{Line 300}
\noindent\colorbox{gray!8}{\parbox{\dimexpr\linewidth-2\fboxsep\relax}{\raggedright\ttfamily\scriptsize     state = load\_checkpoint(args.checkpoint)}}
\par\smallskip
This Python line participates in the export, validation, or report-generation pipeline. It is executed on the host PC to prepare data for RTL or to check RTL outputs.
\paragraph{Line 301}
\noindent\colorbox{gray!8}{\parbox{\dimexpr\linewidth-2\fboxsep\relax}{\raggedright\ttfamily\scriptsize     head\_weight = state["head.weight"].float()}}
\par\smallskip
This Python line participates in the export, validation, or report-generation pipeline. It is executed on the host PC to prepare data for RTL or to check RTL outputs.
\paragraph{Line 302}
\noindent\colorbox{gray!8}{\parbox{\dimexpr\linewidth-2\fboxsep\relax}{\raggedright\ttfamily\scriptsize     head\_bias = state["head.bias"].float()}}
\par\smallskip
This line handles the classifier parameters. Those arrays mirror the RTL memories, so mismatches reveal an export, quantization, or hardware-update bug.
\paragraph{Line 303}
\noindent\colorbox{gray!8}{\parbox{\dimexpr\linewidth-2\fboxsep\relax}{\raggedright\ttfamily\scriptsize \emph{blank line}}}
\par\smallskip
This blank line is not hardware or software behavior. It separates nearby declarations or operations so the file remains readable and reviewable.
\paragraph{Line 304}
\noindent\colorbox{gray!8}{\parbox{\dimexpr\linewidth-2\fboxsep\relax}{\raggedright\ttfamily\scriptsize     feature\_scale = choose\_scale(features, args.feature\_bits)}}
\par\smallskip
This Python line participates in the export, validation, or report-generation pipeline. It is executed on the host PC to prepare data for RTL or to check RTL outputs.
\paragraph{Line 305}
\noindent\colorbox{gray!8}{\parbox{\dimexpr\linewidth-2\fboxsep\relax}{\raggedright\ttfamily\scriptsize     weight\_scale = choose\_scale(head\_weight, args.weight\_bits)}}
\par\smallskip
This Python line participates in the export, validation, or report-generation pipeline. It is executed on the host PC to prepare data for RTL or to check RTL outputs.
\paragraph{Line 306}
\noindent\colorbox{gray!8}{\parbox{\dimexpr\linewidth-2\fboxsep\relax}{\raggedright\ttfamily\scriptsize     product\_scale = feature\_scale * weight\_scale}}
\par\smallskip
This Python line participates in the export, validation, or report-generation pipeline. It is executed on the host PC to prepare data for RTL or to check RTL outputs.
\paragraph{Line 307}
\noindent\colorbox{gray!8}{\parbox{\dimexpr\linewidth-2\fboxsep\relax}{\raggedright\ttfamily\scriptsize     features\_q = quantize\_signed(features, feature\_scale, args.feature\_bits)}}
\par\smallskip
This Python line participates in the export, validation, or report-generation pipeline. It is executed on the host PC to prepare data for RTL or to check RTL outputs.
\paragraph{Line 308}
\noindent\colorbox{gray!8}{\parbox{\dimexpr\linewidth-2\fboxsep\relax}{\raggedright\ttfamily\scriptsize     weights\_q = quantize\_signed(head\_weight, weight\_scale, args.weight\_bits)}}
\par\smallskip
This line handles the classifier parameters. Those arrays mirror the RTL memories, so mismatches reveal an export, quantization, or hardware-update bug.
\paragraph{Line 309}
\noindent\colorbox{gray!8}{\parbox{\dimexpr\linewidth-2\fboxsep\relax}{\raggedright\ttfamily\scriptsize     bias\_q = torch.round(head\_bias / product\_scale).to(torch.int64)}}
\par\smallskip
This line handles the classifier parameters. Those arrays mirror the RTL memories, so mismatches reveal an export, quantization, or hardware-update bug.
\paragraph{Line 310}
\noindent\colorbox{gray!8}{\parbox{\dimexpr\linewidth-2\fboxsep\relax}{\raggedright\ttfamily\scriptsize     features\_eff\_q = features\_q.clone()}}
\par\smallskip
This Python line participates in the export, validation, or report-generation pipeline. It is executed on the host PC to prepare data for RTL or to check RTL outputs.
\paragraph{Line 311}
\noindent\colorbox{gray!8}{\parbox{\dimexpr\linewidth-2\fboxsep\relax}{\raggedright\ttfamily\scriptsize     if args.infer\_skip\_abs\_thresh > 0:}}
\par\smallskip
This Python line participates in the export, validation, or report-generation pipeline. It is executed on the host PC to prepare data for RTL or to check RTL outputs.
\paragraph{Line 312}
\noindent\colorbox{gray!8}{\parbox{\dimexpr\linewidth-2\fboxsep\relax}{\raggedright\ttfamily\scriptsize         features\_eff\_q[features\_eff\_q.abs() <= args.infer\_skip\_abs\_thresh] = 0}}
\par\smallskip
This Python line participates in the export, validation, or report-generation pipeline. It is executed on the host PC to prepare data for RTL or to check RTL outputs.
\paragraph{Line 313}
\noindent\colorbox{gray!8}{\parbox{\dimexpr\linewidth-2\fboxsep\relax}{\raggedright\ttfamily\scriptsize     logits\_q = features\_eff\_q @ weights\_q.t() + bias\_q}}
\par\smallskip
This line handles the classifier parameters. Those arrays mirror the RTL memories, so mismatches reveal an export, quantization, or hardware-update bug.
\paragraph{Line 314}
\noindent\colorbox{gray!8}{\parbox{\dimexpr\linewidth-2\fboxsep\relax}{\raggedright\ttfamily\scriptsize     pred\_float = logits.argmax(dim=1)}}
\par\smallskip
This line handles predicted labels. It checks that argmax behavior in Python and RTL chooses the same class.
\paragraph{Line 315}
\noindent\colorbox{gray!8}{\parbox{\dimexpr\linewidth-2\fboxsep\relax}{\raggedright\ttfamily\scriptsize     pred\_q = logits\_q.argmax(dim=1)}}
\par\smallskip
This line handles predicted labels. It checks that argmax behavior in Python and RTL chooses the same class.
\paragraph{Line 316}
\noindent\colorbox{gray!8}{\parbox{\dimexpr\linewidth-2\fboxsep\relax}{\raggedright\ttfamily\scriptsize     correct\_float = pred\_float.eq(labels)}}
\par\smallskip
This line handles predicted labels. It checks that argmax behavior in Python and RTL chooses the same class.
\paragraph{Line 317}
\noindent\colorbox{gray!8}{\parbox{\dimexpr\linewidth-2\fboxsep\relax}{\raggedright\ttfamily\scriptsize     correct\_q = pred\_q.eq(labels)}}
\par\smallskip
This line handles predicted labels. It checks that argmax behavior in Python and RTL chooses the same class.
\paragraph{Line 318}
\noindent\colorbox{gray!8}{\parbox{\dimexpr\linewidth-2\fboxsep\relax}{\raggedright\ttfamily\scriptsize \emph{blank line}}}
\par\smallskip
This blank line is not hardware or software behavior. It separates nearby declarations or operations so the file remains readable and reviewable.
\paragraph{Line 319}
\noindent\colorbox{gray!8}{\parbox{\dimexpr\linewidth-2\fboxsep\relax}{\raggedright\ttfamily\scriptsize     rtl\_dir = args.rtl\_dir}}
\par\smallskip
This Python line participates in the export, validation, or report-generation pipeline. It is executed on the host PC to prepare data for RTL or to check RTL outputs.
\paragraph{Line 320}
\noindent\colorbox{gray!8}{\parbox{\dimexpr\linewidth-2\fboxsep\relax}{\raggedright\ttfamily\scriptsize     rtl\_dir.mkdir(parents=True, exist\_ok=True)}}
\par\smallskip
This Python line participates in the export, validation, or report-generation pipeline. It is executed on the host PC to prepare data for RTL or to check RTL outputs.
\paragraph{Line 321}
\noindent\colorbox{gray!8}{\parbox{\dimexpr\linewidth-2\fboxsep\relax}{\raggedright\ttfamily\scriptsize     out\_dir = args.out\_dir}}
\par\smallskip
This Python line participates in the export, validation, or report-generation pipeline. It is executed on the host PC to prepare data for RTL or to check RTL outputs.
\paragraph{Line 322}
\noindent\colorbox{gray!8}{\parbox{\dimexpr\linewidth-2\fboxsep\relax}{\raggedright\ttfamily\scriptsize     out\_dir.mkdir(parents=True, exist\_ok=True)}}
\par\smallskip
This Python line participates in the export, validation, or report-generation pipeline. It is executed on the host PC to prepare data for RTL or to check RTL outputs.
\paragraph{Line 323}
\noindent\colorbox{gray!8}{\parbox{\dimexpr\linewidth-2\fboxsep\relax}{\raggedright\ttfamily\scriptsize \emph{blank line}}}
\par\smallskip
This blank line is not hardware or software behavior. It separates nearby declarations or operations so the file remains readable and reviewable.
\paragraph{Line 324}
\noindent\colorbox{gray!8}{\parbox{\dimexpr\linewidth-2\fboxsep\relax}{\raggedright\ttfamily\scriptsize     with (rtl\_dir / "cifar10\_spikingformer\_head\_features.mem").open("w", encoding="ascii") as fh:}}
\par\smallskip
This reads an artifact produced by training, export, simulation, or Vivado. The validator/report depends on files generated by the actual run rather than hidden constants.
\paragraph{Line 325}
\noindent\colorbox{gray!8}{\parbox{\dimexpr\linewidth-2\fboxsep\relax}{\raggedright\ttfamily\scriptsize         for row in features\_q.tolist():}}
\par\smallskip
This Python line participates in the export, validation, or report-generation pipeline. It is executed on the host PC to prepare data for RTL or to check RTL outputs.
\paragraph{Line 326}
\noindent\colorbox{gray!8}{\parbox{\dimexpr\linewidth-2\fboxsep\relax}{\raggedright\ttfamily\scriptsize             for value in row:}}
\par\smallskip
This Python line participates in the export, validation, or report-generation pipeline. It is executed on the host PC to prepare data for RTL or to check RTL outputs.
\paragraph{Line 327}
\noindent\colorbox{gray!8}{\parbox{\dimexpr\linewidth-2\fboxsep\relax}{\raggedright\ttfamily\scriptsize                 fh.write(signed\_hex(int(value), args.feature\_bits) + "\textbackslash{}n")}}
\par\smallskip
This Python line participates in the export, validation, or report-generation pipeline. It is executed on the host PC to prepare data for RTL or to check RTL outputs.
\paragraph{Line 328}
\noindent\colorbox{gray!8}{\parbox{\dimexpr\linewidth-2\fboxsep\relax}{\raggedright\ttfamily\scriptsize     with (rtl\_dir / "cifar10\_spikingformer\_head\_weights.mem").open("w", encoding="ascii") as fh:}}
\par\smallskip
This reads an artifact produced by training, export, simulation, or Vivado. The validator/report depends on files generated by the actual run rather than hidden constants.
\paragraph{Line 329}
\noindent\colorbox{gray!8}{\parbox{\dimexpr\linewidth-2\fboxsep\relax}{\raggedright\ttfamily\scriptsize         for row in weights\_q.tolist():}}
\par\smallskip
This line handles the classifier parameters. Those arrays mirror the RTL memories, so mismatches reveal an export, quantization, or hardware-update bug.
\paragraph{Line 330}
\noindent\colorbox{gray!8}{\parbox{\dimexpr\linewidth-2\fboxsep\relax}{\raggedright\ttfamily\scriptsize             for value in row:}}
\par\smallskip
This Python line participates in the export, validation, or report-generation pipeline. It is executed on the host PC to prepare data for RTL or to check RTL outputs.
\paragraph{Line 331}
\noindent\colorbox{gray!8}{\parbox{\dimexpr\linewidth-2\fboxsep\relax}{\raggedright\ttfamily\scriptsize                 fh.write(signed\_hex(int(value), args.weight\_bits) + "\textbackslash{}n")}}
\par\smallskip
This Python line participates in the export, validation, or report-generation pipeline. It is executed on the host PC to prepare data for RTL or to check RTL outputs.
\paragraph{Line 332}
\noindent\colorbox{gray!8}{\parbox{\dimexpr\linewidth-2\fboxsep\relax}{\raggedright\ttfamily\scriptsize     with (rtl\_dir / "cifar10\_spikingformer\_head\_bias.mem").open("w", encoding="ascii") as fh:}}
\par\smallskip
This reads an artifact produced by training, export, simulation, or Vivado. The validator/report depends on files generated by the actual run rather than hidden constants.
\paragraph{Line 333}
\noindent\colorbox{gray!8}{\parbox{\dimexpr\linewidth-2\fboxsep\relax}{\raggedright\ttfamily\scriptsize         for value in bias\_q.tolist():}}
\par\smallskip
This line handles the classifier parameters. Those arrays mirror the RTL memories, so mismatches reveal an export, quantization, or hardware-update bug.
\paragraph{Line 334}
\noindent\colorbox{gray!8}{\parbox{\dimexpr\linewidth-2\fboxsep\relax}{\raggedright\ttfamily\scriptsize             fh.write(signed\_hex(int(value), 64) + "\textbackslash{}n")}}
\par\smallskip
This Python line participates in the export, validation, or report-generation pipeline. It is executed on the host PC to prepare data for RTL or to check RTL outputs.
\paragraph{Line 335}
\noindent\colorbox{gray!8}{\parbox{\dimexpr\linewidth-2\fboxsep\relax}{\raggedright\ttfamily\scriptsize     with (rtl\_dir / "cifar10\_spikingformer\_head\_labels.mem").open("w", encoding="ascii") as fh:}}
\par\smallskip
This reads an artifact produced by training, export, simulation, or Vivado. The validator/report depends on files generated by the actual run rather than hidden constants.
\paragraph{Line 336}
\noindent\colorbox{gray!8}{\parbox{\dimexpr\linewidth-2\fboxsep\relax}{\raggedright\ttfamily\scriptsize         for value in labels.tolist():}}
\par\smallskip
This Python line participates in the export, validation, or report-generation pipeline. It is executed on the host PC to prepare data for RTL or to check RTL outputs.
\paragraph{Line 337}
\noindent\colorbox{gray!8}{\parbox{\dimexpr\linewidth-2\fboxsep\relax}{\raggedright\ttfamily\scriptsize             fh.write(f"\{int(value):01x\}\textbackslash{}n")}}
\par\smallskip
This Python line participates in the export, validation, or report-generation pipeline. It is executed on the host PC to prepare data for RTL or to check RTL outputs.
\paragraph{Line 338}
\noindent\colorbox{gray!8}{\parbox{\dimexpr\linewidth-2\fboxsep\relax}{\raggedright\ttfamily\scriptsize     params = [}}
\par\smallskip
This Python line participates in the export, validation, or report-generation pipeline. It is executed on the host PC to prepare data for RTL or to check RTL outputs.
\paragraph{Line 339}
\noindent\colorbox{gray!8}{\parbox{\dimexpr\linewidth-2\fboxsep\relax}{\raggedright\ttfamily\scriptsize         f"`define CIFAR10\_SF\_HEAD\_NUM\_SAMPLES \{features.shape[0]\}",}}
\par\smallskip
This Python line participates in the export, validation, or report-generation pipeline. It is executed on the host PC to prepare data for RTL or to check RTL outputs.
\paragraph{Line 340}
\noindent\colorbox{gray!8}{\parbox{\dimexpr\linewidth-2\fboxsep\relax}{\raggedright\ttfamily\scriptsize         f"`define CIFAR10\_SF\_HEAD\_FEATURE\_DIM \{FEATURE\_DIM\}",}}
\par\smallskip
This Python line participates in the export, validation, or report-generation pipeline. It is executed on the host PC to prepare data for RTL or to check RTL outputs.
\paragraph{Line 341}
\noindent\colorbox{gray!8}{\parbox{\dimexpr\linewidth-2\fboxsep\relax}{\raggedright\ttfamily\scriptsize         f"`define CIFAR10\_SF\_HEAD\_NUM\_CLASSES \{NUM\_CLASSES\}",}}
\par\smallskip
This Python line participates in the export, validation, or report-generation pipeline. It is executed on the host PC to prepare data for RTL or to check RTL outputs.
\paragraph{Line 342}
\noindent\colorbox{gray!8}{\parbox{\dimexpr\linewidth-2\fboxsep\relax}{\raggedright\ttfamily\scriptsize         f"`define CIFAR10\_SF\_HEAD\_FEATURE\_BITS \{args.feature\_bits\}",}}
\par\smallskip
This Python line participates in the export, validation, or report-generation pipeline. It is executed on the host PC to prepare data for RTL or to check RTL outputs.
\paragraph{Line 343}
\noindent\colorbox{gray!8}{\parbox{\dimexpr\linewidth-2\fboxsep\relax}{\raggedright\ttfamily\scriptsize         f"`define CIFAR10\_SF\_HEAD\_WEIGHT\_BITS \{args.weight\_bits\}",}}
\par\smallskip
This Python line participates in the export, validation, or report-generation pipeline. It is executed on the host PC to prepare data for RTL or to check RTL outputs.
\paragraph{Line 344}
\noindent\colorbox{gray!8}{\parbox{\dimexpr\linewidth-2\fboxsep\relax}{\raggedright\ttfamily\scriptsize         "`define CIFAR10\_SF\_HEAD\_ACC\_BITS 64",}}
\par\smallskip
This Python line participates in the export, validation, or report-generation pipeline. It is executed on the host PC to prepare data for RTL or to check RTL outputs.
\paragraph{Line 345}
\noindent\colorbox{gray!8}{\parbox{\dimexpr\linewidth-2\fboxsep\relax}{\raggedright\ttfamily\scriptsize         f"`define CIFAR10\_SF\_HEAD\_TRACE\_SAMPLES \{args.trace\_samples\}",}}
\par\smallskip
This Python line participates in the export, validation, or report-generation pipeline. It is executed on the host PC to prepare data for RTL or to check RTL outputs.
\paragraph{Line 346}
\noindent\colorbox{gray!8}{\parbox{\dimexpr\linewidth-2\fboxsep\relax}{\raggedright\ttfamily\scriptsize         f"`define CIFAR10\_SF\_HEAD\_INFER\_SKIP\_ABS\_THRESH \{args.infer\_skip\_abs\_thresh\}",}}
\par\smallskip
This Python line participates in the export, validation, or report-generation pipeline. It is executed on the host PC to prepare data for RTL or to check RTL outputs.
\paragraph{Line 347}
\noindent\colorbox{gray!8}{\parbox{\dimexpr\linewidth-2\fboxsep\relax}{\raggedright\ttfamily\scriptsize     ]}}
\par\smallskip
This Python line participates in the export, validation, or report-generation pipeline. It is executed on the host PC to prepare data for RTL or to check RTL outputs.
\paragraph{Line 348}
\noindent\colorbox{gray!8}{\parbox{\dimexpr\linewidth-2\fboxsep\relax}{\raggedright\ttfamily\scriptsize     (rtl\_dir / "cifar10\_spikingformer\_head\_params.svh").write\_text("\textbackslash{}n".join(params) + "\textbackslash{}n", encoding="ascii")}}
\par\smallskip
This writes machine-readable output for later reports. Keeping JSON/CSV artifacts makes the verification auditable.
\paragraph{Line 349}
\noindent\colorbox{gray!8}{\parbox{\dimexpr\linewidth-2\fboxsep\relax}{\raggedright\ttfamily\scriptsize \emph{blank line}}}
\par\smallskip
This blank line is not hardware or software behavior. It separates nearby declarations or operations so the file remains readable and reviewable.
\paragraph{Line 350}
\noindent\colorbox{gray!8}{\parbox{\dimexpr\linewidth-2\fboxsep\relax}{\raggedright\ttfamily\scriptsize     fixed\_ref = out\_dir / "spikingformer\_head\_fixed\_reference.csv"}}
\par\smallskip
This Python line participates in the export, validation, or report-generation pipeline. It is executed on the host PC to prepare data for RTL or to check RTL outputs.
\paragraph{Line 351}
\noindent\colorbox{gray!8}{\parbox{\dimexpr\linewidth-2\fboxsep\relax}{\raggedright\ttfamily\scriptsize     with fixed\_ref.open("w", newline="", encoding="utf-8") as fh:}}
\par\smallskip
This reads an artifact produced by training, export, simulation, or Vivado. The validator/report depends on files generated by the actual run rather than hidden constants.
\paragraph{Line 352}
\noindent\colorbox{gray!8}{\parbox{\dimexpr\linewidth-2\fboxsep\relax}{\raggedright\ttfamily\scriptsize         fieldnames = ["sample", "label", "float\_pred", "fixed\_pred", "float\_correct", "fixed\_correct"]}}
\par\smallskip
This line handles predicted labels. It checks that argmax behavior in Python and RTL chooses the same class.
\paragraph{Line 353}
\noindent\colorbox{gray!8}{\parbox{\dimexpr\linewidth-2\fboxsep\relax}{\raggedright\ttfamily\scriptsize         fieldnames.extend([f"logit\_q\{c\}" for c in range(NUM\_CLASSES)])}}
\par\smallskip
This Python line participates in the export, validation, or report-generation pipeline. It is executed on the host PC to prepare data for RTL or to check RTL outputs.
\paragraph{Line 354}
\noindent\colorbox{gray!8}{\parbox{\dimexpr\linewidth-2\fboxsep\relax}{\raggedright\ttfamily\scriptsize         writer = csv.DictWriter(fh, fieldnames=fieldnames)}}
\par\smallskip
This Python line participates in the export, validation, or report-generation pipeline. It is executed on the host PC to prepare data for RTL or to check RTL outputs.
\paragraph{Line 355}
\noindent\colorbox{gray!8}{\parbox{\dimexpr\linewidth-2\fboxsep\relax}{\raggedright\ttfamily\scriptsize         writer.writeheader()}}
\par\smallskip
This Python line participates in the export, validation, or report-generation pipeline. It is executed on the host PC to prepare data for RTL or to check RTL outputs.
\paragraph{Line 356}
\noindent\colorbox{gray!8}{\parbox{\dimexpr\linewidth-2\fboxsep\relax}{\raggedright\ttfamily\scriptsize         for idx in range(features.shape[0]):}}
\par\smallskip
This Python line participates in the export, validation, or report-generation pipeline. It is executed on the host PC to prepare data for RTL or to check RTL outputs.
\paragraph{Line 357}
\noindent\colorbox{gray!8}{\parbox{\dimexpr\linewidth-2\fboxsep\relax}{\raggedright\ttfamily\scriptsize             row = \{}}
\par\smallskip
This Python line participates in the export, validation, or report-generation pipeline. It is executed on the host PC to prepare data for RTL or to check RTL outputs.
\paragraph{Line 358}
\noindent\colorbox{gray!8}{\parbox{\dimexpr\linewidth-2\fboxsep\relax}{\raggedright\ttfamily\scriptsize                 "sample": idx,}}
\par\smallskip
This Python line participates in the export, validation, or report-generation pipeline. It is executed on the host PC to prepare data for RTL or to check RTL outputs.
\paragraph{Line 359}
\noindent\colorbox{gray!8}{\parbox{\dimexpr\linewidth-2\fboxsep\relax}{\raggedright\ttfamily\scriptsize                 "label": int(labels[idx]),}}
\par\smallskip
This Python line participates in the export, validation, or report-generation pipeline. It is executed on the host PC to prepare data for RTL or to check RTL outputs.
\paragraph{Line 360}
\noindent\colorbox{gray!8}{\parbox{\dimexpr\linewidth-2\fboxsep\relax}{\raggedright\ttfamily\scriptsize                 "float\_pred": int(pred\_float[idx]),}}
\par\smallskip
This line handles predicted labels. It checks that argmax behavior in Python and RTL chooses the same class.
\paragraph{Line 361}
\noindent\colorbox{gray!8}{\parbox{\dimexpr\linewidth-2\fboxsep\relax}{\raggedright\ttfamily\scriptsize                 "fixed\_pred": int(pred\_q[idx]),}}
\par\smallskip
This line handles predicted labels. It checks that argmax behavior in Python and RTL chooses the same class.
\paragraph{Line 362}
\noindent\colorbox{gray!8}{\parbox{\dimexpr\linewidth-2\fboxsep\relax}{\raggedright\ttfamily\scriptsize                 "float\_correct": int(correct\_float[idx]),}}
\par\smallskip
This Python line participates in the export, validation, or report-generation pipeline. It is executed on the host PC to prepare data for RTL or to check RTL outputs.
\paragraph{Line 363}
\noindent\colorbox{gray!8}{\parbox{\dimexpr\linewidth-2\fboxsep\relax}{\raggedright\ttfamily\scriptsize                 "fixed\_correct": int(correct\_q[idx]),}}
\par\smallskip
This Python line participates in the export, validation, or report-generation pipeline. It is executed on the host PC to prepare data for RTL or to check RTL outputs.
\paragraph{Line 364}
\noindent\colorbox{gray!8}{\parbox{\dimexpr\linewidth-2\fboxsep\relax}{\raggedright\ttfamily\scriptsize             \}}}
\par\smallskip
This Python line participates in the export, validation, or report-generation pipeline. It is executed on the host PC to prepare data for RTL or to check RTL outputs.
\paragraph{Line 365}
\noindent\colorbox{gray!8}{\parbox{\dimexpr\linewidth-2\fboxsep\relax}{\raggedright\ttfamily\scriptsize             for cls in range(NUM\_CLASSES):}}
\par\smallskip
This Python line participates in the export, validation, or report-generation pipeline. It is executed on the host PC to prepare data for RTL or to check RTL outputs.
\paragraph{Line 366}
\noindent\colorbox{gray!8}{\parbox{\dimexpr\linewidth-2\fboxsep\relax}{\raggedright\ttfamily\scriptsize                 row[f"logit\_q\{cls\}"] = int(logits\_q[idx, cls])}}
\par\smallskip
This Python line participates in the export, validation, or report-generation pipeline. It is executed on the host PC to prepare data for RTL or to check RTL outputs.
\paragraph{Line 367}
\noindent\colorbox{gray!8}{\parbox{\dimexpr\linewidth-2\fboxsep\relax}{\raggedright\ttfamily\scriptsize             writer.writerow(row)}}
\par\smallskip
This Python line participates in the export, validation, or report-generation pipeline. It is executed on the host PC to prepare data for RTL or to check RTL outputs.
\paragraph{Line 368}
\noindent\colorbox{gray!8}{\parbox{\dimexpr\linewidth-2\fboxsep\relax}{\raggedright\ttfamily\scriptsize \emph{blank line}}}
\par\smallskip
This blank line is not hardware or software behavior. It separates nearby declarations or operations so the file remains readable and reviewable.
\paragraph{Line 369}
\noindent\colorbox{gray!8}{\parbox{\dimexpr\linewidth-2\fboxsep\relax}{\raggedright\ttfamily\scriptsize     confusion = np.zeros((NUM\_CLASSES, NUM\_CLASSES), dtype=np.int64)}}
\par\smallskip
This performs vectorized numerical work. The fixed-point golden models use NumPy so score, margin, and update calculations are explicit and repeatable.
\paragraph{Line 370}
\noindent\colorbox{gray!8}{\parbox{\dimexpr\linewidth-2\fboxsep\relax}{\raggedright\ttfamily\scriptsize     for label, pred in zip(labels.numpy(), pred\_q.numpy()):}}
\par\smallskip
This performs vectorized numerical work. The fixed-point golden models use NumPy so score, margin, and update calculations are explicit and repeatable.
\paragraph{Line 371}
\noindent\colorbox{gray!8}{\parbox{\dimexpr\linewidth-2\fboxsep\relax}{\raggedright\ttfamily\scriptsize         confusion[int(label), int(pred)] += 1}}
\par\smallskip
This line handles predicted labels. It checks that argmax behavior in Python and RTL chooses the same class.
\paragraph{Line 372}
\noindent\colorbox{gray!8}{\parbox{\dimexpr\linewidth-2\fboxsep\relax}{\raggedright\ttfamily\scriptsize \emph{blank line}}}
\par\smallskip
This blank line is not hardware or software behavior. It separates nearby declarations or operations so the file remains readable and reviewable.
\paragraph{Line 373}
\noindent\colorbox{gray!8}{\parbox{\dimexpr\linewidth-2\fboxsep\relax}{\raggedright\ttfamily\scriptsize     metrics: dict[str, object] = \{}}
\par\smallskip
This Python line participates in the export, validation, or report-generation pipeline. It is executed on the host PC to prepare data for RTL or to check RTL outputs.
\paragraph{Line 374}
\noindent\colorbox{gray!8}{\parbox{\dimexpr\linewidth-2\fboxsep\relax}{\raggedright\ttfamily\scriptsize         "checkpoint": str(args.checkpoint),}}
\par\smallskip
This Python line participates in the export, validation, or report-generation pipeline. It is executed on the host PC to prepare data for RTL or to check RTL outputs.
\paragraph{Line 375}
\noindent\colorbox{gray!8}{\parbox{\dimexpr\linewidth-2\fboxsep\relax}{\raggedright\ttfamily\scriptsize         "samples": int(features.shape[0]),}}
\par\smallskip
This Python line participates in the export, validation, or report-generation pipeline. It is executed on the host PC to prepare data for RTL or to check RTL outputs.
\paragraph{Line 376}
\noindent\colorbox{gray!8}{\parbox{\dimexpr\linewidth-2\fboxsep\relax}{\raggedright\ttfamily\scriptsize         "feature\_dim": FEATURE\_DIM,}}
\par\smallskip
This Python line participates in the export, validation, or report-generation pipeline. It is executed on the host PC to prepare data for RTL or to check RTL outputs.
\paragraph{Line 377}
\noindent\colorbox{gray!8}{\parbox{\dimexpr\linewidth-2\fboxsep\relax}{\raggedright\ttfamily\scriptsize         "float\_accuracy": float(correct\_float.float().mean()),}}
\par\smallskip
This Python line participates in the export, validation, or report-generation pipeline. It is executed on the host PC to prepare data for RTL or to check RTL outputs.
\paragraph{Line 378}
\noindent\colorbox{gray!8}{\parbox{\dimexpr\linewidth-2\fboxsep\relax}{\raggedright\ttfamily\scriptsize         "fixed\_accuracy": float(correct\_q.float().mean()),}}
\par\smallskip
This Python line participates in the export, validation, or report-generation pipeline. It is executed on the host PC to prepare data for RTL or to check RTL outputs.
\paragraph{Line 379}
\noindent\colorbox{gray!8}{\parbox{\dimexpr\linewidth-2\fboxsep\relax}{\raggedright\ttfamily\scriptsize         "float\_fixed\_prediction\_agreement": float(pred\_float.eq(pred\_q).float().mean()),}}
\par\smallskip
This line handles predicted labels. It checks that argmax behavior in Python and RTL chooses the same class.
\paragraph{Line 380}
\noindent\colorbox{gray!8}{\parbox{\dimexpr\linewidth-2\fboxsep\relax}{\raggedright\ttfamily\scriptsize         "feature\_scale": feature\_scale,}}
\par\smallskip
This Python line participates in the export, validation, or report-generation pipeline. It is executed on the host PC to prepare data for RTL or to check RTL outputs.
\paragraph{Line 381}
\noindent\colorbox{gray!8}{\parbox{\dimexpr\linewidth-2\fboxsep\relax}{\raggedright\ttfamily\scriptsize         "weight\_scale": weight\_scale,}}
\par\smallskip
This Python line participates in the export, validation, or report-generation pipeline. It is executed on the host PC to prepare data for RTL or to check RTL outputs.
\paragraph{Line 382}
\noindent\colorbox{gray!8}{\parbox{\dimexpr\linewidth-2\fboxsep\relax}{\raggedright\ttfamily\scriptsize         "product\_scale": product\_scale,}}
\par\smallskip
This Python line participates in the export, validation, or report-generation pipeline. It is executed on the host PC to prepare data for RTL or to check RTL outputs.
\paragraph{Line 383}
\noindent\colorbox{gray!8}{\parbox{\dimexpr\linewidth-2\fboxsep\relax}{\raggedright\ttfamily\scriptsize         "feature\_bits": args.feature\_bits,}}
\par\smallskip
This Python line participates in the export, validation, or report-generation pipeline. It is executed on the host PC to prepare data for RTL or to check RTL outputs.
\paragraph{Line 384}
\noindent\colorbox{gray!8}{\parbox{\dimexpr\linewidth-2\fboxsep\relax}{\raggedright\ttfamily\scriptsize         "weight\_bits": args.weight\_bits,}}
\par\smallskip
This Python line participates in the export, validation, or report-generation pipeline. It is executed on the host PC to prepare data for RTL or to check RTL outputs.
\paragraph{Line 385}
\noindent\colorbox{gray!8}{\parbox{\dimexpr\linewidth-2\fboxsep\relax}{\raggedright\ttfamily\scriptsize         "bias\_bits": 64,}}
\par\smallskip
This line handles the classifier parameters. Those arrays mirror the RTL memories, so mismatches reveal an export, quantization, or hardware-update bug.
\paragraph{Line 386}
\noindent\colorbox{gray!8}{\parbox{\dimexpr\linewidth-2\fboxsep\relax}{\raggedright\ttfamily\scriptsize         "infer\_skip\_abs\_thresh": args.infer\_skip\_abs\_thresh,}}
\par\smallskip
This Python line participates in the export, validation, or report-generation pipeline. It is executed on the host PC to prepare data for RTL or to check RTL outputs.
\paragraph{Line 387}
\noindent\colorbox{gray!8}{\parbox{\dimexpr\linewidth-2\fboxsep\relax}{\raggedright\ttfamily\scriptsize         "feature\_mac\_skip\_fraction": float((features\_q.abs() <= args.infer\_skip\_abs\_thresh).float().mean()) if args.infer\_skip\_abs\_thresh > 0 else 0.0,}}
\par\smallskip
This Python line participates in the export, validation, or report-generation pipeline. It is executed on the host PC to prepare data for RTL or to check RTL outputs.
\paragraph{Line 388}
\noindent\colorbox{gray!8}{\parbox{\dimexpr\linewidth-2\fboxsep\relax}{\raggedright\ttfamily\scriptsize         "confusion": confusion.tolist(),}}
\par\smallskip
This Python line participates in the export, validation, or report-generation pipeline. It is executed on the host PC to prepare data for RTL or to check RTL outputs.
\paragraph{Line 389}
\noindent\colorbox{gray!8}{\parbox{\dimexpr\linewidth-2\fboxsep\relax}{\raggedright\ttfamily\scriptsize         "files": \{}}
\par\smallskip
This Python line participates in the export, validation, or report-generation pipeline. It is executed on the host PC to prepare data for RTL or to check RTL outputs.
\paragraph{Line 390}
\noindent\colorbox{gray!8}{\parbox{\dimexpr\linewidth-2\fboxsep\relax}{\raggedright\ttfamily\scriptsize             "features": str(rtl\_dir / "cifar10\_spikingformer\_head\_features.mem"),}}
\par\smallskip
This Python line participates in the export, validation, or report-generation pipeline. It is executed on the host PC to prepare data for RTL or to check RTL outputs.
\paragraph{Line 391}
\noindent\colorbox{gray!8}{\parbox{\dimexpr\linewidth-2\fboxsep\relax}{\raggedright\ttfamily\scriptsize             "weights": str(rtl\_dir / "cifar10\_spikingformer\_head\_weights.mem"),}}
\par\smallskip
This line handles the classifier parameters. Those arrays mirror the RTL memories, so mismatches reveal an export, quantization, or hardware-update bug.
\paragraph{Line 392}
\noindent\colorbox{gray!8}{\parbox{\dimexpr\linewidth-2\fboxsep\relax}{\raggedright\ttfamily\scriptsize             "bias": str(rtl\_dir / "cifar10\_spikingformer\_head\_bias.mem"),}}
\par\smallskip
This line handles the classifier parameters. Those arrays mirror the RTL memories, so mismatches reveal an export, quantization, or hardware-update bug.
\paragraph{Line 393}
\noindent\colorbox{gray!8}{\parbox{\dimexpr\linewidth-2\fboxsep\relax}{\raggedright\ttfamily\scriptsize             "labels": str(rtl\_dir / "cifar10\_spikingformer\_head\_labels.mem"),}}
\par\smallskip
This Python line participates in the export, validation, or report-generation pipeline. It is executed on the host PC to prepare data for RTL or to check RTL outputs.
\paragraph{Line 394}
\noindent\colorbox{gray!8}{\parbox{\dimexpr\linewidth-2\fboxsep\relax}{\raggedright\ttfamily\scriptsize             "params": str(rtl\_dir / "cifar10\_spikingformer\_head\_params.svh"),}}
\par\smallskip
This Python line participates in the export, validation, or report-generation pipeline. It is executed on the host PC to prepare data for RTL or to check RTL outputs.
\paragraph{Line 395}
\noindent\colorbox{gray!8}{\parbox{\dimexpr\linewidth-2\fboxsep\relax}{\raggedright\ttfamily\scriptsize             "fixed\_reference": str(fixed\_ref),}}
\par\smallskip
This Python line participates in the export, validation, or report-generation pipeline. It is executed on the host PC to prepare data for RTL or to check RTL outputs.
\paragraph{Line 396}
\noindent\colorbox{gray!8}{\parbox{\dimexpr\linewidth-2\fboxsep\relax}{\raggedright\ttfamily\scriptsize         \},}}
\par\smallskip
This Python line participates in the export, validation, or report-generation pipeline. It is executed on the host PC to prepare data for RTL or to check RTL outputs.
\paragraph{Line 397}
\noindent\colorbox{gray!8}{\parbox{\dimexpr\linewidth-2\fboxsep\relax}{\raggedright\ttfamily\scriptsize     \}}}
\par\smallskip
This Python line participates in the export, validation, or report-generation pipeline. It is executed on the host PC to prepare data for RTL or to check RTL outputs.
\paragraph{Line 398}
\noindent\colorbox{gray!8}{\parbox{\dimexpr\linewidth-2\fboxsep\relax}{\raggedright\ttfamily\scriptsize     (out\_dir / "spikingformer\_head\_mapping\_metrics.json").write\_text(json.dumps(metrics, indent=2), encoding="utf-8")}}
\par\smallskip
This writes machine-readable output for later reports. Keeping JSON/CSV artifacts makes the verification auditable.
\paragraph{Line 399}
\noindent\colorbox{gray!8}{\parbox{\dimexpr\linewidth-2\fboxsep\relax}{\raggedright\ttfamily\scriptsize     return metrics}}
\par\smallskip
This Python line participates in the export, validation, or report-generation pipeline. It is executed on the host PC to prepare data for RTL or to check RTL outputs.
\paragraph{Line 400}
\noindent\colorbox{gray!8}{\parbox{\dimexpr\linewidth-2\fboxsep\relax}{\raggedright\ttfamily\scriptsize \emph{blank line}}}
\par\smallskip
This blank line is not hardware or software behavior. It separates nearby declarations or operations so the file remains readable and reviewable.
\paragraph{Line 401}
\noindent\colorbox{gray!8}{\parbox{\dimexpr\linewidth-2\fboxsep\relax}{\raggedright\ttfamily\scriptsize \emph{blank line}}}
\par\smallskip
This blank line is not hardware or software behavior. It separates nearby declarations or operations so the file remains readable and reviewable.
\paragraph{Line 402}
\noindent\colorbox{gray!8}{\parbox{\dimexpr\linewidth-2\fboxsep\relax}{\raggedright\ttfamily\scriptsize def write\_report(args: argparse.Namespace, metrics: dict[str, object], mapping: dict[str, object]) -> None:}}
\par\smallskip
This defines a Python function. The function packages one repeatable step such as loading a memory file, simulating the golden learning rule, validating CSVs, or generating plots.
\paragraph{Line 403}
\noindent\colorbox{gray!8}{\parbox{\dimexpr\linewidth-2\fboxsep\relax}{\raggedright\ttfamily\scriptsize     out\_dir = args.out\_dir}}
\par\smallskip
This Python line participates in the export, validation, or report-generation pipeline. It is executed on the host PC to prepare data for RTL or to check RTL outputs.
\paragraph{Line 404}
\noindent\colorbox{gray!8}{\parbox{\dimexpr\linewidth-2\fboxsep\relax}{\raggedright\ttfamily\scriptsize     confusion = np.array(metrics["confusion"], dtype=np.float64)}}
\par\smallskip
This performs vectorized numerical work. The fixed-point golden models use NumPy so score, margin, and update calculations are explicit and repeatable.
\paragraph{Line 405}
\noindent\colorbox{gray!8}{\parbox{\dimexpr\linewidth-2\fboxsep\relax}{\raggedright\ttfamily\scriptsize     norm\_confusion = confusion / np.maximum(confusion.sum(axis=1, keepdims=True), 1.0)}}
\par\smallskip
This performs vectorized numerical work. The fixed-point golden models use NumPy so score, margin, and update calculations are explicit and repeatable.
\paragraph{Line 406}
\noindent\colorbox{gray!8}{\parbox{\dimexpr\linewidth-2\fboxsep\relax}{\raggedright\ttfamily\scriptsize \emph{blank line}}}
\par\smallskip
This blank line is not hardware or software behavior. It separates nearby declarations or operations so the file remains readable and reviewable.
\paragraph{Line 407}
\noindent\colorbox{gray!8}{\parbox{\dimexpr\linewidth-2\fboxsep\relax}{\raggedright\ttfamily\scriptsize     plt.figure(figsize=(7, 6), dpi=160)}}
\par\smallskip
This creates a debug visualization. The plot is derived from generated CSV or Vivado report data so it can be regenerated after RTL changes.
\paragraph{Line 408}
\noindent\colorbox{gray!8}{\parbox{\dimexpr\linewidth-2\fboxsep\relax}{\raggedright\ttfamily\scriptsize     plt.imshow(norm\_confusion, cmap="viridis", vmin=0, vmax=1)}}
\par\smallskip
This creates a debug visualization. The plot is derived from generated CSV or Vivado report data so it can be regenerated after RTL changes.
\paragraph{Line 409}
\noindent\colorbox{gray!8}{\parbox{\dimexpr\linewidth-2\fboxsep\relax}{\raggedright\ttfamily\scriptsize     plt.xticks(range(NUM\_CLASSES), CIFAR10\_CLASSES, rotation=45, ha="right")}}
\par\smallskip
This creates a debug visualization. The plot is derived from generated CSV or Vivado report data so it can be regenerated after RTL changes.
\paragraph{Line 410}
\noindent\colorbox{gray!8}{\parbox{\dimexpr\linewidth-2\fboxsep\relax}{\raggedright\ttfamily\scriptsize     plt.yticks(range(NUM\_CLASSES), CIFAR10\_CLASSES)}}
\par\smallskip
This creates a debug visualization. The plot is derived from generated CSV or Vivado report data so it can be regenerated after RTL changes.
\paragraph{Line 411}
\noindent\colorbox{gray!8}{\parbox{\dimexpr\linewidth-2\fboxsep\relax}{\raggedright\ttfamily\scriptsize     plt.colorbar(label="row-normalized count")}}
\par\smallskip
This creates a debug visualization. The plot is derived from generated CSV or Vivado report data so it can be regenerated after RTL changes.
\paragraph{Line 412}
\noindent\colorbox{gray!8}{\parbox{\dimexpr\linewidth-2\fboxsep\relax}{\raggedright\ttfamily\scriptsize     plt.tight\_layout()}}
\par\smallskip
This creates a debug visualization. The plot is derived from generated CSV or Vivado report data so it can be regenerated after RTL changes.
\paragraph{Line 413}
\noindent\colorbox{gray!8}{\parbox{\dimexpr\linewidth-2\fboxsep\relax}{\raggedright\ttfamily\scriptsize     plt.savefig(out\_dir / "spikingformer\_head\_fixed\_confusion.png")}}
\par\smallskip
This creates a debug visualization. The plot is derived from generated CSV or Vivado report data so it can be regenerated after RTL changes.
\paragraph{Line 414}
\noindent\colorbox{gray!8}{\parbox{\dimexpr\linewidth-2\fboxsep\relax}{\raggedright\ttfamily\scriptsize     plt.close()}}
\par\smallskip
This creates a debug visualization. The plot is derived from generated CSV or Vivado report data so it can be regenerated after RTL changes.
\paragraph{Line 415}
\noindent\colorbox{gray!8}{\parbox{\dimexpr\linewidth-2\fboxsep\relax}{\raggedright\ttfamily\scriptsize \emph{blank line}}}
\par\smallskip
This blank line is not hardware or software behavior. It separates nearby declarations or operations so the file remains readable and reviewable.
\paragraph{Line 416}
\noindent\colorbox{gray!8}{\parbox{\dimexpr\linewidth-2\fboxsep\relax}{\raggedright\ttfamily\scriptsize     layer\_names = [layer["name"] for layer in mapping["mapped\_layers"]]}}
\par\smallskip
This Python line participates in the export, validation, or report-generation pipeline. It is executed on the host PC to prepare data for RTL or to check RTL outputs.
\paragraph{Line 417}
\noindent\colorbox{gray!8}{\parbox{\dimexpr\linewidth-2\fboxsep\relax}{\raggedright\ttfamily\scriptsize     layer\_counts = [layer["weight\_values"] for layer in mapping["mapped\_layers"]]}}
\par\smallskip
This Python line participates in the export, validation, or report-generation pipeline. It is executed on the host PC to prepare data for RTL or to check RTL outputs.
\paragraph{Line 418}
\noindent\colorbox{gray!8}{\parbox{\dimexpr\linewidth-2\fboxsep\relax}{\raggedright\ttfamily\scriptsize     plt.figure(figsize=(12, 5), dpi=160)}}
\par\smallskip
This creates a debug visualization. The plot is derived from generated CSV or Vivado report data so it can be regenerated after RTL changes.
\paragraph{Line 419}
\noindent\colorbox{gray!8}{\parbox{\dimexpr\linewidth-2\fboxsep\relax}{\raggedright\ttfamily\scriptsize     plt.bar(range(len(layer\_counts)), layer\_counts)}}
\par\smallskip
This creates a debug visualization. The plot is derived from generated CSV or Vivado report data so it can be regenerated after RTL changes.
\paragraph{Line 420}
\noindent\colorbox{gray!8}{\parbox{\dimexpr\linewidth-2\fboxsep\relax}{\raggedright\ttfamily\scriptsize     plt.yscale("log")}}
\par\smallskip
This creates a debug visualization. The plot is derived from generated CSV or Vivado report data so it can be regenerated after RTL changes.
\paragraph{Line 421}
\noindent\colorbox{gray!8}{\parbox{\dimexpr\linewidth-2\fboxsep\relax}{\raggedright\ttfamily\scriptsize     plt.xticks(range(len(layer\_names)), layer\_names, rotation=90, fontsize=6)}}
\par\smallskip
This creates a debug visualization. The plot is derived from generated CSV or Vivado report data so it can be regenerated after RTL changes.
\paragraph{Line 422}
\noindent\colorbox{gray!8}{\parbox{\dimexpr\linewidth-2\fboxsep\relax}{\raggedright\ttfamily\scriptsize     plt.ylabel("trained weight values (log scale)")}}
\par\smallskip
This creates a debug visualization. The plot is derived from generated CSV or Vivado report data so it can be regenerated after RTL changes.
\paragraph{Line 423}
\noindent\colorbox{gray!8}{\parbox{\dimexpr\linewidth-2\fboxsep\relax}{\raggedright\ttfamily\scriptsize     plt.tight\_layout()}}
\par\smallskip
This creates a debug visualization. The plot is derived from generated CSV or Vivado report data so it can be regenerated after RTL changes.
\paragraph{Line 424}
\noindent\colorbox{gray!8}{\parbox{\dimexpr\linewidth-2\fboxsep\relax}{\raggedright\ttfamily\scriptsize     plt.savefig(out\_dir / "spikingformer\_weight\_inventory.png")}}
\par\smallskip
This creates a debug visualization. The plot is derived from generated CSV or Vivado report data so it can be regenerated after RTL changes.
\paragraph{Line 425}
\noindent\colorbox{gray!8}{\parbox{\dimexpr\linewidth-2\fboxsep\relax}{\raggedright\ttfamily\scriptsize     plt.close()}}
\par\smallskip
This creates a debug visualization. The plot is derived from generated CSV or Vivado report data so it can be regenerated after RTL changes.
\paragraph{Line 426}
\noindent\colorbox{gray!8}{\parbox{\dimexpr\linewidth-2\fboxsep\relax}{\raggedright\ttfamily\scriptsize \emph{blank line}}}
\par\smallskip
This blank line is not hardware or software behavior. It separates nearby declarations or operations so the file remains readable and reviewable.
\paragraph{Line 427}
\noindent\colorbox{gray!8}{\parbox{\dimexpr\linewidth-2\fboxsep\relax}{\raggedright\ttfamily\scriptsize     with PdfPages(out\_dir / "spikingformer\_head\_accelerator\_report.pdf") as pdf:}}
\par\smallskip
This creates a debug visualization. The plot is derived from generated CSV or Vivado report data so it can be regenerated after RTL changes.
\paragraph{Line 428}
\noindent\colorbox{gray!8}{\parbox{\dimexpr\linewidth-2\fboxsep\relax}{\raggedright\ttfamily\scriptsize         fig = plt.figure(figsize=(10, 7), dpi=160)}}
\par\smallskip
This creates a debug visualization. The plot is derived from generated CSV or Vivado report data so it can be regenerated after RTL changes.
\paragraph{Line 429}
\noindent\colorbox{gray!8}{\parbox{\dimexpr\linewidth-2\fboxsep\relax}{\raggedright\ttfamily\scriptsize         fig.text(0.06, 0.93, "Trained Spikingformer Weight Mapping", fontsize=18, weight="bold")}}
\par\smallskip
This Python line participates in the export, validation, or report-generation pipeline. It is executed on the host PC to prepare data for RTL or to check RTL outputs.
\paragraph{Line 430}
\noindent\colorbox{gray!8}{\parbox{\dimexpr\linewidth-2\fboxsep\relax}{\raggedright\ttfamily\scriptsize         lines = [}}
\par\smallskip
This Python line participates in the export, validation, or report-generation pipeline. It is executed on the host PC to prepare data for RTL or to check RTL outputs.
\paragraph{Line 431}
\noindent\colorbox{gray!8}{\parbox{\dimexpr\linewidth-2\fboxsep\relax}{\raggedright\ttfamily\scriptsize             f"Checkpoint: \{Path(str(metrics['checkpoint'])).name\}",}}
\par\smallskip
This uses path objects so generated artifacts are addressed consistently across Windows shells and Python code.
\paragraph{Line 432}
\noindent\colorbox{gray!8}{\parbox{\dimexpr\linewidth-2\fboxsep\relax}{\raggedright\ttfamily\scriptsize             f"Samples exported: \{metrics['samples']\}",}}
\par\smallskip
This Python line participates in the export, validation, or report-generation pipeline. It is executed on the host PC to prepare data for RTL or to check RTL outputs.
\paragraph{Line 433}
\noindent\colorbox{gray!8}{\parbox{\dimexpr\linewidth-2\fboxsep\relax}{\raggedright\ttfamily\scriptsize             f"Float top-1 on exported set: \{float(metrics['float\_accuracy']) * 100:.2f\}\%",}}
\par\smallskip
This Python line participates in the export, validation, or report-generation pipeline. It is executed on the host PC to prepare data for RTL or to check RTL outputs.
\paragraph{Line 434}
\noindent\colorbox{gray!8}{\parbox{\dimexpr\linewidth-2\fboxsep\relax}{\raggedright\ttfamily\scriptsize             f"Fixed-point head top-1: \{float(metrics['fixed\_accuracy']) * 100:.2f\}\%",}}
\par\smallskip
This Python line participates in the export, validation, or report-generation pipeline. It is executed on the host PC to prepare data for RTL or to check RTL outputs.
\paragraph{Line 435}
\noindent\colorbox{gray!8}{\parbox{\dimexpr\linewidth-2\fboxsep\relax}{\raggedright\ttfamily\scriptsize             f"Float/fixed prediction agreement: \{float(metrics['float\_fixed\_prediction\_agreement']) * 100:.2f\}\%",}}
\par\smallskip
This line handles predicted labels. It checks that argmax behavior in Python and RTL chooses the same class.
\paragraph{Line 436}
\noindent\colorbox{gray!8}{\parbox{\dimexpr\linewidth-2\fboxsep\relax}{\raggedright\ttfamily\scriptsize             f"Feature gate: abs(feature) <= \{metrics['infer\_skip\_abs\_thresh']\} maps to zero",}}
\par\smallskip
This Python line participates in the export, validation, or report-generation pipeline. It is executed on the host PC to prepare data for RTL or to check RTL outputs.
\paragraph{Line 437}
\noindent\colorbox{gray!8}{\parbox{\dimexpr\linewidth-2\fboxsep\relax}{\raggedright\ttfamily\scriptsize             f"Feature-MAC skip opportunity: \{float(metrics['feature\_mac\_skip\_fraction']) * 100:.2f\}\%",}}
\par\smallskip
This Python line participates in the export, validation, or report-generation pipeline. It is executed on the host PC to prepare data for RTL or to check RTL outputs.
\paragraph{Line 438}
\noindent\colorbox{gray!8}{\parbox{\dimexpr\linewidth-2\fboxsep\relax}{\raggedright\ttfamily\scriptsize             f"Feature scale: \{metrics['feature\_scale']:.8g\}",}}
\par\smallskip
This Python line participates in the export, validation, or report-generation pipeline. It is executed on the host PC to prepare data for RTL or to check RTL outputs.
\paragraph{Line 439}
\noindent\colorbox{gray!8}{\parbox{\dimexpr\linewidth-2\fboxsep\relax}{\raggedright\ttfamily\scriptsize             f"Weight scale: \{metrics['weight\_scale']:.8g\}",}}
\par\smallskip
This Python line participates in the export, validation, or report-generation pipeline. It is executed on the host PC to prepare data for RTL or to check RTL outputs.
\paragraph{Line 440}
\noindent\colorbox{gray!8}{\parbox{\dimexpr\linewidth-2\fboxsep\relax}{\raggedright\ttfamily\scriptsize             f"Mapped trained parameters inventoried: \{mapping['parameter\_counts']['total']:,\}",}}
\par\smallskip
This Python line participates in the export, validation, or report-generation pipeline. It is executed on the host PC to prepare data for RTL or to check RTL outputs.
\paragraph{Line 441}
\noindent\colorbox{gray!8}{\parbox{\dimexpr\linewidth-2\fboxsep\relax}{\raggedright\ttfamily\scriptsize             "RTL scope in this run: trained 384-to-10 Spikingformer classifier head over exported backbone features.",}}
\par\smallskip
This Python line participates in the export, validation, or report-generation pipeline. It is executed on the host PC to prepare data for RTL or to check RTL outputs.
\paragraph{Line 442}
\noindent\colorbox{gray!8}{\parbox{\dimexpr\linewidth-2\fboxsep\relax}{\raggedright\ttfamily\scriptsize             "Full backbone scope captured in manifest: tokenizer convs, q/k/v/proj attention, MLP, BN-folded weights.",}}
\par\smallskip
This line handles the classifier parameters. Those arrays mirror the RTL memories, so mismatches reveal an export, quantization, or hardware-update bug.
\paragraph{Line 443}
\noindent\colorbox{gray!8}{\parbox{\dimexpr\linewidth-2\fboxsep\relax}{\raggedright\ttfamily\scriptsize         ]}}
\par\smallskip
This Python line participates in the export, validation, or report-generation pipeline. It is executed on the host PC to prepare data for RTL or to check RTL outputs.
\paragraph{Line 444}
\noindent\colorbox{gray!8}{\parbox{\dimexpr\linewidth-2\fboxsep\relax}{\raggedright\ttfamily\scriptsize         fig.text(0.06, 0.83, "\textbackslash{}n".join(lines), fontsize=11, va="top")}}
\par\smallskip
This Python line participates in the export, validation, or report-generation pipeline. It is executed on the host PC to prepare data for RTL or to check RTL outputs.
\paragraph{Line 445}
\noindent\colorbox{gray!8}{\parbox{\dimexpr\linewidth-2\fboxsep\relax}{\raggedright\ttfamily\scriptsize         pdf.savefig(fig)}}
\par\smallskip
This Python line participates in the export, validation, or report-generation pipeline. It is executed on the host PC to prepare data for RTL or to check RTL outputs.
\paragraph{Line 446}
\noindent\colorbox{gray!8}{\parbox{\dimexpr\linewidth-2\fboxsep\relax}{\raggedright\ttfamily\scriptsize         plt.close(fig)}}
\par\smallskip
This creates a debug visualization. The plot is derived from generated CSV or Vivado report data so it can be regenerated after RTL changes.
\paragraph{Line 447}
\noindent\colorbox{gray!8}{\parbox{\dimexpr\linewidth-2\fboxsep\relax}{\raggedright\ttfamily\scriptsize         for name in ["spikingformer\_head\_fixed\_confusion.png", "spikingformer\_weight\_inventory.png"]:}}
\par\smallskip
This Python line participates in the export, validation, or report-generation pipeline. It is executed on the host PC to prepare data for RTL or to check RTL outputs.
\paragraph{Line 448}
\noindent\colorbox{gray!8}{\parbox{\dimexpr\linewidth-2\fboxsep\relax}{\raggedright\ttfamily\scriptsize             img = plt.imread(out\_dir / name)}}
\par\smallskip
This creates a debug visualization. The plot is derived from generated CSV or Vivado report data so it can be regenerated after RTL changes.
\paragraph{Line 449}
\noindent\colorbox{gray!8}{\parbox{\dimexpr\linewidth-2\fboxsep\relax}{\raggedright\ttfamily\scriptsize             fig, ax = plt.subplots(figsize=(10, 7), dpi=160)}}
\par\smallskip
This creates a debug visualization. The plot is derived from generated CSV or Vivado report data so it can be regenerated after RTL changes.
\paragraph{Line 450}
\noindent\colorbox{gray!8}{\parbox{\dimexpr\linewidth-2\fboxsep\relax}{\raggedright\ttfamily\scriptsize             ax.imshow(img)}}
\par\smallskip
This Python line participates in the export, validation, or report-generation pipeline. It is executed on the host PC to prepare data for RTL or to check RTL outputs.
\paragraph{Line 451}
\noindent\colorbox{gray!8}{\parbox{\dimexpr\linewidth-2\fboxsep\relax}{\raggedright\ttfamily\scriptsize             ax.axis("off")}}
\par\smallskip
This Python line participates in the export, validation, or report-generation pipeline. It is executed on the host PC to prepare data for RTL or to check RTL outputs.
\paragraph{Line 452}
\noindent\colorbox{gray!8}{\parbox{\dimexpr\linewidth-2\fboxsep\relax}{\raggedright\ttfamily\scriptsize             pdf.savefig(fig)}}
\par\smallskip
This Python line participates in the export, validation, or report-generation pipeline. It is executed on the host PC to prepare data for RTL or to check RTL outputs.
\paragraph{Line 453}
\noindent\colorbox{gray!8}{\parbox{\dimexpr\linewidth-2\fboxsep\relax}{\raggedright\ttfamily\scriptsize             plt.close(fig)}}
\par\smallskip
This creates a debug visualization. The plot is derived from generated CSV or Vivado report data so it can be regenerated after RTL changes.
\paragraph{Line 454}
\noindent\colorbox{gray!8}{\parbox{\dimexpr\linewidth-2\fboxsep\relax}{\raggedright\ttfamily\scriptsize \emph{blank line}}}
\par\smallskip
This blank line is not hardware or software behavior. It separates nearby declarations or operations so the file remains readable and reviewable.
\paragraph{Line 455}
\noindent\colorbox{gray!8}{\parbox{\dimexpr\linewidth-2\fboxsep\relax}{\raggedright\ttfamily\scriptsize \emph{blank line}}}
\par\smallskip
This blank line is not hardware or software behavior. It separates nearby declarations or operations so the file remains readable and reviewable.
\paragraph{Line 456}
\noindent\colorbox{gray!8}{\parbox{\dimexpr\linewidth-2\fboxsep\relax}{\raggedright\ttfamily\scriptsize def parse\_args() -> argparse.Namespace:}}
\par\smallskip
This defines a Python function. The function packages one repeatable step such as loading a memory file, simulating the golden learning rule, validating CSVs, or generating plots.
\paragraph{Line 457}
\noindent\colorbox{gray!8}{\parbox{\dimexpr\linewidth-2\fboxsep\relax}{\raggedright\ttfamily\scriptsize     root = repo\_root()}}
\par\smallskip
This Python line participates in the export, validation, or report-generation pipeline. It is executed on the host PC to prepare data for RTL or to check RTL outputs.
\paragraph{Line 458}
\noindent\colorbox{gray!8}{\parbox{\dimexpr\linewidth-2\fboxsep\relax}{\raggedright\ttfamily\scriptsize     parser = argparse.ArgumentParser(description="Export trained Spikingformer CIFAR-10 weights/features for RTL acceleration.")}}
\par\smallskip
This contributes to the command-line interface. It lets the rerun flow choose paths, sample counts, or output locations without editing code.
\paragraph{Line 459}
\noindent\colorbox{gray!8}{\parbox{\dimexpr\linewidth-2\fboxsep\relax}{\raggedright\ttfamily\scriptsize     parser.add\_argument("--checkpoint", type=Path, default=root / "cifar10\_neuromorphic" / "outputs\_upstream\_spikingformer" / "spikingformer\_4\_384\_400e" / "model\_best.pth.tar")}}
\par\smallskip
This contributes to the command-line interface. It lets the rerun flow choose paths, sample counts, or output locations without editing code.
\paragraph{Line 460}
\noindent\colorbox{gray!8}{\parbox{\dimexpr\linewidth-2\fboxsep\relax}{\raggedright\ttfamily\scriptsize     parser.add\_argument("--data-dir", type=Path, default=root / "datasets")}}
\par\smallskip
This contributes to the command-line interface. It lets the rerun flow choose paths, sample counts, or output locations without editing code.
\paragraph{Line 461}
\noindent\colorbox{gray!8}{\parbox{\dimexpr\linewidth-2\fboxsep\relax}{\raggedright\ttfamily\scriptsize     parser.add\_argument("--rtl-dir", type=Path, default=root / "cifar10\_neuromorphic" / "rtl")}}
\par\smallskip
This contributes to the command-line interface. It lets the rerun flow choose paths, sample counts, or output locations without editing code.
\paragraph{Line 462}
\noindent\colorbox{gray!8}{\parbox{\dimexpr\linewidth-2\fboxsep\relax}{\raggedright\ttfamily\scriptsize     parser.add\_argument("--out-dir", type=Path, default=root / "cifar10\_neuromorphic" / "outputs\_spikingformer\_head\_accelerator")}}
\par\smallskip
This contributes to the command-line interface. It lets the rerun flow choose paths, sample counts, or output locations without editing code.
\paragraph{Line 463}
\noindent\colorbox{gray!8}{\parbox{\dimexpr\linewidth-2\fboxsep\relax}{\raggedright\ttfamily\scriptsize     parser.add\_argument("--samples", type=int, default=10000)}}
\par\smallskip
This contributes to the command-line interface. It lets the rerun flow choose paths, sample counts, or output locations without editing code.
\paragraph{Line 464}
\noindent\colorbox{gray!8}{\parbox{\dimexpr\linewidth-2\fboxsep\relax}{\raggedright\ttfamily\scriptsize     parser.add\_argument("--trace-samples", type=int, default=8)}}
\par\smallskip
This contributes to the command-line interface. It lets the rerun flow choose paths, sample counts, or output locations without editing code.
\paragraph{Line 465}
\noindent\colorbox{gray!8}{\parbox{\dimexpr\linewidth-2\fboxsep\relax}{\raggedright\ttfamily\scriptsize     parser.add\_argument("--batch-size", type=int, default=64)}}
\par\smallskip
This contributes to the command-line interface. It lets the rerun flow choose paths, sample counts, or output locations without editing code.
\paragraph{Line 466}
\noindent\colorbox{gray!8}{\parbox{\dimexpr\linewidth-2\fboxsep\relax}{\raggedright\ttfamily\scriptsize     parser.add\_argument("--workers", type=int, default=4)}}
\par\smallskip
This contributes to the command-line interface. It lets the rerun flow choose paths, sample counts, or output locations without editing code.
\paragraph{Line 467}
\noindent\colorbox{gray!8}{\parbox{\dimexpr\linewidth-2\fboxsep\relax}{\raggedright\ttfamily\scriptsize     parser.add\_argument("--device", type=str, default="auto")}}
\par\smallskip
This contributes to the command-line interface. It lets the rerun flow choose paths, sample counts, or output locations without editing code.
\paragraph{Line 468}
\noindent\colorbox{gray!8}{\parbox{\dimexpr\linewidth-2\fboxsep\relax}{\raggedright\ttfamily\scriptsize     parser.add\_argument("--amp", action=argparse.BooleanOptionalAction, default=True)}}
\par\smallskip
This contributes to the command-line interface. It lets the rerun flow choose paths, sample counts, or output locations without editing code.
\paragraph{Line 469}
\noindent\colorbox{gray!8}{\parbox{\dimexpr\linewidth-2\fboxsep\relax}{\raggedright\ttfamily\scriptsize     parser.add\_argument("--feature-bits", type=int, default=16)}}
\par\smallskip
This contributes to the command-line interface. It lets the rerun flow choose paths, sample counts, or output locations without editing code.
\paragraph{Line 470}
\noindent\colorbox{gray!8}{\parbox{\dimexpr\linewidth-2\fboxsep\relax}{\raggedright\ttfamily\scriptsize     parser.add\_argument("--weight-bits", type=int, default=16)}}
\par\smallskip
This contributes to the command-line interface. It lets the rerun flow choose paths, sample counts, or output locations without editing code.
\paragraph{Line 471}
\noindent\colorbox{gray!8}{\parbox{\dimexpr\linewidth-2\fboxsep\relax}{\raggedright\ttfamily\scriptsize     parser.add\_argument("--infer-skip-abs-thresh", type=int, default=2048)}}
\par\smallskip
This contributes to the command-line interface. It lets the rerun flow choose paths, sample counts, or output locations without editing code.
\paragraph{Line 472}
\noindent\colorbox{gray!8}{\parbox{\dimexpr\linewidth-2\fboxsep\relax}{\raggedright\ttfamily\scriptsize     return parser.parse\_args()}}
\par\smallskip
This Python line participates in the export, validation, or report-generation pipeline. It is executed on the host PC to prepare data for RTL or to check RTL outputs.
\paragraph{Line 473}
\noindent\colorbox{gray!8}{\parbox{\dimexpr\linewidth-2\fboxsep\relax}{\raggedright\ttfamily\scriptsize \emph{blank line}}}
\par\smallskip
This blank line is not hardware or software behavior. It separates nearby declarations or operations so the file remains readable and reviewable.
\paragraph{Line 474}
\noindent\colorbox{gray!8}{\parbox{\dimexpr\linewidth-2\fboxsep\relax}{\raggedright\ttfamily\scriptsize \emph{blank line}}}
\par\smallskip
This blank line is not hardware or software behavior. It separates nearby declarations or operations so the file remains readable and reviewable.
\paragraph{Line 475}
\noindent\colorbox{gray!8}{\parbox{\dimexpr\linewidth-2\fboxsep\relax}{\raggedright\ttfamily\scriptsize def main() -> None:}}
\par\smallskip
This defines a Python function. The function packages one repeatable step such as loading a memory file, simulating the golden learning rule, validating CSVs, or generating plots.
\paragraph{Line 476}
\noindent\colorbox{gray!8}{\parbox{\dimexpr\linewidth-2\fboxsep\relax}{\raggedright\ttfamily\scriptsize     args = parse\_args()}}
\par\smallskip
This Python line participates in the export, validation, or report-generation pipeline. It is executed on the host PC to prepare data for RTL or to check RTL outputs.
\paragraph{Line 477}
\noindent\colorbox{gray!8}{\parbox{\dimexpr\linewidth-2\fboxsep\relax}{\raggedright\ttfamily\scriptsize     features, logits, labels = collect\_features(args)}}
\par\smallskip
This Python line participates in the export, validation, or report-generation pipeline. It is executed on the host PC to prepare data for RTL or to check RTL outputs.
\paragraph{Line 478}
\noindent\colorbox{gray!8}{\parbox{\dimexpr\linewidth-2\fboxsep\relax}{\raggedright\ttfamily\scriptsize     mapping = full\_checkpoint\_mapping(args)}}
\par\smallskip
This Python line participates in the export, validation, or report-generation pipeline. It is executed on the host PC to prepare data for RTL or to check RTL outputs.
\paragraph{Line 479}
\noindent\colorbox{gray!8}{\parbox{\dimexpr\linewidth-2\fboxsep\relax}{\raggedright\ttfamily\scriptsize     args.out\_dir.mkdir(parents=True, exist\_ok=True)}}
\par\smallskip
This Python line participates in the export, validation, or report-generation pipeline. It is executed on the host PC to prepare data for RTL or to check RTL outputs.
\paragraph{Line 480}
\noindent\colorbox{gray!8}{\parbox{\dimexpr\linewidth-2\fboxsep\relax}{\raggedright\ttfamily\scriptsize     metrics = write\_head\_memories(args, features, logits, labels)}}
\par\smallskip
This Python line participates in the export, validation, or report-generation pipeline. It is executed on the host PC to prepare data for RTL or to check RTL outputs.
\paragraph{Line 481}
\noindent\colorbox{gray!8}{\parbox{\dimexpr\linewidth-2\fboxsep\relax}{\raggedright\ttfamily\scriptsize     mapping["exported\_float\_top1\_pct"] = float(metrics["float\_accuracy"]) * 100.0}}
\par\smallskip
This Python line participates in the export, validation, or report-generation pipeline. It is executed on the host PC to prepare data for RTL or to check RTL outputs.
\paragraph{Line 482}
\noindent\colorbox{gray!8}{\parbox{\dimexpr\linewidth-2\fboxsep\relax}{\raggedright\ttfamily\scriptsize     mapping["exported\_fixed\_top1\_pct"] = float(metrics["fixed\_accuracy"]) * 100.0}}
\par\smallskip
This Python line participates in the export, validation, or report-generation pipeline. It is executed on the host PC to prepare data for RTL or to check RTL outputs.
\paragraph{Line 483}
\noindent\colorbox{gray!8}{\parbox{\dimexpr\linewidth-2\fboxsep\relax}{\raggedright\ttfamily\scriptsize     mapping["infer\_skip\_abs\_thresh"] = int(metrics["infer\_skip\_abs\_thresh"])}}
\par\smallskip
This Python line participates in the export, validation, or report-generation pipeline. It is executed on the host PC to prepare data for RTL or to check RTL outputs.
\paragraph{Line 484}
\noindent\colorbox{gray!8}{\parbox{\dimexpr\linewidth-2\fboxsep\relax}{\raggedright\ttfamily\scriptsize     mapping["feature\_mac\_skip\_fraction\_pct"] = float(metrics["feature\_mac\_skip\_fraction"]) * 100.0}}
\par\smallskip
This Python line participates in the export, validation, or report-generation pipeline. It is executed on the host PC to prepare data for RTL or to check RTL outputs.
\paragraph{Line 485}
\noindent\colorbox{gray!8}{\parbox{\dimexpr\linewidth-2\fboxsep\relax}{\raggedright\ttfamily\scriptsize     (args.out\_dir / "spikingformer\_full\_checkpoint\_mapping.json").write\_text(json.dumps(mapping, indent=2), encoding="utf-8")}}
\par\smallskip
This writes machine-readable output for later reports. Keeping JSON/CSV artifacts makes the verification auditable.
\paragraph{Line 486}
\noindent\colorbox{gray!8}{\parbox{\dimexpr\linewidth-2\fboxsep\relax}{\raggedright\ttfamily\scriptsize     write\_report(args, metrics, mapping)}}
\par\smallskip
This Python line participates in the export, validation, or report-generation pipeline. It is executed on the host PC to prepare data for RTL or to check RTL outputs.
\paragraph{Line 487}
\noindent\colorbox{gray!8}{\parbox{\dimexpr\linewidth-2\fboxsep\relax}{\raggedright\ttfamily\scriptsize     print(json.dumps(\{k: v for k, v in metrics.items() if k not in \{"confusion", "files"\}\}, indent=2))}}
\par\smallskip
This writes machine-readable output for later reports. Keeping JSON/CSV artifacts makes the verification auditable.
\paragraph{Line 488}
\noindent\colorbox{gray!8}{\parbox{\dimexpr\linewidth-2\fboxsep\relax}{\raggedright\ttfamily\scriptsize     print(f"outputs=\{args.out\_dir\}")}}
\par\smallskip
This Python line participates in the export, validation, or report-generation pipeline. It is executed on the host PC to prepare data for RTL or to check RTL outputs.
\paragraph{Line 489}
\noindent\colorbox{gray!8}{\parbox{\dimexpr\linewidth-2\fboxsep\relax}{\raggedright\ttfamily\scriptsize \emph{blank line}}}
\par\smallskip
This blank line is not hardware or software behavior. It separates nearby declarations or operations so the file remains readable and reviewable.
\paragraph{Line 490}
\noindent\colorbox{gray!8}{\parbox{\dimexpr\linewidth-2\fboxsep\relax}{\raggedright\ttfamily\scriptsize \emph{blank line}}}
\par\smallskip
This blank line is not hardware or software behavior. It separates nearby declarations or operations so the file remains readable and reviewable.
\paragraph{Line 491}
\noindent\colorbox{gray!8}{\parbox{\dimexpr\linewidth-2\fboxsep\relax}{\raggedright\ttfamily\scriptsize if \_\_name\_\_ == "\_\_main\_\_":}}
\par\smallskip
This Python idiom runs main() only when the file is executed as a script. It still allows the file's functions to be imported for tests or reuse.
\paragraph{Line 492}
\noindent\colorbox{gray!8}{\parbox{\dimexpr\linewidth-2\fboxsep\relax}{\raggedright\ttfamily\scriptsize     main()}}
\par\smallskip
This Python line participates in the export, validation, or report-generation pipeline. It is executed on the host PC to prepare data for RTL or to check RTL outputs.
\clearpage
\subsection{cifar10\_neuromorphic\textbackslash{}validate\_spikingformer\_head\_rtl.py}
Each entry below gives the exact source line followed by its hardware/software role in the implemented accelerator pipeline.
\paragraph{Line 1}
\noindent\colorbox{gray!8}{\parbox{\dimexpr\linewidth-2\fboxsep\relax}{\raggedright\ttfamily\scriptsize from \_\_future\_\_ import annotations}}
\par\smallskip
This imports a Python library or symbol. The scripts use standard parsing, numeric arrays, plotting, and file APIs to make the hardware artifacts reproducible.
\paragraph{Line 2}
\noindent\colorbox{gray!8}{\parbox{\dimexpr\linewidth-2\fboxsep\relax}{\raggedright\ttfamily\scriptsize \emph{blank line}}}
\par\smallskip
This blank line is not hardware or software behavior. It separates nearby declarations or operations so the file remains readable and reviewable.
\paragraph{Line 3}
\noindent\colorbox{gray!8}{\parbox{\dimexpr\linewidth-2\fboxsep\relax}{\raggedright\ttfamily\scriptsize import argparse}}
\par\smallskip
This imports a Python library or symbol. The scripts use standard parsing, numeric arrays, plotting, and file APIs to make the hardware artifacts reproducible.
\paragraph{Line 4}
\noindent\colorbox{gray!8}{\parbox{\dimexpr\linewidth-2\fboxsep\relax}{\raggedright\ttfamily\scriptsize import csv}}
\par\smallskip
This imports a Python library or symbol. The scripts use standard parsing, numeric arrays, plotting, and file APIs to make the hardware artifacts reproducible.
\paragraph{Line 5}
\noindent\colorbox{gray!8}{\parbox{\dimexpr\linewidth-2\fboxsep\relax}{\raggedright\ttfamily\scriptsize import json}}
\par\smallskip
This imports a Python library or symbol. The scripts use standard parsing, numeric arrays, plotting, and file APIs to make the hardware artifacts reproducible.
\paragraph{Line 6}
\noindent\colorbox{gray!8}{\parbox{\dimexpr\linewidth-2\fboxsep\relax}{\raggedright\ttfamily\scriptsize from pathlib import Path}}
\par\smallskip
This imports a Python library or symbol. The scripts use standard parsing, numeric arrays, plotting, and file APIs to make the hardware artifacts reproducible.
\paragraph{Line 7}
\noindent\colorbox{gray!8}{\parbox{\dimexpr\linewidth-2\fboxsep\relax}{\raggedright\ttfamily\scriptsize \emph{blank line}}}
\par\smallskip
This blank line is not hardware or software behavior. It separates nearby declarations or operations so the file remains readable and reviewable.
\paragraph{Line 8}
\noindent\colorbox{gray!8}{\parbox{\dimexpr\linewidth-2\fboxsep\relax}{\raggedright\ttfamily\scriptsize import matplotlib.pyplot as plt}}
\par\smallskip
This imports a Python library or symbol. The scripts use standard parsing, numeric arrays, plotting, and file APIs to make the hardware artifacts reproducible.
\paragraph{Line 9}
\noindent\colorbox{gray!8}{\parbox{\dimexpr\linewidth-2\fboxsep\relax}{\raggedright\ttfamily\scriptsize import numpy as np}}
\par\smallskip
This imports a Python library or symbol. The scripts use standard parsing, numeric arrays, plotting, and file APIs to make the hardware artifacts reproducible.
\paragraph{Line 10}
\noindent\colorbox{gray!8}{\parbox{\dimexpr\linewidth-2\fboxsep\relax}{\raggedright\ttfamily\scriptsize from matplotlib.backends.backend\_pdf import PdfPages}}
\par\smallskip
This imports a Python library or symbol. The scripts use standard parsing, numeric arrays, plotting, and file APIs to make the hardware artifacts reproducible.
\paragraph{Line 11}
\noindent\colorbox{gray!8}{\parbox{\dimexpr\linewidth-2\fboxsep\relax}{\raggedright\ttfamily\scriptsize \emph{blank line}}}
\par\smallskip
This blank line is not hardware or software behavior. It separates nearby declarations or operations so the file remains readable and reviewable.
\paragraph{Line 12}
\noindent\colorbox{gray!8}{\parbox{\dimexpr\linewidth-2\fboxsep\relax}{\raggedright\ttfamily\scriptsize \emph{blank line}}}
\par\smallskip
This blank line is not hardware or software behavior. It separates nearby declarations or operations so the file remains readable and reviewable.
\paragraph{Line 13}
\noindent\colorbox{gray!8}{\parbox{\dimexpr\linewidth-2\fboxsep\relax}{\raggedright\ttfamily\scriptsize NUM\_CLASSES = 10}}
\par\smallskip
This Python line participates in the export, validation, or report-generation pipeline. It is executed on the host PC to prepare data for RTL or to check RTL outputs.
\paragraph{Line 14}
\noindent\colorbox{gray!8}{\parbox{\dimexpr\linewidth-2\fboxsep\relax}{\raggedright\ttfamily\scriptsize FEATURE\_DIM = 384}}
\par\smallskip
This Python line participates in the export, validation, or report-generation pipeline. It is executed on the host PC to prepare data for RTL or to check RTL outputs.
\paragraph{Line 15}
\noindent\colorbox{gray!8}{\parbox{\dimexpr\linewidth-2\fboxsep\relax}{\raggedright\ttfamily\scriptsize INFER\_SKIP\_ABS\_THRESH = 2048}}
\par\smallskip
This Python line participates in the export, validation, or report-generation pipeline. It is executed on the host PC to prepare data for RTL or to check RTL outputs.
\paragraph{Line 16}
\noindent\colorbox{gray!8}{\parbox{\dimexpr\linewidth-2\fboxsep\relax}{\raggedright\ttfamily\scriptsize \emph{blank line}}}
\par\smallskip
This blank line is not hardware or software behavior. It separates nearby declarations or operations so the file remains readable and reviewable.
\paragraph{Line 17}
\noindent\colorbox{gray!8}{\parbox{\dimexpr\linewidth-2\fboxsep\relax}{\raggedright\ttfamily\scriptsize \emph{blank line}}}
\par\smallskip
This blank line is not hardware or software behavior. It separates nearby declarations or operations so the file remains readable and reviewable.
\paragraph{Line 18}
\noindent\colorbox{gray!8}{\parbox{\dimexpr\linewidth-2\fboxsep\relax}{\raggedright\ttfamily\scriptsize def repo\_root() -> Path:}}
\par\smallskip
This defines a Python function. The function packages one repeatable step such as loading a memory file, simulating the golden learning rule, validating CSVs, or generating plots.
\paragraph{Line 19}
\noindent\colorbox{gray!8}{\parbox{\dimexpr\linewidth-2\fboxsep\relax}{\raggedright\ttfamily\scriptsize     return Path(\_\_file\_\_).resolve().parents[1]}}
\par\smallskip
This uses path objects so generated artifacts are addressed consistently across Windows shells and Python code.
\paragraph{Line 20}
\noindent\colorbox{gray!8}{\parbox{\dimexpr\linewidth-2\fboxsep\relax}{\raggedright\ttfamily\scriptsize \emph{blank line}}}
\par\smallskip
This blank line is not hardware or software behavior. It separates nearby declarations or operations so the file remains readable and reviewable.
\paragraph{Line 21}
\noindent\colorbox{gray!8}{\parbox{\dimexpr\linewidth-2\fboxsep\relax}{\raggedright\ttfamily\scriptsize \emph{blank line}}}
\par\smallskip
This blank line is not hardware or software behavior. It separates nearby declarations or operations so the file remains readable and reviewable.
\paragraph{Line 22}
\noindent\colorbox{gray!8}{\parbox{\dimexpr\linewidth-2\fboxsep\relax}{\raggedright\ttfamily\scriptsize def load\_rows(path: Path) -> list[dict[str, str]]:}}
\par\smallskip
This defines a Python function. The function packages one repeatable step such as loading a memory file, simulating the golden learning rule, validating CSVs, or generating plots.
\paragraph{Line 23}
\noindent\colorbox{gray!8}{\parbox{\dimexpr\linewidth-2\fboxsep\relax}{\raggedright\ttfamily\scriptsize     with path.open("r", newline="", encoding="utf-8") as fh:}}
\par\smallskip
This reads an artifact produced by training, export, simulation, or Vivado. The validator/report depends on files generated by the actual run rather than hidden constants.
\paragraph{Line 24}
\noindent\colorbox{gray!8}{\parbox{\dimexpr\linewidth-2\fboxsep\relax}{\raggedright\ttfamily\scriptsize         return list(csv.DictReader(fh))}}
\par\smallskip
This Python line participates in the export, validation, or report-generation pipeline. It is executed on the host PC to prepare data for RTL or to check RTL outputs.
\paragraph{Line 25}
\noindent\colorbox{gray!8}{\parbox{\dimexpr\linewidth-2\fboxsep\relax}{\raggedright\ttfamily\scriptsize \emph{blank line}}}
\par\smallskip
This blank line is not hardware or software behavior. It separates nearby declarations or operations so the file remains readable and reviewable.
\paragraph{Line 26}
\noindent\colorbox{gray!8}{\parbox{\dimexpr\linewidth-2\fboxsep\relax}{\raggedright\ttfamily\scriptsize \emph{blank line}}}
\par\smallskip
This blank line is not hardware or software behavior. It separates nearby declarations or operations so the file remains readable and reviewable.
\paragraph{Line 27}
\noindent\colorbox{gray!8}{\parbox{\dimexpr\linewidth-2\fboxsep\relax}{\raggedright\ttfamily\scriptsize def signed\_hex\_to\_int(token: str, bits: int) -> int:}}
\par\smallskip
This defines a Python function. The function packages one repeatable step such as loading a memory file, simulating the golden learning rule, validating CSVs, or generating plots.
\paragraph{Line 28}
\noindent\colorbox{gray!8}{\parbox{\dimexpr\linewidth-2\fboxsep\relax}{\raggedright\ttfamily\scriptsize     value = int(token.strip(), 16)}}
\par\smallskip
This Python line participates in the export, validation, or report-generation pipeline. It is executed on the host PC to prepare data for RTL or to check RTL outputs.
\paragraph{Line 29}
\noindent\colorbox{gray!8}{\parbox{\dimexpr\linewidth-2\fboxsep\relax}{\raggedright\ttfamily\scriptsize     sign = 1 << (bits - 1)}}
\par\smallskip
This Python line participates in the export, validation, or report-generation pipeline. It is executed on the host PC to prepare data for RTL or to check RTL outputs.
\paragraph{Line 30}
\noindent\colorbox{gray!8}{\parbox{\dimexpr\linewidth-2\fboxsep\relax}{\raggedright\ttfamily\scriptsize     return value - (1 << bits) if value \& sign else value}}
\par\smallskip
This Python line participates in the export, validation, or report-generation pipeline. It is executed on the host PC to prepare data for RTL or to check RTL outputs.
\paragraph{Line 31}
\noindent\colorbox{gray!8}{\parbox{\dimexpr\linewidth-2\fboxsep\relax}{\raggedright\ttfamily\scriptsize \emph{blank line}}}
\par\smallskip
This blank line is not hardware or software behavior. It separates nearby declarations or operations so the file remains readable and reviewable.
\paragraph{Line 32}
\noindent\colorbox{gray!8}{\parbox{\dimexpr\linewidth-2\fboxsep\relax}{\raggedright\ttfamily\scriptsize \emph{blank line}}}
\par\smallskip
This blank line is not hardware or software behavior. It separates nearby declarations or operations so the file remains readable and reviewable.
\paragraph{Line 33}
\noindent\colorbox{gray!8}{\parbox{\dimexpr\linewidth-2\fboxsep\relax}{\raggedright\ttfamily\scriptsize def load\_mem(path: Path, bits: int) -> list[int]:}}
\par\smallskip
This defines a Python function. The function packages one repeatable step such as loading a memory file, simulating the golden learning rule, validating CSVs, or generating plots.
\paragraph{Line 34}
\noindent\colorbox{gray!8}{\parbox{\dimexpr\linewidth-2\fboxsep\relax}{\raggedright\ttfamily\scriptsize     return [signed\_hex\_to\_int(line, bits) for line in path.read\_text(encoding="ascii").splitlines() if line.strip()]}}
\par\smallskip
This reads an artifact produced by training, export, simulation, or Vivado. The validator/report depends on files generated by the actual run rather than hidden constants.
\paragraph{Line 35}
\noindent\colorbox{gray!8}{\parbox{\dimexpr\linewidth-2\fboxsep\relax}{\raggedright\ttfamily\scriptsize \emph{blank line}}}
\par\smallskip
This blank line is not hardware or software behavior. It separates nearby declarations or operations so the file remains readable and reviewable.
\paragraph{Line 36}
\noindent\colorbox{gray!8}{\parbox{\dimexpr\linewidth-2\fboxsep\relax}{\raggedright\ttfamily\scriptsize \emph{blank line}}}
\par\smallskip
This blank line is not hardware or software behavior. It separates nearby declarations or operations so the file remains readable and reviewable.
\paragraph{Line 37}
\noindent\colorbox{gray!8}{\parbox{\dimexpr\linewidth-2\fboxsep\relax}{\raggedright\ttfamily\scriptsize def fixed\_reference\_from\_mem(rtl\_dir: Path, sample\_count: int) -> tuple[np.ndarray, np.ndarray, np.ndarray, float]:}}
\par\smallskip
This defines a Python function. The function packages one repeatable step such as loading a memory file, simulating the golden learning rule, validating CSVs, or generating plots.
\paragraph{Line 38}
\noindent\colorbox{gray!8}{\parbox{\dimexpr\linewidth-2\fboxsep\relax}{\raggedright\ttfamily\scriptsize     features\_flat = np.array(load\_mem(rtl\_dir / "cifar10\_spikingformer\_head\_features.mem", 16), dtype=np.int64)}}
\par\smallskip
This performs vectorized numerical work. The fixed-point golden models use NumPy so score, margin, and update calculations are explicit and repeatable.
\paragraph{Line 39}
\noindent\colorbox{gray!8}{\parbox{\dimexpr\linewidth-2\fboxsep\relax}{\raggedright\ttfamily\scriptsize     weights = np.array(load\_mem(rtl\_dir / "cifar10\_spikingformer\_head\_weights.mem", 16), dtype=np.int64).reshape(NUM\_CLASSES, FEATURE\_DIM)}}
\par\smallskip
This performs vectorized numerical work. The fixed-point golden models use NumPy so score, margin, and update calculations are explicit and repeatable.
\paragraph{Line 40}
\noindent\colorbox{gray!8}{\parbox{\dimexpr\linewidth-2\fboxsep\relax}{\raggedright\ttfamily\scriptsize     bias = np.array(load\_mem(rtl\_dir / "cifar10\_spikingformer\_head\_bias.mem", 64), dtype=np.int64)}}
\par\smallskip
This performs vectorized numerical work. The fixed-point golden models use NumPy so score, margin, and update calculations are explicit and repeatable.
\paragraph{Line 41}
\noindent\colorbox{gray!8}{\parbox{\dimexpr\linewidth-2\fboxsep\relax}{\raggedright\ttfamily\scriptsize     labels = np.array([int(line.strip(), 16) for line in (rtl\_dir / "cifar10\_spikingformer\_head\_labels.mem").read\_text(encoding="ascii").splitlines() if line.strip()], dtype=np.int64)}}
\par\smallskip
This reads an artifact produced by training, export, simulation, or Vivado. The validator/report depends on files generated by the actual run rather than hidden constants.
\paragraph{Line 42}
\noindent\colorbox{gray!8}{\parbox{\dimexpr\linewidth-2\fboxsep\relax}{\raggedright\ttfamily\scriptsize     features = features\_flat[: sample\_count * FEATURE\_DIM].reshape(sample\_count, FEATURE\_DIM)}}
\par\smallskip
This Python line participates in the export, validation, or report-generation pipeline. It is executed on the host PC to prepare data for RTL or to check RTL outputs.
\paragraph{Line 43}
\noindent\colorbox{gray!8}{\parbox{\dimexpr\linewidth-2\fboxsep\relax}{\raggedright\ttfamily\scriptsize     if INFER\_SKIP\_ABS\_THRESH > 0:}}
\par\smallskip
This Python line participates in the export, validation, or report-generation pipeline. It is executed on the host PC to prepare data for RTL or to check RTL outputs.
\paragraph{Line 44}
\noindent\colorbox{gray!8}{\parbox{\dimexpr\linewidth-2\fboxsep\relax}{\raggedright\ttfamily\scriptsize         skipped = np.abs(features) <= INFER\_SKIP\_ABS\_THRESH}}
\par\smallskip
This performs vectorized numerical work. The fixed-point golden models use NumPy so score, margin, and update calculations are explicit and repeatable.
\paragraph{Line 45}
\noindent\colorbox{gray!8}{\parbox{\dimexpr\linewidth-2\fboxsep\relax}{\raggedright\ttfamily\scriptsize         features = features.copy()}}
\par\smallskip
This Python line participates in the export, validation, or report-generation pipeline. It is executed on the host PC to prepare data for RTL or to check RTL outputs.
\paragraph{Line 46}
\noindent\colorbox{gray!8}{\parbox{\dimexpr\linewidth-2\fboxsep\relax}{\raggedright\ttfamily\scriptsize         features[skipped] = 0}}
\par\smallskip
This Python line participates in the export, validation, or report-generation pipeline. It is executed on the host PC to prepare data for RTL or to check RTL outputs.
\paragraph{Line 47}
\noindent\colorbox{gray!8}{\parbox{\dimexpr\linewidth-2\fboxsep\relax}{\raggedright\ttfamily\scriptsize         skip\_fraction = float(skipped.mean())}}
\par\smallskip
This Python line participates in the export, validation, or report-generation pipeline. It is executed on the host PC to prepare data for RTL or to check RTL outputs.
\paragraph{Line 48}
\noindent\colorbox{gray!8}{\parbox{\dimexpr\linewidth-2\fboxsep\relax}{\raggedright\ttfamily\scriptsize     else:}}
\par\smallskip
This Python line participates in the export, validation, or report-generation pipeline. It is executed on the host PC to prepare data for RTL or to check RTL outputs.
\paragraph{Line 49}
\noindent\colorbox{gray!8}{\parbox{\dimexpr\linewidth-2\fboxsep\relax}{\raggedright\ttfamily\scriptsize         skip\_fraction = 0.0}}
\par\smallskip
This Python line participates in the export, validation, or report-generation pipeline. It is executed on the host PC to prepare data for RTL or to check RTL outputs.
\paragraph{Line 50}
\noindent\colorbox{gray!8}{\parbox{\dimexpr\linewidth-2\fboxsep\relax}{\raggedright\ttfamily\scriptsize     scores = features @ weights.T + bias}}
\par\smallskip
This line handles the classifier parameters. Those arrays mirror the RTL memories, so mismatches reveal an export, quantization, or hardware-update bug.
\paragraph{Line 51}
\noindent\colorbox{gray!8}{\parbox{\dimexpr\linewidth-2\fboxsep\relax}{\raggedright\ttfamily\scriptsize     preds = scores.argmax(axis=1)}}
\par\smallskip
This line handles class scores. The score vector is the direct software equivalent of the RTL MAC loop over weights and features.
\paragraph{Line 52}
\noindent\colorbox{gray!8}{\parbox{\dimexpr\linewidth-2\fboxsep\relax}{\raggedright\ttfamily\scriptsize     return labels[:sample\_count], scores, preds, skip\_fraction}}
\par\smallskip
This line handles class scores. The score vector is the direct software equivalent of the RTL MAC loop over weights and features.
\paragraph{Line 53}
\noindent\colorbox{gray!8}{\parbox{\dimexpr\linewidth-2\fboxsep\relax}{\raggedright\ttfamily\scriptsize \emph{blank line}}}
\par\smallskip
This blank line is not hardware or software behavior. It separates nearby declarations or operations so the file remains readable and reviewable.
\paragraph{Line 54}
\noindent\colorbox{gray!8}{\parbox{\dimexpr\linewidth-2\fboxsep\relax}{\raggedright\ttfamily\scriptsize \emph{blank line}}}
\par\smallskip
This blank line is not hardware or software behavior. It separates nearby declarations or operations so the file remains readable and reviewable.
\paragraph{Line 55}
\noindent\colorbox{gray!8}{\parbox{\dimexpr\linewidth-2\fboxsep\relax}{\raggedright\ttfamily\scriptsize def write\_report(args: argparse.Namespace) -> dict[str, object]:}}
\par\smallskip
This defines a Python function. The function packages one repeatable step such as loading a memory file, simulating the golden learning rule, validating CSVs, or generating plots.
\paragraph{Line 56}
\noindent\colorbox{gray!8}{\parbox{\dimexpr\linewidth-2\fboxsep\relax}{\raggedright\ttfamily\scriptsize     rtl\_rows = load\_rows(args.results\_csv)}}
\par\smallskip
This Python line participates in the export, validation, or report-generation pipeline. It is executed on the host PC to prepare data for RTL or to check RTL outputs.
\paragraph{Line 57}
\noindent\colorbox{gray!8}{\parbox{\dimexpr\linewidth-2\fboxsep\relax}{\raggedright\ttfamily\scriptsize \emph{blank line}}}
\par\smallskip
This blank line is not hardware or software behavior. It separates nearby declarations or operations so the file remains readable and reviewable.
\paragraph{Line 58}
\noindent\colorbox{gray!8}{\parbox{\dimexpr\linewidth-2\fboxsep\relax}{\raggedright\ttfamily\scriptsize     labels = np.array([int(r["label"]) for r in rtl\_rows], dtype=np.int64)}}
\par\smallskip
This performs vectorized numerical work. The fixed-point golden models use NumPy so score, margin, and update calculations are explicit and repeatable.
\paragraph{Line 59}
\noindent\colorbox{gray!8}{\parbox{\dimexpr\linewidth-2\fboxsep\relax}{\raggedright\ttfamily\scriptsize     preds = np.array([int(r["predicted"]) for r in rtl\_rows], dtype=np.int64)}}
\par\smallskip
This performs vectorized numerical work. The fixed-point golden models use NumPy so score, margin, and update calculations are explicit and repeatable.
\paragraph{Line 60}
\noindent\colorbox{gray!8}{\parbox{\dimexpr\linewidth-2\fboxsep\relax}{\raggedright\ttfamily\scriptsize     correct = preds == labels}}
\par\smallskip
This line handles predicted labels. It checks that argmax behavior in Python and RTL chooses the same class.
\paragraph{Line 61}
\noindent\colorbox{gray!8}{\parbox{\dimexpr\linewidth-2\fboxsep\relax}{\raggedright\ttfamily\scriptsize     rtl\_logits = np.array([[int(r[f"score\{c\}"]) for c in range(NUM\_CLASSES)] for r in rtl\_rows], dtype=np.int64)}}
\par\smallskip
This performs vectorized numerical work. The fixed-point golden models use NumPy so score, margin, and update calculations are explicit and repeatable.
\paragraph{Line 62}
\noindent\colorbox{gray!8}{\parbox{\dimexpr\linewidth-2\fboxsep\relax}{\raggedright\ttfamily\scriptsize     labels\_ref, ref\_logits, pred\_ref, skip\_fraction = fixed\_reference\_from\_mem(args.rtl\_dir, len(rtl\_rows))}}
\par\smallskip
This line handles predicted labels. It checks that argmax behavior in Python and RTL chooses the same class.
\paragraph{Line 63}
\noindent\colorbox{gray!8}{\parbox{\dimexpr\linewidth-2\fboxsep\relax}{\raggedright\ttfamily\scriptsize     if not np.array\_equal(labels, labels\_ref):}}
\par\smallskip
This performs vectorized numerical work. The fixed-point golden models use NumPy so score, margin, and update calculations are explicit and repeatable.
\paragraph{Line 64}
\noindent\colorbox{gray!8}{\parbox{\dimexpr\linewidth-2\fboxsep\relax}{\raggedright\ttfamily\scriptsize         raise ValueError("RTL result labels do not match labels memory")}}
\par\smallskip
This Python line participates in the export, validation, or report-generation pipeline. It is executed on the host PC to prepare data for RTL or to check RTL outputs.
\paragraph{Line 65}
\noindent\colorbox{gray!8}{\parbox{\dimexpr\linewidth-2\fboxsep\relax}{\raggedright\ttfamily\scriptsize     error = rtl\_logits - ref\_logits}}
\par\smallskip
This Python line participates in the export, validation, or report-generation pipeline. It is executed on the host PC to prepare data for RTL or to check RTL outputs.
\paragraph{Line 66}
\noindent\colorbox{gray!8}{\parbox{\dimexpr\linewidth-2\fboxsep\relax}{\raggedright\ttfamily\scriptsize     mismatch\_count = int((error != 0).sum())}}
\par\smallskip
This Python line participates in the export, validation, or report-generation pipeline. It is executed on the host PC to prepare data for RTL or to check RTL outputs.
\paragraph{Line 67}
\noindent\colorbox{gray!8}{\parbox{\dimexpr\linewidth-2\fboxsep\relax}{\raggedright\ttfamily\scriptsize     pred\_mismatch = int((preds != pred\_ref).sum())}}
\par\smallskip
This line handles predicted labels. It checks that argmax behavior in Python and RTL chooses the same class.
\paragraph{Line 68}
\noindent\colorbox{gray!8}{\parbox{\dimexpr\linewidth-2\fboxsep\relax}{\raggedright\ttfamily\scriptsize \emph{blank line}}}
\par\smallskip
This blank line is not hardware or software behavior. It separates nearby declarations or operations so the file remains readable and reviewable.
\paragraph{Line 69}
\noindent\colorbox{gray!8}{\parbox{\dimexpr\linewidth-2\fboxsep\relax}{\raggedright\ttfamily\scriptsize     confusion = np.zeros((NUM\_CLASSES, NUM\_CLASSES), dtype=np.int64)}}
\par\smallskip
This performs vectorized numerical work. The fixed-point golden models use NumPy so score, margin, and update calculations are explicit and repeatable.
\paragraph{Line 70}
\noindent\colorbox{gray!8}{\parbox{\dimexpr\linewidth-2\fboxsep\relax}{\raggedright\ttfamily\scriptsize     for label, pred in zip(labels, preds):}}
\par\smallskip
This line handles predicted labels. It checks that argmax behavior in Python and RTL chooses the same class.
\paragraph{Line 71}
\noindent\colorbox{gray!8}{\parbox{\dimexpr\linewidth-2\fboxsep\relax}{\raggedright\ttfamily\scriptsize         confusion[label, pred] += 1}}
\par\smallskip
This line handles predicted labels. It checks that argmax behavior in Python and RTL chooses the same class.
\paragraph{Line 72}
\noindent\colorbox{gray!8}{\parbox{\dimexpr\linewidth-2\fboxsep\relax}{\raggedright\ttfamily\scriptsize \emph{blank line}}}
\par\smallskip
This blank line is not hardware or software behavior. It separates nearby declarations or operations so the file remains readable and reviewable.
\paragraph{Line 73}
\noindent\colorbox{gray!8}{\parbox{\dimexpr\linewidth-2\fboxsep\relax}{\raggedright\ttfamily\scriptsize     metrics = \{}}
\par\smallskip
This Python line participates in the export, validation, or report-generation pipeline. It is executed on the host PC to prepare data for RTL or to check RTL outputs.
\paragraph{Line 74}
\noindent\colorbox{gray!8}{\parbox{\dimexpr\linewidth-2\fboxsep\relax}{\raggedright\ttfamily\scriptsize         "samples": int(len(rtl\_rows)),}}
\par\smallskip
This Python line participates in the export, validation, or report-generation pipeline. It is executed on the host PC to prepare data for RTL or to check RTL outputs.
\paragraph{Line 75}
\noindent\colorbox{gray!8}{\parbox{\dimexpr\linewidth-2\fboxsep\relax}{\raggedright\ttfamily\scriptsize         "accuracy": float(correct.mean()) if len(rtl\_rows) else 0.0,}}
\par\smallskip
This Python line participates in the export, validation, or report-generation pipeline. It is executed on the host PC to prepare data for RTL or to check RTL outputs.
\paragraph{Line 76}
\noindent\colorbox{gray!8}{\parbox{\dimexpr\linewidth-2\fboxsep\relax}{\raggedright\ttfamily\scriptsize         "correct": int(correct.sum()),}}
\par\smallskip
This Python line participates in the export, validation, or report-generation pipeline. It is executed on the host PC to prepare data for RTL or to check RTL outputs.
\paragraph{Line 77}
\noindent\colorbox{gray!8}{\parbox{\dimexpr\linewidth-2\fboxsep\relax}{\raggedright\ttfamily\scriptsize         "score\_mismatch\_count": mismatch\_count,}}
\par\smallskip
This Python line participates in the export, validation, or report-generation pipeline. It is executed on the host PC to prepare data for RTL or to check RTL outputs.
\paragraph{Line 78}
\noindent\colorbox{gray!8}{\parbox{\dimexpr\linewidth-2\fboxsep\relax}{\raggedright\ttfamily\scriptsize         "score\_exact\_match": mismatch\_count == 0,}}
\par\smallskip
This Python line participates in the export, validation, or report-generation pipeline. It is executed on the host PC to prepare data for RTL or to check RTL outputs.
\paragraph{Line 79}
\noindent\colorbox{gray!8}{\parbox{\dimexpr\linewidth-2\fboxsep\relax}{\raggedright\ttfamily\scriptsize         "prediction\_mismatch\_count": pred\_mismatch,}}
\par\smallskip
This line handles predicted labels. It checks that argmax behavior in Python and RTL chooses the same class.
\paragraph{Line 80}
\noindent\colorbox{gray!8}{\parbox{\dimexpr\linewidth-2\fboxsep\relax}{\raggedright\ttfamily\scriptsize         "prediction\_exact\_match": pred\_mismatch == 0,}}
\par\smallskip
This line handles predicted labels. It checks that argmax behavior in Python and RTL chooses the same class.
\paragraph{Line 81}
\noindent\colorbox{gray!8}{\parbox{\dimexpr\linewidth-2\fboxsep\relax}{\raggedright\ttfamily\scriptsize         "max\_abs\_score\_error": int(np.abs(error).max()) if error.size else 0,}}
\par\smallskip
This performs vectorized numerical work. The fixed-point golden models use NumPy so score, margin, and update calculations are explicit and repeatable.
\paragraph{Line 82}
\noindent\colorbox{gray!8}{\parbox{\dimexpr\linewidth-2\fboxsep\relax}{\raggedright\ttfamily\scriptsize         "infer\_skip\_abs\_thresh": INFER\_SKIP\_ABS\_THRESH,}}
\par\smallskip
This Python line participates in the export, validation, or report-generation pipeline. It is executed on the host PC to prepare data for RTL or to check RTL outputs.
\paragraph{Line 83}
\noindent\colorbox{gray!8}{\parbox{\dimexpr\linewidth-2\fboxsep\relax}{\raggedright\ttfamily\scriptsize         "feature\_mac\_skip\_fraction": skip\_fraction,}}
\par\smallskip
This Python line participates in the export, validation, or report-generation pipeline. It is executed on the host PC to prepare data for RTL or to check RTL outputs.
\paragraph{Line 84}
\noindent\colorbox{gray!8}{\parbox{\dimexpr\linewidth-2\fboxsep\relax}{\raggedright\ttfamily\scriptsize         "results\_csv": str(args.results\_csv),}}
\par\smallskip
This Python line participates in the export, validation, or report-generation pipeline. It is executed on the host PC to prepare data for RTL or to check RTL outputs.
\paragraph{Line 85}
\noindent\colorbox{gray!8}{\parbox{\dimexpr\linewidth-2\fboxsep\relax}{\raggedright\ttfamily\scriptsize         "reference": "recomputed from current RTL .mem files",}}
\par\smallskip
This Python line participates in the export, validation, or report-generation pipeline. It is executed on the host PC to prepare data for RTL or to check RTL outputs.
\paragraph{Line 86}
\noindent\colorbox{gray!8}{\parbox{\dimexpr\linewidth-2\fboxsep\relax}{\raggedright\ttfamily\scriptsize     \}}}
\par\smallskip
This Python line participates in the export, validation, or report-generation pipeline. It is executed on the host PC to prepare data for RTL or to check RTL outputs.
\paragraph{Line 87}
\noindent\colorbox{gray!8}{\parbox{\dimexpr\linewidth-2\fboxsep\relax}{\raggedright\ttfamily\scriptsize     args.out\_dir.mkdir(parents=True, exist\_ok=True)}}
\par\smallskip
This Python line participates in the export, validation, or report-generation pipeline. It is executed on the host PC to prepare data for RTL or to check RTL outputs.
\paragraph{Line 88}
\noindent\colorbox{gray!8}{\parbox{\dimexpr\linewidth-2\fboxsep\relax}{\raggedright\ttfamily\scriptsize     (args.out\_dir / "spikingformer\_head\_rtl\_validation.json").write\_text(json.dumps(metrics, indent=2), encoding="utf-8")}}
\par\smallskip
This writes machine-readable output for later reports. Keeping JSON/CSV artifacts makes the verification auditable.
\paragraph{Line 89}
\noindent\colorbox{gray!8}{\parbox{\dimexpr\linewidth-2\fboxsep\relax}{\raggedright\ttfamily\scriptsize \emph{blank line}}}
\par\smallskip
This blank line is not hardware or software behavior. It separates nearby declarations or operations so the file remains readable and reviewable.
\paragraph{Line 90}
\noindent\colorbox{gray!8}{\parbox{\dimexpr\linewidth-2\fboxsep\relax}{\raggedright\ttfamily\scriptsize     plt.figure(figsize=(10, 5), dpi=160)}}
\par\smallskip
This creates a debug visualization. The plot is derived from generated CSV or Vivado report data so it can be regenerated after RTL changes.
\paragraph{Line 91}
\noindent\colorbox{gray!8}{\parbox{\dimexpr\linewidth-2\fboxsep\relax}{\raggedright\ttfamily\scriptsize     plt.imshow(error, aspect="auto", cmap="coolwarm")}}
\par\smallskip
This creates a debug visualization. The plot is derived from generated CSV or Vivado report data so it can be regenerated after RTL changes.
\paragraph{Line 92}
\noindent\colorbox{gray!8}{\parbox{\dimexpr\linewidth-2\fboxsep\relax}{\raggedright\ttfamily\scriptsize     plt.xlabel("class")}}
\par\smallskip
This creates a debug visualization. The plot is derived from generated CSV or Vivado report data so it can be regenerated after RTL changes.
\paragraph{Line 93}
\noindent\colorbox{gray!8}{\parbox{\dimexpr\linewidth-2\fboxsep\relax}{\raggedright\ttfamily\scriptsize     plt.ylabel("sample")}}
\par\smallskip
This creates a debug visualization. The plot is derived from generated CSV or Vivado report data so it can be regenerated after RTL changes.
\paragraph{Line 94}
\noindent\colorbox{gray!8}{\parbox{\dimexpr\linewidth-2\fboxsep\relax}{\raggedright\ttfamily\scriptsize     plt.colorbar(label="RTL score - fixed reference")}}
\par\smallskip
This creates a debug visualization. The plot is derived from generated CSV or Vivado report data so it can be regenerated after RTL changes.
\paragraph{Line 95}
\noindent\colorbox{gray!8}{\parbox{\dimexpr\linewidth-2\fboxsep\relax}{\raggedright\ttfamily\scriptsize     plt.tight\_layout()}}
\par\smallskip
This creates a debug visualization. The plot is derived from generated CSV or Vivado report data so it can be regenerated after RTL changes.
\paragraph{Line 96}
\noindent\colorbox{gray!8}{\parbox{\dimexpr\linewidth-2\fboxsep\relax}{\raggedright\ttfamily\scriptsize     plt.savefig(args.out\_dir / "spikingformer\_head\_rtl\_score\_error.png")}}
\par\smallskip
This creates a debug visualization. The plot is derived from generated CSV or Vivado report data so it can be regenerated after RTL changes.
\paragraph{Line 97}
\noindent\colorbox{gray!8}{\parbox{\dimexpr\linewidth-2\fboxsep\relax}{\raggedright\ttfamily\scriptsize     plt.close()}}
\par\smallskip
This creates a debug visualization. The plot is derived from generated CSV or Vivado report data so it can be regenerated after RTL changes.
\paragraph{Line 98}
\noindent\colorbox{gray!8}{\parbox{\dimexpr\linewidth-2\fboxsep\relax}{\raggedright\ttfamily\scriptsize \emph{blank line}}}
\par\smallskip
This blank line is not hardware or software behavior. It separates nearby declarations or operations so the file remains readable and reviewable.
\paragraph{Line 99}
\noindent\colorbox{gray!8}{\parbox{\dimexpr\linewidth-2\fboxsep\relax}{\raggedright\ttfamily\scriptsize     norm\_confusion = confusion / np.maximum(confusion.sum(axis=1, keepdims=True), 1)}}
\par\smallskip
This performs vectorized numerical work. The fixed-point golden models use NumPy so score, margin, and update calculations are explicit and repeatable.
\paragraph{Line 100}
\noindent\colorbox{gray!8}{\parbox{\dimexpr\linewidth-2\fboxsep\relax}{\raggedright\ttfamily\scriptsize     plt.figure(figsize=(7, 6), dpi=160)}}
\par\smallskip
This creates a debug visualization. The plot is derived from generated CSV or Vivado report data so it can be regenerated after RTL changes.
\paragraph{Line 101}
\noindent\colorbox{gray!8}{\parbox{\dimexpr\linewidth-2\fboxsep\relax}{\raggedright\ttfamily\scriptsize     plt.imshow(norm\_confusion, cmap="viridis", vmin=0, vmax=1)}}
\par\smallskip
This creates a debug visualization. The plot is derived from generated CSV or Vivado report data so it can be regenerated after RTL changes.
\paragraph{Line 102}
\noindent\colorbox{gray!8}{\parbox{\dimexpr\linewidth-2\fboxsep\relax}{\raggedright\ttfamily\scriptsize     plt.xlabel("predicted")}}
\par\smallskip
This line handles predicted labels. It checks that argmax behavior in Python and RTL chooses the same class.
\paragraph{Line 103}
\noindent\colorbox{gray!8}{\parbox{\dimexpr\linewidth-2\fboxsep\relax}{\raggedright\ttfamily\scriptsize     plt.ylabel("label")}}
\par\smallskip
This creates a debug visualization. The plot is derived from generated CSV or Vivado report data so it can be regenerated after RTL changes.
\paragraph{Line 104}
\noindent\colorbox{gray!8}{\parbox{\dimexpr\linewidth-2\fboxsep\relax}{\raggedright\ttfamily\scriptsize     plt.colorbar(label="row-normalized count")}}
\par\smallskip
This creates a debug visualization. The plot is derived from generated CSV or Vivado report data so it can be regenerated after RTL changes.
\paragraph{Line 105}
\noindent\colorbox{gray!8}{\parbox{\dimexpr\linewidth-2\fboxsep\relax}{\raggedright\ttfamily\scriptsize     plt.tight\_layout()}}
\par\smallskip
This creates a debug visualization. The plot is derived from generated CSV or Vivado report data so it can be regenerated after RTL changes.
\paragraph{Line 106}
\noindent\colorbox{gray!8}{\parbox{\dimexpr\linewidth-2\fboxsep\relax}{\raggedright\ttfamily\scriptsize     plt.savefig(args.out\_dir / "spikingformer\_head\_rtl\_confusion.png")}}
\par\smallskip
This creates a debug visualization. The plot is derived from generated CSV or Vivado report data so it can be regenerated after RTL changes.
\paragraph{Line 107}
\noindent\colorbox{gray!8}{\parbox{\dimexpr\linewidth-2\fboxsep\relax}{\raggedright\ttfamily\scriptsize     plt.close()}}
\par\smallskip
This creates a debug visualization. The plot is derived from generated CSV or Vivado report data so it can be regenerated after RTL changes.
\paragraph{Line 108}
\noindent\colorbox{gray!8}{\parbox{\dimexpr\linewidth-2\fboxsep\relax}{\raggedright\ttfamily\scriptsize \emph{blank line}}}
\par\smallskip
This blank line is not hardware or software behavior. It separates nearby declarations or operations so the file remains readable and reviewable.
\paragraph{Line 109}
\noindent\colorbox{gray!8}{\parbox{\dimexpr\linewidth-2\fboxsep\relax}{\raggedright\ttfamily\scriptsize     with PdfPages(args.out\_dir / "spikingformer\_head\_rtl\_validation.pdf") as pdf:}}
\par\smallskip
This creates a debug visualization. The plot is derived from generated CSV or Vivado report data so it can be regenerated after RTL changes.
\paragraph{Line 110}
\noindent\colorbox{gray!8}{\parbox{\dimexpr\linewidth-2\fboxsep\relax}{\raggedright\ttfamily\scriptsize         fig = plt.figure(figsize=(10, 7), dpi=160)}}
\par\smallskip
This creates a debug visualization. The plot is derived from generated CSV or Vivado report data so it can be regenerated after RTL changes.
\paragraph{Line 111}
\noindent\colorbox{gray!8}{\parbox{\dimexpr\linewidth-2\fboxsep\relax}{\raggedright\ttfamily\scriptsize         fig.text(0.06, 0.92, "Spikingformer Head RTL Validation", fontsize=18, weight="bold")}}
\par\smallskip
This Python line participates in the export, validation, or report-generation pipeline. It is executed on the host PC to prepare data for RTL or to check RTL outputs.
\paragraph{Line 112}
\noindent\colorbox{gray!8}{\parbox{\dimexpr\linewidth-2\fboxsep\relax}{\raggedright\ttfamily\scriptsize         lines = [}}
\par\smallskip
This Python line participates in the export, validation, or report-generation pipeline. It is executed on the host PC to prepare data for RTL or to check RTL outputs.
\paragraph{Line 113}
\noindent\colorbox{gray!8}{\parbox{\dimexpr\linewidth-2\fboxsep\relax}{\raggedright\ttfamily\scriptsize             f"Samples: \{metrics['samples']\}",}}
\par\smallskip
This Python line participates in the export, validation, or report-generation pipeline. It is executed on the host PC to prepare data for RTL or to check RTL outputs.
\paragraph{Line 114}
\noindent\colorbox{gray!8}{\parbox{\dimexpr\linewidth-2\fboxsep\relax}{\raggedright\ttfamily\scriptsize             f"RTL top-1: \{metrics['accuracy'] * 100:.2f\}\%",}}
\par\smallskip
This Python line participates in the export, validation, or report-generation pipeline. It is executed on the host PC to prepare data for RTL or to check RTL outputs.
\paragraph{Line 115}
\noindent\colorbox{gray!8}{\parbox{\dimexpr\linewidth-2\fboxsep\relax}{\raggedright\ttfamily\scriptsize             f"Correct: \{metrics['correct']\} / \{metrics['samples']\}",}}
\par\smallskip
This Python line participates in the export, validation, or report-generation pipeline. It is executed on the host PC to prepare data for RTL or to check RTL outputs.
\paragraph{Line 116}
\noindent\colorbox{gray!8}{\parbox{\dimexpr\linewidth-2\fboxsep\relax}{\raggedright\ttfamily\scriptsize             f"Score exact match: \{metrics['score\_exact\_match']\}",}}
\par\smallskip
This Python line participates in the export, validation, or report-generation pipeline. It is executed on the host PC to prepare data for RTL or to check RTL outputs.
\paragraph{Line 117}
\noindent\colorbox{gray!8}{\parbox{\dimexpr\linewidth-2\fboxsep\relax}{\raggedright\ttfamily\scriptsize             f"Score mismatches: \{metrics['score\_mismatch\_count']\}",}}
\par\smallskip
This Python line participates in the export, validation, or report-generation pipeline. It is executed on the host PC to prepare data for RTL or to check RTL outputs.
\paragraph{Line 118}
\noindent\colorbox{gray!8}{\parbox{\dimexpr\linewidth-2\fboxsep\relax}{\raggedright\ttfamily\scriptsize             f"Prediction exact match: \{metrics['prediction\_exact\_match']\}",}}
\par\smallskip
This line handles predicted labels. It checks that argmax behavior in Python and RTL chooses the same class.
\paragraph{Line 119}
\noindent\colorbox{gray!8}{\parbox{\dimexpr\linewidth-2\fboxsep\relax}{\raggedright\ttfamily\scriptsize             f"Prediction mismatches: \{metrics['prediction\_mismatch\_count']\}",}}
\par\smallskip
This line handles predicted labels. It checks that argmax behavior in Python and RTL chooses the same class.
\paragraph{Line 120}
\noindent\colorbox{gray!8}{\parbox{\dimexpr\linewidth-2\fboxsep\relax}{\raggedright\ttfamily\scriptsize             f"Max absolute score error: \{metrics['max\_abs\_score\_error']\}",}}
\par\smallskip
This Python line participates in the export, validation, or report-generation pipeline. It is executed on the host PC to prepare data for RTL or to check RTL outputs.
\paragraph{Line 121}
\noindent\colorbox{gray!8}{\parbox{\dimexpr\linewidth-2\fboxsep\relax}{\raggedright\ttfamily\scriptsize         ]}}
\par\smallskip
This Python line participates in the export, validation, or report-generation pipeline. It is executed on the host PC to prepare data for RTL or to check RTL outputs.
\paragraph{Line 122}
\noindent\colorbox{gray!8}{\parbox{\dimexpr\linewidth-2\fboxsep\relax}{\raggedright\ttfamily\scriptsize         fig.text(0.06, 0.82, "\textbackslash{}n".join(lines), fontsize=12, va="top")}}
\par\smallskip
This Python line participates in the export, validation, or report-generation pipeline. It is executed on the host PC to prepare data for RTL or to check RTL outputs.
\paragraph{Line 123}
\noindent\colorbox{gray!8}{\parbox{\dimexpr\linewidth-2\fboxsep\relax}{\raggedright\ttfamily\scriptsize         pdf.savefig(fig)}}
\par\smallskip
This Python line participates in the export, validation, or report-generation pipeline. It is executed on the host PC to prepare data for RTL or to check RTL outputs.
\paragraph{Line 124}
\noindent\colorbox{gray!8}{\parbox{\dimexpr\linewidth-2\fboxsep\relax}{\raggedright\ttfamily\scriptsize         plt.close(fig)}}
\par\smallskip
This creates a debug visualization. The plot is derived from generated CSV or Vivado report data so it can be regenerated after RTL changes.
\paragraph{Line 125}
\noindent\colorbox{gray!8}{\parbox{\dimexpr\linewidth-2\fboxsep\relax}{\raggedright\ttfamily\scriptsize         for name in ["spikingformer\_head\_rtl\_score\_error.png", "spikingformer\_head\_rtl\_confusion.png"]:}}
\par\smallskip
This Python line participates in the export, validation, or report-generation pipeline. It is executed on the host PC to prepare data for RTL or to check RTL outputs.
\paragraph{Line 126}
\noindent\colorbox{gray!8}{\parbox{\dimexpr\linewidth-2\fboxsep\relax}{\raggedright\ttfamily\scriptsize             img = plt.imread(args.out\_dir / name)}}
\par\smallskip
This creates a debug visualization. The plot is derived from generated CSV or Vivado report data so it can be regenerated after RTL changes.
\paragraph{Line 127}
\noindent\colorbox{gray!8}{\parbox{\dimexpr\linewidth-2\fboxsep\relax}{\raggedright\ttfamily\scriptsize             fig, ax = plt.subplots(figsize=(10, 7), dpi=160)}}
\par\smallskip
This creates a debug visualization. The plot is derived from generated CSV or Vivado report data so it can be regenerated after RTL changes.
\paragraph{Line 128}
\noindent\colorbox{gray!8}{\parbox{\dimexpr\linewidth-2\fboxsep\relax}{\raggedright\ttfamily\scriptsize             ax.imshow(img)}}
\par\smallskip
This Python line participates in the export, validation, or report-generation pipeline. It is executed on the host PC to prepare data for RTL or to check RTL outputs.
\paragraph{Line 129}
\noindent\colorbox{gray!8}{\parbox{\dimexpr\linewidth-2\fboxsep\relax}{\raggedright\ttfamily\scriptsize             ax.axis("off")}}
\par\smallskip
This Python line participates in the export, validation, or report-generation pipeline. It is executed on the host PC to prepare data for RTL or to check RTL outputs.
\paragraph{Line 130}
\noindent\colorbox{gray!8}{\parbox{\dimexpr\linewidth-2\fboxsep\relax}{\raggedright\ttfamily\scriptsize             pdf.savefig(fig)}}
\par\smallskip
This Python line participates in the export, validation, or report-generation pipeline. It is executed on the host PC to prepare data for RTL or to check RTL outputs.
\paragraph{Line 131}
\noindent\colorbox{gray!8}{\parbox{\dimexpr\linewidth-2\fboxsep\relax}{\raggedright\ttfamily\scriptsize             plt.close(fig)}}
\par\smallskip
This creates a debug visualization. The plot is derived from generated CSV or Vivado report data so it can be regenerated after RTL changes.
\paragraph{Line 132}
\noindent\colorbox{gray!8}{\parbox{\dimexpr\linewidth-2\fboxsep\relax}{\raggedright\ttfamily\scriptsize \emph{blank line}}}
\par\smallskip
This blank line is not hardware or software behavior. It separates nearby declarations or operations so the file remains readable and reviewable.
\paragraph{Line 133}
\noindent\colorbox{gray!8}{\parbox{\dimexpr\linewidth-2\fboxsep\relax}{\raggedright\ttfamily\scriptsize     return metrics}}
\par\smallskip
This Python line participates in the export, validation, or report-generation pipeline. It is executed on the host PC to prepare data for RTL or to check RTL outputs.
\paragraph{Line 134}
\noindent\colorbox{gray!8}{\parbox{\dimexpr\linewidth-2\fboxsep\relax}{\raggedright\ttfamily\scriptsize \emph{blank line}}}
\par\smallskip
This blank line is not hardware or software behavior. It separates nearby declarations or operations so the file remains readable and reviewable.
\paragraph{Line 135}
\noindent\colorbox{gray!8}{\parbox{\dimexpr\linewidth-2\fboxsep\relax}{\raggedright\ttfamily\scriptsize \emph{blank line}}}
\par\smallskip
This blank line is not hardware or software behavior. It separates nearby declarations or operations so the file remains readable and reviewable.
\paragraph{Line 136}
\noindent\colorbox{gray!8}{\parbox{\dimexpr\linewidth-2\fboxsep\relax}{\raggedright\ttfamily\scriptsize def parse\_args() -> argparse.Namespace:}}
\par\smallskip
This defines a Python function. The function packages one repeatable step such as loading a memory file, simulating the golden learning rule, validating CSVs, or generating plots.
\paragraph{Line 137}
\noindent\colorbox{gray!8}{\parbox{\dimexpr\linewidth-2\fboxsep\relax}{\raggedright\ttfamily\scriptsize     root = repo\_root()}}
\par\smallskip
This Python line participates in the export, validation, or report-generation pipeline. It is executed on the host PC to prepare data for RTL or to check RTL outputs.
\paragraph{Line 138}
\noindent\colorbox{gray!8}{\parbox{\dimexpr\linewidth-2\fboxsep\relax}{\raggedright\ttfamily\scriptsize     parser = argparse.ArgumentParser(description="Validate Spikingformer head RTL against exported fixed-point reference.")}}
\par\smallskip
This contributes to the command-line interface. It lets the rerun flow choose paths, sample counts, or output locations without editing code.
\paragraph{Line 139}
\noindent\colorbox{gray!8}{\parbox{\dimexpr\linewidth-2\fboxsep\relax}{\raggedright\ttfamily\scriptsize     parser.add\_argument("--results-csv", type=Path, default=root / "cifar10\_neuromorphic" / "cifar10\_spikingformer\_head\_results.csv")}}
\par\smallskip
This contributes to the command-line interface. It lets the rerun flow choose paths, sample counts, or output locations without editing code.
\paragraph{Line 140}
\noindent\colorbox{gray!8}{\parbox{\dimexpr\linewidth-2\fboxsep\relax}{\raggedright\ttfamily\scriptsize     parser.add\_argument("--rtl-dir", type=Path, default=root / "cifar10\_neuromorphic" / "rtl")}}
\par\smallskip
This contributes to the command-line interface. It lets the rerun flow choose paths, sample counts, or output locations without editing code.
\paragraph{Line 141}
\noindent\colorbox{gray!8}{\parbox{\dimexpr\linewidth-2\fboxsep\relax}{\raggedright\ttfamily\scriptsize     parser.add\_argument("--out-dir", type=Path, default=root / "cifar10\_neuromorphic" / "outputs\_spikingformer\_head\_accelerator")}}
\par\smallskip
This contributes to the command-line interface. It lets the rerun flow choose paths, sample counts, or output locations without editing code.
\paragraph{Line 142}
\noindent\colorbox{gray!8}{\parbox{\dimexpr\linewidth-2\fboxsep\relax}{\raggedright\ttfamily\scriptsize     return parser.parse\_args()}}
\par\smallskip
This Python line participates in the export, validation, or report-generation pipeline. It is executed on the host PC to prepare data for RTL or to check RTL outputs.
\paragraph{Line 143}
\noindent\colorbox{gray!8}{\parbox{\dimexpr\linewidth-2\fboxsep\relax}{\raggedright\ttfamily\scriptsize \emph{blank line}}}
\par\smallskip
This blank line is not hardware or software behavior. It separates nearby declarations or operations so the file remains readable and reviewable.
\paragraph{Line 144}
\noindent\colorbox{gray!8}{\parbox{\dimexpr\linewidth-2\fboxsep\relax}{\raggedright\ttfamily\scriptsize \emph{blank line}}}
\par\smallskip
This blank line is not hardware or software behavior. It separates nearby declarations or operations so the file remains readable and reviewable.
\paragraph{Line 145}
\noindent\colorbox{gray!8}{\parbox{\dimexpr\linewidth-2\fboxsep\relax}{\raggedright\ttfamily\scriptsize def main() -> None:}}
\par\smallskip
This defines a Python function. The function packages one repeatable step such as loading a memory file, simulating the golden learning rule, validating CSVs, or generating plots.
\paragraph{Line 146}
\noindent\colorbox{gray!8}{\parbox{\dimexpr\linewidth-2\fboxsep\relax}{\raggedright\ttfamily\scriptsize     metrics = write\_report(parse\_args())}}
\par\smallskip
This Python line participates in the export, validation, or report-generation pipeline. It is executed on the host PC to prepare data for RTL or to check RTL outputs.
\paragraph{Line 147}
\noindent\colorbox{gray!8}{\parbox{\dimexpr\linewidth-2\fboxsep\relax}{\raggedright\ttfamily\scriptsize     print(json.dumps(metrics, indent=2))}}
\par\smallskip
This writes machine-readable output for later reports. Keeping JSON/CSV artifacts makes the verification auditable.
\paragraph{Line 148}
\noindent\colorbox{gray!8}{\parbox{\dimexpr\linewidth-2\fboxsep\relax}{\raggedright\ttfamily\scriptsize \emph{blank line}}}
\par\smallskip
This blank line is not hardware or software behavior. It separates nearby declarations or operations so the file remains readable and reviewable.
\paragraph{Line 149}
\noindent\colorbox{gray!8}{\parbox{\dimexpr\linewidth-2\fboxsep\relax}{\raggedright\ttfamily\scriptsize \emph{blank line}}}
\par\smallskip
This blank line is not hardware or software behavior. It separates nearby declarations or operations so the file remains readable and reviewable.
\paragraph{Line 150}
\noindent\colorbox{gray!8}{\parbox{\dimexpr\linewidth-2\fboxsep\relax}{\raggedright\ttfamily\scriptsize if \_\_name\_\_ == "\_\_main\_\_":}}
\par\smallskip
This Python idiom runs main() only when the file is executed as a script. It still allows the file's functions to be imported for tests or reuse.
\paragraph{Line 151}
\noindent\colorbox{gray!8}{\parbox{\dimexpr\linewidth-2\fboxsep\relax}{\raggedright\ttfamily\scriptsize     main()}}
\par\smallskip
This Python line participates in the export, validation, or report-generation pipeline. It is executed on the host PC to prepare data for RTL or to check RTL outputs.
\clearpage
\subsection{cifar10\_neuromorphic\textbackslash{}validate\_spikingformer\_head\_onchip\_learning.py}
Each entry below gives the exact source line followed by its hardware/software role in the implemented accelerator pipeline.
\paragraph{Line 1}
\noindent\colorbox{gray!8}{\parbox{\dimexpr\linewidth-2\fboxsep\relax}{\raggedright\ttfamily\scriptsize from \_\_future\_\_ import annotations}}
\par\smallskip
This imports a Python library or symbol. The scripts use standard parsing, numeric arrays, plotting, and file APIs to make the hardware artifacts reproducible.
\paragraph{Line 2}
\noindent\colorbox{gray!8}{\parbox{\dimexpr\linewidth-2\fboxsep\relax}{\raggedright\ttfamily\scriptsize \emph{blank line}}}
\par\smallskip
This blank line is not hardware or software behavior. It separates nearby declarations or operations so the file remains readable and reviewable.
\paragraph{Line 3}
\noindent\colorbox{gray!8}{\parbox{\dimexpr\linewidth-2\fboxsep\relax}{\raggedright\ttfamily\scriptsize import argparse}}
\par\smallskip
This imports a Python library or symbol. The scripts use standard parsing, numeric arrays, plotting, and file APIs to make the hardware artifacts reproducible.
\paragraph{Line 4}
\noindent\colorbox{gray!8}{\parbox{\dimexpr\linewidth-2\fboxsep\relax}{\raggedright\ttfamily\scriptsize import csv}}
\par\smallskip
This imports a Python library or symbol. The scripts use standard parsing, numeric arrays, plotting, and file APIs to make the hardware artifacts reproducible.
\paragraph{Line 5}
\noindent\colorbox{gray!8}{\parbox{\dimexpr\linewidth-2\fboxsep\relax}{\raggedright\ttfamily\scriptsize import json}}
\par\smallskip
This imports a Python library or symbol. The scripts use standard parsing, numeric arrays, plotting, and file APIs to make the hardware artifacts reproducible.
\paragraph{Line 6}
\noindent\colorbox{gray!8}{\parbox{\dimexpr\linewidth-2\fboxsep\relax}{\raggedright\ttfamily\scriptsize from pathlib import Path}}
\par\smallskip
This imports a Python library or symbol. The scripts use standard parsing, numeric arrays, plotting, and file APIs to make the hardware artifacts reproducible.
\paragraph{Line 7}
\noindent\colorbox{gray!8}{\parbox{\dimexpr\linewidth-2\fboxsep\relax}{\raggedright\ttfamily\scriptsize \emph{blank line}}}
\par\smallskip
This blank line is not hardware or software behavior. It separates nearby declarations or operations so the file remains readable and reviewable.
\paragraph{Line 8}
\noindent\colorbox{gray!8}{\parbox{\dimexpr\linewidth-2\fboxsep\relax}{\raggedright\ttfamily\scriptsize import numpy as np}}
\par\smallskip
This imports a Python library or symbol. The scripts use standard parsing, numeric arrays, plotting, and file APIs to make the hardware artifacts reproducible.
\paragraph{Line 9}
\noindent\colorbox{gray!8}{\parbox{\dimexpr\linewidth-2\fboxsep\relax}{\raggedright\ttfamily\scriptsize \emph{blank line}}}
\par\smallskip
This blank line is not hardware or software behavior. It separates nearby declarations or operations so the file remains readable and reviewable.
\paragraph{Line 10}
\noindent\colorbox{gray!8}{\parbox{\dimexpr\linewidth-2\fboxsep\relax}{\raggedright\ttfamily\scriptsize \emph{blank line}}}
\par\smallskip
This blank line is not hardware or software behavior. It separates nearby declarations or operations so the file remains readable and reviewable.
\paragraph{Line 11}
\noindent\colorbox{gray!8}{\parbox{\dimexpr\linewidth-2\fboxsep\relax}{\raggedright\ttfamily\scriptsize NUM\_CLASSES = 10}}
\par\smallskip
This Python line participates in the export, validation, or report-generation pipeline. It is executed on the host PC to prepare data for RTL or to check RTL outputs.
\paragraph{Line 12}
\noindent\colorbox{gray!8}{\parbox{\dimexpr\linewidth-2\fboxsep\relax}{\raggedright\ttfamily\scriptsize FEATURE\_DIM = 384}}
\par\smallskip
This Python line participates in the export, validation, or report-generation pipeline. It is executed on the host PC to prepare data for RTL or to check RTL outputs.
\paragraph{Line 13}
\noindent\colorbox{gray!8}{\parbox{\dimexpr\linewidth-2\fboxsep\relax}{\raggedright\ttfamily\scriptsize LEARN\_SHIFT = 14}}
\par\smallskip
This Python line participates in the export, validation, or report-generation pipeline. It is executed on the host PC to prepare data for RTL or to check RTL outputs.
\paragraph{Line 14}
\noindent\colorbox{gray!8}{\parbox{\dimexpr\linewidth-2\fboxsep\relax}{\raggedright\ttfamily\scriptsize BIAS\_LEARN\_DELTA = 1}}
\par\smallskip
This Python line participates in the export, validation, or report-generation pipeline. It is executed on the host PC to prepare data for RTL or to check RTL outputs.
\paragraph{Line 15}
\noindent\colorbox{gray!8}{\parbox{\dimexpr\linewidth-2\fboxsep\relax}{\raggedright\ttfamily\scriptsize LEARN\_MARGIN\_Q = 268435456}}
\par\smallskip
This Python line participates in the export, validation, or report-generation pipeline. It is executed on the host PC to prepare data for RTL or to check RTL outputs.
\paragraph{Line 16}
\noindent\colorbox{gray!8}{\parbox{\dimexpr\linewidth-2\fboxsep\relax}{\raggedright\ttfamily\scriptsize INFER\_SKIP\_ABS\_THRESH = 2048}}
\par\smallskip
This Python line participates in the export, validation, or report-generation pipeline. It is executed on the host PC to prepare data for RTL or to check RTL outputs.
\paragraph{Line 17}
\noindent\colorbox{gray!8}{\parbox{\dimexpr\linewidth-2\fboxsep\relax}{\raggedright\ttfamily\scriptsize \emph{blank line}}}
\par\smallskip
This blank line is not hardware or software behavior. It separates nearby declarations or operations so the file remains readable and reviewable.
\paragraph{Line 18}
\noindent\colorbox{gray!8}{\parbox{\dimexpr\linewidth-2\fboxsep\relax}{\raggedright\ttfamily\scriptsize \emph{blank line}}}
\par\smallskip
This blank line is not hardware or software behavior. It separates nearby declarations or operations so the file remains readable and reviewable.
\paragraph{Line 19}
\noindent\colorbox{gray!8}{\parbox{\dimexpr\linewidth-2\fboxsep\relax}{\raggedright\ttfamily\scriptsize def repo\_root() -> Path:}}
\par\smallskip
This defines a Python function. The function packages one repeatable step such as loading a memory file, simulating the golden learning rule, validating CSVs, or generating plots.
\paragraph{Line 20}
\noindent\colorbox{gray!8}{\parbox{\dimexpr\linewidth-2\fboxsep\relax}{\raggedright\ttfamily\scriptsize     return Path(\_\_file\_\_).resolve().parents[1]}}
\par\smallskip
This uses path objects so generated artifacts are addressed consistently across Windows shells and Python code.
\paragraph{Line 21}
\noindent\colorbox{gray!8}{\parbox{\dimexpr\linewidth-2\fboxsep\relax}{\raggedright\ttfamily\scriptsize \emph{blank line}}}
\par\smallskip
This blank line is not hardware or software behavior. It separates nearby declarations or operations so the file remains readable and reviewable.
\paragraph{Line 22}
\noindent\colorbox{gray!8}{\parbox{\dimexpr\linewidth-2\fboxsep\relax}{\raggedright\ttfamily\scriptsize \emph{blank line}}}
\par\smallskip
This blank line is not hardware or software behavior. It separates nearby declarations or operations so the file remains readable and reviewable.
\paragraph{Line 23}
\noindent\colorbox{gray!8}{\parbox{\dimexpr\linewidth-2\fboxsep\relax}{\raggedright\ttfamily\scriptsize def signed\_hex\_to\_int(token: str, bits: int) -> int:}}
\par\smallskip
This defines a Python function. The function packages one repeatable step such as loading a memory file, simulating the golden learning rule, validating CSVs, or generating plots.
\paragraph{Line 24}
\noindent\colorbox{gray!8}{\parbox{\dimexpr\linewidth-2\fboxsep\relax}{\raggedright\ttfamily\scriptsize     value = int(token.strip(), 16)}}
\par\smallskip
This Python line participates in the export, validation, or report-generation pipeline. It is executed on the host PC to prepare data for RTL or to check RTL outputs.
\paragraph{Line 25}
\noindent\colorbox{gray!8}{\parbox{\dimexpr\linewidth-2\fboxsep\relax}{\raggedright\ttfamily\scriptsize     sign = 1 << (bits - 1)}}
\par\smallskip
This Python line participates in the export, validation, or report-generation pipeline. It is executed on the host PC to prepare data for RTL or to check RTL outputs.
\paragraph{Line 26}
\noindent\colorbox{gray!8}{\parbox{\dimexpr\linewidth-2\fboxsep\relax}{\raggedright\ttfamily\scriptsize     return value - (1 << bits) if value \& sign else value}}
\par\smallskip
This Python line participates in the export, validation, or report-generation pipeline. It is executed on the host PC to prepare data for RTL or to check RTL outputs.
\paragraph{Line 27}
\noindent\colorbox{gray!8}{\parbox{\dimexpr\linewidth-2\fboxsep\relax}{\raggedright\ttfamily\scriptsize \emph{blank line}}}
\par\smallskip
This blank line is not hardware or software behavior. It separates nearby declarations or operations so the file remains readable and reviewable.
\paragraph{Line 28}
\noindent\colorbox{gray!8}{\parbox{\dimexpr\linewidth-2\fboxsep\relax}{\raggedright\ttfamily\scriptsize \emph{blank line}}}
\par\smallskip
This blank line is not hardware or software behavior. It separates nearby declarations or operations so the file remains readable and reviewable.
\paragraph{Line 29}
\noindent\colorbox{gray!8}{\parbox{\dimexpr\linewidth-2\fboxsep\relax}{\raggedright\ttfamily\scriptsize def load\_mem(path: Path, bits: int) -> list[int]:}}
\par\smallskip
This defines a Python function. The function packages one repeatable step such as loading a memory file, simulating the golden learning rule, validating CSVs, or generating plots.
\paragraph{Line 30}
\noindent\colorbox{gray!8}{\parbox{\dimexpr\linewidth-2\fboxsep\relax}{\raggedright\ttfamily\scriptsize     return [signed\_hex\_to\_int(line, bits) for line in path.read\_text(encoding="ascii").splitlines() if line.strip()]}}
\par\smallskip
This reads an artifact produced by training, export, simulation, or Vivado. The validator/report depends on files generated by the actual run rather than hidden constants.
\paragraph{Line 31}
\noindent\colorbox{gray!8}{\parbox{\dimexpr\linewidth-2\fboxsep\relax}{\raggedright\ttfamily\scriptsize \emph{blank line}}}
\par\smallskip
This blank line is not hardware or software behavior. It separates nearby declarations or operations so the file remains readable and reviewable.
\paragraph{Line 32}
\noindent\colorbox{gray!8}{\parbox{\dimexpr\linewidth-2\fboxsep\relax}{\raggedright\ttfamily\scriptsize \emph{blank line}}}
\par\smallskip
This blank line is not hardware or software behavior. It separates nearby declarations or operations so the file remains readable and reviewable.
\paragraph{Line 33}
\noindent\colorbox{gray!8}{\parbox{\dimexpr\linewidth-2\fboxsep\relax}{\raggedright\ttfamily\scriptsize def clamp\_i16(value: int) -> int:}}
\par\smallskip
This defines a Python function. The function packages one repeatable step such as loading a memory file, simulating the golden learning rule, validating CSVs, or generating plots.
\paragraph{Line 34}
\noindent\colorbox{gray!8}{\parbox{\dimexpr\linewidth-2\fboxsep\relax}{\raggedright\ttfamily\scriptsize     return max(-32768, min(32767, value))}}
\par\smallskip
This Python line participates in the export, validation, or report-generation pipeline. It is executed on the host PC to prepare data for RTL or to check RTL outputs.
\paragraph{Line 35}
\noindent\colorbox{gray!8}{\parbox{\dimexpr\linewidth-2\fboxsep\relax}{\raggedright\ttfamily\scriptsize \emph{blank line}}}
\par\smallskip
This blank line is not hardware or software behavior. It separates nearby declarations or operations so the file remains readable and reviewable.
\paragraph{Line 36}
\noindent\colorbox{gray!8}{\parbox{\dimexpr\linewidth-2\fboxsep\relax}{\raggedright\ttfamily\scriptsize \emph{blank line}}}
\par\smallskip
This blank line is not hardware or software behavior. It separates nearby declarations or operations so the file remains readable and reviewable.
\paragraph{Line 37}
\noindent\colorbox{gray!8}{\parbox{\dimexpr\linewidth-2\fboxsep\relax}{\raggedright\ttfamily\scriptsize def effective\_features(features: np.ndarray) -> np.ndarray:}}
\par\smallskip
This defines a Python function. The function packages one repeatable step such as loading a memory file, simulating the golden learning rule, validating CSVs, or generating plots.
\paragraph{Line 38}
\noindent\colorbox{gray!8}{\parbox{\dimexpr\linewidth-2\fboxsep\relax}{\raggedright\ttfamily\scriptsize     if INFER\_SKIP\_ABS\_THRESH <= 0:}}
\par\smallskip
This Python line participates in the export, validation, or report-generation pipeline. It is executed on the host PC to prepare data for RTL or to check RTL outputs.
\paragraph{Line 39}
\noindent\colorbox{gray!8}{\parbox{\dimexpr\linewidth-2\fboxsep\relax}{\raggedright\ttfamily\scriptsize         return features}}
\par\smallskip
This Python line participates in the export, validation, or report-generation pipeline. It is executed on the host PC to prepare data for RTL or to check RTL outputs.
\paragraph{Line 40}
\noindent\colorbox{gray!8}{\parbox{\dimexpr\linewidth-2\fboxsep\relax}{\raggedright\ttfamily\scriptsize     out = features.copy()}}
\par\smallskip
This Python line participates in the export, validation, or report-generation pipeline. It is executed on the host PC to prepare data for RTL or to check RTL outputs.
\paragraph{Line 41}
\noindent\colorbox{gray!8}{\parbox{\dimexpr\linewidth-2\fboxsep\relax}{\raggedright\ttfamily\scriptsize     out[np.abs(out) <= INFER\_SKIP\_ABS\_THRESH] = 0}}
\par\smallskip
This performs vectorized numerical work. The fixed-point golden models use NumPy so score, margin, and update calculations are explicit and repeatable.
\paragraph{Line 42}
\noindent\colorbox{gray!8}{\parbox{\dimexpr\linewidth-2\fboxsep\relax}{\raggedright\ttfamily\scriptsize     return out}}
\par\smallskip
This Python line participates in the export, validation, or report-generation pipeline. It is executed on the host PC to prepare data for RTL or to check RTL outputs.
\paragraph{Line 43}
\noindent\colorbox{gray!8}{\parbox{\dimexpr\linewidth-2\fboxsep\relax}{\raggedright\ttfamily\scriptsize \emph{blank line}}}
\par\smallskip
This blank line is not hardware or software behavior. It separates nearby declarations or operations so the file remains readable and reviewable.
\paragraph{Line 44}
\noindent\colorbox{gray!8}{\parbox{\dimexpr\linewidth-2\fboxsep\relax}{\raggedright\ttfamily\scriptsize \emph{blank line}}}
\par\smallskip
This blank line is not hardware or software behavior. It separates nearby declarations or operations so the file remains readable and reviewable.
\paragraph{Line 45}
\noindent\colorbox{gray!8}{\parbox{\dimexpr\linewidth-2\fboxsep\relax}{\raggedright\ttfamily\scriptsize def load\_rows(path: Path) -> list[dict[str, str]]:}}
\par\smallskip
This defines a Python function. The function packages one repeatable step such as loading a memory file, simulating the golden learning rule, validating CSVs, or generating plots.
\paragraph{Line 46}
\noindent\colorbox{gray!8}{\parbox{\dimexpr\linewidth-2\fboxsep\relax}{\raggedright\ttfamily\scriptsize     with path.open("r", newline="", encoding="utf-8") as fh:}}
\par\smallskip
This reads an artifact produced by training, export, simulation, or Vivado. The validator/report depends on files generated by the actual run rather than hidden constants.
\paragraph{Line 47}
\noindent\colorbox{gray!8}{\parbox{\dimexpr\linewidth-2\fboxsep\relax}{\raggedright\ttfamily\scriptsize         return list(csv.DictReader(fh))}}
\par\smallskip
This Python line participates in the export, validation, or report-generation pipeline. It is executed on the host PC to prepare data for RTL or to check RTL outputs.
\paragraph{Line 48}
\noindent\colorbox{gray!8}{\parbox{\dimexpr\linewidth-2\fboxsep\relax}{\raggedright\ttfamily\scriptsize \emph{blank line}}}
\par\smallskip
This blank line is not hardware or software behavior. It separates nearby declarations or operations so the file remains readable and reviewable.
\paragraph{Line 49}
\noindent\colorbox{gray!8}{\parbox{\dimexpr\linewidth-2\fboxsep\relax}{\raggedright\ttfamily\scriptsize \emph{blank line}}}
\par\smallskip
This blank line is not hardware or software behavior. It separates nearby declarations or operations so the file remains readable and reviewable.
\paragraph{Line 50}
\noindent\colorbox{gray!8}{\parbox{\dimexpr\linewidth-2\fboxsep\relax}{\raggedright\ttfamily\scriptsize def simulate\_online\_learning(rtl\_dir: Path, sample\_count: int) -> tuple[np.ndarray, np.ndarray, np.ndarray, np.ndarray]:}}
\par\smallskip
This defines a Python function. The function packages one repeatable step such as loading a memory file, simulating the golden learning rule, validating CSVs, or generating plots.
\paragraph{Line 51}
\noindent\colorbox{gray!8}{\parbox{\dimexpr\linewidth-2\fboxsep\relax}{\raggedright\ttfamily\scriptsize     features\_flat = load\_mem(rtl\_dir / "cifar10\_spikingformer\_head\_features.mem", 16)}}
\par\smallskip
This Python line participates in the export, validation, or report-generation pipeline. It is executed on the host PC to prepare data for RTL or to check RTL outputs.
\paragraph{Line 52}
\noindent\colorbox{gray!8}{\parbox{\dimexpr\linewidth-2\fboxsep\relax}{\raggedright\ttfamily\scriptsize     labels = np.array([int(line.strip(), 16) for line in (rtl\_dir / "cifar10\_spikingformer\_head\_labels.mem").read\_text(encoding="ascii").splitlines() if line.strip()], dtype=np.int64)}}
\par\smallskip
This reads an artifact produced by training, export, simulation, or Vivado. The validator/report depends on files generated by the actual run rather than hidden constants.
\paragraph{Line 53}
\noindent\colorbox{gray!8}{\parbox{\dimexpr\linewidth-2\fboxsep\relax}{\raggedright\ttfamily\scriptsize     weights = np.array(load\_mem(rtl\_dir / "cifar10\_spikingformer\_head\_weights.mem", 16), dtype=np.int64).reshape(NUM\_CLASSES, FEATURE\_DIM)}}
\par\smallskip
This performs vectorized numerical work. The fixed-point golden models use NumPy so score, margin, and update calculations are explicit and repeatable.
\paragraph{Line 54}
\noindent\colorbox{gray!8}{\parbox{\dimexpr\linewidth-2\fboxsep\relax}{\raggedright\ttfamily\scriptsize     bias = np.array(load\_mem(rtl\_dir / "cifar10\_spikingformer\_head\_bias.mem", 64), dtype=np.int64)}}
\par\smallskip
This performs vectorized numerical work. The fixed-point golden models use NumPy so score, margin, and update calculations are explicit and repeatable.
\paragraph{Line 55}
\noindent\colorbox{gray!8}{\parbox{\dimexpr\linewidth-2\fboxsep\relax}{\raggedright\ttfamily\scriptsize     features = np.array(features\_flat[: sample\_count * FEATURE\_DIM], dtype=np.int64).reshape(sample\_count, FEATURE\_DIM)}}
\par\smallskip
This performs vectorized numerical work. The fixed-point golden models use NumPy so score, margin, and update calculations are explicit and repeatable.
\paragraph{Line 56}
\noindent\colorbox{gray!8}{\parbox{\dimexpr\linewidth-2\fboxsep\relax}{\raggedright\ttfamily\scriptsize \emph{blank line}}}
\par\smallskip
This blank line is not hardware or software behavior. It separates nearby declarations or operations so the file remains readable and reviewable.
\paragraph{Line 57}
\noindent\colorbox{gray!8}{\parbox{\dimexpr\linewidth-2\fboxsep\relax}{\raggedright\ttfamily\scriptsize     scores\_out = np.zeros((sample\_count, NUM\_CLASSES), dtype=np.int64)}}
\par\smallskip
This performs vectorized numerical work. The fixed-point golden models use NumPy so score, margin, and update calculations are explicit and repeatable.
\paragraph{Line 58}
\noindent\colorbox{gray!8}{\parbox{\dimexpr\linewidth-2\fboxsep\relax}{\raggedright\ttfamily\scriptsize     preds = np.zeros(sample\_count, dtype=np.int64)}}
\par\smallskip
This performs vectorized numerical work. The fixed-point golden models use NumPy so score, margin, and update calculations are explicit and repeatable.
\paragraph{Line 59}
\noindent\colorbox{gray!8}{\parbox{\dimexpr\linewidth-2\fboxsep\relax}{\raggedright\ttfamily\scriptsize     update\_applied = np.zeros(sample\_count, dtype=np.int64)}}
\par\smallskip
This performs vectorized numerical work. The fixed-point golden models use NumPy so score, margin, and update calculations are explicit and repeatable.
\paragraph{Line 60}
\noindent\colorbox{gray!8}{\parbox{\dimexpr\linewidth-2\fboxsep\relax}{\raggedright\ttfamily\scriptsize     update\_count = np.zeros(sample\_count, dtype=np.int64)}}
\par\smallskip
This performs vectorized numerical work. The fixed-point golden models use NumPy so score, margin, and update calculations are explicit and repeatable.
\paragraph{Line 61}
\noindent\colorbox{gray!8}{\parbox{\dimexpr\linewidth-2\fboxsep\relax}{\raggedright\ttfamily\scriptsize     count = 0}}
\par\smallskip
This Python line participates in the export, validation, or report-generation pipeline. It is executed on the host PC to prepare data for RTL or to check RTL outputs.
\paragraph{Line 62}
\noindent\colorbox{gray!8}{\parbox{\dimexpr\linewidth-2\fboxsep\relax}{\raggedright\ttfamily\scriptsize \emph{blank line}}}
\par\smallskip
This blank line is not hardware or software behavior. It separates nearby declarations or operations so the file remains readable and reviewable.
\paragraph{Line 63}
\noindent\colorbox{gray!8}{\parbox{\dimexpr\linewidth-2\fboxsep\relax}{\raggedright\ttfamily\scriptsize     for sample\_idx in range(sample\_count):}}
\par\smallskip
This Python line participates in the export, validation, or report-generation pipeline. It is executed on the host PC to prepare data for RTL or to check RTL outputs.
\paragraph{Line 64}
\noindent\colorbox{gray!8}{\parbox{\dimexpr\linewidth-2\fboxsep\relax}{\raggedright\ttfamily\scriptsize         x = effective\_features(features[sample\_idx])}}
\par\smallskip
This Python line participates in the export, validation, or report-generation pipeline. It is executed on the host PC to prepare data for RTL or to check RTL outputs.
\paragraph{Line 65}
\noindent\colorbox{gray!8}{\parbox{\dimexpr\linewidth-2\fboxsep\relax}{\raggedright\ttfamily\scriptsize         scores = bias + weights @ x}}
\par\smallskip
This line handles the classifier parameters. Those arrays mirror the RTL memories, so mismatches reveal an export, quantization, or hardware-update bug.
\paragraph{Line 66}
\noindent\colorbox{gray!8}{\parbox{\dimexpr\linewidth-2\fboxsep\relax}{\raggedright\ttfamily\scriptsize         pred = int(scores.argmax())}}
\par\smallskip
This line handles class scores. The score vector is the direct software equivalent of the RTL MAC loop over weights and features.
\paragraph{Line 67}
\noindent\colorbox{gray!8}{\parbox{\dimexpr\linewidth-2\fboxsep\relax}{\raggedright\ttfamily\scriptsize         label = int(labels[sample\_idx])}}
\par\smallskip
This Python line participates in the export, validation, or report-generation pipeline. It is executed on the host PC to prepare data for RTL or to check RTL outputs.
\paragraph{Line 68}
\noindent\colorbox{gray!8}{\parbox{\dimexpr\linewidth-2\fboxsep\relax}{\raggedright\ttfamily\scriptsize         scores\_out[sample\_idx] = scores}}
\par\smallskip
This line handles class scores. The score vector is the direct software equivalent of the RTL MAC loop over weights and features.
\paragraph{Line 69}
\noindent\colorbox{gray!8}{\parbox{\dimexpr\linewidth-2\fboxsep\relax}{\raggedright\ttfamily\scriptsize         preds[sample\_idx] = pred}}
\par\smallskip
This line handles predicted labels. It checks that argmax behavior in Python and RTL chooses the same class.
\paragraph{Line 70}
\noindent\colorbox{gray!8}{\parbox{\dimexpr\linewidth-2\fboxsep\relax}{\raggedright\ttfamily\scriptsize         if pred != label and (int(scores[pred]) - int(scores[label])) <= LEARN\_MARGIN\_Q:}}
\par\smallskip
This line handles class scores. The score vector is the direct software equivalent of the RTL MAC loop over weights and features.
\paragraph{Line 71}
\noindent\colorbox{gray!8}{\parbox{\dimexpr\linewidth-2\fboxsep\relax}{\raggedright\ttfamily\scriptsize             delta = np.array([int(v) >> LEARN\_SHIFT for v in x], dtype=np.int64)}}
\par\smallskip
This performs vectorized numerical work. The fixed-point golden models use NumPy so score, margin, and update calculations are explicit and repeatable.
\paragraph{Line 72}
\noindent\colorbox{gray!8}{\parbox{\dimexpr\linewidth-2\fboxsep\relax}{\raggedright\ttfamily\scriptsize             weights[label] = np.array([clamp\_i16(int(w) + int(d)) for w, d in zip(weights[label], delta)], dtype=np.int64)}}
\par\smallskip
This performs vectorized numerical work. The fixed-point golden models use NumPy so score, margin, and update calculations are explicit and repeatable.
\paragraph{Line 73}
\noindent\colorbox{gray!8}{\parbox{\dimexpr\linewidth-2\fboxsep\relax}{\raggedright\ttfamily\scriptsize             weights[pred] = np.array([clamp\_i16(int(w) - int(d)) for w, d in zip(weights[pred], delta)], dtype=np.int64)}}
\par\smallskip
This performs vectorized numerical work. The fixed-point golden models use NumPy so score, margin, and update calculations are explicit and repeatable.
\paragraph{Line 74}
\noindent\colorbox{gray!8}{\parbox{\dimexpr\linewidth-2\fboxsep\relax}{\raggedright\ttfamily\scriptsize             bias[label] += BIAS\_LEARN\_DELTA}}
\par\smallskip
This line handles the classifier parameters. Those arrays mirror the RTL memories, so mismatches reveal an export, quantization, or hardware-update bug.
\paragraph{Line 75}
\noindent\colorbox{gray!8}{\parbox{\dimexpr\linewidth-2\fboxsep\relax}{\raggedright\ttfamily\scriptsize             bias[pred] -= BIAS\_LEARN\_DELTA}}
\par\smallskip
This line handles the classifier parameters. Those arrays mirror the RTL memories, so mismatches reveal an export, quantization, or hardware-update bug.
\paragraph{Line 76}
\noindent\colorbox{gray!8}{\parbox{\dimexpr\linewidth-2\fboxsep\relax}{\raggedright\ttfamily\scriptsize             count += 1}}
\par\smallskip
This Python line participates in the export, validation, or report-generation pipeline. It is executed on the host PC to prepare data for RTL or to check RTL outputs.
\paragraph{Line 77}
\noindent\colorbox{gray!8}{\parbox{\dimexpr\linewidth-2\fboxsep\relax}{\raggedright\ttfamily\scriptsize             update\_applied[sample\_idx] = 1}}
\par\smallskip
This line handles the online-learning rule or its validation fields. It is central to proving that the FPGA update path is not bypassed.
\paragraph{Line 78}
\noindent\colorbox{gray!8}{\parbox{\dimexpr\linewidth-2\fboxsep\relax}{\raggedright\ttfamily\scriptsize         update\_count[sample\_idx] = count}}
\par\smallskip
This line handles the online-learning rule or its validation fields. It is central to proving that the FPGA update path is not bypassed.
\paragraph{Line 79}
\noindent\colorbox{gray!8}{\parbox{\dimexpr\linewidth-2\fboxsep\relax}{\raggedright\ttfamily\scriptsize     return scores\_out, preds, update\_applied, update\_count}}
\par\smallskip
This line handles class scores. The score vector is the direct software equivalent of the RTL MAC loop over weights and features.
\paragraph{Line 80}
\noindent\colorbox{gray!8}{\parbox{\dimexpr\linewidth-2\fboxsep\relax}{\raggedright\ttfamily\scriptsize \emph{blank line}}}
\par\smallskip
This blank line is not hardware or software behavior. It separates nearby declarations or operations so the file remains readable and reviewable.
\paragraph{Line 81}
\noindent\colorbox{gray!8}{\parbox{\dimexpr\linewidth-2\fboxsep\relax}{\raggedright\ttfamily\scriptsize \emph{blank line}}}
\par\smallskip
This blank line is not hardware or software behavior. It separates nearby declarations or operations so the file remains readable and reviewable.
\paragraph{Line 82}
\noindent\colorbox{gray!8}{\parbox{\dimexpr\linewidth-2\fboxsep\relax}{\raggedright\ttfamily\scriptsize def main() -> None:}}
\par\smallskip
This defines a Python function. The function packages one repeatable step such as loading a memory file, simulating the golden learning rule, validating CSVs, or generating plots.
\paragraph{Line 83}
\noindent\colorbox{gray!8}{\parbox{\dimexpr\linewidth-2\fboxsep\relax}{\raggedright\ttfamily\scriptsize     parser = argparse.ArgumentParser(description="Validate on-chip learning RTL against fixed-point online-learning golden model.")}}
\par\smallskip
This contributes to the command-line interface. It lets the rerun flow choose paths, sample counts, or output locations without editing code.
\paragraph{Line 84}
\noindent\colorbox{gray!8}{\parbox{\dimexpr\linewidth-2\fboxsep\relax}{\raggedright\ttfamily\scriptsize     parser.add\_argument("--rtl-dir", type=Path, default=repo\_root() / "cifar10\_neuromorphic" / "rtl")}}
\par\smallskip
This contributes to the command-line interface. It lets the rerun flow choose paths, sample counts, or output locations without editing code.
\paragraph{Line 85}
\noindent\colorbox{gray!8}{\parbox{\dimexpr\linewidth-2\fboxsep\relax}{\raggedright\ttfamily\scriptsize     parser.add\_argument("--results-csv", type=Path, default=repo\_root() / "cifar10\_neuromorphic" / "cifar10\_spikingformer\_head\_onchip\_learning\_results.csv")}}
\par\smallskip
This contributes to the command-line interface. It lets the rerun flow choose paths, sample counts, or output locations without editing code.
\paragraph{Line 86}
\noindent\colorbox{gray!8}{\parbox{\dimexpr\linewidth-2\fboxsep\relax}{\raggedright\ttfamily\scriptsize     parser.add\_argument("--out-dir", type=Path, default=repo\_root() / "cifar10\_neuromorphic" / "outputs\_spikingformer\_head\_accelerator")}}
\par\smallskip
This contributes to the command-line interface. It lets the rerun flow choose paths, sample counts, or output locations without editing code.
\paragraph{Line 87}
\noindent\colorbox{gray!8}{\parbox{\dimexpr\linewidth-2\fboxsep\relax}{\raggedright\ttfamily\scriptsize     args = parser.parse\_args()}}
\par\smallskip
This Python line participates in the export, validation, or report-generation pipeline. It is executed on the host PC to prepare data for RTL or to check RTL outputs.
\paragraph{Line 88}
\noindent\colorbox{gray!8}{\parbox{\dimexpr\linewidth-2\fboxsep\relax}{\raggedright\ttfamily\scriptsize \emph{blank line}}}
\par\smallskip
This blank line is not hardware or software behavior. It separates nearby declarations or operations so the file remains readable and reviewable.
\paragraph{Line 89}
\noindent\colorbox{gray!8}{\parbox{\dimexpr\linewidth-2\fboxsep\relax}{\raggedright\ttfamily\scriptsize     rows = load\_rows(args.results\_csv)}}
\par\smallskip
This Python line participates in the export, validation, or report-generation pipeline. It is executed on the host PC to prepare data for RTL or to check RTL outputs.
\paragraph{Line 90}
\noindent\colorbox{gray!8}{\parbox{\dimexpr\linewidth-2\fboxsep\relax}{\raggedright\ttfamily\scriptsize     scores\_ref, preds\_ref, updates\_ref, update\_count\_ref = simulate\_online\_learning(args.rtl\_dir, len(rows))}}
\par\smallskip
This line handles class scores. The score vector is the direct software equivalent of the RTL MAC loop over weights and features.
\paragraph{Line 91}
\noindent\colorbox{gray!8}{\parbox{\dimexpr\linewidth-2\fboxsep\relax}{\raggedright\ttfamily\scriptsize     scores\_rtl = np.array([[int(row[f"score\{c\}"]) for c in range(NUM\_CLASSES)] for row in rows], dtype=np.int64)}}
\par\smallskip
This performs vectorized numerical work. The fixed-point golden models use NumPy so score, margin, and update calculations are explicit and repeatable.
\paragraph{Line 92}
\noindent\colorbox{gray!8}{\parbox{\dimexpr\linewidth-2\fboxsep\relax}{\raggedright\ttfamily\scriptsize     preds\_rtl = np.array([int(row["predicted"]) for row in rows], dtype=np.int64)}}
\par\smallskip
This performs vectorized numerical work. The fixed-point golden models use NumPy so score, margin, and update calculations are explicit and repeatable.
\paragraph{Line 93}
\noindent\colorbox{gray!8}{\parbox{\dimexpr\linewidth-2\fboxsep\relax}{\raggedright\ttfamily\scriptsize     labels = np.array([int(row["label"]) for row in rows], dtype=np.int64)}}
\par\smallskip
This performs vectorized numerical work. The fixed-point golden models use NumPy so score, margin, and update calculations are explicit and repeatable.
\paragraph{Line 94}
\noindent\colorbox{gray!8}{\parbox{\dimexpr\linewidth-2\fboxsep\relax}{\raggedright\ttfamily\scriptsize     updates\_rtl = np.array([int(row["update\_applied"]) for row in rows], dtype=np.int64)}}
\par\smallskip
This performs vectorized numerical work. The fixed-point golden models use NumPy so score, margin, and update calculations are explicit and repeatable.
\paragraph{Line 95}
\noindent\colorbox{gray!8}{\parbox{\dimexpr\linewidth-2\fboxsep\relax}{\raggedright\ttfamily\scriptsize     update\_count\_rtl = np.array([int(row["update\_count"]) for row in rows], dtype=np.int64)}}
\par\smallskip
This performs vectorized numerical work. The fixed-point golden models use NumPy so score, margin, and update calculations are explicit and repeatable.
\paragraph{Line 96}
\noindent\colorbox{gray!8}{\parbox{\dimexpr\linewidth-2\fboxsep\relax}{\raggedright\ttfamily\scriptsize \emph{blank line}}}
\par\smallskip
This blank line is not hardware or software behavior. It separates nearby declarations or operations so the file remains readable and reviewable.
\paragraph{Line 97}
\noindent\colorbox{gray!8}{\parbox{\dimexpr\linewidth-2\fboxsep\relax}{\raggedright\ttfamily\scriptsize     score\_error = scores\_rtl - scores\_ref}}
\par\smallskip
This line handles class scores. The score vector is the direct software equivalent of the RTL MAC loop over weights and features.
\paragraph{Line 98}
\noindent\colorbox{gray!8}{\parbox{\dimexpr\linewidth-2\fboxsep\relax}{\raggedright\ttfamily\scriptsize     metrics = \{}}
\par\smallskip
This Python line participates in the export, validation, or report-generation pipeline. It is executed on the host PC to prepare data for RTL or to check RTL outputs.
\paragraph{Line 99}
\noindent\colorbox{gray!8}{\parbox{\dimexpr\linewidth-2\fboxsep\relax}{\raggedright\ttfamily\scriptsize         "samples": int(len(rows)),}}
\par\smallskip
This Python line participates in the export, validation, or report-generation pipeline. It is executed on the host PC to prepare data for RTL or to check RTL outputs.
\paragraph{Line 100}
\noindent\colorbox{gray!8}{\parbox{\dimexpr\linewidth-2\fboxsep\relax}{\raggedright\ttfamily\scriptsize         "accuracy": float((preds\_rtl == labels).mean()) if rows else 0.0,}}
\par\smallskip
This line handles predicted labels. It checks that argmax behavior in Python and RTL chooses the same class.
\paragraph{Line 101}
\noindent\colorbox{gray!8}{\parbox{\dimexpr\linewidth-2\fboxsep\relax}{\raggedright\ttfamily\scriptsize         "correct": int((preds\_rtl == labels).sum()),}}
\par\smallskip
This line handles predicted labels. It checks that argmax behavior in Python and RTL chooses the same class.
\paragraph{Line 102}
\noindent\colorbox{gray!8}{\parbox{\dimexpr\linewidth-2\fboxsep\relax}{\raggedright\ttfamily\scriptsize         "updates": int(updates\_rtl.sum()),}}
\par\smallskip
This line handles the online-learning rule or its validation fields. It is central to proving that the FPGA update path is not bypassed.
\paragraph{Line 103}
\noindent\colorbox{gray!8}{\parbox{\dimexpr\linewidth-2\fboxsep\relax}{\raggedright\ttfamily\scriptsize         "score\_mismatch\_count": int((score\_error != 0).sum()),}}
\par\smallskip
This Python line participates in the export, validation, or report-generation pipeline. It is executed on the host PC to prepare data for RTL or to check RTL outputs.
\paragraph{Line 104}
\noindent\colorbox{gray!8}{\parbox{\dimexpr\linewidth-2\fboxsep\relax}{\raggedright\ttfamily\scriptsize         "max\_abs\_score\_error": int(np.abs(score\_error).max()) if score\_error.size else 0,}}
\par\smallskip
This performs vectorized numerical work. The fixed-point golden models use NumPy so score, margin, and update calculations are explicit and repeatable.
\paragraph{Line 105}
\noindent\colorbox{gray!8}{\parbox{\dimexpr\linewidth-2\fboxsep\relax}{\raggedright\ttfamily\scriptsize         "prediction\_mismatch\_count": int((preds\_rtl != preds\_ref).sum()),}}
\par\smallskip
This line handles predicted labels. It checks that argmax behavior in Python and RTL chooses the same class.
\paragraph{Line 106}
\noindent\colorbox{gray!8}{\parbox{\dimexpr\linewidth-2\fboxsep\relax}{\raggedright\ttfamily\scriptsize         "update\_flag\_mismatch\_count": int((updates\_rtl != updates\_ref).sum()),}}
\par\smallskip
This line handles the online-learning rule or its validation fields. It is central to proving that the FPGA update path is not bypassed.
\paragraph{Line 107}
\noindent\colorbox{gray!8}{\parbox{\dimexpr\linewidth-2\fboxsep\relax}{\raggedright\ttfamily\scriptsize         "update\_count\_mismatch\_count": int((update\_count\_rtl != update\_count\_ref).sum()),}}
\par\smallskip
This line handles the online-learning rule or its validation fields. It is central to proving that the FPGA update path is not bypassed.
\paragraph{Line 108}
\noindent\colorbox{gray!8}{\parbox{\dimexpr\linewidth-2\fboxsep\relax}{\raggedright\ttfamily\scriptsize         "learn\_shift": LEARN\_SHIFT,}}
\par\smallskip
This Python line participates in the export, validation, or report-generation pipeline. It is executed on the host PC to prepare data for RTL or to check RTL outputs.
\paragraph{Line 109}
\noindent\colorbox{gray!8}{\parbox{\dimexpr\linewidth-2\fboxsep\relax}{\raggedright\ttfamily\scriptsize         "bias\_learn\_delta": BIAS\_LEARN\_DELTA,}}
\par\smallskip
This line handles the classifier parameters. Those arrays mirror the RTL memories, so mismatches reveal an export, quantization, or hardware-update bug.
\paragraph{Line 110}
\noindent\colorbox{gray!8}{\parbox{\dimexpr\linewidth-2\fboxsep\relax}{\raggedright\ttfamily\scriptsize         "learn\_margin\_q": LEARN\_MARGIN\_Q,}}
\par\smallskip
This Python line participates in the export, validation, or report-generation pipeline. It is executed on the host PC to prepare data for RTL or to check RTL outputs.
\paragraph{Line 111}
\noindent\colorbox{gray!8}{\parbox{\dimexpr\linewidth-2\fboxsep\relax}{\raggedright\ttfamily\scriptsize         "infer\_skip\_abs\_thresh": INFER\_SKIP\_ABS\_THRESH,}}
\par\smallskip
This Python line participates in the export, validation, or report-generation pipeline. It is executed on the host PC to prepare data for RTL or to check RTL outputs.
\paragraph{Line 112}
\noindent\colorbox{gray!8}{\parbox{\dimexpr\linewidth-2\fboxsep\relax}{\raggedright\ttfamily\scriptsize     \}}}
\par\smallskip
This Python line participates in the export, validation, or report-generation pipeline. It is executed on the host PC to prepare data for RTL or to check RTL outputs.
\paragraph{Line 113}
\noindent\colorbox{gray!8}{\parbox{\dimexpr\linewidth-2\fboxsep\relax}{\raggedright\ttfamily\scriptsize     args.out\_dir.mkdir(parents=True, exist\_ok=True)}}
\par\smallskip
This Python line participates in the export, validation, or report-generation pipeline. It is executed on the host PC to prepare data for RTL or to check RTL outputs.
\paragraph{Line 114}
\noindent\colorbox{gray!8}{\parbox{\dimexpr\linewidth-2\fboxsep\relax}{\raggedright\ttfamily\scriptsize     (args.out\_dir / "spikingformer\_head\_onchip\_learning\_validation.json").write\_text(json.dumps(metrics, indent=2), encoding="utf-8")}}
\par\smallskip
This writes machine-readable output for later reports. Keeping JSON/CSV artifacts makes the verification auditable.
\paragraph{Line 115}
\noindent\colorbox{gray!8}{\parbox{\dimexpr\linewidth-2\fboxsep\relax}{\raggedright\ttfamily\scriptsize     print(json.dumps(metrics, indent=2))}}
\par\smallskip
This writes machine-readable output for later reports. Keeping JSON/CSV artifacts makes the verification auditable.
\paragraph{Line 116}
\noindent\colorbox{gray!8}{\parbox{\dimexpr\linewidth-2\fboxsep\relax}{\raggedright\ttfamily\scriptsize \emph{blank line}}}
\par\smallskip
This blank line is not hardware or software behavior. It separates nearby declarations or operations so the file remains readable and reviewable.
\paragraph{Line 117}
\noindent\colorbox{gray!8}{\parbox{\dimexpr\linewidth-2\fboxsep\relax}{\raggedright\ttfamily\scriptsize \emph{blank line}}}
\par\smallskip
This blank line is not hardware or software behavior. It separates nearby declarations or operations so the file remains readable and reviewable.
\paragraph{Line 118}
\noindent\colorbox{gray!8}{\parbox{\dimexpr\linewidth-2\fboxsep\relax}{\raggedright\ttfamily\scriptsize if \_\_name\_\_ == "\_\_main\_\_":}}
\par\smallskip
This Python idiom runs main() only when the file is executed as a script. It still allows the file's functions to be imported for tests or reuse.
\paragraph{Line 119}
\noindent\colorbox{gray!8}{\parbox{\dimexpr\linewidth-2\fboxsep\relax}{\raggedright\ttfamily\scriptsize     main()}}
\par\smallskip
This Python line participates in the export, validation, or report-generation pipeline. It is executed on the host PC to prepare data for RTL or to check RTL outputs.
\clearpage
\subsection{cifar10\_neuromorphic\textbackslash{}generate\_spikingformer\_debug\_reports.py}
Each entry below gives the exact source line followed by its hardware/software role in the implemented accelerator pipeline.
\paragraph{Line 1}
\noindent\colorbox{gray!8}{\parbox{\dimexpr\linewidth-2\fboxsep\relax}{\raggedright\ttfamily\scriptsize from \_\_future\_\_ import annotations}}
\par\smallskip
This imports a Python library or symbol. The scripts use standard parsing, numeric arrays, plotting, and file APIs to make the hardware artifacts reproducible.
\paragraph{Line 2}
\noindent\colorbox{gray!8}{\parbox{\dimexpr\linewidth-2\fboxsep\relax}{\raggedright\ttfamily\scriptsize \emph{blank line}}}
\par\smallskip
This blank line is not hardware or software behavior. It separates nearby declarations or operations so the file remains readable and reviewable.
\paragraph{Line 3}
\noindent\colorbox{gray!8}{\parbox{\dimexpr\linewidth-2\fboxsep\relax}{\raggedright\ttfamily\scriptsize import csv}}
\par\smallskip
This imports a Python library or symbol. The scripts use standard parsing, numeric arrays, plotting, and file APIs to make the hardware artifacts reproducible.
\paragraph{Line 4}
\noindent\colorbox{gray!8}{\parbox{\dimexpr\linewidth-2\fboxsep\relax}{\raggedright\ttfamily\scriptsize import json}}
\par\smallskip
This imports a Python library or symbol. The scripts use standard parsing, numeric arrays, plotting, and file APIs to make the hardware artifacts reproducible.
\paragraph{Line 5}
\noindent\colorbox{gray!8}{\parbox{\dimexpr\linewidth-2\fboxsep\relax}{\raggedright\ttfamily\scriptsize import math}}
\par\smallskip
This imports a Python library or symbol. The scripts use standard parsing, numeric arrays, plotting, and file APIs to make the hardware artifacts reproducible.
\paragraph{Line 6}
\noindent\colorbox{gray!8}{\parbox{\dimexpr\linewidth-2\fboxsep\relax}{\raggedright\ttfamily\scriptsize import re}}
\par\smallskip
This imports a Python library or symbol. The scripts use standard parsing, numeric arrays, plotting, and file APIs to make the hardware artifacts reproducible.
\paragraph{Line 7}
\noindent\colorbox{gray!8}{\parbox{\dimexpr\linewidth-2\fboxsep\relax}{\raggedright\ttfamily\scriptsize import shutil}}
\par\smallskip
This imports a Python library or symbol. The scripts use standard parsing, numeric arrays, plotting, and file APIs to make the hardware artifacts reproducible.
\paragraph{Line 8}
\noindent\colorbox{gray!8}{\parbox{\dimexpr\linewidth-2\fboxsep\relax}{\raggedright\ttfamily\scriptsize import subprocess}}
\par\smallskip
This imports a Python library or symbol. The scripts use standard parsing, numeric arrays, plotting, and file APIs to make the hardware artifacts reproducible.
\paragraph{Line 9}
\noindent\colorbox{gray!8}{\parbox{\dimexpr\linewidth-2\fboxsep\relax}{\raggedright\ttfamily\scriptsize from collections import Counter}}
\par\smallskip
This imports a Python library or symbol. The scripts use standard parsing, numeric arrays, plotting, and file APIs to make the hardware artifacts reproducible.
\paragraph{Line 10}
\noindent\colorbox{gray!8}{\parbox{\dimexpr\linewidth-2\fboxsep\relax}{\raggedright\ttfamily\scriptsize from pathlib import Path}}
\par\smallskip
This imports a Python library or symbol. The scripts use standard parsing, numeric arrays, plotting, and file APIs to make the hardware artifacts reproducible.
\paragraph{Line 11}
\noindent\colorbox{gray!8}{\parbox{\dimexpr\linewidth-2\fboxsep\relax}{\raggedright\ttfamily\scriptsize from typing import Iterable}}
\par\smallskip
This imports a Python library or symbol. The scripts use standard parsing, numeric arrays, plotting, and file APIs to make the hardware artifacts reproducible.
\paragraph{Line 12}
\noindent\colorbox{gray!8}{\parbox{\dimexpr\linewidth-2\fboxsep\relax}{\raggedright\ttfamily\scriptsize \emph{blank line}}}
\par\smallskip
This blank line is not hardware or software behavior. It separates nearby declarations or operations so the file remains readable and reviewable.
\paragraph{Line 13}
\noindent\colorbox{gray!8}{\parbox{\dimexpr\linewidth-2\fboxsep\relax}{\raggedright\ttfamily\scriptsize import matplotlib.pyplot as plt}}
\par\smallskip
This imports a Python library or symbol. The scripts use standard parsing, numeric arrays, plotting, and file APIs to make the hardware artifacts reproducible.
\paragraph{Line 14}
\noindent\colorbox{gray!8}{\parbox{\dimexpr\linewidth-2\fboxsep\relax}{\raggedright\ttfamily\scriptsize import numpy as np}}
\par\smallskip
This imports a Python library or symbol. The scripts use standard parsing, numeric arrays, plotting, and file APIs to make the hardware artifacts reproducible.
\paragraph{Line 15}
\noindent\colorbox{gray!8}{\parbox{\dimexpr\linewidth-2\fboxsep\relax}{\raggedright\ttfamily\scriptsize from matplotlib.backends.backend\_pdf import PdfPages}}
\par\smallskip
This imports a Python library or symbol. The scripts use standard parsing, numeric arrays, plotting, and file APIs to make the hardware artifacts reproducible.
\paragraph{Line 16}
\noindent\colorbox{gray!8}{\parbox{\dimexpr\linewidth-2\fboxsep\relax}{\raggedright\ttfamily\scriptsize \emph{blank line}}}
\par\smallskip
This blank line is not hardware or software behavior. It separates nearby declarations or operations so the file remains readable and reviewable.
\paragraph{Line 17}
\noindent\colorbox{gray!8}{\parbox{\dimexpr\linewidth-2\fboxsep\relax}{\raggedright\ttfamily\scriptsize \emph{blank line}}}
\par\smallskip
This blank line is not hardware or software behavior. It separates nearby declarations or operations so the file remains readable and reviewable.
\paragraph{Line 18}
\noindent\colorbox{gray!8}{\parbox{\dimexpr\linewidth-2\fboxsep\relax}{\raggedright\ttfamily\scriptsize NUM\_CLASSES = 10}}
\par\smallskip
This Python line participates in the export, validation, or report-generation pipeline. It is executed on the host PC to prepare data for RTL or to check RTL outputs.
\paragraph{Line 19}
\noindent\colorbox{gray!8}{\parbox{\dimexpr\linewidth-2\fboxsep\relax}{\raggedright\ttfamily\scriptsize FEATURE\_DIM = 384}}
\par\smallskip
This Python line participates in the export, validation, or report-generation pipeline. It is executed on the host PC to prepare data for RTL or to check RTL outputs.
\paragraph{Line 20}
\noindent\colorbox{gray!8}{\parbox{\dimexpr\linewidth-2\fboxsep\relax}{\raggedright\ttfamily\scriptsize CLASS\_NAMES = [str(i) for i in range(NUM\_CLASSES)]}}
\par\smallskip
This Python line participates in the export, validation, or report-generation pipeline. It is executed on the host PC to prepare data for RTL or to check RTL outputs.
\paragraph{Line 21}
\noindent\colorbox{gray!8}{\parbox{\dimexpr\linewidth-2\fboxsep\relax}{\raggedright\ttfamily\scriptsize \emph{blank line}}}
\par\smallskip
This blank line is not hardware or software behavior. It separates nearby declarations or operations so the file remains readable and reviewable.
\paragraph{Line 22}
\noindent\colorbox{gray!8}{\parbox{\dimexpr\linewidth-2\fboxsep\relax}{\raggedright\ttfamily\scriptsize STATE\_NAMES = \{}}
\par\smallskip
This Python line participates in the export, validation, or report-generation pipeline. It is executed on the host PC to prepare data for RTL or to check RTL outputs.
\paragraph{Line 23}
\noindent\colorbox{gray!8}{\parbox{\dimexpr\linewidth-2\fboxsep\relax}{\raggedright\ttfamily\scriptsize     0: "IDLE",}}
\par\smallskip
This Python line participates in the export, validation, or report-generation pipeline. It is executed on the host PC to prepare data for RTL or to check RTL outputs.
\paragraph{Line 24}
\noindent\colorbox{gray!8}{\parbox{\dimexpr\linewidth-2\fboxsep\relax}{\raggedright\ttfamily\scriptsize     1: "LOAD",}}
\par\smallskip
This Python line participates in the export, validation, or report-generation pipeline. It is executed on the host PC to prepare data for RTL or to check RTL outputs.
\paragraph{Line 25}
\noindent\colorbox{gray!8}{\parbox{\dimexpr\linewidth-2\fboxsep\relax}{\raggedright\ttfamily\scriptsize     2: "CLASS\_INIT",}}
\par\smallskip
This Python line participates in the export, validation, or report-generation pipeline. It is executed on the host PC to prepare data for RTL or to check RTL outputs.
\paragraph{Line 26}
\noindent\colorbox{gray!8}{\parbox{\dimexpr\linewidth-2\fboxsep\relax}{\raggedright\ttfamily\scriptsize     3: "MAC\_READ",}}
\par\smallskip
This Python line participates in the export, validation, or report-generation pipeline. It is executed on the host PC to prepare data for RTL or to check RTL outputs.
\paragraph{Line 27}
\noindent\colorbox{gray!8}{\parbox{\dimexpr\linewidth-2\fboxsep\relax}{\raggedright\ttfamily\scriptsize     4: "MAC\_MUL",}}
\par\smallskip
This Python line participates in the export, validation, or report-generation pipeline. It is executed on the host PC to prepare data for RTL or to check RTL outputs.
\paragraph{Line 28}
\noindent\colorbox{gray!8}{\parbox{\dimexpr\linewidth-2\fboxsep\relax}{\raggedright\ttfamily\scriptsize     5: "MAC\_ACCUM",}}
\par\smallskip
This Python line participates in the export, validation, or report-generation pipeline. It is executed on the host PC to prepare data for RTL or to check RTL outputs.
\paragraph{Line 29}
\noindent\colorbox{gray!8}{\parbox{\dimexpr\linewidth-2\fboxsep\relax}{\raggedright\ttfamily\scriptsize     6: "CLASS\_DONE",}}
\par\smallskip
This Python line participates in the export, validation, or report-generation pipeline. It is executed on the host PC to prepare data for RTL or to check RTL outputs.
\paragraph{Line 30}
\noindent\colorbox{gray!8}{\parbox{\dimexpr\linewidth-2\fboxsep\relax}{\raggedright\ttfamily\scriptsize     7: "RESULT",}}
\par\smallskip
This Python line participates in the export, validation, or report-generation pipeline. It is executed on the host PC to prepare data for RTL or to check RTL outputs.
\paragraph{Line 31}
\noindent\colorbox{gray!8}{\parbox{\dimexpr\linewidth-2\fboxsep\relax}{\raggedright\ttfamily\scriptsize     8: "TRAIN\_CHECK",}}
\par\smallskip
This Python line participates in the export, validation, or report-generation pipeline. It is executed on the host PC to prepare data for RTL or to check RTL outputs.
\paragraph{Line 32}
\noindent\colorbox{gray!8}{\parbox{\dimexpr\linewidth-2\fboxsep\relax}{\raggedright\ttfamily\scriptsize     9: "TRAIN\_READ\_TRUE",}}
\par\smallskip
This Python line participates in the export, validation, or report-generation pipeline. It is executed on the host PC to prepare data for RTL or to check RTL outputs.
\paragraph{Line 33}
\noindent\colorbox{gray!8}{\parbox{\dimexpr\linewidth-2\fboxsep\relax}{\raggedright\ttfamily\scriptsize     10: "TRAIN\_CALC\_TRUE",}}
\par\smallskip
This Python line participates in the export, validation, or report-generation pipeline. It is executed on the host PC to prepare data for RTL or to check RTL outputs.
\paragraph{Line 34}
\noindent\colorbox{gray!8}{\parbox{\dimexpr\linewidth-2\fboxsep\relax}{\raggedright\ttfamily\scriptsize     11: "TRAIN\_WRITE\_TRUE",}}
\par\smallskip
This Python line participates in the export, validation, or report-generation pipeline. It is executed on the host PC to prepare data for RTL or to check RTL outputs.
\paragraph{Line 35}
\noindent\colorbox{gray!8}{\parbox{\dimexpr\linewidth-2\fboxsep\relax}{\raggedright\ttfamily\scriptsize     12: "TRAIN\_READ\_PRED",}}
\par\smallskip
This Python line participates in the export, validation, or report-generation pipeline. It is executed on the host PC to prepare data for RTL or to check RTL outputs.
\paragraph{Line 36}
\noindent\colorbox{gray!8}{\parbox{\dimexpr\linewidth-2\fboxsep\relax}{\raggedright\ttfamily\scriptsize     13: "TRAIN\_CALC\_PRED",}}
\par\smallskip
This Python line participates in the export, validation, or report-generation pipeline. It is executed on the host PC to prepare data for RTL or to check RTL outputs.
\paragraph{Line 37}
\noindent\colorbox{gray!8}{\parbox{\dimexpr\linewidth-2\fboxsep\relax}{\raggedright\ttfamily\scriptsize     14: "TRAIN\_WRITE\_PRED",}}
\par\smallskip
This Python line participates in the export, validation, or report-generation pipeline. It is executed on the host PC to prepare data for RTL or to check RTL outputs.
\paragraph{Line 38}
\noindent\colorbox{gray!8}{\parbox{\dimexpr\linewidth-2\fboxsep\relax}{\raggedright\ttfamily\scriptsize     15: "TRAIN\_BIAS\_TRUE",}}
\par\smallskip
This Python line participates in the export, validation, or report-generation pipeline. It is executed on the host PC to prepare data for RTL or to check RTL outputs.
\paragraph{Line 39}
\noindent\colorbox{gray!8}{\parbox{\dimexpr\linewidth-2\fboxsep\relax}{\raggedright\ttfamily\scriptsize     16: "TRAIN\_BIAS\_PRED",}}
\par\smallskip
This Python line participates in the export, validation, or report-generation pipeline. It is executed on the host PC to prepare data for RTL or to check RTL outputs.
\paragraph{Line 40}
\noindent\colorbox{gray!8}{\parbox{\dimexpr\linewidth-2\fboxsep\relax}{\raggedright\ttfamily\scriptsize \}}}
\par\smallskip
This Python line participates in the export, validation, or report-generation pipeline. It is executed on the host PC to prepare data for RTL or to check RTL outputs.
\paragraph{Line 41}
\noindent\colorbox{gray!8}{\parbox{\dimexpr\linewidth-2\fboxsep\relax}{\raggedright\ttfamily\scriptsize \emph{blank line}}}
\par\smallskip
This blank line is not hardware or software behavior. It separates nearby declarations or operations so the file remains readable and reviewable.
\paragraph{Line 42}
\noindent\colorbox{gray!8}{\parbox{\dimexpr\linewidth-2\fboxsep\relax}{\raggedright\ttfamily\scriptsize \emph{blank line}}}
\par\smallskip
This blank line is not hardware or software behavior. It separates nearby declarations or operations so the file remains readable and reviewable.
\paragraph{Line 43}
\noindent\colorbox{gray!8}{\parbox{\dimexpr\linewidth-2\fboxsep\relax}{\raggedright\ttfamily\scriptsize def repo\_root() -> Path:}}
\par\smallskip
This defines a Python function. The function packages one repeatable step such as loading a memory file, simulating the golden learning rule, validating CSVs, or generating plots.
\paragraph{Line 44}
\noindent\colorbox{gray!8}{\parbox{\dimexpr\linewidth-2\fboxsep\relax}{\raggedright\ttfamily\scriptsize     return Path(\_\_file\_\_).resolve().parents[1]}}
\par\smallskip
This uses path objects so generated artifacts are addressed consistently across Windows shells and Python code.
\paragraph{Line 45}
\noindent\colorbox{gray!8}{\parbox{\dimexpr\linewidth-2\fboxsep\relax}{\raggedright\ttfamily\scriptsize \emph{blank line}}}
\par\smallskip
This blank line is not hardware or software behavior. It separates nearby declarations or operations so the file remains readable and reviewable.
\paragraph{Line 46}
\noindent\colorbox{gray!8}{\parbox{\dimexpr\linewidth-2\fboxsep\relax}{\raggedright\ttfamily\scriptsize \emph{blank line}}}
\par\smallskip
This blank line is not hardware or software behavior. It separates nearby declarations or operations so the file remains readable and reviewable.
\paragraph{Line 47}
\noindent\colorbox{gray!8}{\parbox{\dimexpr\linewidth-2\fboxsep\relax}{\raggedright\ttfamily\scriptsize def project\_dir() -> Path:}}
\par\smallskip
This defines a Python function. The function packages one repeatable step such as loading a memory file, simulating the golden learning rule, validating CSVs, or generating plots.
\paragraph{Line 48}
\noindent\colorbox{gray!8}{\parbox{\dimexpr\linewidth-2\fboxsep\relax}{\raggedright\ttfamily\scriptsize     return Path(\_\_file\_\_).resolve().parent}}
\par\smallskip
This uses path objects so generated artifacts are addressed consistently across Windows shells and Python code.
\paragraph{Line 49}
\noindent\colorbox{gray!8}{\parbox{\dimexpr\linewidth-2\fboxsep\relax}{\raggedright\ttfamily\scriptsize \emph{blank line}}}
\par\smallskip
This blank line is not hardware or software behavior. It separates nearby declarations or operations so the file remains readable and reviewable.
\paragraph{Line 50}
\noindent\colorbox{gray!8}{\parbox{\dimexpr\linewidth-2\fboxsep\relax}{\raggedright\ttfamily\scriptsize \emph{blank line}}}
\par\smallskip
This blank line is not hardware or software behavior. It separates nearby declarations or operations so the file remains readable and reviewable.
\paragraph{Line 51}
\noindent\colorbox{gray!8}{\parbox{\dimexpr\linewidth-2\fboxsep\relax}{\raggedright\ttfamily\scriptsize def out\_dir() -> Path:}}
\par\smallskip
This defines a Python function. The function packages one repeatable step such as loading a memory file, simulating the golden learning rule, validating CSVs, or generating plots.
\paragraph{Line 52}
\noindent\colorbox{gray!8}{\parbox{\dimexpr\linewidth-2\fboxsep\relax}{\raggedright\ttfamily\scriptsize     return project\_dir() / "outputs\_spikingformer\_head\_accelerator"}}
\par\smallskip
This Python line participates in the export, validation, or report-generation pipeline. It is executed on the host PC to prepare data for RTL or to check RTL outputs.
\paragraph{Line 53}
\noindent\colorbox{gray!8}{\parbox{\dimexpr\linewidth-2\fboxsep\relax}{\raggedright\ttfamily\scriptsize \emph{blank line}}}
\par\smallskip
This blank line is not hardware or software behavior. It separates nearby declarations or operations so the file remains readable and reviewable.
\paragraph{Line 54}
\noindent\colorbox{gray!8}{\parbox{\dimexpr\linewidth-2\fboxsep\relax}{\raggedright\ttfamily\scriptsize \emph{blank line}}}
\par\smallskip
This blank line is not hardware or software behavior. It separates nearby declarations or operations so the file remains readable and reviewable.
\paragraph{Line 55}
\noindent\colorbox{gray!8}{\parbox{\dimexpr\linewidth-2\fboxsep\relax}{\raggedright\ttfamily\scriptsize def debug\_dir() -> Path:}}
\par\smallskip
This defines a Python function. The function packages one repeatable step such as loading a memory file, simulating the golden learning rule, validating CSVs, or generating plots.
\paragraph{Line 56}
\noindent\colorbox{gray!8}{\parbox{\dimexpr\linewidth-2\fboxsep\relax}{\raggedright\ttfamily\scriptsize     return out\_dir() / "deep\_debug"}}
\par\smallskip
This Python line participates in the export, validation, or report-generation pipeline. It is executed on the host PC to prepare data for RTL or to check RTL outputs.
\paragraph{Line 57}
\noindent\colorbox{gray!8}{\parbox{\dimexpr\linewidth-2\fboxsep\relax}{\raggedright\ttfamily\scriptsize \emph{blank line}}}
\par\smallskip
This blank line is not hardware or software behavior. It separates nearby declarations or operations so the file remains readable and reviewable.
\paragraph{Line 58}
\noindent\colorbox{gray!8}{\parbox{\dimexpr\linewidth-2\fboxsep\relax}{\raggedright\ttfamily\scriptsize \emph{blank line}}}
\par\smallskip
This blank line is not hardware or software behavior. It separates nearby declarations or operations so the file remains readable and reviewable.
\paragraph{Line 59}
\noindent\colorbox{gray!8}{\parbox{\dimexpr\linewidth-2\fboxsep\relax}{\raggedright\ttfamily\scriptsize def read\_csv\_rows(path: Path) -> list[dict[str, str]]:}}
\par\smallskip
This defines a Python function. The function packages one repeatable step such as loading a memory file, simulating the golden learning rule, validating CSVs, or generating plots.
\paragraph{Line 60}
\noindent\colorbox{gray!8}{\parbox{\dimexpr\linewidth-2\fboxsep\relax}{\raggedright\ttfamily\scriptsize     with path.open("r", newline="", encoding="utf-8") as fh:}}
\par\smallskip
This reads an artifact produced by training, export, simulation, or Vivado. The validator/report depends on files generated by the actual run rather than hidden constants.
\paragraph{Line 61}
\noindent\colorbox{gray!8}{\parbox{\dimexpr\linewidth-2\fboxsep\relax}{\raggedright\ttfamily\scriptsize         return list(csv.DictReader(fh))}}
\par\smallskip
This Python line participates in the export, validation, or report-generation pipeline. It is executed on the host PC to prepare data for RTL or to check RTL outputs.
\paragraph{Line 62}
\noindent\colorbox{gray!8}{\parbox{\dimexpr\linewidth-2\fboxsep\relax}{\raggedright\ttfamily\scriptsize \emph{blank line}}}
\par\smallskip
This blank line is not hardware or software behavior. It separates nearby declarations or operations so the file remains readable and reviewable.
\paragraph{Line 63}
\noindent\colorbox{gray!8}{\parbox{\dimexpr\linewidth-2\fboxsep\relax}{\raggedright\ttfamily\scriptsize \emph{blank line}}}
\par\smallskip
This blank line is not hardware or software behavior. It separates nearby declarations or operations so the file remains readable and reviewable.
\paragraph{Line 64}
\noindent\colorbox{gray!8}{\parbox{\dimexpr\linewidth-2\fboxsep\relax}{\raggedright\ttfamily\scriptsize def col(rows: list[dict[str, str]], name: str, dtype=float) -> np.ndarray:}}
\par\smallskip
This defines a Python function. The function packages one repeatable step such as loading a memory file, simulating the golden learning rule, validating CSVs, or generating plots.
\paragraph{Line 65}
\noindent\colorbox{gray!8}{\parbox{\dimexpr\linewidth-2\fboxsep\relax}{\raggedright\ttfamily\scriptsize     if not rows:}}
\par\smallskip
This Python line participates in the export, validation, or report-generation pipeline. It is executed on the host PC to prepare data for RTL or to check RTL outputs.
\paragraph{Line 66}
\noindent\colorbox{gray!8}{\parbox{\dimexpr\linewidth-2\fboxsep\relax}{\raggedright\ttfamily\scriptsize         return np.array([], dtype=dtype)}}
\par\smallskip
This performs vectorized numerical work. The fixed-point golden models use NumPy so score, margin, and update calculations are explicit and repeatable.
\paragraph{Line 67}
\noindent\colorbox{gray!8}{\parbox{\dimexpr\linewidth-2\fboxsep\relax}{\raggedright\ttfamily\scriptsize     return np.array([dtype(row[name]) for row in rows])}}
\par\smallskip
This performs vectorized numerical work. The fixed-point golden models use NumPy so score, margin, and update calculations are explicit and repeatable.
\paragraph{Line 68}
\noindent\colorbox{gray!8}{\parbox{\dimexpr\linewidth-2\fboxsep\relax}{\raggedright\ttfamily\scriptsize \emph{blank line}}}
\par\smallskip
This blank line is not hardware or software behavior. It separates nearby declarations or operations so the file remains readable and reviewable.
\paragraph{Line 69}
\noindent\colorbox{gray!8}{\parbox{\dimexpr\linewidth-2\fboxsep\relax}{\raggedright\ttfamily\scriptsize \emph{blank line}}}
\par\smallskip
This blank line is not hardware or software behavior. It separates nearby declarations or operations so the file remains readable and reviewable.
\paragraph{Line 70}
\noindent\colorbox{gray!8}{\parbox{\dimexpr\linewidth-2\fboxsep\relax}{\raggedright\ttfamily\scriptsize def score\_matrix(rows: list[dict[str, str]]) -> np.ndarray:}}
\par\smallskip
This defines a Python function. The function packages one repeatable step such as loading a memory file, simulating the golden learning rule, validating CSVs, or generating plots.
\paragraph{Line 71}
\noindent\colorbox{gray!8}{\parbox{\dimexpr\linewidth-2\fboxsep\relax}{\raggedright\ttfamily\scriptsize     return np.array([[int(row[f"score\{i\}"]) for i in range(NUM\_CLASSES)] for row in rows], dtype=np.int64)}}
\par\smallskip
This performs vectorized numerical work. The fixed-point golden models use NumPy so score, margin, and update calculations are explicit and repeatable.
\paragraph{Line 72}
\noindent\colorbox{gray!8}{\parbox{\dimexpr\linewidth-2\fboxsep\relax}{\raggedright\ttfamily\scriptsize \emph{blank line}}}
\par\smallskip
This blank line is not hardware or software behavior. It separates nearby declarations or operations so the file remains readable and reviewable.
\paragraph{Line 73}
\noindent\colorbox{gray!8}{\parbox{\dimexpr\linewidth-2\fboxsep\relax}{\raggedright\ttfamily\scriptsize \emph{blank line}}}
\par\smallskip
This blank line is not hardware or software behavior. It separates nearby declarations or operations so the file remains readable and reviewable.
\paragraph{Line 74}
\noindent\colorbox{gray!8}{\parbox{\dimexpr\linewidth-2\fboxsep\relax}{\raggedright\ttfamily\scriptsize def signed\_hex\_to\_int(token: str, bits: int) -> int:}}
\par\smallskip
This defines a Python function. The function packages one repeatable step such as loading a memory file, simulating the golden learning rule, validating CSVs, or generating plots.
\paragraph{Line 75}
\noindent\colorbox{gray!8}{\parbox{\dimexpr\linewidth-2\fboxsep\relax}{\raggedright\ttfamily\scriptsize     value = int(token.strip(), 16)}}
\par\smallskip
This Python line participates in the export, validation, or report-generation pipeline. It is executed on the host PC to prepare data for RTL or to check RTL outputs.
\paragraph{Line 76}
\noindent\colorbox{gray!8}{\parbox{\dimexpr\linewidth-2\fboxsep\relax}{\raggedright\ttfamily\scriptsize     sign = 1 << (bits - 1)}}
\par\smallskip
This Python line participates in the export, validation, or report-generation pipeline. It is executed on the host PC to prepare data for RTL or to check RTL outputs.
\paragraph{Line 77}
\noindent\colorbox{gray!8}{\parbox{\dimexpr\linewidth-2\fboxsep\relax}{\raggedright\ttfamily\scriptsize     return value - (1 << bits) if value \& sign else value}}
\par\smallskip
This Python line participates in the export, validation, or report-generation pipeline. It is executed on the host PC to prepare data for RTL or to check RTL outputs.
\paragraph{Line 78}
\noindent\colorbox{gray!8}{\parbox{\dimexpr\linewidth-2\fboxsep\relax}{\raggedright\ttfamily\scriptsize \emph{blank line}}}
\par\smallskip
This blank line is not hardware or software behavior. It separates nearby declarations or operations so the file remains readable and reviewable.
\paragraph{Line 79}
\noindent\colorbox{gray!8}{\parbox{\dimexpr\linewidth-2\fboxsep\relax}{\raggedright\ttfamily\scriptsize \emph{blank line}}}
\par\smallskip
This blank line is not hardware or software behavior. It separates nearby declarations or operations so the file remains readable and reviewable.
\paragraph{Line 80}
\noindent\colorbox{gray!8}{\parbox{\dimexpr\linewidth-2\fboxsep\relax}{\raggedright\ttfamily\scriptsize def read\_hex\_mem(path: Path, bits: int) -> np.ndarray:}}
\par\smallskip
This defines a Python function. The function packages one repeatable step such as loading a memory file, simulating the golden learning rule, validating CSVs, or generating plots.
\paragraph{Line 81}
\noindent\colorbox{gray!8}{\parbox{\dimexpr\linewidth-2\fboxsep\relax}{\raggedright\ttfamily\scriptsize     values = []}}
\par\smallskip
This Python line participates in the export, validation, or report-generation pipeline. It is executed on the host PC to prepare data for RTL or to check RTL outputs.
\paragraph{Line 82}
\noindent\colorbox{gray!8}{\parbox{\dimexpr\linewidth-2\fboxsep\relax}{\raggedright\ttfamily\scriptsize     for line in path.read\_text(encoding="ascii").splitlines():}}
\par\smallskip
This reads an artifact produced by training, export, simulation, or Vivado. The validator/report depends on files generated by the actual run rather than hidden constants.
\paragraph{Line 83}
\noindent\colorbox{gray!8}{\parbox{\dimexpr\linewidth-2\fboxsep\relax}{\raggedright\ttfamily\scriptsize         token = line.strip()}}
\par\smallskip
This Python line participates in the export, validation, or report-generation pipeline. It is executed on the host PC to prepare data for RTL or to check RTL outputs.
\paragraph{Line 84}
\noindent\colorbox{gray!8}{\parbox{\dimexpr\linewidth-2\fboxsep\relax}{\raggedright\ttfamily\scriptsize         if token:}}
\par\smallskip
This Python line participates in the export, validation, or report-generation pipeline. It is executed on the host PC to prepare data for RTL or to check RTL outputs.
\paragraph{Line 85}
\noindent\colorbox{gray!8}{\parbox{\dimexpr\linewidth-2\fboxsep\relax}{\raggedright\ttfamily\scriptsize             values.append(signed\_hex\_to\_int(token, bits))}}
\par\smallskip
This Python line participates in the export, validation, or report-generation pipeline. It is executed on the host PC to prepare data for RTL or to check RTL outputs.
\paragraph{Line 86}
\noindent\colorbox{gray!8}{\parbox{\dimexpr\linewidth-2\fboxsep\relax}{\raggedright\ttfamily\scriptsize     return np.array(values, dtype=np.int64)}}
\par\smallskip
This performs vectorized numerical work. The fixed-point golden models use NumPy so score, margin, and update calculations are explicit and repeatable.
\paragraph{Line 87}
\noindent\colorbox{gray!8}{\parbox{\dimexpr\linewidth-2\fboxsep\relax}{\raggedright\ttfamily\scriptsize \emph{blank line}}}
\par\smallskip
This blank line is not hardware or software behavior. It separates nearby declarations or operations so the file remains readable and reviewable.
\paragraph{Line 88}
\noindent\colorbox{gray!8}{\parbox{\dimexpr\linewidth-2\fboxsep\relax}{\raggedright\ttfamily\scriptsize \emph{blank line}}}
\par\smallskip
This blank line is not hardware or software behavior. It separates nearby declarations or operations so the file remains readable and reviewable.
\paragraph{Line 89}
\noindent\colorbox{gray!8}{\parbox{\dimexpr\linewidth-2\fboxsep\relax}{\raggedright\ttfamily\scriptsize def read\_decimal\_mem(path: Path) -> np.ndarray:}}
\par\smallskip
This defines a Python function. The function packages one repeatable step such as loading a memory file, simulating the golden learning rule, validating CSVs, or generating plots.
\paragraph{Line 90}
\noindent\colorbox{gray!8}{\parbox{\dimexpr\linewidth-2\fboxsep\relax}{\raggedright\ttfamily\scriptsize     values = []}}
\par\smallskip
This Python line participates in the export, validation, or report-generation pipeline. It is executed on the host PC to prepare data for RTL or to check RTL outputs.
\paragraph{Line 91}
\noindent\colorbox{gray!8}{\parbox{\dimexpr\linewidth-2\fboxsep\relax}{\raggedright\ttfamily\scriptsize     for line in path.read\_text(encoding="ascii").splitlines():}}
\par\smallskip
This reads an artifact produced by training, export, simulation, or Vivado. The validator/report depends on files generated by the actual run rather than hidden constants.
\paragraph{Line 92}
\noindent\colorbox{gray!8}{\parbox{\dimexpr\linewidth-2\fboxsep\relax}{\raggedright\ttfamily\scriptsize         token = line.strip()}}
\par\smallskip
This Python line participates in the export, validation, or report-generation pipeline. It is executed on the host PC to prepare data for RTL or to check RTL outputs.
\paragraph{Line 93}
\noindent\colorbox{gray!8}{\parbox{\dimexpr\linewidth-2\fboxsep\relax}{\raggedright\ttfamily\scriptsize         if token:}}
\par\smallskip
This Python line participates in the export, validation, or report-generation pipeline. It is executed on the host PC to prepare data for RTL or to check RTL outputs.
\paragraph{Line 94}
\noindent\colorbox{gray!8}{\parbox{\dimexpr\linewidth-2\fboxsep\relax}{\raggedright\ttfamily\scriptsize             values.append(int(token))}}
\par\smallskip
This Python line participates in the export, validation, or report-generation pipeline. It is executed on the host PC to prepare data for RTL or to check RTL outputs.
\paragraph{Line 95}
\noindent\colorbox{gray!8}{\parbox{\dimexpr\linewidth-2\fboxsep\relax}{\raggedright\ttfamily\scriptsize     return np.array(values, dtype=np.int64)}}
\par\smallskip
This performs vectorized numerical work. The fixed-point golden models use NumPy so score, margin, and update calculations are explicit and repeatable.
\paragraph{Line 96}
\noindent\colorbox{gray!8}{\parbox{\dimexpr\linewidth-2\fboxsep\relax}{\raggedright\ttfamily\scriptsize \emph{blank line}}}
\par\smallskip
This blank line is not hardware or software behavior. It separates nearby declarations or operations so the file remains readable and reviewable.
\paragraph{Line 97}
\noindent\colorbox{gray!8}{\parbox{\dimexpr\linewidth-2\fboxsep\relax}{\raggedright\ttfamily\scriptsize \emph{blank line}}}
\par\smallskip
This blank line is not hardware or software behavior. It separates nearby declarations or operations so the file remains readable and reviewable.
\paragraph{Line 98}
\noindent\colorbox{gray!8}{\parbox{\dimexpr\linewidth-2\fboxsep\relax}{\raggedright\ttfamily\scriptsize def read\_json(path: Path) -> dict:}}
\par\smallskip
This defines a Python function. The function packages one repeatable step such as loading a memory file, simulating the golden learning rule, validating CSVs, or generating plots.
\paragraph{Line 99}
\noindent\colorbox{gray!8}{\parbox{\dimexpr\linewidth-2\fboxsep\relax}{\raggedright\ttfamily\scriptsize     return json.loads(path.read\_text(encoding="utf-8")) if path.exists() else \{\}}}
\par\smallskip
This reads an artifact produced by training, export, simulation, or Vivado. The validator/report depends on files generated by the actual run rather than hidden constants.
\paragraph{Line 100}
\noindent\colorbox{gray!8}{\parbox{\dimexpr\linewidth-2\fboxsep\relax}{\raggedright\ttfamily\scriptsize \emph{blank line}}}
\par\smallskip
This blank line is not hardware or software behavior. It separates nearby declarations or operations so the file remains readable and reviewable.
\paragraph{Line 101}
\noindent\colorbox{gray!8}{\parbox{\dimexpr\linewidth-2\fboxsep\relax}{\raggedright\ttfamily\scriptsize \emph{blank line}}}
\par\smallskip
This blank line is not hardware or software behavior. It separates nearby declarations or operations so the file remains readable and reviewable.
\paragraph{Line 102}
\noindent\colorbox{gray!8}{\parbox{\dimexpr\linewidth-2\fboxsep\relax}{\raggedright\ttfamily\scriptsize def parse\_first\_float(pattern: str, text: str, default: float = float("nan")) -> float:}}
\par\smallskip
This defines a Python function. The function packages one repeatable step such as loading a memory file, simulating the golden learning rule, validating CSVs, or generating plots.
\paragraph{Line 103}
\noindent\colorbox{gray!8}{\parbox{\dimexpr\linewidth-2\fboxsep\relax}{\raggedright\ttfamily\scriptsize     match = re.search(pattern, text, flags=re.IGNORECASE | re.MULTILINE)}}
\par\smallskip
This Python line participates in the export, validation, or report-generation pipeline. It is executed on the host PC to prepare data for RTL or to check RTL outputs.
\paragraph{Line 104}
\noindent\colorbox{gray!8}{\parbox{\dimexpr\linewidth-2\fboxsep\relax}{\raggedright\ttfamily\scriptsize     return float(match.group(1)) if match else default}}
\par\smallskip
This Python line participates in the export, validation, or report-generation pipeline. It is executed on the host PC to prepare data for RTL or to check RTL outputs.
\paragraph{Line 105}
\noindent\colorbox{gray!8}{\parbox{\dimexpr\linewidth-2\fboxsep\relax}{\raggedright\ttfamily\scriptsize \emph{blank line}}}
\par\smallskip
This blank line is not hardware or software behavior. It separates nearby declarations or operations so the file remains readable and reviewable.
\paragraph{Line 106}
\noindent\colorbox{gray!8}{\parbox{\dimexpr\linewidth-2\fboxsep\relax}{\raggedright\ttfamily\scriptsize \emph{blank line}}}
\par\smallskip
This blank line is not hardware or software behavior. It separates nearby declarations or operations so the file remains readable and reviewable.
\paragraph{Line 107}
\noindent\colorbox{gray!8}{\parbox{\dimexpr\linewidth-2\fboxsep\relax}{\raggedright\ttfamily\scriptsize def parse\_first\_int(pattern: str, text: str, default: int = 0) -> int:}}
\par\smallskip
This defines a Python function. The function packages one repeatable step such as loading a memory file, simulating the golden learning rule, validating CSVs, or generating plots.
\paragraph{Line 108}
\noindent\colorbox{gray!8}{\parbox{\dimexpr\linewidth-2\fboxsep\relax}{\raggedright\ttfamily\scriptsize     match = re.search(pattern, text, flags=re.IGNORECASE | re.MULTILINE)}}
\par\smallskip
This Python line participates in the export, validation, or report-generation pipeline. It is executed on the host PC to prepare data for RTL or to check RTL outputs.
\paragraph{Line 109}
\noindent\colorbox{gray!8}{\parbox{\dimexpr\linewidth-2\fboxsep\relax}{\raggedright\ttfamily\scriptsize     return int(match.group(1)) if match else default}}
\par\smallskip
This Python line participates in the export, validation, or report-generation pipeline. It is executed on the host PC to prepare data for RTL or to check RTL outputs.
\paragraph{Line 110}
\noindent\colorbox{gray!8}{\parbox{\dimexpr\linewidth-2\fboxsep\relax}{\raggedright\ttfamily\scriptsize \emph{blank line}}}
\par\smallskip
This blank line is not hardware or software behavior. It separates nearby declarations or operations so the file remains readable and reviewable.
\paragraph{Line 111}
\noindent\colorbox{gray!8}{\parbox{\dimexpr\linewidth-2\fboxsep\relax}{\raggedright\ttfamily\scriptsize \emph{blank line}}}
\par\smallskip
This blank line is not hardware or software behavior. It separates nearby declarations or operations so the file remains readable and reviewable.
\paragraph{Line 112}
\noindent\colorbox{gray!8}{\parbox{\dimexpr\linewidth-2\fboxsep\relax}{\raggedright\ttfamily\scriptsize def parse\_vivado\_reports(vivado\_dir: Path) -> dict[str, float | int | str]:}}
\par\smallskip
This defines a Python function. The function packages one repeatable step such as loading a memory file, simulating the golden learning rule, validating CSVs, or generating plots.
\paragraph{Line 113}
\noindent\colorbox{gray!8}{\parbox{\dimexpr\linewidth-2\fboxsep\relax}{\raggedright\ttfamily\scriptsize     timing = (vivado\_dir / "spikingformer\_head\_timing\_summary.rpt").read\_text(encoding="utf-8", errors="ignore")}}
\par\smallskip
This reads an artifact produced by training, export, simulation, or Vivado. The validator/report depends on files generated by the actual run rather than hidden constants.
\paragraph{Line 114}
\noindent\colorbox{gray!8}{\parbox{\dimexpr\linewidth-2\fboxsep\relax}{\raggedright\ttfamily\scriptsize     util = (vivado\_dir / "spikingformer\_head\_utilization.rpt").read\_text(encoding="utf-8", errors="ignore")}}
\par\smallskip
This reads an artifact produced by training, export, simulation, or Vivado. The validator/report depends on files generated by the actual run rather than hidden constants.
\paragraph{Line 115}
\noindent\colorbox{gray!8}{\parbox{\dimexpr\linewidth-2\fboxsep\relax}{\raggedright\ttfamily\scriptsize     power = (vivado\_dir / "spikingformer\_head\_power.rpt").read\_text(encoding="utf-8", errors="ignore")}}
\par\smallskip
This reads an artifact produced by training, export, simulation, or Vivado. The validator/report depends on files generated by the actual run rather than hidden constants.
\paragraph{Line 116}
\noindent\colorbox{gray!8}{\parbox{\dimexpr\linewidth-2\fboxsep\relax}{\raggedright\ttfamily\scriptsize     summary = (vivado\_dir / "spikingformer\_head\_synth\_summary.txt").read\_text(encoding="utf-8", errors="ignore")}}
\par\smallskip
This reads an artifact produced by training, export, simulation, or Vivado. The validator/report depends on files generated by the actual run rather than hidden constants.
\paragraph{Line 117}
\noindent\colorbox{gray!8}{\parbox{\dimexpr\linewidth-2\fboxsep\relax}{\raggedright\ttfamily\scriptsize     return \{}}
\par\smallskip
This Python line participates in the export, validation, or report-generation pipeline. It is executed on the host PC to prepare data for RTL or to check RTL outputs.
\paragraph{Line 118}
\noindent\colorbox{gray!8}{\parbox{\dimexpr\linewidth-2\fboxsep\relax}{\raggedright\ttfamily\scriptsize         "summary": summary,}}
\par\smallskip
This Python line participates in the export, validation, or report-generation pipeline. It is executed on the host PC to prepare data for RTL or to check RTL outputs.
\paragraph{Line 119}
\noindent\colorbox{gray!8}{\parbox{\dimexpr\linewidth-2\fboxsep\relax}{\raggedright\ttfamily\scriptsize         "wns\_ns": parse\_first\_float(r"\textbackslash{}n\textbackslash{}s*([+-]?\textbackslash{}d+\textbackslash{}.\textbackslash{}d+)\textbackslash{}s+\textbackslash{}d+\textbackslash{}.\textbackslash{}d+\textbackslash{}s+\textbackslash{}d+\textbackslash{}s+\textbackslash{}d+\textbackslash{}s+[+-]?\textbackslash{}d+\textbackslash{}.\textbackslash{}d+", timing),}}
\par\smallskip
This Python line participates in the export, validation, or report-generation pipeline. It is executed on the host PC to prepare data for RTL or to check RTL outputs.
\paragraph{Line 120}
\noindent\colorbox{gray!8}{\parbox{\dimexpr\linewidth-2\fboxsep\relax}{\raggedright\ttfamily\scriptsize         "whs\_ns": parse\_first\_float(r"\textbackslash{}n\textbackslash{}s*[+-]?\textbackslash{}d+\textbackslash{}.\textbackslash{}d+\textbackslash{}s+\textbackslash{}d+\textbackslash{}.\textbackslash{}d+\textbackslash{}s+\textbackslash{}d+\textbackslash{}s+\textbackslash{}d+\textbackslash{}s+([+-]?\textbackslash{}d+\textbackslash{}.\textbackslash{}d+)", timing),}}
\par\smallskip
This Python line participates in the export, validation, or report-generation pipeline. It is executed on the host PC to prepare data for RTL or to check RTL outputs.
\paragraph{Line 121}
\noindent\colorbox{gray!8}{\parbox{\dimexpr\linewidth-2\fboxsep\relax}{\raggedright\ttfamily\scriptsize         "lut": parse\_first\_int(r"\textbackslash{}|\textbackslash{}s*Slice LUTs\textbackslash{}s*\textbackslash{}|\textbackslash{}s*(\textbackslash{}d+)", util),}}
\par\smallskip
This Python line participates in the export, validation, or report-generation pipeline. It is executed on the host PC to prepare data for RTL or to check RTL outputs.
\paragraph{Line 122}
\noindent\colorbox{gray!8}{\parbox{\dimexpr\linewidth-2\fboxsep\relax}{\raggedright\ttfamily\scriptsize         "regs": parse\_first\_int(r"\textbackslash{}|\textbackslash{}s*Slice Registers\textbackslash{}s*\textbackslash{}|\textbackslash{}s*(\textbackslash{}d+)", util),}}
\par\smallskip
This Python line participates in the export, validation, or report-generation pipeline. It is executed on the host PC to prepare data for RTL or to check RTL outputs.
\paragraph{Line 123}
\noindent\colorbox{gray!8}{\parbox{\dimexpr\linewidth-2\fboxsep\relax}{\raggedright\ttfamily\scriptsize         "bram": parse\_first\_int(r"\textbackslash{}|\textbackslash{}s*Block RAM Tile\textbackslash{}s*\textbackslash{}|\textbackslash{}s*(\textbackslash{}d+)", util),}}
\par\smallskip
This Python line participates in the export, validation, or report-generation pipeline. It is executed on the host PC to prepare data for RTL or to check RTL outputs.
\paragraph{Line 124}
\noindent\colorbox{gray!8}{\parbox{\dimexpr\linewidth-2\fboxsep\relax}{\raggedright\ttfamily\scriptsize         "dsp": parse\_first\_int(r"\textbackslash{}|\textbackslash{}s*DSPs\textbackslash{}s*\textbackslash{}|\textbackslash{}s*(\textbackslash{}d+)", util),}}
\par\smallskip
This Python line participates in the export, validation, or report-generation pipeline. It is executed on the host PC to prepare data for RTL or to check RTL outputs.
\paragraph{Line 125}
\noindent\colorbox{gray!8}{\parbox{\dimexpr\linewidth-2\fboxsep\relax}{\raggedright\ttfamily\scriptsize         "lut\_mem": parse\_first\_int(r"\textbackslash{}|\textbackslash{}s*LUT as Memory\textbackslash{}s*\textbackslash{}|\textbackslash{}s*(\textbackslash{}d+)", util),}}
\par\smallskip
This Python line participates in the export, validation, or report-generation pipeline. It is executed on the host PC to prepare data for RTL or to check RTL outputs.
\paragraph{Line 126}
\noindent\colorbox{gray!8}{\parbox{\dimexpr\linewidth-2\fboxsep\relax}{\raggedright\ttfamily\scriptsize         "total\_power\_w": parse\_first\_float(r"\textbackslash{}|\textbackslash{}s*Total On-Chip Power \textbackslash{}(W\textbackslash{})\textbackslash{}s*\textbackslash{}|\textbackslash{}s*([0-9.]+)", power),}}
\par\smallskip
This Python line participates in the export, validation, or report-generation pipeline. It is executed on the host PC to prepare data for RTL or to check RTL outputs.
\paragraph{Line 127}
\noindent\colorbox{gray!8}{\parbox{\dimexpr\linewidth-2\fboxsep\relax}{\raggedright\ttfamily\scriptsize         "dynamic\_power\_w": parse\_first\_float(r"\textbackslash{}|\textbackslash{}s*Dynamic \textbackslash{}(W\textbackslash{})\textbackslash{}s*\textbackslash{}|\textbackslash{}s*([0-9.]+)", power),}}
\par\smallskip
This Python line participates in the export, validation, or report-generation pipeline. It is executed on the host PC to prepare data for RTL or to check RTL outputs.
\paragraph{Line 128}
\noindent\colorbox{gray!8}{\parbox{\dimexpr\linewidth-2\fboxsep\relax}{\raggedright\ttfamily\scriptsize         "static\_power\_w": parse\_first\_float(r"\textbackslash{}|\textbackslash{}s*Device Static \textbackslash{}(W\textbackslash{})\textbackslash{}s*\textbackslash{}|\textbackslash{}s*([0-9.]+)", power),}}
\par\smallskip
This Python line participates in the export, validation, or report-generation pipeline. It is executed on the host PC to prepare data for RTL or to check RTL outputs.
\paragraph{Line 129}
\noindent\colorbox{gray!8}{\parbox{\dimexpr\linewidth-2\fboxsep\relax}{\raggedright\ttfamily\scriptsize         "clock\_power\_w": parse\_first\_float(r"\textbackslash{}|\textbackslash{}s*Clocks\textbackslash{}s*\textbackslash{}|\textbackslash{}s*([0-9.<]+)", power.replace("<0.001", "0.0005")),}}
\par\smallskip
This Python line participates in the export, validation, or report-generation pipeline. It is executed on the host PC to prepare data for RTL or to check RTL outputs.
\paragraph{Line 130}
\noindent\colorbox{gray!8}{\parbox{\dimexpr\linewidth-2\fboxsep\relax}{\raggedright\ttfamily\scriptsize         "logic\_power\_w": parse\_first\_float(r"\textbackslash{}|\textbackslash{}s*Slice Logic\textbackslash{}s*\textbackslash{}|\textbackslash{}s*([0-9.<]+)", power.replace("<0.001", "0.0005")),}}
\par\smallskip
This Python line participates in the export, validation, or report-generation pipeline. It is executed on the host PC to prepare data for RTL or to check RTL outputs.
\paragraph{Line 131}
\noindent\colorbox{gray!8}{\parbox{\dimexpr\linewidth-2\fboxsep\relax}{\raggedright\ttfamily\scriptsize         "signal\_power\_w": parse\_first\_float(r"\textbackslash{}|\textbackslash{}s*Signals\textbackslash{}s*\textbackslash{}|\textbackslash{}s*([0-9.<]+)", power.replace("<0.001", "0.0005")),}}
\par\smallskip
This Python line participates in the export, validation, or report-generation pipeline. It is executed on the host PC to prepare data for RTL or to check RTL outputs.
\paragraph{Line 132}
\noindent\colorbox{gray!8}{\parbox{\dimexpr\linewidth-2\fboxsep\relax}{\raggedright\ttfamily\scriptsize         "bram\_power\_w": parse\_first\_float(r"\textbackslash{}|\textbackslash{}s*Block RAM\textbackslash{}s*\textbackslash{}|\textbackslash{}s*([0-9.<]+)", power.replace("<0.001", "0.0005")),}}
\par\smallskip
This Python line participates in the export, validation, or report-generation pipeline. It is executed on the host PC to prepare data for RTL or to check RTL outputs.
\paragraph{Line 133}
\noindent\colorbox{gray!8}{\parbox{\dimexpr\linewidth-2\fboxsep\relax}{\raggedright\ttfamily\scriptsize         "dsp\_power\_w": parse\_first\_float(r"\textbackslash{}|\textbackslash{}s*DSPs\textbackslash{}s*\textbackslash{}|\textbackslash{}s*([0-9.<]+)", power.replace("<0.001", "0.0005")),}}
\par\smallskip
This Python line participates in the export, validation, or report-generation pipeline. It is executed on the host PC to prepare data for RTL or to check RTL outputs.
\paragraph{Line 134}
\noindent\colorbox{gray!8}{\parbox{\dimexpr\linewidth-2\fboxsep\relax}{\raggedright\ttfamily\scriptsize         "io\_power\_w": parse\_first\_float(r"\textbackslash{}|\textbackslash{}s*I/O\textbackslash{}s*\textbackslash{}|\textbackslash{}s*([0-9.<]+)", power.replace("<0.001", "0.0005")),}}
\par\smallskip
This Python line participates in the export, validation, or report-generation pipeline. It is executed on the host PC to prepare data for RTL or to check RTL outputs.
\paragraph{Line 135}
\noindent\colorbox{gray!8}{\parbox{\dimexpr\linewidth-2\fboxsep\relax}{\raggedright\ttfamily\scriptsize     \}}}
\par\smallskip
This Python line participates in the export, validation, or report-generation pipeline. It is executed on the host PC to prepare data for RTL or to check RTL outputs.
\paragraph{Line 136}
\noindent\colorbox{gray!8}{\parbox{\dimexpr\linewidth-2\fboxsep\relax}{\raggedright\ttfamily\scriptsize \emph{blank line}}}
\par\smallskip
This blank line is not hardware or software behavior. It separates nearby declarations or operations so the file remains readable and reviewable.
\paragraph{Line 137}
\noindent\colorbox{gray!8}{\parbox{\dimexpr\linewidth-2\fboxsep\relax}{\raggedright\ttfamily\scriptsize \emph{blank line}}}
\par\smallskip
This blank line is not hardware or software behavior. It separates nearby declarations or operations so the file remains readable and reviewable.
\paragraph{Line 138}
\noindent\colorbox{gray!8}{\parbox{\dimexpr\linewidth-2\fboxsep\relax}{\raggedright\ttfamily\scriptsize def save\_pdf\_page(pdf: PdfPages, fig: plt.Figure, png\_dir: Path, name: str) -> None:}}
\par\smallskip
This defines a Python function. The function packages one repeatable step such as loading a memory file, simulating the golden learning rule, validating CSVs, or generating plots.
\paragraph{Line 139}
\noindent\colorbox{gray!8}{\parbox{\dimexpr\linewidth-2\fboxsep\relax}{\raggedright\ttfamily\scriptsize     png\_dir.mkdir(parents=True, exist\_ok=True)}}
\par\smallskip
This Python line participates in the export, validation, or report-generation pipeline. It is executed on the host PC to prepare data for RTL or to check RTL outputs.
\paragraph{Line 140}
\noindent\colorbox{gray!8}{\parbox{\dimexpr\linewidth-2\fboxsep\relax}{\raggedright\ttfamily\scriptsize     fig.tight\_layout()}}
\par\smallskip
This Python line participates in the export, validation, or report-generation pipeline. It is executed on the host PC to prepare data for RTL or to check RTL outputs.
\paragraph{Line 141}
\noindent\colorbox{gray!8}{\parbox{\dimexpr\linewidth-2\fboxsep\relax}{\raggedright\ttfamily\scriptsize     fig.savefig(png\_dir / f"\{name\}.png", dpi=180, bbox\_inches="tight")}}
\par\smallskip
This Python line participates in the export, validation, or report-generation pipeline. It is executed on the host PC to prepare data for RTL or to check RTL outputs.
\paragraph{Line 142}
\noindent\colorbox{gray!8}{\parbox{\dimexpr\linewidth-2\fboxsep\relax}{\raggedright\ttfamily\scriptsize     pdf.savefig(fig, bbox\_inches="tight")}}
\par\smallskip
This Python line participates in the export, validation, or report-generation pipeline. It is executed on the host PC to prepare data for RTL or to check RTL outputs.
\paragraph{Line 143}
\noindent\colorbox{gray!8}{\parbox{\dimexpr\linewidth-2\fboxsep\relax}{\raggedright\ttfamily\scriptsize     plt.close(fig)}}
\par\smallskip
This creates a debug visualization. The plot is derived from generated CSV or Vivado report data so it can be regenerated after RTL changes.
\paragraph{Line 144}
\noindent\colorbox{gray!8}{\parbox{\dimexpr\linewidth-2\fboxsep\relax}{\raggedright\ttfamily\scriptsize \emph{blank line}}}
\par\smallskip
This blank line is not hardware or software behavior. It separates nearby declarations or operations so the file remains readable and reviewable.
\paragraph{Line 145}
\noindent\colorbox{gray!8}{\parbox{\dimexpr\linewidth-2\fboxsep\relax}{\raggedright\ttfamily\scriptsize \emph{blank line}}}
\par\smallskip
This blank line is not hardware or software behavior. It separates nearby declarations or operations so the file remains readable and reviewable.
\paragraph{Line 146}
\noindent\colorbox{gray!8}{\parbox{\dimexpr\linewidth-2\fboxsep\relax}{\raggedright\ttfamily\scriptsize def confusion(labels: np.ndarray, preds: np.ndarray) -> np.ndarray:}}
\par\smallskip
This defines a Python function. The function packages one repeatable step such as loading a memory file, simulating the golden learning rule, validating CSVs, or generating plots.
\paragraph{Line 147}
\noindent\colorbox{gray!8}{\parbox{\dimexpr\linewidth-2\fboxsep\relax}{\raggedright\ttfamily\scriptsize     mat = np.zeros((NUM\_CLASSES, NUM\_CLASSES), dtype=np.int64)}}
\par\smallskip
This performs vectorized numerical work. The fixed-point golden models use NumPy so score, margin, and update calculations are explicit and repeatable.
\paragraph{Line 148}
\noindent\colorbox{gray!8}{\parbox{\dimexpr\linewidth-2\fboxsep\relax}{\raggedright\ttfamily\scriptsize     for label, pred in zip(labels, preds):}}
\par\smallskip
This line handles predicted labels. It checks that argmax behavior in Python and RTL chooses the same class.
\paragraph{Line 149}
\noindent\colorbox{gray!8}{\parbox{\dimexpr\linewidth-2\fboxsep\relax}{\raggedright\ttfamily\scriptsize         mat[int(label), int(pred)] += 1}}
\par\smallskip
This line handles predicted labels. It checks that argmax behavior in Python and RTL chooses the same class.
\paragraph{Line 150}
\noindent\colorbox{gray!8}{\parbox{\dimexpr\linewidth-2\fboxsep\relax}{\raggedright\ttfamily\scriptsize     return mat}}
\par\smallskip
This Python line participates in the export, validation, or report-generation pipeline. It is executed on the host PC to prepare data for RTL or to check RTL outputs.
\paragraph{Line 151}
\noindent\colorbox{gray!8}{\parbox{\dimexpr\linewidth-2\fboxsep\relax}{\raggedright\ttfamily\scriptsize \emph{blank line}}}
\par\smallskip
This blank line is not hardware or software behavior. It separates nearby declarations or operations so the file remains readable and reviewable.
\paragraph{Line 152}
\noindent\colorbox{gray!8}{\parbox{\dimexpr\linewidth-2\fboxsep\relax}{\raggedright\ttfamily\scriptsize \emph{blank line}}}
\par\smallskip
This blank line is not hardware or software behavior. It separates nearby declarations or operations so the file remains readable and reviewable.
\paragraph{Line 153}
\noindent\colorbox{gray!8}{\parbox{\dimexpr\linewidth-2\fboxsep\relax}{\raggedright\ttfamily\scriptsize def precision\_recall\_f1(conf: np.ndarray) -> tuple[np.ndarray, np.ndarray, np.ndarray]:}}
\par\smallskip
This defines a Python function. The function packages one repeatable step such as loading a memory file, simulating the golden learning rule, validating CSVs, or generating plots.
\paragraph{Line 154}
\noindent\colorbox{gray!8}{\parbox{\dimexpr\linewidth-2\fboxsep\relax}{\raggedright\ttfamily\scriptsize     tp = np.diag(conf).astype(float)}}
\par\smallskip
This performs vectorized numerical work. The fixed-point golden models use NumPy so score, margin, and update calculations are explicit and repeatable.
\paragraph{Line 155}
\noindent\colorbox{gray!8}{\parbox{\dimexpr\linewidth-2\fboxsep\relax}{\raggedright\ttfamily\scriptsize     precision = tp / np.maximum(conf.sum(axis=0), 1)}}
\par\smallskip
This performs vectorized numerical work. The fixed-point golden models use NumPy so score, margin, and update calculations are explicit and repeatable.
\paragraph{Line 156}
\noindent\colorbox{gray!8}{\parbox{\dimexpr\linewidth-2\fboxsep\relax}{\raggedright\ttfamily\scriptsize     recall = tp / np.maximum(conf.sum(axis=1), 1)}}
\par\smallskip
This performs vectorized numerical work. The fixed-point golden models use NumPy so score, margin, and update calculations are explicit and repeatable.
\paragraph{Line 157}
\noindent\colorbox{gray!8}{\parbox{\dimexpr\linewidth-2\fboxsep\relax}{\raggedright\ttfamily\scriptsize     f1 = 2 * precision * recall / np.maximum(precision + recall, 1e-12)}}
\par\smallskip
This performs vectorized numerical work. The fixed-point golden models use NumPy so score, margin, and update calculations are explicit and repeatable.
\paragraph{Line 158}
\noindent\colorbox{gray!8}{\parbox{\dimexpr\linewidth-2\fboxsep\relax}{\raggedright\ttfamily\scriptsize     return precision, recall, f1}}
\par\smallskip
This Python line participates in the export, validation, or report-generation pipeline. It is executed on the host PC to prepare data for RTL or to check RTL outputs.
\paragraph{Line 159}
\noindent\colorbox{gray!8}{\parbox{\dimexpr\linewidth-2\fboxsep\relax}{\raggedright\ttfamily\scriptsize \emph{blank line}}}
\par\smallskip
This blank line is not hardware or software behavior. It separates nearby declarations or operations so the file remains readable and reviewable.
\paragraph{Line 160}
\noindent\colorbox{gray!8}{\parbox{\dimexpr\linewidth-2\fboxsep\relax}{\raggedright\ttfamily\scriptsize \emph{blank line}}}
\par\smallskip
This blank line is not hardware or software behavior. It separates nearby declarations or operations so the file remains readable and reviewable.
\paragraph{Line 161}
\noindent\colorbox{gray!8}{\parbox{\dimexpr\linewidth-2\fboxsep\relax}{\raggedright\ttfamily\scriptsize def margin(scores: np.ndarray, labels: np.ndarray) -> tuple[np.ndarray, np.ndarray]:}}
\par\smallskip
This defines a Python function. The function packages one repeatable step such as loading a memory file, simulating the golden learning rule, validating CSVs, or generating plots.
\paragraph{Line 162}
\noindent\colorbox{gray!8}{\parbox{\dimexpr\linewidth-2\fboxsep\relax}{\raggedright\ttfamily\scriptsize     sorted\_scores = np.sort(scores, axis=1)}}
\par\smallskip
This performs vectorized numerical work. The fixed-point golden models use NumPy so score, margin, and update calculations are explicit and repeatable.
\paragraph{Line 163}
\noindent\colorbox{gray!8}{\parbox{\dimexpr\linewidth-2\fboxsep\relax}{\raggedright\ttfamily\scriptsize     top1 = sorted\_scores[:, -1]}}
\par\smallskip
This line handles class scores. The score vector is the direct software equivalent of the RTL MAC loop over weights and features.
\paragraph{Line 164}
\noindent\colorbox{gray!8}{\parbox{\dimexpr\linewidth-2\fboxsep\relax}{\raggedright\ttfamily\scriptsize     top2 = sorted\_scores[:, -2]}}
\par\smallskip
This line handles class scores. The score vector is the direct software equivalent of the RTL MAC loop over weights and features.
\paragraph{Line 165}
\noindent\colorbox{gray!8}{\parbox{\dimexpr\linewidth-2\fboxsep\relax}{\raggedright\ttfamily\scriptsize     true\_scores = scores[np.arange(scores.shape[0]), labels.astype(int)]}}
\par\smallskip
This performs vectorized numerical work. The fixed-point golden models use NumPy so score, margin, and update calculations are explicit and repeatable.
\paragraph{Line 166}
\noindent\colorbox{gray!8}{\parbox{\dimexpr\linewidth-2\fboxsep\relax}{\raggedright\ttfamily\scriptsize     return top1 - top2, top1 - true\_scores}}
\par\smallskip
This line handles class scores. The score vector is the direct software equivalent of the RTL MAC loop over weights and features.
\paragraph{Line 167}
\noindent\colorbox{gray!8}{\parbox{\dimexpr\linewidth-2\fboxsep\relax}{\raggedright\ttfamily\scriptsize \emph{blank line}}}
\par\smallskip
This blank line is not hardware or software behavior. It separates nearby declarations or operations so the file remains readable and reviewable.
\paragraph{Line 168}
\noindent\colorbox{gray!8}{\parbox{\dimexpr\linewidth-2\fboxsep\relax}{\raggedright\ttfamily\scriptsize \emph{blank line}}}
\par\smallskip
This blank line is not hardware or software behavior. It separates nearby declarations or operations so the file remains readable and reviewable.
\paragraph{Line 169}
\noindent\colorbox{gray!8}{\parbox{\dimexpr\linewidth-2\fboxsep\relax}{\raggedright\ttfamily\scriptsize def plot\_text\_page(title: str, lines: Iterable[str]) -> plt.Figure:}}
\par\smallskip
This defines a Python function. The function packages one repeatable step such as loading a memory file, simulating the golden learning rule, validating CSVs, or generating plots.
\paragraph{Line 170}
\noindent\colorbox{gray!8}{\parbox{\dimexpr\linewidth-2\fboxsep\relax}{\raggedright\ttfamily\scriptsize     fig = plt.figure(figsize=(11, 8.5))}}
\par\smallskip
This creates a debug visualization. The plot is derived from generated CSV or Vivado report data so it can be regenerated after RTL changes.
\paragraph{Line 171}
\noindent\colorbox{gray!8}{\parbox{\dimexpr\linewidth-2\fboxsep\relax}{\raggedright\ttfamily\scriptsize     fig.text(0.04, 0.95, title, fontsize=18, weight="bold", va="top")}}
\par\smallskip
This Python line participates in the export, validation, or report-generation pipeline. It is executed on the host PC to prepare data for RTL or to check RTL outputs.
\paragraph{Line 172}
\noindent\colorbox{gray!8}{\parbox{\dimexpr\linewidth-2\fboxsep\relax}{\raggedright\ttfamily\scriptsize     y = 0.90}}
\par\smallskip
This Python line participates in the export, validation, or report-generation pipeline. It is executed on the host PC to prepare data for RTL or to check RTL outputs.
\paragraph{Line 173}
\noindent\colorbox{gray!8}{\parbox{\dimexpr\linewidth-2\fboxsep\relax}{\raggedright\ttfamily\scriptsize     for line in lines:}}
\par\smallskip
This Python line participates in the export, validation, or report-generation pipeline. It is executed on the host PC to prepare data for RTL or to check RTL outputs.
\paragraph{Line 174}
\noindent\colorbox{gray!8}{\parbox{\dimexpr\linewidth-2\fboxsep\relax}{\raggedright\ttfamily\scriptsize         fig.text(0.05, y, line, fontsize=10.5, family="monospace", va="top")}}
\par\smallskip
This Python line participates in the export, validation, or report-generation pipeline. It is executed on the host PC to prepare data for RTL or to check RTL outputs.
\paragraph{Line 175}
\noindent\colorbox{gray!8}{\parbox{\dimexpr\linewidth-2\fboxsep\relax}{\raggedright\ttfamily\scriptsize         y -= 0.035}}
\par\smallskip
This Python line participates in the export, validation, or report-generation pipeline. It is executed on the host PC to prepare data for RTL or to check RTL outputs.
\paragraph{Line 176}
\noindent\colorbox{gray!8}{\parbox{\dimexpr\linewidth-2\fboxsep\relax}{\raggedright\ttfamily\scriptsize         if y < 0.05:}}
\par\smallskip
This Python line participates in the export, validation, or report-generation pipeline. It is executed on the host PC to prepare data for RTL or to check RTL outputs.
\paragraph{Line 177}
\noindent\colorbox{gray!8}{\parbox{\dimexpr\linewidth-2\fboxsep\relax}{\raggedright\ttfamily\scriptsize             break}}
\par\smallskip
This Python line participates in the export, validation, or report-generation pipeline. It is executed on the host PC to prepare data for RTL or to check RTL outputs.
\paragraph{Line 178}
\noindent\colorbox{gray!8}{\parbox{\dimexpr\linewidth-2\fboxsep\relax}{\raggedright\ttfamily\scriptsize     return fig}}
\par\smallskip
This Python line participates in the export, validation, or report-generation pipeline. It is executed on the host PC to prepare data for RTL or to check RTL outputs.
\paragraph{Line 179}
\noindent\colorbox{gray!8}{\parbox{\dimexpr\linewidth-2\fboxsep\relax}{\raggedright\ttfamily\scriptsize \emph{blank line}}}
\par\smallskip
This blank line is not hardware or software behavior. It separates nearby declarations or operations so the file remains readable and reviewable.
\paragraph{Line 180}
\noindent\colorbox{gray!8}{\parbox{\dimexpr\linewidth-2\fboxsep\relax}{\raggedright\ttfamily\scriptsize \emph{blank line}}}
\par\smallskip
This blank line is not hardware or software behavior. It separates nearby declarations or operations so the file remains readable and reviewable.
\paragraph{Line 181}
\noindent\colorbox{gray!8}{\parbox{\dimexpr\linewidth-2\fboxsep\relax}{\raggedright\ttfamily\scriptsize def draw\_box(ax, xy, w, h, text, color="\#e8f1ff"):}}
\par\smallskip
This defines a Python function. The function packages one repeatable step such as loading a memory file, simulating the golden learning rule, validating CSVs, or generating plots.
\paragraph{Line 182}
\noindent\colorbox{gray!8}{\parbox{\dimexpr\linewidth-2\fboxsep\relax}{\raggedright\ttfamily\scriptsize     rect = plt.Rectangle(xy, w, h, facecolor=color, edgecolor="\#1d3557", linewidth=1.6)}}
\par\smallskip
This creates a debug visualization. The plot is derived from generated CSV or Vivado report data so it can be regenerated after RTL changes.
\paragraph{Line 183}
\noindent\colorbox{gray!8}{\parbox{\dimexpr\linewidth-2\fboxsep\relax}{\raggedright\ttfamily\scriptsize     ax.add\_patch(rect)}}
\par\smallskip
This Python line participates in the export, validation, or report-generation pipeline. It is executed on the host PC to prepare data for RTL or to check RTL outputs.
\paragraph{Line 184}
\noindent\colorbox{gray!8}{\parbox{\dimexpr\linewidth-2\fboxsep\relax}{\raggedright\ttfamily\scriptsize     ax.text(xy[0] + w / 2, xy[1] + h / 2, text, ha="center", va="center", fontsize=9, wrap=True)}}
\par\smallskip
This Python line participates in the export, validation, or report-generation pipeline. It is executed on the host PC to prepare data for RTL or to check RTL outputs.
\paragraph{Line 185}
\noindent\colorbox{gray!8}{\parbox{\dimexpr\linewidth-2\fboxsep\relax}{\raggedright\ttfamily\scriptsize \emph{blank line}}}
\par\smallskip
This blank line is not hardware or software behavior. It separates nearby declarations or operations so the file remains readable and reviewable.
\paragraph{Line 186}
\noindent\colorbox{gray!8}{\parbox{\dimexpr\linewidth-2\fboxsep\relax}{\raggedright\ttfamily\scriptsize \emph{blank line}}}
\par\smallskip
This blank line is not hardware or software behavior. It separates nearby declarations or operations so the file remains readable and reviewable.
\paragraph{Line 187}
\noindent\colorbox{gray!8}{\parbox{\dimexpr\linewidth-2\fboxsep\relax}{\raggedright\ttfamily\scriptsize def plot\_architecture() -> plt.Figure:}}
\par\smallskip
This defines a Python function. The function packages one repeatable step such as loading a memory file, simulating the golden learning rule, validating CSVs, or generating plots.
\paragraph{Line 188}
\noindent\colorbox{gray!8}{\parbox{\dimexpr\linewidth-2\fboxsep\relax}{\raggedright\ttfamily\scriptsize     fig, ax = plt.subplots(figsize=(11, 8.5))}}
\par\smallskip
This creates a debug visualization. The plot is derived from generated CSV or Vivado report data so it can be regenerated after RTL changes.
\paragraph{Line 189}
\noindent\colorbox{gray!8}{\parbox{\dimexpr\linewidth-2\fboxsep\relax}{\raggedright\ttfamily\scriptsize     ax.set\_axis\_off()}}
\par\smallskip
This Python line participates in the export, validation, or report-generation pipeline. It is executed on the host PC to prepare data for RTL or to check RTL outputs.
\paragraph{Line 190}
\noindent\colorbox{gray!8}{\parbox{\dimexpr\linewidth-2\fboxsep\relax}{\raggedright\ttfamily\scriptsize     ax.set\_xlim(0, 10)}}
\par\smallskip
This Python line participates in the export, validation, or report-generation pipeline. It is executed on the host PC to prepare data for RTL or to check RTL outputs.
\paragraph{Line 191}
\noindent\colorbox{gray!8}{\parbox{\dimexpr\linewidth-2\fboxsep\relax}{\raggedright\ttfamily\scriptsize     ax.set\_ylim(0, 7)}}
\par\smallskip
This Python line participates in the export, validation, or report-generation pipeline. It is executed on the host PC to prepare data for RTL or to check RTL outputs.
\paragraph{Line 192}
\noindent\colorbox{gray!8}{\parbox{\dimexpr\linewidth-2\fboxsep\relax}{\raggedright\ttfamily\scriptsize     blocks = [}}
\par\smallskip
This Python line participates in the export, validation, or report-generation pipeline. It is executed on the host PC to prepare data for RTL or to check RTL outputs.
\paragraph{Line 193}
\noindent\colorbox{gray!8}{\parbox{\dimexpr\linewidth-2\fboxsep\relax}{\raggedright\ttfamily\scriptsize         ((0.4, 5.4), 1.6, 0.7, "PYNQ PS / DMA\textbackslash{}nAXI-stream features", "\#eef6e8"),}}
\par\smallskip
This Python line participates in the export, validation, or report-generation pipeline. It is executed on the host PC to prepare data for RTL or to check RTL outputs.
\paragraph{Line 194}
\noindent\colorbox{gray!8}{\parbox{\dimexpr\linewidth-2\fboxsep\relax}{\raggedright\ttfamily\scriptsize         ((2.5, 5.4), 1.7, 0.7, "AXI-stream wrapper\textbackslash{}nstart, TLAST, result pack", "\#e8f1ff"),}}
\par\smallskip
This Python line participates in the export, validation, or report-generation pipeline. It is executed on the host PC to prepare data for RTL or to check RTL outputs.
\paragraph{Line 195}
\noindent\colorbox{gray!8}{\parbox{\dimexpr\linewidth-2\fboxsep\relax}{\raggedright\ttfamily\scriptsize         ((4.8, 5.4), 1.6, 0.7, "Feature SRAM\textbackslash{}n384 x int16", "\#fff0d9"),}}
\par\smallskip
This Python line participates in the export, validation, or report-generation pipeline. It is executed on the host PC to prepare data for RTL or to check RTL outputs.
\paragraph{Line 196}
\noindent\colorbox{gray!8}{\parbox{\dimexpr\linewidth-2\fboxsep\relax}{\raggedright\ttfamily\scriptsize         ((7.0, 5.4), 1.8, 0.7, "Score output\textbackslash{}nlabel + update status", "\#eef6e8"),}}
\par\smallskip
This line handles the online-learning rule or its validation fields. It is central to proving that the FPGA update path is not bypassed.
\paragraph{Line 197}
\noindent\colorbox{gray!8}{\parbox{\dimexpr\linewidth-2\fboxsep\relax}{\raggedright\ttfamily\scriptsize         ((1.0, 3.4), 1.7, 0.7, "Class loop\textbackslash{}n10 classes", "\#f6e8ff"),}}
\par\smallskip
This Python line participates in the export, validation, or report-generation pipeline. It is executed on the host PC to prepare data for RTL or to check RTL outputs.
\paragraph{Line 198}
\noindent\colorbox{gray!8}{\parbox{\dimexpr\linewidth-2\fboxsep\relax}{\raggedright\ttfamily\scriptsize         ((3.2, 3.4), 1.7, 0.7, "Weight BRAM\textbackslash{}n10 x 384 x int16", "\#fff0d9"),}}
\par\smallskip
This Python line participates in the export, validation, or report-generation pipeline. It is executed on the host PC to prepare data for RTL or to check RTL outputs.
\paragraph{Line 199}
\noindent\colorbox{gray!8}{\parbox{\dimexpr\linewidth-2\fboxsep\relax}{\raggedright\ttfamily\scriptsize         ((5.4, 3.4), 1.7, 0.7, "DSP MAC\textbackslash{}nint16 x int16 -> acc64", "\#e8f1ff"),}}
\par\smallskip
This Python line participates in the export, validation, or report-generation pipeline. It is executed on the host PC to prepare data for RTL or to check RTL outputs.
\paragraph{Line 200}
\noindent\colorbox{gray!8}{\parbox{\dimexpr\linewidth-2\fboxsep\relax}{\raggedright\ttfamily\scriptsize         ((7.6, 3.4), 1.7, 0.7, "Argmax score\textbackslash{}nprediction", "\#e8f1ff"),}}
\par\smallskip
This line handles predicted labels. It checks that argmax behavior in Python and RTL chooses the same class.
\paragraph{Line 201}
\noindent\colorbox{gray!8}{\parbox{\dimexpr\linewidth-2\fboxsep\relax}{\raggedright\ttfamily\scriptsize         ((2.1, 1.3), 1.8, 0.7, "Margin-gated\textbackslash{}non-chip update", "\#ffe8e8"),}}
\par\smallskip
This line handles the online-learning rule or its validation fields. It is central to proving that the FPGA update path is not bypassed.
\paragraph{Line 202}
\noindent\colorbox{gray!8}{\parbox{\dimexpr\linewidth-2\fboxsep\relax}{\raggedright\ttfamily\scriptsize         ((4.5, 1.3), 1.8, 0.7, "Registered\textbackslash{}nsat add/sub", "\#ffe8e8"),}}
\par\smallskip
This Python line participates in the export, validation, or report-generation pipeline. It is executed on the host PC to prepare data for RTL or to check RTL outputs.
\paragraph{Line 203}
\noindent\colorbox{gray!8}{\parbox{\dimexpr\linewidth-2\fboxsep\relax}{\raggedright\ttfamily\scriptsize         ((6.9, 1.3), 1.8, 0.7, "Writeback\textbackslash{}ntrue/pred rows + bias", "\#ffe8e8"),}}
\par\smallskip
This line handles the classifier parameters. Those arrays mirror the RTL memories, so mismatches reveal an export, quantization, or hardware-update bug.
\paragraph{Line 204}
\noindent\colorbox{gray!8}{\parbox{\dimexpr\linewidth-2\fboxsep\relax}{\raggedright\ttfamily\scriptsize     ]}}
\par\smallskip
This Python line participates in the export, validation, or report-generation pipeline. It is executed on the host PC to prepare data for RTL or to check RTL outputs.
\paragraph{Line 205}
\noindent\colorbox{gray!8}{\parbox{\dimexpr\linewidth-2\fboxsep\relax}{\raggedright\ttfamily\scriptsize     for args in blocks:}}
\par\smallskip
This Python line participates in the export, validation, or report-generation pipeline. It is executed on the host PC to prepare data for RTL or to check RTL outputs.
\paragraph{Line 206}
\noindent\colorbox{gray!8}{\parbox{\dimexpr\linewidth-2\fboxsep\relax}{\raggedright\ttfamily\scriptsize         draw\_box(ax, *args)}}
\par\smallskip
This Python line participates in the export, validation, or report-generation pipeline. It is executed on the host PC to prepare data for RTL or to check RTL outputs.
\paragraph{Line 207}
\noindent\colorbox{gray!8}{\parbox{\dimexpr\linewidth-2\fboxsep\relax}{\raggedright\ttfamily\scriptsize     arrows = [}}
\par\smallskip
This Python line participates in the export, validation, or report-generation pipeline. It is executed on the host PC to prepare data for RTL or to check RTL outputs.
\paragraph{Line 208}
\noindent\colorbox{gray!8}{\parbox{\dimexpr\linewidth-2\fboxsep\relax}{\raggedright\ttfamily\scriptsize         ((2.0, 5.75), (2.5, 5.75)),}}
\par\smallskip
This Python line participates in the export, validation, or report-generation pipeline. It is executed on the host PC to prepare data for RTL or to check RTL outputs.
\paragraph{Line 209}
\noindent\colorbox{gray!8}{\parbox{\dimexpr\linewidth-2\fboxsep\relax}{\raggedright\ttfamily\scriptsize         ((4.2, 5.75), (4.8, 5.75)),}}
\par\smallskip
This Python line participates in the export, validation, or report-generation pipeline. It is executed on the host PC to prepare data for RTL or to check RTL outputs.
\paragraph{Line 210}
\noindent\colorbox{gray!8}{\parbox{\dimexpr\linewidth-2\fboxsep\relax}{\raggedright\ttfamily\scriptsize         ((6.4, 5.75), (7.0, 5.75)),}}
\par\smallskip
This Python line participates in the export, validation, or report-generation pipeline. It is executed on the host PC to prepare data for RTL or to check RTL outputs.
\paragraph{Line 211}
\noindent\colorbox{gray!8}{\parbox{\dimexpr\linewidth-2\fboxsep\relax}{\raggedright\ttfamily\scriptsize         ((4.0, 3.75), (5.4, 3.75)),}}
\par\smallskip
This Python line participates in the export, validation, or report-generation pipeline. It is executed on the host PC to prepare data for RTL or to check RTL outputs.
\paragraph{Line 212}
\noindent\colorbox{gray!8}{\parbox{\dimexpr\linewidth-2\fboxsep\relax}{\raggedright\ttfamily\scriptsize         ((7.1, 3.75), (7.6, 3.75)),}}
\par\smallskip
This Python line participates in the export, validation, or report-generation pipeline. It is executed on the host PC to prepare data for RTL or to check RTL outputs.
\paragraph{Line 213}
\noindent\colorbox{gray!8}{\parbox{\dimexpr\linewidth-2\fboxsep\relax}{\raggedright\ttfamily\scriptsize         ((7.6, 3.35), (3.0, 2.0)),}}
\par\smallskip
This Python line participates in the export, validation, or report-generation pipeline. It is executed on the host PC to prepare data for RTL or to check RTL outputs.
\paragraph{Line 214}
\noindent\colorbox{gray!8}{\parbox{\dimexpr\linewidth-2\fboxsep\relax}{\raggedright\ttfamily\scriptsize         ((3.9, 1.65), (4.5, 1.65)),}}
\par\smallskip
This Python line participates in the export, validation, or report-generation pipeline. It is executed on the host PC to prepare data for RTL or to check RTL outputs.
\paragraph{Line 215}
\noindent\colorbox{gray!8}{\parbox{\dimexpr\linewidth-2\fboxsep\relax}{\raggedright\ttfamily\scriptsize         ((6.3, 1.65), (6.9, 1.65)),}}
\par\smallskip
This Python line participates in the export, validation, or report-generation pipeline. It is executed on the host PC to prepare data for RTL or to check RTL outputs.
\paragraph{Line 216}
\noindent\colorbox{gray!8}{\parbox{\dimexpr\linewidth-2\fboxsep\relax}{\raggedright\ttfamily\scriptsize         ((7.8, 2.0), (4.1, 3.35)),}}
\par\smallskip
This Python line participates in the export, validation, or report-generation pipeline. It is executed on the host PC to prepare data for RTL or to check RTL outputs.
\paragraph{Line 217}
\noindent\colorbox{gray!8}{\parbox{\dimexpr\linewidth-2\fboxsep\relax}{\raggedright\ttfamily\scriptsize     ]}}
\par\smallskip
This Python line participates in the export, validation, or report-generation pipeline. It is executed on the host PC to prepare data for RTL or to check RTL outputs.
\paragraph{Line 218}
\noindent\colorbox{gray!8}{\parbox{\dimexpr\linewidth-2\fboxsep\relax}{\raggedright\ttfamily\scriptsize     for start, end in arrows:}}
\par\smallskip
This Python line participates in the export, validation, or report-generation pipeline. It is executed on the host PC to prepare data for RTL or to check RTL outputs.
\paragraph{Line 219}
\noindent\colorbox{gray!8}{\parbox{\dimexpr\linewidth-2\fboxsep\relax}{\raggedright\ttfamily\scriptsize         ax.annotate("", xy=end, xytext=start, arrowprops=\{"arrowstyle": "->", "lw": 1.5, "color": "\#333"\})}}
\par\smallskip
This Python line participates in the export, validation, or report-generation pipeline. It is executed on the host PC to prepare data for RTL or to check RTL outputs.
\paragraph{Line 220}
\noindent\colorbox{gray!8}{\parbox{\dimexpr\linewidth-2\fboxsep\relax}{\raggedright\ttfamily\scriptsize     ax.set\_title("NeuraEdge CIFAR-10 Spikingformer Head Accelerator Dataflow", fontsize=16, weight="bold")}}
\par\smallskip
This Python line participates in the export, validation, or report-generation pipeline. It is executed on the host PC to prepare data for RTL or to check RTL outputs.
\paragraph{Line 221}
\noindent\colorbox{gray!8}{\parbox{\dimexpr\linewidth-2\fboxsep\relax}{\raggedright\ttfamily\scriptsize     ax.text(0.5, 0.3, "All numbers in this report come from the generated RTL memories, XSim CSV traces, and Vivado reports in outputs\_spikingformer\_head\_accelerator.", fontsize=9)}}
\par\smallskip
This Python line participates in the export, validation, or report-generation pipeline. It is executed on the host PC to prepare data for RTL or to check RTL outputs.
\paragraph{Line 222}
\noindent\colorbox{gray!8}{\parbox{\dimexpr\linewidth-2\fboxsep\relax}{\raggedright\ttfamily\scriptsize     return fig}}
\par\smallskip
This Python line participates in the export, validation, or report-generation pipeline. It is executed on the host PC to prepare data for RTL or to check RTL outputs.
\paragraph{Line 223}
\noindent\colorbox{gray!8}{\parbox{\dimexpr\linewidth-2\fboxsep\relax}{\raggedright\ttfamily\scriptsize \emph{blank line}}}
\par\smallskip
This blank line is not hardware or software behavior. It separates nearby declarations or operations so the file remains readable and reviewable.
\paragraph{Line 224}
\noindent\colorbox{gray!8}{\parbox{\dimexpr\linewidth-2\fboxsep\relax}{\raggedright\ttfamily\scriptsize \emph{blank line}}}
\par\smallskip
This blank line is not hardware or software behavior. It separates nearby declarations or operations so the file remains readable and reviewable.
\paragraph{Line 225}
\noindent\colorbox{gray!8}{\parbox{\dimexpr\linewidth-2\fboxsep\relax}{\raggedright\ttfamily\scriptsize def plot\_fsm() -> plt.Figure:}}
\par\smallskip
This defines a Python function. The function packages one repeatable step such as loading a memory file, simulating the golden learning rule, validating CSVs, or generating plots.
\paragraph{Line 226}
\noindent\colorbox{gray!8}{\parbox{\dimexpr\linewidth-2\fboxsep\relax}{\raggedright\ttfamily\scriptsize     fig, ax = plt.subplots(figsize=(11, 8.5))}}
\par\smallskip
This creates a debug visualization. The plot is derived from generated CSV or Vivado report data so it can be regenerated after RTL changes.
\paragraph{Line 227}
\noindent\colorbox{gray!8}{\parbox{\dimexpr\linewidth-2\fboxsep\relax}{\raggedright\ttfamily\scriptsize     ax.set\_axis\_off()}}
\par\smallskip
This Python line participates in the export, validation, or report-generation pipeline. It is executed on the host PC to prepare data for RTL or to check RTL outputs.
\paragraph{Line 228}
\noindent\colorbox{gray!8}{\parbox{\dimexpr\linewidth-2\fboxsep\relax}{\raggedright\ttfamily\scriptsize     ax.set\_xlim(0, 11)}}
\par\smallskip
This Python line participates in the export, validation, or report-generation pipeline. It is executed on the host PC to prepare data for RTL or to check RTL outputs.
\paragraph{Line 229}
\noindent\colorbox{gray!8}{\parbox{\dimexpr\linewidth-2\fboxsep\relax}{\raggedright\ttfamily\scriptsize     ax.set\_ylim(0, 8)}}
\par\smallskip
This Python line participates in the export, validation, or report-generation pipeline. It is executed on the host PC to prepare data for RTL or to check RTL outputs.
\paragraph{Line 230}
\noindent\colorbox{gray!8}{\parbox{\dimexpr\linewidth-2\fboxsep\relax}{\raggedright\ttfamily\scriptsize     fsm = [}}
\par\smallskip
This Python line participates in the export, validation, or report-generation pipeline. It is executed on the host PC to prepare data for RTL or to check RTL outputs.
\paragraph{Line 231}
\noindent\colorbox{gray!8}{\parbox{\dimexpr\linewidth-2\fboxsep\relax}{\raggedright\ttfamily\scriptsize         ("IDLE", 0.5, 6.5), ("LOAD", 2.0, 6.5), ("CLASS\_INIT", 3.7, 6.5),}}
\par\smallskip
This Python line participates in the export, validation, or report-generation pipeline. It is executed on the host PC to prepare data for RTL or to check RTL outputs.
\paragraph{Line 232}
\noindent\colorbox{gray!8}{\parbox{\dimexpr\linewidth-2\fboxsep\relax}{\raggedright\ttfamily\scriptsize         ("MAC\_READ", 5.6, 6.5), ("MAC\_MUL", 7.3, 6.5), ("MAC\_ACCUM", 9.0, 6.5),}}
\par\smallskip
This Python line participates in the export, validation, or report-generation pipeline. It is executed on the host PC to prepare data for RTL or to check RTL outputs.
\paragraph{Line 233}
\noindent\colorbox{gray!8}{\parbox{\dimexpr\linewidth-2\fboxsep\relax}{\raggedright\ttfamily\scriptsize         ("CLASS\_DONE", 5.5, 5.0), ("TRAIN\_CHECK", 3.2, 3.6), ("RESULT", 8.0, 3.6),}}
\par\smallskip
This Python line participates in the export, validation, or report-generation pipeline. It is executed on the host PC to prepare data for RTL or to check RTL outputs.
\paragraph{Line 234}
\noindent\colorbox{gray!8}{\parbox{\dimexpr\linewidth-2\fboxsep\relax}{\raggedright\ttfamily\scriptsize         ("READ\_TRUE", 0.7, 2.2), ("CALC\_TRUE", 2.4, 2.2), ("WRITE\_TRUE", 4.1, 2.2),}}
\par\smallskip
This Python line participates in the export, validation, or report-generation pipeline. It is executed on the host PC to prepare data for RTL or to check RTL outputs.
\paragraph{Line 235}
\noindent\colorbox{gray!8}{\parbox{\dimexpr\linewidth-2\fboxsep\relax}{\raggedright\ttfamily\scriptsize         ("READ\_PRED", 5.8, 2.2), ("CALC\_PRED", 7.5, 2.2), ("WRITE\_PRED", 9.2, 2.2),}}
\par\smallskip
This Python line participates in the export, validation, or report-generation pipeline. It is executed on the host PC to prepare data for RTL or to check RTL outputs.
\paragraph{Line 236}
\noindent\colorbox{gray!8}{\parbox{\dimexpr\linewidth-2\fboxsep\relax}{\raggedright\ttfamily\scriptsize         ("BIAS\_TRUE", 4.2, 0.9), ("BIAS\_PRED", 6.2, 0.9),}}
\par\smallskip
This Python line participates in the export, validation, or report-generation pipeline. It is executed on the host PC to prepare data for RTL or to check RTL outputs.
\paragraph{Line 237}
\noindent\colorbox{gray!8}{\parbox{\dimexpr\linewidth-2\fboxsep\relax}{\raggedright\ttfamily\scriptsize     ]}}
\par\smallskip
This Python line participates in the export, validation, or report-generation pipeline. It is executed on the host PC to prepare data for RTL or to check RTL outputs.
\paragraph{Line 238}
\noindent\colorbox{gray!8}{\parbox{\dimexpr\linewidth-2\fboxsep\relax}{\raggedright\ttfamily\scriptsize     positions = \{\}}}
\par\smallskip
This Python line participates in the export, validation, or report-generation pipeline. It is executed on the host PC to prepare data for RTL or to check RTL outputs.
\paragraph{Line 239}
\noindent\colorbox{gray!8}{\parbox{\dimexpr\linewidth-2\fboxsep\relax}{\raggedright\ttfamily\scriptsize     for name, x, y in fsm:}}
\par\smallskip
This Python line participates in the export, validation, or report-generation pipeline. It is executed on the host PC to prepare data for RTL or to check RTL outputs.
\paragraph{Line 240}
\noindent\colorbox{gray!8}{\parbox{\dimexpr\linewidth-2\fboxsep\relax}{\raggedright\ttfamily\scriptsize         positions[name] = (x, y)}}
\par\smallskip
This Python line participates in the export, validation, or report-generation pipeline. It is executed on the host PC to prepare data for RTL or to check RTL outputs.
\paragraph{Line 241}
\noindent\colorbox{gray!8}{\parbox{\dimexpr\linewidth-2\fboxsep\relax}{\raggedright\ttfamily\scriptsize         draw\_box(ax, (x, y), 1.35, 0.55, name, "\#f6fbff" if "TRAIN" not in name and "TRUE" not in name and "PRED" not in name and "BIAS" not in name else "\#fff1f1")}}
\par\smallskip
This Python line participates in the export, validation, or report-generation pipeline. It is executed on the host PC to prepare data for RTL or to check RTL outputs.
\paragraph{Line 242}
\noindent\colorbox{gray!8}{\parbox{\dimexpr\linewidth-2\fboxsep\relax}{\raggedright\ttfamily\scriptsize     edges = [}}
\par\smallskip
This Python line participates in the export, validation, or report-generation pipeline. It is executed on the host PC to prepare data for RTL or to check RTL outputs.
\paragraph{Line 243}
\noindent\colorbox{gray!8}{\parbox{\dimexpr\linewidth-2\fboxsep\relax}{\raggedright\ttfamily\scriptsize         ("IDLE", "LOAD"), ("LOAD", "CLASS\_INIT"), ("CLASS\_INIT", "MAC\_READ"), ("MAC\_READ", "MAC\_MUL"),}}
\par\smallskip
This Python line participates in the export, validation, or report-generation pipeline. It is executed on the host PC to prepare data for RTL or to check RTL outputs.
\paragraph{Line 244}
\noindent\colorbox{gray!8}{\parbox{\dimexpr\linewidth-2\fboxsep\relax}{\raggedright\ttfamily\scriptsize         ("MAC\_MUL", "MAC\_ACCUM"), ("MAC\_ACCUM", "CLASS\_DONE"), ("CLASS\_DONE", "MAC\_READ"),}}
\par\smallskip
This Python line participates in the export, validation, or report-generation pipeline. It is executed on the host PC to prepare data for RTL or to check RTL outputs.
\paragraph{Line 245}
\noindent\colorbox{gray!8}{\parbox{\dimexpr\linewidth-2\fboxsep\relax}{\raggedright\ttfamily\scriptsize         ("CLASS\_DONE", "TRAIN\_CHECK"), ("TRAIN\_CHECK", "RESULT"), ("RESULT", "IDLE"),}}
\par\smallskip
This Python line participates in the export, validation, or report-generation pipeline. It is executed on the host PC to prepare data for RTL or to check RTL outputs.
\paragraph{Line 246}
\noindent\colorbox{gray!8}{\parbox{\dimexpr\linewidth-2\fboxsep\relax}{\raggedright\ttfamily\scriptsize         ("TRAIN\_CHECK", "READ\_TRUE"), ("READ\_TRUE", "CALC\_TRUE"), ("CALC\_TRUE", "WRITE\_TRUE"),}}
\par\smallskip
This Python line participates in the export, validation, or report-generation pipeline. It is executed on the host PC to prepare data for RTL or to check RTL outputs.
\paragraph{Line 247}
\noindent\colorbox{gray!8}{\parbox{\dimexpr\linewidth-2\fboxsep\relax}{\raggedright\ttfamily\scriptsize         ("WRITE\_TRUE", "READ\_PRED"), ("READ\_PRED", "CALC\_PRED"), ("CALC\_PRED", "WRITE\_PRED"),}}
\par\smallskip
This Python line participates in the export, validation, or report-generation pipeline. It is executed on the host PC to prepare data for RTL or to check RTL outputs.
\paragraph{Line 248}
\noindent\colorbox{gray!8}{\parbox{\dimexpr\linewidth-2\fboxsep\relax}{\raggedright\ttfamily\scriptsize         ("WRITE\_PRED", "READ\_TRUE"), ("WRITE\_PRED", "BIAS\_TRUE"), ("BIAS\_TRUE", "BIAS\_PRED"), ("BIAS\_PRED", "RESULT"),}}
\par\smallskip
This Python line participates in the export, validation, or report-generation pipeline. It is executed on the host PC to prepare data for RTL or to check RTL outputs.
\paragraph{Line 249}
\noindent\colorbox{gray!8}{\parbox{\dimexpr\linewidth-2\fboxsep\relax}{\raggedright\ttfamily\scriptsize     ]}}
\par\smallskip
This Python line participates in the export, validation, or report-generation pipeline. It is executed on the host PC to prepare data for RTL or to check RTL outputs.
\paragraph{Line 250}
\noindent\colorbox{gray!8}{\parbox{\dimexpr\linewidth-2\fboxsep\relax}{\raggedright\ttfamily\scriptsize     for a, b in edges:}}
\par\smallskip
This Python line participates in the export, validation, or report-generation pipeline. It is executed on the host PC to prepare data for RTL or to check RTL outputs.
\paragraph{Line 251}
\noindent\colorbox{gray!8}{\parbox{\dimexpr\linewidth-2\fboxsep\relax}{\raggedright\ttfamily\scriptsize         ax.annotate("", xy=(positions[b][0], positions[b][1] + 0.27), xytext=(positions[a][0] + 1.35, positions[a][1] + 0.27), arrowprops=\{"arrowstyle": "->", "lw": 1.0\})}}
\par\smallskip
This Python line participates in the export, validation, or report-generation pipeline. It is executed on the host PC to prepare data for RTL or to check RTL outputs.
\paragraph{Line 252}
\noindent\colorbox{gray!8}{\parbox{\dimexpr\linewidth-2\fboxsep\relax}{\raggedright\ttfamily\scriptsize     ax.set\_title("RTL State Machine Including On-Chip Learning", fontsize=16, weight="bold")}}
\par\smallskip
This Python line participates in the export, validation, or report-generation pipeline. It is executed on the host PC to prepare data for RTL or to check RTL outputs.
\paragraph{Line 253}
\noindent\colorbox{gray!8}{\parbox{\dimexpr\linewidth-2\fboxsep\relax}{\raggedright\ttfamily\scriptsize     return fig}}
\par\smallskip
This Python line participates in the export, validation, or report-generation pipeline. It is executed on the host PC to prepare data for RTL or to check RTL outputs.
\paragraph{Line 254}
\noindent\colorbox{gray!8}{\parbox{\dimexpr\linewidth-2\fboxsep\relax}{\raggedright\ttfamily\scriptsize \emph{blank line}}}
\par\smallskip
This blank line is not hardware or software behavior. It separates nearby declarations or operations so the file remains readable and reviewable.
\paragraph{Line 255}
\noindent\colorbox{gray!8}{\parbox{\dimexpr\linewidth-2\fboxsep\relax}{\raggedright\ttfamily\scriptsize \emph{blank line}}}
\par\smallskip
This blank line is not hardware or software behavior. It separates nearby declarations or operations so the file remains readable and reviewable.
\paragraph{Line 256}
\noindent\colorbox{gray!8}{\parbox{\dimexpr\linewidth-2\fboxsep\relax}{\raggedright\ttfamily\scriptsize def generate\_debug\_pdf() -> None:}}
\par\smallskip
This defines a Python function. The function packages one repeatable step such as loading a memory file, simulating the golden learning rule, validating CSVs, or generating plots.
\paragraph{Line 257}
\noindent\colorbox{gray!8}{\parbox{\dimexpr\linewidth-2\fboxsep\relax}{\raggedright\ttfamily\scriptsize     pdir = project\_dir()}}
\par\smallskip
This Python line participates in the export, validation, or report-generation pipeline. It is executed on the host PC to prepare data for RTL or to check RTL outputs.
\paragraph{Line 258}
\noindent\colorbox{gray!8}{\parbox{\dimexpr\linewidth-2\fboxsep\relax}{\raggedright\ttfamily\scriptsize     odir = out\_dir()}}
\par\smallskip
This Python line participates in the export, validation, or report-generation pipeline. It is executed on the host PC to prepare data for RTL or to check RTL outputs.
\paragraph{Line 259}
\noindent\colorbox{gray!8}{\parbox{\dimexpr\linewidth-2\fboxsep\relax}{\raggedright\ttfamily\scriptsize     ddir = debug\_dir()}}
\par\smallskip
This Python line participates in the export, validation, or report-generation pipeline. It is executed on the host PC to prepare data for RTL or to check RTL outputs.
\paragraph{Line 260}
\noindent\colorbox{gray!8}{\parbox{\dimexpr\linewidth-2\fboxsep\relax}{\raggedright\ttfamily\scriptsize     png\_dir = ddir / "plots"}}
\par\smallskip
This Python line participates in the export, validation, or report-generation pipeline. It is executed on the host PC to prepare data for RTL or to check RTL outputs.
\paragraph{Line 261}
\noindent\colorbox{gray!8}{\parbox{\dimexpr\linewidth-2\fboxsep\relax}{\raggedright\ttfamily\scriptsize     ddir.mkdir(parents=True, exist\_ok=True)}}
\par\smallskip
This Python line participates in the export, validation, or report-generation pipeline. It is executed on the host PC to prepare data for RTL or to check RTL outputs.
\paragraph{Line 262}
\noindent\colorbox{gray!8}{\parbox{\dimexpr\linewidth-2\fboxsep\relax}{\raggedright\ttfamily\scriptsize \emph{blank line}}}
\par\smallskip
This blank line is not hardware or software behavior. It separates nearby declarations or operations so the file remains readable and reviewable.
\paragraph{Line 263}
\noindent\colorbox{gray!8}{\parbox{\dimexpr\linewidth-2\fboxsep\relax}{\raggedright\ttfamily\scriptsize     inference\_rows = read\_csv\_rows(pdir / "cifar10\_spikingformer\_head\_results.csv")}}
\par\smallskip
This Python line participates in the export, validation, or report-generation pipeline. It is executed on the host PC to prepare data for RTL or to check RTL outputs.
\paragraph{Line 264}
\noindent\colorbox{gray!8}{\parbox{\dimexpr\linewidth-2\fboxsep\relax}{\raggedright\ttfamily\scriptsize     learn\_rows = read\_csv\_rows(pdir / "cifar10\_spikingformer\_head\_onchip\_learning\_results.csv")}}
\par\smallskip
This Python line participates in the export, validation, or report-generation pipeline. It is executed on the host PC to prepare data for RTL or to check RTL outputs.
\paragraph{Line 265}
\noindent\colorbox{gray!8}{\parbox{\dimexpr\linewidth-2\fboxsep\relax}{\raggedright\ttfamily\scriptsize     trace\_rows = read\_csv\_rows(pdir / "cifar10\_spikingformer\_head\_trace.csv")}}
\par\smallskip
This Python line participates in the export, validation, or report-generation pipeline. It is executed on the host PC to prepare data for RTL or to check RTL outputs.
\paragraph{Line 266}
\noindent\colorbox{gray!8}{\parbox{\dimexpr\linewidth-2\fboxsep\relax}{\raggedright\ttfamily\scriptsize     learn\_trace\_rows = read\_csv\_rows(pdir / "cifar10\_spikingformer\_head\_onchip\_learning\_trace.csv")}}
\par\smallskip
This Python line participates in the export, validation, or report-generation pipeline. It is executed on the host PC to prepare data for RTL or to check RTL outputs.
\paragraph{Line 267}
\noindent\colorbox{gray!8}{\parbox{\dimexpr\linewidth-2\fboxsep\relax}{\raggedright\ttfamily\scriptsize     rtl\_val = read\_json(odir / "spikingformer\_head\_rtl\_validation.json")}}
\par\smallskip
This Python line participates in the export, validation, or report-generation pipeline. It is executed on the host PC to prepare data for RTL or to check RTL outputs.
\paragraph{Line 268}
\noindent\colorbox{gray!8}{\parbox{\dimexpr\linewidth-2\fboxsep\relax}{\raggedright\ttfamily\scriptsize     learn\_val = read\_json(odir / "spikingformer\_head\_onchip\_learning\_validation.json")}}
\par\smallskip
This Python line participates in the export, validation, or report-generation pipeline. It is executed on the host PC to prepare data for RTL or to check RTL outputs.
\paragraph{Line 269}
\noindent\colorbox{gray!8}{\parbox{\dimexpr\linewidth-2\fboxsep\relax}{\raggedright\ttfamily\scriptsize     mapping = read\_json(odir / "spikingformer\_head\_mapping\_metrics.json")}}
\par\smallskip
This Python line participates in the export, validation, or report-generation pipeline. It is executed on the host PC to prepare data for RTL or to check RTL outputs.
\paragraph{Line 270}
\noindent\colorbox{gray!8}{\parbox{\dimexpr\linewidth-2\fboxsep\relax}{\raggedright\ttfamily\scriptsize     vivado = parse\_vivado\_reports(odir / "vivado")}}
\par\smallskip
This Python line participates in the export, validation, or report-generation pipeline. It is executed on the host PC to prepare data for RTL or to check RTL outputs.
\paragraph{Line 271}
\noindent\colorbox{gray!8}{\parbox{\dimexpr\linewidth-2\fboxsep\relax}{\raggedright\ttfamily\scriptsize     infer\_skip\_abs\_thresh = int(rtl\_val.get("infer\_skip\_abs\_thresh", mapping.get("infer\_skip\_abs\_thresh", 0)) or 0)}}
\par\smallskip
This Python line participates in the export, validation, or report-generation pipeline. It is executed on the host PC to prepare data for RTL or to check RTL outputs.
\paragraph{Line 272}
\noindent\colorbox{gray!8}{\parbox{\dimexpr\linewidth-2\fboxsep\relax}{\raggedright\ttfamily\scriptsize     feature\_mac\_skip\_fraction = float(rtl\_val.get("feature\_mac\_skip\_fraction", mapping.get("feature\_mac\_skip\_fraction", 0.0)) or 0.0)}}
\par\smallskip
This Python line participates in the export, validation, or report-generation pipeline. It is executed on the host PC to prepare data for RTL or to check RTL outputs.
\paragraph{Line 273}
\noindent\colorbox{gray!8}{\parbox{\dimexpr\linewidth-2\fboxsep\relax}{\raggedright\ttfamily\scriptsize     rtl\_infer\_skip\_count = int(inference\_rows[-1].get("infer\_skip\_count", 0)) if inference\_rows else 0}}
\par\smallskip
This Python line participates in the export, validation, or report-generation pipeline. It is executed on the host PC to prepare data for RTL or to check RTL outputs.
\paragraph{Line 274}
\noindent\colorbox{gray!8}{\parbox{\dimexpr\linewidth-2\fboxsep\relax}{\raggedright\ttfamily\scriptsize     learn\_infer\_skip\_count = int(learn\_rows[-1].get("infer\_skip\_count", 0)) if learn\_rows else 0}}
\par\smallskip
This Python line participates in the export, validation, or report-generation pipeline. It is executed on the host PC to prepare data for RTL or to check RTL outputs.
\paragraph{Line 275}
\noindent\colorbox{gray!8}{\parbox{\dimexpr\linewidth-2\fboxsep\relax}{\raggedright\ttfamily\scriptsize     learn\_train\_skip\_count = int(learn\_rows[-1].get("train\_skip\_count", 0)) if learn\_rows else 0}}
\par\smallskip
This Python line participates in the export, validation, or report-generation pipeline. It is executed on the host PC to prepare data for RTL or to check RTL outputs.
\paragraph{Line 276}
\noindent\colorbox{gray!8}{\parbox{\dimexpr\linewidth-2\fboxsep\relax}{\raggedright\ttfamily\scriptsize \emph{blank line}}}
\par\smallskip
This blank line is not hardware or software behavior. It separates nearby declarations or operations so the file remains readable and reviewable.
\paragraph{Line 277}
\noindent\colorbox{gray!8}{\parbox{\dimexpr\linewidth-2\fboxsep\relax}{\raggedright\ttfamily\scriptsize     labels = col(inference\_rows, "label", int)}}
\par\smallskip
This Python line participates in the export, validation, or report-generation pipeline. It is executed on the host PC to prepare data for RTL or to check RTL outputs.
\paragraph{Line 278}
\noindent\colorbox{gray!8}{\parbox{\dimexpr\linewidth-2\fboxsep\relax}{\raggedright\ttfamily\scriptsize     preds = col(inference\_rows, "predicted", int)}}
\par\smallskip
This line handles predicted labels. It checks that argmax behavior in Python and RTL chooses the same class.
\paragraph{Line 279}
\noindent\colorbox{gray!8}{\parbox{\dimexpr\linewidth-2\fboxsep\relax}{\raggedright\ttfamily\scriptsize     correct = col(inference\_rows, "correct", int)}}
\par\smallskip
This Python line participates in the export, validation, or report-generation pipeline. It is executed on the host PC to prepare data for RTL or to check RTL outputs.
\paragraph{Line 280}
\noindent\colorbox{gray!8}{\parbox{\dimexpr\linewidth-2\fboxsep\relax}{\raggedright\ttfamily\scriptsize     learn\_preds = col(learn\_rows, "predicted", int)}}
\par\smallskip
This line handles predicted labels. It checks that argmax behavior in Python and RTL chooses the same class.
\paragraph{Line 281}
\noindent\colorbox{gray!8}{\parbox{\dimexpr\linewidth-2\fboxsep\relax}{\raggedright\ttfamily\scriptsize     learn\_correct = col(learn\_rows, "correct", int)}}
\par\smallskip
This Python line participates in the export, validation, or report-generation pipeline. It is executed on the host PC to prepare data for RTL or to check RTL outputs.
\paragraph{Line 282}
\noindent\colorbox{gray!8}{\parbox{\dimexpr\linewidth-2\fboxsep\relax}{\raggedright\ttfamily\scriptsize     updates = col(learn\_rows, "update\_applied", int)}}
\par\smallskip
This line handles the online-learning rule or its validation fields. It is central to proving that the FPGA update path is not bypassed.
\paragraph{Line 283}
\noindent\colorbox{gray!8}{\parbox{\dimexpr\linewidth-2\fboxsep\relax}{\raggedright\ttfamily\scriptsize     update\_count = col(learn\_rows, "update\_count", int)}}
\par\smallskip
This line handles the online-learning rule or its validation fields. It is central to proving that the FPGA update path is not bypassed.
\paragraph{Line 284}
\noindent\colorbox{gray!8}{\parbox{\dimexpr\linewidth-2\fboxsep\relax}{\raggedright\ttfamily\scriptsize     scores = score\_matrix(inference\_rows)}}
\par\smallskip
This line handles class scores. The score vector is the direct software equivalent of the RTL MAC loop over weights and features.
\paragraph{Line 285}
\noindent\colorbox{gray!8}{\parbox{\dimexpr\linewidth-2\fboxsep\relax}{\raggedright\ttfamily\scriptsize     learn\_scores = score\_matrix(learn\_rows)}}
\par\smallskip
This line handles class scores. The score vector is the direct software equivalent of the RTL MAC loop over weights and features.
\paragraph{Line 286}
\noindent\colorbox{gray!8}{\parbox{\dimexpr\linewidth-2\fboxsep\relax}{\raggedright\ttfamily\scriptsize     conf = confusion(labels, preds)}}
\par\smallskip
This line handles predicted labels. It checks that argmax behavior in Python and RTL chooses the same class.
\paragraph{Line 287}
\noindent\colorbox{gray!8}{\parbox{\dimexpr\linewidth-2\fboxsep\relax}{\raggedright\ttfamily\scriptsize     precision, recall, f1 = precision\_recall\_f1(conf)}}
\par\smallskip
This Python line participates in the export, validation, or report-generation pipeline. It is executed on the host PC to prepare data for RTL or to check RTL outputs.
\paragraph{Line 288}
\noindent\colorbox{gray!8}{\parbox{\dimexpr\linewidth-2\fboxsep\relax}{\raggedright\ttfamily\scriptsize     margins, true\_gaps = margin(scores, labels)}}
\par\smallskip
This line handles class scores. The score vector is the direct software equivalent of the RTL MAC loop over weights and features.
\paragraph{Line 289}
\noindent\colorbox{gray!8}{\parbox{\dimexpr\linewidth-2\fboxsep\relax}{\raggedright\ttfamily\scriptsize \emph{blank line}}}
\par\smallskip
This blank line is not hardware or software behavior. It separates nearby declarations or operations so the file remains readable and reviewable.
\paragraph{Line 290}
\noindent\colorbox{gray!8}{\parbox{\dimexpr\linewidth-2\fboxsep\relax}{\raggedright\ttfamily\scriptsize     rtl\_dir = pdir / "rtl"}}
\par\smallskip
This Python line participates in the export, validation, or report-generation pipeline. It is executed on the host PC to prepare data for RTL or to check RTL outputs.
\paragraph{Line 291}
\noindent\colorbox{gray!8}{\parbox{\dimexpr\linewidth-2\fboxsep\relax}{\raggedright\ttfamily\scriptsize     features = read\_hex\_mem(rtl\_dir / "cifar10\_spikingformer\_head\_features.mem", 16)}}
\par\smallskip
This Python line participates in the export, validation, or report-generation pipeline. It is executed on the host PC to prepare data for RTL or to check RTL outputs.
\paragraph{Line 292}
\noindent\colorbox{gray!8}{\parbox{\dimexpr\linewidth-2\fboxsep\relax}{\raggedright\ttfamily\scriptsize     weights = read\_hex\_mem(rtl\_dir / "cifar10\_spikingformer\_head\_weights.mem", 16).reshape(NUM\_CLASSES, FEATURE\_DIM)}}
\par\smallskip
This line handles the classifier parameters. Those arrays mirror the RTL memories, so mismatches reveal an export, quantization, or hardware-update bug.
\paragraph{Line 293}
\noindent\colorbox{gray!8}{\parbox{\dimexpr\linewidth-2\fboxsep\relax}{\raggedright\ttfamily\scriptsize     bias = read\_hex\_mem(rtl\_dir / "cifar10\_spikingformer\_head\_bias.mem", 64)}}
\par\smallskip
This line handles the classifier parameters. Those arrays mirror the RTL memories, so mismatches reveal an export, quantization, or hardware-update bug.
\paragraph{Line 294}
\noindent\colorbox{gray!8}{\parbox{\dimexpr\linewidth-2\fboxsep\relax}{\raggedright\ttfamily\scriptsize     after\_weights\_path = pdir / "cifar10\_spikingformer\_head\_onchip\_learning\_weights\_after.mem"}}
\par\smallskip
This line handles the classifier parameters. Those arrays mirror the RTL memories, so mismatches reveal an export, quantization, or hardware-update bug.
\paragraph{Line 295}
\noindent\colorbox{gray!8}{\parbox{\dimexpr\linewidth-2\fboxsep\relax}{\raggedright\ttfamily\scriptsize     after\_bias\_path = pdir / "cifar10\_spikingformer\_head\_onchip\_learning\_bias\_after.mem"}}
\par\smallskip
This line handles the classifier parameters. Those arrays mirror the RTL memories, so mismatches reveal an export, quantization, or hardware-update bug.
\paragraph{Line 296}
\noindent\colorbox{gray!8}{\parbox{\dimexpr\linewidth-2\fboxsep\relax}{\raggedright\ttfamily\scriptsize     after\_weights = read\_decimal\_mem(after\_weights\_path).reshape(NUM\_CLASSES, FEATURE\_DIM) if after\_weights\_path.exists() else weights.copy()}}
\par\smallskip
This line handles the classifier parameters. Those arrays mirror the RTL memories, so mismatches reveal an export, quantization, or hardware-update bug.
\paragraph{Line 297}
\noindent\colorbox{gray!8}{\parbox{\dimexpr\linewidth-2\fboxsep\relax}{\raggedright\ttfamily\scriptsize     after\_bias = read\_decimal\_mem(after\_bias\_path) if after\_bias\_path.exists() else bias.copy()}}
\par\smallskip
This line handles the classifier parameters. Those arrays mirror the RTL memories, so mismatches reveal an export, quantization, or hardware-update bug.
\paragraph{Line 298}
\noindent\colorbox{gray!8}{\parbox{\dimexpr\linewidth-2\fboxsep\relax}{\raggedright\ttfamily\scriptsize \emph{blank line}}}
\par\smallskip
This blank line is not hardware or software behavior. It separates nearby declarations or operations so the file remains readable and reviewable.
\paragraph{Line 299}
\noindent\colorbox{gray!8}{\parbox{\dimexpr\linewidth-2\fboxsep\relax}{\raggedright\ttfamily\scriptsize     pdf\_path = ddir / "neuraedge\_spikingformer\_deep\_debug\_report.pdf"}}
\par\smallskip
This Python line participates in the export, validation, or report-generation pipeline. It is executed on the host PC to prepare data for RTL or to check RTL outputs.
\paragraph{Line 300}
\noindent\colorbox{gray!8}{\parbox{\dimexpr\linewidth-2\fboxsep\relax}{\raggedright\ttfamily\scriptsize     with PdfPages(pdf\_path) as pdf:}}
\par\smallskip
This creates a debug visualization. The plot is derived from generated CSV or Vivado report data so it can be regenerated after RTL changes.
\paragraph{Line 301}
\noindent\colorbox{gray!8}{\parbox{\dimexpr\linewidth-2\fboxsep\relax}{\raggedright\ttfamily\scriptsize         summary\_lines = [}}
\par\smallskip
This Python line participates in the export, validation, or report-generation pipeline. It is executed on the host PC to prepare data for RTL or to check RTL outputs.
\paragraph{Line 302}
\noindent\colorbox{gray!8}{\parbox{\dimexpr\linewidth-2\fboxsep\relax}{\raggedright\ttfamily\scriptsize             f"RTL inference samples: \{rtl\_val.get('samples')\}  accuracy: \{float(rtl\_val.get('accuracy', float('nan'))):.4f\}  exact score match: \{rtl\_val.get('score\_exact\_match')\}",}}
\par\smallskip
This Python line participates in the export, validation, or report-generation pipeline. It is executed on the host PC to prepare data for RTL or to check RTL outputs.
\paragraph{Line 303}
\noindent\colorbox{gray!8}{\parbox{\dimexpr\linewidth-2\fboxsep\relax}{\raggedright\ttfamily\scriptsize             f"On-chip learning samples: \{learn\_val.get('samples')\}  updates: \{learn\_val.get('updates')\}  score mismatches: \{learn\_val.get('score\_mismatch\_count')\}",}}
\par\smallskip
This line handles the online-learning rule or its validation fields. It is central to proving that the FPGA update path is not bypassed.
\paragraph{Line 304}
\noindent\colorbox{gray!8}{\parbox{\dimexpr\linewidth-2\fboxsep\relax}{\raggedright\ttfamily\scriptsize             f"Activity gate: abs(feature) <= \{infer\_skip\_abs\_thresh\} is forced to zero; measured feature-MAC skip fraction: \{feature\_mac\_skip\_fraction * 100:.2f\}\%",}}
\par\smallskip
This Python line participates in the export, validation, or report-generation pipeline. It is executed on the host PC to prepare data for RTL or to check RTL outputs.
\paragraph{Line 305}
\noindent\colorbox{gray!8}{\parbox{\dimexpr\linewidth-2\fboxsep\relax}{\raggedright\ttfamily\scriptsize             f"RTL debug counters: inference skips=\{rtl\_infer\_skip\_count:,\}; learning-run inference skips=\{learn\_infer\_skip\_count:,\}; learning feature skips=\{learn\_train\_skip\_count:,\}",}}
\par\smallskip
This Python line participates in the export, validation, or report-generation pipeline. It is executed on the host PC to prepare data for RTL or to check RTL outputs.
\paragraph{Line 306}
\noindent\colorbox{gray!8}{\parbox{\dimexpr\linewidth-2\fboxsep\relax}{\raggedright\ttfamily\scriptsize             f"Vivado top: cifar10\_spikingformer\_head\_pynq\_axi\_stream\_top  part: xc7z020clg400-1  clock: 10 ns",}}
\par\smallskip
This Python line participates in the export, validation, or report-generation pipeline. It is executed on the host PC to prepare data for RTL or to check RTL outputs.
\paragraph{Line 307}
\noindent\colorbox{gray!8}{\parbox{\dimexpr\linewidth-2\fboxsep\relax}{\raggedright\ttfamily\scriptsize             f"Timing: WNS=\{vivado['wns\_ns']:.3f\} ns  WHS=\{vivado['whs\_ns']:.3f\} ns  timing\_met=\{'timing\_met=1' in vivado['summary']\}",}}
\par\smallskip
This Python line participates in the export, validation, or report-generation pipeline. It is executed on the host PC to prepare data for RTL or to check RTL outputs.
\paragraph{Line 308}
\noindent\colorbox{gray!8}{\parbox{\dimexpr\linewidth-2\fboxsep\relax}{\raggedright\ttfamily\scriptsize             f"Utilization: LUT=\{vivado['lut']\}  FF=\{vivado['regs']\}  BRAM36=\{vivado['bram']\}  DSP=\{vivado['dsp']\}  LUTRAM=\{vivado['lut\_mem']\}",}}
\par\smallskip
This Python line participates in the export, validation, or report-generation pipeline. It is executed on the host PC to prepare data for RTL or to check RTL outputs.
\paragraph{Line 309}
\noindent\colorbox{gray!8}{\parbox{\dimexpr\linewidth-2\fboxsep\relax}{\raggedright\ttfamily\scriptsize             f"Power: total=\{vivado['total\_power\_w']:.3f\} W  dynamic=\{vivado['dynamic\_power\_w']:.3f\} W  static=\{vivado['static\_power\_w']:.3f\} W",}}
\par\smallskip
This Python line participates in the export, validation, or report-generation pipeline. It is executed on the host PC to prepare data for RTL or to check RTL outputs.
\paragraph{Line 310}
\noindent\colorbox{gray!8}{\parbox{\dimexpr\linewidth-2\fboxsep\relax}{\raggedright\ttfamily\scriptsize             f"Mapping metrics file fields: \{', '.join(sorted(mapping.keys())[:10])\}",}}
\par\smallskip
This Python line participates in the export, validation, or report-generation pipeline. It is executed on the host PC to prepare data for RTL or to check RTL outputs.
\paragraph{Line 311}
\noindent\colorbox{gray!8}{\parbox{\dimexpr\linewidth-2\fboxsep\relax}{\raggedright\ttfamily\scriptsize             "Important interpretation: the RTL accelerates the trained Spikingformer feature-head path, gates low-magnitude feature activity, and implements a margin-gated on-chip classifier update.",}}
\par\smallskip
This line handles the online-learning rule or its validation fields. It is central to proving that the FPGA update path is not bypassed.
\paragraph{Line 312}
\noindent\colorbox{gray!8}{\parbox{\dimexpr\linewidth-2\fboxsep\relax}{\raggedright\ttfamily\scriptsize         ]}}
\par\smallskip
This Python line participates in the export, validation, or report-generation pipeline. It is executed on the host PC to prepare data for RTL or to check RTL outputs.
\paragraph{Line 313}
\noindent\colorbox{gray!8}{\parbox{\dimexpr\linewidth-2\fboxsep\relax}{\raggedright\ttfamily\scriptsize         save\_pdf\_page(pdf, plot\_text\_page("Measured Build Summary", summary\_lines), png\_dir, "00\_summary")}}
\par\smallskip
This Python line participates in the export, validation, or report-generation pipeline. It is executed on the host PC to prepare data for RTL or to check RTL outputs.
\paragraph{Line 314}
\noindent\colorbox{gray!8}{\parbox{\dimexpr\linewidth-2\fboxsep\relax}{\raggedright\ttfamily\scriptsize         save\_pdf\_page(pdf, plot\_architecture(), png\_dir, "01\_architecture")}}
\par\smallskip
This Python line participates in the export, validation, or report-generation pipeline. It is executed on the host PC to prepare data for RTL or to check RTL outputs.
\paragraph{Line 315}
\noindent\colorbox{gray!8}{\parbox{\dimexpr\linewidth-2\fboxsep\relax}{\raggedright\ttfamily\scriptsize         save\_pdf\_page(pdf, plot\_fsm(), png\_dir, "02\_fsm")}}
\par\smallskip
This Python line participates in the export, validation, or report-generation pipeline. It is executed on the host PC to prepare data for RTL or to check RTL outputs.
\paragraph{Line 316}
\noindent\colorbox{gray!8}{\parbox{\dimexpr\linewidth-2\fboxsep\relax}{\raggedright\ttfamily\scriptsize \emph{blank line}}}
\par\smallskip
This blank line is not hardware or software behavior. It separates nearby declarations or operations so the file remains readable and reviewable.
\paragraph{Line 317}
\noindent\colorbox{gray!8}{\parbox{\dimexpr\linewidth-2\fboxsep\relax}{\raggedright\ttfamily\scriptsize         for name, rows in [("inference", trace\_rows), ("onchip\_learning", learn\_trace\_rows)]:}}
\par\smallskip
This Python line participates in the export, validation, or report-generation pipeline. It is executed on the host PC to prepare data for RTL or to check RTL outputs.
\paragraph{Line 318}
\noindent\colorbox{gray!8}{\parbox{\dimexpr\linewidth-2\fboxsep\relax}{\raggedright\ttfamily\scriptsize             cycles = col(rows, "cycle", int)}}
\par\smallskip
This Python line participates in the export, validation, or report-generation pipeline. It is executed on the host PC to prepare data for RTL or to check RTL outputs.
\paragraph{Line 319}
\noindent\colorbox{gray!8}{\parbox{\dimexpr\linewidth-2\fboxsep\relax}{\raggedright\ttfamily\scriptsize             fig, axes = plt.subplots(5, 1, figsize=(12, 9), sharex=True)}}
\par\smallskip
This creates a debug visualization. The plot is derived from generated CSV or Vivado report data so it can be regenerated after RTL changes.
\paragraph{Line 320}
\noindent\colorbox{gray!8}{\parbox{\dimexpr\linewidth-2\fboxsep\relax}{\raggedright\ttfamily\scriptsize             axes[0].step(cycles, col(rows, "state", int), where="post")}}
\par\smallskip
This Python line participates in the export, validation, or report-generation pipeline. It is executed on the host PC to prepare data for RTL or to check RTL outputs.
\paragraph{Line 321}
\noindent\colorbox{gray!8}{\parbox{\dimexpr\linewidth-2\fboxsep\relax}{\raggedright\ttfamily\scriptsize             axes[0].set\_ylabel("state")}}
\par\smallskip
This Python line participates in the export, validation, or report-generation pipeline. It is executed on the host PC to prepare data for RTL or to check RTL outputs.
\paragraph{Line 322}
\noindent\colorbox{gray!8}{\parbox{\dimexpr\linewidth-2\fboxsep\relax}{\raggedright\ttfamily\scriptsize             axes[1].step(cycles, col(rows, "feature\_valid", int), where="post", label="valid")}}
\par\smallskip
This Python line participates in the export, validation, or report-generation pipeline. It is executed on the host PC to prepare data for RTL or to check RTL outputs.
\paragraph{Line 323}
\noindent\colorbox{gray!8}{\parbox{\dimexpr\linewidth-2\fboxsep\relax}{\raggedright\ttfamily\scriptsize             axes[1].step(cycles, col(rows, "feature\_ready", int), where="post", label="ready")}}
\par\smallskip
This Python line participates in the export, validation, or report-generation pipeline. It is executed on the host PC to prepare data for RTL or to check RTL outputs.
\paragraph{Line 324}
\noindent\colorbox{gray!8}{\parbox{\dimexpr\linewidth-2\fboxsep\relax}{\raggedright\ttfamily\scriptsize             axes[1].legend(loc="upper right")}}
\par\smallskip
This Python line participates in the export, validation, or report-generation pipeline. It is executed on the host PC to prepare data for RTL or to check RTL outputs.
\paragraph{Line 325}
\noindent\colorbox{gray!8}{\parbox{\dimexpr\linewidth-2\fboxsep\relax}{\raggedright\ttfamily\scriptsize             axes[1].set\_ylabel("feature")}}
\par\smallskip
This Python line participates in the export, validation, or report-generation pipeline. It is executed on the host PC to prepare data for RTL or to check RTL outputs.
\paragraph{Line 326}
\noindent\colorbox{gray!8}{\parbox{\dimexpr\linewidth-2\fboxsep\relax}{\raggedright\ttfamily\scriptsize             axes[2].plot(cycles, col(rows, "feature\_idx", int))}}
\par\smallskip
This Python line participates in the export, validation, or report-generation pipeline. It is executed on the host PC to prepare data for RTL or to check RTL outputs.
\paragraph{Line 327}
\noindent\colorbox{gray!8}{\parbox{\dimexpr\linewidth-2\fboxsep\relax}{\raggedright\ttfamily\scriptsize             axes[2].set\_ylabel("feature\_idx")}}
\par\smallskip
This Python line participates in the export, validation, or report-generation pipeline. It is executed on the host PC to prepare data for RTL or to check RTL outputs.
\paragraph{Line 328}
\noindent\colorbox{gray!8}{\parbox{\dimexpr\linewidth-2\fboxsep\relax}{\raggedright\ttfamily\scriptsize             axes[3].plot(cycles, col(rows, "class\_idx", int))}}
\par\smallskip
This Python line participates in the export, validation, or report-generation pipeline. It is executed on the host PC to prepare data for RTL or to check RTL outputs.
\paragraph{Line 329}
\noindent\colorbox{gray!8}{\parbox{\dimexpr\linewidth-2\fboxsep\relax}{\raggedright\ttfamily\scriptsize             axes[3].set\_ylabel("class\_idx")}}
\par\smallskip
This Python line participates in the export, validation, or report-generation pipeline. It is executed on the host PC to prepare data for RTL or to check RTL outputs.
\paragraph{Line 330}
\noindent\colorbox{gray!8}{\parbox{\dimexpr\linewidth-2\fboxsep\relax}{\raggedright\ttfamily\scriptsize             axes[4].step(cycles, col(rows, "result\_valid", int), where="post", label="result")}}
\par\smallskip
This Python line participates in the export, validation, or report-generation pipeline. It is executed on the host PC to prepare data for RTL or to check RTL outputs.
\paragraph{Line 331}
\noindent\colorbox{gray!8}{\parbox{\dimexpr\linewidth-2\fboxsep\relax}{\raggedright\ttfamily\scriptsize             if "train\_update\_applied" in rows[0]:}}
\par\smallskip
This line handles the online-learning rule or its validation fields. It is central to proving that the FPGA update path is not bypassed.
\paragraph{Line 332}
\noindent\colorbox{gray!8}{\parbox{\dimexpr\linewidth-2\fboxsep\relax}{\raggedright\ttfamily\scriptsize                 axes[4].step(cycles, col(rows, "train\_update\_applied", int), where="post", label="update")}}
\par\smallskip
This line handles the online-learning rule or its validation fields. It is central to proving that the FPGA update path is not bypassed.
\paragraph{Line 333}
\noindent\colorbox{gray!8}{\parbox{\dimexpr\linewidth-2\fboxsep\relax}{\raggedright\ttfamily\scriptsize             axes[4].legend(loc="upper right")}}
\par\smallskip
This Python line participates in the export, validation, or report-generation pipeline. It is executed on the host PC to prepare data for RTL or to check RTL outputs.
\paragraph{Line 334}
\noindent\colorbox{gray!8}{\parbox{\dimexpr\linewidth-2\fboxsep\relax}{\raggedright\ttfamily\scriptsize             axes[4].set\_xlabel("cycle")}}
\par\smallskip
This Python line participates in the export, validation, or report-generation pipeline. It is executed on the host PC to prepare data for RTL or to check RTL outputs.
\paragraph{Line 335}
\noindent\colorbox{gray!8}{\parbox{\dimexpr\linewidth-2\fboxsep\relax}{\raggedright\ttfamily\scriptsize             axes[4].set\_ylabel("flags")}}
\par\smallskip
This Python line participates in the export, validation, or report-generation pipeline. It is executed on the host PC to prepare data for RTL or to check RTL outputs.
\paragraph{Line 336}
\noindent\colorbox{gray!8}{\parbox{\dimexpr\linewidth-2\fboxsep\relax}{\raggedright\ttfamily\scriptsize             fig.suptitle(f"Cycle-Level Trace From XSim CSV: \{name\}")}}
\par\smallskip
This Python line participates in the export, validation, or report-generation pipeline. It is executed on the host PC to prepare data for RTL or to check RTL outputs.
\paragraph{Line 337}
\noindent\colorbox{gray!8}{\parbox{\dimexpr\linewidth-2\fboxsep\relax}{\raggedright\ttfamily\scriptsize             save\_pdf\_page(pdf, fig, png\_dir, f"03\_trace\_\{name\}")}}
\par\smallskip
This Python line participates in the export, validation, or report-generation pipeline. It is executed on the host PC to prepare data for RTL or to check RTL outputs.
\paragraph{Line 338}
\noindent\colorbox{gray!8}{\parbox{\dimexpr\linewidth-2\fboxsep\relax}{\raggedright\ttfamily\scriptsize \emph{blank line}}}
\par\smallskip
This blank line is not hardware or software behavior. It separates nearby declarations or operations so the file remains readable and reviewable.
\paragraph{Line 339}
\noindent\colorbox{gray!8}{\parbox{\dimexpr\linewidth-2\fboxsep\relax}{\raggedright\ttfamily\scriptsize             counts = Counter(col(rows, "state", int))}}
\par\smallskip
This Python line participates in the export, validation, or report-generation pipeline. It is executed on the host PC to prepare data for RTL or to check RTL outputs.
\paragraph{Line 340}
\noindent\colorbox{gray!8}{\parbox{\dimexpr\linewidth-2\fboxsep\relax}{\raggedright\ttfamily\scriptsize             states = sorted(counts)}}
\par\smallskip
This Python line participates in the export, validation, or report-generation pipeline. It is executed on the host PC to prepare data for RTL or to check RTL outputs.
\paragraph{Line 341}
\noindent\colorbox{gray!8}{\parbox{\dimexpr\linewidth-2\fboxsep\relax}{\raggedright\ttfamily\scriptsize             fig, ax = plt.subplots(figsize=(11, 5))}}
\par\smallskip
This creates a debug visualization. The plot is derived from generated CSV or Vivado report data so it can be regenerated after RTL changes.
\paragraph{Line 342}
\noindent\colorbox{gray!8}{\parbox{\dimexpr\linewidth-2\fboxsep\relax}{\raggedright\ttfamily\scriptsize             ax.bar([STATE\_NAMES.get(int(s), str(s)) for s in states], [counts[s] for s in states])}}
\par\smallskip
This Python line participates in the export, validation, or report-generation pipeline. It is executed on the host PC to prepare data for RTL or to check RTL outputs.
\paragraph{Line 343}
\noindent\colorbox{gray!8}{\parbox{\dimexpr\linewidth-2\fboxsep\relax}{\raggedright\ttfamily\scriptsize             ax.tick\_params(axis="x", rotation=45)}}
\par\smallskip
This Python line participates in the export, validation, or report-generation pipeline. It is executed on the host PC to prepare data for RTL or to check RTL outputs.
\paragraph{Line 344}
\noindent\colorbox{gray!8}{\parbox{\dimexpr\linewidth-2\fboxsep\relax}{\raggedright\ttfamily\scriptsize             ax.set\_ylabel("cycles in traced samples")}}
\par\smallskip
This Python line participates in the export, validation, or report-generation pipeline. It is executed on the host PC to prepare data for RTL or to check RTL outputs.
\paragraph{Line 345}
\noindent\colorbox{gray!8}{\parbox{\dimexpr\linewidth-2\fboxsep\relax}{\raggedright\ttfamily\scriptsize             ax.set\_title(f"FSM State Occupancy: \{name\}")}}
\par\smallskip
This Python line participates in the export, validation, or report-generation pipeline. It is executed on the host PC to prepare data for RTL or to check RTL outputs.
\paragraph{Line 346}
\noindent\colorbox{gray!8}{\parbox{\dimexpr\linewidth-2\fboxsep\relax}{\raggedright\ttfamily\scriptsize             save\_pdf\_page(pdf, fig, png\_dir, f"04\_state\_occupancy\_\{name\}")}}
\par\smallskip
This Python line participates in the export, validation, or report-generation pipeline. It is executed on the host PC to prepare data for RTL or to check RTL outputs.
\paragraph{Line 347}
\noindent\colorbox{gray!8}{\parbox{\dimexpr\linewidth-2\fboxsep\relax}{\raggedright\ttfamily\scriptsize \emph{blank line}}}
\par\smallskip
This blank line is not hardware or software behavior. It separates nearby declarations or operations so the file remains readable and reviewable.
\paragraph{Line 348}
\noindent\colorbox{gray!8}{\parbox{\dimexpr\linewidth-2\fboxsep\relax}{\raggedright\ttfamily\scriptsize         sample\_idx = np.arange(len(labels))}}
\par\smallskip
This performs vectorized numerical work. The fixed-point golden models use NumPy so score, margin, and update calculations are explicit and repeatable.
\paragraph{Line 349}
\noindent\colorbox{gray!8}{\parbox{\dimexpr\linewidth-2\fboxsep\relax}{\raggedright\ttfamily\scriptsize         fig, ax = plt.subplots(figsize=(11, 5))}}
\par\smallskip
This creates a debug visualization. The plot is derived from generated CSV or Vivado report data so it can be regenerated after RTL changes.
\paragraph{Line 350}
\noindent\colorbox{gray!8}{\parbox{\dimexpr\linewidth-2\fboxsep\relax}{\raggedright\ttfamily\scriptsize         ax.plot(sample\_idx, np.cumsum(correct) / np.maximum(sample\_idx + 1, 1), label="inference")}}
\par\smallskip
This performs vectorized numerical work. The fixed-point golden models use NumPy so score, margin, and update calculations are explicit and repeatable.
\paragraph{Line 351}
\noindent\colorbox{gray!8}{\parbox{\dimexpr\linewidth-2\fboxsep\relax}{\raggedright\ttfamily\scriptsize         ax.plot(sample\_idx, np.cumsum(learn\_correct) / np.maximum(sample\_idx + 1, 1), label="on-chip learning")}}
\par\smallskip
This performs vectorized numerical work. The fixed-point golden models use NumPy so score, margin, and update calculations are explicit and repeatable.
\paragraph{Line 352}
\noindent\colorbox{gray!8}{\parbox{\dimexpr\linewidth-2\fboxsep\relax}{\raggedright\ttfamily\scriptsize         ax.set\_xlabel("sample")}}
\par\smallskip
This Python line participates in the export, validation, or report-generation pipeline. It is executed on the host PC to prepare data for RTL or to check RTL outputs.
\paragraph{Line 353}
\noindent\colorbox{gray!8}{\parbox{\dimexpr\linewidth-2\fboxsep\relax}{\raggedright\ttfamily\scriptsize         ax.set\_ylabel("cumulative accuracy")}}
\par\smallskip
This Python line participates in the export, validation, or report-generation pipeline. It is executed on the host PC to prepare data for RTL or to check RTL outputs.
\paragraph{Line 354}
\noindent\colorbox{gray!8}{\parbox{\dimexpr\linewidth-2\fboxsep\relax}{\raggedright\ttfamily\scriptsize         ax.set\_title("Cumulative Accuracy Over Full CIFAR-10 Test Stream")}}
\par\smallskip
This Python line participates in the export, validation, or report-generation pipeline. It is executed on the host PC to prepare data for RTL or to check RTL outputs.
\paragraph{Line 355}
\noindent\colorbox{gray!8}{\parbox{\dimexpr\linewidth-2\fboxsep\relax}{\raggedright\ttfamily\scriptsize         ax.legend()}}
\par\smallskip
This Python line participates in the export, validation, or report-generation pipeline. It is executed on the host PC to prepare data for RTL or to check RTL outputs.
\paragraph{Line 356}
\noindent\colorbox{gray!8}{\parbox{\dimexpr\linewidth-2\fboxsep\relax}{\raggedright\ttfamily\scriptsize         save\_pdf\_page(pdf, fig, png\_dir, "05\_cumulative\_accuracy")}}
\par\smallskip
This Python line participates in the export, validation, or report-generation pipeline. It is executed on the host PC to prepare data for RTL or to check RTL outputs.
\paragraph{Line 357}
\noindent\colorbox{gray!8}{\parbox{\dimexpr\linewidth-2\fboxsep\relax}{\raggedright\ttfamily\scriptsize \emph{blank line}}}
\par\smallskip
This blank line is not hardware or software behavior. It separates nearby declarations or operations so the file remains readable and reviewable.
\paragraph{Line 358}
\noindent\colorbox{gray!8}{\parbox{\dimexpr\linewidth-2\fboxsep\relax}{\raggedright\ttfamily\scriptsize         window = 200}}
\par\smallskip
This Python line participates in the export, validation, or report-generation pipeline. It is executed on the host PC to prepare data for RTL or to check RTL outputs.
\paragraph{Line 359}
\noindent\colorbox{gray!8}{\parbox{\dimexpr\linewidth-2\fboxsep\relax}{\raggedright\ttfamily\scriptsize         roll = np.convolve(correct, np.ones(window) / window, mode="valid")}}
\par\smallskip
This performs vectorized numerical work. The fixed-point golden models use NumPy so score, margin, and update calculations are explicit and repeatable.
\paragraph{Line 360}
\noindent\colorbox{gray!8}{\parbox{\dimexpr\linewidth-2\fboxsep\relax}{\raggedright\ttfamily\scriptsize         fig, ax = plt.subplots(figsize=(11, 5))}}
\par\smallskip
This creates a debug visualization. The plot is derived from generated CSV or Vivado report data so it can be regenerated after RTL changes.
\paragraph{Line 361}
\noindent\colorbox{gray!8}{\parbox{\dimexpr\linewidth-2\fboxsep\relax}{\raggedright\ttfamily\scriptsize         ax.plot(np.arange(len(roll)) + window, roll)}}
\par\smallskip
This performs vectorized numerical work. The fixed-point golden models use NumPy so score, margin, and update calculations are explicit and repeatable.
\paragraph{Line 362}
\noindent\colorbox{gray!8}{\parbox{\dimexpr\linewidth-2\fboxsep\relax}{\raggedright\ttfamily\scriptsize         ax.set\_xlabel("sample")}}
\par\smallskip
This Python line participates in the export, validation, or report-generation pipeline. It is executed on the host PC to prepare data for RTL or to check RTL outputs.
\paragraph{Line 363}
\noindent\colorbox{gray!8}{\parbox{\dimexpr\linewidth-2\fboxsep\relax}{\raggedright\ttfamily\scriptsize         ax.set\_ylabel(f"rolling accuracy (\{window\})")}}
\par\smallskip
This Python line participates in the export, validation, or report-generation pipeline. It is executed on the host PC to prepare data for RTL or to check RTL outputs.
\paragraph{Line 364}
\noindent\colorbox{gray!8}{\parbox{\dimexpr\linewidth-2\fboxsep\relax}{\raggedright\ttfamily\scriptsize         ax.set\_title("Rolling Accuracy Debug View")}}
\par\smallskip
This Python line participates in the export, validation, or report-generation pipeline. It is executed on the host PC to prepare data for RTL or to check RTL outputs.
\paragraph{Line 365}
\noindent\colorbox{gray!8}{\parbox{\dimexpr\linewidth-2\fboxsep\relax}{\raggedright\ttfamily\scriptsize         save\_pdf\_page(pdf, fig, png\_dir, "06\_rolling\_accuracy")}}
\par\smallskip
This Python line participates in the export, validation, or report-generation pipeline. It is executed on the host PC to prepare data for RTL or to check RTL outputs.
\paragraph{Line 366}
\noindent\colorbox{gray!8}{\parbox{\dimexpr\linewidth-2\fboxsep\relax}{\raggedright\ttfamily\scriptsize \emph{blank line}}}
\par\smallskip
This blank line is not hardware or software behavior. It separates nearby declarations or operations so the file remains readable and reviewable.
\paragraph{Line 367}
\noindent\colorbox{gray!8}{\parbox{\dimexpr\linewidth-2\fboxsep\relax}{\raggedright\ttfamily\scriptsize         fig, ax = plt.subplots(figsize=(8, 7))}}
\par\smallskip
This creates a debug visualization. The plot is derived from generated CSV or Vivado report data so it can be regenerated after RTL changes.
\paragraph{Line 368}
\noindent\colorbox{gray!8}{\parbox{\dimexpr\linewidth-2\fboxsep\relax}{\raggedright\ttfamily\scriptsize         im = ax.imshow(conf, cmap="Blues")}}
\par\smallskip
This Python line participates in the export, validation, or report-generation pipeline. It is executed on the host PC to prepare data for RTL or to check RTL outputs.
\paragraph{Line 369}
\noindent\colorbox{gray!8}{\parbox{\dimexpr\linewidth-2\fboxsep\relax}{\raggedright\ttfamily\scriptsize         ax.set\_xlabel("predicted")}}
\par\smallskip
This line handles predicted labels. It checks that argmax behavior in Python and RTL chooses the same class.
\paragraph{Line 370}
\noindent\colorbox{gray!8}{\parbox{\dimexpr\linewidth-2\fboxsep\relax}{\raggedright\ttfamily\scriptsize         ax.set\_ylabel("label")}}
\par\smallskip
This Python line participates in the export, validation, or report-generation pipeline. It is executed on the host PC to prepare data for RTL or to check RTL outputs.
\paragraph{Line 371}
\noindent\colorbox{gray!8}{\parbox{\dimexpr\linewidth-2\fboxsep\relax}{\raggedright\ttfamily\scriptsize         ax.set\_title("Confusion Matrix")}}
\par\smallskip
This Python line participates in the export, validation, or report-generation pipeline. It is executed on the host PC to prepare data for RTL or to check RTL outputs.
\paragraph{Line 372}
\noindent\colorbox{gray!8}{\parbox{\dimexpr\linewidth-2\fboxsep\relax}{\raggedright\ttfamily\scriptsize         ax.set\_xticks(range(NUM\_CLASSES), CLASS\_NAMES)}}
\par\smallskip
This Python line participates in the export, validation, or report-generation pipeline. It is executed on the host PC to prepare data for RTL or to check RTL outputs.
\paragraph{Line 373}
\noindent\colorbox{gray!8}{\parbox{\dimexpr\linewidth-2\fboxsep\relax}{\raggedright\ttfamily\scriptsize         ax.set\_yticks(range(NUM\_CLASSES), CLASS\_NAMES)}}
\par\smallskip
This Python line participates in the export, validation, or report-generation pipeline. It is executed on the host PC to prepare data for RTL or to check RTL outputs.
\paragraph{Line 374}
\noindent\colorbox{gray!8}{\parbox{\dimexpr\linewidth-2\fboxsep\relax}{\raggedright\ttfamily\scriptsize         for i in range(NUM\_CLASSES):}}
\par\smallskip
This Python line participates in the export, validation, or report-generation pipeline. It is executed on the host PC to prepare data for RTL or to check RTL outputs.
\paragraph{Line 375}
\noindent\colorbox{gray!8}{\parbox{\dimexpr\linewidth-2\fboxsep\relax}{\raggedright\ttfamily\scriptsize             for j in range(NUM\_CLASSES):}}
\par\smallskip
This Python line participates in the export, validation, or report-generation pipeline. It is executed on the host PC to prepare data for RTL or to check RTL outputs.
\paragraph{Line 376}
\noindent\colorbox{gray!8}{\parbox{\dimexpr\linewidth-2\fboxsep\relax}{\raggedright\ttfamily\scriptsize                 ax.text(j, i, str(conf[i, j]), ha="center", va="center", fontsize=7)}}
\par\smallskip
This Python line participates in the export, validation, or report-generation pipeline. It is executed on the host PC to prepare data for RTL or to check RTL outputs.
\paragraph{Line 377}
\noindent\colorbox{gray!8}{\parbox{\dimexpr\linewidth-2\fboxsep\relax}{\raggedright\ttfamily\scriptsize         fig.colorbar(im, ax=ax)}}
\par\smallskip
This Python line participates in the export, validation, or report-generation pipeline. It is executed on the host PC to prepare data for RTL or to check RTL outputs.
\paragraph{Line 378}
\noindent\colorbox{gray!8}{\parbox{\dimexpr\linewidth-2\fboxsep\relax}{\raggedright\ttfamily\scriptsize         save\_pdf\_page(pdf, fig, png\_dir, "07\_confusion")}}
\par\smallskip
This Python line participates in the export, validation, or report-generation pipeline. It is executed on the host PC to prepare data for RTL or to check RTL outputs.
\paragraph{Line 379}
\noindent\colorbox{gray!8}{\parbox{\dimexpr\linewidth-2\fboxsep\relax}{\raggedright\ttfamily\scriptsize \emph{blank line}}}
\par\smallskip
This blank line is not hardware or software behavior. It separates nearby declarations or operations so the file remains readable and reviewable.
\paragraph{Line 380}
\noindent\colorbox{gray!8}{\parbox{\dimexpr\linewidth-2\fboxsep\relax}{\raggedright\ttfamily\scriptsize         norm\_conf = conf / np.maximum(conf.sum(axis=1, keepdims=True), 1)}}
\par\smallskip
This performs vectorized numerical work. The fixed-point golden models use NumPy so score, margin, and update calculations are explicit and repeatable.
\paragraph{Line 381}
\noindent\colorbox{gray!8}{\parbox{\dimexpr\linewidth-2\fboxsep\relax}{\raggedright\ttfamily\scriptsize         fig, ax = plt.subplots(figsize=(8, 7))}}
\par\smallskip
This creates a debug visualization. The plot is derived from generated CSV or Vivado report data so it can be regenerated after RTL changes.
\paragraph{Line 382}
\noindent\colorbox{gray!8}{\parbox{\dimexpr\linewidth-2\fboxsep\relax}{\raggedright\ttfamily\scriptsize         im = ax.imshow(norm\_conf, cmap="viridis", vmin=0, vmax=1)}}
\par\smallskip
This Python line participates in the export, validation, or report-generation pipeline. It is executed on the host PC to prepare data for RTL or to check RTL outputs.
\paragraph{Line 383}
\noindent\colorbox{gray!8}{\parbox{\dimexpr\linewidth-2\fboxsep\relax}{\raggedright\ttfamily\scriptsize         ax.set\_xlabel("predicted")}}
\par\smallskip
This line handles predicted labels. It checks that argmax behavior in Python and RTL chooses the same class.
\paragraph{Line 384}
\noindent\colorbox{gray!8}{\parbox{\dimexpr\linewidth-2\fboxsep\relax}{\raggedright\ttfamily\scriptsize         ax.set\_ylabel("label")}}
\par\smallskip
This Python line participates in the export, validation, or report-generation pipeline. It is executed on the host PC to prepare data for RTL or to check RTL outputs.
\paragraph{Line 385}
\noindent\colorbox{gray!8}{\parbox{\dimexpr\linewidth-2\fboxsep\relax}{\raggedright\ttfamily\scriptsize         ax.set\_title("Normalized Confusion Matrix")}}
\par\smallskip
This Python line participates in the export, validation, or report-generation pipeline. It is executed on the host PC to prepare data for RTL or to check RTL outputs.
\paragraph{Line 386}
\noindent\colorbox{gray!8}{\parbox{\dimexpr\linewidth-2\fboxsep\relax}{\raggedright\ttfamily\scriptsize         ax.set\_xticks(range(NUM\_CLASSES), CLASS\_NAMES)}}
\par\smallskip
This Python line participates in the export, validation, or report-generation pipeline. It is executed on the host PC to prepare data for RTL or to check RTL outputs.
\paragraph{Line 387}
\noindent\colorbox{gray!8}{\parbox{\dimexpr\linewidth-2\fboxsep\relax}{\raggedright\ttfamily\scriptsize         ax.set\_yticks(range(NUM\_CLASSES), CLASS\_NAMES)}}
\par\smallskip
This Python line participates in the export, validation, or report-generation pipeline. It is executed on the host PC to prepare data for RTL or to check RTL outputs.
\paragraph{Line 388}
\noindent\colorbox{gray!8}{\parbox{\dimexpr\linewidth-2\fboxsep\relax}{\raggedright\ttfamily\scriptsize         fig.colorbar(im, ax=ax)}}
\par\smallskip
This Python line participates in the export, validation, or report-generation pipeline. It is executed on the host PC to prepare data for RTL or to check RTL outputs.
\paragraph{Line 389}
\noindent\colorbox{gray!8}{\parbox{\dimexpr\linewidth-2\fboxsep\relax}{\raggedright\ttfamily\scriptsize         save\_pdf\_page(pdf, fig, png\_dir, "08\_confusion\_normalized")}}
\par\smallskip
This Python line participates in the export, validation, or report-generation pipeline. It is executed on the host PC to prepare data for RTL or to check RTL outputs.
\paragraph{Line 390}
\noindent\colorbox{gray!8}{\parbox{\dimexpr\linewidth-2\fboxsep\relax}{\raggedright\ttfamily\scriptsize \emph{blank line}}}
\par\smallskip
This blank line is not hardware or software behavior. It separates nearby declarations or operations so the file remains readable and reviewable.
\paragraph{Line 391}
\noindent\colorbox{gray!8}{\parbox{\dimexpr\linewidth-2\fboxsep\relax}{\raggedright\ttfamily\scriptsize         fig, ax = plt.subplots(figsize=(11, 5))}}
\par\smallskip
This creates a debug visualization. The plot is derived from generated CSV or Vivado report data so it can be regenerated after RTL changes.
\paragraph{Line 392}
\noindent\colorbox{gray!8}{\parbox{\dimexpr\linewidth-2\fboxsep\relax}{\raggedright\ttfamily\scriptsize         x = np.arange(NUM\_CLASSES)}}
\par\smallskip
This performs vectorized numerical work. The fixed-point golden models use NumPy so score, margin, and update calculations are explicit and repeatable.
\paragraph{Line 393}
\noindent\colorbox{gray!8}{\parbox{\dimexpr\linewidth-2\fboxsep\relax}{\raggedright\ttfamily\scriptsize         ax.bar(x - 0.25, precision, width=0.25, label="precision")}}
\par\smallskip
This Python line participates in the export, validation, or report-generation pipeline. It is executed on the host PC to prepare data for RTL or to check RTL outputs.
\paragraph{Line 394}
\noindent\colorbox{gray!8}{\parbox{\dimexpr\linewidth-2\fboxsep\relax}{\raggedright\ttfamily\scriptsize         ax.bar(x, recall, width=0.25, label="recall")}}
\par\smallskip
This Python line participates in the export, validation, or report-generation pipeline. It is executed on the host PC to prepare data for RTL or to check RTL outputs.
\paragraph{Line 395}
\noindent\colorbox{gray!8}{\parbox{\dimexpr\linewidth-2\fboxsep\relax}{\raggedright\ttfamily\scriptsize         ax.bar(x + 0.25, f1, width=0.25, label="F1")}}
\par\smallskip
This Python line participates in the export, validation, or report-generation pipeline. It is executed on the host PC to prepare data for RTL or to check RTL outputs.
\paragraph{Line 396}
\noindent\colorbox{gray!8}{\parbox{\dimexpr\linewidth-2\fboxsep\relax}{\raggedright\ttfamily\scriptsize         ax.set\_xticks(x, CLASS\_NAMES)}}
\par\smallskip
This Python line participates in the export, validation, or report-generation pipeline. It is executed on the host PC to prepare data for RTL or to check RTL outputs.
\paragraph{Line 397}
\noindent\colorbox{gray!8}{\parbox{\dimexpr\linewidth-2\fboxsep\relax}{\raggedright\ttfamily\scriptsize         ax.set\_ylim(0, 1)}}
\par\smallskip
This Python line participates in the export, validation, or report-generation pipeline. It is executed on the host PC to prepare data for RTL or to check RTL outputs.
\paragraph{Line 398}
\noindent\colorbox{gray!8}{\parbox{\dimexpr\linewidth-2\fboxsep\relax}{\raggedright\ttfamily\scriptsize         ax.set\_title("Per-Class Precision / Recall / F1")}}
\par\smallskip
This Python line participates in the export, validation, or report-generation pipeline. It is executed on the host PC to prepare data for RTL or to check RTL outputs.
\paragraph{Line 399}
\noindent\colorbox{gray!8}{\parbox{\dimexpr\linewidth-2\fboxsep\relax}{\raggedright\ttfamily\scriptsize         ax.legend()}}
\par\smallskip
This Python line participates in the export, validation, or report-generation pipeline. It is executed on the host PC to prepare data for RTL or to check RTL outputs.
\paragraph{Line 400}
\noindent\colorbox{gray!8}{\parbox{\dimexpr\linewidth-2\fboxsep\relax}{\raggedright\ttfamily\scriptsize         save\_pdf\_page(pdf, fig, png\_dir, "09\_precision\_recall\_f1")}}
\par\smallskip
This Python line participates in the export, validation, or report-generation pipeline. It is executed on the host PC to prepare data for RTL or to check RTL outputs.
\paragraph{Line 401}
\noindent\colorbox{gray!8}{\parbox{\dimexpr\linewidth-2\fboxsep\relax}{\raggedright\ttfamily\scriptsize \emph{blank line}}}
\par\smallskip
This blank line is not hardware or software behavior. It separates nearby declarations or operations so the file remains readable and reviewable.
\paragraph{Line 402}
\noindent\colorbox{gray!8}{\parbox{\dimexpr\linewidth-2\fboxsep\relax}{\raggedright\ttfamily\scriptsize         fig, ax = plt.subplots(figsize=(11, 5))}}
\par\smallskip
This creates a debug visualization. The plot is derived from generated CSV or Vivado report data so it can be regenerated after RTL changes.
\paragraph{Line 403}
\noindent\colorbox{gray!8}{\parbox{\dimexpr\linewidth-2\fboxsep\relax}{\raggedright\ttfamily\scriptsize         ax.hist(margins, bins=80, alpha=0.75)}}
\par\smallskip
This Python line participates in the export, validation, or report-generation pipeline. It is executed on the host PC to prepare data for RTL or to check RTL outputs.
\paragraph{Line 404}
\noindent\colorbox{gray!8}{\parbox{\dimexpr\linewidth-2\fboxsep\relax}{\raggedright\ttfamily\scriptsize         ax.set\_title("Top-1 Minus Top-2 Score Margin Histogram")}}
\par\smallskip
This Python line participates in the export, validation, or report-generation pipeline. It is executed on the host PC to prepare data for RTL or to check RTL outputs.
\paragraph{Line 405}
\noindent\colorbox{gray!8}{\parbox{\dimexpr\linewidth-2\fboxsep\relax}{\raggedright\ttfamily\scriptsize         ax.set\_xlabel("fixed-point score margin")}}
\par\smallskip
This Python line participates in the export, validation, or report-generation pipeline. It is executed on the host PC to prepare data for RTL or to check RTL outputs.
\paragraph{Line 406}
\noindent\colorbox{gray!8}{\parbox{\dimexpr\linewidth-2\fboxsep\relax}{\raggedright\ttfamily\scriptsize         ax.set\_ylabel("samples")}}
\par\smallskip
This Python line participates in the export, validation, or report-generation pipeline. It is executed on the host PC to prepare data for RTL or to check RTL outputs.
\paragraph{Line 407}
\noindent\colorbox{gray!8}{\parbox{\dimexpr\linewidth-2\fboxsep\relax}{\raggedright\ttfamily\scriptsize         save\_pdf\_page(pdf, fig, png\_dir, "10\_top12\_margin\_hist")}}
\par\smallskip
This Python line participates in the export, validation, or report-generation pipeline. It is executed on the host PC to prepare data for RTL or to check RTL outputs.
\paragraph{Line 408}
\noindent\colorbox{gray!8}{\parbox{\dimexpr\linewidth-2\fboxsep\relax}{\raggedright\ttfamily\scriptsize \emph{blank line}}}
\par\smallskip
This blank line is not hardware or software behavior. It separates nearby declarations or operations so the file remains readable and reviewable.
\paragraph{Line 409}
\noindent\colorbox{gray!8}{\parbox{\dimexpr\linewidth-2\fboxsep\relax}{\raggedright\ttfamily\scriptsize         fig, ax = plt.subplots(figsize=(11, 5))}}
\par\smallskip
This creates a debug visualization. The plot is derived from generated CSV or Vivado report data so it can be regenerated after RTL changes.
\paragraph{Line 410}
\noindent\colorbox{gray!8}{\parbox{\dimexpr\linewidth-2\fboxsep\relax}{\raggedright\ttfamily\scriptsize         ax.scatter(sample\_idx[correct == 1], margins[correct == 1], s=3, alpha=0.35, label="correct")}}
\par\smallskip
This Python line participates in the export, validation, or report-generation pipeline. It is executed on the host PC to prepare data for RTL or to check RTL outputs.
\paragraph{Line 411}
\noindent\colorbox{gray!8}{\parbox{\dimexpr\linewidth-2\fboxsep\relax}{\raggedright\ttfamily\scriptsize         ax.scatter(sample\_idx[correct == 0], margins[correct == 0], s=8, alpha=0.75, label="wrong")}}
\par\smallskip
This Python line participates in the export, validation, or report-generation pipeline. It is executed on the host PC to prepare data for RTL or to check RTL outputs.
\paragraph{Line 412}
\noindent\colorbox{gray!8}{\parbox{\dimexpr\linewidth-2\fboxsep\relax}{\raggedright\ttfamily\scriptsize         ax.set\_xlabel("sample")}}
\par\smallskip
This Python line participates in the export, validation, or report-generation pipeline. It is executed on the host PC to prepare data for RTL or to check RTL outputs.
\paragraph{Line 413}
\noindent\colorbox{gray!8}{\parbox{\dimexpr\linewidth-2\fboxsep\relax}{\raggedright\ttfamily\scriptsize         ax.set\_ylabel("top1-top2 margin")}}
\par\smallskip
This Python line participates in the export, validation, or report-generation pipeline. It is executed on the host PC to prepare data for RTL or to check RTL outputs.
\paragraph{Line 414}
\noindent\colorbox{gray!8}{\parbox{\dimexpr\linewidth-2\fboxsep\relax}{\raggedright\ttfamily\scriptsize         ax.set\_title("Margin Timeline With Misclassification Overlay")}}
\par\smallskip
This Python line participates in the export, validation, or report-generation pipeline. It is executed on the host PC to prepare data for RTL or to check RTL outputs.
\paragraph{Line 415}
\noindent\colorbox{gray!8}{\parbox{\dimexpr\linewidth-2\fboxsep\relax}{\raggedright\ttfamily\scriptsize         ax.legend()}}
\par\smallskip
This Python line participates in the export, validation, or report-generation pipeline. It is executed on the host PC to prepare data for RTL or to check RTL outputs.
\paragraph{Line 416}
\noindent\colorbox{gray!8}{\parbox{\dimexpr\linewidth-2\fboxsep\relax}{\raggedright\ttfamily\scriptsize         save\_pdf\_page(pdf, fig, png\_dir, "11\_margin\_timeline")}}
\par\smallskip
This Python line participates in the export, validation, or report-generation pipeline. It is executed on the host PC to prepare data for RTL or to check RTL outputs.
\paragraph{Line 417}
\noindent\colorbox{gray!8}{\parbox{\dimexpr\linewidth-2\fboxsep\relax}{\raggedright\ttfamily\scriptsize \emph{blank line}}}
\par\smallskip
This blank line is not hardware or software behavior. It separates nearby declarations or operations so the file remains readable and reviewable.
\paragraph{Line 418}
\noindent\colorbox{gray!8}{\parbox{\dimexpr\linewidth-2\fboxsep\relax}{\raggedright\ttfamily\scriptsize         fig, ax = plt.subplots(figsize=(8, 5))}}
\par\smallskip
This creates a debug visualization. The plot is derived from generated CSV or Vivado report data so it can be regenerated after RTL changes.
\paragraph{Line 419}
\noindent\colorbox{gray!8}{\parbox{\dimexpr\linewidth-2\fboxsep\relax}{\raggedright\ttfamily\scriptsize         ax.boxplot([margins[correct == 1], margins[correct == 0]], tick\_labels=["correct", "wrong"], showfliers=False)}}
\par\smallskip
This Python line participates in the export, validation, or report-generation pipeline. It is executed on the host PC to prepare data for RTL or to check RTL outputs.
\paragraph{Line 420}
\noindent\colorbox{gray!8}{\parbox{\dimexpr\linewidth-2\fboxsep\relax}{\raggedright\ttfamily\scriptsize         ax.set\_ylabel("top1-top2 margin")}}
\par\smallskip
This Python line participates in the export, validation, or report-generation pipeline. It is executed on the host PC to prepare data for RTL or to check RTL outputs.
\paragraph{Line 421}
\noindent\colorbox{gray!8}{\parbox{\dimexpr\linewidth-2\fboxsep\relax}{\raggedright\ttfamily\scriptsize         ax.set\_title("Margins Split By Correctness")}}
\par\smallskip
This Python line participates in the export, validation, or report-generation pipeline. It is executed on the host PC to prepare data for RTL or to check RTL outputs.
\paragraph{Line 422}
\noindent\colorbox{gray!8}{\parbox{\dimexpr\linewidth-2\fboxsep\relax}{\raggedright\ttfamily\scriptsize         save\_pdf\_page(pdf, fig, png\_dir, "12\_margin\_box")}}
\par\smallskip
This Python line participates in the export, validation, or report-generation pipeline. It is executed on the host PC to prepare data for RTL or to check RTL outputs.
\paragraph{Line 423}
\noindent\colorbox{gray!8}{\parbox{\dimexpr\linewidth-2\fboxsep\relax}{\raggedright\ttfamily\scriptsize \emph{blank line}}}
\par\smallskip
This blank line is not hardware or software behavior. It separates nearby declarations or operations so the file remains readable and reviewable.
\paragraph{Line 424}
\noindent\colorbox{gray!8}{\parbox{\dimexpr\linewidth-2\fboxsep\relax}{\raggedright\ttfamily\scriptsize         fig, ax = plt.subplots(figsize=(12, 6))}}
\par\smallskip
This creates a debug visualization. The plot is derived from generated CSV or Vivado report data so it can be regenerated after RTL changes.
\paragraph{Line 425}
\noindent\colorbox{gray!8}{\parbox{\dimexpr\linewidth-2\fboxsep\relax}{\raggedright\ttfamily\scriptsize         im = ax.imshow(scores[:80] / 1e9, aspect="auto", cmap="coolwarm")}}
\par\smallskip
This line handles class scores. The score vector is the direct software equivalent of the RTL MAC loop over weights and features.
\paragraph{Line 426}
\noindent\colorbox{gray!8}{\parbox{\dimexpr\linewidth-2\fboxsep\relax}{\raggedright\ttfamily\scriptsize         ax.set\_xlabel("class")}}
\par\smallskip
This Python line participates in the export, validation, or report-generation pipeline. It is executed on the host PC to prepare data for RTL or to check RTL outputs.
\paragraph{Line 427}
\noindent\colorbox{gray!8}{\parbox{\dimexpr\linewidth-2\fboxsep\relax}{\raggedright\ttfamily\scriptsize         ax.set\_ylabel("sample")}}
\par\smallskip
This Python line participates in the export, validation, or report-generation pipeline. It is executed on the host PC to prepare data for RTL or to check RTL outputs.
\paragraph{Line 428}
\noindent\colorbox{gray!8}{\parbox{\dimexpr\linewidth-2\fboxsep\relax}{\raggedright\ttfamily\scriptsize         ax.set\_title("Score Heatmap: First 80 Samples (billions of fixed-point units)")}}
\par\smallskip
This Python line participates in the export, validation, or report-generation pipeline. It is executed on the host PC to prepare data for RTL or to check RTL outputs.
\paragraph{Line 429}
\noindent\colorbox{gray!8}{\parbox{\dimexpr\linewidth-2\fboxsep\relax}{\raggedright\ttfamily\scriptsize         ax.set\_xticks(range(NUM\_CLASSES), CLASS\_NAMES)}}
\par\smallskip
This Python line participates in the export, validation, or report-generation pipeline. It is executed on the host PC to prepare data for RTL or to check RTL outputs.
\paragraph{Line 430}
\noindent\colorbox{gray!8}{\parbox{\dimexpr\linewidth-2\fboxsep\relax}{\raggedright\ttfamily\scriptsize         fig.colorbar(im, ax=ax)}}
\par\smallskip
This Python line participates in the export, validation, or report-generation pipeline. It is executed on the host PC to prepare data for RTL or to check RTL outputs.
\paragraph{Line 431}
\noindent\colorbox{gray!8}{\parbox{\dimexpr\linewidth-2\fboxsep\relax}{\raggedright\ttfamily\scriptsize         save\_pdf\_page(pdf, fig, png\_dir, "13\_score\_heatmap")}}
\par\smallskip
This Python line participates in the export, validation, or report-generation pipeline. It is executed on the host PC to prepare data for RTL or to check RTL outputs.
\paragraph{Line 432}
\noindent\colorbox{gray!8}{\parbox{\dimexpr\linewidth-2\fboxsep\relax}{\raggedright\ttfamily\scriptsize \emph{blank line}}}
\par\smallskip
This blank line is not hardware or software behavior. It separates nearby declarations or operations so the file remains readable and reviewable.
\paragraph{Line 433}
\noindent\colorbox{gray!8}{\parbox{\dimexpr\linewidth-2\fboxsep\relax}{\raggedright\ttfamily\scriptsize         diff = learn\_scores - scores}}
\par\smallskip
This line handles class scores. The score vector is the direct software equivalent of the RTL MAC loop over weights and features.
\paragraph{Line 434}
\noindent\colorbox{gray!8}{\parbox{\dimexpr\linewidth-2\fboxsep\relax}{\raggedright\ttfamily\scriptsize         fig, ax = plt.subplots(figsize=(11, 5))}}
\par\smallskip
This creates a debug visualization. The plot is derived from generated CSV or Vivado report data so it can be regenerated after RTL changes.
\paragraph{Line 435}
\noindent\colorbox{gray!8}{\parbox{\dimexpr\linewidth-2\fboxsep\relax}{\raggedright\ttfamily\scriptsize         ax.hist(diff.reshape(-1), bins=80)}}
\par\smallskip
This Python line participates in the export, validation, or report-generation pipeline. It is executed on the host PC to prepare data for RTL or to check RTL outputs.
\paragraph{Line 436}
\noindent\colorbox{gray!8}{\parbox{\dimexpr\linewidth-2\fboxsep\relax}{\raggedright\ttfamily\scriptsize         ax.set\_title("Score Difference Caused By Online Weight Updates")}}
\par\smallskip
This Python line participates in the export, validation, or report-generation pipeline. It is executed on the host PC to prepare data for RTL or to check RTL outputs.
\paragraph{Line 437}
\noindent\colorbox{gray!8}{\parbox{\dimexpr\linewidth-2\fboxsep\relax}{\raggedright\ttfamily\scriptsize         ax.set\_xlabel("on-chip-learning score - fixed inference score")}}
\par\smallskip
This Python line participates in the export, validation, or report-generation pipeline. It is executed on the host PC to prepare data for RTL or to check RTL outputs.
\paragraph{Line 438}
\noindent\colorbox{gray!8}{\parbox{\dimexpr\linewidth-2\fboxsep\relax}{\raggedright\ttfamily\scriptsize         save\_pdf\_page(pdf, fig, png\_dir, "14\_learning\_score\_diff")}}
\par\smallskip
This Python line participates in the export, validation, or report-generation pipeline. It is executed on the host PC to prepare data for RTL or to check RTL outputs.
\paragraph{Line 439}
\noindent\colorbox{gray!8}{\parbox{\dimexpr\linewidth-2\fboxsep\relax}{\raggedright\ttfamily\scriptsize \emph{blank line}}}
\par\smallskip
This blank line is not hardware or software behavior. It separates nearby declarations or operations so the file remains readable and reviewable.
\paragraph{Line 440}
\noindent\colorbox{gray!8}{\parbox{\dimexpr\linewidth-2\fboxsep\relax}{\raggedright\ttfamily\scriptsize         update\_samples = sample\_idx[updates == 1]}}
\par\smallskip
This line handles the online-learning rule or its validation fields. It is central to proving that the FPGA update path is not bypassed.
\paragraph{Line 441}
\noindent\colorbox{gray!8}{\parbox{\dimexpr\linewidth-2\fboxsep\relax}{\raggedright\ttfamily\scriptsize         fig, ax = plt.subplots(figsize=(11, 5))}}
\par\smallskip
This creates a debug visualization. The plot is derived from generated CSV or Vivado report data so it can be regenerated after RTL changes.
\paragraph{Line 442}
\noindent\colorbox{gray!8}{\parbox{\dimexpr\linewidth-2\fboxsep\relax}{\raggedright\ttfamily\scriptsize         ax.vlines(update\_samples, 0, 1, color="\#c1121f")}}
\par\smallskip
This line handles the online-learning rule or its validation fields. It is central to proving that the FPGA update path is not bypassed.
\paragraph{Line 443}
\noindent\colorbox{gray!8}{\parbox{\dimexpr\linewidth-2\fboxsep\relax}{\raggedright\ttfamily\scriptsize         ax.set\_ylim(0, 1.05)}}
\par\smallskip
This Python line participates in the export, validation, or report-generation pipeline. It is executed on the host PC to prepare data for RTL or to check RTL outputs.
\paragraph{Line 444}
\noindent\colorbox{gray!8}{\parbox{\dimexpr\linewidth-2\fboxsep\relax}{\raggedright\ttfamily\scriptsize         ax.set\_xlabel("sample")}}
\par\smallskip
This Python line participates in the export, validation, or report-generation pipeline. It is executed on the host PC to prepare data for RTL or to check RTL outputs.
\paragraph{Line 445}
\noindent\colorbox{gray!8}{\parbox{\dimexpr\linewidth-2\fboxsep\relax}{\raggedright\ttfamily\scriptsize         ax.set\_yticks([])}}
\par\smallskip
This Python line participates in the export, validation, or report-generation pipeline. It is executed on the host PC to prepare data for RTL or to check RTL outputs.
\paragraph{Line 446}
\noindent\colorbox{gray!8}{\parbox{\dimexpr\linewidth-2\fboxsep\relax}{\raggedright\ttfamily\scriptsize         ax.set\_title("On-Chip Learning Update Events")}}
\par\smallskip
This Python line participates in the export, validation, or report-generation pipeline. It is executed on the host PC to prepare data for RTL or to check RTL outputs.
\paragraph{Line 447}
\noindent\colorbox{gray!8}{\parbox{\dimexpr\linewidth-2\fboxsep\relax}{\raggedright\ttfamily\scriptsize         save\_pdf\_page(pdf, fig, png\_dir, "15\_update\_events")}}
\par\smallskip
This line handles the online-learning rule or its validation fields. It is central to proving that the FPGA update path is not bypassed.
\paragraph{Line 448}
\noindent\colorbox{gray!8}{\parbox{\dimexpr\linewidth-2\fboxsep\relax}{\raggedright\ttfamily\scriptsize \emph{blank line}}}
\par\smallskip
This blank line is not hardware or software behavior. It separates nearby declarations or operations so the file remains readable and reviewable.
\paragraph{Line 449}
\noindent\colorbox{gray!8}{\parbox{\dimexpr\linewidth-2\fboxsep\relax}{\raggedright\ttfamily\scriptsize         fig, ax = plt.subplots(figsize=(11, 5))}}
\par\smallskip
This creates a debug visualization. The plot is derived from generated CSV or Vivado report data so it can be regenerated after RTL changes.
\paragraph{Line 450}
\noindent\colorbox{gray!8}{\parbox{\dimexpr\linewidth-2\fboxsep\relax}{\raggedright\ttfamily\scriptsize         ax.step(sample\_idx, update\_count, where="post")}}
\par\smallskip
This line handles the online-learning rule or its validation fields. It is central to proving that the FPGA update path is not bypassed.
\paragraph{Line 451}
\noindent\colorbox{gray!8}{\parbox{\dimexpr\linewidth-2\fboxsep\relax}{\raggedright\ttfamily\scriptsize         ax.set\_xlabel("sample")}}
\par\smallskip
This Python line participates in the export, validation, or report-generation pipeline. It is executed on the host PC to prepare data for RTL or to check RTL outputs.
\paragraph{Line 452}
\noindent\colorbox{gray!8}{\parbox{\dimexpr\linewidth-2\fboxsep\relax}{\raggedright\ttfamily\scriptsize         ax.set\_ylabel("cumulative updates")}}
\par\smallskip
This line handles the online-learning rule or its validation fields. It is central to proving that the FPGA update path is not bypassed.
\paragraph{Line 453}
\noindent\colorbox{gray!8}{\parbox{\dimexpr\linewidth-2\fboxsep\relax}{\raggedright\ttfamily\scriptsize         ax.set\_title("Hardware Update Counter")}}
\par\smallskip
This Python line participates in the export, validation, or report-generation pipeline. It is executed on the host PC to prepare data for RTL or to check RTL outputs.
\paragraph{Line 454}
\noindent\colorbox{gray!8}{\parbox{\dimexpr\linewidth-2\fboxsep\relax}{\raggedright\ttfamily\scriptsize         save\_pdf\_page(pdf, fig, png\_dir, "16\_update\_count")}}
\par\smallskip
This line handles the online-learning rule or its validation fields. It is central to proving that the FPGA update path is not bypassed.
\paragraph{Line 455}
\noindent\colorbox{gray!8}{\parbox{\dimexpr\linewidth-2\fboxsep\relax}{\raggedright\ttfamily\scriptsize \emph{blank line}}}
\par\smallskip
This blank line is not hardware or software behavior. It separates nearby declarations or operations so the file remains readable and reviewable.
\paragraph{Line 456}
\noindent\colorbox{gray!8}{\parbox{\dimexpr\linewidth-2\fboxsep\relax}{\raggedright\ttfamily\scriptsize         if "infer\_skip\_count" in inference\_rows[0]:}}
\par\smallskip
This Python line participates in the export, validation, or report-generation pipeline. It is executed on the host PC to prepare data for RTL or to check RTL outputs.
\paragraph{Line 457}
\noindent\colorbox{gray!8}{\parbox{\dimexpr\linewidth-2\fboxsep\relax}{\raggedright\ttfamily\scriptsize             fig, ax = plt.subplots(figsize=(11, 5))}}
\par\smallskip
This creates a debug visualization. The plot is derived from generated CSV or Vivado report data so it can be regenerated after RTL changes.
\paragraph{Line 458}
\noindent\colorbox{gray!8}{\parbox{\dimexpr\linewidth-2\fboxsep\relax}{\raggedright\ttfamily\scriptsize             ax.step(sample\_idx, col(inference\_rows, "infer\_skip\_count", int), where="post", label="inference run")}}
\par\smallskip
This Python line participates in the export, validation, or report-generation pipeline. It is executed on the host PC to prepare data for RTL or to check RTL outputs.
\paragraph{Line 459}
\noindent\colorbox{gray!8}{\parbox{\dimexpr\linewidth-2\fboxsep\relax}{\raggedright\ttfamily\scriptsize             if "infer\_skip\_count" in learn\_rows[0]:}}
\par\smallskip
This Python line participates in the export, validation, or report-generation pipeline. It is executed on the host PC to prepare data for RTL or to check RTL outputs.
\paragraph{Line 460}
\noindent\colorbox{gray!8}{\parbox{\dimexpr\linewidth-2\fboxsep\relax}{\raggedright\ttfamily\scriptsize                 ax.step(sample\_idx, col(learn\_rows, "infer\_skip\_count", int), where="post", label="learning run inference counter", alpha=0.75)}}
\par\smallskip
This Python line participates in the export, validation, or report-generation pipeline. It is executed on the host PC to prepare data for RTL or to check RTL outputs.
\paragraph{Line 461}
\noindent\colorbox{gray!8}{\parbox{\dimexpr\linewidth-2\fboxsep\relax}{\raggedright\ttfamily\scriptsize             ax.set\_xlabel("sample")}}
\par\smallskip
This Python line participates in the export, validation, or report-generation pipeline. It is executed on the host PC to prepare data for RTL or to check RTL outputs.
\paragraph{Line 462}
\noindent\colorbox{gray!8}{\parbox{\dimexpr\linewidth-2\fboxsep\relax}{\raggedright\ttfamily\scriptsize             ax.set\_ylabel("cumulative skipped feature-MACs")}}
\par\smallskip
This Python line participates in the export, validation, or report-generation pipeline. It is executed on the host PC to prepare data for RTL or to check RTL outputs.
\paragraph{Line 463}
\noindent\colorbox{gray!8}{\parbox{\dimexpr\linewidth-2\fboxsep\relax}{\raggedright\ttfamily\scriptsize             ax.set\_title("RTL Activity-Gate Skip Counter")}}
\par\smallskip
This Python line participates in the export, validation, or report-generation pipeline. It is executed on the host PC to prepare data for RTL or to check RTL outputs.
\paragraph{Line 464}
\noindent\colorbox{gray!8}{\parbox{\dimexpr\linewidth-2\fboxsep\relax}{\raggedright\ttfamily\scriptsize             ax.legend()}}
\par\smallskip
This Python line participates in the export, validation, or report-generation pipeline. It is executed on the host PC to prepare data for RTL or to check RTL outputs.
\paragraph{Line 465}
\noindent\colorbox{gray!8}{\parbox{\dimexpr\linewidth-2\fboxsep\relax}{\raggedright\ttfamily\scriptsize             save\_pdf\_page(pdf, fig, png\_dir, "16b\_skip\_count")}}
\par\smallskip
This Python line participates in the export, validation, or report-generation pipeline. It is executed on the host PC to prepare data for RTL or to check RTL outputs.
\paragraph{Line 466}
\noindent\colorbox{gray!8}{\parbox{\dimexpr\linewidth-2\fboxsep\relax}{\raggedright\ttfamily\scriptsize \emph{blank line}}}
\par\smallskip
This blank line is not hardware or software behavior. It separates nearby declarations or operations so the file remains readable and reviewable.
\paragraph{Line 467}
\noindent\colorbox{gray!8}{\parbox{\dimexpr\linewidth-2\fboxsep\relax}{\raggedright\ttfamily\scriptsize         if "train\_skip\_count" in learn\_rows[0]:}}
\par\smallskip
This Python line participates in the export, validation, or report-generation pipeline. It is executed on the host PC to prepare data for RTL or to check RTL outputs.
\paragraph{Line 468}
\noindent\colorbox{gray!8}{\parbox{\dimexpr\linewidth-2\fboxsep\relax}{\raggedright\ttfamily\scriptsize             fig, ax = plt.subplots(figsize=(11, 5))}}
\par\smallskip
This creates a debug visualization. The plot is derived from generated CSV or Vivado report data so it can be regenerated after RTL changes.
\paragraph{Line 469}
\noindent\colorbox{gray!8}{\parbox{\dimexpr\linewidth-2\fboxsep\relax}{\raggedright\ttfamily\scriptsize             ax.step(sample\_idx, col(learn\_rows, "train\_skip\_count", int), where="post")}}
\par\smallskip
This Python line participates in the export, validation, or report-generation pipeline. It is executed on the host PC to prepare data for RTL or to check RTL outputs.
\paragraph{Line 470}
\noindent\colorbox{gray!8}{\parbox{\dimexpr\linewidth-2\fboxsep\relax}{\raggedright\ttfamily\scriptsize             ax.set\_xlabel("sample")}}
\par\smallskip
This Python line participates in the export, validation, or report-generation pipeline. It is executed on the host PC to prepare data for RTL or to check RTL outputs.
\paragraph{Line 471}
\noindent\colorbox{gray!8}{\parbox{\dimexpr\linewidth-2\fboxsep\relax}{\raggedright\ttfamily\scriptsize             ax.set\_ylabel("cumulative skipped learning features")}}
\par\smallskip
This Python line participates in the export, validation, or report-generation pipeline. It is executed on the host PC to prepare data for RTL or to check RTL outputs.
\paragraph{Line 472}
\noindent\colorbox{gray!8}{\parbox{\dimexpr\linewidth-2\fboxsep\relax}{\raggedright\ttfamily\scriptsize             ax.set\_title("RTL On-Chip-Learning Feature Gate Counter")}}
\par\smallskip
This Python line participates in the export, validation, or report-generation pipeline. It is executed on the host PC to prepare data for RTL or to check RTL outputs.
\paragraph{Line 473}
\noindent\colorbox{gray!8}{\parbox{\dimexpr\linewidth-2\fboxsep\relax}{\raggedright\ttfamily\scriptsize             save\_pdf\_page(pdf, fig, png\_dir, "16c\_train\_skip\_count")}}
\par\smallskip
This Python line participates in the export, validation, or report-generation pipeline. It is executed on the host PC to prepare data for RTL or to check RTL outputs.
\paragraph{Line 474}
\noindent\colorbox{gray!8}{\parbox{\dimexpr\linewidth-2\fboxsep\relax}{\raggedright\ttfamily\scriptsize \emph{blank line}}}
\par\smallskip
This blank line is not hardware or software behavior. It separates nearby declarations or operations so the file remains readable and reviewable.
\paragraph{Line 475}
\noindent\colorbox{gray!8}{\parbox{\dimexpr\linewidth-2\fboxsep\relax}{\raggedright\ttfamily\scriptsize         if len(update\_samples):}}
\par\smallskip
This line handles the online-learning rule or its validation fields. It is central to proving that the FPGA update path is not bypassed.
\paragraph{Line 476}
\noindent\colorbox{gray!8}{\parbox{\dimexpr\linewidth-2\fboxsep\relax}{\raggedright\ttfamily\scriptsize             fig, ax = plt.subplots(figsize=(11, 5))}}
\par\smallskip
This creates a debug visualization. The plot is derived from generated CSV or Vivado report data so it can be regenerated after RTL changes.
\paragraph{Line 477}
\noindent\colorbox{gray!8}{\parbox{\dimexpr\linewidth-2\fboxsep\relax}{\raggedright\ttfamily\scriptsize             width = 0.35}}
\par\smallskip
This Python line participates in the export, validation, or report-generation pipeline. It is executed on the host PC to prepare data for RTL or to check RTL outputs.
\paragraph{Line 478}
\noindent\colorbox{gray!8}{\parbox{\dimexpr\linewidth-2\fboxsep\relax}{\raggedright\ttfamily\scriptsize             upd\_labels = labels[updates == 1]}}
\par\smallskip
This line handles the online-learning rule or its validation fields. It is central to proving that the FPGA update path is not bypassed.
\paragraph{Line 479}
\noindent\colorbox{gray!8}{\parbox{\dimexpr\linewidth-2\fboxsep\relax}{\raggedright\ttfamily\scriptsize             upd\_preds = learn\_preds[updates == 1]}}
\par\smallskip
This line handles predicted labels. It checks that argmax behavior in Python and RTL chooses the same class.
\paragraph{Line 480}
\noindent\colorbox{gray!8}{\parbox{\dimexpr\linewidth-2\fboxsep\relax}{\raggedright\ttfamily\scriptsize             ax.bar(np.arange(len(update\_samples)) - width / 2, upd\_labels, width, label="true label")}}
\par\smallskip
This performs vectorized numerical work. The fixed-point golden models use NumPy so score, margin, and update calculations are explicit and repeatable.
\paragraph{Line 481}
\noindent\colorbox{gray!8}{\parbox{\dimexpr\linewidth-2\fboxsep\relax}{\raggedright\ttfamily\scriptsize             ax.bar(np.arange(len(update\_samples)) + width / 2, upd\_preds, width, label="predicted")}}
\par\smallskip
This performs vectorized numerical work. The fixed-point golden models use NumPy so score, margin, and update calculations are explicit and repeatable.
\paragraph{Line 482}
\noindent\colorbox{gray!8}{\parbox{\dimexpr\linewidth-2\fboxsep\relax}{\raggedright\ttfamily\scriptsize             ax.set\_xticks(np.arange(len(update\_samples)), update\_samples, rotation=90)}}
\par\smallskip
This performs vectorized numerical work. The fixed-point golden models use NumPy so score, margin, and update calculations are explicit and repeatable.
\paragraph{Line 483}
\noindent\colorbox{gray!8}{\parbox{\dimexpr\linewidth-2\fboxsep\relax}{\raggedright\ttfamily\scriptsize             ax.set\_ylabel("class")}}
\par\smallskip
This Python line participates in the export, validation, or report-generation pipeline. It is executed on the host PC to prepare data for RTL or to check RTL outputs.
\paragraph{Line 484}
\noindent\colorbox{gray!8}{\parbox{\dimexpr\linewidth-2\fboxsep\relax}{\raggedright\ttfamily\scriptsize             ax.set\_xlabel("sample where update occurred")}}
\par\smallskip
This line handles the online-learning rule or its validation fields. It is central to proving that the FPGA update path is not bypassed.
\paragraph{Line 485}
\noindent\colorbox{gray!8}{\parbox{\dimexpr\linewidth-2\fboxsep\relax}{\raggedright\ttfamily\scriptsize             ax.set\_title("Classes Involved In Each On-Chip Update")}}
\par\smallskip
This Python line participates in the export, validation, or report-generation pipeline. It is executed on the host PC to prepare data for RTL or to check RTL outputs.
\paragraph{Line 486}
\noindent\colorbox{gray!8}{\parbox{\dimexpr\linewidth-2\fboxsep\relax}{\raggedright\ttfamily\scriptsize             ax.legend()}}
\par\smallskip
This Python line participates in the export, validation, or report-generation pipeline. It is executed on the host PC to prepare data for RTL or to check RTL outputs.
\paragraph{Line 487}
\noindent\colorbox{gray!8}{\parbox{\dimexpr\linewidth-2\fboxsep\relax}{\raggedright\ttfamily\scriptsize             save\_pdf\_page(pdf, fig, png\_dir, "17\_update\_classes")}}
\par\smallskip
This line handles the online-learning rule or its validation fields. It is central to proving that the FPGA update path is not bypassed.
\paragraph{Line 488}
\noindent\colorbox{gray!8}{\parbox{\dimexpr\linewidth-2\fboxsep\relax}{\raggedright\ttfamily\scriptsize \emph{blank line}}}
\par\smallskip
This blank line is not hardware or software behavior. It separates nearby declarations or operations so the file remains readable and reviewable.
\paragraph{Line 489}
\noindent\colorbox{gray!8}{\parbox{\dimexpr\linewidth-2\fboxsep\relax}{\raggedright\ttfamily\scriptsize         fig, ax = plt.subplots(figsize=(11, 5))}}
\par\smallskip
This creates a debug visualization. The plot is derived from generated CSV or Vivado report data so it can be regenerated after RTL changes.
\paragraph{Line 490}
\noindent\colorbox{gray!8}{\parbox{\dimexpr\linewidth-2\fboxsep\relax}{\raggedright\ttfamily\scriptsize         ax.hist(features, bins=100)}}
\par\smallskip
This Python line participates in the export, validation, or report-generation pipeline. It is executed on the host PC to prepare data for RTL or to check RTL outputs.
\paragraph{Line 491}
\noindent\colorbox{gray!8}{\parbox{\dimexpr\linewidth-2\fboxsep\relax}{\raggedright\ttfamily\scriptsize         if infer\_skip\_abs\_thresh > 0:}}
\par\smallskip
This Python line participates in the export, validation, or report-generation pipeline. It is executed on the host PC to prepare data for RTL or to check RTL outputs.
\paragraph{Line 492}
\noindent\colorbox{gray!8}{\parbox{\dimexpr\linewidth-2\fboxsep\relax}{\raggedright\ttfamily\scriptsize             ax.axvline(-infer\_skip\_abs\_thresh, color="\#c1121f", linestyle="--", label="-gate threshold")}}
\par\smallskip
This Python line participates in the export, validation, or report-generation pipeline. It is executed on the host PC to prepare data for RTL or to check RTL outputs.
\paragraph{Line 493}
\noindent\colorbox{gray!8}{\parbox{\dimexpr\linewidth-2\fboxsep\relax}{\raggedright\ttfamily\scriptsize             ax.axvline(infer\_skip\_abs\_thresh, color="\#c1121f", linestyle="--", label="+gate threshold")}}
\par\smallskip
This Python line participates in the export, validation, or report-generation pipeline. It is executed on the host PC to prepare data for RTL or to check RTL outputs.
\paragraph{Line 494}
\noindent\colorbox{gray!8}{\parbox{\dimexpr\linewidth-2\fboxsep\relax}{\raggedright\ttfamily\scriptsize             ax.legend()}}
\par\smallskip
This Python line participates in the export, validation, or report-generation pipeline. It is executed on the host PC to prepare data for RTL or to check RTL outputs.
\paragraph{Line 495}
\noindent\colorbox{gray!8}{\parbox{\dimexpr\linewidth-2\fboxsep\relax}{\raggedright\ttfamily\scriptsize         ax.set\_title("Exported Spikingformer Feature Distribution")}}
\par\smallskip
This Python line participates in the export, validation, or report-generation pipeline. It is executed on the host PC to prepare data for RTL or to check RTL outputs.
\paragraph{Line 496}
\noindent\colorbox{gray!8}{\parbox{\dimexpr\linewidth-2\fboxsep\relax}{\raggedright\ttfamily\scriptsize         ax.set\_xlabel("int16 feature value")}}
\par\smallskip
This Python line participates in the export, validation, or report-generation pipeline. It is executed on the host PC to prepare data for RTL or to check RTL outputs.
\paragraph{Line 497}
\noindent\colorbox{gray!8}{\parbox{\dimexpr\linewidth-2\fboxsep\relax}{\raggedright\ttfamily\scriptsize         save\_pdf\_page(pdf, fig, png\_dir, "18\_feature\_hist")}}
\par\smallskip
This Python line participates in the export, validation, or report-generation pipeline. It is executed on the host PC to prepare data for RTL or to check RTL outputs.
\paragraph{Line 498}
\noindent\colorbox{gray!8}{\parbox{\dimexpr\linewidth-2\fboxsep\relax}{\raggedright\ttfamily\scriptsize \emph{blank line}}}
\par\smallskip
This blank line is not hardware or software behavior. It separates nearby declarations or operations so the file remains readable and reviewable.
\paragraph{Line 499}
\noindent\colorbox{gray!8}{\parbox{\dimexpr\linewidth-2\fboxsep\relax}{\raggedright\ttfamily\scriptsize         abs\_features = np.abs(features)}}
\par\smallskip
This performs vectorized numerical work. The fixed-point golden models use NumPy so score, margin, and update calculations are explicit and repeatable.
\paragraph{Line 500}
\noindent\colorbox{gray!8}{\parbox{\dimexpr\linewidth-2\fboxsep\relax}{\raggedright\ttfamily\scriptsize         thresholds = np.array([0, 16, 32, 64, 128, 256, 512, 1024, 1536, 2048, 4096, 8192, 16384])}}
\par\smallskip
This performs vectorized numerical work. The fixed-point golden models use NumPy so score, margin, and update calculations are explicit and repeatable.
\paragraph{Line 501}
\noindent\colorbox{gray!8}{\parbox{\dimexpr\linewidth-2\fboxsep\relax}{\raggedright\ttfamily\scriptsize         skip\_fracs = np.array([(abs\_features <= t).mean() for t in thresholds], dtype=float)}}
\par\smallskip
This performs vectorized numerical work. The fixed-point golden models use NumPy so score, margin, and update calculations are explicit and repeatable.
\paragraph{Line 502}
\noindent\colorbox{gray!8}{\parbox{\dimexpr\linewidth-2\fboxsep\relax}{\raggedright\ttfamily\scriptsize         fig, ax = plt.subplots(figsize=(11, 5))}}
\par\smallskip
This creates a debug visualization. The plot is derived from generated CSV or Vivado report data so it can be regenerated after RTL changes.
\paragraph{Line 503}
\noindent\colorbox{gray!8}{\parbox{\dimexpr\linewidth-2\fboxsep\relax}{\raggedright\ttfamily\scriptsize         ax.plot(thresholds, skip\_fracs * 100.0, marker="o")}}
\par\smallskip
This Python line participates in the export, validation, or report-generation pipeline. It is executed on the host PC to prepare data for RTL or to check RTL outputs.
\paragraph{Line 504}
\noindent\colorbox{gray!8}{\parbox{\dimexpr\linewidth-2\fboxsep\relax}{\raggedright\ttfamily\scriptsize         if infer\_skip\_abs\_thresh > 0:}}
\par\smallskip
This Python line participates in the export, validation, or report-generation pipeline. It is executed on the host PC to prepare data for RTL or to check RTL outputs.
\paragraph{Line 505}
\noindent\colorbox{gray!8}{\parbox{\dimexpr\linewidth-2\fboxsep\relax}{\raggedright\ttfamily\scriptsize             ax.axvline(infer\_skip\_abs\_thresh, color="\#c1121f", linestyle="--")}}
\par\smallskip
This Python line participates in the export, validation, or report-generation pipeline. It is executed on the host PC to prepare data for RTL or to check RTL outputs.
\paragraph{Line 506}
\noindent\colorbox{gray!8}{\parbox{\dimexpr\linewidth-2\fboxsep\relax}{\raggedright\ttfamily\scriptsize             ax.text(infer\_skip\_abs\_thresh, feature\_mac\_skip\_fraction * 100.0, f"  selected \{feature\_mac\_skip\_fraction * 100:.2f\}\%", va="center")}}
\par\smallskip
This Python line participates in the export, validation, or report-generation pipeline. It is executed on the host PC to prepare data for RTL or to check RTL outputs.
\paragraph{Line 507}
\noindent\colorbox{gray!8}{\parbox{\dimexpr\linewidth-2\fboxsep\relax}{\raggedright\ttfamily\scriptsize         ax.set\_xscale("symlog", linthresh=16)}}
\par\smallskip
This Python line participates in the export, validation, or report-generation pipeline. It is executed on the host PC to prepare data for RTL or to check RTL outputs.
\paragraph{Line 508}
\noindent\colorbox{gray!8}{\parbox{\dimexpr\linewidth-2\fboxsep\relax}{\raggedright\ttfamily\scriptsize         ax.set\_xlabel("abs(feature) gate threshold")}}
\par\smallskip
This Python line participates in the export, validation, or report-generation pipeline. It is executed on the host PC to prepare data for RTL or to check RTL outputs.
\paragraph{Line 509}
\noindent\colorbox{gray!8}{\parbox{\dimexpr\linewidth-2\fboxsep\relax}{\raggedright\ttfamily\scriptsize         ax.set\_ylabel("feature-MACs skipped (\%)")}}
\par\smallskip
This Python line participates in the export, validation, or report-generation pipeline. It is executed on the host PC to prepare data for RTL or to check RTL outputs.
\paragraph{Line 510}
\noindent\colorbox{gray!8}{\parbox{\dimexpr\linewidth-2\fboxsep\relax}{\raggedright\ttfamily\scriptsize         ax.set\_title("Measured Feature Activity Gating Curve")}}
\par\smallskip
This Python line participates in the export, validation, or report-generation pipeline. It is executed on the host PC to prepare data for RTL or to check RTL outputs.
\paragraph{Line 511}
\noindent\colorbox{gray!8}{\parbox{\dimexpr\linewidth-2\fboxsep\relax}{\raggedright\ttfamily\scriptsize         save\_pdf\_page(pdf, fig, png\_dir, "18b\_feature\_gate\_sweep")}}
\par\smallskip
This Python line participates in the export, validation, or report-generation pipeline. It is executed on the host PC to prepare data for RTL or to check RTL outputs.
\paragraph{Line 512}
\noindent\colorbox{gray!8}{\parbox{\dimexpr\linewidth-2\fboxsep\relax}{\raggedright\ttfamily\scriptsize \emph{blank line}}}
\par\smallskip
This blank line is not hardware or software behavior. It separates nearby declarations or operations so the file remains readable and reviewable.
\paragraph{Line 513}
\noindent\colorbox{gray!8}{\parbox{\dimexpr\linewidth-2\fboxsep\relax}{\raggedright\ttfamily\scriptsize         first\_eff = features[:FEATURE\_DIM].copy()}}
\par\smallskip
This Python line participates in the export, validation, or report-generation pipeline. It is executed on the host PC to prepare data for RTL or to check RTL outputs.
\paragraph{Line 514}
\noindent\colorbox{gray!8}{\parbox{\dimexpr\linewidth-2\fboxsep\relax}{\raggedright\ttfamily\scriptsize         if infer\_skip\_abs\_thresh > 0:}}
\par\smallskip
This Python line participates in the export, validation, or report-generation pipeline. It is executed on the host PC to prepare data for RTL or to check RTL outputs.
\paragraph{Line 515}
\noindent\colorbox{gray!8}{\parbox{\dimexpr\linewidth-2\fboxsep\relax}{\raggedright\ttfamily\scriptsize             first\_eff[np.abs(first\_eff) <= infer\_skip\_abs\_thresh] = 0}}
\par\smallskip
This performs vectorized numerical work. The fixed-point golden models use NumPy so score, margin, and update calculations are explicit and repeatable.
\paragraph{Line 516}
\noindent\colorbox{gray!8}{\parbox{\dimexpr\linewidth-2\fboxsep\relax}{\raggedright\ttfamily\scriptsize         fig, ax = plt.subplots(figsize=(12, 5))}}
\par\smallskip
This creates a debug visualization. The plot is derived from generated CSV or Vivado report data so it can be regenerated after RTL changes.
\paragraph{Line 517}
\noindent\colorbox{gray!8}{\parbox{\dimexpr\linewidth-2\fboxsep\relax}{\raggedright\ttfamily\scriptsize         ax.plot(features[:FEATURE\_DIM], label="raw")}}
\par\smallskip
This Python line participates in the export, validation, or report-generation pipeline. It is executed on the host PC to prepare data for RTL or to check RTL outputs.
\paragraph{Line 518}
\noindent\colorbox{gray!8}{\parbox{\dimexpr\linewidth-2\fboxsep\relax}{\raggedright\ttfamily\scriptsize         ax.plot(first\_eff, label="after gate", alpha=0.8)}}
\par\smallskip
This Python line participates in the export, validation, or report-generation pipeline. It is executed on the host PC to prepare data for RTL or to check RTL outputs.
\paragraph{Line 519}
\noindent\colorbox{gray!8}{\parbox{\dimexpr\linewidth-2\fboxsep\relax}{\raggedright\ttfamily\scriptsize         ax.axhline(0, color="black", linewidth=0.6)}}
\par\smallskip
This Python line participates in the export, validation, or report-generation pipeline. It is executed on the host PC to prepare data for RTL or to check RTL outputs.
\paragraph{Line 520}
\noindent\colorbox{gray!8}{\parbox{\dimexpr\linewidth-2\fboxsep\relax}{\raggedright\ttfamily\scriptsize         ax.set\_title("First Feature Vector Before/After Activity Gate")}}
\par\smallskip
This Python line participates in the export, validation, or report-generation pipeline. It is executed on the host PC to prepare data for RTL or to check RTL outputs.
\paragraph{Line 521}
\noindent\colorbox{gray!8}{\parbox{\dimexpr\linewidth-2\fboxsep\relax}{\raggedright\ttfamily\scriptsize         ax.set\_xlabel("feature index")}}
\par\smallskip
This Python line participates in the export, validation, or report-generation pipeline. It is executed on the host PC to prepare data for RTL or to check RTL outputs.
\paragraph{Line 522}
\noindent\colorbox{gray!8}{\parbox{\dimexpr\linewidth-2\fboxsep\relax}{\raggedright\ttfamily\scriptsize         ax.set\_ylabel("int16 value")}}
\par\smallskip
This Python line participates in the export, validation, or report-generation pipeline. It is executed on the host PC to prepare data for RTL or to check RTL outputs.
\paragraph{Line 523}
\noindent\colorbox{gray!8}{\parbox{\dimexpr\linewidth-2\fboxsep\relax}{\raggedright\ttfamily\scriptsize         ax.legend()}}
\par\smallskip
This Python line participates in the export, validation, or report-generation pipeline. It is executed on the host PC to prepare data for RTL or to check RTL outputs.
\paragraph{Line 524}
\noindent\colorbox{gray!8}{\parbox{\dimexpr\linewidth-2\fboxsep\relax}{\raggedright\ttfamily\scriptsize         save\_pdf\_page(pdf, fig, png\_dir, "18c\_first\_feature\_gate\_overlay")}}
\par\smallskip
This Python line participates in the export, validation, or report-generation pipeline. It is executed on the host PC to prepare data for RTL or to check RTL outputs.
\paragraph{Line 525}
\noindent\colorbox{gray!8}{\parbox{\dimexpr\linewidth-2\fboxsep\relax}{\raggedright\ttfamily\scriptsize \emph{blank line}}}
\par\smallskip
This blank line is not hardware or software behavior. It separates nearby declarations or operations so the file remains readable and reviewable.
\paragraph{Line 526}
\noindent\colorbox{gray!8}{\parbox{\dimexpr\linewidth-2\fboxsep\relax}{\raggedright\ttfamily\scriptsize         fig, ax = plt.subplots(figsize=(12, 5))}}
\par\smallskip
This creates a debug visualization. The plot is derived from generated CSV or Vivado report data so it can be regenerated after RTL changes.
\paragraph{Line 527}
\noindent\colorbox{gray!8}{\parbox{\dimexpr\linewidth-2\fboxsep\relax}{\raggedright\ttfamily\scriptsize         ax.plot(features[:FEATURE\_DIM])}}
\par\smallskip
This Python line participates in the export, validation, or report-generation pipeline. It is executed on the host PC to prepare data for RTL or to check RTL outputs.
\paragraph{Line 528}
\noindent\colorbox{gray!8}{\parbox{\dimexpr\linewidth-2\fboxsep\relax}{\raggedright\ttfamily\scriptsize         ax.set\_title("First Exported 384-D Feature Vector")}}
\par\smallskip
This Python line participates in the export, validation, or report-generation pipeline. It is executed on the host PC to prepare data for RTL or to check RTL outputs.
\paragraph{Line 529}
\noindent\colorbox{gray!8}{\parbox{\dimexpr\linewidth-2\fboxsep\relax}{\raggedright\ttfamily\scriptsize         ax.set\_xlabel("feature index")}}
\par\smallskip
This Python line participates in the export, validation, or report-generation pipeline. It is executed on the host PC to prepare data for RTL or to check RTL outputs.
\paragraph{Line 530}
\noindent\colorbox{gray!8}{\parbox{\dimexpr\linewidth-2\fboxsep\relax}{\raggedright\ttfamily\scriptsize         ax.set\_ylabel("int16 value")}}
\par\smallskip
This Python line participates in the export, validation, or report-generation pipeline. It is executed on the host PC to prepare data for RTL or to check RTL outputs.
\paragraph{Line 531}
\noindent\colorbox{gray!8}{\parbox{\dimexpr\linewidth-2\fboxsep\relax}{\raggedright\ttfamily\scriptsize         save\_pdf\_page(pdf, fig, png\_dir, "19\_first\_feature\_vector")}}
\par\smallskip
This Python line participates in the export, validation, or report-generation pipeline. It is executed on the host PC to prepare data for RTL or to check RTL outputs.
\paragraph{Line 532}
\noindent\colorbox{gray!8}{\parbox{\dimexpr\linewidth-2\fboxsep\relax}{\raggedright\ttfamily\scriptsize \emph{blank line}}}
\par\smallskip
This blank line is not hardware or software behavior. It separates nearby declarations or operations so the file remains readable and reviewable.
\paragraph{Line 533}
\noindent\colorbox{gray!8}{\parbox{\dimexpr\linewidth-2\fboxsep\relax}{\raggedright\ttfamily\scriptsize         fig, ax = plt.subplots(figsize=(12, 5))}}
\par\smallskip
This creates a debug visualization. The plot is derived from generated CSV or Vivado report data so it can be regenerated after RTL changes.
\paragraph{Line 534}
\noindent\colorbox{gray!8}{\parbox{\dimexpr\linewidth-2\fboxsep\relax}{\raggedright\ttfamily\scriptsize         im = ax.imshow(features[:FEATURE\_DIM].reshape(16, 24), aspect="auto", cmap="coolwarm")}}
\par\smallskip
This Python line participates in the export, validation, or report-generation pipeline. It is executed on the host PC to prepare data for RTL or to check RTL outputs.
\paragraph{Line 535}
\noindent\colorbox{gray!8}{\parbox{\dimexpr\linewidth-2\fboxsep\relax}{\raggedright\ttfamily\scriptsize         ax.set\_title("First Feature Vector Reshaped 16 x 24 For Inspection")}}
\par\smallskip
This Python line participates in the export, validation, or report-generation pipeline. It is executed on the host PC to prepare data for RTL or to check RTL outputs.
\paragraph{Line 536}
\noindent\colorbox{gray!8}{\parbox{\dimexpr\linewidth-2\fboxsep\relax}{\raggedright\ttfamily\scriptsize         fig.colorbar(im, ax=ax)}}
\par\smallskip
This Python line participates in the export, validation, or report-generation pipeline. It is executed on the host PC to prepare data for RTL or to check RTL outputs.
\paragraph{Line 537}
\noindent\colorbox{gray!8}{\parbox{\dimexpr\linewidth-2\fboxsep\relax}{\raggedright\ttfamily\scriptsize         save\_pdf\_page(pdf, fig, png\_dir, "20\_first\_feature\_grid")}}
\par\smallskip
This Python line participates in the export, validation, or report-generation pipeline. It is executed on the host PC to prepare data for RTL or to check RTL outputs.
\paragraph{Line 538}
\noindent\colorbox{gray!8}{\parbox{\dimexpr\linewidth-2\fboxsep\relax}{\raggedright\ttfamily\scriptsize \emph{blank line}}}
\par\smallskip
This blank line is not hardware or software behavior. It separates nearby declarations or operations so the file remains readable and reviewable.
\paragraph{Line 539}
\noindent\colorbox{gray!8}{\parbox{\dimexpr\linewidth-2\fboxsep\relax}{\raggedright\ttfamily\scriptsize         fig, ax = plt.subplots(figsize=(11, 5))}}
\par\smallskip
This creates a debug visualization. The plot is derived from generated CSV or Vivado report data so it can be regenerated after RTL changes.
\paragraph{Line 540}
\noindent\colorbox{gray!8}{\parbox{\dimexpr\linewidth-2\fboxsep\relax}{\raggedright\ttfamily\scriptsize         ax.hist(weights.reshape(-1), bins=100)}}
\par\smallskip
This line handles the classifier parameters. Those arrays mirror the RTL memories, so mismatches reveal an export, quantization, or hardware-update bug.
\paragraph{Line 541}
\noindent\colorbox{gray!8}{\parbox{\dimexpr\linewidth-2\fboxsep\relax}{\raggedright\ttfamily\scriptsize         ax.set\_title("Classifier Weight Distribution Before On-Chip Updates")}}
\par\smallskip
This Python line participates in the export, validation, or report-generation pipeline. It is executed on the host PC to prepare data for RTL or to check RTL outputs.
\paragraph{Line 542}
\noindent\colorbox{gray!8}{\parbox{\dimexpr\linewidth-2\fboxsep\relax}{\raggedright\ttfamily\scriptsize         ax.set\_xlabel("int16 weight")}}
\par\smallskip
This Python line participates in the export, validation, or report-generation pipeline. It is executed on the host PC to prepare data for RTL or to check RTL outputs.
\paragraph{Line 543}
\noindent\colorbox{gray!8}{\parbox{\dimexpr\linewidth-2\fboxsep\relax}{\raggedright\ttfamily\scriptsize         save\_pdf\_page(pdf, fig, png\_dir, "21\_weight\_hist")}}
\par\smallskip
This Python line participates in the export, validation, or report-generation pipeline. It is executed on the host PC to prepare data for RTL or to check RTL outputs.
\paragraph{Line 544}
\noindent\colorbox{gray!8}{\parbox{\dimexpr\linewidth-2\fboxsep\relax}{\raggedright\ttfamily\scriptsize \emph{blank line}}}
\par\smallskip
This blank line is not hardware or software behavior. It separates nearby declarations or operations so the file remains readable and reviewable.
\paragraph{Line 545}
\noindent\colorbox{gray!8}{\parbox{\dimexpr\linewidth-2\fboxsep\relax}{\raggedright\ttfamily\scriptsize         fig, ax = plt.subplots(figsize=(12, 5))}}
\par\smallskip
This creates a debug visualization. The plot is derived from generated CSV or Vivado report data so it can be regenerated after RTL changes.
\paragraph{Line 546}
\noindent\colorbox{gray!8}{\parbox{\dimexpr\linewidth-2\fboxsep\relax}{\raggedright\ttfamily\scriptsize         im = ax.imshow(weights, aspect="auto", cmap="coolwarm")}}
\par\smallskip
This line handles the classifier parameters. Those arrays mirror the RTL memories, so mismatches reveal an export, quantization, or hardware-update bug.
\paragraph{Line 547}
\noindent\colorbox{gray!8}{\parbox{\dimexpr\linewidth-2\fboxsep\relax}{\raggedright\ttfamily\scriptsize         ax.set\_xlabel("feature index")}}
\par\smallskip
This Python line participates in the export, validation, or report-generation pipeline. It is executed on the host PC to prepare data for RTL or to check RTL outputs.
\paragraph{Line 548}
\noindent\colorbox{gray!8}{\parbox{\dimexpr\linewidth-2\fboxsep\relax}{\raggedright\ttfamily\scriptsize         ax.set\_ylabel("class")}}
\par\smallskip
This Python line participates in the export, validation, or report-generation pipeline. It is executed on the host PC to prepare data for RTL or to check RTL outputs.
\paragraph{Line 549}
\noindent\colorbox{gray!8}{\parbox{\dimexpr\linewidth-2\fboxsep\relax}{\raggedright\ttfamily\scriptsize         ax.set\_title("Classifier Weight Matrix: Class x Feature")}}
\par\smallskip
This Python line participates in the export, validation, or report-generation pipeline. It is executed on the host PC to prepare data for RTL or to check RTL outputs.
\paragraph{Line 550}
\noindent\colorbox{gray!8}{\parbox{\dimexpr\linewidth-2\fboxsep\relax}{\raggedright\ttfamily\scriptsize         fig.colorbar(im, ax=ax)}}
\par\smallskip
This Python line participates in the export, validation, or report-generation pipeline. It is executed on the host PC to prepare data for RTL or to check RTL outputs.
\paragraph{Line 551}
\noindent\colorbox{gray!8}{\parbox{\dimexpr\linewidth-2\fboxsep\relax}{\raggedright\ttfamily\scriptsize         save\_pdf\_page(pdf, fig, png\_dir, "22\_weight\_heatmap")}}
\par\smallskip
This Python line participates in the export, validation, or report-generation pipeline. It is executed on the host PC to prepare data for RTL or to check RTL outputs.
\paragraph{Line 552}
\noindent\colorbox{gray!8}{\parbox{\dimexpr\linewidth-2\fboxsep\relax}{\raggedright\ttfamily\scriptsize \emph{blank line}}}
\par\smallskip
This blank line is not hardware or software behavior. It separates nearby declarations or operations so the file remains readable and reviewable.
\paragraph{Line 553}
\noindent\colorbox{gray!8}{\parbox{\dimexpr\linewidth-2\fboxsep\relax}{\raggedright\ttfamily\scriptsize         fig, ax = plt.subplots(figsize=(11, 5))}}
\par\smallskip
This creates a debug visualization. The plot is derived from generated CSV or Vivado report data so it can be regenerated after RTL changes.
\paragraph{Line 554}
\noindent\colorbox{gray!8}{\parbox{\dimexpr\linewidth-2\fboxsep\relax}{\raggedright\ttfamily\scriptsize         ax.bar(np.arange(NUM\_CLASSES), np.linalg.norm(weights, axis=1))}}
\par\smallskip
This performs vectorized numerical work. The fixed-point golden models use NumPy so score, margin, and update calculations are explicit and repeatable.
\paragraph{Line 555}
\noindent\colorbox{gray!8}{\parbox{\dimexpr\linewidth-2\fboxsep\relax}{\raggedright\ttfamily\scriptsize         ax.set\_xticks(range(NUM\_CLASSES), CLASS\_NAMES)}}
\par\smallskip
This Python line participates in the export, validation, or report-generation pipeline. It is executed on the host PC to prepare data for RTL or to check RTL outputs.
\paragraph{Line 556}
\noindent\colorbox{gray!8}{\parbox{\dimexpr\linewidth-2\fboxsep\relax}{\raggedright\ttfamily\scriptsize         ax.set\_ylabel("L2 norm")}}
\par\smallskip
This Python line participates in the export, validation, or report-generation pipeline. It is executed on the host PC to prepare data for RTL or to check RTL outputs.
\paragraph{Line 557}
\noindent\colorbox{gray!8}{\parbox{\dimexpr\linewidth-2\fboxsep\relax}{\raggedright\ttfamily\scriptsize         ax.set\_title("Per-Class Weight Norm")}}
\par\smallskip
This Python line participates in the export, validation, or report-generation pipeline. It is executed on the host PC to prepare data for RTL or to check RTL outputs.
\paragraph{Line 558}
\noindent\colorbox{gray!8}{\parbox{\dimexpr\linewidth-2\fboxsep\relax}{\raggedright\ttfamily\scriptsize         save\_pdf\_page(pdf, fig, png\_dir, "23\_weight\_norms")}}
\par\smallskip
This Python line participates in the export, validation, or report-generation pipeline. It is executed on the host PC to prepare data for RTL or to check RTL outputs.
\paragraph{Line 559}
\noindent\colorbox{gray!8}{\parbox{\dimexpr\linewidth-2\fboxsep\relax}{\raggedright\ttfamily\scriptsize \emph{blank line}}}
\par\smallskip
This blank line is not hardware or software behavior. It separates nearby declarations or operations so the file remains readable and reviewable.
\paragraph{Line 560}
\noindent\colorbox{gray!8}{\parbox{\dimexpr\linewidth-2\fboxsep\relax}{\raggedright\ttfamily\scriptsize         fig, ax = plt.subplots(figsize=(11, 5))}}
\par\smallskip
This creates a debug visualization. The plot is derived from generated CSV or Vivado report data so it can be regenerated after RTL changes.
\paragraph{Line 561}
\noindent\colorbox{gray!8}{\parbox{\dimexpr\linewidth-2\fboxsep\relax}{\raggedright\ttfamily\scriptsize         ax.bar(np.arange(NUM\_CLASSES) - 0.18, bias / 1e9, width=0.36, label="before")}}
\par\smallskip
This performs vectorized numerical work. The fixed-point golden models use NumPy so score, margin, and update calculations are explicit and repeatable.
\paragraph{Line 562}
\noindent\colorbox{gray!8}{\parbox{\dimexpr\linewidth-2\fboxsep\relax}{\raggedright\ttfamily\scriptsize         ax.bar(np.arange(NUM\_CLASSES) + 0.18, after\_bias / 1e9, width=0.36, label="after")}}
\par\smallskip
This performs vectorized numerical work. The fixed-point golden models use NumPy so score, margin, and update calculations are explicit and repeatable.
\paragraph{Line 563}
\noindent\colorbox{gray!8}{\parbox{\dimexpr\linewidth-2\fboxsep\relax}{\raggedright\ttfamily\scriptsize         ax.set\_xticks(range(NUM\_CLASSES), CLASS\_NAMES)}}
\par\smallskip
This Python line participates in the export, validation, or report-generation pipeline. It is executed on the host PC to prepare data for RTL or to check RTL outputs.
\paragraph{Line 564}
\noindent\colorbox{gray!8}{\parbox{\dimexpr\linewidth-2\fboxsep\relax}{\raggedright\ttfamily\scriptsize         ax.set\_ylabel("bias / 1e9")}}
\par\smallskip
This line handles the classifier parameters. Those arrays mirror the RTL memories, so mismatches reveal an export, quantization, or hardware-update bug.
\paragraph{Line 565}
\noindent\colorbox{gray!8}{\parbox{\dimexpr\linewidth-2\fboxsep\relax}{\raggedright\ttfamily\scriptsize         ax.set\_title("Classifier Bias Before/After On-Chip Updates")}}
\par\smallskip
This Python line participates in the export, validation, or report-generation pipeline. It is executed on the host PC to prepare data for RTL or to check RTL outputs.
\paragraph{Line 566}
\noindent\colorbox{gray!8}{\parbox{\dimexpr\linewidth-2\fboxsep\relax}{\raggedright\ttfamily\scriptsize         ax.legend()}}
\par\smallskip
This Python line participates in the export, validation, or report-generation pipeline. It is executed on the host PC to prepare data for RTL or to check RTL outputs.
\paragraph{Line 567}
\noindent\colorbox{gray!8}{\parbox{\dimexpr\linewidth-2\fboxsep\relax}{\raggedright\ttfamily\scriptsize         save\_pdf\_page(pdf, fig, png\_dir, "24\_bias\_before\_after")}}
\par\smallskip
This line handles the classifier parameters. Those arrays mirror the RTL memories, so mismatches reveal an export, quantization, or hardware-update bug.
\paragraph{Line 568}
\noindent\colorbox{gray!8}{\parbox{\dimexpr\linewidth-2\fboxsep\relax}{\raggedright\ttfamily\scriptsize \emph{blank line}}}
\par\smallskip
This blank line is not hardware or software behavior. It separates nearby declarations or operations so the file remains readable and reviewable.
\paragraph{Line 569}
\noindent\colorbox{gray!8}{\parbox{\dimexpr\linewidth-2\fboxsep\relax}{\raggedright\ttfamily\scriptsize         wdiff = after\_weights - weights}}
\par\smallskip
This line handles the classifier parameters. Those arrays mirror the RTL memories, so mismatches reveal an export, quantization, or hardware-update bug.
\paragraph{Line 570}
\noindent\colorbox{gray!8}{\parbox{\dimexpr\linewidth-2\fboxsep\relax}{\raggedright\ttfamily\scriptsize         fig, ax = plt.subplots(figsize=(11, 5))}}
\par\smallskip
This creates a debug visualization. The plot is derived from generated CSV or Vivado report data so it can be regenerated after RTL changes.
\paragraph{Line 571}
\noindent\colorbox{gray!8}{\parbox{\dimexpr\linewidth-2\fboxsep\relax}{\raggedright\ttfamily\scriptsize         ax.hist(wdiff.reshape(-1), bins=np.arange(wdiff.min() - 1, wdiff.max() + 2) if wdiff.size and wdiff.max() - wdiff.min() < 200 else 100)}}
\par\smallskip
This performs vectorized numerical work. The fixed-point golden models use NumPy so score, margin, and update calculations are explicit and repeatable.
\paragraph{Line 572}
\noindent\colorbox{gray!8}{\parbox{\dimexpr\linewidth-2\fboxsep\relax}{\raggedright\ttfamily\scriptsize         ax.set\_title("Weight Delta Distribution After Online Learning")}}
\par\smallskip
This Python line participates in the export, validation, or report-generation pipeline. It is executed on the host PC to prepare data for RTL or to check RTL outputs.
\paragraph{Line 573}
\noindent\colorbox{gray!8}{\parbox{\dimexpr\linewidth-2\fboxsep\relax}{\raggedright\ttfamily\scriptsize         ax.set\_xlabel("after - before")}}
\par\smallskip
This Python line participates in the export, validation, or report-generation pipeline. It is executed on the host PC to prepare data for RTL or to check RTL outputs.
\paragraph{Line 574}
\noindent\colorbox{gray!8}{\parbox{\dimexpr\linewidth-2\fboxsep\relax}{\raggedright\ttfamily\scriptsize         save\_pdf\_page(pdf, fig, png\_dir, "25\_weight\_delta\_hist")}}
\par\smallskip
This Python line participates in the export, validation, or report-generation pipeline. It is executed on the host PC to prepare data for RTL or to check RTL outputs.
\paragraph{Line 575}
\noindent\colorbox{gray!8}{\parbox{\dimexpr\linewidth-2\fboxsep\relax}{\raggedright\ttfamily\scriptsize \emph{blank line}}}
\par\smallskip
This blank line is not hardware or software behavior. It separates nearby declarations or operations so the file remains readable and reviewable.
\paragraph{Line 576}
\noindent\colorbox{gray!8}{\parbox{\dimexpr\linewidth-2\fboxsep\relax}{\raggedright\ttfamily\scriptsize         fig, ax = plt.subplots(figsize=(12, 5))}}
\par\smallskip
This creates a debug visualization. The plot is derived from generated CSV or Vivado report data so it can be regenerated after RTL changes.
\paragraph{Line 577}
\noindent\colorbox{gray!8}{\parbox{\dimexpr\linewidth-2\fboxsep\relax}{\raggedright\ttfamily\scriptsize         im = ax.imshow(wdiff, aspect="auto", cmap="coolwarm")}}
\par\smallskip
This Python line participates in the export, validation, or report-generation pipeline. It is executed on the host PC to prepare data for RTL or to check RTL outputs.
\paragraph{Line 578}
\noindent\colorbox{gray!8}{\parbox{\dimexpr\linewidth-2\fboxsep\relax}{\raggedright\ttfamily\scriptsize         ax.set\_xlabel("feature index")}}
\par\smallskip
This Python line participates in the export, validation, or report-generation pipeline. It is executed on the host PC to prepare data for RTL or to check RTL outputs.
\paragraph{Line 579}
\noindent\colorbox{gray!8}{\parbox{\dimexpr\linewidth-2\fboxsep\relax}{\raggedright\ttfamily\scriptsize         ax.set\_ylabel("class")}}
\par\smallskip
This Python line participates in the export, validation, or report-generation pipeline. It is executed on the host PC to prepare data for RTL or to check RTL outputs.
\paragraph{Line 580}
\noindent\colorbox{gray!8}{\parbox{\dimexpr\linewidth-2\fboxsep\relax}{\raggedright\ttfamily\scriptsize         ax.set\_title("Weight Delta Matrix After On-Chip Updates")}}
\par\smallskip
This Python line participates in the export, validation, or report-generation pipeline. It is executed on the host PC to prepare data for RTL or to check RTL outputs.
\paragraph{Line 581}
\noindent\colorbox{gray!8}{\parbox{\dimexpr\linewidth-2\fboxsep\relax}{\raggedright\ttfamily\scriptsize         fig.colorbar(im, ax=ax)}}
\par\smallskip
This Python line participates in the export, validation, or report-generation pipeline. It is executed on the host PC to prepare data for RTL or to check RTL outputs.
\paragraph{Line 582}
\noindent\colorbox{gray!8}{\parbox{\dimexpr\linewidth-2\fboxsep\relax}{\raggedright\ttfamily\scriptsize         save\_pdf\_page(pdf, fig, png\_dir, "26\_weight\_delta\_heatmap")}}
\par\smallskip
This Python line participates in the export, validation, or report-generation pipeline. It is executed on the host PC to prepare data for RTL or to check RTL outputs.
\paragraph{Line 583}
\noindent\colorbox{gray!8}{\parbox{\dimexpr\linewidth-2\fboxsep\relax}{\raggedright\ttfamily\scriptsize \emph{blank line}}}
\par\smallskip
This blank line is not hardware or software behavior. It separates nearby declarations or operations so the file remains readable and reviewable.
\paragraph{Line 584}
\noindent\colorbox{gray!8}{\parbox{\dimexpr\linewidth-2\fboxsep\relax}{\raggedright\ttfamily\scriptsize         fig, ax = plt.subplots(figsize=(11, 5))}}
\par\smallskip
This creates a debug visualization. The plot is derived from generated CSV or Vivado report data so it can be regenerated after RTL changes.
\paragraph{Line 585}
\noindent\colorbox{gray!8}{\parbox{\dimexpr\linewidth-2\fboxsep\relax}{\raggedright\ttfamily\scriptsize         names = ["LUT", "FF", "BRAM36", "DSP", "LUTRAM"]}}
\par\smallskip
This Python line participates in the export, validation, or report-generation pipeline. It is executed on the host PC to prepare data for RTL or to check RTL outputs.
\paragraph{Line 586}
\noindent\colorbox{gray!8}{\parbox{\dimexpr\linewidth-2\fboxsep\relax}{\raggedright\ttfamily\scriptsize         vals = [vivado["lut"], vivado["regs"], vivado["bram"], vivado["dsp"], vivado["lut\_mem"]]}}
\par\smallskip
This Python line participates in the export, validation, or report-generation pipeline. It is executed on the host PC to prepare data for RTL or to check RTL outputs.
\paragraph{Line 587}
\noindent\colorbox{gray!8}{\parbox{\dimexpr\linewidth-2\fboxsep\relax}{\raggedright\ttfamily\scriptsize         ax.bar(names, vals)}}
\par\smallskip
This Python line participates in the export, validation, or report-generation pipeline. It is executed on the host PC to prepare data for RTL or to check RTL outputs.
\paragraph{Line 588}
\noindent\colorbox{gray!8}{\parbox{\dimexpr\linewidth-2\fboxsep\relax}{\raggedright\ttfamily\scriptsize         ax.set\_title("Vivado Utilization Counts")}}
\par\smallskip
This Python line participates in the export, validation, or report-generation pipeline. It is executed on the host PC to prepare data for RTL or to check RTL outputs.
\paragraph{Line 589}
\noindent\colorbox{gray!8}{\parbox{\dimexpr\linewidth-2\fboxsep\relax}{\raggedright\ttfamily\scriptsize         for i, value in enumerate(vals):}}
\par\smallskip
This Python line participates in the export, validation, or report-generation pipeline. It is executed on the host PC to prepare data for RTL or to check RTL outputs.
\paragraph{Line 590}
\noindent\colorbox{gray!8}{\parbox{\dimexpr\linewidth-2\fboxsep\relax}{\raggedright\ttfamily\scriptsize             ax.text(i, value, str(value), ha="center", va="bottom")}}
\par\smallskip
This Python line participates in the export, validation, or report-generation pipeline. It is executed on the host PC to prepare data for RTL or to check RTL outputs.
\paragraph{Line 591}
\noindent\colorbox{gray!8}{\parbox{\dimexpr\linewidth-2\fboxsep\relax}{\raggedright\ttfamily\scriptsize         save\_pdf\_page(pdf, fig, png\_dir, "27\_utilization")}}
\par\smallskip
This Python line participates in the export, validation, or report-generation pipeline. It is executed on the host PC to prepare data for RTL or to check RTL outputs.
\paragraph{Line 592}
\noindent\colorbox{gray!8}{\parbox{\dimexpr\linewidth-2\fboxsep\relax}{\raggedright\ttfamily\scriptsize \emph{blank line}}}
\par\smallskip
This blank line is not hardware or software behavior. It separates nearby declarations or operations so the file remains readable and reviewable.
\paragraph{Line 593}
\noindent\colorbox{gray!8}{\parbox{\dimexpr\linewidth-2\fboxsep\relax}{\raggedright\ttfamily\scriptsize         fig, ax = plt.subplots(figsize=(8, 6))}}
\par\smallskip
This creates a debug visualization. The plot is derived from generated CSV or Vivado report data so it can be regenerated after RTL changes.
\paragraph{Line 594}
\noindent\colorbox{gray!8}{\parbox{\dimexpr\linewidth-2\fboxsep\relax}{\raggedright\ttfamily\scriptsize         power\_names = ["clock", "logic", "signals", "BRAM", "DSP", "IO"]}}
\par\smallskip
This Python line participates in the export, validation, or report-generation pipeline. It is executed on the host PC to prepare data for RTL or to check RTL outputs.
\paragraph{Line 595}
\noindent\colorbox{gray!8}{\parbox{\dimexpr\linewidth-2\fboxsep\relax}{\raggedright\ttfamily\scriptsize         power\_vals = [vivado["clock\_power\_w"], vivado["logic\_power\_w"], vivado["signal\_power\_w"], vivado["bram\_power\_w"], vivado["dsp\_power\_w"], vivado["io\_power\_w"]]}}
\par\smallskip
This Python line participates in the export, validation, or report-generation pipeline. It is executed on the host PC to prepare data for RTL or to check RTL outputs.
\paragraph{Line 596}
\noindent\colorbox{gray!8}{\parbox{\dimexpr\linewidth-2\fboxsep\relax}{\raggedright\ttfamily\scriptsize         ax.pie(power\_vals, labels=power\_names, autopct="\%1.1f\%\%")}}
\par\smallskip
This Python line participates in the export, validation, or report-generation pipeline. It is executed on the host PC to prepare data for RTL or to check RTL outputs.
\paragraph{Line 597}
\noindent\colorbox{gray!8}{\parbox{\dimexpr\linewidth-2\fboxsep\relax}{\raggedright\ttfamily\scriptsize         ax.set\_title("Vivado Dynamic Power Breakdown")}}
\par\smallskip
This Python line participates in the export, validation, or report-generation pipeline. It is executed on the host PC to prepare data for RTL or to check RTL outputs.
\paragraph{Line 598}
\noindent\colorbox{gray!8}{\parbox{\dimexpr\linewidth-2\fboxsep\relax}{\raggedright\ttfamily\scriptsize         save\_pdf\_page(pdf, fig, png\_dir, "28\_power\_pie")}}
\par\smallskip
This Python line participates in the export, validation, or report-generation pipeline. It is executed on the host PC to prepare data for RTL or to check RTL outputs.
\paragraph{Line 599}
\noindent\colorbox{gray!8}{\parbox{\dimexpr\linewidth-2\fboxsep\relax}{\raggedright\ttfamily\scriptsize \emph{blank line}}}
\par\smallskip
This blank line is not hardware or software behavior. It separates nearby declarations or operations so the file remains readable and reviewable.
\paragraph{Line 600}
\noindent\colorbox{gray!8}{\parbox{\dimexpr\linewidth-2\fboxsep\relax}{\raggedright\ttfamily\scriptsize         fig, ax = plt.subplots(figsize=(9, 5))}}
\par\smallskip
This creates a debug visualization. The plot is derived from generated CSV or Vivado report data so it can be regenerated after RTL changes.
\paragraph{Line 601}
\noindent\colorbox{gray!8}{\parbox{\dimexpr\linewidth-2\fboxsep\relax}{\raggedright\ttfamily\scriptsize         ax.bar(["WNS", "WHS"], [vivado["wns\_ns"], vivado["whs\_ns"]])}}
\par\smallskip
This Python line participates in the export, validation, or report-generation pipeline. It is executed on the host PC to prepare data for RTL or to check RTL outputs.
\paragraph{Line 602}
\noindent\colorbox{gray!8}{\parbox{\dimexpr\linewidth-2\fboxsep\relax}{\raggedright\ttfamily\scriptsize         ax.axhline(0, color="black", linewidth=0.8)}}
\par\smallskip
This Python line participates in the export, validation, or report-generation pipeline. It is executed on the host PC to prepare data for RTL or to check RTL outputs.
\paragraph{Line 603}
\noindent\colorbox{gray!8}{\parbox{\dimexpr\linewidth-2\fboxsep\relax}{\raggedright\ttfamily\scriptsize         ax.set\_ylabel("ns")}}
\par\smallskip
This Python line participates in the export, validation, or report-generation pipeline. It is executed on the host PC to prepare data for RTL or to check RTL outputs.
\paragraph{Line 604}
\noindent\colorbox{gray!8}{\parbox{\dimexpr\linewidth-2\fboxsep\relax}{\raggedright\ttfamily\scriptsize         ax.set\_title("Routed Timing Slack At 100 MHz")}}
\par\smallskip
This Python line participates in the export, validation, or report-generation pipeline. It is executed on the host PC to prepare data for RTL or to check RTL outputs.
\paragraph{Line 605}
\noindent\colorbox{gray!8}{\parbox{\dimexpr\linewidth-2\fboxsep\relax}{\raggedright\ttfamily\scriptsize         save\_pdf\_page(pdf, fig, png\_dir, "29\_timing\_slack")}}
\par\smallskip
This Python line participates in the export, validation, or report-generation pipeline. It is executed on the host PC to prepare data for RTL or to check RTL outputs.
\paragraph{Line 606}
\noindent\colorbox{gray!8}{\parbox{\dimexpr\linewidth-2\fboxsep\relax}{\raggedright\ttfamily\scriptsize \emph{blank line}}}
\par\smallskip
This blank line is not hardware or software behavior. It separates nearby declarations or operations so the file remains readable and reviewable.
\paragraph{Line 607}
\noindent\colorbox{gray!8}{\parbox{\dimexpr\linewidth-2\fboxsep\relax}{\raggedright\ttfamily\scriptsize         saved\_infer\_cycles = 2 * rtl\_infer\_skip\_count}}
\par\smallskip
This Python line participates in the export, validation, or report-generation pipeline. It is executed on the host PC to prepare data for RTL or to check RTL outputs.
\paragraph{Line 608}
\noindent\colorbox{gray!8}{\parbox{\dimexpr\linewidth-2\fboxsep\relax}{\raggedright\ttfamily\scriptsize         saved\_train\_cycles = 5 * learn\_train\_skip\_count}}
\par\smallskip
This Python line participates in the export, validation, or report-generation pipeline. It is executed on the host PC to prepare data for RTL or to check RTL outputs.
\paragraph{Line 609}
\noindent\colorbox{gray!8}{\parbox{\dimexpr\linewidth-2\fboxsep\relax}{\raggedright\ttfamily\scriptsize         latency\_lines = [}}
\par\smallskip
This Python line participates in the export, validation, or report-generation pipeline. It is executed on the host PC to prepare data for RTL or to check RTL outputs.
\paragraph{Line 610}
\noindent\colorbox{gray!8}{\parbox{\dimexpr\linewidth-2\fboxsep\relax}{\raggedright\ttfamily\scriptsize             "Feature load: 384 accepted feature words per frame through AXI-stream wrapper.",}}
\par\smallskip
This Python line participates in the export, validation, or report-generation pipeline. It is executed on the host PC to prepare data for RTL or to check RTL outputs.
\paragraph{Line 611}
\noindent\colorbox{gray!8}{\parbox{\dimexpr\linewidth-2\fboxsep\relax}{\raggedright\ttfamily\scriptsize             "Per class: CLASS\_INIT plus one MAC\_READ per feature; active features run MAC\_MUL and MAC\_ACCUM, while gated features jump directly to CLASS\_DONE.",}}
\par\smallskip
This Python line participates in the export, validation, or report-generation pipeline. It is executed on the host PC to prepare data for RTL or to check RTL outputs.
\paragraph{Line 612}
\noindent\colorbox{gray!8}{\parbox{\dimexpr\linewidth-2\fboxsep\relax}{\raggedright\ttfamily\scriptsize             f"Activity gate: abs(feature) <= \{infer\_skip\_abs\_thresh\} bypasses DSP multiply/accumulate for inference; measured exported-set skip fraction is \{feature\_mac\_skip\_fraction * 100:.2f\}\%.",}}
\par\smallskip
This Python line participates in the export, validation, or report-generation pipeline. It is executed on the host PC to prepare data for RTL or to check RTL outputs.
\paragraph{Line 613}
\noindent\colorbox{gray!8}{\parbox{\dimexpr\linewidth-2\fboxsep\relax}{\raggedright\ttfamily\scriptsize             f"RTL inference skip counter \{rtl\_infer\_skip\_count:,\} removes \{saved\_infer\_cycles:,\} MAC\_MUL/MAC\_ACCUM cycles over the 10k fixture.",}}
\par\smallskip
This Python line participates in the export, validation, or report-generation pipeline. It is executed on the host PC to prepare data for RTL or to check RTL outputs.
\paragraph{Line 614}
\noindent\colorbox{gray!8}{\parbox{\dimexpr\linewidth-2\fboxsep\relax}{\raggedright\ttfamily\scriptsize             "Learning update, when applied: active features update true and predicted class rows; gated features skip the no-op read/calc/write sequence.",}}
\par\smallskip
This line handles predicted labels. It checks that argmax behavior in Python and RTL chooses the same class.
\paragraph{Line 615}
\noindent\colorbox{gray!8}{\parbox{\dimexpr\linewidth-2\fboxsep\relax}{\raggedright\ttfamily\scriptsize             f"RTL learning skip counter \{learn\_train\_skip\_count:,\} removes approximately \{saved\_train\_cycles:,\} true/pred update cycles in the full learning run.",}}
\par\smallskip
This line handles predicted labels. It checks that argmax behavior in Python and RTL chooses the same class.
\paragraph{Line 616}
\noindent\colorbox{gray!8}{\parbox{\dimexpr\linewidth-2\fboxsep\relax}{\raggedright\ttfamily\scriptsize             f"Learning is margin-gated, so the full 10k validation stream applied \{learn\_val.get('updates')\} updates.",}}
\par\smallskip
This line handles the online-learning rule or its validation fields. It is central to proving that the FPGA update path is not bypassed.
\paragraph{Line 617}
\noindent\colorbox{gray!8}{\parbox{\dimexpr\linewidth-2\fboxsep\relax}{\raggedright\ttfamily\scriptsize         ]}}
\par\smallskip
This Python line participates in the export, validation, or report-generation pipeline. It is executed on the host PC to prepare data for RTL or to check RTL outputs.
\paragraph{Line 618}
\noindent\colorbox{gray!8}{\parbox{\dimexpr\linewidth-2\fboxsep\relax}{\raggedright\ttfamily\scriptsize         save\_pdf\_page(pdf, plot\_text\_page("Pipeline / Latency Accounting", latency\_lines), png\_dir, "30\_latency")}}
\par\smallskip
This Python line participates in the export, validation, or report-generation pipeline. It is executed on the host PC to prepare data for RTL or to check RTL outputs.
\paragraph{Line 619}
\noindent\colorbox{gray!8}{\parbox{\dimexpr\linewidth-2\fboxsep\relax}{\raggedright\ttfamily\scriptsize \emph{blank line}}}
\par\smallskip
This blank line is not hardware or software behavior. It separates nearby declarations or operations so the file remains readable and reviewable.
\paragraph{Line 620}
\noindent\colorbox{gray!8}{\parbox{\dimexpr\linewidth-2\fboxsep\relax}{\raggedright\ttfamily\scriptsize     print(f"Wrote \{pdf\_path\}")}}
\par\smallskip
This Python line participates in the export, validation, or report-generation pipeline. It is executed on the host PC to prepare data for RTL or to check RTL outputs.
\paragraph{Line 621}
\noindent\colorbox{gray!8}{\parbox{\dimexpr\linewidth-2\fboxsep\relax}{\raggedright\ttfamily\scriptsize     print(f"Wrote plot PNGs in \{png\_dir\}")}}
\par\smallskip
This Python line participates in the export, validation, or report-generation pipeline. It is executed on the host PC to prepare data for RTL or to check RTL outputs.
\paragraph{Line 622}
\noindent\colorbox{gray!8}{\parbox{\dimexpr\linewidth-2\fboxsep\relax}{\raggedright\ttfamily\scriptsize \emph{blank line}}}
\par\smallskip
This blank line is not hardware or software behavior. It separates nearby declarations or operations so the file remains readable and reviewable.
\paragraph{Line 623}
\noindent\colorbox{gray!8}{\parbox{\dimexpr\linewidth-2\fboxsep\relax}{\raggedright\ttfamily\scriptsize \emph{blank line}}}
\par\smallskip
This blank line is not hardware or software behavior. It separates nearby declarations or operations so the file remains readable and reviewable.
\paragraph{Line 624}
\noindent\colorbox{gray!8}{\parbox{\dimexpr\linewidth-2\fboxsep\relax}{\raggedright\ttfamily\scriptsize def latex\_escape(text: str) -> str:}}
\par\smallskip
This defines a Python function. The function packages one repeatable step such as loading a memory file, simulating the golden learning rule, validating CSVs, or generating plots.
\paragraph{Line 625}
\noindent\colorbox{gray!8}{\parbox{\dimexpr\linewidth-2\fboxsep\relax}{\raggedright\ttfamily\scriptsize     replacements = \{}}
\par\smallskip
This Python line participates in the export, validation, or report-generation pipeline. It is executed on the host PC to prepare data for RTL or to check RTL outputs.
\paragraph{Line 626}
\noindent\colorbox{gray!8}{\parbox{\dimexpr\linewidth-2\fboxsep\relax}{\raggedright\ttfamily\scriptsize         "\textbackslash{}\textbackslash{}": r"\textbackslash{}textbackslash\{\}",}}
\par\smallskip
This Python line participates in the export, validation, or report-generation pipeline. It is executed on the host PC to prepare data for RTL or to check RTL outputs.
\paragraph{Line 627}
\noindent\colorbox{gray!8}{\parbox{\dimexpr\linewidth-2\fboxsep\relax}{\raggedright\ttfamily\scriptsize         "\&": r"\textbackslash{}\&",}}
\par\smallskip
This Python line participates in the export, validation, or report-generation pipeline. It is executed on the host PC to prepare data for RTL or to check RTL outputs.
\paragraph{Line 628}
\noindent\colorbox{gray!8}{\parbox{\dimexpr\linewidth-2\fboxsep\relax}{\raggedright\ttfamily\scriptsize         "\%": r"\textbackslash{}\%",}}
\par\smallskip
This Python line participates in the export, validation, or report-generation pipeline. It is executed on the host PC to prepare data for RTL or to check RTL outputs.
\paragraph{Line 629}
\noindent\colorbox{gray!8}{\parbox{\dimexpr\linewidth-2\fboxsep\relax}{\raggedright\ttfamily\scriptsize         "\$": r"\textbackslash{}\$",}}
\par\smallskip
This Python line participates in the export, validation, or report-generation pipeline. It is executed on the host PC to prepare data for RTL or to check RTL outputs.
\paragraph{Line 630}
\noindent\colorbox{gray!8}{\parbox{\dimexpr\linewidth-2\fboxsep\relax}{\raggedright\ttfamily\scriptsize         "\#": r"\textbackslash{}\#",}}
\par\smallskip
This Python line participates in the export, validation, or report-generation pipeline. It is executed on the host PC to prepare data for RTL or to check RTL outputs.
\paragraph{Line 631}
\noindent\colorbox{gray!8}{\parbox{\dimexpr\linewidth-2\fboxsep\relax}{\raggedright\ttfamily\scriptsize         "\_": r"\textbackslash{}\_",}}
\par\smallskip
This Python line participates in the export, validation, or report-generation pipeline. It is executed on the host PC to prepare data for RTL or to check RTL outputs.
\paragraph{Line 632}
\noindent\colorbox{gray!8}{\parbox{\dimexpr\linewidth-2\fboxsep\relax}{\raggedright\ttfamily\scriptsize         "\{": r"\textbackslash{}\{",}}
\par\smallskip
This Python line participates in the export, validation, or report-generation pipeline. It is executed on the host PC to prepare data for RTL or to check RTL outputs.
\paragraph{Line 633}
\noindent\colorbox{gray!8}{\parbox{\dimexpr\linewidth-2\fboxsep\relax}{\raggedright\ttfamily\scriptsize         "\}": r"\textbackslash{}\}",}}
\par\smallskip
This Python line participates in the export, validation, or report-generation pipeline. It is executed on the host PC to prepare data for RTL or to check RTL outputs.
\paragraph{Line 634}
\noindent\colorbox{gray!8}{\parbox{\dimexpr\linewidth-2\fboxsep\relax}{\raggedright\ttfamily\scriptsize         "\textasciitilde{}": r"\textbackslash{}textasciitilde\{\}",}}
\par\smallskip
This Python line participates in the export, validation, or report-generation pipeline. It is executed on the host PC to prepare data for RTL or to check RTL outputs.
\paragraph{Line 635}
\noindent\colorbox{gray!8}{\parbox{\dimexpr\linewidth-2\fboxsep\relax}{\raggedright\ttfamily\scriptsize         "\textasciicircum{}": r"\textbackslash{}textasciicircum\{\}",}}
\par\smallskip
This Python line participates in the export, validation, or report-generation pipeline. It is executed on the host PC to prepare data for RTL or to check RTL outputs.
\paragraph{Line 636}
\noindent\colorbox{gray!8}{\parbox{\dimexpr\linewidth-2\fboxsep\relax}{\raggedright\ttfamily\scriptsize     \}}}
\par\smallskip
This Python line participates in the export, validation, or report-generation pipeline. It is executed on the host PC to prepare data for RTL or to check RTL outputs.
\paragraph{Line 637}
\noindent\colorbox{gray!8}{\parbox{\dimexpr\linewidth-2\fboxsep\relax}{\raggedright\ttfamily\scriptsize     return "".join(replacements.get(ch, ch) for ch in text)}}
\par\smallskip
This Python line participates in the export, validation, or report-generation pipeline. It is executed on the host PC to prepare data for RTL or to check RTL outputs.
\paragraph{Line 638}
\noindent\colorbox{gray!8}{\parbox{\dimexpr\linewidth-2\fboxsep\relax}{\raggedright\ttfamily\scriptsize \emph{blank line}}}
\par\smallskip
This blank line is not hardware or software behavior. It separates nearby declarations or operations so the file remains readable and reviewable.
\paragraph{Line 639}
\noindent\colorbox{gray!8}{\parbox{\dimexpr\linewidth-2\fboxsep\relax}{\raggedright\ttfamily\scriptsize \emph{blank line}}}
\par\smallskip
This blank line is not hardware or software behavior. It separates nearby declarations or operations so the file remains readable and reviewable.
\paragraph{Line 640}
\noindent\colorbox{gray!8}{\parbox{\dimexpr\linewidth-2\fboxsep\relax}{\raggedright\ttfamily\scriptsize def explain\_line(path: Path, line: str, index: int, previous\_nonblank: str) -> str:}}
\par\smallskip
This defines a Python function. The function packages one repeatable step such as loading a memory file, simulating the golden learning rule, validating CSVs, or generating plots.
\paragraph{Line 641}
\noindent\colorbox{gray!8}{\parbox{\dimexpr\linewidth-2\fboxsep\relax}{\raggedright\ttfamily\scriptsize     stripped = line.strip()}}
\par\smallskip
This Python line participates in the export, validation, or report-generation pipeline. It is executed on the host PC to prepare data for RTL or to check RTL outputs.
\paragraph{Line 642}
\noindent\colorbox{gray!8}{\parbox{\dimexpr\linewidth-2\fboxsep\relax}{\raggedright\ttfamily\scriptsize     suffix = path.suffix.lower()}}
\par\smallskip
This Python line participates in the export, validation, or report-generation pipeline. It is executed on the host PC to prepare data for RTL or to check RTL outputs.
\paragraph{Line 643}
\noindent\colorbox{gray!8}{\parbox{\dimexpr\linewidth-2\fboxsep\relax}{\raggedright\ttfamily\scriptsize     name = path.name}}
\par\smallskip
This Python line participates in the export, validation, or report-generation pipeline. It is executed on the host PC to prepare data for RTL or to check RTL outputs.
\paragraph{Line 644}
\noindent\colorbox{gray!8}{\parbox{\dimexpr\linewidth-2\fboxsep\relax}{\raggedright\ttfamily\scriptsize     if not stripped:}}
\par\smallskip
This Python line participates in the export, validation, or report-generation pipeline. It is executed on the host PC to prepare data for RTL or to check RTL outputs.
\paragraph{Line 645}
\noindent\colorbox{gray!8}{\parbox{\dimexpr\linewidth-2\fboxsep\relax}{\raggedright\ttfamily\scriptsize         return "This blank line is not hardware or software behavior. It separates nearby declarations or operations so the file remains readable and reviewable."}}
\par\smallskip
This Python line participates in the export, validation, or report-generation pipeline. It is executed on the host PC to prepare data for RTL or to check RTL outputs.
\paragraph{Line 646}
\noindent\colorbox{gray!8}{\parbox{\dimexpr\linewidth-2\fboxsep\relax}{\raggedright\ttfamily\scriptsize     if stripped.startswith("//") or stripped.startswith("\#"):}}
\par\smallskip
This Python line participates in the export, validation, or report-generation pipeline. It is executed on the host PC to prepare data for RTL or to check RTL outputs.
\paragraph{Line 647}
\noindent\colorbox{gray!8}{\parbox{\dimexpr\linewidth-2\fboxsep\relax}{\raggedright\ttfamily\scriptsize         return "This is a human-facing comment. It documents intent for the reader and is ignored by the compiler or interpreter."}}
\par\smallskip
This Python line participates in the export, validation, or report-generation pipeline. It is executed on the host PC to prepare data for RTL or to check RTL outputs.
\paragraph{Line 648}
\noindent\colorbox{gray!8}{\parbox{\dimexpr\linewidth-2\fboxsep\relax}{\raggedright\ttfamily\scriptsize     if suffix in \{".sv", ".v", ".svh"\}:}}
\par\smallskip
This Python line participates in the export, validation, or report-generation pipeline. It is executed on the host PC to prepare data for RTL or to check RTL outputs.
\paragraph{Line 649}
\noindent\colorbox{gray!8}{\parbox{\dimexpr\linewidth-2\fboxsep\relax}{\raggedright\ttfamily\scriptsize         if stripped.startswith("`timescale"):}}
\par\smallskip
This Python line participates in the export, validation, or report-generation pipeline. It is executed on the host PC to prepare data for RTL or to check RTL outputs.
\paragraph{Line 650}
\noindent\colorbox{gray!8}{\parbox{\dimexpr\linewidth-2\fboxsep\relax}{\raggedright\ttfamily\scriptsize             return "This simulator directive sets the time unit and precision used by testbenches and waveform timestamps. It keeps delays such as clock half-periods numerically meaningful in XSim."}}
\par\smallskip
This Python line participates in the export, validation, or report-generation pipeline. It is executed on the host PC to prepare data for RTL or to check RTL outputs.
\paragraph{Line 651}
\noindent\colorbox{gray!8}{\parbox{\dimexpr\linewidth-2\fboxsep\relax}{\raggedright\ttfamily\scriptsize         if stripped.startswith("`default\_nettype none"):}}
\par\smallskip
This Python line participates in the export, validation, or report-generation pipeline. It is executed on the host PC to prepare data for RTL or to check RTL outputs.
\paragraph{Line 652}
\noindent\colorbox{gray!8}{\parbox{\dimexpr\linewidth-2\fboxsep\relax}{\raggedright\ttfamily\scriptsize             return "This disables implicit one-bit wires. It prevents spelling mistakes from silently creating unwanted hardware nets."}}
\par\smallskip
This Python line participates in the export, validation, or report-generation pipeline. It is executed on the host PC to prepare data for RTL or to check RTL outputs.
\paragraph{Line 653}
\noindent\colorbox{gray!8}{\parbox{\dimexpr\linewidth-2\fboxsep\relax}{\raggedright\ttfamily\scriptsize         if stripped.startswith("`include"):}}
\par\smallskip
This Python line participates in the export, validation, or report-generation pipeline. It is executed on the host PC to prepare data for RTL or to check RTL outputs.
\paragraph{Line 654}
\noindent\colorbox{gray!8}{\parbox{\dimexpr\linewidth-2\fboxsep\relax}{\raggedright\ttfamily\scriptsize             return "This includes shared compile-time constants such as feature dimensions and fixed-point widths. The RTL and testbenches therefore agree on the same trained-model shape."}}
\par\smallskip
This Python line participates in the export, validation, or report-generation pipeline. It is executed on the host PC to prepare data for RTL or to check RTL outputs.
\paragraph{Line 655}
\noindent\colorbox{gray!8}{\parbox{\dimexpr\linewidth-2\fboxsep\relax}{\raggedright\ttfamily\scriptsize         if stripped.startswith("module "):}}
\par\smallskip
This Python line participates in the export, validation, or report-generation pipeline. It is executed on the host PC to prepare data for RTL or to check RTL outputs.
\paragraph{Line 656}
\noindent\colorbox{gray!8}{\parbox{\dimexpr\linewidth-2\fboxsep\relax}{\raggedright\ttfamily\scriptsize             return "This begins a hardware module, which is the FPGA equivalent of a reusable circuit block. Everything until the matching endmodule describes ports, storage, combinational equations, and clocked behavior for this block."}}
\par\smallskip
This Python line participates in the export, validation, or report-generation pipeline. It is executed on the host PC to prepare data for RTL or to check RTL outputs.
\paragraph{Line 657}
\noindent\colorbox{gray!8}{\parbox{\dimexpr\linewidth-2\fboxsep\relax}{\raggedright\ttfamily\scriptsize         if stripped.startswith("endmodule"):}}
\par\smallskip
This Python line participates in the export, validation, or report-generation pipeline. It is executed on the host PC to prepare data for RTL or to check RTL outputs.
\paragraph{Line 658}
\noindent\colorbox{gray!8}{\parbox{\dimexpr\linewidth-2\fboxsep\relax}{\raggedright\ttfamily\scriptsize             return "This closes the module definition. Synthesis stops adding logic to this circuit block at this point."}}
\par\smallskip
This Python line participates in the export, validation, or report-generation pipeline. It is executed on the host PC to prepare data for RTL or to check RTL outputs.
\paragraph{Line 659}
\noindent\colorbox{gray!8}{\parbox{\dimexpr\linewidth-2\fboxsep\relax}{\raggedright\ttfamily\scriptsize         if "INFER\_SKIP\_ABS\_THRESH" in stripped:}}
\par\smallskip
This Python line participates in the export, validation, or report-generation pipeline. It is executed on the host PC to prepare data for RTL or to check RTL outputs.
\paragraph{Line 660}
\noindent\colorbox{gray!8}{\parbox{\dimexpr\linewidth-2\fboxsep\relax}{\raggedright\ttfamily\scriptsize             return "This line controls the feature-activity gate. When the absolute fixed-point feature is at or below this threshold, the RTL treats it as zero so the exported Spikingformer head ignores small noisy activations and records a measurable MAC-skip opportunity."}}
\par\smallskip
This Python line participates in the export, validation, or report-generation pipeline. It is executed on the host PC to prepare data for RTL or to check RTL outputs.
\paragraph{Line 661}
\noindent\colorbox{gray!8}{\parbox{\dimexpr\linewidth-2\fboxsep\relax}{\raggedright\ttfamily\scriptsize         if "skip\_feature" in stripped:}}
\par\smallskip
This Python line participates in the export, validation, or report-generation pipeline. It is executed on the host PC to prepare data for RTL or to check RTL outputs.
\paragraph{Line 662}
\noindent\colorbox{gray!8}{\parbox{\dimexpr\linewidth-2\fboxsep\relax}{\raggedright\ttfamily\scriptsize             return "This line is part of the combinational decision that detects low-magnitude features. It is the hardware equivalent of asking whether a small activation should be silent for this inference/update pass."}}
\par\smallskip
This line handles the online-learning rule or its validation fields. It is central to proving that the FPGA update path is not bypassed.
\paragraph{Line 663}
\noindent\colorbox{gray!8}{\parbox{\dimexpr\linewidth-2\fboxsep\relax}{\raggedright\ttfamily\scriptsize         if "effective\_feature" in stripped:}}
\par\smallskip
This Python line participates in the export, validation, or report-generation pipeline. It is executed on the host PC to prepare data for RTL or to check RTL outputs.
\paragraph{Line 664}
\noindent\colorbox{gray!8}{\parbox{\dimexpr\linewidth-2\fboxsep\relax}{\raggedright\ttfamily\scriptsize             return "This line is part of the denoised feature path. The returned value is either the original signed feature or exactly zero, and that registered value is what feeds both the DSP multiply and the online-learning delta."}}
\par\smallskip
This Python line participates in the export, validation, or report-generation pipeline. It is executed on the host PC to prepare data for RTL or to check RTL outputs.
\paragraph{Line 665}
\noindent\colorbox{gray!8}{\parbox{\dimexpr\linewidth-2\fboxsep\relax}{\raggedright\ttfamily\scriptsize         if "skip\_count" in stripped:}}
\par\smallskip
This Python line participates in the export, validation, or report-generation pipeline. It is executed on the host PC to prepare data for RTL or to check RTL outputs.
\paragraph{Line 666}
\noindent\colorbox{gray!8}{\parbox{\dimexpr\linewidth-2\fboxsep\relax}{\raggedright\ttfamily\scriptsize             return "This line exposes or updates a hardware activity counter. The counter makes the feature gate visible in waveforms and CSV traces, so the power-saving opportunity is observable from RTL rather than only from Python statistics."}}
\par\smallskip
This line handles the online-learning rule or its validation fields. It is central to proving that the FPGA update path is not bypassed.
\paragraph{Line 667}
\noindent\colorbox{gray!8}{\parbox{\dimexpr\linewidth-2\fboxsep\relax}{\raggedright\ttfamily\scriptsize         if "parameter" in stripped:}}
\par\smallskip
This Python line participates in the export, validation, or report-generation pipeline. It is executed on the host PC to prepare data for RTL or to check RTL outputs.
\paragraph{Line 668}
\noindent\colorbox{gray!8}{\parbox{\dimexpr\linewidth-2\fboxsep\relax}{\raggedright\ttfamily\scriptsize             return "This is a compile-time hardware setting. Vivado resolves it before synthesis, so it controls the generated circuit size, fixed-point width, or learning rule without runtime overhead."}}
\par\smallskip
This Python line participates in the export, validation, or report-generation pipeline. It is executed on the host PC to prepare data for RTL or to check RTL outputs.
\paragraph{Line 669}
\noindent\colorbox{gray!8}{\parbox{\dimexpr\linewidth-2\fboxsep\relax}{\raggedright\ttfamily\scriptsize         if re.search(r"\textbackslash{}binput\textbackslash{}b|\textbackslash{}boutput\textbackslash{}b", stripped):}}
\par\smallskip
This Python line participates in the export, validation, or report-generation pipeline. It is executed on the host PC to prepare data for RTL or to check RTL outputs.
\paragraph{Line 670}
\noindent\colorbox{gray!8}{\parbox{\dimexpr\linewidth-2\fboxsep\relax}{\raggedright\ttfamily\scriptsize             return "This line declares a module port. A port is a physical or wrapper-visible signal used to move clocks, control pulses, features, labels, predictions, or debug information across module boundaries."}}
\par\smallskip
This line handles predicted labels. It checks that argmax behavior in Python and RTL chooses the same class.
\paragraph{Line 671}
\noindent\colorbox{gray!8}{\parbox{\dimexpr\linewidth-2\fboxsep\relax}{\raggedright\ttfamily\scriptsize         if "typedef enum" in stripped:}}
\par\smallskip
This Python line participates in the export, validation, or report-generation pipeline. It is executed on the host PC to prepare data for RTL or to check RTL outputs.
\paragraph{Line 672}
\noindent\colorbox{gray!8}{\parbox{\dimexpr\linewidth-2\fboxsep\relax}{\raggedright\ttfamily\scriptsize             return "This declares the finite-state-machine state type. Using named states makes the cycle-by-cycle control flow explicit instead of relying on anonymous numeric constants."}}
\par\smallskip
This Python line participates in the export, validation, or report-generation pipeline. It is executed on the host PC to prepare data for RTL or to check RTL outputs.
\paragraph{Line 673}
\noindent\colorbox{gray!8}{\parbox{\dimexpr\linewidth-2\fboxsep\relax}{\raggedright\ttfamily\scriptsize         if stripped.startswith("ST\_") or stripped.startswith("WR\_"):}}
\par\smallskip
This Python line participates in the export, validation, or report-generation pipeline. It is executed on the host PC to prepare data for RTL or to check RTL outputs.
\paragraph{Line 674}
\noindent\colorbox{gray!8}{\parbox{\dimexpr\linewidth-2\fboxsep\relax}{\raggedright\ttfamily\scriptsize             return "This names one state in the control FSM. The numeric encoding is what appears in waveform traces, while the name explains the pipeline phase being executed."}}
\par\smallskip
This Python line participates in the export, validation, or report-generation pipeline. It is executed on the host PC to prepare data for RTL or to check RTL outputs.
\paragraph{Line 675}
\noindent\colorbox{gray!8}{\parbox{\dimexpr\linewidth-2\fboxsep\relax}{\raggedright\ttfamily\scriptsize         if "feature\_mem" in stripped:}}
\par\smallskip
This Python line participates in the export, validation, or report-generation pipeline. It is executed on the host PC to prepare data for RTL or to check RTL outputs.
\paragraph{Line 676}
\noindent\colorbox{gray!8}{\parbox{\dimexpr\linewidth-2\fboxsep\relax}{\raggedright\ttfamily\scriptsize             return "This line touches the feature memory that holds one 384-dimensional fixed-point Spikingformer feature vector for the current image. Keeping the feature vector local avoids repeatedly streaming the same values for every class score."}}
\par\smallskip
This Python line participates in the export, validation, or report-generation pipeline. It is executed on the host PC to prepare data for RTL or to check RTL outputs.
\paragraph{Line 677}
\noindent\colorbox{gray!8}{\parbox{\dimexpr\linewidth-2\fboxsep\relax}{\raggedright\ttfamily\scriptsize         if "weight\_mem" in stripped:}}
\par\smallskip
This Python line participates in the export, validation, or report-generation pipeline. It is executed on the host PC to prepare data for RTL or to check RTL outputs.
\paragraph{Line 678}
\noindent\colorbox{gray!8}{\parbox{\dimexpr\linewidth-2\fboxsep\relax}{\raggedright\ttfamily\scriptsize             return "This line touches the trained classifier weight store. In the final RTL the large 10 x 384 int16 matrix is inferred as true dual-port BRAM, so weights remain on chip and can also be updated by the learning state machine."}}
\par\smallskip
This line handles the classifier parameters. Those arrays mirror the RTL memories, so mismatches reveal an export, quantization, or hardware-update bug.
\paragraph{Line 679}
\noindent\colorbox{gray!8}{\parbox{\dimexpr\linewidth-2\fboxsep\relax}{\raggedright\ttfamily\scriptsize         if "bias\_mem" in stripped:}}
\par\smallskip
This line handles the classifier parameters. Those arrays mirror the RTL memories, so mismatches reveal an export, quantization, or hardware-update bug.
\paragraph{Line 680}
\noindent\colorbox{gray!8}{\parbox{\dimexpr\linewidth-2\fboxsep\relax}{\raggedright\ttfamily\scriptsize             return "This line touches the class-bias store. Bias values are part of the affine classifier score and are updated on chip in separate one-write states to keep synthesis clean."}}
\par\smallskip
This line handles the classifier parameters. Those arrays mirror the RTL memories, so mismatches reveal an export, quantization, or hardware-update bug.
\paragraph{Line 681}
\noindent\colorbox{gray!8}{\parbox{\dimexpr\linewidth-2\fboxsep\relax}{\raggedright\ttfamily\scriptsize         if "class\_score\_mem" in stripped:}}
\par\smallskip
This Python line participates in the export, validation, or report-generation pipeline. It is executed on the host PC to prepare data for RTL or to check RTL outputs.
\paragraph{Line 682}
\noindent\colorbox{gray!8}{\parbox{\dimexpr\linewidth-2\fboxsep\relax}{\raggedright\ttfamily\scriptsize             return "This stores or reads the ten class scores for the current sample. Those scores are needed for argmax prediction, CSV testbench reporting, and the margin test used before an online update."}}
\par\smallskip
This line handles class scores. The score vector is the direct software equivalent of the RTL MAC loop over weights and features.
\paragraph{Line 683}
\noindent\colorbox{gray!8}{\parbox{\dimexpr\linewidth-2\fboxsep\relax}{\raggedright\ttfamily\scriptsize         if stripped.startswith("assign "):}}
\par\smallskip
This Python line participates in the export, validation, or report-generation pipeline. It is executed on the host PC to prepare data for RTL or to check RTL outputs.
\paragraph{Line 684}
\noindent\colorbox{gray!8}{\parbox{\dimexpr\linewidth-2\fboxsep\relax}{\raggedright\ttfamily\scriptsize             return "This is continuous combinational wiring. The output follows its source expression immediately in hardware rather than waiting for a clock edge."}}
\par\smallskip
This Python line participates in the export, validation, or report-generation pipeline. It is executed on the host PC to prepare data for RTL or to check RTL outputs.
\paragraph{Line 685}
\noindent\colorbox{gray!8}{\parbox{\dimexpr\linewidth-2\fboxsep\relax}{\raggedright\ttfamily\scriptsize         if stripped.startswith("initial"):}}
\par\smallskip
This Python line participates in the export, validation, or report-generation pipeline. It is executed on the host PC to prepare data for RTL or to check RTL outputs.
\paragraph{Line 686}
\noindent\colorbox{gray!8}{\parbox{\dimexpr\linewidth-2\fboxsep\relax}{\raggedright\ttfamily\scriptsize             return "This simulation/synthesis initialization block loads trained memories from generated .mem files. On FPGA, Vivado maps these initial values into BRAM/LUTRAM initialization data."}}
\par\smallskip
This Python line participates in the export, validation, or report-generation pipeline. It is executed on the host PC to prepare data for RTL or to check RTL outputs.
\paragraph{Line 687}
\noindent\colorbox{gray!8}{\parbox{\dimexpr\linewidth-2\fboxsep\relax}{\raggedright\ttfamily\scriptsize         if "\$readmemh" in stripped:}}
\par\smallskip
This Python line participates in the export, validation, or report-generation pipeline. It is executed on the host PC to prepare data for RTL or to check RTL outputs.
\paragraph{Line 688}
\noindent\colorbox{gray!8}{\parbox{\dimexpr\linewidth-2\fboxsep\relax}{\raggedright\ttfamily\scriptsize             return "This reads a hexadecimal memory file into an RTL array. It is how exported trained weights and biases become deterministic FPGA memory contents."}}
\par\smallskip
This line handles the classifier parameters. Those arrays mirror the RTL memories, so mismatches reveal an export, quantization, or hardware-update bug.
\paragraph{Line 689}
\noindent\colorbox{gray!8}{\parbox{\dimexpr\linewidth-2\fboxsep\relax}{\raggedright\ttfamily\scriptsize         if stripped.startswith("function"):}}
\par\smallskip
This Python line participates in the export, validation, or report-generation pipeline. It is executed on the host PC to prepare data for RTL or to check RTL outputs.
\paragraph{Line 690}
\noindent\colorbox{gray!8}{\parbox{\dimexpr\linewidth-2\fboxsep\relax}{\raggedright\ttfamily\scriptsize             return "This begins a pure helper function used inside the RTL. It describes fixed combinational math, not a software function running on a CPU."}}
\par\smallskip
This Python line participates in the export, validation, or report-generation pipeline. It is executed on the host PC to prepare data for RTL or to check RTL outputs.
\paragraph{Line 691}
\noindent\colorbox{gray!8}{\parbox{\dimexpr\linewidth-2\fboxsep\relax}{\raggedright\ttfamily\scriptsize         if "sat\_weight" in stripped:}}
\par\smallskip
This Python line participates in the export, validation, or report-generation pipeline. It is executed on the host PC to prepare data for RTL or to check RTL outputs.
\paragraph{Line 692}
\noindent\colorbox{gray!8}{\parbox{\dimexpr\linewidth-2\fboxsep\relax}{\raggedright\ttfamily\scriptsize             return "This line is part of signed saturation for 16-bit weights. Saturation prevents online learning from wrapping around when an update would exceed the representable fixed-point range."}}
\par\smallskip
This line handles the classifier parameters. Those arrays mirror the RTL memories, so mismatches reveal an export, quantization, or hardware-update bug.
\paragraph{Line 693}
\noindent\colorbox{gray!8}{\parbox{\dimexpr\linewidth-2\fboxsep\relax}{\raggedright\ttfamily\scriptsize         if "learn\_delta" in stripped:}}
\par\smallskip
This Python line participates in the export, validation, or report-generation pipeline. It is executed on the host PC to prepare data for RTL or to check RTL outputs.
\paragraph{Line 694}
\noindent\colorbox{gray!8}{\parbox{\dimexpr\linewidth-2\fboxsep\relax}{\raggedright\ttfamily\scriptsize             return "This line implements the hardware learning step size. The feature is sign-extended and arithmetically shifted so positive and negative features update weights symmetrically."}}
\par\smallskip
This line handles the classifier parameters. Those arrays mirror the RTL memories, so mismatches reveal an export, quantization, or hardware-update bug.
\paragraph{Line 695}
\noindent\colorbox{gray!8}{\parbox{\dimexpr\linewidth-2\fboxsep\relax}{\raggedright\ttfamily\scriptsize         if "always\_ff" in stripped:}}
\par\smallskip
This Python line participates in the export, validation, or report-generation pipeline. It is executed on the host PC to prepare data for RTL or to check RTL outputs.
\paragraph{Line 696}
\noindent\colorbox{gray!8}{\parbox{\dimexpr\linewidth-2\fboxsep\relax}{\raggedright\ttfamily\scriptsize             return "This begins clocked sequential logic. Registers and RAM writes inside this block update only on a clock edge, which is what makes the design synthesizable and waveform-predictable."}}
\par\smallskip
This line handles predicted labels. It checks that argmax behavior in Python and RTL chooses the same class.
\paragraph{Line 697}
\noindent\colorbox{gray!8}{\parbox{\dimexpr\linewidth-2\fboxsep\relax}{\raggedright\ttfamily\scriptsize         if "if (!rst\_n)" in stripped:}}
\par\smallskip
This Python line participates in the export, validation, or report-generation pipeline. It is executed on the host PC to prepare data for RTL or to check RTL outputs.
\paragraph{Line 698}
\noindent\colorbox{gray!8}{\parbox{\dimexpr\linewidth-2\fboxsep\relax}{\raggedright\ttfamily\scriptsize             return "This is the active-low reset branch. It puts control registers and visible outputs into a known starting state before a frame is processed."}}
\par\smallskip
This Python line participates in the export, validation, or report-generation pipeline. It is executed on the host PC to prepare data for RTL or to check RTL outputs.
\paragraph{Line 699}
\noindent\colorbox{gray!8}{\parbox{\dimexpr\linewidth-2\fboxsep\relax}{\raggedright\ttfamily\scriptsize         if "unique case" in stripped:}}
\par\smallskip
This Python line participates in the export, validation, or report-generation pipeline. It is executed on the host PC to prepare data for RTL or to check RTL outputs.
\paragraph{Line 700}
\noindent\colorbox{gray!8}{\parbox{\dimexpr\linewidth-2\fboxsep\relax}{\raggedright\ttfamily\scriptsize             return "This starts a state-machine branch table. The unique qualifier tells tools that one legal state should match at a time, improving diagnostics and optimization."}}
\par\smallskip
This Python line participates in the export, validation, or report-generation pipeline. It is executed on the host PC to prepare data for RTL or to check RTL outputs.
\paragraph{Line 701}
\noindent\colorbox{gray!8}{\parbox{\dimexpr\linewidth-2\fboxsep\relax}{\raggedright\ttfamily\scriptsize         if "state <=" in stripped:}}
\par\smallskip
This Python line participates in the export, validation, or report-generation pipeline. It is executed on the host PC to prepare data for RTL or to check RTL outputs.
\paragraph{Line 702}
\noindent\colorbox{gray!8}{\parbox{\dimexpr\linewidth-2\fboxsep\relax}{\raggedright\ttfamily\scriptsize             return "This schedules the next FSM state for the following clock edge. That transition is the hardware pipeline control path."}}
\par\smallskip
This Python line participates in the export, validation, or report-generation pipeline. It is executed on the host PC to prepare data for RTL or to check RTL outputs.
\paragraph{Line 703}
\noindent\colorbox{gray!8}{\parbox{\dimexpr\linewidth-2\fboxsep\relax}{\raggedright\ttfamily\scriptsize         if "<=" in stripped:}}
\par\smallskip
This Python line participates in the export, validation, or report-generation pipeline. It is executed on the host PC to prepare data for RTL or to check RTL outputs.
\paragraph{Line 704}
\noindent\colorbox{gray!8}{\parbox{\dimexpr\linewidth-2\fboxsep\relax}{\raggedright\ttfamily\scriptsize             return "This is a nonblocking register or memory assignment. The right side is sampled now and the left side updates together with other clocked assignments at the end of the clock edge."}}
\par\smallskip
This line handles the online-learning rule or its validation fields. It is central to proving that the FPGA update path is not bypassed.
\paragraph{Line 705}
\noindent\colorbox{gray!8}{\parbox{\dimexpr\linewidth-2\fboxsep\relax}{\raggedright\ttfamily\scriptsize         if "=" in stripped:}}
\par\smallskip
This Python line participates in the export, validation, or report-generation pipeline. It is executed on the host PC to prepare data for RTL or to check RTL outputs.
\paragraph{Line 706}
\noindent\colorbox{gray!8}{\parbox{\dimexpr\linewidth-2\fboxsep\relax}{\raggedright\ttfamily\scriptsize             return "This is a local compile-time or procedural assignment. In context it either defines a helper value or connects a temporary expression used by nearby hardware logic."}}
\par\smallskip
This Python line participates in the export, validation, or report-generation pipeline. It is executed on the host PC to prepare data for RTL or to check RTL outputs.
\paragraph{Line 707}
\noindent\colorbox{gray!8}{\parbox{\dimexpr\linewidth-2\fboxsep\relax}{\raggedright\ttfamily\scriptsize         return "This SystemVerilog line contributes to the structural or sequential description of the accelerator. Its effect is resolved by XSim for verification and by Vivado during synthesis."}}
\par\smallskip
This Python line participates in the export, validation, or report-generation pipeline. It is executed on the host PC to prepare data for RTL or to check RTL outputs.
\paragraph{Line 708}
\noindent\colorbox{gray!8}{\parbox{\dimexpr\linewidth-2\fboxsep\relax}{\raggedright\ttfamily\scriptsize     if suffix == ".py":}}
\par\smallskip
This Python line participates in the export, validation, or report-generation pipeline. It is executed on the host PC to prepare data for RTL or to check RTL outputs.
\paragraph{Line 709}
\noindent\colorbox{gray!8}{\parbox{\dimexpr\linewidth-2\fboxsep\relax}{\raggedright\ttfamily\scriptsize         if stripped.startswith("from ") or stripped.startswith("import "):}}
\par\smallskip
This Python line participates in the export, validation, or report-generation pipeline. It is executed on the host PC to prepare data for RTL or to check RTL outputs.
\paragraph{Line 710}
\noindent\colorbox{gray!8}{\parbox{\dimexpr\linewidth-2\fboxsep\relax}{\raggedright\ttfamily\scriptsize             return "This imports a Python library or symbol. The scripts use standard parsing, numeric arrays, plotting, and file APIs to make the hardware artifacts reproducible."}}
\par\smallskip
This Python line participates in the export, validation, or report-generation pipeline. It is executed on the host PC to prepare data for RTL or to check RTL outputs.
\paragraph{Line 711}
\noindent\colorbox{gray!8}{\parbox{\dimexpr\linewidth-2\fboxsep\relax}{\raggedright\ttfamily\scriptsize         if stripped.startswith("def "):}}
\par\smallskip
This Python line participates in the export, validation, or report-generation pipeline. It is executed on the host PC to prepare data for RTL or to check RTL outputs.
\paragraph{Line 712}
\noindent\colorbox{gray!8}{\parbox{\dimexpr\linewidth-2\fboxsep\relax}{\raggedright\ttfamily\scriptsize             return "This defines a Python function. The function packages one repeatable step such as loading a memory file, simulating the golden learning rule, validating CSVs, or generating plots."}}
\par\smallskip
This Python line participates in the export, validation, or report-generation pipeline. It is executed on the host PC to prepare data for RTL or to check RTL outputs.
\paragraph{Line 713}
\noindent\colorbox{gray!8}{\parbox{\dimexpr\linewidth-2\fboxsep\relax}{\raggedright\ttfamily\scriptsize         if "argparse" in stripped or "add\_argument" in stripped:}}
\par\smallskip
This contributes to the command-line interface. It lets the rerun flow choose paths, sample counts, or output locations without editing code.
\paragraph{Line 714}
\noindent\colorbox{gray!8}{\parbox{\dimexpr\linewidth-2\fboxsep\relax}{\raggedright\ttfamily\scriptsize             return "This contributes to the command-line interface. It lets the rerun flow choose paths, sample counts, or output locations without editing code."}}
\par\smallskip
This Python line participates in the export, validation, or report-generation pipeline. It is executed on the host PC to prepare data for RTL or to check RTL outputs.
\paragraph{Line 715}
\noindent\colorbox{gray!8}{\parbox{\dimexpr\linewidth-2\fboxsep\relax}{\raggedright\ttfamily\scriptsize         if "Path" in stripped:}}
\par\smallskip
This uses path objects so generated artifacts are addressed consistently across Windows shells and Python code.
\paragraph{Line 716}
\noindent\colorbox{gray!8}{\parbox{\dimexpr\linewidth-2\fboxsep\relax}{\raggedright\ttfamily\scriptsize             return "This uses path objects so generated artifacts are addressed consistently across Windows shells and Python code."}}
\par\smallskip
This Python line participates in the export, validation, or report-generation pipeline. It is executed on the host PC to prepare data for RTL or to check RTL outputs.
\paragraph{Line 717}
\noindent\colorbox{gray!8}{\parbox{\dimexpr\linewidth-2\fboxsep\relax}{\raggedright\ttfamily\scriptsize         if "read\_text" in stripped or "open(" in stripped:}}
\par\smallskip
This reads an artifact produced by training, export, simulation, or Vivado. The validator/report depends on files generated by the actual run rather than hidden constants.
\paragraph{Line 718}
\noindent\colorbox{gray!8}{\parbox{\dimexpr\linewidth-2\fboxsep\relax}{\raggedright\ttfamily\scriptsize             return "This reads an artifact produced by training, export, simulation, or Vivado. The validator/report depends on files generated by the actual run rather than hidden constants."}}
\par\smallskip
This Python line participates in the export, validation, or report-generation pipeline. It is executed on the host PC to prepare data for RTL or to check RTL outputs.
\paragraph{Line 719}
\noindent\colorbox{gray!8}{\parbox{\dimexpr\linewidth-2\fboxsep\relax}{\raggedright\ttfamily\scriptsize         if "write\_text" in stripped or "json.dump" in stripped:}}
\par\smallskip
This writes machine-readable output for later reports. Keeping JSON/CSV artifacts makes the verification auditable.
\paragraph{Line 720}
\noindent\colorbox{gray!8}{\parbox{\dimexpr\linewidth-2\fboxsep\relax}{\raggedright\ttfamily\scriptsize             return "This writes machine-readable output for later reports. Keeping JSON/CSV artifacts makes the verification auditable."}}
\par\smallskip
This Python line participates in the export, validation, or report-generation pipeline. It is executed on the host PC to prepare data for RTL or to check RTL outputs.
\paragraph{Line 721}
\noindent\colorbox{gray!8}{\parbox{\dimexpr\linewidth-2\fboxsep\relax}{\raggedright\ttfamily\scriptsize         if "np." in stripped or "numpy" in stripped:}}
\par\smallskip
This performs vectorized numerical work. The fixed-point golden models use NumPy so score, margin, and update calculations are explicit and repeatable.
\paragraph{Line 722}
\noindent\colorbox{gray!8}{\parbox{\dimexpr\linewidth-2\fboxsep\relax}{\raggedright\ttfamily\scriptsize             return "This performs vectorized numerical work. The fixed-point golden models use NumPy so score, margin, and update calculations are explicit and repeatable."}}
\par\smallskip
This line handles the online-learning rule or its validation fields. It is central to proving that the FPGA update path is not bypassed.
\paragraph{Line 723}
\noindent\colorbox{gray!8}{\parbox{\dimexpr\linewidth-2\fboxsep\relax}{\raggedright\ttfamily\scriptsize         if "weights" in stripped or "bias" in stripped:}}
\par\smallskip
This line handles the classifier parameters. Those arrays mirror the RTL memories, so mismatches reveal an export, quantization, or hardware-update bug.
\paragraph{Line 724}
\noindent\colorbox{gray!8}{\parbox{\dimexpr\linewidth-2\fboxsep\relax}{\raggedright\ttfamily\scriptsize             return "This line handles the classifier parameters. Those arrays mirror the RTL memories, so mismatches reveal an export, quantization, or hardware-update bug."}}
\par\smallskip
This line handles the online-learning rule or its validation fields. It is central to proving that the FPGA update path is not bypassed.
\paragraph{Line 725}
\noindent\colorbox{gray!8}{\parbox{\dimexpr\linewidth-2\fboxsep\relax}{\raggedright\ttfamily\scriptsize         if "scores" in stripped:}}
\par\smallskip
This line handles class scores. The score vector is the direct software equivalent of the RTL MAC loop over weights and features.
\paragraph{Line 726}
\noindent\colorbox{gray!8}{\parbox{\dimexpr\linewidth-2\fboxsep\relax}{\raggedright\ttfamily\scriptsize             return "This line handles class scores. The score vector is the direct software equivalent of the RTL MAC loop over weights and features."}}
\par\smallskip
This line handles the classifier parameters. Those arrays mirror the RTL memories, so mismatches reveal an export, quantization, or hardware-update bug.
\paragraph{Line 727}
\noindent\colorbox{gray!8}{\parbox{\dimexpr\linewidth-2\fboxsep\relax}{\raggedright\ttfamily\scriptsize         if "pred" in stripped:}}
\par\smallskip
This line handles predicted labels. It checks that argmax behavior in Python and RTL chooses the same class.
\paragraph{Line 728}
\noindent\colorbox{gray!8}{\parbox{\dimexpr\linewidth-2\fboxsep\relax}{\raggedright\ttfamily\scriptsize             return "This line handles predicted labels. It checks that argmax behavior in Python and RTL chooses the same class."}}
\par\smallskip
This line handles predicted labels. It checks that argmax behavior in Python and RTL chooses the same class.
\paragraph{Line 729}
\noindent\colorbox{gray!8}{\parbox{\dimexpr\linewidth-2\fboxsep\relax}{\raggedright\ttfamily\scriptsize         if "update" in stripped:}}
\par\smallskip
This line handles the online-learning rule or its validation fields. It is central to proving that the FPGA update path is not bypassed.
\paragraph{Line 730}
\noindent\colorbox{gray!8}{\parbox{\dimexpr\linewidth-2\fboxsep\relax}{\raggedright\ttfamily\scriptsize             return "This line handles the online-learning rule or its validation fields. It is central to proving that the FPGA update path is not bypassed."}}
\par\smallskip
This line handles the online-learning rule or its validation fields. It is central to proving that the FPGA update path is not bypassed.
\paragraph{Line 731}
\noindent\colorbox{gray!8}{\parbox{\dimexpr\linewidth-2\fboxsep\relax}{\raggedright\ttfamily\scriptsize         if "plt." in stripped or "PdfPages" in stripped:}}
\par\smallskip
This creates a debug visualization. The plot is derived from generated CSV or Vivado report data so it can be regenerated after RTL changes.
\paragraph{Line 732}
\noindent\colorbox{gray!8}{\parbox{\dimexpr\linewidth-2\fboxsep\relax}{\raggedright\ttfamily\scriptsize             return "This creates a debug visualization. The plot is derived from generated CSV or Vivado report data so it can be regenerated after RTL changes."}}
\par\smallskip
This Python line participates in the export, validation, or report-generation pipeline. It is executed on the host PC to prepare data for RTL or to check RTL outputs.
\paragraph{Line 733}
\noindent\colorbox{gray!8}{\parbox{\dimexpr\linewidth-2\fboxsep\relax}{\raggedright\ttfamily\scriptsize         if stripped.startswith("if \_\_name\_\_"):}}
\par\smallskip
This Python line participates in the export, validation, or report-generation pipeline. It is executed on the host PC to prepare data for RTL or to check RTL outputs.
\paragraph{Line 734}
\noindent\colorbox{gray!8}{\parbox{\dimexpr\linewidth-2\fboxsep\relax}{\raggedright\ttfamily\scriptsize             return "This Python idiom runs main() only when the file is executed as a script. It still allows the file's functions to be imported for tests or reuse."}}
\par\smallskip
This Python line participates in the export, validation, or report-generation pipeline. It is executed on the host PC to prepare data for RTL or to check RTL outputs.
\paragraph{Line 735}
\noindent\colorbox{gray!8}{\parbox{\dimexpr\linewidth-2\fboxsep\relax}{\raggedright\ttfamily\scriptsize         return "This Python line participates in the export, validation, or report-generation pipeline. It is executed on the host PC to prepare data for RTL or to check RTL outputs."}}
\par\smallskip
This Python line participates in the export, validation, or report-generation pipeline. It is executed on the host PC to prepare data for RTL or to check RTL outputs.
\paragraph{Line 736}
\noindent\colorbox{gray!8}{\parbox{\dimexpr\linewidth-2\fboxsep\relax}{\raggedright\ttfamily\scriptsize     if suffix == ".ps1":}}
\par\smallskip
This Python line participates in the export, validation, or report-generation pipeline. It is executed on the host PC to prepare data for RTL or to check RTL outputs.
\paragraph{Line 737}
\noindent\colorbox{gray!8}{\parbox{\dimexpr\linewidth-2\fboxsep\relax}{\raggedright\ttfamily\scriptsize         if stripped.startswith("param"):}}
\par\smallskip
This Python line participates in the export, validation, or report-generation pipeline. It is executed on the host PC to prepare data for RTL or to check RTL outputs.
\paragraph{Line 738}
\noindent\colorbox{gray!8}{\parbox{\dimexpr\linewidth-2\fboxsep\relax}{\raggedright\ttfamily\scriptsize             return "This begins the PowerShell parameter block. It exposes repeatable knobs such as sample count, Vivado path, board part, and skip switches."}}
\par\smallskip
This Python line participates in the export, validation, or report-generation pipeline. It is executed on the host PC to prepare data for RTL or to check RTL outputs.
\paragraph{Line 739}
\noindent\colorbox{gray!8}{\parbox{\dimexpr\linewidth-2\fboxsep\relax}{\raggedright\ttfamily\scriptsize         if stripped.startswith("["):}}
\par\smallskip
This Python line participates in the export, validation, or report-generation pipeline. It is executed on the host PC to prepare data for RTL or to check RTL outputs.
\paragraph{Line 740}
\noindent\colorbox{gray!8}{\parbox{\dimexpr\linewidth-2\fboxsep\relax}{\raggedright\ttfamily\scriptsize             return "This declares one typed PowerShell parameter. Strong typing catches accidental string/numeric mistakes before the long training or Vivado flow starts."}}
\par\smallskip
This Python line participates in the export, validation, or report-generation pipeline. It is executed on the host PC to prepare data for RTL or to check RTL outputs.
\paragraph{Line 741}
\noindent\colorbox{gray!8}{\parbox{\dimexpr\linewidth-2\fboxsep\relax}{\raggedright\ttfamily\scriptsize         if "\$ErrorActionPreference" in stripped:}}
\par\smallskip
This Python line participates in the export, validation, or report-generation pipeline. It is executed on the host PC to prepare data for RTL or to check RTL outputs.
\paragraph{Line 742}
\noindent\colorbox{gray!8}{\parbox{\dimexpr\linewidth-2\fboxsep\relax}{\raggedright\ttfamily\scriptsize             return "This makes failures stop the script instead of continuing after a broken export, simulation, or synthesis step."}}
\par\smallskip
This Python line participates in the export, validation, or report-generation pipeline. It is executed on the host PC to prepare data for RTL or to check RTL outputs.
\paragraph{Line 743}
\noindent\colorbox{gray!8}{\parbox{\dimexpr\linewidth-2\fboxsep\relax}{\raggedright\ttfamily\scriptsize         if "Assert-LastExitCode" in stripped:}}
\par\smallskip
This Python line participates in the export, validation, or report-generation pipeline. It is executed on the host PC to prepare data for RTL or to check RTL outputs.
\paragraph{Line 744}
\noindent\colorbox{gray!8}{\parbox{\dimexpr\linewidth-2\fboxsep\relax}{\raggedright\ttfamily\scriptsize             return "This calls the helper that turns a failed external tool exit code into a hard pipeline failure."}}
\par\smallskip
This Python line participates in the export, validation, or report-generation pipeline. It is executed on the host PC to prepare data for RTL or to check RTL outputs.
\paragraph{Line 745}
\noindent\colorbox{gray!8}{\parbox{\dimexpr\linewidth-2\fboxsep\relax}{\raggedright\ttfamily\scriptsize         if "Vivado" in stripped or "xvlog" in stripped or "xelab" in stripped or "xsim" in stripped:}}
\par\smallskip
This Python line participates in the export, validation, or report-generation pipeline. It is executed on the host PC to prepare data for RTL or to check RTL outputs.
\paragraph{Line 746}
\noindent\colorbox{gray!8}{\parbox{\dimexpr\linewidth-2\fboxsep\relax}{\raggedright\ttfamily\scriptsize             return "This line invokes or locates Vivado simulator/implementation tools. The flow uses the installed AMD toolchain on D: rather than a substitute simulator."}}
\par\smallskip
This Python line participates in the export, validation, or report-generation pipeline. It is executed on the host PC to prepare data for RTL or to check RTL outputs.
\paragraph{Line 747}
\noindent\colorbox{gray!8}{\parbox{\dimexpr\linewidth-2\fboxsep\relax}{\raggedright\ttfamily\scriptsize         if "python" in stripped.lower():}}
\par\smallskip
This Python line participates in the export, validation, or report-generation pipeline. It is executed on the host PC to prepare data for RTL or to check RTL outputs.
\paragraph{Line 748}
\noindent\colorbox{gray!8}{\parbox{\dimexpr\linewidth-2\fboxsep\relax}{\raggedright\ttfamily\scriptsize             return "This line invokes the Python environment used for export, validation, or report generation. It ties host-side math to the hardware artifacts."}}
\par\smallskip
This Python line participates in the export, validation, or report-generation pipeline. It is executed on the host PC to prepare data for RTL or to check RTL outputs.
\paragraph{Line 749}
\noindent\colorbox{gray!8}{\parbox{\dimexpr\linewidth-2\fboxsep\relax}{\raggedright\ttfamily\scriptsize         return "This PowerShell line orchestrates the reproducible full flow. It exists so changes can be rerun without manually retyping many commands."}}
\par\smallskip
This Python line participates in the export, validation, or report-generation pipeline. It is executed on the host PC to prepare data for RTL or to check RTL outputs.
\paragraph{Line 750}
\noindent\colorbox{gray!8}{\parbox{\dimexpr\linewidth-2\fboxsep\relax}{\raggedright\ttfamily\scriptsize     if suffix == ".tcl":}}
\par\smallskip
This Python line participates in the export, validation, or report-generation pipeline. It is executed on the host PC to prepare data for RTL or to check RTL outputs.
\paragraph{Line 751}
\noindent\colorbox{gray!8}{\parbox{\dimexpr\linewidth-2\fboxsep\relax}{\raggedright\ttfamily\scriptsize         if stripped.startswith("read\_verilog"):}}
\par\smallskip
This Python line participates in the export, validation, or report-generation pipeline. It is executed on the host PC to prepare data for RTL or to check RTL outputs.
\paragraph{Line 752}
\noindent\colorbox{gray!8}{\parbox{\dimexpr\linewidth-2\fboxsep\relax}{\raggedright\ttfamily\scriptsize             return "This tells Vivado which SystemVerilog files to synthesize. The wrapper and accelerator are both read so the routed report matches the PYNQ-facing design."}}
\par\smallskip
This Python line participates in the export, validation, or report-generation pipeline. It is executed on the host PC to prepare data for RTL or to check RTL outputs.
\paragraph{Line 753}
\noindent\colorbox{gray!8}{\parbox{\dimexpr\linewidth-2\fboxsep\relax}{\raggedright\ttfamily\scriptsize         if stripped.startswith("synth\_design"):}}
\par\smallskip
This Python line participates in the export, validation, or report-generation pipeline. It is executed on the host PC to prepare data for RTL or to check RTL outputs.
\paragraph{Line 754}
\noindent\colorbox{gray!8}{\parbox{\dimexpr\linewidth-2\fboxsep\relax}{\raggedright\ttfamily\scriptsize             return "This starts synthesis for the selected FPGA part and top module. Synthesis maps RTL arrays, arithmetic, and FSMs into FPGA primitives."}}
\par\smallskip
This Python line participates in the export, validation, or report-generation pipeline. It is executed on the host PC to prepare data for RTL or to check RTL outputs.
\paragraph{Line 755}
\noindent\colorbox{gray!8}{\parbox{\dimexpr\linewidth-2\fboxsep\relax}{\raggedright\ttfamily\scriptsize         if stripped.startswith("create\_clock"):}}
\par\smallskip
This Python line participates in the export, validation, or report-generation pipeline. It is executed on the host PC to prepare data for RTL or to check RTL outputs.
\paragraph{Line 756}
\noindent\colorbox{gray!8}{\parbox{\dimexpr\linewidth-2\fboxsep\relax}{\raggedright\ttfamily\scriptsize             return "This creates the 100 MHz timing constraint. Vivado uses it to decide whether every register-to-register path is fast enough."}}
\par\smallskip
This Python line participates in the export, validation, or report-generation pipeline. It is executed on the host PC to prepare data for RTL or to check RTL outputs.
\paragraph{Line 757}
\noindent\colorbox{gray!8}{\parbox{\dimexpr\linewidth-2\fboxsep\relax}{\raggedright\ttfamily\scriptsize         if stripped.startswith("place\_design") or stripped.startswith("route\_design") or stripped.startswith("opt\_design"):}}
\par\smallskip
This Python line participates in the export, validation, or report-generation pipeline. It is executed on the host PC to prepare data for RTL or to check RTL outputs.
\paragraph{Line 758}
\noindent\colorbox{gray!8}{\parbox{\dimexpr\linewidth-2\fboxsep\relax}{\raggedright\ttfamily\scriptsize             return "This runs an implementation stage. These commands turn the synthesized netlist into a placed and routed design with real timing and power estimates."}}
\par\smallskip
This Python line participates in the export, validation, or report-generation pipeline. It is executed on the host PC to prepare data for RTL or to check RTL outputs.
\paragraph{Line 759}
\noindent\colorbox{gray!8}{\parbox{\dimexpr\linewidth-2\fboxsep\relax}{\raggedright\ttfamily\scriptsize         if stripped.startswith("report\_") or stripped.startswith("write\_checkpoint"):}}
\par\smallskip
This Python line participates in the export, validation, or report-generation pipeline. It is executed on the host PC to prepare data for RTL or to check RTL outputs.
\paragraph{Line 760}
\noindent\colorbox{gray!8}{\parbox{\dimexpr\linewidth-2\fboxsep\relax}{\raggedright\ttfamily\scriptsize             return "This emits a signoff artifact. The generated reports are the source for utilization, timing, and power numbers in the debug PDF."}}
\par\smallskip
This Python line participates in the export, validation, or report-generation pipeline. It is executed on the host PC to prepare data for RTL or to check RTL outputs.
\paragraph{Line 761}
\noindent\colorbox{gray!8}{\parbox{\dimexpr\linewidth-2\fboxsep\relax}{\raggedright\ttfamily\scriptsize         return "This TCL line controls the Vivado batch implementation flow. It makes the hardware report reproducible from the command line."}}
\par\smallskip
This Python line participates in the export, validation, or report-generation pipeline. It is executed on the host PC to prepare data for RTL or to check RTL outputs.
\paragraph{Line 762}
\noindent\colorbox{gray!8}{\parbox{\dimexpr\linewidth-2\fboxsep\relax}{\raggedright\ttfamily\scriptsize     return "This line is part of the accelerator support files and is included so the full rerun flow is documented, not just the RTL."}}
\par\smallskip
This Python line participates in the export, validation, or report-generation pipeline. It is executed on the host PC to prepare data for RTL or to check RTL outputs.
\paragraph{Line 763}
\noindent\colorbox{gray!8}{\parbox{\dimexpr\linewidth-2\fboxsep\relax}{\raggedright\ttfamily\scriptsize \emph{blank line}}}
\par\smallskip
This blank line is not hardware or software behavior. It separates nearby declarations or operations so the file remains readable and reviewable.
\paragraph{Line 764}
\noindent\colorbox{gray!8}{\parbox{\dimexpr\linewidth-2\fboxsep\relax}{\raggedright\ttfamily\scriptsize \emph{blank line}}}
\par\smallskip
This blank line is not hardware or software behavior. It separates nearby declarations or operations so the file remains readable and reviewable.
\paragraph{Line 765}
\noindent\colorbox{gray!8}{\parbox{\dimexpr\linewidth-2\fboxsep\relax}{\raggedright\ttfamily\scriptsize def generate\_explainer\_tex() -> Path:}}
\par\smallskip
This defines a Python function. The function packages one repeatable step such as loading a memory file, simulating the golden learning rule, validating CSVs, or generating plots.
\paragraph{Line 766}
\noindent\colorbox{gray!8}{\parbox{\dimexpr\linewidth-2\fboxsep\relax}{\raggedright\ttfamily\scriptsize     pdir = project\_dir()}}
\par\smallskip
This Python line participates in the export, validation, or report-generation pipeline. It is executed on the host PC to prepare data for RTL or to check RTL outputs.
\paragraph{Line 767}
\noindent\colorbox{gray!8}{\parbox{\dimexpr\linewidth-2\fboxsep\relax}{\raggedright\ttfamily\scriptsize     ddir = debug\_dir()}}
\par\smallskip
This Python line participates in the export, validation, or report-generation pipeline. It is executed on the host PC to prepare data for RTL or to check RTL outputs.
\paragraph{Line 768}
\noindent\colorbox{gray!8}{\parbox{\dimexpr\linewidth-2\fboxsep\relax}{\raggedright\ttfamily\scriptsize     ddir.mkdir(parents=True, exist\_ok=True)}}
\par\smallskip
This Python line participates in the export, validation, or report-generation pipeline. It is executed on the host PC to prepare data for RTL or to check RTL outputs.
\paragraph{Line 769}
\noindent\colorbox{gray!8}{\parbox{\dimexpr\linewidth-2\fboxsep\relax}{\raggedright\ttfamily\scriptsize     odir = out\_dir()}}
\par\smallskip
This Python line participates in the export, validation, or report-generation pipeline. It is executed on the host PC to prepare data for RTL or to check RTL outputs.
\paragraph{Line 770}
\noindent\colorbox{gray!8}{\parbox{\dimexpr\linewidth-2\fboxsep\relax}{\raggedright\ttfamily\scriptsize     rtl\_val = read\_json(odir / "spikingformer\_head\_rtl\_validation.json")}}
\par\smallskip
This Python line participates in the export, validation, or report-generation pipeline. It is executed on the host PC to prepare data for RTL or to check RTL outputs.
\paragraph{Line 771}
\noindent\colorbox{gray!8}{\parbox{\dimexpr\linewidth-2\fboxsep\relax}{\raggedright\ttfamily\scriptsize     learn\_val = read\_json(odir / "spikingformer\_head\_onchip\_learning\_validation.json")}}
\par\smallskip
This Python line participates in the export, validation, or report-generation pipeline. It is executed on the host PC to prepare data for RTL or to check RTL outputs.
\paragraph{Line 772}
\noindent\colorbox{gray!8}{\parbox{\dimexpr\linewidth-2\fboxsep\relax}{\raggedright\ttfamily\scriptsize     mapping = read\_json(odir / "spikingformer\_head\_mapping\_metrics.json")}}
\par\smallskip
This Python line participates in the export, validation, or report-generation pipeline. It is executed on the host PC to prepare data for RTL or to check RTL outputs.
\paragraph{Line 773}
\noindent\colorbox{gray!8}{\parbox{\dimexpr\linewidth-2\fboxsep\relax}{\raggedright\ttfamily\scriptsize     vivado = parse\_vivado\_reports(odir / "vivado")}}
\par\smallskip
This Python line participates in the export, validation, or report-generation pipeline. It is executed on the host PC to prepare data for RTL or to check RTL outputs.
\paragraph{Line 774}
\noindent\colorbox{gray!8}{\parbox{\dimexpr\linewidth-2\fboxsep\relax}{\raggedright\ttfamily\scriptsize     acc\_pct = float(rtl\_val.get("accuracy", float("nan"))) * 100.0}}
\par\smallskip
This Python line participates in the export, validation, or report-generation pipeline. It is executed on the host PC to prepare data for RTL or to check RTL outputs.
\paragraph{Line 775}
\noindent\colorbox{gray!8}{\parbox{\dimexpr\linewidth-2\fboxsep\relax}{\raggedright\ttfamily\scriptsize     sample\_count = int(rtl\_val.get("samples", 0) or 0)}}
\par\smallskip
This Python line participates in the export, validation, or report-generation pipeline. It is executed on the host PC to prepare data for RTL or to check RTL outputs.
\paragraph{Line 776}
\noindent\colorbox{gray!8}{\parbox{\dimexpr\linewidth-2\fboxsep\relax}{\raggedright\ttfamily\scriptsize     update\_count = int(learn\_val.get("updates", 0) or 0)}}
\par\smallskip
This line handles the online-learning rule or its validation fields. It is central to proving that the FPGA update path is not bypassed.
\paragraph{Line 777}
\noindent\colorbox{gray!8}{\parbox{\dimexpr\linewidth-2\fboxsep\relax}{\raggedright\ttfamily\scriptsize     infer\_skip\_abs\_thresh = int(rtl\_val.get("infer\_skip\_abs\_thresh", mapping.get("infer\_skip\_abs\_thresh", 0)) or 0)}}
\par\smallskip
This Python line participates in the export, validation, or report-generation pipeline. It is executed on the host PC to prepare data for RTL or to check RTL outputs.
\paragraph{Line 778}
\noindent\colorbox{gray!8}{\parbox{\dimexpr\linewidth-2\fboxsep\relax}{\raggedright\ttfamily\scriptsize     feature\_mac\_skip\_fraction = float(rtl\_val.get("feature\_mac\_skip\_fraction", mapping.get("feature\_mac\_skip\_fraction", 0.0)) or 0.0) * 100.0}}
\par\smallskip
This Python line participates in the export, validation, or report-generation pipeline. It is executed on the host PC to prepare data for RTL or to check RTL outputs.
\paragraph{Line 779}
\noindent\colorbox{gray!8}{\parbox{\dimexpr\linewidth-2\fboxsep\relax}{\raggedright\ttfamily\scriptsize     files = [}}
\par\smallskip
This Python line participates in the export, validation, or report-generation pipeline. It is executed on the host PC to prepare data for RTL or to check RTL outputs.
\paragraph{Line 780}
\noindent\colorbox{gray!8}{\parbox{\dimexpr\linewidth-2\fboxsep\relax}{\raggedright\ttfamily\scriptsize         pdir / "rtl" / "cifar10\_spikingformer\_head\_accelerator.sv",}}
\par\smallskip
This Python line participates in the export, validation, or report-generation pipeline. It is executed on the host PC to prepare data for RTL or to check RTL outputs.
\paragraph{Line 781}
\noindent\colorbox{gray!8}{\parbox{\dimexpr\linewidth-2\fboxsep\relax}{\raggedright\ttfamily\scriptsize         pdir / "rtl" / "cifar10\_spikingformer\_head\_pynq\_axi\_stream\_top.sv",}}
\par\smallskip
This Python line participates in the export, validation, or report-generation pipeline. It is executed on the host PC to prepare data for RTL or to check RTL outputs.
\paragraph{Line 782}
\noindent\colorbox{gray!8}{\parbox{\dimexpr\linewidth-2\fboxsep\relax}{\raggedright\ttfamily\scriptsize         pdir / "rtl" / "tb\_spikingformer\_head\_learning\_unit.sv",}}
\par\smallskip
This Python line participates in the export, validation, or report-generation pipeline. It is executed on the host PC to prepare data for RTL or to check RTL outputs.
\paragraph{Line 783}
\noindent\colorbox{gray!8}{\parbox{\dimexpr\linewidth-2\fboxsep\relax}{\raggedright\ttfamily\scriptsize         pdir / "rtl" / "tb\_cifar10\_spikingformer\_head\_accelerator.sv",}}
\par\smallskip
This Python line participates in the export, validation, or report-generation pipeline. It is executed on the host PC to prepare data for RTL or to check RTL outputs.
\paragraph{Line 784}
\noindent\colorbox{gray!8}{\parbox{\dimexpr\linewidth-2\fboxsep\relax}{\raggedright\ttfamily\scriptsize         pdir / "rtl" / "tb\_cifar10\_spikingformer\_head\_onchip\_learning.sv",}}
\par\smallskip
This Python line participates in the export, validation, or report-generation pipeline. It is executed on the host PC to prepare data for RTL or to check RTL outputs.
\paragraph{Line 785}
\noindent\colorbox{gray!8}{\parbox{\dimexpr\linewidth-2\fboxsep\relax}{\raggedright\ttfamily\scriptsize         pdir / "export\_spikingformer\_head\_accelerator.py",}}
\par\smallskip
This Python line participates in the export, validation, or report-generation pipeline. It is executed on the host PC to prepare data for RTL or to check RTL outputs.
\paragraph{Line 786}
\noindent\colorbox{gray!8}{\parbox{\dimexpr\linewidth-2\fboxsep\relax}{\raggedright\ttfamily\scriptsize         pdir / "validate\_spikingformer\_head\_rtl.py",}}
\par\smallskip
This Python line participates in the export, validation, or report-generation pipeline. It is executed on the host PC to prepare data for RTL or to check RTL outputs.
\paragraph{Line 787}
\noindent\colorbox{gray!8}{\parbox{\dimexpr\linewidth-2\fboxsep\relax}{\raggedright\ttfamily\scriptsize         pdir / "validate\_spikingformer\_head\_onchip\_learning.py",}}
\par\smallskip
This Python line participates in the export, validation, or report-generation pipeline. It is executed on the host PC to prepare data for RTL or to check RTL outputs.
\paragraph{Line 788}
\noindent\colorbox{gray!8}{\parbox{\dimexpr\linewidth-2\fboxsep\relax}{\raggedright\ttfamily\scriptsize         pdir / "generate\_spikingformer\_debug\_reports.py",}}
\par\smallskip
This Python line participates in the export, validation, or report-generation pipeline. It is executed on the host PC to prepare data for RTL or to check RTL outputs.
\paragraph{Line 789}
\noindent\colorbox{gray!8}{\parbox{\dimexpr\linewidth-2\fboxsep\relax}{\raggedright\ttfamily\scriptsize         pdir / "run\_spikingformer\_head\_accelerator.ps1",}}
\par\smallskip
This Python line participates in the export, validation, or report-generation pipeline. It is executed on the host PC to prepare data for RTL or to check RTL outputs.
\paragraph{Line 790}
\noindent\colorbox{gray!8}{\parbox{\dimexpr\linewidth-2\fboxsep\relax}{\raggedright\ttfamily\scriptsize         pdir / "spikingformer\_head\_synth.tcl",}}
\par\smallskip
This Python line participates in the export, validation, or report-generation pipeline. It is executed on the host PC to prepare data for RTL or to check RTL outputs.
\paragraph{Line 791}
\noindent\colorbox{gray!8}{\parbox{\dimexpr\linewidth-2\fboxsep\relax}{\raggedright\ttfamily\scriptsize     ]}}
\par\smallskip
This Python line participates in the export, validation, or report-generation pipeline. It is executed on the host PC to prepare data for RTL or to check RTL outputs.
\paragraph{Line 792}
\noindent\colorbox{gray!8}{\parbox{\dimexpr\linewidth-2\fboxsep\relax}{\raggedright\ttfamily\scriptsize     tex\_path = ddir / "neuraedge\_full\_accelerator\_line\_by\_line\_explainer.tex"}}
\par\smallskip
This Python line participates in the export, validation, or report-generation pipeline. It is executed on the host PC to prepare data for RTL or to check RTL outputs.
\paragraph{Line 793}
\noindent\colorbox{gray!8}{\parbox{\dimexpr\linewidth-2\fboxsep\relax}{\raggedright\ttfamily\scriptsize     lines: list[str] = []}}
\par\smallskip
This Python line participates in the export, validation, or report-generation pipeline. It is executed on the host PC to prepare data for RTL or to check RTL outputs.
\paragraph{Line 794}
\noindent\colorbox{gray!8}{\parbox{\dimexpr\linewidth-2\fboxsep\relax}{\raggedright\ttfamily\scriptsize     lines.extend(}}
\par\smallskip
This Python line participates in the export, validation, or report-generation pipeline. It is executed on the host PC to prepare data for RTL or to check RTL outputs.
\paragraph{Line 795}
\noindent\colorbox{gray!8}{\parbox{\dimexpr\linewidth-2\fboxsep\relax}{\raggedright\ttfamily\scriptsize         [}}
\par\smallskip
This Python line participates in the export, validation, or report-generation pipeline. It is executed on the host PC to prepare data for RTL or to check RTL outputs.
\paragraph{Line 796}
\noindent\colorbox{gray!8}{\parbox{\dimexpr\linewidth-2\fboxsep\relax}{\raggedright\ttfamily\scriptsize             r"\textbackslash{}documentclass[10pt]\{article\}",}}
\par\smallskip
This Python line participates in the export, validation, or report-generation pipeline. It is executed on the host PC to prepare data for RTL or to check RTL outputs.
\paragraph{Line 797}
\noindent\colorbox{gray!8}{\parbox{\dimexpr\linewidth-2\fboxsep\relax}{\raggedright\ttfamily\scriptsize             r"\textbackslash{}usepackage[margin=0.65in]\{geometry\}",}}
\par\smallskip
This Python line participates in the export, validation, or report-generation pipeline. It is executed on the host PC to prepare data for RTL or to check RTL outputs.
\paragraph{Line 798}
\noindent\colorbox{gray!8}{\parbox{\dimexpr\linewidth-2\fboxsep\relax}{\raggedright\ttfamily\scriptsize             r"\textbackslash{}usepackage[T1]\{fontenc\}",}}
\par\smallskip
This Python line participates in the export, validation, or report-generation pipeline. It is executed on the host PC to prepare data for RTL or to check RTL outputs.
\paragraph{Line 799}
\noindent\colorbox{gray!8}{\parbox{\dimexpr\linewidth-2\fboxsep\relax}{\raggedright\ttfamily\scriptsize             r"\textbackslash{}usepackage\{hyperref\}",}}
\par\smallskip
This Python line participates in the export, validation, or report-generation pipeline. It is executed on the host PC to prepare data for RTL or to check RTL outputs.
\paragraph{Line 800}
\noindent\colorbox{gray!8}{\parbox{\dimexpr\linewidth-2\fboxsep\relax}{\raggedright\ttfamily\scriptsize             r"\textbackslash{}usepackage\{xcolor\}",}}
\par\smallskip
This Python line participates in the export, validation, or report-generation pipeline. It is executed on the host PC to prepare data for RTL or to check RTL outputs.
\paragraph{Line 801}
\noindent\colorbox{gray!8}{\parbox{\dimexpr\linewidth-2\fboxsep\relax}{\raggedright\ttfamily\scriptsize             r"\textbackslash{}usepackage\{listings\}",}}
\par\smallskip
This Python line participates in the export, validation, or report-generation pipeline. It is executed on the host PC to prepare data for RTL or to check RTL outputs.
\paragraph{Line 802}
\noindent\colorbox{gray!8}{\parbox{\dimexpr\linewidth-2\fboxsep\relax}{\raggedright\ttfamily\scriptsize             r"\textbackslash{}lstset\{basicstyle=\textbackslash{}ttfamily\textbackslash{}scriptsize,breaklines=true,columns=fullflexible,keepspaces=true,showstringspaces=false\}",}}
\par\smallskip
This Python line participates in the export, validation, or report-generation pipeline. It is executed on the host PC to prepare data for RTL or to check RTL outputs.
\paragraph{Line 803}
\noindent\colorbox{gray!8}{\parbox{\dimexpr\linewidth-2\fboxsep\relax}{\raggedright\ttfamily\scriptsize             r"\textbackslash{}title\{NeuraEdge CIFAR-10 Neuromorphic Accelerator: Full Pipeline and Line-by-Line Explainer\}",}}
\par\smallskip
This Python line participates in the export, validation, or report-generation pipeline. It is executed on the host PC to prepare data for RTL or to check RTL outputs.
\paragraph{Line 804}
\noindent\colorbox{gray!8}{\parbox{\dimexpr\linewidth-2\fboxsep\relax}{\raggedright\ttfamily\scriptsize             r"\textbackslash{}author\{Generated from repository source and measured Vivado/XSim artifacts\}",}}
\par\smallskip
This Python line participates in the export, validation, or report-generation pipeline. It is executed on the host PC to prepare data for RTL or to check RTL outputs.
\paragraph{Line 805}
\noindent\colorbox{gray!8}{\parbox{\dimexpr\linewidth-2\fboxsep\relax}{\raggedright\ttfamily\scriptsize             r"\textbackslash{}date\{\textbackslash{}today\}",}}
\par\smallskip
This Python line participates in the export, validation, or report-generation pipeline. It is executed on the host PC to prepare data for RTL or to check RTL outputs.
\paragraph{Line 806}
\noindent\colorbox{gray!8}{\parbox{\dimexpr\linewidth-2\fboxsep\relax}{\raggedright\ttfamily\scriptsize             r"\textbackslash{}begin\{document\}",}}
\par\smallskip
This Python line participates in the export, validation, or report-generation pipeline. It is executed on the host PC to prepare data for RTL or to check RTL outputs.
\paragraph{Line 807}
\noindent\colorbox{gray!8}{\parbox{\dimexpr\linewidth-2\fboxsep\relax}{\raggedright\ttfamily\scriptsize             r"\textbackslash{}maketitle",}}
\par\smallskip
This Python line participates in the export, validation, or report-generation pipeline. It is executed on the host PC to prepare data for RTL or to check RTL outputs.
\paragraph{Line 808}
\noindent\colorbox{gray!8}{\parbox{\dimexpr\linewidth-2\fboxsep\relax}{\raggedright\ttfamily\scriptsize             r"\textbackslash{}section\{What Was Built\}",}}
\par\smallskip
This Python line participates in the export, validation, or report-generation pipeline. It is executed on the host PC to prepare data for RTL or to check RTL outputs.
\paragraph{Line 809}
\noindent\colorbox{gray!8}{\parbox{\dimexpr\linewidth-2\fboxsep\relax}{\raggedright\ttfamily\scriptsize             "This document describes the implemented CIFAR-10 accelerator path. The trained Spikingformer backbone is used to export 384-dimensional fixed-point features and a trained classifier head. The FPGA accelerator stores one feature vector, applies a measured low-magnitude feature activity gate, computes ten fixed-point class scores with an on-chip weight memory and one DSP MAC datapath, returns the predicted label, and includes a margin-gated on-chip learning path that updates the true-label and predicted-label classifier rows in FPGA memory.",}}
\par\smallskip
This line handles class scores. The score vector is the direct software equivalent of the RTL MAC loop over weights and features.
\paragraph{Line 810}
\noindent\colorbox{gray!8}{\parbox{\dimexpr\linewidth-2\fboxsep\relax}{\raggedright\ttfamily\scriptsize             r"\textbackslash{}section\{Measured Verification Summary\}",}}
\par\smallskip
This Python line participates in the export, validation, or report-generation pipeline. It is executed on the host PC to prepare data for RTL or to check RTL outputs.
\paragraph{Line 811}
\noindent\colorbox{gray!8}{\parbox{\dimexpr\linewidth-2\fboxsep\relax}{\raggedright\ttfamily\scriptsize             r"\textbackslash{}begin\{itemize\}",}}
\par\smallskip
This Python line participates in the export, validation, or report-generation pipeline. It is executed on the host PC to prepare data for RTL or to check RTL outputs.
\paragraph{Line 812}
\noindent\colorbox{gray!8}{\parbox{\dimexpr\linewidth-2\fboxsep\relax}{\raggedright\ttfamily\scriptsize             rf"\textbackslash{}item Full inference XSim run: \{sample\_count:,\} CIFAR-10 samples, \{acc\_pct:.2f\}\textbackslash{}\% top-1 accuracy, exact score match against the fixed-point Python reference.",}}
\par\smallskip
This Python line participates in the export, validation, or report-generation pipeline. It is executed on the host PC to prepare data for RTL or to check RTL outputs.
\paragraph{Line 813}
\noindent\colorbox{gray!8}{\parbox{\dimexpr\linewidth-2\fboxsep\relax}{\raggedright\ttfamily\scriptsize             rf"\textbackslash{}item Feature activity gate: abs(feature) \$\textbackslash{}leq\$ \{infer\_skip\_abs\_thresh\} is treated as zero before inference and learning; the exported test set skips \{feature\_mac\_skip\_fraction:.2f\}\textbackslash{}\% of feature-MAC opportunities.",}}
\par\smallskip
This Python line participates in the export, validation, or report-generation pipeline. It is executed on the host PC to prepare data for RTL or to check RTL outputs.
\paragraph{Line 814}
\noindent\colorbox{gray!8}{\parbox{\dimexpr\linewidth-2\fboxsep\relax}{\raggedright\ttfamily\scriptsize             rf"\textbackslash{}item Full on-chip-learning XSim run: \{sample\_count:,\} samples, \{update\_count\} RTL weight/bias updates, zero score mismatches, zero prediction mismatches, zero update-flag mismatches, and zero update-count mismatches against the independent Python online-learning golden model.",}}
\par\smallskip
This line handles the classifier parameters. Those arrays mirror the RTL memories, so mismatches reveal an export, quantization, or hardware-update bug.
\paragraph{Line 815}
\noindent\colorbox{gray!8}{\parbox{\dimexpr\linewidth-2\fboxsep\relax}{\raggedright\ttfamily\scriptsize             rf"\textbackslash{}item Vivado implementation target: PYNQ Zynq-7020, xc7z020clg400-1, 100 MHz. Routed timing met with WNS \{vivado['wns\_ns']:.3f\} ns and WHS \{vivado['whs\_ns']:.3f\} ns.",}}
\par\smallskip
This Python line participates in the export, validation, or report-generation pipeline. It is executed on the host PC to prepare data for RTL or to check RTL outputs.
\paragraph{Line 816}
\noindent\colorbox{gray!8}{\parbox{\dimexpr\linewidth-2\fboxsep\relax}{\raggedright\ttfamily\scriptsize             rf"\textbackslash{}item Resource result: \{vivado['lut']\} LUTs, \{vivado['regs']\} FFs, \{vivado['bram']\} RAMB36E1 blocks, and \{vivado['dsp']\} DSP48E1.",}}
\par\smallskip
This Python line participates in the export, validation, or report-generation pipeline. It is executed on the host PC to prepare data for RTL or to check RTL outputs.
\paragraph{Line 817}
\noindent\colorbox{gray!8}{\parbox{\dimexpr\linewidth-2\fboxsep\relax}{\raggedright\ttfamily\scriptsize             rf"\textbackslash{}item Vectorless Vivado power: \{vivado['total\_power\_w']:.3f\} W total, \{vivado['dynamic\_power\_w']:.3f\} W dynamic, \{vivado['static\_power\_w']:.3f\} W static.",}}
\par\smallskip
This Python line participates in the export, validation, or report-generation pipeline. It is executed on the host PC to prepare data for RTL or to check RTL outputs.
\paragraph{Line 818}
\noindent\colorbox{gray!8}{\parbox{\dimexpr\linewidth-2\fboxsep\relax}{\raggedright\ttfamily\scriptsize             r"\textbackslash{}end\{itemize\}",}}
\par\smallskip
This Python line participates in the export, validation, or report-generation pipeline. It is executed on the host PC to prepare data for RTL or to check RTL outputs.
\paragraph{Line 819}
\noindent\colorbox{gray!8}{\parbox{\dimexpr\linewidth-2\fboxsep\relax}{\raggedright\ttfamily\scriptsize             r"\textbackslash{}section\{Research Context And Comparison\}",}}
\par\smallskip
This Python line participates in the export, validation, or report-generation pipeline. It is executed on the host PC to prepare data for RTL or to check RTL outputs.
\paragraph{Line 820}
\noindent\colorbox{gray!8}{\parbox{\dimexpr\linewidth-2\fboxsep\relax}{\raggedright\ttfamily\scriptsize             rf"The model side follows the Spikingformer family of spike-driven transformer SNNs, whose reported CIFAR-10 accuracy is about 95.81--95.95\textbackslash{}\% in the cited paper/repository. The local fixed-point exported checkpoint reaches \{acc\_pct:.2f\}\textbackslash{}\% after mapping into this routed hardware head with the measured feature gate. Recent SNN-transformer targets to beat include Max-Former at 97.04\textbackslash{}\% on CIFAR-10 and alpha-XNOR spiking self-attention at 96.15\textbackslash{}\% for the reported multi-stage ViT setting. The implemented hardware contribution here is the verified FPGA path: local feature and weight memories, exact fixed-point RTL/golden agreement, activity-gated feature silence, and on-chip classifier learning.",}}
\par\smallskip
This Python line participates in the export, validation, or report-generation pipeline. It is executed on the host PC to prepare data for RTL or to check RTL outputs.
\paragraph{Line 821}
\noindent\colorbox{gray!8}{\parbox{\dimexpr\linewidth-2\fboxsep\relax}{\raggedright\ttfamily\scriptsize             r"\textbackslash{}begin\{itemize\}",}}
\par\smallskip
This Python line participates in the export, validation, or report-generation pipeline. It is executed on the host PC to prepare data for RTL or to check RTL outputs.
\paragraph{Line 822}
\noindent\colorbox{gray!8}{\parbox{\dimexpr\linewidth-2\fboxsep\relax}{\raggedright\ttfamily\scriptsize             r"\textbackslash{}item Spikingformer paper/repository: \textbackslash{}url\{https://arxiv.org/html/2304.11954v5\}, \textbackslash{}url\{https://github.com/TheBrainLab/Spikingformer\}.",}}
\par\smallskip
This Python line participates in the export, validation, or report-generation pipeline. It is executed on the host PC to prepare data for RTL or to check RTL outputs.
\paragraph{Line 823}
\noindent\colorbox{gray!8}{\parbox{\dimexpr\linewidth-2\fboxsep\relax}{\raggedright\ttfamily\scriptsize             r"\textbackslash{}item Max-Former high-frequency spiking transformer reference: \textbackslash{}url\{https://arxiv.org/html/2505.18608v4\}.",}}
\par\smallskip
This Python line participates in the export, validation, or report-generation pipeline. It is executed on the host PC to prepare data for RTL or to check RTL outputs.
\paragraph{Line 824}
\noindent\colorbox{gray!8}{\parbox{\dimexpr\linewidth-2\fboxsep\relax}{\raggedright\ttfamily\scriptsize             r"\textbackslash{}item Alpha-XNOR spiking self-attention reference: \textbackslash{}url\{https://openaccess.thecvf.com/content/CVPR2025/papers/Xiao\_Rethinking\_Spiking\_Self-Attention\_Mechanism\_Implementing\_a-XNOR\_Similarity\_Calculation\_in\_Spiking\_CVPR\_2025\_paper.pdf\}.",}}
\par\smallskip
This Python line participates in the export, validation, or report-generation pipeline. It is executed on the host PC to prepare data for RTL or to check RTL outputs.
\paragraph{Line 825}
\noindent\colorbox{gray!8}{\parbox{\dimexpr\linewidth-2\fboxsep\relax}{\raggedright\ttfamily\scriptsize             r"\textbackslash{}item TrueNorth reference: \textbackslash{}url\{https://research.ibm.com/publications/truenorth-design-and-tool-flow-of-a-65-mw-1-million-neuron-programmable-neurosynaptic-chip\}.",}}
\par\smallskip
This Python line participates in the export, validation, or report-generation pipeline. It is executed on the host PC to prepare data for RTL or to check RTL outputs.
\paragraph{Line 826}
\noindent\colorbox{gray!8}{\parbox{\dimexpr\linewidth-2\fboxsep\relax}{\raggedright\ttfamily\scriptsize             r"\textbackslash{}item Loihi on-chip learning reference: \textbackslash{}url\{https://redwood.berkeley.edu/wp-content/uploads/2021/08/Davies2018.pdf\}.",}}
\par\smallskip
This Python line participates in the export, validation, or report-generation pipeline. It is executed on the host PC to prepare data for RTL or to check RTL outputs.
\paragraph{Line 827}
\noindent\colorbox{gray!8}{\parbox{\dimexpr\linewidth-2\fboxsep\relax}{\raggedright\ttfamily\scriptsize             r"\textbackslash{}end\{itemize\}",}}
\par\smallskip
This Python line participates in the export, validation, or report-generation pipeline. It is executed on the host PC to prepare data for RTL or to check RTL outputs.
\paragraph{Line 828}
\noindent\colorbox{gray!8}{\parbox{\dimexpr\linewidth-2\fboxsep\relax}{\raggedright\ttfamily\scriptsize             r"\textbackslash{}section\{End-to-End Dataflow\}",}}
\par\smallskip
This Python line participates in the export, validation, or report-generation pipeline. It is executed on the host PC to prepare data for RTL or to check RTL outputs.
\paragraph{Line 829}
\noindent\colorbox{gray!8}{\parbox{\dimexpr\linewidth-2\fboxsep\relax}{\raggedright\ttfamily\scriptsize             r"\textbackslash{}begin\{enumerate\}",}}
\par\smallskip
This Python line participates in the export, validation, or report-generation pipeline. It is executed on the host PC to prepare data for RTL or to check RTL outputs.
\paragraph{Line 830}
\noindent\colorbox{gray!8}{\parbox{\dimexpr\linewidth-2\fboxsep\relax}{\raggedright\ttfamily\scriptsize             r"\textbackslash{}item Python loads the trained Spikingformer checkpoint and CIFAR-10 data, runs the trained network, and exports fixed-point feature, label, weight, and bias memories.",}}
\par\smallskip
This line handles the classifier parameters. Those arrays mirror the RTL memories, so mismatches reveal an export, quantization, or hardware-update bug.
\paragraph{Line 831}
\noindent\colorbox{gray!8}{\parbox{\dimexpr\linewidth-2\fboxsep\relax}{\raggedright\ttfamily\scriptsize             r"\textbackslash{}item XSim streams each 384-word feature vector into the RTL through the same handshake the PYNQ AXI-stream wrapper uses.",}}
\par\smallskip
This Python line participates in the export, validation, or report-generation pipeline. It is executed on the host PC to prepare data for RTL or to check RTL outputs.
\paragraph{Line 832}
\noindent\colorbox{gray!8}{\parbox{\dimexpr\linewidth-2\fboxsep\relax}{\raggedright\ttfamily\scriptsize             r"\textbackslash{}item The RTL stores the feature vector locally, zeroes features whose magnitude is below the selected gate threshold, loops over the ten classes, reads each class row from BRAM, multiplies effective feature and weight in a DSP, accumulates a 64-bit score, and chooses argmax.",}}
\par\smallskip
This Python line participates in the export, validation, or report-generation pipeline. It is executed on the host PC to prepare data for RTL or to check RTL outputs.
\paragraph{Line 833}
\noindent\colorbox{gray!8}{\parbox{\dimexpr\linewidth-2\fboxsep\relax}{\raggedright\ttfamily\scriptsize             r"\textbackslash{}item If training is enabled and the prediction is wrong but close enough by margin, the RTL latches the label, reads the true and predicted class rows, computes signed shifted deltas from the same effective feature path, saturates int16 weight updates, writes updated weights, and updates the two class biases.",}}
\par\smallskip
This line handles the classifier parameters. Those arrays mirror the RTL memories, so mismatches reveal an export, quantization, or hardware-update bug.
\paragraph{Line 834}
\noindent\colorbox{gray!8}{\parbox{\dimexpr\linewidth-2\fboxsep\relax}{\raggedright\ttfamily\scriptsize             r"\textbackslash{}item Python validators independently replay the same fixed-point inference and online-learning rule from the exported memories, then compare every RTL score, prediction, update flag, and update count.",}}
\par\smallskip
This line handles predicted labels. It checks that argmax behavior in Python and RTL chooses the same class.
\paragraph{Line 835}
\noindent\colorbox{gray!8}{\parbox{\dimexpr\linewidth-2\fboxsep\relax}{\raggedright\ttfamily\scriptsize             r"\textbackslash{}item Vivado synthesizes and routes the PYNQ AXI-stream top, producing timing, utilization, power, and checkpoint artifacts.",}}
\par\smallskip
This Python line participates in the export, validation, or report-generation pipeline. It is executed on the host PC to prepare data for RTL or to check RTL outputs.
\paragraph{Line 836}
\noindent\colorbox{gray!8}{\parbox{\dimexpr\linewidth-2\fboxsep\relax}{\raggedright\ttfamily\scriptsize             r"\textbackslash{}end\{enumerate\}",}}
\par\smallskip
This Python line participates in the export, validation, or report-generation pipeline. It is executed on the host PC to prepare data for RTL or to check RTL outputs.
\paragraph{Line 837}
\noindent\colorbox{gray!8}{\parbox{\dimexpr\linewidth-2\fboxsep\relax}{\raggedright\ttfamily\scriptsize             r"\textbackslash{}section\{Line-by-Line Source Explanation\}",}}
\par\smallskip
This Python line participates in the export, validation, or report-generation pipeline. It is executed on the host PC to prepare data for RTL or to check RTL outputs.
\paragraph{Line 838}
\noindent\colorbox{gray!8}{\parbox{\dimexpr\linewidth-2\fboxsep\relax}{\raggedright\ttfamily\scriptsize         ]}}
\par\smallskip
This Python line participates in the export, validation, or report-generation pipeline. It is executed on the host PC to prepare data for RTL or to check RTL outputs.
\paragraph{Line 839}
\noindent\colorbox{gray!8}{\parbox{\dimexpr\linewidth-2\fboxsep\relax}{\raggedright\ttfamily\scriptsize     )}}
\par\smallskip
This Python line participates in the export, validation, or report-generation pipeline. It is executed on the host PC to prepare data for RTL or to check RTL outputs.
\paragraph{Line 840}
\noindent\colorbox{gray!8}{\parbox{\dimexpr\linewidth-2\fboxsep\relax}{\raggedright\ttfamily\scriptsize     for path in files:}}
\par\smallskip
This Python line participates in the export, validation, or report-generation pipeline. It is executed on the host PC to prepare data for RTL or to check RTL outputs.
\paragraph{Line 841}
\noindent\colorbox{gray!8}{\parbox{\dimexpr\linewidth-2\fboxsep\relax}{\raggedright\ttfamily\scriptsize         rel = path.relative\_to(repo\_root())}}
\par\smallskip
This Python line participates in the export, validation, or report-generation pipeline. It is executed on the host PC to prepare data for RTL or to check RTL outputs.
\paragraph{Line 842}
\noindent\colorbox{gray!8}{\parbox{\dimexpr\linewidth-2\fboxsep\relax}{\raggedright\ttfamily\scriptsize         lines.append(r"\textbackslash{}clearpage")}}
\par\smallskip
This Python line participates in the export, validation, or report-generation pipeline. It is executed on the host PC to prepare data for RTL or to check RTL outputs.
\paragraph{Line 843}
\noindent\colorbox{gray!8}{\parbox{\dimexpr\linewidth-2\fboxsep\relax}{\raggedright\ttfamily\scriptsize         lines.append(r"\textbackslash{}subsection\{" + latex\_escape(str(rel)) + "\}")}}
\par\smallskip
This Python line participates in the export, validation, or report-generation pipeline. It is executed on the host PC to prepare data for RTL or to check RTL outputs.
\paragraph{Line 844}
\noindent\colorbox{gray!8}{\parbox{\dimexpr\linewidth-2\fboxsep\relax}{\raggedright\ttfamily\scriptsize         lines.append("Each entry below gives the exact source line followed by its hardware/software role in the implemented accelerator pipeline.")}}
\par\smallskip
This Python line participates in the export, validation, or report-generation pipeline. It is executed on the host PC to prepare data for RTL or to check RTL outputs.
\paragraph{Line 845}
\noindent\colorbox{gray!8}{\parbox{\dimexpr\linewidth-2\fboxsep\relax}{\raggedright\ttfamily\scriptsize         previous = ""}}
\par\smallskip
This Python line participates in the export, validation, or report-generation pipeline. It is executed on the host PC to prepare data for RTL or to check RTL outputs.
\paragraph{Line 846}
\noindent\colorbox{gray!8}{\parbox{\dimexpr\linewidth-2\fboxsep\relax}{\raggedright\ttfamily\scriptsize         for i, raw in enumerate(path.read\_text(encoding="utf-8", errors="ignore").splitlines(), start=1):}}
\par\smallskip
This reads an artifact produced by training, export, simulation, or Vivado. The validator/report depends on files generated by the actual run rather than hidden constants.
\paragraph{Line 847}
\noindent\colorbox{gray!8}{\parbox{\dimexpr\linewidth-2\fboxsep\relax}{\raggedright\ttfamily\scriptsize             explanation = explain\_line(path, raw, i, previous)}}
\par\smallskip
This Python line participates in the export, validation, or report-generation pipeline. It is executed on the host PC to prepare data for RTL or to check RTL outputs.
\paragraph{Line 848}
\noindent\colorbox{gray!8}{\parbox{\dimexpr\linewidth-2\fboxsep\relax}{\raggedright\ttfamily\scriptsize             if raw.strip():}}
\par\smallskip
This Python line participates in the export, validation, or report-generation pipeline. It is executed on the host PC to prepare data for RTL or to check RTL outputs.
\paragraph{Line 849}
\noindent\colorbox{gray!8}{\parbox{\dimexpr\linewidth-2\fboxsep\relax}{\raggedright\ttfamily\scriptsize                 previous = raw.strip()}}
\par\smallskip
This Python line participates in the export, validation, or report-generation pipeline. It is executed on the host PC to prepare data for RTL or to check RTL outputs.
\paragraph{Line 850}
\noindent\colorbox{gray!8}{\parbox{\dimexpr\linewidth-2\fboxsep\relax}{\raggedright\ttfamily\scriptsize             lines.append(r"\textbackslash{}paragraph\{Line " + str(i) + "\}")}}
\par\smallskip
This Python line participates in the export, validation, or report-generation pipeline. It is executed on the host PC to prepare data for RTL or to check RTL outputs.
\paragraph{Line 851}
\noindent\colorbox{gray!8}{\parbox{\dimexpr\linewidth-2\fboxsep\relax}{\raggedright\ttfamily\scriptsize             code\_text = latex\_escape(raw) if raw else r"\textbackslash{}emph\{blank line\}"}}
\par\smallskip
This Python line participates in the export, validation, or report-generation pipeline. It is executed on the host PC to prepare data for RTL or to check RTL outputs.
\paragraph{Line 852}
\noindent\colorbox{gray!8}{\parbox{\dimexpr\linewidth-2\fboxsep\relax}{\raggedright\ttfamily\scriptsize             lines.append(}}
\par\smallskip
This Python line participates in the export, validation, or report-generation pipeline. It is executed on the host PC to prepare data for RTL or to check RTL outputs.
\paragraph{Line 853}
\noindent\colorbox{gray!8}{\parbox{\dimexpr\linewidth-2\fboxsep\relax}{\raggedright\ttfamily\scriptsize                 r"\textbackslash{}noindent\textbackslash{}colorbox\{gray!8\}\{\textbackslash{}parbox\{\textbackslash{}dimexpr\textbackslash{}linewidth-2\textbackslash{}fboxsep\textbackslash{}relax\}\{\textbackslash{}raggedright\textbackslash{}ttfamily\textbackslash{}scriptsize "}}
\par\smallskip
This Python line participates in the export, validation, or report-generation pipeline. It is executed on the host PC to prepare data for RTL or to check RTL outputs.
\paragraph{Line 854}
\noindent\colorbox{gray!8}{\parbox{\dimexpr\linewidth-2\fboxsep\relax}{\raggedright\ttfamily\scriptsize                 + code\_text}}
\par\smallskip
This Python line participates in the export, validation, or report-generation pipeline. It is executed on the host PC to prepare data for RTL or to check RTL outputs.
\paragraph{Line 855}
\noindent\colorbox{gray!8}{\parbox{\dimexpr\linewidth-2\fboxsep\relax}{\raggedright\ttfamily\scriptsize                 + r"\}\}"}}
\par\smallskip
This Python line participates in the export, validation, or report-generation pipeline. It is executed on the host PC to prepare data for RTL or to check RTL outputs.
\paragraph{Line 856}
\noindent\colorbox{gray!8}{\parbox{\dimexpr\linewidth-2\fboxsep\relax}{\raggedright\ttfamily\scriptsize             )}}
\par\smallskip
This Python line participates in the export, validation, or report-generation pipeline. It is executed on the host PC to prepare data for RTL or to check RTL outputs.
\paragraph{Line 857}
\noindent\colorbox{gray!8}{\parbox{\dimexpr\linewidth-2\fboxsep\relax}{\raggedright\ttfamily\scriptsize             lines.append(r"\textbackslash{}par\textbackslash{}smallskip")}}
\par\smallskip
This Python line participates in the export, validation, or report-generation pipeline. It is executed on the host PC to prepare data for RTL or to check RTL outputs.
\paragraph{Line 858}
\noindent\colorbox{gray!8}{\parbox{\dimexpr\linewidth-2\fboxsep\relax}{\raggedright\ttfamily\scriptsize             lines.append(latex\_escape(explanation))}}
\par\smallskip
This Python line participates in the export, validation, or report-generation pipeline. It is executed on the host PC to prepare data for RTL or to check RTL outputs.
\paragraph{Line 859}
\noindent\colorbox{gray!8}{\parbox{\dimexpr\linewidth-2\fboxsep\relax}{\raggedright\ttfamily\scriptsize     lines.append(r"\textbackslash{}end\{document\}")}}
\par\smallskip
This Python line participates in the export, validation, or report-generation pipeline. It is executed on the host PC to prepare data for RTL or to check RTL outputs.
\paragraph{Line 860}
\noindent\colorbox{gray!8}{\parbox{\dimexpr\linewidth-2\fboxsep\relax}{\raggedright\ttfamily\scriptsize     tex\_path.write\_text("\textbackslash{}n".join(lines), encoding="utf-8")}}
\par\smallskip
This writes machine-readable output for later reports. Keeping JSON/CSV artifacts makes the verification auditable.
\paragraph{Line 861}
\noindent\colorbox{gray!8}{\parbox{\dimexpr\linewidth-2\fboxsep\relax}{\raggedright\ttfamily\scriptsize     return tex\_path}}
\par\smallskip
This Python line participates in the export, validation, or report-generation pipeline. It is executed on the host PC to prepare data for RTL or to check RTL outputs.
\paragraph{Line 862}
\noindent\colorbox{gray!8}{\parbox{\dimexpr\linewidth-2\fboxsep\relax}{\raggedright\ttfamily\scriptsize \emph{blank line}}}
\par\smallskip
This blank line is not hardware or software behavior. It separates nearby declarations or operations so the file remains readable and reviewable.
\paragraph{Line 863}
\noindent\colorbox{gray!8}{\parbox{\dimexpr\linewidth-2\fboxsep\relax}{\raggedright\ttfamily\scriptsize \emph{blank line}}}
\par\smallskip
This blank line is not hardware or software behavior. It separates nearby declarations or operations so the file remains readable and reviewable.
\paragraph{Line 864}
\noindent\colorbox{gray!8}{\parbox{\dimexpr\linewidth-2\fboxsep\relax}{\raggedright\ttfamily\scriptsize def compile\_tex(tex\_path: Path) -> None:}}
\par\smallskip
This defines a Python function. The function packages one repeatable step such as loading a memory file, simulating the golden learning rule, validating CSVs, or generating plots.
\paragraph{Line 865}
\noindent\colorbox{gray!8}{\parbox{\dimexpr\linewidth-2\fboxsep\relax}{\raggedright\ttfamily\scriptsize     pdflatex = shutil.which("pdflatex")}}
\par\smallskip
This Python line participates in the export, validation, or report-generation pipeline. It is executed on the host PC to prepare data for RTL or to check RTL outputs.
\paragraph{Line 866}
\noindent\colorbox{gray!8}{\parbox{\dimexpr\linewidth-2\fboxsep\relax}{\raggedright\ttfamily\scriptsize     if not pdflatex:}}
\par\smallskip
This Python line participates in the export, validation, or report-generation pipeline. It is executed on the host PC to prepare data for RTL or to check RTL outputs.
\paragraph{Line 867}
\noindent\colorbox{gray!8}{\parbox{\dimexpr\linewidth-2\fboxsep\relax}{\raggedright\ttfamily\scriptsize         print("pdflatex not found; wrote TeX only.")}}
\par\smallskip
This Python line participates in the export, validation, or report-generation pipeline. It is executed on the host PC to prepare data for RTL or to check RTL outputs.
\paragraph{Line 868}
\noindent\colorbox{gray!8}{\parbox{\dimexpr\linewidth-2\fboxsep\relax}{\raggedright\ttfamily\scriptsize         return}}
\par\smallskip
This Python line participates in the export, validation, or report-generation pipeline. It is executed on the host PC to prepare data for RTL or to check RTL outputs.
\paragraph{Line 869}
\noindent\colorbox{gray!8}{\parbox{\dimexpr\linewidth-2\fboxsep\relax}{\raggedright\ttfamily\scriptsize     for \_ in range(2):}}
\par\smallskip
This Python line participates in the export, validation, or report-generation pipeline. It is executed on the host PC to prepare data for RTL or to check RTL outputs.
\paragraph{Line 870}
\noindent\colorbox{gray!8}{\parbox{\dimexpr\linewidth-2\fboxsep\relax}{\raggedright\ttfamily\scriptsize         subprocess.run(}}
\par\smallskip
This Python line participates in the export, validation, or report-generation pipeline. It is executed on the host PC to prepare data for RTL or to check RTL outputs.
\paragraph{Line 871}
\noindent\colorbox{gray!8}{\parbox{\dimexpr\linewidth-2\fboxsep\relax}{\raggedright\ttfamily\scriptsize             [pdflatex, "-interaction=nonstopmode", "-halt-on-error", tex\_path.name],}}
\par\smallskip
This Python line participates in the export, validation, or report-generation pipeline. It is executed on the host PC to prepare data for RTL or to check RTL outputs.
\paragraph{Line 872}
\noindent\colorbox{gray!8}{\parbox{\dimexpr\linewidth-2\fboxsep\relax}{\raggedright\ttfamily\scriptsize             cwd=tex\_path.parent,}}
\par\smallskip
This Python line participates in the export, validation, or report-generation pipeline. It is executed on the host PC to prepare data for RTL or to check RTL outputs.
\paragraph{Line 873}
\noindent\colorbox{gray!8}{\parbox{\dimexpr\linewidth-2\fboxsep\relax}{\raggedright\ttfamily\scriptsize             check=True,}}
\par\smallskip
This Python line participates in the export, validation, or report-generation pipeline. It is executed on the host PC to prepare data for RTL or to check RTL outputs.
\paragraph{Line 874}
\noindent\colorbox{gray!8}{\parbox{\dimexpr\linewidth-2\fboxsep\relax}{\raggedright\ttfamily\scriptsize             stdout=subprocess.PIPE,}}
\par\smallskip
This Python line participates in the export, validation, or report-generation pipeline. It is executed on the host PC to prepare data for RTL or to check RTL outputs.
\paragraph{Line 875}
\noindent\colorbox{gray!8}{\parbox{\dimexpr\linewidth-2\fboxsep\relax}{\raggedright\ttfamily\scriptsize             stderr=subprocess.STDOUT,}}
\par\smallskip
This Python line participates in the export, validation, or report-generation pipeline. It is executed on the host PC to prepare data for RTL or to check RTL outputs.
\paragraph{Line 876}
\noindent\colorbox{gray!8}{\parbox{\dimexpr\linewidth-2\fboxsep\relax}{\raggedright\ttfamily\scriptsize             text=True,}}
\par\smallskip
This Python line participates in the export, validation, or report-generation pipeline. It is executed on the host PC to prepare data for RTL or to check RTL outputs.
\paragraph{Line 877}
\noindent\colorbox{gray!8}{\parbox{\dimexpr\linewidth-2\fboxsep\relax}{\raggedright\ttfamily\scriptsize         )}}
\par\smallskip
This Python line participates in the export, validation, or report-generation pipeline. It is executed on the host PC to prepare data for RTL or to check RTL outputs.
\paragraph{Line 878}
\noindent\colorbox{gray!8}{\parbox{\dimexpr\linewidth-2\fboxsep\relax}{\raggedright\ttfamily\scriptsize     print(f"Wrote \{tex\_path.with\_suffix('.pdf')\}")}}
\par\smallskip
This Python line participates in the export, validation, or report-generation pipeline. It is executed on the host PC to prepare data for RTL or to check RTL outputs.
\paragraph{Line 879}
\noindent\colorbox{gray!8}{\parbox{\dimexpr\linewidth-2\fboxsep\relax}{\raggedright\ttfamily\scriptsize \emph{blank line}}}
\par\smallskip
This blank line is not hardware or software behavior. It separates nearby declarations or operations so the file remains readable and reviewable.
\paragraph{Line 880}
\noindent\colorbox{gray!8}{\parbox{\dimexpr\linewidth-2\fboxsep\relax}{\raggedright\ttfamily\scriptsize \emph{blank line}}}
\par\smallskip
This blank line is not hardware or software behavior. It separates nearby declarations or operations so the file remains readable and reviewable.
\paragraph{Line 881}
\noindent\colorbox{gray!8}{\parbox{\dimexpr\linewidth-2\fboxsep\relax}{\raggedright\ttfamily\scriptsize def main() -> None:}}
\par\smallskip
This defines a Python function. The function packages one repeatable step such as loading a memory file, simulating the golden learning rule, validating CSVs, or generating plots.
\paragraph{Line 882}
\noindent\colorbox{gray!8}{\parbox{\dimexpr\linewidth-2\fboxsep\relax}{\raggedright\ttfamily\scriptsize     generate\_debug\_pdf()}}
\par\smallskip
This Python line participates in the export, validation, or report-generation pipeline. It is executed on the host PC to prepare data for RTL or to check RTL outputs.
\paragraph{Line 883}
\noindent\colorbox{gray!8}{\parbox{\dimexpr\linewidth-2\fboxsep\relax}{\raggedright\ttfamily\scriptsize     tex\_path = generate\_explainer\_tex()}}
\par\smallskip
This Python line participates in the export, validation, or report-generation pipeline. It is executed on the host PC to prepare data for RTL or to check RTL outputs.
\paragraph{Line 884}
\noindent\colorbox{gray!8}{\parbox{\dimexpr\linewidth-2\fboxsep\relax}{\raggedright\ttfamily\scriptsize     print(f"Wrote \{tex\_path\}")}}
\par\smallskip
This Python line participates in the export, validation, or report-generation pipeline. It is executed on the host PC to prepare data for RTL or to check RTL outputs.
\paragraph{Line 885}
\noindent\colorbox{gray!8}{\parbox{\dimexpr\linewidth-2\fboxsep\relax}{\raggedright\ttfamily\scriptsize     compile\_tex(tex\_path)}}
\par\smallskip
This Python line participates in the export, validation, or report-generation pipeline. It is executed on the host PC to prepare data for RTL or to check RTL outputs.
\paragraph{Line 886}
\noindent\colorbox{gray!8}{\parbox{\dimexpr\linewidth-2\fboxsep\relax}{\raggedright\ttfamily\scriptsize \emph{blank line}}}
\par\smallskip
This blank line is not hardware or software behavior. It separates nearby declarations or operations so the file remains readable and reviewable.
\paragraph{Line 887}
\noindent\colorbox{gray!8}{\parbox{\dimexpr\linewidth-2\fboxsep\relax}{\raggedright\ttfamily\scriptsize \emph{blank line}}}
\par\smallskip
This blank line is not hardware or software behavior. It separates nearby declarations or operations so the file remains readable and reviewable.
\paragraph{Line 888}
\noindent\colorbox{gray!8}{\parbox{\dimexpr\linewidth-2\fboxsep\relax}{\raggedright\ttfamily\scriptsize if \_\_name\_\_ == "\_\_main\_\_":}}
\par\smallskip
This Python idiom runs main() only when the file is executed as a script. It still allows the file's functions to be imported for tests or reuse.
\paragraph{Line 889}
\noindent\colorbox{gray!8}{\parbox{\dimexpr\linewidth-2\fboxsep\relax}{\raggedright\ttfamily\scriptsize     main()}}
\par\smallskip
This Python line participates in the export, validation, or report-generation pipeline. It is executed on the host PC to prepare data for RTL or to check RTL outputs.
\clearpage
\subsection{cifar10\_neuromorphic\textbackslash{}run\_spikingformer\_head\_accelerator.ps1}
Each entry below gives the exact source line followed by its hardware/software role in the implemented accelerator pipeline.
\paragraph{Line 1}
\noindent\colorbox{gray!8}{\parbox{\dimexpr\linewidth-2\fboxsep\relax}{\raggedright\ttfamily\scriptsize param(}}
\par\smallskip
This begins the PowerShell parameter block. It exposes repeatable knobs such as sample count, Vivado path, board part, and skip switches.
\paragraph{Line 2}
\noindent\colorbox{gray!8}{\parbox{\dimexpr\linewidth-2\fboxsep\relax}{\raggedright\ttfamily\scriptsize     [int]\$Samples = 10000,}}
\par\smallskip
This declares one typed PowerShell parameter. Strong typing catches accidental string/numeric mistakes before the long training or Vivado flow starts.
\paragraph{Line 3}
\noindent\colorbox{gray!8}{\parbox{\dimexpr\linewidth-2\fboxsep\relax}{\raggedright\ttfamily\scriptsize     [int]\$BatchSize = 64,}}
\par\smallskip
This declares one typed PowerShell parameter. Strong typing catches accidental string/numeric mistakes before the long training or Vivado flow starts.
\paragraph{Line 4}
\noindent\colorbox{gray!8}{\parbox{\dimexpr\linewidth-2\fboxsep\relax}{\raggedright\ttfamily\scriptsize     [int]\$Workers = 4,}}
\par\smallskip
This declares one typed PowerShell parameter. Strong typing catches accidental string/numeric mistakes before the long training or Vivado flow starts.
\paragraph{Line 5}
\noindent\colorbox{gray!8}{\parbox{\dimexpr\linewidth-2\fboxsep\relax}{\raggedright\ttfamily\scriptsize     [int]\$TraceSamples = 8,}}
\par\smallskip
This declares one typed PowerShell parameter. Strong typing catches accidental string/numeric mistakes before the long training or Vivado flow starts.
\paragraph{Line 6}
\noindent\colorbox{gray!8}{\parbox{\dimexpr\linewidth-2\fboxsep\relax}{\raggedright\ttfamily\scriptsize     [int]\$InferSkipAbsThresh = 2048,}}
\par\smallskip
This declares one typed PowerShell parameter. Strong typing catches accidental string/numeric mistakes before the long training or Vivado flow starts.
\paragraph{Line 7}
\noindent\colorbox{gray!8}{\parbox{\dimexpr\linewidth-2\fboxsep\relax}{\raggedright\ttfamily\scriptsize     [string]\$Checkpoint = "cifar10\_neuromorphic\textbackslash{}outputs\_upstream\_spikingformer\textbackslash{}spikingformer\_4\_384\_400e\textbackslash{}model\_best.pth.tar",}}
\par\smallskip
This declares one typed PowerShell parameter. Strong typing catches accidental string/numeric mistakes before the long training or Vivado flow starts.
\paragraph{Line 8}
\noindent\colorbox{gray!8}{\parbox{\dimexpr\linewidth-2\fboxsep\relax}{\raggedright\ttfamily\scriptsize     [string]\$OutDir = "cifar10\_neuromorphic\textbackslash{}outputs\_spikingformer\_head\_accelerator",}}
\par\smallskip
This declares one typed PowerShell parameter. Strong typing catches accidental string/numeric mistakes before the long training or Vivado flow starts.
\paragraph{Line 9}
\noindent\colorbox{gray!8}{\parbox{\dimexpr\linewidth-2\fboxsep\relax}{\raggedright\ttfamily\scriptsize     [string]\$VivadoBin = "D:\textbackslash{}AMDDesignTools\textbackslash{}2026.1\textbackslash{}Vivado\textbackslash{}bin",}}
\par\smallskip
This declares one typed PowerShell parameter. Strong typing catches accidental string/numeric mistakes before the long training or Vivado flow starts.
\paragraph{Line 10}
\noindent\colorbox{gray!8}{\parbox{\dimexpr\linewidth-2\fboxsep\relax}{\raggedright\ttfamily\scriptsize     [string]\$Part = "xc7z020clg400-1",}}
\par\smallskip
This declares one typed PowerShell parameter. Strong typing catches accidental string/numeric mistakes before the long training or Vivado flow starts.
\paragraph{Line 11}
\noindent\colorbox{gray!8}{\parbox{\dimexpr\linewidth-2\fboxsep\relax}{\raggedright\ttfamily\scriptsize     [double]\$ClockPeriodNs = 10.0,}}
\par\smallskip
This declares one typed PowerShell parameter. Strong typing catches accidental string/numeric mistakes before the long training or Vivado flow starts.
\paragraph{Line 12}
\noindent\colorbox{gray!8}{\parbox{\dimexpr\linewidth-2\fboxsep\relax}{\raggedright\ttfamily\scriptsize     [switch]\$SkipExport,}}
\par\smallskip
This declares one typed PowerShell parameter. Strong typing catches accidental string/numeric mistakes before the long training or Vivado flow starts.
\paragraph{Line 13}
\noindent\colorbox{gray!8}{\parbox{\dimexpr\linewidth-2\fboxsep\relax}{\raggedright\ttfamily\scriptsize     [switch]\$SkipVivado,}}
\par\smallskip
This declares one typed PowerShell parameter. Strong typing catches accidental string/numeric mistakes before the long training or Vivado flow starts.
\paragraph{Line 14}
\noindent\colorbox{gray!8}{\parbox{\dimexpr\linewidth-2\fboxsep\relax}{\raggedright\ttfamily\scriptsize     [switch]\$SkipOnChipLearning,}}
\par\smallskip
This declares one typed PowerShell parameter. Strong typing catches accidental string/numeric mistakes before the long training or Vivado flow starts.
\paragraph{Line 15}
\noindent\colorbox{gray!8}{\parbox{\dimexpr\linewidth-2\fboxsep\relax}{\raggedright\ttfamily\scriptsize     [switch]\$SkipSynthesis}}
\par\smallskip
This declares one typed PowerShell parameter. Strong typing catches accidental string/numeric mistakes before the long training or Vivado flow starts.
\paragraph{Line 16}
\noindent\colorbox{gray!8}{\parbox{\dimexpr\linewidth-2\fboxsep\relax}{\raggedright\ttfamily\scriptsize )}}
\par\smallskip
This PowerShell line orchestrates the reproducible full flow. It exists so changes can be rerun without manually retyping many commands.
\paragraph{Line 17}
\noindent\colorbox{gray!8}{\parbox{\dimexpr\linewidth-2\fboxsep\relax}{\raggedright\ttfamily\scriptsize \emph{blank line}}}
\par\smallskip
This blank line is not hardware or software behavior. It separates nearby declarations or operations so the file remains readable and reviewable.
\paragraph{Line 18}
\noindent\colorbox{gray!8}{\parbox{\dimexpr\linewidth-2\fboxsep\relax}{\raggedright\ttfamily\scriptsize \$ErrorActionPreference = "Stop"}}
\par\smallskip
This makes failures stop the script instead of continuing after a broken export, simulation, or synthesis step.
\paragraph{Line 19}
\noindent\colorbox{gray!8}{\parbox{\dimexpr\linewidth-2\fboxsep\relax}{\raggedright\ttfamily\scriptsize \$ScriptBoundParameters = @\{\} + \$PSBoundParameters}}
\par\smallskip
This PowerShell line orchestrates the reproducible full flow. It exists so changes can be rerun without manually retyping many commands.
\paragraph{Line 20}
\noindent\colorbox{gray!8}{\parbox{\dimexpr\linewidth-2\fboxsep\relax}{\raggedright\ttfamily\scriptsize \emph{blank line}}}
\par\smallskip
This blank line is not hardware or software behavior. It separates nearby declarations or operations so the file remains readable and reviewable.
\paragraph{Line 21}
\noindent\colorbox{gray!8}{\parbox{\dimexpr\linewidth-2\fboxsep\relax}{\raggedright\ttfamily\scriptsize function Use-PowerShell7 \{}}
\par\smallskip
This PowerShell line orchestrates the reproducible full flow. It exists so changes can be rerun without manually retyping many commands.
\paragraph{Line 22}
\noindent\colorbox{gray!8}{\parbox{\dimexpr\linewidth-2\fboxsep\relax}{\raggedright\ttfamily\scriptsize     param([hashtable]\$BoundParameters)}}
\par\smallskip
This begins the PowerShell parameter block. It exposes repeatable knobs such as sample count, Vivado path, board part, and skip switches.
\paragraph{Line 23}
\noindent\colorbox{gray!8}{\parbox{\dimexpr\linewidth-2\fboxsep\relax}{\raggedright\ttfamily\scriptsize     if (\$PSVersionTable.PSEdition -eq "Core" -and \$PSVersionTable.PSVersion.Major -ge 7) \{ return \}}}
\par\smallskip
This PowerShell line orchestrates the reproducible full flow. It exists so changes can be rerun without manually retyping many commands.
\paragraph{Line 24}
\noindent\colorbox{gray!8}{\parbox{\dimexpr\linewidth-2\fboxsep\relax}{\raggedright\ttfamily\scriptsize     \$pwsh = (Get-Command pwsh -CommandType Application -ErrorAction SilentlyContinue).Source}}
\par\smallskip
This PowerShell line orchestrates the reproducible full flow. It exists so changes can be rerun without manually retyping many commands.
\paragraph{Line 25}
\noindent\colorbox{gray!8}{\parbox{\dimexpr\linewidth-2\fboxsep\relax}{\raggedright\ttfamily\scriptsize     if (!\$pwsh) \{ throw "PowerShell 7+ (pwsh) is required. Install it with: winget install --id Microsoft.PowerShell --source winget" \}}}
\par\smallskip
This PowerShell line orchestrates the reproducible full flow. It exists so changes can be rerun without manually retyping many commands.
\paragraph{Line 26}
\noindent\colorbox{gray!8}{\parbox{\dimexpr\linewidth-2\fboxsep\relax}{\raggedright\ttfamily\scriptsize     \$argList = @("-NoLogo", "-NoProfile", "-ExecutionPolicy", "Bypass", "-File", \$PSCommandPath)}}
\par\smallskip
This PowerShell line orchestrates the reproducible full flow. It exists so changes can be rerun without manually retyping many commands.
\paragraph{Line 27}
\noindent\colorbox{gray!8}{\parbox{\dimexpr\linewidth-2\fboxsep\relax}{\raggedright\ttfamily\scriptsize     foreach (\$key in \$BoundParameters.Keys) \{}}
\par\smallskip
This PowerShell line orchestrates the reproducible full flow. It exists so changes can be rerun without manually retyping many commands.
\paragraph{Line 28}
\noindent\colorbox{gray!8}{\parbox{\dimexpr\linewidth-2\fboxsep\relax}{\raggedright\ttfamily\scriptsize         \$value = \$BoundParameters[\$key]}}
\par\smallskip
This PowerShell line orchestrates the reproducible full flow. It exists so changes can be rerun without manually retyping many commands.
\paragraph{Line 29}
\noindent\colorbox{gray!8}{\parbox{\dimexpr\linewidth-2\fboxsep\relax}{\raggedright\ttfamily\scriptsize         if (\$value -is [System.Management.Automation.SwitchParameter]) \{}}
\par\smallskip
This PowerShell line orchestrates the reproducible full flow. It exists so changes can be rerun without manually retyping many commands.
\paragraph{Line 30}
\noindent\colorbox{gray!8}{\parbox{\dimexpr\linewidth-2\fboxsep\relax}{\raggedright\ttfamily\scriptsize             if (\$value.IsPresent) \{ \$argList += "-\$key" \}}}
\par\smallskip
This PowerShell line orchestrates the reproducible full flow. It exists so changes can be rerun without manually retyping many commands.
\paragraph{Line 31}
\noindent\colorbox{gray!8}{\parbox{\dimexpr\linewidth-2\fboxsep\relax}{\raggedright\ttfamily\scriptsize         \} elseif (\$null -ne \$value) \{}}
\par\smallskip
This PowerShell line orchestrates the reproducible full flow. It exists so changes can be rerun without manually retyping many commands.
\paragraph{Line 32}
\noindent\colorbox{gray!8}{\parbox{\dimexpr\linewidth-2\fboxsep\relax}{\raggedright\ttfamily\scriptsize             \$argList += "-\$key"}}
\par\smallskip
This PowerShell line orchestrates the reproducible full flow. It exists so changes can be rerun without manually retyping many commands.
\paragraph{Line 33}
\noindent\colorbox{gray!8}{\parbox{\dimexpr\linewidth-2\fboxsep\relax}{\raggedright\ttfamily\scriptsize             \$argList += "\$value"}}
\par\smallskip
This PowerShell line orchestrates the reproducible full flow. It exists so changes can be rerun without manually retyping many commands.
\paragraph{Line 34}
\noindent\colorbox{gray!8}{\parbox{\dimexpr\linewidth-2\fboxsep\relax}{\raggedright\ttfamily\scriptsize         \}}}
\par\smallskip
This PowerShell line orchestrates the reproducible full flow. It exists so changes can be rerun without manually retyping many commands.
\paragraph{Line 35}
\noindent\colorbox{gray!8}{\parbox{\dimexpr\linewidth-2\fboxsep\relax}{\raggedright\ttfamily\scriptsize     \}}}
\par\smallskip
This PowerShell line orchestrates the reproducible full flow. It exists so changes can be rerun without manually retyping many commands.
\paragraph{Line 36}
\noindent\colorbox{gray!8}{\parbox{\dimexpr\linewidth-2\fboxsep\relax}{\raggedright\ttfamily\scriptsize     \& \$pwsh @argList}}
\par\smallskip
This PowerShell line orchestrates the reproducible full flow. It exists so changes can be rerun without manually retyping many commands.
\paragraph{Line 37}
\noindent\colorbox{gray!8}{\parbox{\dimexpr\linewidth-2\fboxsep\relax}{\raggedright\ttfamily\scriptsize     exit \$LASTEXITCODE}}
\par\smallskip
This PowerShell line orchestrates the reproducible full flow. It exists so changes can be rerun without manually retyping many commands.
\paragraph{Line 38}
\noindent\colorbox{gray!8}{\parbox{\dimexpr\linewidth-2\fboxsep\relax}{\raggedright\ttfamily\scriptsize \}}}
\par\smallskip
This PowerShell line orchestrates the reproducible full flow. It exists so changes can be rerun without manually retyping many commands.
\paragraph{Line 39}
\noindent\colorbox{gray!8}{\parbox{\dimexpr\linewidth-2\fboxsep\relax}{\raggedright\ttfamily\scriptsize Use-PowerShell7 -BoundParameters \$ScriptBoundParameters}}
\par\smallskip
This PowerShell line orchestrates the reproducible full flow. It exists so changes can be rerun without manually retyping many commands.
\paragraph{Line 40}
\noindent\colorbox{gray!8}{\parbox{\dimexpr\linewidth-2\fboxsep\relax}{\raggedright\ttfamily\scriptsize \emph{blank line}}}
\par\smallskip
This blank line is not hardware or software behavior. It separates nearby declarations or operations so the file remains readable and reviewable.
\paragraph{Line 41}
\noindent\colorbox{gray!8}{\parbox{\dimexpr\linewidth-2\fboxsep\relax}{\raggedright\ttfamily\scriptsize \$ProjectDir = Split-Path -Parent \$MyInvocation.MyCommand.Path}}
\par\smallskip
This PowerShell line orchestrates the reproducible full flow. It exists so changes can be rerun without manually retyping many commands.
\paragraph{Line 42}
\noindent\colorbox{gray!8}{\parbox{\dimexpr\linewidth-2\fboxsep\relax}{\raggedright\ttfamily\scriptsize \$RepoRoot = Resolve-Path (Join-Path \$ProjectDir "..")}}
\par\smallskip
This PowerShell line orchestrates the reproducible full flow. It exists so changes can be rerun without manually retyping many commands.
\paragraph{Line 43}
\noindent\colorbox{gray!8}{\parbox{\dimexpr\linewidth-2\fboxsep\relax}{\raggedright\ttfamily\scriptsize \$Python = Join-Path \$RepoRoot "external\_research\textbackslash{}spikingformer\_py310\textbackslash{}Scripts\textbackslash{}python.exe"}}
\par\smallskip
This line invokes the Python environment used for export, validation, or report generation. It ties host-side math to the hardware artifacts.
\paragraph{Line 44}
\noindent\colorbox{gray!8}{\parbox{\dimexpr\linewidth-2\fboxsep\relax}{\raggedright\ttfamily\scriptsize \$TorchLib = Join-Path \$RepoRoot "external\_research\textbackslash{}spikingformer\_py310\textbackslash{}Lib\textbackslash{}site-packages\textbackslash{}torch\textbackslash{}lib"}}
\par\smallskip
This PowerShell line orchestrates the reproducible full flow. It exists so changes can be rerun without manually retyping many commands.
\paragraph{Line 45}
\noindent\colorbox{gray!8}{\parbox{\dimexpr\linewidth-2\fboxsep\relax}{\raggedright\ttfamily\scriptsize \$VivadoXvlog = Join-Path \$VivadoBin "xvlog.bat"}}
\par\smallskip
This line invokes or locates Vivado simulator/implementation tools. The flow uses the installed AMD toolchain on D: rather than a substitute simulator.
\paragraph{Line 46}
\noindent\colorbox{gray!8}{\parbox{\dimexpr\linewidth-2\fboxsep\relax}{\raggedright\ttfamily\scriptsize \$VivadoXelab = Join-Path \$VivadoBin "xelab.bat"}}
\par\smallskip
This line invokes or locates Vivado simulator/implementation tools. The flow uses the installed AMD toolchain on D: rather than a substitute simulator.
\paragraph{Line 47}
\noindent\colorbox{gray!8}{\parbox{\dimexpr\linewidth-2\fboxsep\relax}{\raggedright\ttfamily\scriptsize \$VivadoXsim = Join-Path \$VivadoBin "xsim.bat"}}
\par\smallskip
This line invokes or locates Vivado simulator/implementation tools. The flow uses the installed AMD toolchain on D: rather than a substitute simulator.
\paragraph{Line 48}
\noindent\colorbox{gray!8}{\parbox{\dimexpr\linewidth-2\fboxsep\relax}{\raggedright\ttfamily\scriptsize \$Vivado = Join-Path \$VivadoBin "vivado.bat"}}
\par\smallskip
This line invokes or locates Vivado simulator/implementation tools. The flow uses the installed AMD toolchain on D: rather than a substitute simulator.
\paragraph{Line 49}
\noindent\colorbox{gray!8}{\parbox{\dimexpr\linewidth-2\fboxsep\relax}{\raggedright\ttfamily\scriptsize \emph{blank line}}}
\par\smallskip
This blank line is not hardware or software behavior. It separates nearby declarations or operations so the file remains readable and reviewable.
\paragraph{Line 50}
\noindent\colorbox{gray!8}{\parbox{\dimexpr\linewidth-2\fboxsep\relax}{\raggedright\ttfamily\scriptsize if (!(Test-Path \$Python)) \{ throw "Missing Python env: \$Python" \}}}
\par\smallskip
This line invokes the Python environment used for export, validation, or report generation. It ties host-side math to the hardware artifacts.
\paragraph{Line 51}
\noindent\colorbox{gray!8}{\parbox{\dimexpr\linewidth-2\fboxsep\relax}{\raggedright\ttfamily\scriptsize if (Test-Path \$TorchLib) \{ \$env:PATH = "\$TorchLib;\$env:PATH" \}}}
\par\smallskip
This PowerShell line orchestrates the reproducible full flow. It exists so changes can be rerun without manually retyping many commands.
\paragraph{Line 52}
\noindent\colorbox{gray!8}{\parbox{\dimexpr\linewidth-2\fboxsep\relax}{\raggedright\ttfamily\scriptsize \emph{blank line}}}
\par\smallskip
This blank line is not hardware or software behavior. It separates nearby declarations or operations so the file remains readable and reviewable.
\paragraph{Line 53}
\noindent\colorbox{gray!8}{\parbox{\dimexpr\linewidth-2\fboxsep\relax}{\raggedright\ttfamily\scriptsize function Assert-LastExitCode \{}}
\par\smallskip
This calls the helper that turns a failed external tool exit code into a hard pipeline failure.
\paragraph{Line 54}
\noindent\colorbox{gray!8}{\parbox{\dimexpr\linewidth-2\fboxsep\relax}{\raggedright\ttfamily\scriptsize     param([string]\$Step)}}
\par\smallskip
This begins the PowerShell parameter block. It exposes repeatable knobs such as sample count, Vivado path, board part, and skip switches.
\paragraph{Line 55}
\noindent\colorbox{gray!8}{\parbox{\dimexpr\linewidth-2\fboxsep\relax}{\raggedright\ttfamily\scriptsize     if (\$LASTEXITCODE -ne 0) \{}}
\par\smallskip
This PowerShell line orchestrates the reproducible full flow. It exists so changes can be rerun without manually retyping many commands.
\paragraph{Line 56}
\noindent\colorbox{gray!8}{\parbox{\dimexpr\linewidth-2\fboxsep\relax}{\raggedright\ttfamily\scriptsize         throw "\$Step failed with exit code \$LASTEXITCODE"}}
\par\smallskip
This PowerShell line orchestrates the reproducible full flow. It exists so changes can be rerun without manually retyping many commands.
\paragraph{Line 57}
\noindent\colorbox{gray!8}{\parbox{\dimexpr\linewidth-2\fboxsep\relax}{\raggedright\ttfamily\scriptsize     \}}}
\par\smallskip
This PowerShell line orchestrates the reproducible full flow. It exists so changes can be rerun without manually retyping many commands.
\paragraph{Line 58}
\noindent\colorbox{gray!8}{\parbox{\dimexpr\linewidth-2\fboxsep\relax}{\raggedright\ttfamily\scriptsize \}}}
\par\smallskip
This PowerShell line orchestrates the reproducible full flow. It exists so changes can be rerun without manually retyping many commands.
\paragraph{Line 59}
\noindent\colorbox{gray!8}{\parbox{\dimexpr\linewidth-2\fboxsep\relax}{\raggedright\ttfamily\scriptsize \emph{blank line}}}
\par\smallskip
This blank line is not hardware or software behavior. It separates nearby declarations or operations so the file remains readable and reviewable.
\paragraph{Line 60}
\noindent\colorbox{gray!8}{\parbox{\dimexpr\linewidth-2\fboxsep\relax}{\raggedright\ttfamily\scriptsize Push-Location \$RepoRoot}}
\par\smallskip
This PowerShell line orchestrates the reproducible full flow. It exists so changes can be rerun without manually retyping many commands.
\paragraph{Line 61}
\noindent\colorbox{gray!8}{\parbox{\dimexpr\linewidth-2\fboxsep\relax}{\raggedright\ttfamily\scriptsize try \{}}
\par\smallskip
This PowerShell line orchestrates the reproducible full flow. It exists so changes can be rerun without manually retyping many commands.
\paragraph{Line 62}
\noindent\colorbox{gray!8}{\parbox{\dimexpr\linewidth-2\fboxsep\relax}{\raggedright\ttfamily\scriptsize     if (!\$SkipExport) \{}}
\par\smallskip
This PowerShell line orchestrates the reproducible full flow. It exists so changes can be rerun without manually retyping many commands.
\paragraph{Line 63}
\noindent\colorbox{gray!8}{\parbox{\dimexpr\linewidth-2\fboxsep\relax}{\raggedright\ttfamily\scriptsize         Write-Host "[export] trained Spikingformer checkpoint to fixed-point accelerator memories"}}
\par\smallskip
This PowerShell line orchestrates the reproducible full flow. It exists so changes can be rerun without manually retyping many commands.
\paragraph{Line 64}
\noindent\colorbox{gray!8}{\parbox{\dimexpr\linewidth-2\fboxsep\relax}{\raggedright\ttfamily\scriptsize         \& \$Python "cifar10\_neuromorphic\textbackslash{}export\_spikingformer\_head\_accelerator.py" `}}
\par\smallskip
This line invokes the Python environment used for export, validation, or report generation. It ties host-side math to the hardware artifacts.
\paragraph{Line 65}
\noindent\colorbox{gray!8}{\parbox{\dimexpr\linewidth-2\fboxsep\relax}{\raggedright\ttfamily\scriptsize             --checkpoint "\$Checkpoint" `}}
\par\smallskip
This PowerShell line orchestrates the reproducible full flow. It exists so changes can be rerun without manually retyping many commands.
\paragraph{Line 66}
\noindent\colorbox{gray!8}{\parbox{\dimexpr\linewidth-2\fboxsep\relax}{\raggedright\ttfamily\scriptsize             --samples "\$Samples" `}}
\par\smallskip
This PowerShell line orchestrates the reproducible full flow. It exists so changes can be rerun without manually retyping many commands.
\paragraph{Line 67}
\noindent\colorbox{gray!8}{\parbox{\dimexpr\linewidth-2\fboxsep\relax}{\raggedright\ttfamily\scriptsize             --batch-size "\$BatchSize" `}}
\par\smallskip
This PowerShell line orchestrates the reproducible full flow. It exists so changes can be rerun without manually retyping many commands.
\paragraph{Line 68}
\noindent\colorbox{gray!8}{\parbox{\dimexpr\linewidth-2\fboxsep\relax}{\raggedright\ttfamily\scriptsize             --workers "\$Workers" `}}
\par\smallskip
This PowerShell line orchestrates the reproducible full flow. It exists so changes can be rerun without manually retyping many commands.
\paragraph{Line 69}
\noindent\colorbox{gray!8}{\parbox{\dimexpr\linewidth-2\fboxsep\relax}{\raggedright\ttfamily\scriptsize             --trace-samples "\$TraceSamples" `}}
\par\smallskip
This PowerShell line orchestrates the reproducible full flow. It exists so changes can be rerun without manually retyping many commands.
\paragraph{Line 70}
\noindent\colorbox{gray!8}{\parbox{\dimexpr\linewidth-2\fboxsep\relax}{\raggedright\ttfamily\scriptsize             --infer-skip-abs-thresh "\$InferSkipAbsThresh" `}}
\par\smallskip
This PowerShell line orchestrates the reproducible full flow. It exists so changes can be rerun without manually retyping many commands.
\paragraph{Line 71}
\noindent\colorbox{gray!8}{\parbox{\dimexpr\linewidth-2\fboxsep\relax}{\raggedright\ttfamily\scriptsize             --out-dir "\$OutDir"}}
\par\smallskip
This PowerShell line orchestrates the reproducible full flow. It exists so changes can be rerun without manually retyping many commands.
\paragraph{Line 72}
\noindent\colorbox{gray!8}{\parbox{\dimexpr\linewidth-2\fboxsep\relax}{\raggedright\ttfamily\scriptsize         Assert-LastExitCode "Spikingformer head export"}}
\par\smallskip
This calls the helper that turns a failed external tool exit code into a hard pipeline failure.
\paragraph{Line 73}
\noindent\colorbox{gray!8}{\parbox{\dimexpr\linewidth-2\fboxsep\relax}{\raggedright\ttfamily\scriptsize     \}}}
\par\smallskip
This PowerShell line orchestrates the reproducible full flow. It exists so changes can be rerun without manually retyping many commands.
\paragraph{Line 74}
\noindent\colorbox{gray!8}{\parbox{\dimexpr\linewidth-2\fboxsep\relax}{\raggedright\ttfamily\scriptsize \} finally \{}}
\par\smallskip
This PowerShell line orchestrates the reproducible full flow. It exists so changes can be rerun without manually retyping many commands.
\paragraph{Line 75}
\noindent\colorbox{gray!8}{\parbox{\dimexpr\linewidth-2\fboxsep\relax}{\raggedright\ttfamily\scriptsize     Pop-Location}}
\par\smallskip
This PowerShell line orchestrates the reproducible full flow. It exists so changes can be rerun without manually retyping many commands.
\paragraph{Line 76}
\noindent\colorbox{gray!8}{\parbox{\dimexpr\linewidth-2\fboxsep\relax}{\raggedright\ttfamily\scriptsize \}}}
\par\smallskip
This PowerShell line orchestrates the reproducible full flow. It exists so changes can be rerun without manually retyping many commands.
\paragraph{Line 77}
\noindent\colorbox{gray!8}{\parbox{\dimexpr\linewidth-2\fboxsep\relax}{\raggedright\ttfamily\scriptsize \emph{blank line}}}
\par\smallskip
This blank line is not hardware or software behavior. It separates nearby declarations or operations so the file remains readable and reviewable.
\paragraph{Line 78}
\noindent\colorbox{gray!8}{\parbox{\dimexpr\linewidth-2\fboxsep\relax}{\raggedright\ttfamily\scriptsize if (!\$SkipVivado) \{}}
\par\smallskip
This line invokes or locates Vivado simulator/implementation tools. The flow uses the installed AMD toolchain on D: rather than a substitute simulator.
\paragraph{Line 79}
\noindent\colorbox{gray!8}{\parbox{\dimexpr\linewidth-2\fboxsep\relax}{\raggedright\ttfamily\scriptsize     if (!(Test-Path \$VivadoXvlog)) \{ throw "Vivado xvlog not found at \$VivadoXvlog" \}}}
\par\smallskip
This line invokes or locates Vivado simulator/implementation tools. The flow uses the installed AMD toolchain on D: rather than a substitute simulator.
\paragraph{Line 80}
\noindent\colorbox{gray!8}{\parbox{\dimexpr\linewidth-2\fboxsep\relax}{\raggedright\ttfamily\scriptsize     if (!(Test-Path \$VivadoXelab)) \{ throw "Vivado xelab not found at \$VivadoXelab" \}}}
\par\smallskip
This line invokes or locates Vivado simulator/implementation tools. The flow uses the installed AMD toolchain on D: rather than a substitute simulator.
\paragraph{Line 81}
\noindent\colorbox{gray!8}{\parbox{\dimexpr\linewidth-2\fboxsep\relax}{\raggedright\ttfamily\scriptsize     if (!(Test-Path \$VivadoXsim)) \{ throw "Vivado xsim not found at \$VivadoXsim" \}}}
\par\smallskip
This line invokes or locates Vivado simulator/implementation tools. The flow uses the installed AMD toolchain on D: rather than a substitute simulator.
\paragraph{Line 82}
\noindent\colorbox{gray!8}{\parbox{\dimexpr\linewidth-2\fboxsep\relax}{\raggedright\ttfamily\scriptsize \emph{blank line}}}
\par\smallskip
This blank line is not hardware or software behavior. It separates nearby declarations or operations so the file remains readable and reviewable.
\paragraph{Line 83}
\noindent\colorbox{gray!8}{\parbox{\dimexpr\linewidth-2\fboxsep\relax}{\raggedright\ttfamily\scriptsize     Push-Location \$ProjectDir}}
\par\smallskip
This PowerShell line orchestrates the reproducible full flow. It exists so changes can be rerun without manually retyping many commands.
\paragraph{Line 84}
\noindent\colorbox{gray!8}{\parbox{\dimexpr\linewidth-2\fboxsep\relax}{\raggedright\ttfamily\scriptsize     try \{}}
\par\smallskip
This PowerShell line orchestrates the reproducible full flow. It exists so changes can be rerun without manually retyping many commands.
\paragraph{Line 85}
\noindent\colorbox{gray!8}{\parbox{\dimexpr\linewidth-2\fboxsep\relax}{\raggedright\ttfamily\scriptsize         Write-Host "[vivado] compiling trained Spikingformer head accelerator"}}
\par\smallskip
This PowerShell line orchestrates the reproducible full flow. It exists so changes can be rerun without manually retyping many commands.
\paragraph{Line 86}
\noindent\colorbox{gray!8}{\parbox{\dimexpr\linewidth-2\fboxsep\relax}{\raggedright\ttfamily\scriptsize         \& \$VivadoXvlog -sv `}}
\par\smallskip
This line invokes or locates Vivado simulator/implementation tools. The flow uses the installed AMD toolchain on D: rather than a substitute simulator.
\paragraph{Line 87}
\noindent\colorbox{gray!8}{\parbox{\dimexpr\linewidth-2\fboxsep\relax}{\raggedright\ttfamily\scriptsize             rtl\textbackslash{}cifar10\_spikingformer\_head\_accelerator.sv `}}
\par\smallskip
This PowerShell line orchestrates the reproducible full flow. It exists so changes can be rerun without manually retyping many commands.
\paragraph{Line 88}
\noindent\colorbox{gray!8}{\parbox{\dimexpr\linewidth-2\fboxsep\relax}{\raggedright\ttfamily\scriptsize             rtl\textbackslash{}tb\_cifar10\_spikingformer\_head\_accelerator.sv}}
\par\smallskip
This PowerShell line orchestrates the reproducible full flow. It exists so changes can be rerun without manually retyping many commands.
\paragraph{Line 89}
\noindent\colorbox{gray!8}{\parbox{\dimexpr\linewidth-2\fboxsep\relax}{\raggedright\ttfamily\scriptsize         Assert-LastExitCode "Vivado xvlog"}}
\par\smallskip
This calls the helper that turns a failed external tool exit code into a hard pipeline failure.
\paragraph{Line 90}
\noindent\colorbox{gray!8}{\parbox{\dimexpr\linewidth-2\fboxsep\relax}{\raggedright\ttfamily\scriptsize         \& \$VivadoXelab tb\_cifar10\_spikingformer\_head\_accelerator -debug typical}}
\par\smallskip
This line invokes or locates Vivado simulator/implementation tools. The flow uses the installed AMD toolchain on D: rather than a substitute simulator.
\paragraph{Line 91}
\noindent\colorbox{gray!8}{\parbox{\dimexpr\linewidth-2\fboxsep\relax}{\raggedright\ttfamily\scriptsize         Assert-LastExitCode "Vivado xelab"}}
\par\smallskip
This calls the helper that turns a failed external tool exit code into a hard pipeline failure.
\paragraph{Line 92}
\noindent\colorbox{gray!8}{\parbox{\dimexpr\linewidth-2\fboxsep\relax}{\raggedright\ttfamily\scriptsize         \& \$VivadoXsim tb\_cifar10\_spikingformer\_head\_accelerator -runall}}
\par\smallskip
This line invokes or locates Vivado simulator/implementation tools. The flow uses the installed AMD toolchain on D: rather than a substitute simulator.
\paragraph{Line 93}
\noindent\colorbox{gray!8}{\parbox{\dimexpr\linewidth-2\fboxsep\relax}{\raggedright\ttfamily\scriptsize         Assert-LastExitCode "Vivado xsim"}}
\par\smallskip
This calls the helper that turns a failed external tool exit code into a hard pipeline failure.
\paragraph{Line 94}
\noindent\colorbox{gray!8}{\parbox{\dimexpr\linewidth-2\fboxsep\relax}{\raggedright\ttfamily\scriptsize     \} finally \{}}
\par\smallskip
This PowerShell line orchestrates the reproducible full flow. It exists so changes can be rerun without manually retyping many commands.
\paragraph{Line 95}
\noindent\colorbox{gray!8}{\parbox{\dimexpr\linewidth-2\fboxsep\relax}{\raggedright\ttfamily\scriptsize         Pop-Location}}
\par\smallskip
This PowerShell line orchestrates the reproducible full flow. It exists so changes can be rerun without manually retyping many commands.
\paragraph{Line 96}
\noindent\colorbox{gray!8}{\parbox{\dimexpr\linewidth-2\fboxsep\relax}{\raggedright\ttfamily\scriptsize     \}}}
\par\smallskip
This PowerShell line orchestrates the reproducible full flow. It exists so changes can be rerun without manually retyping many commands.
\paragraph{Line 97}
\noindent\colorbox{gray!8}{\parbox{\dimexpr\linewidth-2\fboxsep\relax}{\raggedright\ttfamily\scriptsize \emph{blank line}}}
\par\smallskip
This blank line is not hardware or software behavior. It separates nearby declarations or operations so the file remains readable and reviewable.
\paragraph{Line 98}
\noindent\colorbox{gray!8}{\parbox{\dimexpr\linewidth-2\fboxsep\relax}{\raggedright\ttfamily\scriptsize     Push-Location \$RepoRoot}}
\par\smallskip
This PowerShell line orchestrates the reproducible full flow. It exists so changes can be rerun without manually retyping many commands.
\paragraph{Line 99}
\noindent\colorbox{gray!8}{\parbox{\dimexpr\linewidth-2\fboxsep\relax}{\raggedright\ttfamily\scriptsize     try \{}}
\par\smallskip
This PowerShell line orchestrates the reproducible full flow. It exists so changes can be rerun without manually retyping many commands.
\paragraph{Line 100}
\noindent\colorbox{gray!8}{\parbox{\dimexpr\linewidth-2\fboxsep\relax}{\raggedright\ttfamily\scriptsize         Write-Host "[validate] checking RTL scores against fixed-point reference"}}
\par\smallskip
This PowerShell line orchestrates the reproducible full flow. It exists so changes can be rerun without manually retyping many commands.
\paragraph{Line 101}
\noindent\colorbox{gray!8}{\parbox{\dimexpr\linewidth-2\fboxsep\relax}{\raggedright\ttfamily\scriptsize         \& \$Python "cifar10\_neuromorphic\textbackslash{}validate\_spikingformer\_head\_rtl.py" --out-dir "\$OutDir"}}
\par\smallskip
This line invokes the Python environment used for export, validation, or report generation. It ties host-side math to the hardware artifacts.
\paragraph{Line 102}
\noindent\colorbox{gray!8}{\parbox{\dimexpr\linewidth-2\fboxsep\relax}{\raggedright\ttfamily\scriptsize         Assert-LastExitCode "Spikingformer head RTL validation"}}
\par\smallskip
This calls the helper that turns a failed external tool exit code into a hard pipeline failure.
\paragraph{Line 103}
\noindent\colorbox{gray!8}{\parbox{\dimexpr\linewidth-2\fboxsep\relax}{\raggedright\ttfamily\scriptsize     \} finally \{}}
\par\smallskip
This PowerShell line orchestrates the reproducible full flow. It exists so changes can be rerun without manually retyping many commands.
\paragraph{Line 104}
\noindent\colorbox{gray!8}{\parbox{\dimexpr\linewidth-2\fboxsep\relax}{\raggedright\ttfamily\scriptsize         Pop-Location}}
\par\smallskip
This PowerShell line orchestrates the reproducible full flow. It exists so changes can be rerun without manually retyping many commands.
\paragraph{Line 105}
\noindent\colorbox{gray!8}{\parbox{\dimexpr\linewidth-2\fboxsep\relax}{\raggedright\ttfamily\scriptsize     \}}}
\par\smallskip
This PowerShell line orchestrates the reproducible full flow. It exists so changes can be rerun without manually retyping many commands.
\paragraph{Line 106}
\noindent\colorbox{gray!8}{\parbox{\dimexpr\linewidth-2\fboxsep\relax}{\raggedright\ttfamily\scriptsize \emph{blank line}}}
\par\smallskip
This blank line is not hardware or software behavior. It separates nearby declarations or operations so the file remains readable and reviewable.
\paragraph{Line 107}
\noindent\colorbox{gray!8}{\parbox{\dimexpr\linewidth-2\fboxsep\relax}{\raggedright\ttfamily\scriptsize     if (!\$SkipOnChipLearning) \{}}
\par\smallskip
This PowerShell line orchestrates the reproducible full flow. It exists so changes can be rerun without manually retyping many commands.
\paragraph{Line 108}
\noindent\colorbox{gray!8}{\parbox{\dimexpr\linewidth-2\fboxsep\relax}{\raggedright\ttfamily\scriptsize         Push-Location \$ProjectDir}}
\par\smallskip
This PowerShell line orchestrates the reproducible full flow. It exists so changes can be rerun without manually retyping many commands.
\paragraph{Line 109}
\noindent\colorbox{gray!8}{\parbox{\dimexpr\linewidth-2\fboxsep\relax}{\raggedright\ttfamily\scriptsize         try \{}}
\par\smallskip
This PowerShell line orchestrates the reproducible full flow. It exists so changes can be rerun without manually retyping many commands.
\paragraph{Line 110}
\noindent\colorbox{gray!8}{\parbox{\dimexpr\linewidth-2\fboxsep\relax}{\raggedright\ttfamily\scriptsize             Write-Host "[vivado] verifying on-chip learning update datapath"}}
\par\smallskip
This PowerShell line orchestrates the reproducible full flow. It exists so changes can be rerun without manually retyping many commands.
\paragraph{Line 111}
\noindent\colorbox{gray!8}{\parbox{\dimexpr\linewidth-2\fboxsep\relax}{\raggedright\ttfamily\scriptsize             \& \$VivadoXvlog -sv `}}
\par\smallskip
This line invokes or locates Vivado simulator/implementation tools. The flow uses the installed AMD toolchain on D: rather than a substitute simulator.
\paragraph{Line 112}
\noindent\colorbox{gray!8}{\parbox{\dimexpr\linewidth-2\fboxsep\relax}{\raggedright\ttfamily\scriptsize                 rtl\textbackslash{}cifar10\_spikingformer\_head\_accelerator.sv `}}
\par\smallskip
This PowerShell line orchestrates the reproducible full flow. It exists so changes can be rerun without manually retyping many commands.
\paragraph{Line 113}
\noindent\colorbox{gray!8}{\parbox{\dimexpr\linewidth-2\fboxsep\relax}{\raggedright\ttfamily\scriptsize                 rtl\textbackslash{}tb\_spikingformer\_head\_learning\_unit.sv}}
\par\smallskip
This PowerShell line orchestrates the reproducible full flow. It exists so changes can be rerun without manually retyping many commands.
\paragraph{Line 114}
\noindent\colorbox{gray!8}{\parbox{\dimexpr\linewidth-2\fboxsep\relax}{\raggedright\ttfamily\scriptsize             Assert-LastExitCode "Vivado xvlog on-chip learning unit"}}
\par\smallskip
This calls the helper that turns a failed external tool exit code into a hard pipeline failure.
\paragraph{Line 115}
\noindent\colorbox{gray!8}{\parbox{\dimexpr\linewidth-2\fboxsep\relax}{\raggedright\ttfamily\scriptsize             \& \$VivadoXelab tb\_spikingformer\_head\_learning\_unit -debug typical}}
\par\smallskip
This line invokes or locates Vivado simulator/implementation tools. The flow uses the installed AMD toolchain on D: rather than a substitute simulator.
\paragraph{Line 116}
\noindent\colorbox{gray!8}{\parbox{\dimexpr\linewidth-2\fboxsep\relax}{\raggedright\ttfamily\scriptsize             Assert-LastExitCode "Vivado xelab on-chip learning unit"}}
\par\smallskip
This calls the helper that turns a failed external tool exit code into a hard pipeline failure.
\paragraph{Line 117}
\noindent\colorbox{gray!8}{\parbox{\dimexpr\linewidth-2\fboxsep\relax}{\raggedright\ttfamily\scriptsize             \& \$VivadoXsim tb\_spikingformer\_head\_learning\_unit -runall}}
\par\smallskip
This line invokes or locates Vivado simulator/implementation tools. The flow uses the installed AMD toolchain on D: rather than a substitute simulator.
\paragraph{Line 118}
\noindent\colorbox{gray!8}{\parbox{\dimexpr\linewidth-2\fboxsep\relax}{\raggedright\ttfamily\scriptsize             Assert-LastExitCode "Vivado xsim on-chip learning unit"}}
\par\smallskip
This calls the helper that turns a failed external tool exit code into a hard pipeline failure.
\paragraph{Line 119}
\noindent\colorbox{gray!8}{\parbox{\dimexpr\linewidth-2\fboxsep\relax}{\raggedright\ttfamily\scriptsize \emph{blank line}}}
\par\smallskip
This blank line is not hardware or software behavior. It separates nearby declarations or operations so the file remains readable and reviewable.
\paragraph{Line 120}
\noindent\colorbox{gray!8}{\parbox{\dimexpr\linewidth-2\fboxsep\relax}{\raggedright\ttfamily\scriptsize             Write-Host "[vivado] running full CIFAR-10 on-chip learning simulation"}}
\par\smallskip
This PowerShell line orchestrates the reproducible full flow. It exists so changes can be rerun without manually retyping many commands.
\paragraph{Line 121}
\noindent\colorbox{gray!8}{\parbox{\dimexpr\linewidth-2\fboxsep\relax}{\raggedright\ttfamily\scriptsize             \& \$VivadoXvlog -sv `}}
\par\smallskip
This line invokes or locates Vivado simulator/implementation tools. The flow uses the installed AMD toolchain on D: rather than a substitute simulator.
\paragraph{Line 122}
\noindent\colorbox{gray!8}{\parbox{\dimexpr\linewidth-2\fboxsep\relax}{\raggedright\ttfamily\scriptsize                 rtl\textbackslash{}cifar10\_spikingformer\_head\_accelerator.sv `}}
\par\smallskip
This PowerShell line orchestrates the reproducible full flow. It exists so changes can be rerun without manually retyping many commands.
\paragraph{Line 123}
\noindent\colorbox{gray!8}{\parbox{\dimexpr\linewidth-2\fboxsep\relax}{\raggedright\ttfamily\scriptsize                 rtl\textbackslash{}tb\_cifar10\_spikingformer\_head\_onchip\_learning.sv}}
\par\smallskip
This PowerShell line orchestrates the reproducible full flow. It exists so changes can be rerun without manually retyping many commands.
\paragraph{Line 124}
\noindent\colorbox{gray!8}{\parbox{\dimexpr\linewidth-2\fboxsep\relax}{\raggedright\ttfamily\scriptsize             Assert-LastExitCode "Vivado xvlog on-chip learning full"}}
\par\smallskip
This calls the helper that turns a failed external tool exit code into a hard pipeline failure.
\paragraph{Line 125}
\noindent\colorbox{gray!8}{\parbox{\dimexpr\linewidth-2\fboxsep\relax}{\raggedright\ttfamily\scriptsize             \& \$VivadoXelab tb\_cifar10\_spikingformer\_head\_onchip\_learning -debug typical}}
\par\smallskip
This line invokes or locates Vivado simulator/implementation tools. The flow uses the installed AMD toolchain on D: rather than a substitute simulator.
\paragraph{Line 126}
\noindent\colorbox{gray!8}{\parbox{\dimexpr\linewidth-2\fboxsep\relax}{\raggedright\ttfamily\scriptsize             Assert-LastExitCode "Vivado xelab on-chip learning full"}}
\par\smallskip
This calls the helper that turns a failed external tool exit code into a hard pipeline failure.
\paragraph{Line 127}
\noindent\colorbox{gray!8}{\parbox{\dimexpr\linewidth-2\fboxsep\relax}{\raggedright\ttfamily\scriptsize             \& \$VivadoXsim tb\_cifar10\_spikingformer\_head\_onchip\_learning -runall}}
\par\smallskip
This line invokes or locates Vivado simulator/implementation tools. The flow uses the installed AMD toolchain on D: rather than a substitute simulator.
\paragraph{Line 128}
\noindent\colorbox{gray!8}{\parbox{\dimexpr\linewidth-2\fboxsep\relax}{\raggedright\ttfamily\scriptsize             Assert-LastExitCode "Vivado xsim on-chip learning full"}}
\par\smallskip
This calls the helper that turns a failed external tool exit code into a hard pipeline failure.
\paragraph{Line 129}
\noindent\colorbox{gray!8}{\parbox{\dimexpr\linewidth-2\fboxsep\relax}{\raggedright\ttfamily\scriptsize         \} finally \{}}
\par\smallskip
This PowerShell line orchestrates the reproducible full flow. It exists so changes can be rerun without manually retyping many commands.
\paragraph{Line 130}
\noindent\colorbox{gray!8}{\parbox{\dimexpr\linewidth-2\fboxsep\relax}{\raggedright\ttfamily\scriptsize             Pop-Location}}
\par\smallskip
This PowerShell line orchestrates the reproducible full flow. It exists so changes can be rerun without manually retyping many commands.
\paragraph{Line 131}
\noindent\colorbox{gray!8}{\parbox{\dimexpr\linewidth-2\fboxsep\relax}{\raggedright\ttfamily\scriptsize         \}}}
\par\smallskip
This PowerShell line orchestrates the reproducible full flow. It exists so changes can be rerun without manually retyping many commands.
\paragraph{Line 132}
\noindent\colorbox{gray!8}{\parbox{\dimexpr\linewidth-2\fboxsep\relax}{\raggedright\ttfamily\scriptsize \emph{blank line}}}
\par\smallskip
This blank line is not hardware or software behavior. It separates nearby declarations or operations so the file remains readable and reviewable.
\paragraph{Line 133}
\noindent\colorbox{gray!8}{\parbox{\dimexpr\linewidth-2\fboxsep\relax}{\raggedright\ttfamily\scriptsize         Push-Location \$RepoRoot}}
\par\smallskip
This PowerShell line orchestrates the reproducible full flow. It exists so changes can be rerun without manually retyping many commands.
\paragraph{Line 134}
\noindent\colorbox{gray!8}{\parbox{\dimexpr\linewidth-2\fboxsep\relax}{\raggedright\ttfamily\scriptsize         try \{}}
\par\smallskip
This PowerShell line orchestrates the reproducible full flow. It exists so changes can be rerun without manually retyping many commands.
\paragraph{Line 135}
\noindent\colorbox{gray!8}{\parbox{\dimexpr\linewidth-2\fboxsep\relax}{\raggedright\ttfamily\scriptsize             Write-Host "[validate] checking RTL on-chip learning against golden online update model"}}
\par\smallskip
This PowerShell line orchestrates the reproducible full flow. It exists so changes can be rerun without manually retyping many commands.
\paragraph{Line 136}
\noindent\colorbox{gray!8}{\parbox{\dimexpr\linewidth-2\fboxsep\relax}{\raggedright\ttfamily\scriptsize             \& \$Python "cifar10\_neuromorphic\textbackslash{}validate\_spikingformer\_head\_onchip\_learning.py" --out-dir "\$OutDir"}}
\par\smallskip
This line invokes the Python environment used for export, validation, or report generation. It ties host-side math to the hardware artifacts.
\paragraph{Line 137}
\noindent\colorbox{gray!8}{\parbox{\dimexpr\linewidth-2\fboxsep\relax}{\raggedright\ttfamily\scriptsize             Assert-LastExitCode "Spikingformer head on-chip learning validation"}}
\par\smallskip
This calls the helper that turns a failed external tool exit code into a hard pipeline failure.
\paragraph{Line 138}
\noindent\colorbox{gray!8}{\parbox{\dimexpr\linewidth-2\fboxsep\relax}{\raggedright\ttfamily\scriptsize         \} finally \{}}
\par\smallskip
This PowerShell line orchestrates the reproducible full flow. It exists so changes can be rerun without manually retyping many commands.
\paragraph{Line 139}
\noindent\colorbox{gray!8}{\parbox{\dimexpr\linewidth-2\fboxsep\relax}{\raggedright\ttfamily\scriptsize             Pop-Location}}
\par\smallskip
This PowerShell line orchestrates the reproducible full flow. It exists so changes can be rerun without manually retyping many commands.
\paragraph{Line 140}
\noindent\colorbox{gray!8}{\parbox{\dimexpr\linewidth-2\fboxsep\relax}{\raggedright\ttfamily\scriptsize         \}}}
\par\smallskip
This PowerShell line orchestrates the reproducible full flow. It exists so changes can be rerun without manually retyping many commands.
\paragraph{Line 141}
\noindent\colorbox{gray!8}{\parbox{\dimexpr\linewidth-2\fboxsep\relax}{\raggedright\ttfamily\scriptsize     \}}}
\par\smallskip
This PowerShell line orchestrates the reproducible full flow. It exists so changes can be rerun without manually retyping many commands.
\paragraph{Line 142}
\noindent\colorbox{gray!8}{\parbox{\dimexpr\linewidth-2\fboxsep\relax}{\raggedright\ttfamily\scriptsize \}}}
\par\smallskip
This PowerShell line orchestrates the reproducible full flow. It exists so changes can be rerun without manually retyping many commands.
\paragraph{Line 143}
\noindent\colorbox{gray!8}{\parbox{\dimexpr\linewidth-2\fboxsep\relax}{\raggedright\ttfamily\scriptsize \emph{blank line}}}
\par\smallskip
This blank line is not hardware or software behavior. It separates nearby declarations or operations so the file remains readable and reviewable.
\paragraph{Line 144}
\noindent\colorbox{gray!8}{\parbox{\dimexpr\linewidth-2\fboxsep\relax}{\raggedright\ttfamily\scriptsize if (!\$SkipSynthesis) \{}}
\par\smallskip
This PowerShell line orchestrates the reproducible full flow. It exists so changes can be rerun without manually retyping many commands.
\paragraph{Line 145}
\noindent\colorbox{gray!8}{\parbox{\dimexpr\linewidth-2\fboxsep\relax}{\raggedright\ttfamily\scriptsize     if (!(Test-Path \$Vivado)) \{ throw "Vivado not found at \$Vivado" \}}}
\par\smallskip
This line invokes or locates Vivado simulator/implementation tools. The flow uses the installed AMD toolchain on D: rather than a substitute simulator.
\paragraph{Line 146}
\noindent\colorbox{gray!8}{\parbox{\dimexpr\linewidth-2\fboxsep\relax}{\raggedright\ttfamily\scriptsize     \$SynthOut = Join-Path (Resolve-Path \$OutDir) "vivado"}}
\par\smallskip
This PowerShell line orchestrates the reproducible full flow. It exists so changes can be rerun without manually retyping many commands.
\paragraph{Line 147}
\noindent\colorbox{gray!8}{\parbox{\dimexpr\linewidth-2\fboxsep\relax}{\raggedright\ttfamily\scriptsize     New-Item -ItemType Directory -Force -Path \$SynthOut | Out-Null}}
\par\smallskip
This PowerShell line orchestrates the reproducible full flow. It exists so changes can be rerun without manually retyping many commands.
\paragraph{Line 148}
\noindent\colorbox{gray!8}{\parbox{\dimexpr\linewidth-2\fboxsep\relax}{\raggedright\ttfamily\scriptsize     \$env:NEURAEDGE\_PART = \$Part}}
\par\smallskip
This PowerShell line orchestrates the reproducible full flow. It exists so changes can be rerun without manually retyping many commands.
\paragraph{Line 149}
\noindent\colorbox{gray!8}{\parbox{\dimexpr\linewidth-2\fboxsep\relax}{\raggedright\ttfamily\scriptsize     \$env:NEURAEDGE\_CLOCK\_NS = "\$ClockPeriodNs"}}
\par\smallskip
This PowerShell line orchestrates the reproducible full flow. It exists so changes can be rerun without manually retyping many commands.
\paragraph{Line 150}
\noindent\colorbox{gray!8}{\parbox{\dimexpr\linewidth-2\fboxsep\relax}{\raggedright\ttfamily\scriptsize     \$env:NEURAEDGE\_OUT\_DIR = "\$SynthOut"}}
\par\smallskip
This PowerShell line orchestrates the reproducible full flow. It exists so changes can be rerun without manually retyping many commands.
\paragraph{Line 151}
\noindent\colorbox{gray!8}{\parbox{\dimexpr\linewidth-2\fboxsep\relax}{\raggedright\ttfamily\scriptsize     Push-Location \$ProjectDir}}
\par\smallskip
This PowerShell line orchestrates the reproducible full flow. It exists so changes can be rerun without manually retyping many commands.
\paragraph{Line 152}
\noindent\colorbox{gray!8}{\parbox{\dimexpr\linewidth-2\fboxsep\relax}{\raggedright\ttfamily\scriptsize     try \{}}
\par\smallskip
This PowerShell line orchestrates the reproducible full flow. It exists so changes can be rerun without manually retyping many commands.
\paragraph{Line 153}
\noindent\colorbox{gray!8}{\parbox{\dimexpr\linewidth-2\fboxsep\relax}{\raggedright\ttfamily\scriptsize         Write-Host "[synth] routing trained Spikingformer head accelerator for \$Part at \$\{ClockPeriodNs\}ns"}}
\par\smallskip
This PowerShell line orchestrates the reproducible full flow. It exists so changes can be rerun without manually retyping many commands.
\paragraph{Line 154}
\noindent\colorbox{gray!8}{\parbox{\dimexpr\linewidth-2\fboxsep\relax}{\raggedright\ttfamily\scriptsize         \& \$Vivado -mode batch -source "spikingformer\_head\_synth.tcl"}}
\par\smallskip
This line invokes or locates Vivado simulator/implementation tools. The flow uses the installed AMD toolchain on D: rather than a substitute simulator.
\paragraph{Line 155}
\noindent\colorbox{gray!8}{\parbox{\dimexpr\linewidth-2\fboxsep\relax}{\raggedright\ttfamily\scriptsize         Assert-LastExitCode "Spikingformer head synthesis/implementation"}}
\par\smallskip
This calls the helper that turns a failed external tool exit code into a hard pipeline failure.
\paragraph{Line 156}
\noindent\colorbox{gray!8}{\parbox{\dimexpr\linewidth-2\fboxsep\relax}{\raggedright\ttfamily\scriptsize     \} finally \{}}
\par\smallskip
This PowerShell line orchestrates the reproducible full flow. It exists so changes can be rerun without manually retyping many commands.
\paragraph{Line 157}
\noindent\colorbox{gray!8}{\parbox{\dimexpr\linewidth-2\fboxsep\relax}{\raggedright\ttfamily\scriptsize         Pop-Location}}
\par\smallskip
This PowerShell line orchestrates the reproducible full flow. It exists so changes can be rerun without manually retyping many commands.
\paragraph{Line 158}
\noindent\colorbox{gray!8}{\parbox{\dimexpr\linewidth-2\fboxsep\relax}{\raggedright\ttfamily\scriptsize         Remove-Item Env:\textbackslash{}NEURAEDGE\_PART -ErrorAction SilentlyContinue}}
\par\smallskip
This PowerShell line orchestrates the reproducible full flow. It exists so changes can be rerun without manually retyping many commands.
\paragraph{Line 159}
\noindent\colorbox{gray!8}{\parbox{\dimexpr\linewidth-2\fboxsep\relax}{\raggedright\ttfamily\scriptsize         Remove-Item Env:\textbackslash{}NEURAEDGE\_CLOCK\_NS -ErrorAction SilentlyContinue}}
\par\smallskip
This PowerShell line orchestrates the reproducible full flow. It exists so changes can be rerun without manually retyping many commands.
\paragraph{Line 160}
\noindent\colorbox{gray!8}{\parbox{\dimexpr\linewidth-2\fboxsep\relax}{\raggedright\ttfamily\scriptsize         Remove-Item Env:\textbackslash{}NEURAEDGE\_OUT\_DIR -ErrorAction SilentlyContinue}}
\par\smallskip
This PowerShell line orchestrates the reproducible full flow. It exists so changes can be rerun without manually retyping many commands.
\paragraph{Line 161}
\noindent\colorbox{gray!8}{\parbox{\dimexpr\linewidth-2\fboxsep\relax}{\raggedright\ttfamily\scriptsize     \}}}
\par\smallskip
This PowerShell line orchestrates the reproducible full flow. It exists so changes can be rerun without manually retyping many commands.
\paragraph{Line 162}
\noindent\colorbox{gray!8}{\parbox{\dimexpr\linewidth-2\fboxsep\relax}{\raggedright\ttfamily\scriptsize \}}}
\par\smallskip
This PowerShell line orchestrates the reproducible full flow. It exists so changes can be rerun without manually retyping many commands.
\paragraph{Line 163}
\noindent\colorbox{gray!8}{\parbox{\dimexpr\linewidth-2\fboxsep\relax}{\raggedright\ttfamily\scriptsize \emph{blank line}}}
\par\smallskip
This blank line is not hardware or software behavior. It separates nearby declarations or operations so the file remains readable and reviewable.
\paragraph{Line 164}
\noindent\colorbox{gray!8}{\parbox{\dimexpr\linewidth-2\fboxsep\relax}{\raggedright\ttfamily\scriptsize Write-Host "[done] trained Spikingformer head accelerator flow completed"}}
\par\smallskip
This PowerShell line orchestrates the reproducible full flow. It exists so changes can be rerun without manually retyping many commands.
\paragraph{Line 165}
\noindent\colorbox{gray!8}{\parbox{\dimexpr\linewidth-2\fboxsep\relax}{\raggedright\ttfamily\scriptsize Write-Host "       outputs: \$OutDir"}}
\par\smallskip
This PowerShell line orchestrates the reproducible full flow. It exists so changes can be rerun without manually retyping many commands.
\paragraph{Line 166}
\noindent\colorbox{gray!8}{\parbox{\dimexpr\linewidth-2\fboxsep\relax}{\raggedright\ttfamily\scriptsize Write-Host "       rtl csv: cifar10\_neuromorphic\textbackslash{}cifar10\_spikingformer\_head\_results.csv"}}
\par\smallskip
This PowerShell line orchestrates the reproducible full flow. It exists so changes can be rerun without manually retyping many commands.
\clearpage
\subsection{cifar10\_neuromorphic\textbackslash{}spikingformer\_head\_synth.tcl}
Each entry below gives the exact source line followed by its hardware/software role in the implemented accelerator pipeline.
\paragraph{Line 1}
\noindent\colorbox{gray!8}{\parbox{\dimexpr\linewidth-2\fboxsep\relax}{\raggedright\ttfamily\scriptsize set project\_dir [file normalize [file dirname [info script]]]}}
\par\smallskip
This TCL line controls the Vivado batch implementation flow. It makes the hardware report reproducible from the command line.
\paragraph{Line 2}
\noindent\colorbox{gray!8}{\parbox{\dimexpr\linewidth-2\fboxsep\relax}{\raggedright\ttfamily\scriptsize set repo\_root [file normalize [file join \$project\_dir ".."]]}}
\par\smallskip
This TCL line controls the Vivado batch implementation flow. It makes the hardware report reproducible from the command line.
\paragraph{Line 3}
\noindent\colorbox{gray!8}{\parbox{\dimexpr\linewidth-2\fboxsep\relax}{\raggedright\ttfamily\scriptsize \emph{blank line}}}
\par\smallskip
This blank line is not hardware or software behavior. It separates nearby declarations or operations so the file remains readable and reviewable.
\paragraph{Line 4}
\noindent\colorbox{gray!8}{\parbox{\dimexpr\linewidth-2\fboxsep\relax}{\raggedright\ttfamily\scriptsize if \{![info exists ::env(NEURAEDGE\_PART)]\} \{}}
\par\smallskip
This TCL line controls the Vivado batch implementation flow. It makes the hardware report reproducible from the command line.
\paragraph{Line 5}
\noindent\colorbox{gray!8}{\parbox{\dimexpr\linewidth-2\fboxsep\relax}{\raggedright\ttfamily\scriptsize     set part\_name "xc7z020clg400-1"}}
\par\smallskip
This TCL line controls the Vivado batch implementation flow. It makes the hardware report reproducible from the command line.
\paragraph{Line 6}
\noindent\colorbox{gray!8}{\parbox{\dimexpr\linewidth-2\fboxsep\relax}{\raggedright\ttfamily\scriptsize \} else \{}}
\par\smallskip
This TCL line controls the Vivado batch implementation flow. It makes the hardware report reproducible from the command line.
\paragraph{Line 7}
\noindent\colorbox{gray!8}{\parbox{\dimexpr\linewidth-2\fboxsep\relax}{\raggedright\ttfamily\scriptsize     set part\_name \$::env(NEURAEDGE\_PART)}}
\par\smallskip
This TCL line controls the Vivado batch implementation flow. It makes the hardware report reproducible from the command line.
\paragraph{Line 8}
\noindent\colorbox{gray!8}{\parbox{\dimexpr\linewidth-2\fboxsep\relax}{\raggedright\ttfamily\scriptsize \}}}
\par\smallskip
This TCL line controls the Vivado batch implementation flow. It makes the hardware report reproducible from the command line.
\paragraph{Line 9}
\noindent\colorbox{gray!8}{\parbox{\dimexpr\linewidth-2\fboxsep\relax}{\raggedright\ttfamily\scriptsize \emph{blank line}}}
\par\smallskip
This blank line is not hardware or software behavior. It separates nearby declarations or operations so the file remains readable and reviewable.
\paragraph{Line 10}
\noindent\colorbox{gray!8}{\parbox{\dimexpr\linewidth-2\fboxsep\relax}{\raggedright\ttfamily\scriptsize if \{![info exists ::env(NEURAEDGE\_CLOCK\_NS)]\} \{}}
\par\smallskip
This TCL line controls the Vivado batch implementation flow. It makes the hardware report reproducible from the command line.
\paragraph{Line 11}
\noindent\colorbox{gray!8}{\parbox{\dimexpr\linewidth-2\fboxsep\relax}{\raggedright\ttfamily\scriptsize     set clock\_ns 10.0}}
\par\smallskip
This TCL line controls the Vivado batch implementation flow. It makes the hardware report reproducible from the command line.
\paragraph{Line 12}
\noindent\colorbox{gray!8}{\parbox{\dimexpr\linewidth-2\fboxsep\relax}{\raggedright\ttfamily\scriptsize \} else \{}}
\par\smallskip
This TCL line controls the Vivado batch implementation flow. It makes the hardware report reproducible from the command line.
\paragraph{Line 13}
\noindent\colorbox{gray!8}{\parbox{\dimexpr\linewidth-2\fboxsep\relax}{\raggedright\ttfamily\scriptsize     set clock\_ns \$::env(NEURAEDGE\_CLOCK\_NS)}}
\par\smallskip
This TCL line controls the Vivado batch implementation flow. It makes the hardware report reproducible from the command line.
\paragraph{Line 14}
\noindent\colorbox{gray!8}{\parbox{\dimexpr\linewidth-2\fboxsep\relax}{\raggedright\ttfamily\scriptsize \}}}
\par\smallskip
This TCL line controls the Vivado batch implementation flow. It makes the hardware report reproducible from the command line.
\paragraph{Line 15}
\noindent\colorbox{gray!8}{\parbox{\dimexpr\linewidth-2\fboxsep\relax}{\raggedright\ttfamily\scriptsize \emph{blank line}}}
\par\smallskip
This blank line is not hardware or software behavior. It separates nearby declarations or operations so the file remains readable and reviewable.
\paragraph{Line 16}
\noindent\colorbox{gray!8}{\parbox{\dimexpr\linewidth-2\fboxsep\relax}{\raggedright\ttfamily\scriptsize if \{![info exists ::env(NEURAEDGE\_OUT\_DIR)]\} \{}}
\par\smallskip
This TCL line controls the Vivado batch implementation flow. It makes the hardware report reproducible from the command line.
\paragraph{Line 17}
\noindent\colorbox{gray!8}{\parbox{\dimexpr\linewidth-2\fboxsep\relax}{\raggedright\ttfamily\scriptsize     set out\_dir [file join \$project\_dir "outputs\_spikingformer\_head\_accelerator" "vivado"]}}
\par\smallskip
This TCL line controls the Vivado batch implementation flow. It makes the hardware report reproducible from the command line.
\paragraph{Line 18}
\noindent\colorbox{gray!8}{\parbox{\dimexpr\linewidth-2\fboxsep\relax}{\raggedright\ttfamily\scriptsize \} else \{}}
\par\smallskip
This TCL line controls the Vivado batch implementation flow. It makes the hardware report reproducible from the command line.
\paragraph{Line 19}
\noindent\colorbox{gray!8}{\parbox{\dimexpr\linewidth-2\fboxsep\relax}{\raggedright\ttfamily\scriptsize     set out\_dir [file normalize \$::env(NEURAEDGE\_OUT\_DIR)]}}
\par\smallskip
This TCL line controls the Vivado batch implementation flow. It makes the hardware report reproducible from the command line.
\paragraph{Line 20}
\noindent\colorbox{gray!8}{\parbox{\dimexpr\linewidth-2\fboxsep\relax}{\raggedright\ttfamily\scriptsize \}}}
\par\smallskip
This TCL line controls the Vivado batch implementation flow. It makes the hardware report reproducible from the command line.
\paragraph{Line 21}
\noindent\colorbox{gray!8}{\parbox{\dimexpr\linewidth-2\fboxsep\relax}{\raggedright\ttfamily\scriptsize \emph{blank line}}}
\par\smallskip
This blank line is not hardware or software behavior. It separates nearby declarations or operations so the file remains readable and reviewable.
\paragraph{Line 22}
\noindent\colorbox{gray!8}{\parbox{\dimexpr\linewidth-2\fboxsep\relax}{\raggedright\ttfamily\scriptsize file mkdir \$out\_dir}}
\par\smallskip
This TCL line controls the Vivado batch implementation flow. It makes the hardware report reproducible from the command line.
\paragraph{Line 23}
\noindent\colorbox{gray!8}{\parbox{\dimexpr\linewidth-2\fboxsep\relax}{\raggedright\ttfamily\scriptsize cd \$project\_dir}}
\par\smallskip
This TCL line controls the Vivado batch implementation flow. It makes the hardware report reproducible from the command line.
\paragraph{Line 24}
\noindent\colorbox{gray!8}{\parbox{\dimexpr\linewidth-2\fboxsep\relax}{\raggedright\ttfamily\scriptsize \emph{blank line}}}
\par\smallskip
This blank line is not hardware or software behavior. It separates nearby declarations or operations so the file remains readable and reviewable.
\paragraph{Line 25}
\noindent\colorbox{gray!8}{\parbox{\dimexpr\linewidth-2\fboxsep\relax}{\raggedright\ttfamily\scriptsize read\_verilog -sv [file join \$project\_dir "rtl" "cifar10\_spikingformer\_head\_accelerator.sv"]}}
\par\smallskip
This tells Vivado which SystemVerilog files to synthesize. The wrapper and accelerator are both read so the routed report matches the PYNQ-facing design.
\paragraph{Line 26}
\noindent\colorbox{gray!8}{\parbox{\dimexpr\linewidth-2\fboxsep\relax}{\raggedright\ttfamily\scriptsize read\_verilog -sv [file join \$project\_dir "rtl" "cifar10\_spikingformer\_head\_pynq\_axi\_stream\_top.sv"]}}
\par\smallskip
This tells Vivado which SystemVerilog files to synthesize. The wrapper and accelerator are both read so the routed report matches the PYNQ-facing design.
\paragraph{Line 27}
\noindent\colorbox{gray!8}{\parbox{\dimexpr\linewidth-2\fboxsep\relax}{\raggedright\ttfamily\scriptsize synth\_design -top cifar10\_spikingformer\_head\_pynq\_axi\_stream\_top -part \$part\_name}}
\par\smallskip
This starts synthesis for the selected FPGA part and top module. Synthesis maps RTL arrays, arithmetic, and FSMs into FPGA primitives.
\paragraph{Line 28}
\noindent\colorbox{gray!8}{\parbox{\dimexpr\linewidth-2\fboxsep\relax}{\raggedright\ttfamily\scriptsize create\_clock -period \$clock\_ns -name ap\_clk [get\_ports ap\_clk]}}
\par\smallskip
This creates the 100 MHz timing constraint. Vivado uses it to decide whether every register-to-register path is fast enough.
\paragraph{Line 29}
\noindent\colorbox{gray!8}{\parbox{\dimexpr\linewidth-2\fboxsep\relax}{\raggedright\ttfamily\scriptsize \emph{blank line}}}
\par\smallskip
This blank line is not hardware or software behavior. It separates nearby declarations or operations so the file remains readable and reviewable.
\paragraph{Line 30}
\noindent\colorbox{gray!8}{\parbox{\dimexpr\linewidth-2\fboxsep\relax}{\raggedright\ttfamily\scriptsize opt\_design}}
\par\smallskip
This runs an implementation stage. These commands turn the synthesized netlist into a placed and routed design with real timing and power estimates.
\paragraph{Line 31}
\noindent\colorbox{gray!8}{\parbox{\dimexpr\linewidth-2\fboxsep\relax}{\raggedright\ttfamily\scriptsize place\_design}}
\par\smallskip
This runs an implementation stage. These commands turn the synthesized netlist into a placed and routed design with real timing and power estimates.
\paragraph{Line 32}
\noindent\colorbox{gray!8}{\parbox{\dimexpr\linewidth-2\fboxsep\relax}{\raggedright\ttfamily\scriptsize route\_design}}
\par\smallskip
This runs an implementation stage. These commands turn the synthesized netlist into a placed and routed design with real timing and power estimates.
\paragraph{Line 33}
\noindent\colorbox{gray!8}{\parbox{\dimexpr\linewidth-2\fboxsep\relax}{\raggedright\ttfamily\scriptsize \emph{blank line}}}
\par\smallskip
This blank line is not hardware or software behavior. It separates nearby declarations or operations so the file remains readable and reviewable.
\paragraph{Line 34}
\noindent\colorbox{gray!8}{\parbox{\dimexpr\linewidth-2\fboxsep\relax}{\raggedright\ttfamily\scriptsize report\_utilization -file [file join \$out\_dir "spikingformer\_head\_utilization.rpt"]}}
\par\smallskip
This emits a signoff artifact. The generated reports are the source for utilization, timing, and power numbers in the debug PDF.
\paragraph{Line 35}
\noindent\colorbox{gray!8}{\parbox{\dimexpr\linewidth-2\fboxsep\relax}{\raggedright\ttfamily\scriptsize report\_timing\_summary -file [file join \$out\_dir "spikingformer\_head\_timing\_summary.rpt"]}}
\par\smallskip
This emits a signoff artifact. The generated reports are the source for utilization, timing, and power numbers in the debug PDF.
\paragraph{Line 36}
\noindent\colorbox{gray!8}{\parbox{\dimexpr\linewidth-2\fboxsep\relax}{\raggedright\ttfamily\scriptsize report\_power -file [file join \$out\_dir "spikingformer\_head\_power.rpt"]}}
\par\smallskip
This emits a signoff artifact. The generated reports are the source for utilization, timing, and power numbers in the debug PDF.
\paragraph{Line 37}
\noindent\colorbox{gray!8}{\parbox{\dimexpr\linewidth-2\fboxsep\relax}{\raggedright\ttfamily\scriptsize write\_checkpoint -force [file join \$out\_dir "spikingformer\_head\_routed.dcp"]}}
\par\smallskip
This emits a signoff artifact. The generated reports are the source for utilization, timing, and power numbers in the debug PDF.
\paragraph{Line 38}
\noindent\colorbox{gray!8}{\parbox{\dimexpr\linewidth-2\fboxsep\relax}{\raggedright\ttfamily\scriptsize \emph{blank line}}}
\par\smallskip
This blank line is not hardware or software behavior. It separates nearby declarations or operations so the file remains readable and reviewable.
\paragraph{Line 39}
\noindent\colorbox{gray!8}{\parbox{\dimexpr\linewidth-2\fboxsep\relax}{\raggedright\ttfamily\scriptsize set timing\_ok [expr \{[get\_property SLACK [get\_timing\_paths -max\_paths 1 -setup]] >= 0.0\}]}}
\par\smallskip
This TCL line controls the Vivado batch implementation flow. It makes the hardware report reproducible from the command line.
\paragraph{Line 40}
\noindent\colorbox{gray!8}{\parbox{\dimexpr\linewidth-2\fboxsep\relax}{\raggedright\ttfamily\scriptsize set fh [open [file join \$out\_dir "spikingformer\_head\_synth\_summary.txt"] "w"]}}
\par\smallskip
This TCL line controls the Vivado batch implementation flow. It makes the hardware report reproducible from the command line.
\paragraph{Line 41}
\noindent\colorbox{gray!8}{\parbox{\dimexpr\linewidth-2\fboxsep\relax}{\raggedright\ttfamily\scriptsize puts \$fh "top=cifar10\_spikingformer\_head\_pynq\_axi\_stream\_top"}}
\par\smallskip
This TCL line controls the Vivado batch implementation flow. It makes the hardware report reproducible from the command line.
\paragraph{Line 42}
\noindent\colorbox{gray!8}{\parbox{\dimexpr\linewidth-2\fboxsep\relax}{\raggedright\ttfamily\scriptsize puts \$fh "part=\$part\_name"}}
\par\smallskip
This TCL line controls the Vivado batch implementation flow. It makes the hardware report reproducible from the command line.
\paragraph{Line 43}
\noindent\colorbox{gray!8}{\parbox{\dimexpr\linewidth-2\fboxsep\relax}{\raggedright\ttfamily\scriptsize puts \$fh "clock\_ns=\$clock\_ns"}}
\par\smallskip
This TCL line controls the Vivado batch implementation flow. It makes the hardware report reproducible from the command line.
\paragraph{Line 44}
\noindent\colorbox{gray!8}{\parbox{\dimexpr\linewidth-2\fboxsep\relax}{\raggedright\ttfamily\scriptsize puts \$fh "timing\_met=\$timing\_ok"}}
\par\smallskip
This TCL line controls the Vivado batch implementation flow. It makes the hardware report reproducible from the command line.
\paragraph{Line 45}
\noindent\colorbox{gray!8}{\parbox{\dimexpr\linewidth-2\fboxsep\relax}{\raggedright\ttfamily\scriptsize close \$fh}}
\par\smallskip
This TCL line controls the Vivado batch implementation flow. It makes the hardware report reproducible from the command line.
\end{document}